% !TeX program = pdflatex
\documentclass[11pt]{preprint}
\usepackage[T1]{fontenc}
\usepackage[utf8]{inputenc}
\usepackage{amsmath,amssymb,amsthm}
\usepackage{mathtools}
\usepackage{mathrsfs}
\usepackage[a4paper,margin=28mm]{geometry}
\usepackage{enumitem}
\usepackage{graphicx}
\usepackage[hidelinks]{hyperref}
\newcommand{\dd}{\mathrm{d}}
\newcommand{\Hol}{\operatorname{Hol}}
\newcommand{\tr}{\operatorname{tr}}
\newcommand{\id}{\mathbf 1}
\newcommand{\diam}{\operatorname{diam}}
\newcommand{\dist}{\operatorname{dist}}
\DeclarePairedDelimiter{\abs}{\lvert}{\rvert}
\DeclarePairedDelimiter{\norm}{\lVert}{\rVert}
\DeclarePairedDelimiter{\floor}{\lfloor}{\rfloor}
\DeclarePairedDelimiter{\ceil}{\lceil}{\rceil}
\makeatletter % a fat \cdot which is not as fat as a \bullet
\newcommand*{\bigcdot}{}% Check if undefined
\DeclareRobustCommand*{\bigcdot}{%
  \mathbin{\mathpalette\bigcdot@{}}%
}
\newcommand*{\bigcdot@scalefactor}{.7}
\newcommand*{\bigcdot@widthfactor}{1.15}
\newcommand*{\bigcdot@}[2]{%
  % #1: math style
  % #2: unused
  \sbox0{$#1\vcenter{}$}% math axis
  \sbox2{$#1\cdot\m@th$}%
  \hbox to \bigcdot@widthfactor\wd2{%
    \hss
    \raise\ht0\hbox{%
      \scalebox{\bigcdot@scalefactor}{%
        \lower\ht0\hbox{$#1\bullet\m@th$}%
      }%
    }%
    \hss
  }%
}
\makeatother
\setlength{\parskip}{0.6em}
\emergencystretch=3em
\title{The Continuum Limit of the Correlation Functions of Wilson Loops in Four-Dimensional Lattice Yang--Mills Theory}
\author{}
\date{}

\numberwithin{equation}{section}
\newtheorem{theorem}{Theorem}[section]
\newtheorem{proposition}[theorem]{Proposition}
\newtheorem{lemma}[theorem]{Lemma}
\newtheorem{corollary}[theorem]{Corollary}
\newtheorem{conjecture}[theorem]{Conjecture}
\theoremstyle{definition}
\newtheorem{definition}[theorem]{Definition}
\theoremstyle{remark}
\newtheorem{remark}[theorem]{Remark}
\theoremstyle{plain}
\renewcommand{\qedsymbol}{\textit{q.e.d.}}
\begin{document}
\maketitle

\begin{abstract}
We consider pure Yang--Mills theory with a compact connected simple gauge group on a four-dimensional periodic lattice of spacing $\varepsilon$, with the bare coupling fixed by Ba{\l}aban's renormalisation group so that the running coupling at a fixed unit scale has a prescribed small value. For every finite family of pairwise disjoint simple closed lattice loops, in any planes, each with the character of an arbitrary irreducible representation, we divide the expectation of the product of their Wilson loops by an explicit number for each loop, the exponential of a perimeter term and of a term for each corner, whose coefficients do not depend on the loop. The coupling is fixed independently of the representations, which enter only through the scale at which the number is evaluated. We prove that the resulting renormalised expectation is bounded uniformly in $\varepsilon$ and converges as $\varepsilon \to 0$ along the sequence $\varepsilon = uL^{-K}$. The limits are reflection positive, and they do not vanish for loops whose lengths lie below a bound which depends on their representations and tends to infinity as the coupling tends to zero.
\end{abstract}

\setcounter{tocdepth}{2}
{\setlength{\parskip}{0pt}\tableofcontents}

\section{Introduction}\label{sec:introduction}

The expectation of a Wilson loop in lattice Yang--Mills theory is not expected to have a limit as the lattice spacing $\varepsilon$ tends to zero. Already at the level of perturbation theory the loop acquires a self-energy proportional to its length in lattice units, so that one expects a factor of the form $\exp(-m(\varepsilon)P/\varepsilon)$ for a loop of physical length $P$, together with one further factor for each of its corners. The question that we address here is whether these are the only divergences: if one divides the expectation by a number of this form, with coefficients that do not depend on the loop, does the quotient converge?

We answer this question for the pure gauge theory on a four-dimensional torus, at small values of the renormalised coupling. The construction on which everything rests is Ba{\l}aban's renormalisation group for lattice gauge theories \cite{BalabanRG1,BalabanRG2,BalabanCRE,BalabanLFR1,BalabanLFR2}, which controls the effective densities obtained from the lattice measure by successive block averaging, uniformly in the number of steps. Ba{\l}aban's results concern the partition function and the effective actions, however, and not the expectations of observables, so that the Wilson loops have to be carried through the construction. This is what we do: we insert the loops into the bare density and follow them from step to step, as a line defect whose effect on the density, after each step, is the product of a number and of a function of the block field that is close to the Wilson loop of the coarser lattice. The numbers, collected over the steps, are the normalisation by which we divide.

\subsection{The Lattice Theory}\label{sec:lattice-theory}

Let $G$ be a compact connected simple Lie group in a fixed unitary realisation $G \subset U(N)$, and write $\tr \coloneqq N^{-1}\operatorname{Tr}$ for the normalised trace. We fix an odd integer $L$, a length $u > 0$, and an integer $m_T \geq 0$. The torus is $\mathbb T \coloneqq (\mathbb R/T\mathbb Z)^4$ with the period $T \coloneqq 2L^{m_T}u$, and for an integer $K \geq 0$ the lattice of spacing
\begin{equation}\label{eq:lattice-spacing}
    \varepsilon_K \coloneqq uL^{-K}
\end{equation}
is $\mathbb T_K \coloneqq (\varepsilon_K\mathbb Z/T\mathbb Z)^4 \subset \mathbb T$. A bond of $\mathbb T_K$ is an ordered pair $b = (x, x \pm \varepsilon_Ke_\mu)$ of neighbouring sites, with $e_0, \dots, e_3$ the coordinate unit vectors, and $b^{-1}$ denotes the reversed bond. A gauge field is a map $U$ from the bonds to $G$ with $U(b^{-1}) = U(b)^{-1}$. For a path $\gamma = (b_1, \dots, b_n)$ of consecutive bonds we write $U(\gamma) \coloneqq U(b_1)\cdots U(b_n)$, and $\partial p$ is the boundary of a plaquette $p$. The lattice theory at the bare coupling $g_0 > 0$ is the probability measure
\begin{equation}\label{eq:lattice-measure}
    \dd\mu_{K,g_0}(U) \coloneqq \frac{1}{Z_{K,g_0}}\exp\Bigl(-\frac{1}{g_0^2}\sum_{p}\bigl(1 - \operatorname{Re}\tr U(\partial p)\bigr)\Bigr)\,\dd U\;,
\end{equation}
where the sum runs over the unoriented plaquettes of $\mathbb T_K$ and $\dd U$ is the product of the normalised Haar measures over the positive bonds.

The bare coupling is not a free parameter in the limit $K \to \infty$. Ba{\l}aban's renormalisation group assigns to $g_0$ a sequence of running couplings $g_1, \dots, g_K$, the coupling $g_k$ belonging to the lattice of spacing $L^k\varepsilon_K$, and we fix $g_0$ by the requirement that the running coupling at the unit scale $u$ has a prescribed value $\bar g$. More precisely, Theorem~\ref{thm:coupling-flow} provides, for every sufficiently small $\bar g > 0$ and every $K$, a bare coupling $g_0(K, \bar g)$ for which $g_K = \bar g$, and constants $0 < b \leq b'$ and $C_W$ such that
\begin{equation}\label{eq:bare-coupling-window}
    \bar g^{-2} + bK - C_W \leq g_0(K, \bar g)^{-2} \leq \bar g^{-2} + b'K + C_W\;.
\end{equation}
The bare coupling thus tends to zero like $(\log(1/\varepsilon_K))^{-1/2}$, which is asymptotic freedom in the form in which it enters here. We write $\mathbb E_K$ for the expectation with respect to $\mu_{K,g_0(K,\bar g)}$.

\subsection{Loop Ensembles}\label{sec:loop-ensembles}

The loops that we consider are closed lattice paths without self-intersections. Since the lattices $\mathbb T_K$ are nested, a path of bonds of $\mathbb T_q$ is also a path of bonds of $\mathbb T_K$ for every $K \geq q$, and it is in this sense that a loop is fixed while the cutoff is removed.

\begin{definition}[Loops and Ensembles]\label{def:loop-ensemble}
Let $q \geq 0$ be an integer.
\begin{enumerate}[label=(\roman*)]
\item A \emph{loop at level $q$} is a closed path $C$ of bonds of $\mathbb T_q$ which visits no site twice and whose $\ell^\infty$-diameter $\diam(C)$, measured along a lift of $C$ to $\mathbb R^4$, satisfies $3\diam(C) < T$; its \emph{corners} are the sites at which it changes direction. A loop carries its orientation.
\item A \emph{loop ensemble at level $q$} is a finite family $\mathcal C = ((C_1, \chi_1), \dots, (C_n, \chi_n))$, $n \geq 1$, of pairwise vertex-disjoint loops at level $q$, each with a character $\chi_i = (\dim\pi_i)^{-1}\operatorname{Tr}\pi_i$ of an irreducible unitary representation $\pi_i$ of $G$.
\end{enumerate}
\end{definition}

The loops need not be planar, and no assumption is made on their relative position beyond disjointness. A loop at level $q$ is also a loop at every level $q' \geq q$. For a gauge field $U$ on $\mathbb T_K$, $K \geq q$, the Wilson loop of $(C, \chi)$ is $W_\chi(C; U) \coloneqq \chi(U(C))$, which does not depend on the site at which the closed path $C$ is started, and we write
\begin{equation}\label{eq:ensemble-observable}
    W_{\mathcal C}(U) \coloneqq \prod_{i=1}^{n}W_{\chi_i}(C_i; U)
\end{equation}
for the observable of an ensemble.

Two quantities measure the shape of a loop, and all our constants depend on the loops only through them and through the perimeters. The first is the chord--arc constant $\Lambda(C) \geq 2$ of Definition~\ref{def:chord-arc}, the largest ratio between the distance of two points along the loop and their $\ell^\infty$-distance. The second is the ratio $P(C)/r(C)$ of the perimeter to the corner isolation $r(C)$ of Definition~\ref{def:lengths}, which is the smallest of half the shortest straight segment, the least distance between two non-adjacent segments, and an eighth of the perimeter. Both are invariant under dilations, and for a square $\Lambda(C) = 2$ and $P(C)/r(C) = 8$. For an ensemble we put $\Lambda(\mathcal C) \coloneqq \max_i\Lambda(C_i)$.

\begin{remark}[The Chord--Arc Constant and Self-Contact]\label{rem:chord-arc}
The chord--arc constant measures how close a loop comes to touching itself, relative to the length of the loop between the two places which come close. Near a corner two points are close both in space and along the loop, which is why $\Lambda(C) \geq 2$ for every loop, and the constant becomes large only when two stretches of $C$ which are far apart along the loop come close to each other in space, as the two long sides of a thin rectangle do: for a rectangle with sides $a \geq b$ one finds $\Lambda(C) = \max(2, 1 + a/b)$. Since $\Lambda(C)$ does not change under dilations, the bound $\Lambda(C) \leq \Lambda_*$ keeps a loop away from contact with itself uniformly on all scales, whatever its size. This is the property which our estimates use: any two points of $C$ in a cube of side $s$ are joined by an arc of $C$ of length at most $\Lambda(C)s$, so that the loop cannot leave a cube and return to it except along a short arc, and the number of its bonds which meet the cube is at most $2\Lambda(C)(s + 1)$ (Lemma~\ref{lem:square-properties}).

The restriction is also natural from the point of view of renormalisation, although nothing in the proofs depends on what follows. The perimeter factor and one factor for each corner are what perturbation theory predicts for a loop which does not meet itself \cite{Polyakov,BrandtNeriSato}. A loop which passes twice through the same point is not renormalised by factors of this kind: in perturbation theory its expectation mixes with those of the loops obtained by rerouting the path at the point of self-intersection \cite{BrandtNeriSato}, so that no number attached to the loop alone, of the kind by which we divide, removes its divergences. A lattice loop in the sense of Definition~\ref{def:loop-ensemble} never meets itself, but as $\Lambda(C)$ grows it comes closer, relative to its own length, to a loop which does. The constants of our estimates grow with $\Lambda(C)$, and so does the lower bound $L_0(\Lambda_*)$ on the blocking factor below; we do not show that this growth is necessary.
\end{remark}

\subsection{The Normalisation and the Results}\label{sec:results}

The renormalisation group is run from the lattice $\mathbb T_K$ down to a lattice whose spacing does not depend on $K$, on which the loops of an ensemble at level $q$ are paths of a fixed number of bonds. This lattice has the spacing $uL^{-j_q}$, where the \emph{stopping index} $j_q \geq q$ of Definition~\ref{def:stopping-index} depends on $q$, on $\bar g$, and on the constants of the construction, and not on $K$. The run thus consists of $K' \coloneqq K - j_q$ steps. Each step $k \leq K'$ contributes to a loop $C$ a factor $\kappa_k(C, \chi) \neq 0$, and the \emph{normalisation} of $(C, \chi)$ at the cutoff $K$ is the product
\begin{equation}\label{eq:normalisation-product}
    \kappa_K^{(q)}(C, \chi) \coloneqq \prod_{k=1}^{K - j_q}\kappa_k(C, \chi)\;,
\end{equation}
defined in Definition~\ref{def:normalisation}. Its structure is the one described at the beginning: $\log\kappa_k(C, \chi)$ is the sum of a term proportional to the total length of the segments of $C$ of each orientation and of a term for each corner of $C$, with coefficients that are numbers of the step $k$ and of the character and depend neither on the loop nor on the torus.

\begin{definition}[Renormalised Expectation]\label{def:renormalised-expectation}
For a loop ensemble $\mathcal C$ at level $q$ and $K > j_q$,
\begin{equation}\label{eq:renormalised-expectation}
    W_K(\mathcal C) \coloneqq \frac{\mathbb E_K[W_{\mathcal C}]}{\prod_{i=1}^{n}\kappa_K^{(q)}(C_i, \chi_i)}\;.
\end{equation}
\end{definition}

Since a loop at level $q$ is also a loop at every level $q' \geq q$, the level is part of the data of an ensemble: it fixes the stopping index $j_q$ and with it the number of step factors in \eqref{eq:normalisation-product}, while the step factors themselves do not depend on it.

The blocking factor $L$ and the other constants of the construction are fixed first, in the order of Appendix~\ref{app:constants}, and the coupling $\bar g$ after them. These choices decide which loops the results cover, and they decide it loop by loop and through the shape alone: a loop is covered if it is not too irregular for the chosen $L$, whatever its size and its representation, and then any number of such loops may be combined. The condition on the shape is a bound on the chord--arc constant: there is an increasing function $L_0$ on $[2, \infty)$, defined in \eqref{eq:blocking-factor-condition}, and we fix $\Lambda_* \geq 2$ with $L \geq L_0(\Lambda_*)$.

The representations enter in another way. A character is normalised, $\abs{\chi} \leq 1$, but the insertions which a loop produces are built from the derivatives of $\pi(U(C))$, and their size is measured by
\begin{equation}\label{eq:representation-constant}
    c_\pi \coloneqq \sup\bigl\{ \norm{\dd\pi(X)} \,\big|\, X \in \operatorname{Lie}G,\ \norm{X} \leq 1 \bigr\}\;,
\end{equation}
where $\dd\pi$ is the differential of $\pi$ on the Lie algebra $\operatorname{Lie}G \subset \mathfrak u(N)$ and $\norm{\cdot}$ is the operator norm. This number grows with the highest weight of $\pi$: for $G = SU(2)$ and the representation of spin $j$, $c_\pi = 2j$, and $\chi(\operatorname{diag}(e^{i\phi}, e^{-i\phi})) = \sin((2j + 1)\phi)/((2j + 1)\sin\phi)$ is close to $1$ on fluctuations $\phi$ of the order of a coupling $g$ only if $(2j + 1)g$ is small. Every expansion of the construction in which a loop takes part therefore requires $c_\pi$ times the running coupling of its step to be small, and at a fixed coupling $\bar g$ no single bound on the running couplings achieves this for every $\pi$. The loops, however, take part only in the steps $k \leq K'$, which integrate the scales finer than the stopping scale $uL^{-j_q}$, and at these scales the running coupling is smaller than $\bar g$ by asymptotic freedom: Section~\ref{sec:stopping-index} gives $g_k \leq \gamma_q$ for $k \leq K'$, with a number $\gamma_q$ which does not depend on $K$ and tends to zero as $q \to \infty$, at every $\bar g$. A large representation is therefore accommodated by reading its loop at a fine enough level, that is, by stopping the run and the product of the step factors at a fine enough scale, while the coupling $\bar g$, the construction, and the step factors stay the same. The estimates see a representation only through $c_\pi$, by the bound $\norm{\dd\pi(X)} \leq c_\pi\norm{X}$ and its consequences for the transports $\pi(U(J))$ at complex fields, while the normalised trace of a product is bounded by the product of the operator norms; every constant of the paper which depends on the representations of an ensemble is therefore a nondecreasing function of $c_{\mathcal C} \coloneqq \max_ic_{\pi_i}$. The size of a loop at level $q$ is measured by the number $P_q(C)$ of its bonds in $\mathbb T_q$, so that its physical length is $P_q(C)uL^{-q}$; it enters the constants of our estimates and none of their conditions.

\begin{definition}[Admissible Loops]\label{def:admissible}
Let $L$ and $\Lambda_*$ be fixed. A loop $C$ at level $q$ is \emph{admissible} if
\begin{equation}\label{eq:admissible}
    \Lambda(C) \leq \Lambda_*\;.
\end{equation}
\end{definition}

Admissibility is thus a condition on the shape alone. The size of a loop enters no condition because every estimate of the paper which involves a loop is a sum, over its traversals at some scale, of quantities which are small uniformly in the loop, and only these quantities are ever required to be small: a long loop makes the sums, and with them the constants, large, but it makes none of the summands larger. Our main result is the following.

\begin{theorem}[Uniform Bound and Limit]\label{thm:main}
Let $L$ and $\Lambda_*$ be as above. There are $\gamma_0 > 0$ and a function $q_\sharp \colon [0, \infty) \times (0, \gamma_0] \to \{ 0, 1, 2, \dots \}$, nondecreasing in each argument and, for every $c$, equal to $0$ for all sufficiently small $\bar g$, which depend only on the constants of the construction and on $\Lambda_*$, such that the following holds for every $\bar g \in (0, \gamma_0]$. For every loop ensemble $\mathcal C$ at a level
\begin{equation}\label{eq:level-condition}
    q \geq q_\sharp(c_{\mathcal C}, \bar g)
\end{equation}
whose loops are admissible, whatever its number $n$ of loops and the irreducible representations $\pi_i$ of its characters, there are an integer $K_0$ and $B < \infty$ such that
\begin{equation}\label{eq:main-bound}
    \sup_{K \geq K_0}\abs{W_K(\mathcal C)} \leq B
    \qquad\text{and}\qquad
    \sum_{K \geq K_0}\abs{W_{K+1}(\mathcal C) - W_K(\mathcal C)} < \infty\;.
\end{equation}
In particular the limit
\begin{equation}\label{eq:main-limit}
    W(\mathcal C) \coloneqq \lim_{K \to \infty}W_K(\mathcal C)
\end{equation}
exists. The numbers $K_0$ and $B$ depend on the constants of the construction, on $\bar g$, $q$, and $T/u$, and on $\mathcal C$ through $n$, the characters, the ratios $P(C_i)/r(C_i)$, and the numbers $P_q(C_i)$.
\end{theorem}

The level enters the theorem only through the normalisation, and for the continuum limit it is a choice without physical consequence. The loops are fixed paths in the torus, and the level of an ensemble decides only the scale $uL^{-j_q}$ at which the product \eqref{eq:normalisation-product} of the step factors is stopped, that is, the scale at which the normalisation is fixed. Writing $W^{(q)}_K$ for the renormalised expectation with the normalisation of the level $q$, we have, for levels $q' \geq q$ and every $K > j_{q'}$,
\begin{equation}\label{eq:level-change}
    W^{(q')}_K(\mathcal C) = W^{(q)}_K(\mathcal C)\prod_{i=1}^{n}\prod_{k=K-j_{q'}+1}^{K-j_q}\kappa_k(C_i, \chi_i)\;,
\end{equation}
by \eqref{eq:normalisation-product} and because the step factors do not depend on the level. The product runs over the $j_{q'} - j_q$ steps between the two stopping scales, a number which does not depend on $K$, and each of its factors is the exponential of a perimeter term and of a term for each corner, with coefficients which do not depend on the loop. A change of the level is therefore a finite renormalisation of the perimeter and corner terms, whose finite parts are a matter of convention in any renormalisation of Wilson loops, and a quantity in which the normalisation cancels does not depend on the level at all: for instance $W_K(\mathcal C)/\prod_iW_K((C_i, \chi_i))$ is the quotient $\mathbb E_K[W_{\mathcal C}]/\prod_i\mathbb E_K[W_{\chi_i}(C_i; \cdot)]$ whenever the denominators do not vanish. Since an ensemble at a level $q_0$ is an ensemble at every level $q \geq q_0$, the condition \eqref{eq:level-condition} restricts only this choice and not the ensembles, and the theorem covers, at every coupling $\bar g \leq \gamma_0$, every ensemble of admissible loops with arbitrary irreducible representations. The level required is bounded in the coupling, $q_\sharp(c, \bar g) \leq q_\sharp(c, \gamma_0)$, and it is given by \eqref{eq:level-threshold}. The restriction to irreducible representations concerns the normalisation only: the character of a reducible representation is the convex combination of the characters of its irreducible constituents, weighted by their dimensions, and the theorem applies to each term with its own normalisation.

The requirements on the coupling in Sections~\ref{sec:construction}--\ref{sec:reflection-positivity} are of two kinds. Those which involve no representation are upper bounds on $\bar g$ in terms of the constants of the construction and of $\Lambda_*$, and $\gamma_0$ is the largest number below which all of them hold. Those which involve the representations of an ensemble $\mathcal C$ concern only the steps $k \leq K'$ at which its loops are carried, and they are imposed as upper bounds on the number $\gamma_q$ of \eqref{eq:stopping-coupling} in terms of $c_{\mathcal C}$, of $\Lambda_*$, and of the constants of the construction; we then say that $\gamma_q$ is small, depending on $c_{\mathcal C}$. A coupling and a level ``as in Theorem~\ref{thm:main}'' for $\mathcal C$ are a coupling $\bar g \in (0, \gamma_0]$ and a level $q$ at which $\gamma_q$ satisfies all the requirements of the second kind. The function $q_\sharp$ is computed from these requirements in Section~\ref{sec:cauchy-proof}.

The bound and the limit are proved by different means. The bound follows from an exact identity, Theorem~\ref{thm:stopping-identity}, which expresses $W_K(\mathcal C)$ as a convex combination, over the terms of Ba{\l}aban's expansion of the density at the stopping scale, of expectations of functions which are bounded uniformly in $K$. The limit requires a comparison of the two runs with $K$ and $K + 1$ steps, which share all their lattices but the finest one, and this comparison, Theorem~\ref{thm:cauchy}, is the longest part of the paper.

The limits inherit the symmetries of the lattice measures, as far as the normalisations share them. The torus has the same period in the four directions, and for $K \geq q$ the lattice $\mathbb T_K$ is mapped to itself by every map $h(x) = Ox + t$ of $\mathbb T$ in which $O$ is a signed permutation matrix and $t \in uL^{-q}\mathbb Z^4$. We call these maps the \emph{symmetries of level $q$}, and we write $\mathfrak G_q$ for the group of those whose matrix $O$ is diagonal. Thus $\mathfrak G_q$ is generated by the translations by vectors of $uL^{-q}\mathbb Z^4$ and by the reflections $\theta_\mu$, $\mu \in \{ 0, \dots, 3 \}$, which replace $x_\mu$ by $-x_\mu$, and it contains, for every $\mu$ and every $a$ with $2a \in uL^{-q}\mathbb Z$, the reflection $\theta_{\mu,a}$ which replaces $x_\mu$ by $2a - x_\mu$ modulo $T$; the symmetries of level $q$ which do not belong to $\mathfrak G_q$ are those which exchange two axes. A symmetry $h$ of level $q$ is an isometry of $\mathbb T_q$ for the distance $\abs{\cdot}_\infty$ which maps segments to segments and corners to corners, and it therefore maps a loop $C$ at level $q$ to a loop $hC$ at level $q$, with the image orientation, $\Lambda(hC) = \Lambda(C)$, $P(hC)/r(hC) = P(C)/r(C)$, and $P_q(hC) = P_q(C)$. For an ensemble $\mathcal C$ we write $h\mathcal C \coloneqq ((hC_1, \chi_1), \dots, (hC_n, \chi_n))$, with the same characters, and $\bar{\mathcal C}$ for the ensemble of the loops of $\mathcal C$ with the conjugate characters $\bar\chi_i$. The character $\bar\chi$ is that of the conjugate representation $\bar\pi$, whose differential $\dd\bar\pi(X)$ is the complex conjugate of the matrix $\dd\pi(X)$ and has the same operator norm, so that $c_{\bar\pi} = c_\pi$ and $c_{\bar{\mathcal C}} = c_{\mathcal C}$. Finally $C^{-1}$ denotes the loop $C$ with the opposite orientation.

\begin{corollary}[Symmetries of the Limits]\label{cor:symmetries}
Let $L$, $\Lambda_*$, $\gamma_0$, and $q_\sharp$ be as in Theorem~\ref{thm:main}, let $\bar g \in (0, \gamma_0]$, and let $\mathcal C$ be a loop ensemble at a level $q \geq q_\sharp(c_{\mathcal C}, \bar g)$ whose loops are admissible. Then the limit $W(\mathcal C')$ exists for each of the ensembles $\mathcal C'$ below, and:
\begin{enumerate}[label=(\roman*)]
\item $W(\mathcal C)$ does not depend on the order of the loops of $\mathcal C$;
\item $W(g\mathcal C) = W(\mathcal C)$ for every $g \in \mathfrak G_q$;
\item $W(\bar{\mathcal C}) = \overline{W(\mathcal C)}$, and $W(\mathcal C)$ does not change if a loop $C_i$ of $\mathcal C$ is replaced by $C_i^{-1}$ and its character by $\bar\chi_i$;
\item $W(h\mathcal C) = W(\mathcal C)$ for every symmetry $h$ of level $q$ for which there are $g_1, \dots, g_n \in \mathfrak G_q$ with $hC_i = g_iC_i$ for every $i$.
\end{enumerate}
The same identities hold for $W_K$ at every $K > j_q$.
\end{corollary}

Statement (ii) is the covariance under translations and under the reflections in coordinate hyperplanes, with unchanged characters, and statement (iii) is the invariance under charge conjugation. Since an ensemble at level $q$ is an ensemble at every level $q' \geq q$, statement (ii) gives the covariance under every translation of the dense group $\bigcup_q uL^{-q}\mathbb Z^4$, at a level at which it is a symmetry; the normalisations of two levels differ by the factors \eqref{eq:level-change}, which are the same for an ensemble and its translate. Statement (iv) applies, for instance, to the reflection which exchanges $x_\mu$ and $x_\nu$ and to an ensemble of squares in planes spanned by $e_\mu$ and $e_\nu$, since the image of such a square is also its image under $\theta_\mu$ followed by a translation. For a general ensemble we do not know whether the limits are covariant under the exchange of two axes: the lattice measures are invariant under it, but the normalisations are not shown to be, for the reason explained in Section~\ref{sec:symmetries-proof}.

The limits also inherit the reflection positivity of the lattice measures, for the reflections $\theta_{\mu,a}$ with $a \in uL^{-q}\mathbb Z$, which we call the \emph{reflections of level $q$}. The points fixed by such a reflection form the two hyperplanes $\{ x_\mu \equiv a \}$ and $\{ x_\mu \equiv a + T/2 \}$, coordinates being taken modulo $T$, and their complement consists of two open half-tori, $\{ a < x_\mu < a + T/2 \}$ and $\{ a + T/2 < x_\mu < a + T \}$, which the reflection exchanges. The pairing of reflection positivity combines the reflection of a loop with the conjugation of its character, and so involves the ensembles $\theta\bar{\mathcal D}$ of the reflected loops $\theta C$, with the image orientations and the conjugate characters; we have $c_{\theta\bar{\mathcal D}} = c_{\mathcal D}$. We admit the empty ensemble, with $W(\emptyset) \coloneqq 1$.

\begin{corollary}[Reflection Positivity]\label{cor:reflection-positivity}
Let $L$, $\Lambda_*$, $\gamma_0$, and $q_\sharp$ be as in Theorem~\ref{thm:main}, let $\bar g \in (0, \gamma_0]$ and $q \geq 0$, let $\theta$ be a reflection of level $q$, and let $H_+$ be one of the two open half-tori bounded by the points fixed by $\theta$. Let $\mathcal D_1, \dots, \mathcal D_m$ be loop ensembles at level $q$, possibly empty, whose loops are admissible and lie in $H_+$, and assume $q \geq q_\sharp(c, \bar g)$, where $c$ is the largest of the numbers $c_\pi$ of their representations. Then the limits $W(\mathcal D_a \cup \theta\bar{\mathcal D}_b)$ exist for $0 \leq a, b \leq m$, with $\mathcal D_0 \coloneqq \emptyset$, and:
\begin{enumerate}[label=(\roman*)]
\item the matrix $\bigl(W(\mathcal D_a \cup \theta\bar{\mathcal D}_b)\bigr)_{1 \leq a, b \leq m}$ is Hermitian and positive semidefinite;
\item $W(\theta\bar{\mathcal D}_a) = \overline{W(\mathcal D_a)}$ for every $a$, and the matrix $\bigl(W(\mathcal D_a \cup \theta\bar{\mathcal D}_b) - W(\mathcal D_a)\overline{W(\mathcal D_b)}\bigr)_{1 \leq a, b \leq m}$ is Hermitian and positive semidefinite.
\end{enumerate}
\end{corollary}

The statement (ii) is positivity for the connected functions, and it is not empty: at small coupling the limits do not vanish on the ensembles whose loops are short compared with a length which grows without limit as $\bar g \to 0$ (Proposition~\ref{prop:nonvanishing}), and for a loop $C$ in $H_+$ to which this applies the ensemble $\mathcal D \coloneqq ((C, \chi))$ has $W(\mathcal D \cup \theta\bar{\mathcal D}) \geq \abs{W(\mathcal D)}^2 \geq 1/4$. In the other direction the limits satisfy a bound which is a product over the loops of factors exponential in their lengths, with a rate which tends to zero with the coupling (Proposition~\ref{prop:growth-bound}).

Finally we ask whether the limits are new, that is, whether they could be the Wilson-loop functions of a Gaussian field. The question has to be posed with some care, because the normalisation is a convention: a function which differs from $W$ by a factor $\exp(\sum_v\lambda_vN_v(C) + \sum_c\zeta_{v_c,w_c})$ for each loop, where $N_v(C)$ is the total length of the segments of $C$ of orientation $v$, the second sum runs over the corners, and the coefficients do not depend on the loop, is the same object with another normalisation. For an abelian Gaussian field the answer is then contained in a single number. If $\mathsf B$ is a real Gaussian $1$-form and $\mathsf g$ a charge, then $\log\abs{\mathbb E[e^{i\mathsf g\oint_C\mathsf B}]} = -\frac12\mathsf g^2\operatorname{Var}\oint_C\mathsf B$ is linear in the covariance of $\mathsf B$ and a quadratic form in the current of the loop, a change of normalisation adds to it a function which is linear in the lengths and in the numbers of corners of each type, and a suitable signed combination of these logarithms over a few rectangles and pairs of rectangles annihilates both, whatever the covariance and whatever the normalisation.

To state this precisely, we fix a level $q$, a coordinate plane spanned by $e_\mu$ and $e_\nu$, lengths $t, s \in uL^{-q}\mathbb N$ with $3\max(4t, s) < T$, a point $o \in uL^{-q}\mathbb Z^4$, and a nontrivial irreducible unitary representation $\pi$ of $G$ with the character $\chi$. For $a \in \{ 0, \dots, 3 \}$ and $n \in \{ 1, \dots, 4 \}$ let $I_a$ and $R_n$ be the boundaries of the rectangles $o + [at, (a + 1)t]e_\mu + [0, s]e_\nu$ and $o + [0, nt]e_\mu + [0, s]e_\nu$, traversed in the sense which runs from $e_\mu$ to $e_\nu$, so that $R_1 = I_0$, and let $R^{(1)}_n \coloneqq R_n + te_\mu$ for $n \in \{ 2, 3 \}$. The \emph{eleven diagrams} are the ensembles $((C, \chi))$ of one loop, for $C$ one of the nine rectangles $I_0, \dots, I_3$, $R_2$, $R_3$, $R_4$, $R^{(1)}_2$, $R^{(1)}_3$, and the pairs $P_k \coloneqq ((I_0, \chi), (I_k, \chi))$ for $k \in \{ 2, 3 \}$, whose loops are at distance $(k - 1)t \geq t$; we write $C$ also for the ensemble $((C, \chi))$. They are ensembles at level $q$. For a function $W$ on the eleven diagrams with nonzero values let
\begin{equation}\label{eq:witness}
    \mathfrak Y(W) \coloneqq \log\Bigl\lvert\frac{W(R_1)\,W(R_3)^3}{W(R_2)^3\,W(R_4)}\Bigr\rvert - \log\Bigl\lvert\frac{W(P_2)}{W(R_1)^2}\Bigr\rvert + \log\Bigl\lvert\frac{W(P_3)}{W(R_1)^2}\Bigr\rvert\;.
\end{equation}
A \emph{Gaussian charge model} is a triple $(\mathsf g, \mathsf X, \kappa)$ of a number $\mathsf g \in \mathbb R$, a real Gaussian vector $\mathsf X = (\mathsf X_0, \dots, \mathsf X_3)$ of any mean and covariance, and a function $\kappa$ on loops with values in $\mathbb C \setminus \{ 0 \}$ such that $\log\abs{\kappa(C)} = \sum_v\lambda_vN_v(C) + \sum_c\zeta_{v_c,w_c}$ with real numbers $\lambda_v$ and $\zeta_{v,w}$ which do not depend on $C$. Its function on the eleven diagrams is $W(\mathcal D) \coloneqq \mathbb E[e^{i\mathsf g\mathsf X(\mathcal D)}]/\prod_{C \in \mathcal D}\kappa(C)$, where $\mathsf X(\mathcal D)$ is the sum of the $\mathsf X_a$ over the indices $a$ for which $I_a$ is one of the rectangles that tile the rectangle $\mathcal D$, or one of the two loops of the pair $\mathcal D$. The models with $\mathsf X_a = \oint_{I_a}\mathsf B$ for a real Gaussian $1$-form $\mathsf B$, on a lattice which contains the rectangles or on the continuum torus, are of this kind, with $\mathsf X(C) = \oint_C\mathsf B$ for each of the nine rectangles, since the interior sides of a tiling are traversed once in each direction.

\begin{theorem}[The Limits Are Not Gaussian]\label{thm:non-gaussian}
Let $q$, the plane, $t$, $s$, $o$, and $\pi$ be as above, let $\Lambda_* \geq 1 + \max(s/t, 4t/s)$, and let $L$ be as in Theorem~\ref{thm:main}. There is $\bar g_0 > 0$, depending on $q$, $t/u$, $s/u$, $T/u$, $\pi$, and the constants of the construction, such that for every $\bar g \in (0, \bar g_0]$ the following hold.
\begin{enumerate}[label=(\roman*)]
\item The limits $W(\mathcal D)$ exist on the eleven diagrams, $\abs{W(\mathcal D)} \geq 1/2$, $W(I_a) = W(I_0)$, and $W(R^{(1)}_n) = W(R_n)$.
\item $\mathfrak Y(W) \geq \frac12c_O\bar g^2$, with $c_O \coloneqq C_2(\pi)s^2t^3/(672e^2T^5)$ and the Casimir number $C_2(\pi) > 0$ of \eqref{eq:casimir-number}, and $\mathfrak Y(W)$ is the same for the normalisations of every level $q' \geq q$.
\item No Gaussian charge model has the function $W$ on the eleven diagrams, and $W$ is not the limit on the eleven diagrams of the functions of a sequence of Gaussian charge models.
\end{enumerate}
\end{theorem}

The number $\mathfrak Y$ is a third difference in the width of the rectangles, corrected by the two \emph{pair terms} $\log\abs{W(P_k)/W(R_1)^2}$. In an abelian Gaussian model all three are of the order $\mathsf g^2$, and they cancel: the covariances of separated rectangles which make up the third difference also make up the pair terms. In the limits of Yang--Mills theory the pair terms are of the order $\bar g^4$, up to logarithms, because the term of first order of a non-abelian Wilson loop has trace zero, while the third difference is, at the order $\bar g^2$, the number $C_2(\pi)$ times that of the free field, which is positive. The same mechanism makes the theorem silent on Gaussian fields with values in the Lie algebra and path-ordered holonomies, which agree with Yang--Mills theory at the second order (Section~\ref{sec:gaussian-models}).

\subsection{The Argument}\label{sec:argument}

Let us describe how the proof is organised and where its difficulties lie. Ba{\l}aban's construction is a sequence of steps, each of which integrates the fluctuations of the gauge field at one scale. The integrand is first decomposed by test functions, which decide in every region of the lattice whether the field is small there, and the outcome is a sum of terms, labelled by the regions in which some test function has failed. On the regions where all of them pass, the fluctuation integral is performed by a cluster expansion around a background field, the minimiser of the action with given block averages. On the other regions no expansion is available, and the term carries instead a small factor from the action, which a separate operation, Ba{\l}aban's $\mathbf R$-operation \cite{BalabanLFR1,BalabanLFR2}, later turns into terms of the effective action.

The first difficulty is one of consistency and is independent of the loops. The thresholds of the test functions are functions of the running couplings, while the running couplings are defined by the construction, and it is not known whether the coupling of Ba{\l}aban's construction as printed is a continuous function of the bare coupling, which is what one needs in order to solve $g_K = \bar g$. We therefore run a variant of the construction, in which every threshold, every cube size, and the depth of the $\mathbf R$-operation are functions of the target $(K, \bar g)$ alone, and in which the plaquette test functions are continuous rather than characteristic functions. Section~\ref{sec:construction} defines this variant and shows that Ba{\l}aban's estimates hold for it, and Section~\ref{sec:coupling-flow} proves the continuity of the couplings and solves for the bare coupling.

The loops enter in Section~\ref{sec:first-step}, which treats the first step, and in Section~\ref{sec:inductive-class}, which treats the general one. After $k$ steps the contribution of an ensemble to a term of the density is a \emph{dressed loop}: the product of the normalisations $\kappa_1 \cdots \kappa_k$ of its loops, of the exponential of a small local functional, and of the sum of the Wilson loops of the background field and of a series of insertions along the loops, with activities that decay in the size of their supports. The heart of the paper is the statement that this form is reproduced by a step, with constants that do not depend on $k$, including on the terms with large fields and through the $\mathbf R$-operation. The perimeter and corner coefficients are extracted in Section~\ref{sec:flat-part} from the part of the loop's contribution that survives at the trivial background, and it is here that the geometry of a general loop is used: away from its corners a loop cannot be told from a straight line, and near a corner it cannot be told from two perpendicular rays, up to errors that decay exponentially in the corner isolation.

Stopping the induction at the index $K'$ gives the identity and the bound of Section~\ref{sec:stopping-scale}. For the limit we have to compare the dressed loops, and the densities themselves, of the runs with $K$ and $K + 1$ steps on their common lattices. In the regions of small fields this is a comparison of two expansions whose data differ by the one additional step, and the difference decays geometrically in the number of steps that separate a scale from the finest one. The terms with large fields require a different argument, since their test functions are evaluated on different fields in the two runs, and it is for this comparison that the test functions were made continuous. Sections~\ref{sec:large-field-comparison} and~\ref{sec:walk-factor-difference} prove the two inputs of the comparison, and Section~\ref{sec:cauchy} assembles them. Section~\ref{sec:reflection-positivity} derives the two corollaries, the lower bound on the limits, and the growth bound, Section~\ref{sec:non-gaussianity} proves Theorem~\ref{thm:non-gaussian}, and Section~\ref{sec:notes-on-papers} records where our reading of Ba{\l}aban's papers departs from the printed text.

The results of Ba{\l}aban's which we use in the setting in which he proved them are cited by page and equation. Those which we use outside his hypotheses, or of which he printed no proof, are proved again in a supplement to this paper, cited as ``Supplement, Proposition~S.$n$''.

\section{Loops and Their Geometry}\label{sec:loops}

All statements of this section are about closed paths on the lattice $\mathbb Z^4$ with the distance $\abs{x - y}_\infty \coloneqq \max_\mu\abs{x_\mu - y_\mu}$. They will be applied at every scale of the construction, with the lattice unit equal to the spacing of that scale, and their purpose is to replace the elementary facts about a square that such an argument uses without comment: that a short piece of it is straight or contains one corner, that few of its bonds lie in a small box, and that a piece containing half of it is large.

\subsection{Segments and Corners}\label{sec:segments-corners}

Let $\mathcal V \coloneqq \{ \pm e_0, \pm e_1, \pm e_2, \pm e_3 \}$ be the set of the eight orientations. A \emph{closed path} of length $n$ is a sequence of sites $x_0, x_1, \dots, x_n = x_0$ of $\mathbb Z^4$ with $v_i \coloneqq x_{i+1} - x_i \in \mathcal V$ for $0 \leq i < n$, and it is \emph{simple} if $x_0, \dots, x_{n-1}$ are distinct. We read the indices modulo $n$. A \emph{loop} is a simple closed path, and two loops which differ by a cyclic shift of the indices are the same loop. The bonds of the loop $C$ are the pairs $b_i = (x_i, x_{i+1})$, the \emph{position} of $b_i$ is its initial site $x_i$, and $\abs{C} \subset \mathbb R^4$ is the union of the closed unit intervals $[x_i, x_{i+1}]$. The number $P(C) \coloneqq n$ of bonds is the \emph{perimeter}.

A site $x_i$ of $C$ is a \emph{corner} if $v_i \neq v_{i-1}$, and $\mathcal K(C)$ is the set of the corners. A \emph{segment} of $C$ is a maximal run of consecutive bonds of one orientation, together with its two end sites, so that the segments are the pieces of $C$ between consecutive corners. At a corner $c = x_i$ we write $v_c \coloneqq v_{i-1}$ and $w_c \coloneqq v_i$ for the incoming and the outgoing orientation.

\begin{lemma}[Corners]\label{lem:corners}
Let $C$ be a loop.
\begin{enumerate}[label=(\roman*)]
\item For every $v \in \mathcal V$, $C$ has as many bonds of orientation $v$ as of orientation $-v$. In particular $P(C)$ is even.
\item At every corner $c$ the orientations $v_c$ and $w_c$ are perpendicular.
\item $C$ has at least four segments, and as many corners as segments.
\item $\diam(C) \leq P(C)/2$, where $\diam(C)$ is the $\ell^\infty$-diameter of $\abs{C}$.
\end{enumerate}
\end{lemma}

\begin{proof}
For (i) we sum the orientations around the loop. Since $\sum_{i<n}v_i = x_n - x_0 = 0$, the coefficient of each $e_\mu$ in this sum vanishes, and this coefficient is the number of bonds of orientation $e_\mu$ minus the number of bonds of orientation $-e_\mu$.

For (ii), $w_c \neq v_c$ by the definition of a corner. If $w_c = -v_c$ at $c = x_i$, then $x_{i+1} = x_i - v_{i-1} = x_{i-1}$, so that a site is visited twice, which is excluded because $n \geq 3$ for a simple closed path with $v_i \in \mathcal V$. Two elements of $\mathcal V$ which are neither equal nor opposite are perpendicular.

For (iii), the orientations $v_i$ are not all equal, since otherwise $x_n \neq x_0$. So $C$ has a corner, and by (ii) at least two coordinate axes carry bonds of $C$. By (i) each of these axes carries bonds of both of its orientations, so that at least four orientations occur, and each of them is the orientation of at least one segment. The segments and the corners alternate around the loop, which gives the last assertion.

For (iv), let $R_\mu$ be the difference between the largest and the smallest value of the coordinate $x_\mu$ on $C$. The path passes from a site with the smallest value to a site with the largest value and back, and each bond changes $x_\mu$ by at most one, so that $P(C) \geq 2R_\mu$. Since $\diam(C) = \max_\mu R_\mu$, this is (iv).
\end{proof}

Two segments are \emph{adjacent} if they are consecutive around the loop, that is, if they share a corner. By Lemma~\ref{lem:corners}(iii) two distinct segments share at most one corner, and every segment has a segment which is not adjacent to it.

\begin{definition}[Lengths]\label{def:lengths}
For a loop $C$, the \emph{shortest segment} $s(C)$ is the least number of bonds of a segment of $C$, the \emph{separation} $\delta(C)$ is the least $\ell^\infty$-distance between two distinct non-adjacent segments of $C$, as closed subsets of $\mathbb R^4$, and the \emph{corner isolation} is
\begin{equation}\label{eq:corner-isolation}
    r(C) \coloneqq \min\Bigl(\frac{s(C)}{2},\, \delta(C),\, \frac{P(C)}{8}\Bigr)\;.
\end{equation}
\end{definition}

The separation is positive, and it is at least as large as the mesh of any lattice which contains the corners. This is the content of the next lemma, which also separates the loops of an ensemble from each other.

\begin{lemma}[Separation]\label{lem:separation}
Let $a \geq 1$ be an integer and $x_* \in \mathbb Z^4$.
\begin{enumerate}[label=(\roman*)]
\item If all corners of a loop $C$ lie in $x_* + a\mathbb Z^4$, then every segment of $C$ has a number of bonds which is a positive multiple of $a$, and $s(C) \geq a$ and $\delta(C) \geq a$.
\item If $C$ and $C'$ are loops without a common site, all of whose corners lie in $x_* + a\mathbb Z^4$, then the $\ell^\infty$-distance between $\abs{C}$ and $\abs{C'}$ is at least $a$.
\end{enumerate}
\end{lemma}

\begin{proof}
We may take $x_* = 0$. A segment joins two corners which differ in one coordinate only, by a nonzero multiple of $a$, which gives the first assertion of (i) and $s(C) \geq a$.

Let now $A$ and $B$ be two closed axis-parallel intervals of $\mathbb R^4$ with end points in $a\mathbb Z^4$, $A$ parallel to $e_\mu$ and $B$ parallel to $e_\nu$. We claim that, if $A \cap B = \emptyset$, then their $\ell^\infty$-distance is at least $a$. On each of the two intervals all coordinates but one are constant and lie in $a\mathbb Z$. If a coordinate $x_\lambda$ with $\lambda \notin \{ \mu, \nu \}$ takes different values on $A$ and on $B$, these values differ by at least $a$, and the claim holds. Assume that this is not the case. If $\mu = \nu$, then $A$ and $B$ lie on one line, where they are disjoint intervals with end points in $a\mathbb Z$, and so they are at least $a$ apart. If $\mu \neq \nu$, let $b_\mu$ be the value of $x_\mu$ on $B$ and $a_\nu$ the value of $x_\nu$ on $A$. The intervals meet if and only if $b_\mu$ lies in the range of $x_\mu$ on $A$ and $a_\nu$ lies in the range of $x_\nu$ on $B$. Since they do not meet, one of these fails, say the first, and then $b_\mu \in a\mathbb Z$ lies at distance at least $a$ from that range, whose end points are in $a\mathbb Z$. So $\abs{x_\mu - y_\mu} \geq a$ for all $x \in A$ and $y \in B$, which proves the claim.

It remains to see that the claim applies. If two axis-parallel intervals with end points in $\mathbb Z^4$ meet, then they have a site of $\mathbb Z^4$ in common: on a common line two closed intervals with integer end points which meet share an integer point, and two perpendicular ones meet in the point with the coordinates $b_\mu$ and $a_\nu$ above, which is a site. A site which lies on two segments of the loop $C$ is a site $x_i$ of $C$, and, since $C$ is simple, it lies on two segments only if it is the corner between them. Two non-adjacent segments of $C$ are therefore disjoint, and the claim gives $\delta(C) \geq a$. In (ii) a segment of $C$ and a segment of $C'$ have no site in common by hypothesis, so they are disjoint, and the claim gives (ii).
\end{proof}

\subsection{The Chord--Arc Constant}\label{sec:chord-arc}

For two points $x, y \in \abs{C}$ let $d_C(x, y)$ be the length of the shorter of the two arcs of $\abs{C}$ which join them. Along an arc every coordinate changes by at most the length of the arc, so that
\begin{equation}\label{eq:chord-below-arc}
    \abs{x - y}_\infty \leq d_C(x, y) \leq \frac{P(C)}{2}\;.
\end{equation}
The constant which bounds $d_C$ from above by the $\ell^\infty$-distance is the one geometric quantity through which the shape of a loop enters most of our counts.

\begin{definition}[Chord--Arc Constant]\label{def:chord-arc}
For a loop $C$,
\begin{equation}\label{eq:chord-arc-constant}
    \Lambda(C) \coloneqq \sup\Bigl\{ \frac{d_C(x, y)}{\abs{x - y}_\infty} \, \Big| \, x, y \in \abs{C},\ x \neq y \Bigr\}\;.
\end{equation}
\end{definition}

\begin{lemma}[Bounds on the Chord--Arc Constant]\label{lem:chord-arc}
For every loop $C$,
\begin{equation}\label{eq:chord-arc-bounds}
    2 \leq \Lambda(C) \leq \max\Bigl(2,\, \frac{P(C)}{2\delta(C)}\Bigr) \leq \max\Bigl(2,\, \frac{P(C)}{2r(C)}\Bigr)\;,
    \qquad
    \abs{\mathcal K(C)} \leq \frac{P(C)}{s(C)} \leq \frac{P(C)}{2r(C)}\;.
\end{equation}
For a square, $\Lambda(C) = 2$ and $P(C)/r(C) = 8$.
\end{lemma}

\begin{proof}
For the lower bound let $c$ be a corner, and let $x$ and $y$ be the sites of $C$ before and after $c$. They differ by one in each of two coordinates, by Lemma~\ref{lem:corners}(ii), so that $\abs{x - y}_\infty = 1$, while the arc through $c$ has length $2$ and the other arc has length $P(C) - 2 \geq 2$, because $P(C) \geq 4$ by Lemma~\ref{lem:corners}(iii). Hence $d_C(x, y) = 2$.

For the upper bound let $x \neq y$ be points of $\abs{C}$. If they lie on one segment, the part of the segment between them is an arc of length $\abs{x - y}_\infty$. If they lie on two adjacent segments, which meet at the corner $c$, let $\alpha$ and $\beta$ be their distances from $c$. The two segments are perpendicular, so that $\abs{x - y}_\infty = \max(\alpha, \beta)$, and the arc through $c$ has length $\alpha + \beta \leq 2\max(\alpha, \beta)$. In the remaining case they lie on two non-adjacent segments, so that $\abs{x - y}_\infty \geq \delta(C)$, while $d_C(x, y) \leq P(C)/2$ by \eqref{eq:chord-below-arc}. The third inequality holds because $\delta(C) \geq r(C)$. Every segment has at least $s(C) \geq 2r(C)$ bonds, and there are as many segments as corners, which gives the bound on $\abs{\mathcal K(C)}$.

For a square of side $s$ we have $\delta(C) = s$ and $P(C) = 4s$, so that $\Lambda(C) = 2$ by the two bounds, and $r(C) = \min(s/2, s, s/2) = s/2$.
\end{proof}

All the quantities of Definition~\ref{def:lengths} and the perimeter are multiplied by $m$ when $C$ is replaced by its image $mC$ under the dilation $x \mapsto mx$, $m \geq 1$ an integer, and so are $d_C$ and the $\ell^\infty$-distance. Hence $\Lambda(mC) = \Lambda(C)$ and $P(mC)/r(mC) = P(C)/r(C)$: the two shape constants of a loop do not depend on the scale at which it is read.

\subsection{Arcs around a Set}\label{sec:arcs}

The dressed loops of the later sections are sums over finite sets of positions on a loop, and what one needs to know about such a set is how much of the loop it takes to connect it. On a square, a set of positions of diameter $D < s/2$ lies in an arc of length at most $4D$. The next lemma is the form of this statement for a general loop. By an \emph{arc} we mean a connected closed subset of $\abs{C}$, and a \emph{gap} of a set $S \subset \abs{C}$ is a connected component of $\abs{C} \setminus S$.

\begin{lemma}[Arcs around a Set]\label{lem:arcs-around-set}
Let $C$ be a loop, $S$ a nonempty set of sites of $C$ of $\ell^\infty$-diameter $D$, and $x_0 \in S$. Put $\Lambda \coloneqq \Lambda(C)$ and $A \coloneqq \{ x \in \abs{C} \, | \, d_C(x_0, x) \leq \Lambda D \}$.
\begin{enumerate}[label=(\roman*)]
\item $S \subset A$, and $S$ has at most $2\Lambda D + 1$ elements.
\item If $4\Lambda D < P(C)$, then $S$ has exactly one gap of length greater than $P(C)/2$. Its complement $I(S)$ is the unique shortest arc which contains $S$. It lies in $A$, its length is at most $2\Lambda D$, and each of its points is within $\ell^\infty$-distance $\Lambda D$ of $x_0$.
\item If $4\Lambda D \geq P(C)$, then every point of $\abs{C}$ is within $\ell^\infty$-distance $2\Lambda D$ of every point of $S$.
\end{enumerate}
\end{lemma}

\begin{proof}
For $x \in S$ the definition of $\Lambda$ gives $d_C(x_0, x) \leq \Lambda\abs{x - x_0}_\infty \leq \Lambda D$, so that $S \subset A$. The sites of $C$ in $A$ are those at arc distance at most $\Lambda D$ from $x_0$ on either side, at most $2\floor{\Lambda D} + 1$ in number if $2\Lambda D < P(C)$, and otherwise all the $P(C) \leq 2\Lambda D$ sites of $C$. This is (i).

For (ii), $A$ is the arc of length $2\Lambda D < P(C)/2$ centred at $x_0$, and its complement is an open arc $G_0$ of length greater than $P(C)/2$ which contains no point of $S$ by (i). Let $G$ be the gap of $S$ which contains $G_0$. The gaps of $S$ are pairwise disjoint, so that every other gap has length less than $P(C)/2$, and $G$ is the only gap longer than $P(C)/2$. An arc which contains $S$ and is not all of $\abs{C}$ is the complement of an open arc inside a gap, and it is shortest when that open arc is the longest gap. So $I(S) \coloneqq \abs{C} \setminus G$ is the unique shortest arc containing $S$, and $I(S) \subset \abs{C} \setminus G_0 = A$. The last assertion follows from \eqref{eq:chord-below-arc}.

For (iii), two points of $\abs{C}$ are at $\ell^\infty$-distance at most $\diam(C) \leq P(C)/2$ by Lemma~\ref{lem:corners}(iv), and $P(C)/2 \leq 2\Lambda D$ by hypothesis.
\end{proof}

In both cases of the lemma the shortest arc containing $S$, taken to be all of $\abs{C}$ in the case (iii), has length at most $4\Lambda D$ and lies within $\ell^\infty$-distance $2\Lambda D$ of $S$. For a square these are the constants $8D$ and $4D$.

A cube which is smaller than the separation sees at most one corner of the loop, and the counts become those of a bent line.

\begin{lemma}[Small Cubes]\label{lem:small-cubes}
Let $C$ be a loop and $Q$ a closed cube of side $t < \delta(C)$ with corners in $\mathbb Z^4$. Then $Q \cap \abs{C}$ is empty or an arc contained in one segment or in two adjacent segments. It contains at most $2t + 1$ sites of $C$, and at most $2t + 2$ bonds of $C$ meet $Q$.
\end{lemma}

\begin{proof}
Two segments which meet $Q$ are within $\ell^\infty$-distance $t < \delta(C)$ of each other, so they are equal or adjacent. Three distinct segments cannot be pairwise adjacent, because $C$ has at least four segments and adjacency is the cyclic order, so that at most two segments meet $Q$. A segment meets the convex set $Q$ in an interval, which has length at most $t$ and hence at most $t + 1$ sites. If two segments meet $Q$, in points $x$ and $y$, let $c$ be their common corner, and let the segments be parallel to $e_\mu$ and $e_\nu$. Then $c$ has the coordinate $x_\nu$ of $x$, the coordinate $y_\mu$ of $y$, and the other coordinates in common with both, so that $c$ lies in the coordinate box spanned by $x$ and $y$, which is contained in the cube $Q$. Both intervals thus contain $c$, their union is an arc, and it has at most $2(t + 1) - 1$ sites. Finally, a bond which meets $Q$ has an end site in $Q$, since $Q$ has integer corners, and an arc with $m$ sites of $C$ is met by at most $m + 1$ bonds of $C$.
\end{proof}

The next two statements are used for the pieces of a loop which are long. The first bounds the number of cubes of a partition which a piece meets, and the second shows that a piece which contains half of the loop cannot be small.

\begin{lemma}[Cubes along an Arc]\label{lem:cubes-along-arc}
Let $C$ be a loop, $M' \geq 1$ an integer, and consider the closed cubes $M'z + [0, M']^4$, $z \in \mathbb Z^4$. An arc of $n$ bonds of $C$ consists of at most $n/s(C) + 2$ straight pieces, and it meets at most
\begin{equation}\label{eq:cubes-along-arc}
    \frac{8n}{M'} + \frac{16n}{s(C)} + 32
\end{equation}
of these cubes.
\end{lemma}

\begin{proof}
The straight pieces of an arc are, in the order of the arc, a part of a segment, a number of whole segments, and a part of a segment. The whole segments have at least $s(C)$ bonds each, so there are at most $n/s(C)$ of them. A straight piece of $n_i$ bonds parallel to $e_\mu$ occupies an interval $[\tau, \tau + n_i]$ of the coordinate $x_\mu$, which meets the interval $[M'j, M'(j + 1)]$ only for the integers $j$ with $\tau/M' - 1 \leq j \leq (\tau + n_i)/M'$, of which there are at most $n_i/M' + 2$. Each of its three other coordinates is constant and lies in at most two of the intervals $[M'j, M'(j + 1)]$. The piece therefore meets at most $8(n_i/M' + 2)$ cubes, and the sum over the at most $n/s(C) + 2$ pieces, with $\sum_in_i = n$, is at most \eqref{eq:cubes-along-arc}.
\end{proof}

\begin{lemma}[Long Arcs]\label{lem:long-arcs}
Let $C$ be a loop. The set of the positions of at least $P(C)/2$ consecutive bonds of $C$, together with the final site of the last of these bonds, has $\ell^\infty$-diameter at least $P(C)/(2\Lambda(C))$.
\end{lemma}

\begin{proof}
By Lemma~\ref{lem:corners}(i) the number $P(C)/2$ is an integer, so that the set contains two sites $x$ and $y$ which are $P(C)/2$ bonds apart along the loop. Both arcs between them have length $P(C)/2$, hence $d_C(x, y) = P(C)/2$, and $\abs{x - y}_\infty \geq d_C(x, y)/\Lambda(C)$ by the definition of $\Lambda(C)$.
\end{proof}

\subsection{The Neighbourhood of a Corner}\label{sec:corner-neighbourhood}

The coefficients of the normalisation are defined on two model paths: the bi-infinite straight line, and the bent path which consists of two perpendicular rays. We now make precise the sense in which a loop coincides with one of them near each of its bonds. For a corner $c$ of $C$ let $S_{\mathrm{in}}(c)$ and $S_{\mathrm{out}}(c)$ be the segments which end and begin at $c$, and let the \emph{bent path} at $c$ be
\begin{equation}\label{eq:bent-path}
    B_c \coloneqq \{ c - tv_c \, | \, t \in \mathbb Z,\ t \geq 0 \} \cup \{ c + tw_c \, | \, t \in \mathbb Z,\ t \geq 0 \}\;,
\end{equation}
traversed from the first ray to the second. It contains the sites of $S_{\mathrm{in}}(c)$ and of $S_{\mathrm{out}}(c)$. We distribute the bonds of $C$ among its corners: of a segment of $\ell$ bonds, the first $\ceil{\ell/2}$ bonds are \emph{assigned} to its initial corner and the last $\floor{\ell/2}$ bonds to its final corner. Every bond of $C$ is assigned to exactly one corner.

\begin{lemma}[Corner Isolation]\label{lem:corner-isolation}
Let $C$ be a loop, $c$ a corner of $C$, and $b$ a bond of $C$ assigned to $c$, with position $x$. Every site of $C$ or of $B_c$ which lies neither on $S_{\mathrm{in}}(c)$ nor on $S_{\mathrm{out}}(c)$ is at $\ell^\infty$-distance at least $\min(s(C)/2, \delta(C)) \geq r(C)$ from $x$.
\end{lemma}

\begin{proof}
Let $\ell$ and $\ell'$ be the numbers of bonds of $S_{\mathrm{in}}(c)$ and $S_{\mathrm{out}}(c)$, let $c'$ be the initial corner of $S_{\mathrm{in}}(c)$ and $c''$ the final corner of $S_{\mathrm{out}}(c)$, and write $v \coloneqq v_c$ and $w \coloneqq w_c$. Since $C$ has at least four segments, the segment before $S_{\mathrm{in}}(c)$ is not adjacent to $S_{\mathrm{out}}(c)$, and the segment after $S_{\mathrm{out}}(c)$ is not adjacent to $S_{\mathrm{in}}(c)$.

Assume first that $b$ is one of the last $\floor{\ell/2}$ bonds of $S_{\mathrm{in}}(c)$, so that $x = c' + \alpha v$ with $\alpha \geq \ell - \floor{\ell/2} = \ceil{\ell/2}$. The segment before $S_{\mathrm{in}}(c)$ is perpendicular to $v$ and contains $c'$, so that all its sites have the same coordinate along $v$ as $c'$, which differs from that of $x$ by $\alpha \geq s(C)/2$. All other segments of $C$ apart from $S_{\mathrm{out}}(c)$ are not adjacent to $S_{\mathrm{in}}(c)$, and their sites are at distance at least $\delta(C)$ from $x$. The sites $c' - tv$, $t > 0$, of the first ray of $B_c$ beyond $c'$ differ from $x$ by $\alpha + t$ in the coordinate along $v$. The sites $c + tw$, $t > \ell'$, of the second ray beyond $c''$ differ from $x$ by $t > \ell' \geq s(C)$ in the coordinate along $w$, which is constant on $S_{\mathrm{in}}(c)$.

Assume now that $b$ is one of the first $\ceil{\ell'/2}$ bonds of $S_{\mathrm{out}}(c)$, so that $x = c + iw$ with $0 \leq i \leq \ceil{\ell'/2} - 1$. The segment after $S_{\mathrm{out}}(c)$ is perpendicular to $w$ and contains $c'' = c + \ell'w$, so that its sites differ from $x$ by $\ell' - i \geq \floor{\ell'/2} + 1 > s(C)/2$ in the coordinate along $w$, and the sites $c + tw$, $t > \ell'$, of $B_c$ differ from it by more. All other segments of $C$ apart from $S_{\mathrm{in}}(c)$ are not adjacent to $S_{\mathrm{out}}(c)$, which gives the distance $\delta(C)$. The sites $c - tv$, $t > \ell$, of the first ray beyond $c'$ differ from $x$ by $t > \ell \geq s(C)$ in the coordinate along $v$, which is constant on $S_{\mathrm{out}}(c)$.
\end{proof}

\subsection{A Summation over the Scales}\label{sec:scale-sums}

The lengths of a loop grow by the factor $L$ from each scale to the next finer one, and the errors of the corner analysis decay exponentially in these lengths. Their sum over the scales is controlled by the following elementary bound, which we record here because it is used for every such sum.

\begin{lemma}[Sums over Scales]\label{lem:scale-sums}
For $x > 0$, $c > 0$, and a real number $L \geq 2$,
\begin{equation}\label{eq:scale-sums}
    \sum_{m \geq 0}xL^me^{-cxL^m} \leq xe^{-cx}\Bigl(1 + \frac{2}{cx}\Bigr)
    \qquad\text{and}\qquad
    \sum_{m \geq 0}e^{-cxL^m} \leq e^{-cx}\Bigl(1 + \frac{2}{cx}\Bigr)\;.
\end{equation}
\end{lemma}

\begin{proof}
Put $y_m \coloneqq xL^m$. For $m \geq 1$ we have $y_m - y_{m-1} = y_m(1 - L^{-1}) \geq y_m/2$, and $e^{-ct} \geq e^{-cy_m}$ on the interval $[y_{m-1}, y_m]$, so that
\begin{equation}\label{eq:scale-sum-term}
    y_me^{-cy_m} \leq 2(y_m - y_{m-1})e^{-cy_m} \leq 2\int_{y_{m-1}}^{y_m}e^{-ct}\,\dd t\;.
\end{equation}
The intervals $[y_{m-1}, y_m]$, $m \geq 1$, are disjoint up to their end points and lie in $[x, \infty)$, so that the terms with $m \geq 1$ sum to at most $2\int_x^\infty e^{-ct}\,\dd t = (2/c)e^{-cx}$. Adding the term $xe^{-cx}$ of $m = 0$ gives the first bound. The second follows from the first, because $e^{-cy_m} \leq (y_m/x)e^{-cy_m}$ for $y_m \geq x$.
\end{proof}

\section{The Construction}\label{sec:construction}

This section defines the renormalisation group that the paper runs: Ba{\l}aban's steps, with every threshold, cube size, and depth fixed by the target $(K, \bar g)$, and with continuous plaquette test functions. It shows that Ba{\l}aban's estimates hold for this construction, and it fixes the stopping index of an ensemble.

\subsection{The Tube of a Target}\label{sec:tube}

Ba{\l}aban's construction is governed by a small number of parameters, each of which he defines as a function of the running coupling of its scale. The plaquette threshold of the scale $j$ is
\begin{equation}\label{eq:balaban-threshold}
    \varepsilon_j = g_jp_0(g_j)\;,
    \qquad
    p_0(g) \coloneqq A_0(\log g^{-2})^{p_0}\;,
\end{equation}
and the size of the cubes of that scale is $MR_j$, where $R_j$ is the smallest power of $L$ with $R_j \geq (\log g_j^{-2})^r$ \cite[(2.4), (2.5), p.~255]{BalabanCRE}. Here $A_0$ is a large constant, $r \geq 2$ and $p_0 \geq 5r$ are exponents \cite[pp.~245--246]{BalabanCRE}, and we keep his notation, in which $p_0$ denotes both the function and its exponent. The running couplings, in turn, are produced by the construction: $g_{j+1}$ is computed from the effective action which the first $j$ steps have generated. As long as one studies a given bare coupling this causes no difficulty. We, however, have to find the bare coupling for which $g_K = \bar g$, and for this we need the map $g_0 \mapsto g_K$ to be continuous, which is not clear when a cube size, an integer, is a function of $g_j$.

Ba{\l}aban mentions a way out. In the first step he notes that there are ``two convenient choices'' of the threshold, the second of which replaces the coupling inside the logarithm by an upper bound which is known in advance, and that this choice is ``more convenient in some nonperturbative aspects of the procedure, but it is technically more difficult'', and he assumes the first \cite[p.~246]{BalabanCRE}. We take this second route, in a stronger form: every parameter of the scale $j$ is evaluated at a number $\hat g_j$ which is fixed by the target and which the running coupling $g_j$ is shown, afterwards, to follow within a constant factor.

Throughout, $\beta_0 \in (0, \sqrt2 - 1]$ is a fixed number, the constant of Ba{\l}aban's inequalities (2.6), and $0 < b_* \leq b^*$ are fixed numbers, which Section~\ref{sec:coupling-flow} takes to be multiples of the first coefficient of the $\beta$-function.

\begin{definition}[Tubes]\label{def:tube}
Let $\gamma > 0$. A \emph{tube} consists of an integer $K \geq 1$ and real numbers $b$, $b'$, $C_W$, and $\bar g > 0$ such that, with $t \coloneqq \bar g^{-2}$,
\begin{equation}\label{eq:tube-class}
    b_* \leq b \leq b' \leq \min\bigl(\tfrac{3}{2}b,\, b^*\bigr)\;,
    \qquad
    C_W \geq 0\;,
    \qquad
    t \geq \max\bigl(\gamma^{-2} + C_W,\, 4C_W/\beta_0\bigr)\;.
\end{equation}
For every integer $k \leq K$ we put
\begin{equation}\label{eq:tube-couplings}
    \ell_k \coloneqq t + b(K - k) - C_W\;,
    \qquad
    u_k \coloneqq t + b'(K - k) + C_W\;,
    \qquad
    \hat g_k \coloneqq u_k^{-1/2}\;,
\end{equation}
and for $0 \leq k \leq K$ we let $H_k$ be the set of the $(x_0, \dots, x_k) \in \mathbb R^{k+1}$ with $\ell_j \leq x_j \leq u_j$ for $0 \leq j \leq k$ and
\begin{equation}\label{eq:tube-window}
    b(n - m) - C_W \leq x_m - x_n \leq b'(n - m) + C_W
\end{equation}
for $0 \leq m < n \leq k$. A point of $H_k$ is a \emph{history} of the tube, with the couplings $g_j \coloneqq x_j^{-1/2}$. The \emph{ratio} of the tube is
\begin{equation}\label{eq:tube-ratio}
    \rho \coloneqq \max\Bigl(\frac{b'}{b},\, \frac{t + C_W}{t - 2C_W}\Bigr)\;.
\end{equation}
\end{definition}

The set $H_k$ is closed and convex, being defined by finitely many weak linear inequalities. Since $\beta_0 < 1/2$, the last condition of \eqref{eq:tube-class} gives $C_W \leq \beta_0t/4 < t/8$, so that $(t + C_W)/(t - 2C_W) < (9/8)/(6/8)$ and therefore
\begin{equation}\label{eq:tube-ratio-bound}
    1 \leq \rho \leq \tfrac{3}{2}\;.
\end{equation}
The numbers $\hat g_k$ increase with $k$, and $\hat g_K \leq \bar g$. The role of $x_j$ is that of the inverse square of the running coupling of the scale $j$, and the tube is the set in which we shall confine it: \eqref{eq:tube-window} says that the inverse square of the coupling decreases, from each scale to a coarser one, by an amount between $b$ and $b'$ per step, up to the additive constant $C_W$.

The lemmata of this subsection require $\gamma$ to be small. With $p^* \geq 1$ an upper bound for the exponents $r$, $p_0$, $p_1$, $q_0$, $q_1$ of the construction, the conditions are
\begin{equation}\label{eq:gamma-conditions}
\begin{gathered}
    b^*\gamma^2 \leq \frac{\beta_0}{2}\;,
    \qquad
    \log\gamma^{-2} \geq \frac{p^*}{\beta_0}\;,
    \qquad
    \Bigl(1 + \frac{\log(3/2)}{\log\gamma^{-2}}\Bigr)^{p^*} \leq 2\;,\\
    \Bigl(1 + \frac{b^*\gamma^2}{\log\gamma^{-2}}\Bigr)^{r} < L\;,
    \qquad
    \Bigl(1 + \frac{\log(1 + b^*)}{\log\gamma^{-2}}\Bigr)^{r} \leq 1 + \frac{1}{L}\;.
\end{gathered}
\end{equation}
Each of them holds for all sufficiently small $\gamma$, once $\beta_0$, $b^*$, $p^*$, $r$, and $L$ are fixed.

\begin{lemma}[Couplings of a History]\label{lem:tube-couplings}
Let $\gamma$ satisfy \eqref{eq:gamma-conditions}, and let $(x_0, \dots, x_k)$ be a history of a tube, with the couplings $g_j = x_j^{-1/2}$. Put $\beta' \coloneqq b' + C_W$.
\begin{enumerate}[label=(\roman*)]
\item For $0 \leq j \leq k$,
\begin{equation}\label{eq:tube-coupling-ratio}
    \hat g_j \leq g_j \leq \rho^{1/2}\hat g_j\;,
    \qquad
    g_j \leq \gamma\;,
    \qquad
    \log g_j^{-2} \leq \log\hat g_j^{-2} \leq \log g_j^{-2} + \log\rho\;.
\end{equation}
\item For $0 \leq m < n \leq k$,
\begin{equation}\label{eq:tube-two-six}
    g_n \leq \bigl(1 + g_n^2\beta'(n - m)\bigr)^{1/2}g_m \leq (1 + \beta_0)(n - m)^{1/2}g_m
    \qquad\text{and}\qquad
    g_m \leq (1 + \beta_0)g_n\;,
\end{equation}
and, for every real $p \in (0, p^*]$,
\begin{equation}\label{eq:tube-two-seven}
    (\log g_n^{-2})^p \leq (1 + \beta_0)(\log g_m^{-2})^p
    \qquad\text{and}\qquad
    (\log g_m^{-2})^p \leq \bigl(1 + g_n^2\beta'(n - m)\bigr)^{\beta_0}(\log g_n^{-2})^p\;.
\end{equation}
\item The inequalities \eqref{eq:tube-two-six} and \eqref{eq:tube-two-seven} hold for the numbers $\hat g_j$ in place of the $g_j$, with $b'$ in place of $\beta'$ and for all integers $m < n \leq K$, and moreover $\hat g_m \leq \hat g_n$.
\end{enumerate}
\end{lemma}

The inequalities \eqref{eq:tube-two-six} and \eqref{eq:tube-two-seven} are Ba{\l}aban's (2.6) and (2.7) \cite[p.~255]{BalabanCRE}, which he states for the couplings of his construction and for positive integers $p$. The lemma thus says that a history of a tube, and also the sequence $(\hat g_j)$ itself, satisfy the hypotheses which he places on a sequence of couplings.

\begin{proof}
(i) Since $x_j \leq u_j$ we have $\hat g_j \leq g_j$. With $s \coloneqq K - j \geq 0$,
\begin{equation}\label{eq:tube-mediant}
    \frac{u_j}{x_j} \leq \frac{u_j}{\ell_j} = \frac{(t + C_W) + b's}{(t - C_W) + bs} \leq \max\Bigl(\frac{t + C_W}{t - C_W},\, \frac{b'}{b}\Bigr) \leq \rho\;,
\end{equation}
where the first step uses $x_j \geq \ell_j > 0$, the second that a quotient $(A + A')/(B + B')$ of sums of positive numbers lies between $A/B$ and $A'/B'$, and the third the definition \eqref{eq:tube-ratio}. Hence $g_j^2 = x_j^{-1} \leq \rho u_j^{-1} = \rho\hat g_j^2$. Moreover $x_j \geq \ell_j \geq t - C_W \geq \gamma^{-2}$ by \eqref{eq:tube-class}, so that $g_j \leq \gamma$, and the last chain of \eqref{eq:tube-coupling-ratio} is the logarithm of the first.

(ii) The upper bound of \eqref{eq:tube-window} and $n - m \geq 1$ give $x_m \leq x_n + \beta'(n - m)$, that is, $g_m^{-2} \leq g_n^{-2}(1 + g_n^2\beta'(n - m))$, which is the first inequality of \eqref{eq:tube-two-six}. For the second we use $1 + y(n - m) \leq (1 + y)(n - m)$ for $y \geq 0$, and bound $y = g_n^2\beta'$: by (i) and \eqref{eq:tube-class}, $g_n^2 \leq (t - C_W)^{-1}$ with $C_W \leq \beta_0t/4$ and $t \geq \gamma^{-2}$, so that
\begin{equation}\label{eq:tube-beta-prime}
    g_n^2\beta' \leq \frac{b' + C_W}{t - C_W} \leq \frac{b^*\gamma^2 + \beta_0/4}{1 - \beta_0/4} \leq \frac{3\beta_0/4}{1 - \beta_0/4} \leq \beta_0\;,
\end{equation}
by the first condition of \eqref{eq:gamma-conditions} and $\beta_0 \leq 1$. Thus $(1 + g_n^2\beta')^{1/2} \leq 1 + \beta_0$. For the third inequality, the lower bound of \eqref{eq:tube-window} gives $x_n \leq x_m + C_W$, and $C_W \leq \beta_0t/4 \leq (\beta_0/3)x_m$, because $x_m \geq t - C_W \geq (1 - \beta_0/4)t \geq 3t/4$. Hence $g_m^2 \leq (1 + \beta_0/3)g_n^2$.

We turn to \eqref{eq:tube-two-seven}. Put $\lambda_j \coloneqq \log g_j^{-2}$, which is at least $\log\gamma^{-2} \geq p/\beta_0 \geq 2p$ by (i), the second condition of \eqref{eq:gamma-conditions}, and $\beta_0 \leq 1/2$. From $g_m \leq (1 + \beta_0)g_n$ we get $\lambda_n \leq \lambda_m + 2\log(1 + \beta_0)$, and therefore
\begin{equation}\label{eq:tube-log-first}
    \Bigl(\frac{\lambda_n}{\lambda_m}\Bigr)^p \leq \Bigl(1 + \frac{2\log(1 + \beta_0)}{\lambda_m}\Bigr)^p \leq (1 + \beta_0)^{2p/\lambda_m} \leq 1 + \beta_0\;,
\end{equation}
using $1 + x \leq e^x$ and $2p \leq \lambda_m$. In the same way the first inequality of \eqref{eq:tube-two-six} gives $\lambda_m \leq \lambda_n + \log(1 + y)$ with $y \coloneqq g_n^2\beta'(n - m)$, so that $(\lambda_m/\lambda_n)^p \leq (1 + y)^{p/\lambda_n} \leq (1 + y)^{\beta_0}$, because $p/\lambda_n \leq \beta_0$.

(iii) Since $u_m = u_n + b'(n - m)$, the numbers $\hat g_j$ are increasing and $\hat g_m^{-2} = \hat g_n^{-2}(1 + \hat g_n^2b'(n - m))$, which is the first inequality of \eqref{eq:tube-two-six} with equality. The remaining inequalities follow as in (ii), with $\hat g_n^2b' \leq b^*\gamma^2 \leq \beta_0$ in place of \eqref{eq:tube-beta-prime} and $\log\hat g_j^{-2} \geq \log\gamma^{-2}$, which holds because $u_j \geq t \geq \gamma^{-2}$.
\end{proof}

We can now define the parameters of the construction for a tube. They are Ba{\l}aban's functions of the coupling, evaluated at $\hat g_j$.

\begin{definition}[Parameters of a Tube]\label{def:tube-parameters}
For a tube and every integer $j \leq K$:
\begin{enumerate}[label=(\roman*)]
\item $R_j$ is the smallest power of $L$ with $R_j \geq (\log\hat g_j^{-2})^r$;
\item $\varepsilon_j \coloneqq \hat g_jp_0(\hat g_j)$, with $p_0(g) = A_0(\log g^{-2})^{p_0}$, and $\vartheta_j \coloneqq (A_1/A_0)\varepsilon_j$;
\item $\delta'_j \coloneqq \hat g_jA_1(\log\hat g_j^{-2})^{p_1}$;
\item $\alpha_{0,j} \coloneqq \hat g_jC^\alpha_0(\log\hat g_j^{-2})^{q_0}$ and $\alpha_{1,j} \coloneqq \hat g_jC^\alpha_1(\log\hat g_j^{-2})^{q_1}$.
\end{enumerate}
\end{definition}

The printed definitions are those of \eqref{eq:balaban-threshold} for $\varepsilon_j$ and $R_j$; $\delta_k = g_k(A_1/A_0)p_0(g_k)$ for the threshold of the fluctuation field, which we write $\vartheta_k$ \cite[p.~265, after (3.4)]{BalabanCRE}; $\delta'_j = g_jA_1p_1(g_j)$ with $p_1(g) = (\log g^{-2})^{p_1}$ and $p_1 < p_0$ for the threshold of the $\mathbf R$-operation \cite[p.~183]{BalabanLFR1}; and (2.28) for the radii $\alpha_{0,j}$, $\alpha_{1,j}$ of the domains of analyticity, with integers $q_0, q_1 > 1$ and large constants which he writes $C_0$, $C_1$ \cite[p.~259]{BalabanCRE}. A parameter of the scale $j$ of a tube depends on $K$ and $j$ only through the difference $K - j$, since $u_j$ does. This will be essential in Section~\ref{sec:cauchy}: the runs for the targets $(K, \bar g)$ and $(K + 1, \bar g)$ have the same parameters on each of the lattices which they share.

\begin{lemma}[Cube Sizes of a Tube]\label{lem:tube-cubes}
Let $\gamma$ satisfy \eqref{eq:gamma-conditions}. For every tube and every history $(x_0, \dots, x_k)$ of it:
\begin{enumerate}[label=(\roman*)]
\item $(\log g_j^{-2})^r \leq R_j < L(\log\hat g_j^{-2})^r \leq 2L(\log g_j^{-2})^r$ for $0 \leq j \leq k$;
\item $R_{j+1} \in \{ R_j,\, L^{-1}R_j \}$ for every integer $j < K$, so that $R_n \leq R_m$ for $m < n$;
\item for all integers $m < n \leq K$,
\begin{equation}\label{eq:tube-two-nine}
    R_m \leq L\bigl(1 + \hat g_n^2b'(n - m)\bigr)^{\beta_0}R_n
    \qquad\text{and}\qquad
    R_m \leq (L + 1)(n - m)^{\beta_0}R_n\;.
\end{equation}
\end{enumerate}
\end{lemma}

Part (i) is Ba{\l}aban's (2.5), with a second inequality in which the factor $2L$ stands where his definition would give $L$, and (ii) and (iii) contain his (2.9) \cite[pp.~255--256]{BalabanCRE}, since $\hat g_n^2b' \leq g_n^2\beta'$.

\begin{proof}
(i) The lower bound is the definition of $R_j$ and $\hat g_j \leq g_j$. Since $\log\hat g_j^{-2} \geq \log\gamma^{-2} > 1$, the smallest power of $L$ which is at least $(\log\hat g_j^{-2})^r$ is less than $L$ times this number. Finally $\log\hat g_j^{-2} \leq \log g_j^{-2} + \log\rho$ by Lemma~\ref{lem:tube-couplings}(i), with $\rho \leq 3/2$ and $\log g_j^{-2} \geq \log\gamma^{-2}$, so that the third condition of \eqref{eq:gamma-conditions} gives $(\log\hat g_j^{-2})^r \leq 2(\log g_j^{-2})^r$.

(ii) The numbers $u_j$ decrease, hence $R_{j+1} \leq R_j$. Moreover $u_j = u_{j+1} + b' \leq (1 + b^*\gamma^2)u_{j+1}$, because $u_{j+1} \geq \gamma^{-2}$, so that
\begin{equation}\label{eq:tube-cube-ratio}
    R_j < L(\log u_j)^r \leq L\Bigl(1 + \frac{\log(1 + b^*\gamma^2)}{\log u_{j+1}}\Bigr)^r(\log u_{j+1})^r < L^2R_{j+1}\;,
\end{equation}
by the fourth condition of \eqref{eq:gamma-conditions} and $(\log u_{j+1})^r \leq R_{j+1}$. The quotient $R_j/R_{j+1}$ is a power of $L$ in $[1, L^2)$, so it is $1$ or $L$.

(iii) Put $s \coloneqq n - m$ and $\lambda \coloneqq \log\gamma^{-2}$. As in (i), $R_m < L(\log u_m)^r$ and $(\log u_n)^r \leq R_n$, and $u_m = u_n(1 + \hat g_n^2b's)$, so that
\begin{equation}\label{eq:tube-cube-quotient}
    \frac{R_m}{R_n} < L\Bigl(1 + \frac{\log(1 + \hat g_n^2b's)}{\log u_n}\Bigr)^r\;.
\end{equation}
For the first inequality we bound the bracket by $(1 + \hat g_n^2b's)^{r/\log u_n}$ and use $r/\log u_n \leq r/\lambda \leq \beta_0$. For the second we use $\hat g_n^2 \leq \gamma^2$, $b' \leq b^*$, and $\log u_n \geq \lambda$, which bound the bracket by $(1 + \log(1 + b^*s\gamma^2)/\lambda)^r$, and we distinguish two cases. If $s \leq \gamma^{-2}$, this is at most $(1 + \log(1 + b^*)/\lambda)^r \leq 1 + 1/L$, by the last condition of \eqref{eq:gamma-conditions}, and $1 \leq s^{\beta_0}$. If $s > \gamma^{-2}$, then $1 + b^*s\gamma^2 \leq (1 + b^*)s\gamma^2$, and with $y \coloneqq \log s > \lambda$ the bracket is at most $\varphi(y) \coloneqq ((y + \log(1 + b^*))/\lambda)^r$. We have $\varphi(\lambda) \leq 1 + 1/L$, and the logarithmic derivative of $\varphi$ is $r/(y + \log(1 + b^*)) \leq r/\lambda \leq \beta_0$, so that $\varphi(y) \leq (1 + 1/L)e^{\beta_0(y - \lambda)} \leq (1 + 1/L)s^{\beta_0}$. In both cases $R_m/R_n < L(1 + 1/L)s^{\beta_0}$.
\end{proof}

\begin{lemma}[Thresholds of a Tube]\label{lem:tube-thresholds}
Let $\gamma$ satisfy \eqref{eq:gamma-conditions}, and let $\tau_j$ be one of the parameters $\varepsilon_j$, $\vartheta_j$, $\delta'_j$, $\alpha_{0,j}$, $\alpha_{1,j}$ of a tube, so that $\tau_j = \tau(\hat g_j)$ with $\tau(g) = cg(\log g^{-2})^p$ for a constant $c > 0$ and an exponent $p \in (0, p^*]$.
\begin{enumerate}[label=(\roman*)]
\item For every history of the tube and $0 \leq j \leq k$,
\begin{equation}\label{eq:tube-threshold-ratio}
    \rho^{-1/2}\tau(g_j) \leq \tau_j \leq \tau(g_j)\;.
\end{equation}
\item For all integers $m < n \leq K$,
\begin{equation}\label{eq:tube-two-eight}
    \tau_n \leq (1 + \beta_0)(n - m)^{1/2}\tau_m
    \qquad\text{and}\qquad
    \tau_m \leq \bigl(1 + \hat g_n^2b'(n - m)\bigr)^{\beta_0}\tau_n \leq (1 + \beta_0)(n - m)^{\beta_0}\tau_n\;.
\end{equation}
\end{enumerate}
\end{lemma}

For $\tau = \varepsilon$ the inequalities \eqref{eq:tube-two-eight} are Ba{\l}aban's (2.8) \cite[p.~256]{BalabanCRE}. The upper bound of \eqref{eq:tube-threshold-ratio} says that the thresholds of a tube are the smaller ones: a statement which uses a threshold only as an upper bound on a field which has passed a test function holds for the tube with the printed value. The lower bound controls the loss in the statements which use a threshold from below.

\begin{proof}
(i) The function $\tau$ is increasing on $(0, \gamma]$, since its logarithmic derivative is $g^{-1}(1 - 2p/\log g^{-2})$ and $\log g^{-2} \geq \log\gamma^{-2} \geq p^*/\beta_0 > 2p$. As $\hat g_j \leq g_j \leq \gamma$ by Lemma~\ref{lem:tube-couplings}(i), this gives the upper bound. For the lower bound we write $\tau_j = (\hat g_j/g_j)\,cg_j(\log\hat g_j^{-2})^p$ and use $\hat g_j \geq \rho^{-1/2}g_j$ and $\log\hat g_j^{-2} \geq \log g_j^{-2}$ from the same lemma.

(ii) We have $\tau_n/\tau_m = (\hat g_n/\hat g_m)(\log u_n/\log u_m)^p \leq \hat g_n/\hat g_m$, because $u_n \leq u_m$, and $\hat g_n/\hat g_m \leq (1 + \beta_0)(n - m)^{1/2}$ by Lemma~\ref{lem:tube-couplings}(iii). In the other direction $\tau_m/\tau_n = (\hat g_m/\hat g_n)(\log u_m/\log u_n)^p \leq (\log u_m/\log u_n)^p$, which is at most $(1 + \hat g_n^2b'(n - m))^{\beta_0}$ by the second inequality of \eqref{eq:tube-two-seven} for the sequence $(\hat g_j)$. The last step is $1 + y(n - m) \leq (1 + y)(n - m)$ with $y = \hat g_n^2b' \leq \beta_0$ and $(1 + \beta_0)^{\beta_0} \leq 1 + \beta_0$.
\end{proof}

\subsection{The Density after \texorpdfstring{$k$}{k} Steps}\label{sec:density}

We recall the form which Ba{\l}aban gives to the density of the block field after $k$ steps, since every later section works inside it. The description is his inductive assumption \cite[Sect.~2]{BalabanCRE}, and we cite each of its parts by page and equation. Nothing in this subsection is new.

At the step $k$ we use his units, in which the lattice $T^{(k)}$ of the block field $V_k$ has the spacing one and the fine lattice $T_\eta$, which is our $\mathbb T_K$, has the spacing $\eta = L^{-k}$. A \emph{scale} $j \leq k$ refers to the lattice of spacing $L^{-(k-j)}$ in these units, the lattice on which the block field of the step $j$ lives. For the localisation one fixes $M = L^m$ and, at each scale $j$, the family $\pi_j$ of the closed cubes of side $M$ of that scale. A \emph{localisation domain} of the scale $j$ is the union of a finite family of cubes of $\pi_j$ which is connected, in the sense that any two of its cubes are joined by a chain of cubes of the family in which consecutive cubes have a common wall. The class of these domains is $\mathbf D_j$, and the \emph{linear size} $d_j(X)$ of $X \in \mathbf D_j$ is the length of a shortest tree graph in $X$ which meets all the cubes of $X$, divided by $M$ \cite[p.~257]{BalabanRG1}.

\emph{The terms of the expansion.} A term of the density is labelled by three sequences of sets. The first two are sequences of localisation domains $\Omega_j, \Lambda_j \in \mathbf D_j$ with
\begin{equation}\label{eq:domain-sequences}
    \Omega_1 \supset \Lambda_1 \supset \Omega_2 \supset \Lambda_2 \supset \dots \supset \Omega_k \supset \Lambda_k\;,
\end{equation}
and the third consists of sets $S_j \subset \Omega_j \cap \Lambda_j^c$, each of which is empty or a union of cubes of side $LM_2R_j$ of the scale $j$ \cite[(2.1), p.~254]{BalabanCRE}. Roughly speaking, $\Omega_j$ is the region in which the field of the scale $j$ has passed the test functions of the step $j$, so that the step could be performed there, and $\Lambda_j$ is the smaller region in which the fluctuation field of that step was small as well. The complements
\begin{equation}\label{eq:large-field-regions}
    Z_j \coloneqq \Lambda_j^c
\end{equation}
are the \emph{large-field regions} \cite[(2.3), p.~255]{BalabanCRE}. The sequence $\{ \Omega_j \}$ determines on which lattice the field is still an integration variable at each point: with
\begin{equation}\label{eq:determining-set}
    \Gamma_0 \coloneqq \Omega_1^c\;,
    \qquad
    \Gamma_j \coloneqq \Omega_j^{(j)} \setminus \Omega_{j+1}^{(j)} \quad (1 \leq j < k)\;,
    \qquad
    \Gamma_k \coloneqq \Omega_k^{(k)}\;,
    \qquad
    \mathbf B \coloneqq \bigcup_{j=0}^{k}\Gamma_j\;,
\end{equation}
where the superscript $(j)$ denotes the points, bonds, and plaquettes of the scale $j$ in a set, the \emph{determining set} $\mathbf B$ carries the field $V$ which equals the block field $V_j$ of the step $j$ on $\Gamma_j$ \cite[(2.2), (2.10), pp.~255--256]{BalabanCRE}. The domains $\Omega_j$ and $\Lambda_j$ which are determined by the step $j$, and not by an $\mathbf R$-operation, are unions of cubes of side $MR_j$ of the scale $j$, and the distance between their boundaries is at least $2MR_j$ \cite[p.~256]{BalabanCRE}.

\emph{The background field.} Let $M_{\mathbf B}$ be the averaging operation which maps a gauge field $U$ on $T_\eta$ to the field on $\mathbf B$ equal to the $j$-fold block average $M^j(U)$ on $\Gamma_j$. A regular field $V$ on $\mathbf B$ determines the set of the minima of the action $U \mapsto A^\eta(U)$ on the regular fields $U$ with $M_{\mathbf B}(U) = V$, which is a single gauge orbit, and a minimal configuration is denoted by $U_k(V)$, or by $U_k$ \cite[(2.11), (2.12), p.~256]{BalabanCRE}. The existence, the uniqueness up to gauge transformations, and the regularity of $U_k(V)$ are the subject of \cite{BalabanVAR}. Every term of the effective action depends on $V$ through this \emph{background field}.

\emph{The expansion.} For a cube $\square$ of the partition of $T_\eta$ into cubes of side $LM_2R_k$ let $U_{k,\square}(V_k)$ be the minimal configuration of \cite[(2.16), p.~257]{BalabanCRE}, which is determined by the block field in a neighbourhood of $\square$. The test function of the step $k$ on $\Omega_k$ is the product
\begin{equation}\label{eq:test-function}
    \chi_k(\Omega_k) \coloneqq \prod_{\square \subset \Omega_k}\chi\Bigl(\Bigl\{ \sup_{p \subset \square^{\sim}}\abs{U_{k,\square}(V_k, \partial p) - 1} < \varepsilon_k\eta^2 \Bigr\}\Bigr)
\end{equation}
of the characteristic functions of the sets on which the plaquettes of these local minimisers are below the threshold $\varepsilon_k$, in the units of the fine lattice, and the density is
\begin{equation}\label{eq:density-expansion}
    \varrho_k(V_k) = \sum_{\{ \Omega_j \}, \{ \Lambda_j \}}\chi_k(\Omega_k)\,\mathbf T_k(\{ \Omega_j \}, \{ \Lambda_j \})\exp A_k\Bigl(\frac{1}{g_k^2}, U_k\Bigr)
\end{equation}
\cite[(2.17), (2.18), p.~257]{BalabanCRE}. The sum runs over the admissible sequences of domains, the sum over the sequences $\{ S_j \}$ is included in the operation $\mathbf T_k$, and the effective action $A_k$ depends on all three sequences.

\emph{The operations on the large-field regions.} The operation $\mathbf T_k$ consists of the integrations which the earlier steps have left: over the block fields $V_j$ on $\Omega_{j+1}^c$, and over the fluctuation fields $A_j$ on $Z_{j+1} \cap \Omega_{j+1}$. It factorises over the connected components $X_1, \dots, X_n$ of $Z_k$, and each factor is an ordered product over the scales,
\begin{equation}\label{eq:large-field-operations}
    \mathbf T_k(Z_k) = \prod_{i=1}^{n}\mathbf T_k(X_i)\;,
    \qquad
    \mathbf T_k(X) = \prod_{j=k-1}^{0}\mathbf T^{(j)}(Z_{j+1} \cap X)\;,
\end{equation}
where $\mathbf T^{(j)}$ comes from the step $j + 1$ and, if no $\mathbf R$-operation has changed it, has the form of \cite[(2.21), p.~258]{BalabanCRE}: the integral over $V_j$ against the $\delta$-function of the block averaging and the function $\zeta(\Omega_{j+1}^c)$, which contains the characteristic functions of the large fields, and the integral over $A_j$ against a characteristic function and a Gaussian weight \cite[(2.19), (2.20), p.~258]{BalabanCRE}.

\emph{The effective action.} The effective action has the form
\begin{equation}\label{eq:effective-action}
    A_k\Bigl(\frac{1}{g_k^2}, U_k\Bigr) = -A\Bigl(\frac{1}{g_k^2}, U_k\Bigr) + \mathbf E_k(U_k) + \mathbf R_k(U_k) + \mathbf B_k(U_k, A) - E_k
\end{equation}
\cite[(2.23), p.~258]{BalabanCRE}, in which $A = \{ A_{j-1} \}$ stands for the fluctuation fields in the last but one term. Its five terms are the following.
\begin{enumerate}[label=(\alph*)]
\item The first is the Wilson action of the background field with a coupling which depends on the position, defined by
\begin{equation}\label{eq:position-coupling}
    \frac{1}{g_{j-1}^2(x)} = \frac{1}{g_j^2(x)} + \beta_j(g_{j-1})\phi_j(x)\;,
\end{equation}
where $\phi_j \in C_0^\infty(\Lambda_j)$ equals one on $\Lambda_j^{\sim-1}$, the domain obtained from $\Lambda_j$ by removing a layer of cubes, so that $g_j(x) = g_j$ there \cite[(2.24), p.~259]{BalabanCRE}. The numbers $\beta_j(g_{j-1})$ are the coefficients of the coupling flow, the subject of Section~\ref{sec:coupling-flow}.
\item The \emph{regular part} is
\begin{equation}\label{eq:regular-part}
    \mathbf E_k(U_k) = \sum_{j=1}^{k}\bigl[\mathbf E^{(j)}(\Lambda_j, g_{j-1}, U_k) - \mathbf E^{(j)}(\Lambda_j, g_{j-1}, 1) - \beta_j(g_{j-1})A(\phi_j, U_k)\bigr]\;,
\end{equation}
with $\mathbf E^{(j)}(\Lambda_j, U_k) = \sum_{z \in \Lambda_j^0}\sum_{X \ni z}\mathbf E^{(j)}(X, U_k, z)$, the inner sum over the $X \in \mathbf D_j$ contained in $\Lambda_j$ and $\Lambda_j^0$ obtained from $\Lambda_j^{(j)}$ by removing one layer of the cubes of side $MR_j$ \cite[(2.25)--(2.27), p.~259]{BalabanCRE}. Each term $\mathbf E^{(j)}(X, U_k, z)$ depends on $U_k$ restricted to $X$; it extends to an analytic function of a pair $(\mathbf U, \mathbf J)$ in the space $U_j^c(X, \alpha_{0,j}, \alpha_{1,j})$; the extension is invariant under gauge transformations; and it satisfies the bound
\begin{equation}\label{eq:regular-bound}
    \abs{\mathbf E^{(j)}(X, (\mathbf U, \mathbf J), z)} \leq E_0\exp(-\kappa d_j(X))
\end{equation}
of \cite[(1.18), p.~263]{BalabanRG1}. The space $U_j^c(X, \alpha_0, \alpha_1)$ consists of the gauge orbits of the pairs of a $G^c$-valued field $\mathbf U = U'U$ and a $\mathfrak g^c$-valued field $\mathbf J$ on $X$ in which $U$ is a real field with plaquettes below $\alpha_0$, $U' = \exp i\xi A'$ with $A'$ and its covariant derivative below $\alpha_1$, and $\mathbf J$ is below $\alpha_0$, together with the corresponding bounds for the minimisers built from the block averages of $\mathbf U$ \cite[Definition, (1.10)--(1.16), p.~262]{BalabanRG1}. It is a complex neighbourhood of the real regular fields on $X$, of a width set by $\alpha_1$.
\item The \emph{$\mathbf R$-part} is $\mathbf R_k(U_k) = \sum_{j=1}^{k}[\mathbf R^{(j)}(\Lambda_j, U_k) - \mathbf R^{(j)}(\Lambda_j, 1)]$, where $\mathbf R^{(j)}(\Lambda_j, U_k)$ is a sum of terms $\mathbf R^{(j)}(X, U_k)$ over the domains $X \in \mathbf D_j$ contained in $\Lambda_j^{\sim-1}$ which are unions of cubes of side $MR_j$. These terms have the first three properties of the terms of the regular part and satisfy
\begin{equation}\label{eq:r-bound}
    \abs{\mathbf R^{(j)}(X, (\mathbf U, \mathbf J))} \leq g_j^{\kappa_0}\exp(-\kappa d_j(X))
\end{equation}
\cite[(2.30), (2.31), p.~260]{BalabanCRE}. They are produced by the $\mathbf R$-operation from terms with large fields, and they are so small that they need no counterterm.
\item The \emph{boundary part} is $\mathbf B_k = \sum_{j=1}^{k}\mathbf B^{(j)}(U_k, A, \{ S_i \})$, where $\mathbf B^{(j)}$ is a sum of terms $\mathbf B^{(j)}(X, U_k, A, \{ S_i \cap X \})$ over the $X \in \mathbf D_j$ with $X \cap \Omega_j \neq \emptyset$ and $X \cap Z_j^{\sim} \neq \emptyset$. Each of them depends on $U_k$ and $A$ restricted to $X$, extends analytically to the space $\tilde U_j^c(X, \tilde\alpha_0, \tilde\alpha_1)$, is invariant under $G$-valued gauge transformations, and satisfies
\begin{equation}\label{eq:boundary-bound}
    \abs{\mathbf B^{(j)}(X, (\mathbf U, \mathbf J), A, \{ S_i \cap X \})} < B_0\exp(-\kappa d_j(X))
\end{equation}
\cite[(2.40)--(2.42), p.~261]{BalabanCRE}. The space $\tilde U_j^c(X, \tilde\alpha_0, \tilde\alpha_1)$ is the version of $U_j^c$ for a domain which meets several of the layers $\Omega_n \setminus \Omega_{n+1}$: its conditions (2.34)--(2.39) are those of $U_j^c$ with the radii $\alpha_{0,n}$, $\alpha_{1,n}$ of the scale $n$ on the layer $n$, each multiplied by $1 - \beta(1 - 2^{-(j-n)})$ with a constant $\beta$, for which Ba{\l}aban suggests $\beta = 1/4$ \cite[p.~261]{BalabanCRE}.
\item $E_k$ is a constant, which depends on $\{ \Omega_j \}$ and $\{ \Lambda_j \}$ \cite[p.~262]{BalabanCRE}.
\end{enumerate}
The terms which a step has just created are defined on slightly larger spaces and decay faster, with $(1 + 4\beta)\kappa$ in place of $\kappa$, and this margin is consumed by the $\mathbf R$-operation \cite[p.~262]{BalabanCRE}.

\emph{Ba{\l}aban's theorems.} The main result of \cite{BalabanCRE} is that the form just described is reproduced by a step.

\begin{theorem}[Ba{\l}aban {\cite[Theorem~1, p.~262]{BalabanCRE}}]\label{thm:balaban-inductive}
There exist constants, introduced in the above description, such that if the sequence of coupling constants $\{ g_k \}$, determined by the recursive renormalisation group equations (0.18), (0.20) of \cite{BalabanRG1}, satisfies the inequalities (0.31) of \cite{BalabanRG1}, then the sequence of densities $\{ \varrho_k \}$, generated by successive applications of the operations $\mathbf RT$ to the density $\varrho_0 = \exp[-(1/g_0^2)A - E]$, satisfies all the inductive assumptions.
\end{theorem}

The proof of this theorem occupies \cite{BalabanCRE}, for the step $T$, and \cite{BalabanLFR1,BalabanLFR2}, for the $\mathbf R$-operation. Two consequences are used throughout. The first is a bound on the regular and the $\mathbf R$-parts in a region $\Omega$, \cite[Theorem~2, (2.43), (2.44), p.~263]{BalabanCRE}: with a constant $E_1$ independent of $j$, $k$, $\Omega$, the sequences of domains, and the torus,
\begin{equation}\label{eq:balaban-theorem-two}
    \Bigl\lvert\sum_{z \in \Lambda_j^0 \cap \Omega}\bigl[\mathbf E^{(j)}(\Lambda_j, U_k, z) - \mathbf E^{(j)}(\Lambda_j, 1, z)\bigr] - \beta_j(g_{j-1})A(\phi, U_k)\Bigr\rvert \leq E_1\sum_{n=j}^{k}(L^{j-n})^{\beta}\abs{\Gamma_n \cap \Omega}
\end{equation}
for $\beta < 1$, and the corresponding sum of the $\mathbf R$-terms is at most $R_1g_j^{\kappa_0}\sum_{n=j}^{k}\abs{\Gamma_n \cap \Omega}$. Here $\Omega \in \mathbf D_k$ is a domain contained in the union of the blocks of $\Lambda_j^0$, $\phi \in C_0^\infty(\Omega^{\sim})$ equals one on $\Omega$, the field is restricted by the test functions of \eqref{eq:density-expansion}, and $\abs{\Gamma_n \cap \Omega}$ is the number of points of the scale $n$. The second is ultraviolet stability, \cite[Corollary~3, (2.50), p.~264]{BalabanCRE}: there are constants $E_\pm$, independent of $\eta$ and of the torus but dependent on $g_k$, with
\begin{equation}\label{eq:ultraviolet-stability}
    \chi_k(T_\eta)\exp\Bigl[-\frac{1}{g_k^2}A(U_k(V_k)) - E_-\abs{T_\eta}\Bigr] \leq \varrho_k(V_k) \leq e^{E_+\abs{T_\eta}}\;.
\end{equation}

These statements are for Ba{\l}aban's parameters, which are functions of the running couplings, and for characteristic functions as test functions. The construction of this paper differs from his in both respects, and the remainder of this section shows that the description and the bounds hold for it.

\subsection{Soft Test Functions}\label{sec:soft-test-functions}

The domains of a term are produced by decompositions of unity. In each step the lattice is divided into cubes $\square$, a nonnegative function $f_\square$ of the fields near $\square$ is compared with a threshold $t$, and the integrand is multiplied by
\begin{equation}\label{eq:sharp-decomposition}
    1 = \sum_{P}\,\prod_{\square \not\subset P}\chi\bigl(\{ f_\square < t \}\bigr)\prod_{\square \subset P}\chi\bigl(\{ f_\square \geq t \}\bigr)\;,
\end{equation}
the sum running over the unions $P$ of cubes. In the term $P$ we say that the test function of a cube has \emph{passed} if the cube lies outside $P$, and that it has \emph{failed} if the cube lies in $P$. Ba{\l}aban's steps contain two kinds of such decompositions. In the first, $f_\square$ is the supremum of $\abs{U(\partial p) - 1}$ over the plaquettes $p$ near $\square$, for the fine field in the first step and for a local minimiser $U_{k+1,\square}$ of the action in the later ones, and $t$ is a plaquette threshold $\varepsilon$ in the units of the plaquette \cite[(1.1), (1.4), p.~246; (2.17), p.~257; (3.2), p.~265]{BalabanCRE}. In the second, $f_\square$ is the supremum over the bonds $b$ near $\square$ of the deviation $\abs{V_k(b)(V^{(k)}_{\square}(b))^{-1} - 1}$ of the field from a reference configuration, the block average of a local minimiser, and $t = 2\vartheta_k$ \cite[(1.8), p.~247; (3.3), p.~265]{BalabanCRE}. We call the factors of \eqref{eq:sharp-decomposition} \emph{test functions}, of the plaquettes and of the reference configuration respectively.

A characteristic function is the natural choice as long as one run of the construction is studied. For the limit $K \to \infty$ we shall compare two runs, with $K$ and $K + 1$ steps, on the lattices which they share, and there the functions $f_\square$ of the two runs differ by a small amount, since they are computed from minimisers on two different fine lattices. A characteristic function can change by one under an arbitrarily small change of its argument, and the set on which this happens, a thin shell around the level set $\{ f_\square = t \}$, has no reason to be of small measure for the densities at hand. We therefore replace the characteristic functions by continuous functions, whose change is proportional to that of the argument.

\begin{definition}[Soft Test Functions]\label{def:soft-test-function}
Let $\kappa_s > 0$. For $x \in [0, \infty]$ put
\begin{equation}\label{eq:soft-function}
    \bar\chi(x) \coloneqq \min\Bigl(1,\, \max\Bigl(0,\, \frac{1 + \kappa_s - x}{\kappa_s}\Bigr)\Bigr)\;,
\end{equation}
so that $\bar\chi(x) = 1$ for $x \leq 1$, $\bar\chi(x) = 0$ for $x \geq 1 + \kappa_s$, and $\bar\chi$ is affine in between. The \emph{soft} form of the decomposition \eqref{eq:sharp-decomposition} is
\begin{equation}\label{eq:soft-decomposition}
    1 = \sum_{P}\,\prod_{\square \not\subset P}\bar\chi\Bigl(\frac{f_\square}{t}\Bigr)\prod_{\square \subset P}\Bigl(1 - \bar\chi\Bigl(\frac{f_\square}{t}\Bigr)\Bigr)\;.
\end{equation}
\end{definition}

\begin{lemma}[Properties of the Soft Decomposition]\label{lem:soft-decomposition}
Let $(f_\square)$ be a finite family of functions with values in $[0, \infty]$ and $t > 0$.
\begin{enumerate}[label=(\roman*)]
\item The identity \eqref{eq:soft-decomposition} holds, and each of its terms has values in $[0, 1]$.
\item The term $P$ of \eqref{eq:soft-decomposition} vanishes unless $f_\square < (1 + \kappa_s)t$ for every $\square \not\subset P$ and $f_\square > t$ for every $\square \subset P$. In particular it is bounded by the term $P$ of the decomposition \eqref{eq:sharp-decomposition} in which the passed test functions are taken at the threshold $(1 + \kappa_s)t$ and the failed ones at the threshold $t$, with strict inequality.
\item For $x, y \in [0, \infty)$, $\abs{\bar\chi(x) - \bar\chi(y)} \leq \abs{x - y}/\kappa_s$. Consequently, for two families $(f_\square)$, $(f'_\square)$, with the convention that $\abs{f_\square - f'_\square}$ is $0$ where both are infinite and $\infty$ where exactly one is, and with $\Pi_P$, $\Pi'_P$ their terms $P$ of \eqref{eq:soft-decomposition},
\begin{equation}\label{eq:soft-difference}
    \abs{\Pi_P - \Pi'_P} \leq \frac{1}{\kappa_st}\sum_{\square}\abs{f_\square - f'_\square}\;.
\end{equation}
\end{enumerate}
\end{lemma}

\begin{proof}
(i) Put $a_\square \coloneqq \bar\chi(f_\square/t) \in [0, 1]$. Expanding the product $\prod_\square(a_\square + (1 - a_\square))$, which equals one, gives the sum over the subsets of the family of the products $\prod_{\square \notin P}a_\square\prod_{\square \in P}(1 - a_\square)$, and each of these is a product of numbers in $[0, 1]$.

(ii) The factor $a_\square$ vanishes for $f_\square \geq (1 + \kappa_s)t$, and the factor $1 - a_\square$ vanishes for $f_\square \leq t$. Since all factors are at most one, the term is bounded by the product of the characteristic functions of the sets $\{ f_\square < (1 + \kappa_s)t \}$, $\square \not\subset P$, and $\{ f_\square > t \}$, $\square \subset P$.

(iii) The function $\bar\chi$ is piecewise affine with slopes $0$ and $-1/\kappa_s$, which gives the first assertion. For the second we replace the factors of $\Pi_P$ by those of $\Pi'_P$ one at a time. Each replacement changes one factor by at most $\abs{f_\square - f'_\square}/(\kappa_st)$, both for a factor $a_\square$ and for a factor $1 - a_\square$, while the other factors lie in $[0, 1]$.
\end{proof}

The first property says that nothing is lost: the lattice integrals are the same whichever decomposition is used to expand them, and only the terms of the expansion change. The second is the form in which the soft test functions enter every estimate. A bound which Ba{\l}aban derives from a passed test function is a bound on a field below the threshold, and it holds for the soft term with the threshold enlarged by the factor $1 + \kappa_s$; a bound which he derives from a failed test function, the small factor which the action gives on a region of large fields, holds unchanged. The third property is used only in Section~\ref{sec:large-field-comparison}.

In the construction of this paper all the test functions of the plaquettes and of the reference configuration are soft, with
\begin{equation}\label{eq:soft-width}
    \kappa_s \coloneqq \frac{\beta}{30}\;,
\end{equation}
$\beta$ being the constant of the spaces of the boundary terms in Section~\ref{sec:density}. This value is dictated by the margins which Ba{\l}aban's inequalities leave between a threshold and the bound that is drawn from it, and which the factor $1 + \kappa_s$ must not exhaust. The characteristic functions which restrict the fluctuation fields are not test functions in the above sense, and they are left as they are.

\subsection{The Depth of the \texorpdfstring{$\mathbf R$}{R}-Operation}\label{sec:depth}

The $\mathbf R$-operation of the step $k$ acts on those connected components of the large-field region $Z_k$ which are small and old: a component must be contained in a cube of side $100MR_k$, and in the preceding $N$ steps no new large-field region may have been created inside it, the earlier regions inside it being small in the same sense on their scales \cite[p.~177, (i), (ii)]{BalabanLFR1}. The number $N$ is the \emph{depth} of the operation. Ba{\l}aban does not fix it. He writes that ``conditions on $N$ will be formulated in constructions of this section'' \cite[p.~177]{BalabanLFR1}, and the conditions which he then states are the following four.
\begin{enumerate}[label=(N\arabic*)]
\item $N > N_0$, where $N_0$ is the smallest positive integer with $L^{-N_0+1}MR_{k-N_0+1} = M$ \cite[p.~179]{BalabanLFR1}.
\item $N \leq O(1)(\log g_k^{-2})^{\nu}$ for a positive integer $\nu$ with $\nu/2 \leq p_0 - p_1 - 1$, and $N$ is so large that $L^{-N}N^{1/2}O(1)B_3^3B_5M^6 \leq \alpha/2$ \cite[(1.94), p.~198]{BalabanLFR1}.
\item $N \leq R_k$ \cite[p.~361]{BalabanLFR2}.
\item $L^{-N} \leq (\log g_k^{-2})^{-\nu}$ for a sufficiently large $\nu$, with $\nu \geq 2p_0 + dr$ \cite[p.~363]{BalabanLFR2}.
\end{enumerate}
The last condition recurs as ``the usual conditions on $N$'', which make $p_0(g_k)q_1(g_k)R_k^7L^{-N}$ small \cite[(1.47), p.~368]{BalabanLFR2}, and as ``our assumptions on $N$'', which make $R_j^{1+\beta_0}(p_0(g_j)/p_1(g_j))L^{-N}$ small \cite[p.~382]{BalabanLFR2}. In the large-field induction of \cite[p.~384]{BalabanLFR2} the conditions (i), (ii) are taken ``with $N = R_j$''. We make a different choice, which satisfies (N1)--(N4) with a depth of the order of $\log\log g_k^{-2}$ only. The reason is that in Section~\ref{sec:inductive-class} the factor $L^N$ occurs, and has to be bounded by a power of $\log g_k^{-2}$.

\begin{definition}[Depth]\label{def:depth}
Put $\nu_L \coloneqq p_0 + q_1 + 7r + 1$. For a tube and an integer $k \leq K$:
\begin{enumerate}[label=(\roman*)]
\item $N_0(k)$ is the smallest positive integer $N_0$ with $L^{N_0-1} = R_{k-N_0+1}$;
\item $N'_k$ is the least integer with $L^{-N'_k} \leq (\log\hat g_k^{-2})^{-\nu_L}$;
\item the \emph{depth} is $N_k \coloneqq \max(N_0(k) + 1,\, N'_k)$.
\end{enumerate}
\end{definition}

For the next lemma we assume, in addition to \eqref{eq:gamma-conditions}, that $p^* \geq \nu_L$ and that
\begin{equation}\label{eq:gamma-depth}
    \log_L\bigl(16L^4(L + 1)^2\bigr) + \nu_L\log_L\lambda \leq \lambda^r
    \qquad\text{for every } \lambda \geq \log\gamma^{-2}\;,
\end{equation}
which holds for $\gamma$ small because $r \geq 2$.

\begin{lemma}[Bounds on the Depth]\label{lem:depth}
Let $\gamma$ satisfy \eqref{eq:gamma-conditions} with $p^* \geq \nu_L$ and \eqref{eq:gamma-depth}, and let $L \geq 3$. For every tube and every integer $k \leq K$:
\begin{enumerate}[label=(\roman*)]
\item $N_0(k)$ exists, $N_0(k) \geq 2$, and $N_0(k)$, $N'_k$, and $N_k$ are nonincreasing in $k$;
\item $L^{-N_k} \leq (\log\hat g_k^{-2})^{-\nu_L}$ and $L^{N_k} \leq 4L^4(L + 1)^2(\log\hat g_k^{-2})^{\nu_L}$;
\item for every history $(x_0, \dots, x_{k'})$ of the tube and $0 \leq k \leq k'$, with $C_{\mathrm{dep}} \coloneqq 16L^4(L + 1)^2$,
\begin{equation}\label{eq:depth-bounds}
    L^{-N_k} \leq (\log g_k^{-2})^{-\nu_L}\;,
    \qquad
    L^{N_k} \leq C_{\mathrm{dep}}(\log g_k^{-2})^{\nu_L}\;,
    \qquad
    N_k \leq R_k\;.
\end{equation}
\end{enumerate}
\end{lemma}

\begin{proof}
(i) Put $n \coloneqq K - k \geq 0$, and for an integer $t \geq 0$ let $\varrho(t) \coloneqq \log_LR_{K-t}$ and $h(t) \coloneqq t - \varrho(t)$. By Lemma~\ref{lem:tube-cubes}(i) and (ii), $\varrho$ takes integer values at least $1$ and $\varrho(t + 1) - \varrho(t) \in \{ 0, 1 \}$, so that $h(t + 1) - h(t) \in \{ 0, 1 \}$ and $h(0) \leq -1$. Moreover $h(t) \to \infty$, because $\varrho(t) \leq 1 + r\log_L\log u_{K-t}$ grows like $\log\log t$. A nondecreasing integer-valued function with steps $0$ or $1$ which starts below $n$ and is unbounded takes the value $n$, and we let $T(n)$ be the least $t$ with $h(t) = n$. With $t = n + N_0 - 1$, so that $K - t = k - N_0 + 1$, the equation $L^{N_0-1} = R_{k-N_0+1}$ reads $\varrho(t) = t - n$, that is, $h(t) = n$. Its solutions have $N_0 = \varrho(t) + 1 \geq 2$, and the smallest is $N_0(k) = T(n) - n + 1$. Since $h$ is nondecreasing and $h(T(n)) = n$, we have $h(t) \leq n$ for $t \leq T(n)$, hence $T(n + 1) \geq T(n) + 1$, and $T(n) - n$ is nondecreasing in $n$. So $N_0(k)$ is nonincreasing in $k$. The same holds for $N'_k$, because $\hat g_k$ increases with $k$, and therefore for $N_k$.

(ii) The first inequality is the definition of $N'_k \leq N_k$. By the minimality of $N'_k$ we have $L^{N'_k} < L(\log\hat g_k^{-2})^{\nu_L}$. For $N_0 = N_0(k)$ we use the second inequality of \eqref{eq:tube-two-nine} with $m = k - N_0 + 1$ and $n = k$: with $x \coloneqq N_0 - 1 \geq 1$,
\begin{equation}\label{eq:depth-n-zero}
    L^{x} = R_{k-x} \leq (L + 1)x^{\beta_0}R_k \leq (L + 1)(x + 2)R_k \leq 2(L + 1)L^{x/2}R_k\;,
\end{equation}
where the last step is the inequality $x + 2 \leq 2L^{x/2}$, which holds for $x \geq 0$ and $L \geq 3$, since the two sides agree at $x = 0$ and the derivative $L^{x/2}\log L$ of the right side is at least $\log 3 > 1$. Hence $L^{x/2} \leq 2(L + 1)R_k$, and with $R_k < L(\log\hat g_k^{-2})^r$,
\begin{equation}\label{eq:depth-n-zero-bound}
    L^{N_0+1} = L^2L^{x} \leq 4L^2(L + 1)^2R_k^2 < 4L^4(L + 1)^2(\log\hat g_k^{-2})^{2r}\;.
\end{equation}
Since $\nu_L \geq 2r$ and $\log\hat g_k^{-2} \geq 1$, both $L^{N_0+1}$ and $L^{N'_k}$ are at most $4L^4(L + 1)^2(\log\hat g_k^{-2})^{\nu_L}$.

(iii) By Lemma~\ref{lem:tube-couplings}(i), $\log g_k^{-2} \leq \log\hat g_k^{-2} \leq \log g_k^{-2} + \log\rho$, so that the first inequality follows from (ii), and the second from (ii) and $(\log\hat g_k^{-2})^{\nu_L} \leq 2(\log g_k^{-2})^{\nu_L}$, which is the third condition of \eqref{eq:gamma-conditions} with $\nu_L \leq p^*$. Taking logarithms, $N_k \leq \log_LC_{\mathrm{dep}} + \nu_L\log_L\log g_k^{-2}$, which is at most $(\log g_k^{-2})^r \leq R_k$ by \eqref{eq:gamma-depth} and Lemma~\ref{lem:tube-cubes}(i).
\end{proof}

The depth $N_k$ satisfies Ba{\l}aban's four conditions. (N1) holds by definition and (N3) by \eqref{eq:depth-bounds}. For (N4), \eqref{eq:depth-bounds} gives the exponent $\nu_L$, and $\nu_L - (2p_0 + dr) = q_1 - p_0 + 3r + 1 > 0$ at $d = 4$, because $q_1 \geq p_0$ (Appendix~\ref{app:constants}). The two further uses of (N4) hold as well: with $R_k \leq 2L(\log g_k^{-2})^r$, the product $p_0(g_k)q_1(g_k)R_k^7L^{-N_k}$ is at most a constant times $(\log g_k^{-2})^{p_0+q_1+7r-\nu_L} = (\log g_k^{-2})^{-1}$, and $R_k^{1+\beta_0}(p_0(g_k)/p_1(g_k))L^{-N_k}$ is at most a constant times $(\log g_k^{-2})^{r(1+\beta_0)+p_0-p_1-\nu_L}$, whose exponent is negative. For (N2), the first restriction holds with $\nu = \ceil{r}$, because $N_k \leq R_k \leq 2L(\log g_k^{-2})^r$, under the condition $\ceil{r}/2 \leq p_0 - p_1 - 1$ on the exponents, and the second holds for $\gamma$ small, because $L^{-N_k}N_k^{1/2} \leq (2L)^{1/2}(\log g_k^{-2})^{r/2-\nu_L}$ tends to zero. That these are all the places at which the depth enters the two papers, and that each of them holds for $N_k$, is Supplement, Proposition~S.3.

\subsection{The Construction of This Paper}\label{sec:the-construction}

We can now say which construction we run. It is Ba{\l}aban's, with the parameters of a tube and with soft test functions.

\begin{definition}[The Construction]\label{def:construction}
Let a tube be given, with $\gamma$ as in Lemma~\ref{lem:depth}, and let $g_0 > 0$. The \emph{construction for the tube started at $g_0$} is the sequence of densities $\varrho_0, \varrho_1, \dots$ obtained from $\varrho_0 = \exp[-(1/g_0^2)A - E]$ by Ba{\l}aban's operations $\mathbf RT$ of \cite{BalabanCRE,BalabanLFR1,BalabanLFR2}, with the following changes.
\begin{enumerate}[label=(\roman*)]
\item Every parameter of the scale $j$ which is defined there as a function of the coupling $g_j$ is that function at $\hat g_j$: the cube size $R_j$, the thresholds $\varepsilon_j$, $\vartheta_j$, $\delta'_j$ and the radii $\alpha_{0,j}$, $\alpha_{1,j}$ of Definition~\ref{def:tube-parameters}, and every other threshold, radius, or support of a cutoff function of that scale. The depth of the $\mathbf R$-operation of the step $k$ is $N_k$.
\item The characteristic functions which restrict a fluctuation field are functions of the unscaled field: the fluctuation field $A_j$ of the step $j + 1$ is restricted by $\abs{g_jA_j(b)} < \vartheta_j$, which is the form of \cite[(1.18), p.~250]{BalabanCRE} and \cite[(1.9), p.~178]{BalabanLFR1}, and the field $B_j$ of \cite[(1.55), p.~188]{BalabanLFR1} by $\abs{g_jB_j(b)} < \delta'_j$; the bump functions of the changes of variables \cite[(1.22), p.~251]{BalabanCRE} are functions of the unscaled field as well, with supports fixed by the thresholds (Supplement, Definition~3.5).
\item The test functions of the plaquettes and of the reference configuration are the soft test functions of Definition~\ref{def:soft-test-function}, with the width \eqref{eq:soft-width}; the cascade of test functions of the $\mathbf R$-operation \cite[(1.22)--(1.24), p.~181]{BalabanLFR1} is soft with the nested widths of Supplement, Definition~2.2, and so is the cutoff $\chi(\Lambda_i)$ of \cite[(1.101), p.~201]{BalabanLFR1}.
\end{enumerate}
The \emph{running couplings} are $g_k = x_k^{-1/2}$ with $x_{k+1} \coloneqq x_k - \beta_{k+1}$, where $\beta_{k+1} = \beta_{k+1}(g_0, \dots, g_k)$ is the coefficient which Ba{\l}aban's formula \cite[(3.61), p.~282]{BalabanCRE} assigns to the regular part of the effective action of the step $k + 1$, the coefficient of \eqref{eq:position-coupling}. The construction is defined up to the step $k \leq K$ as long as $(x_0, \dots, x_k)$ is a history of the tube.
\end{definition}

In (ii) the factor $g_j$ which converts the field of the Gaussian integral into the unscaled field is the running coupling, and only the threshold is fixed by the target. For this reason the construction is continuous in the couplings: no cube size, no depth, and no set on which a characteristic or test function is supported depends on $g_0, \dots, g_k$, which enter only through the action, the Gaussian weights, and the terms of the effective action generated before.

\begin{proposition}[Ba{\l}aban's Description for the Construction]\label{prop:construction-admissible}
Let $\gamma$ be sufficiently small. For every tube, every $g_0$, and every $k \leq K$ up to which the construction for the tube started at $g_0$ is defined, the density $\varrho_k$ has the expansion \eqref{eq:density-expansion}, with soft test functions in \eqref{eq:test-function}, and its terms have the properties (a)--(e) of Section~\ref{sec:density}, with constants $E_0$, $B_0$, $\kappa$, $\kappa_0$ which do not depend on the tube, on $g_0$, on $k$ and $K$, or on the torus. The bound \eqref{eq:balaban-theorem-two} holds with a constant which is independent of the same quantities; the bound \eqref{eq:ultraviolet-stability} holds with constants $E_\pm$ which depend on these quantities through the coupling $g_k$ only, as in Ba{\l}aban's statement; and the results of \cite{BalabanLFR1,BalabanLFR2} on the $\mathbf R$-operation hold as well.
\end{proposition}

\begin{proof}
Supplement, Propositions~S.1, S.2, and S.3. Proposition~S.1 goes through the places at which a parameter enters \cite{BalabanCRE,BalabanLFR1,BalabanLFR2} and the papers on which they rest, and shows that each of them is of one of four kinds: a relation among the parameters alone, which holds for the sequence $(\hat g_j)$ by Lemmata~\ref{lem:tube-couplings}--\ref{lem:tube-thresholds}; a lower bound drawn from the action, which holds with the constant in the exponent divided by $\rho \leq 3/2$, because the coefficient of the action is at least $\hat g_m^{-2}/\rho$ at the points which carry a field of the scale $m$ (Supplement, Lemma~2.1); an upper bound relative to the action, which holds because the thresholds of a tube are the smaller ones; or a conversion between a scaled and an unscaled field, which is the change (ii). Proposition~S.2 goes through the places at which a test function enters and shows that every bound drawn from a passed test function holds at the threshold $(1 + \kappa_s)t$ of Lemma~\ref{lem:soft-decomposition}(ii). Proposition~S.3 is the corresponding statement for the depth.
\end{proof}

\subsection{The Stopping Index}\label{sec:stopping-index}

The loops are drawn on a lattice $\mathbb T_q$ whose spacing $uL^{-q}$ does not depend on $K$, and the construction is stopped on a lattice which is finer than $\mathbb T_q$ by a factor that does not depend on $K$ either. This factor has to be so large that, at every step, the straight segments of a loop are long compared with the cubes of the construction. Since the parameters of a tube at the scale $K - n$ depend on $n$ and not on $K$, we may write
\begin{equation}\label{eq:cube-size-below-unit}
    R^{[n]} \coloneqq R_{K-n}\;,
\end{equation}
the smallest power of $L$ which is at least $(\log(\bar g^{-2} + b'n + C_W))^r$.

\begin{definition}[Stopping Index]\label{def:stopping-index}
Let $M = L^m$. For an integer $j \geq 0$ put
\begin{equation}\label{eq:stopping-function}
    \bar q(j) \coloneqq 3 + m + \log_LR^{[j+1]}\;,
    \qquad\text{so that}\qquad
    L^{\bar q(j)} = L^3MR^{[j+1]}\;.
\end{equation}
The \emph{stopping index} $j_q$ of the level $q \geq 0$ is the least integer $j$ with $j \geq q + \bar q(j)$. For a target $(K, \bar g)$ with $K > j_q$, the \emph{stopping step} is $K' \coloneqq K - j_q$.
\end{definition}

\begin{lemma}[The Stopping Index]\label{lem:stopping-index}
Let $\gamma$ be as in Lemma~\ref{lem:tube-cubes}. For every $q \geq 0$ the stopping index $j_q$ exists. It depends on $q$, $L$, $M$, $r$, $\bar g$, $b'$, and $C_W$ only, and
\begin{equation}\label{eq:stopping-identity-index}
    j_q = q + \bar q(j_q)\;,
    \qquad
    L^{j_q-q} = L^3MR_{K'-1} \geq L^3MR_{K'}\;.
\end{equation}
Consequently, for a loop $C$ at level $q$, read on the lattice of the step $k \leq K'$, whose spacing is $uL^{k-K}$: the corners of $C$ lie in a coset of $a_k\mathbb Z^4$ with $a_k \coloneqq L^{j_q-q}L^{K'-k}$, and, in the notation of Section~\ref{sec:loops} for this lattice,
\begin{equation}\label{eq:lengths-at-step}
    P_k(C) = P_q(C)\,a_k\;,
    \qquad
    s_k(C) \geq a_k\;,
    \qquad
    \delta_k(C) \geq a_k\;,
    \qquad
    a_k \geq L^3MR_{K'}L^{K'-k} \geq L^2MR_k\;.
\end{equation}
Two loops of an ensemble at level $q$ are at $\ell^\infty$-distance at least $a_k$ on that lattice.
\end{lemma}

\begin{proof}
The function $\bar q$ is nondecreasing, because $R^{[n]}$ is nondecreasing in $n$ by Lemma~\ref{lem:tube-cubes}(ii), and $\bar q(j) \leq 4 + m + r\log_L\log(\bar g^{-2} + b'(j + 1) + C_W)$ by Lemma~\ref{lem:tube-cubes}(i). Hence $j - q - \bar q(j) \to \infty$ as $j \to \infty$, and the least $j$ with $j \geq q + \bar q(j)$ exists. It is determined by $q$ and the function $\bar q$, which gives the dependence. Since $\bar q \geq 3$, we have $j_q > q$, so that $j_q - 1 \geq q$ and, by minimality, $j_q - 1 < q + \bar q(j_q - 1) \leq q + \bar q(j_q)$. Together with $j_q \geq q + \bar q(j_q)$ this gives the first identity of \eqref{eq:stopping-identity-index}. The second is \eqref{eq:stopping-function} with $R^{[j_q+1]} = R_{K-j_q-1} = R_{K'-1}$, and $R_{K'-1} \geq R_{K'}$ by Lemma~\ref{lem:tube-cubes}(ii).

The lattice of the step $k$ has the spacing $uL^{k-K}$, and $uL^{-q} = L^{K-k-q}uL^{k-K}$ with $K - k - q = (j_q - q) + (K' - k)$. So one bond of $\mathbb T_q$ is $a_k$ bonds of that lattice, which gives the statement on the corners and the first identity of \eqref{eq:lengths-at-step}. The bounds on $s_k(C)$ and $\delta_k(C)$ are Lemma~\ref{lem:separation}(i) with $a = a_k$. For two loops of an ensemble we apply Lemma~\ref{lem:separation}(ii) to a lift of the first and to all the translates of a lift of the second by the periods of the torus, which are multiples of $a_k$ because $T/(uL^{-q}) = 2L^{m_T+q}$ is an integer: these lifts have no common site, since the loops are vertex-disjoint on the torus, and their corners lie in one coset of $a_k\mathbb Z^4$. For the last chain we use $L^{j_q-q} \geq L^3MR_{K'}$ and, by Lemma~\ref{lem:tube-cubes}(ii) applied $K' - k$ times, $R_k \leq L^{K'-k}R_{K'}$.
\end{proof}

The inequality $a_k \geq L^2MR_k$ says that every segment of a loop, and every gap between two of its segments or between two loops, is at least $L^2$ cubes of the step $k$ long, at every step up to the stopping step. On the stopping lattice itself a loop at level $q$ has $P_{K'}(C) = P_q(C)L^3MR_{K'-1}$ bonds, a number which does not depend on $K$.

The steps $k \leq K'$ are those at which the loops will be carried, and their couplings are bounded by a number which does not depend on $K$. Indeed, for a history $(x_0, \dots, x_k)$ of the tube with $k \leq K'$ we have $x_k \geq \ell_k = t + b(K - k) - C_W \geq t + bj_q - C_W$ by \eqref{eq:tube-couplings}, because $K - k \geq K - K' = j_q$, and $\hat g_k \leq g_k$ by Lemma~\ref{lem:tube-couplings}(i). With
\begin{equation}\label{eq:stopping-coupling}
    \gamma_q \coloneqq \bigl(\bar g^{-2} + bj_q - C_W\bigr)^{-1/2}
\end{equation}
we therefore have $\hat g_k \leq g_k \leq \gamma_q$ for $0 \leq k \leq K'$. The number $\gamma_q$ is small if $\bar g$ is small and also if $q$ is large: since $\bar g^{-2} - C_W \geq \gamma^{-2}$ by \eqref{eq:tube-class} and $j_q \geq q$ by \eqref{eq:stopping-identity-index},
\begin{equation}\label{eq:stopping-coupling-bound}
    \gamma_q \leq \bigl(\bar g^{-2} + bq - C_W\bigr)^{-1/2} \leq \min\Bigl(\bigl(\bar g^{-2} - C_W\bigr)^{-1/2},\ \bigl(\gamma^{-2} + bq\bigr)^{-1/2}\Bigr)\;,
\end{equation}
and $(\bar g^{-2} - C_W)^{-1/2} \leq (1 - \beta_0/4)^{-1/2}\bar g \leq (1 + \beta_0)\bar g$, because $C_W \leq \beta_0\bar g^{-2}/4$ by \eqref{eq:tube-class} and $(1 + \beta_0)^2(1 - \beta_0/4) \geq 1$ for $0 < \beta_0 \leq 1$. Every requirement of the paper which involves the representations of an ensemble concerns the steps $k \leq K'$ only, and it is imposed as an upper bound on $\gamma_q$; by \eqref{eq:stopping-coupling-bound} such a bound holds at every coupling $\bar g \leq \gamma$ once the level $q$ is large enough.

\section{The Coupling Flow}\label{sec:coupling-flow}

The construction of Section~\ref{sec:construction} takes a bare coupling and returns the running couplings. What the lattice theory of Section~\ref{sec:lattice-theory} needs is the converse: for a prescribed value $\bar g$ of the coupling at the unit scale and for every number $K$ of steps, a bare coupling $g_0$ with $g_K = \bar g$, together with bounds on the couplings in between which do not depend on $K$. Ba{\l}aban states this as a theorem for $G = SU(2)$ and writes that ``a proof of this theorem, based on perturbative calculations, will be given in a separate paper'' \cite[Theorem~2, p.~259]{BalabanRG1}; in \cite{BalabanCRE} the bounds on the couplings are a hypothesis, as Theorem~\ref{thm:balaban-inductive} shows. In this section we prove the statement, for a simple group and for the construction of this paper.

\begin{theorem}[Coupling Flow]\label{thm:coupling-flow}
There are $\gamma' > 0$ and constants $0 < b \leq b'$ and $C_W \geq 0$ such that, for every integer $K \geq 1$ and every $\bar g \in (0, \gamma']$, there is a bare coupling $g_0 = g_0(K, \bar g)$ for which the running couplings $g_0, g_1, \dots, g_K$ of the construction for the target $(K, \bar g)$ satisfy $g_K = \bar g$ and, for $0 \leq k \leq l \leq K$,
\begin{equation}\label{eq:coupling-window}
    \bar g^{-2} + b(K - k) - C_W \leq g_k^{-2} \leq \bar g^{-2} + b'(K - k) + C_W
    \qquad\text{and}\qquad
    g_k^{-2} \geq g_l^{-2} - C_W\;.
\end{equation}
The constants are $b = 2c\ln L - \delta'$ and $b' = 2c\ln L + \delta'$, where $c$ is the number \eqref{eq:one-loop-coefficient}, the first coefficient of the $\beta$-function, and $0 \leq \delta' \leq (2/5)c\ln L$.
\end{theorem}

The \emph{target} $(K, \bar g)$ stands for the tube of Definition~\ref{def:tube} with these $b$, $b'$, and $C_W$. The proof has three inputs of different kinds. The first is a computation at one loop: at zero coupling the inverse square of the coupling decreases by $2c\ln L$ per step, up to an error which stays bounded when it is summed over any number of steps (Theorem~\ref{thm:one-loop-window}). The second is perturbative as well, but concerns Ba{\l}aban's expansions: the true decrease differs from the one at zero coupling by a multiple of $g_k^2$, uniformly in the step and in the earlier couplings (Proposition~\ref{prop:coefficients}(i), (ii)). With these two facts every bare coupling would produce a sequence that satisfies \eqref{eq:coupling-window} up to its last step, but the last coupling would be whatever it turns out to be. To prescribe it we need the third input, the continuity of the couplings as functions of the bare coupling (Proposition~\ref{prop:coefficients}(iii)), and it is for this that the parameters of the construction were detached from the running couplings in Section~\ref{sec:tube}. Given the three inputs, the bare coupling is found by the intermediate value theorem.

\subsection{The Coefficients of the Flow}\label{sec:flow-coefficients}

We first say what the coefficient $\beta_{k+1}$ of Definition~\ref{def:construction} is. Let $\mathfrak g \subset \mathfrak u(N)$ be the Lie algebra of $G$, a space of anti-Hermitian matrices, with the inner product $\langle X, Y \rangle \coloneqq -\tr(XY)$ and the Killing form $\mathrm K(X, Y) \coloneqq \operatorname{Tr}(\operatorname{ad}X\operatorname{ad}Y)$.

\begin{lemma}[Invariant Forms]\label{lem:invariant-forms}
Every bilinear form on $\mathfrak g$ which is invariant under $\operatorname{Ad}(G)$ is a real multiple of the inner product. In particular there is a number $\kappa > 0$, depending on $G$ and on its realisation, with
\begin{equation}\label{eq:killing-constant}
    \mathrm K(X, Y) = \kappa\tr(XY)\;,
\end{equation}
and $\kappa = 2N^2$ for $G = SU(N)$ in its defining representation.
\end{lemma}

\begin{proof}
An invariant bilinear form can be written as $(X, Y) \mapsto \langle SX, Y \rangle$ with a linear map $S$ of $\mathfrak g$ which commutes with $\operatorname{Ad}(G)$, hence with every $\operatorname{ad}X$, and we treat its symmetric and antisymmetric parts separately, each of which commutes with $\operatorname{ad}$ as well, because the maps $\operatorname{ad}X$ are antisymmetric for the invariant inner product. A symmetric $S$ has a real eigenvalue $\lambda$, and $\ker(S - \lambda)$ is a nonzero ideal of $\mathfrak g$, since $S[X, Y] = [X, SY]$. As $\mathfrak g$ is simple, $S = \lambda$. For an antisymmetric $S$ the square $S^2$ is symmetric and commutes with $\operatorname{ad}$, so that $S^2 = -\mu^2$ with $\mu \geq 0$ by what has been shown and $\langle S^2X, X \rangle = -\abs{SX}^2$. If $\mu$ were positive, $J \coloneqq S/\mu$ would be a complex structure on $\mathfrak g$ with $J[X, Y] = [X, JY]$, so that $\operatorname{ad}(JX) = J\operatorname{ad}X$ and $\mathrm K(JX, JX) = \operatorname{Tr}(J\operatorname{ad}X\,J\operatorname{ad}X) = \operatorname{Tr}(J^2(\operatorname{ad}X)^2) = -\mathrm K(X, X)$. This is impossible, because $\mathrm K(X, X) = -\operatorname{Tr}((\operatorname{ad}X)^{T}\operatorname{ad}X)$ is negative for every $X \neq 0$, the centre of a simple Lie algebra being trivial. Hence $S = 0$. The Killing form is invariant and negative definite, which gives \eqref{eq:killing-constant} with $\kappa > 0$, and for $\mathfrak{su}(N)$ one has $\mathrm K(X, Y) = 2N\operatorname{Tr}(XY) = 2N^2\tr(XY)$.
\end{proof}

The number
\begin{equation}\label{eq:one-loop-coefficient}
    c \coloneqq \frac{11}{48\pi^2}\,\kappa
\end{equation}
is the first coefficient of the $\beta$-function in the normalisation of the action \eqref{eq:lattice-measure}: for $SU(N)$ and the usual coupling $g_{\mathrm s}$, with $g^{-2} = 2Ng_{\mathrm s}^{-2}$, the statement that $g^{-2}$ decreases by $2c\ln L$ when the lattice spacing is multiplied by $L$ is $\dd g_{\mathrm s}^{-2}/\dd\ln\varepsilon = -11N/(24\pi^2)$. From now on the numbers $b_*$ and $b^*$ of Section~\ref{sec:tube} are
\begin{equation}\label{eq:tube-class-constants}
    b_* \coloneqq c\ln L\;,
    \qquad
    b^* \coloneqq 3c\ln L\;,
\end{equation}
and a tube is always one that satisfies \eqref{eq:tube-class} with them.

\emph{The polarisation coefficient.} We use the units of the step $k + 1$, in which the block field of that step lives on the unit lattice. Let $F$ be a function of the background field of the step $k + 1$ on the lattice $\mathbb Z^4$, which is invariant under gauge transformations and under the symmetries of the lattice and is a sum of local terms of the kind \eqref{eq:regular-part}, so that its derivatives at the trivial field are summable. For a block field $V = \exp iB$, with $B$ a field on the bonds with values in $i\mathfrak g$ and components $B^a_\mu(x)$ in an orthonormal basis, the second derivative of $B \mapsto F(U_{k+1}(\exp iB))$ at $B = 0$ with respect to $B^a_\mu(x)$ and $B^{a'}_\nu(x')$ is, for fixed $\mu$, $x$, $\nu$, $x'$, an invariant bilinear form on $\mathfrak g$, because $F$ is invariant under constant gauge transformations. By Lemma~\ref{lem:invariant-forms} and translation invariance it equals $\delta^{aa'}\Pi_{\mu\nu}(F; x - x')$, and we put
\begin{equation}\label{eq:polarisation-coefficient}
    \Lambda(F) \coloneqq \sum_{x \in \mathbb Z^4}\Pi_{\mu\nu}(F; x)\,x_\mu x_\nu
\end{equation}
for a pair $\mu \neq \nu$, on which the value does not depend. This is Ba{\l}aban's definition of the $\beta$-function from the vacuum polarisation \cite[(1.20)--(1.22), p.~264]{BalabanRG1}, where the independence of the pair is part of the statement and the limit of the infinite lattice is taken with the help of the local representation. The coefficient of the step $k + 1$ is $\Lambda$ of the regular part which that step generates on the whole lattice,
\begin{equation}\label{eq:beta-definition}
    \beta_{k+1} \coloneqq \Lambda\bigl(\mathbf E^{(k+1)}\bigr)\;,
\end{equation}
and \cite[(3.57)--(3.64), pp.~281--283]{BalabanCRE} shows that this is the number for which the counterterm $-\beta_{k+1}A(\phi_{k+1}, U)$ of \eqref{eq:regular-part} removes the marginal part of $\mathbf E^{(k+1)}$. The regular part contains the logarithm of the normalisation factor of the Gaussian measure of the fluctuation field of the step,
\begin{equation}\label{eq:gaussian-normalisation}
    Z^{(k)}(U) \coloneqq \int\delta(\tilde QB)\exp\bigl[-\tfrac12\langle B, \Delta^{(k)}(U)B \rangle\bigr]\,\dd B
\end{equation}
\cite[(1.3), (1.4), p.~260]{BalabanRG1}, which is its part of order zero in the couplings, and the rest is the fluctuation integral proper, $\mathbf E^{(k+1)} = \log Z^{(k)} + \mathbf E^{(k+1)}_{\mathrm{fl}}$ up to a constant. Since $\Lambda$ is linear,
\begin{equation}\label{eq:beta-split}
    \beta_{k+1} = b_{k+1} + \Lambda\bigl(\mathbf E^{(k+1)}_{\mathrm{fl}}\bigr)\;,
    \qquad
    b_{k+1} \coloneqq \Lambda\bigl(\log Z^{(k)}\bigr)\;.
\end{equation}
The numbers $b_1, b_2, \dots$ are the \emph{coefficients at zero coupling}. They are determined by the quadratic forms $\Delta^{(k)}(U)$ of \cite[(1.5), p.~261]{BalabanRG1}, that is, by $L$, $G$, the averaging operations, and $k$, and they involve no coupling, threshold, or tube.

\emph{Histories which are not flows.} For the argument that will produce the bare coupling we need $\beta_{k+1}$ as a function on the whole set $H_k$ of Definition~\ref{def:tube}, whereas the construction started at $g_0$ only visits the points of $H_k$ with $x_{j+1} = x_j - \beta_{j+1}$ for all $j < k$. We therefore extend the construction. For $h = (x_0, \dots, x_k) \in H_k$ the densities $\varrho^h_0, \dots, \varrho^h_k$ are defined inductively: $\varrho^h_0 \coloneqq \exp[-x_0A - E]$, and $\varrho^h_{j+1}$ is obtained by applying the operations $\mathbf RT$ of Definition~\ref{def:construction} to $\varrho^h_j$, with $g_j = x_j^{-1/2}$ as the coupling of the scale $j$, and by multiplying each term of the resulting expansion \eqref{eq:density-expansion} by
\begin{equation}\label{eq:history-factor}
    \exp\bigl[-\bigl(x_{j+1} - x_j + \beta_{j+1}(x_0, \dots, x_j)\bigr)A(U_{j+1})\bigr]\;,
\end{equation}
where $A(U_{j+1})$ is the Wilson action of the background field of the term and $\beta_{j+1}(x_0, \dots, x_j)$ is the coefficient \eqref{eq:beta-definition} of the step just performed. The effect is that the Wilson term of every term of $\varrho^h_{j+1}$ has the coefficient function
\begin{equation}\label{eq:history-coupling}
    \frac{1}{g_{j+1}^2(x)} = x_{j+1} + \sum_{i=1}^{j+1}\beta_i(x_0, \dots, x_{i-1})\bigl(1 - \phi_i(x)\bigr)\;,
\end{equation}
which is $x_{j+1}$ on the domain $\Lambda_{j+1}^{\sim-1}$ where all the $\phi_i$ equal one. Indeed, by \eqref{eq:position-coupling} the step turns the coefficient function $y$ of $\varrho^h_j$ into $y - \beta_{j+1}\phi_{j+1}$, and \eqref{eq:history-factor} adds the constant $x_{j+1} - x_j + \beta_{j+1}$. If $h$ is a flow, the factor \eqref{eq:history-factor} is one at every step and $\varrho^h_j$ is the density of the construction started at $g_0 = x_0^{-1/2}$. In all cases $\beta_{k+1}(h)$ denotes the coefficient \eqref{eq:beta-definition} of the step $k + 1$ applied to $\varrho^h_k$.

\begin{proposition}[The Coefficients on a Tube]\label{prop:coefficients}
There are constants $C_\Lambda$, $C_3$, $C_{\mathrm{rem}}$, $C_R$, and $\gamma_1 > 0$ such that, for every $\gamma \in (0, \gamma_1]$, every tube, every $k < K$, and every history $h = (x_0, \dots, x_k) \in H_k$, with $g_k = x_k^{-1/2}$, the following holds. The densities $\varrho^h_0, \dots, \varrho^h_k$ are defined and have the properties stated in Proposition~\ref{prop:construction-admissible}, and:
\begin{enumerate}[label=(\roman*)]
\item $\abs{\beta_{k+1}(h) - b_{k+1}} \leq \delta'$, where $\delta' \coloneqq C_\Lambda C_3\gamma p_0(\gamma) + C_R\gamma^{\kappa_0-4}$;
\item $\abs{\beta_{k+1}(h) - b_{k+1}} \leq (C_{\mathrm{rem}} + C_R\gamma^{\kappa_0-6})\,g_k^2$;
\item the function $\beta_{k+1}$ is continuous on $H_k$.
\end{enumerate}
The constants depend on those of the construction and not on $\gamma$, the tube, $k$, or $h$.
\end{proposition}

\begin{proof}
Supplement, Propositions~S.1, S.4, and~S.5. Let us say where the difficulty lies in each part. That the densities of a history which is not a flow have Ba{\l}aban's properties is part of Proposition~S.1, since the coefficient function \eqref{eq:history-coupling} differs from that of a flow only through coefficients which obey the same bounds. By \eqref{eq:beta-split} the difference in (i) and (ii) is $\Lambda$ of the fluctuation integral, whose local terms Ba{\l}aban bounds by a multiple of the threshold $\varepsilon_1$ of the fluctuation field \cite[(2.41), p.~21]{BalabanRG2}. In the construction of this paper that threshold is $\vartheta_k$, which in the units of the step is at most $\gamma p_0(\gamma)$, and since $\Lambda$ is bounded on local terms with exponential decay this gives the first term of $\delta'$. The second term accounts for the $\mathbf R$-terms of the earlier steps, which enter the regular part through the fluctuation integral and which the bound \eqref{eq:r-bound} makes small. The improvement to $g_k^2$ in (ii) rests on a symmetry: when $g_k$ is replaced by a complex number $z$ in the exponent of the fluctuation integral and the threshold is held fixed, the integral is an even analytic function of $z$ which vanishes at $z = 0$ up to a Gaussian tail. For (iii) one has to follow every term of the density through the step and to check that it is continuous in the history at each configuration, which is where the construction of Section~\ref{sec:construction} is used: no set on which a characteristic function or a test function is supported moves with the history.
\end{proof}

The exponent $\kappa_0$ is that of \eqref{eq:r-bound}, and Ba{\l}aban takes $\kappa_0 \geq 7$ \cite[(2.46), p.~263]{BalabanCRE}, so that $\delta' \to 0$ as $\gamma \to 0$. The first input of the proof, on the other hand, concerns the numbers $b_k$ alone.

\begin{theorem}[The Coefficients at Zero Coupling]\label{thm:one-loop-window}
There is a constant $C_W$, depending on $L$, $G$, and the averaging operations only, such that for all integers $0 \leq m < n$
\begin{equation}\label{eq:one-loop-window}
    \Bigl\lvert\sum_{j=m+1}^{n}b_j - 2c\,(n - m)\ln L\Bigr\rvert \leq C_W\;.
\end{equation}
\end{theorem}

\begin{proof}
Supplement, Proposition~S.6. The statement is the content of Ba{\l}aban's Theorem~2 of \cite{BalabanRG1} at zero coupling, and we describe the structure of its proof, since the number $c$ is computed there. It suffices to show that $\Phi_k \coloneqq \sum_{j=1}^{k}b_j$ satisfies $\abs{\Phi_k - 2ck\ln L} \leq C_W/2$ for every $k$. By \eqref{eq:beta-split} and the local representation of the terms of order zero \cite[(0.28)--(0.30), p.~258]{BalabanRG1}, $\Phi_k$ is, up to a bounded error, the coefficient of the polarisation of the Gaussian factor $\sum_{j<k}\log Z^{(j)}(U_k)$ of $k$ steps, measured against the polarisation of the action of the background field, whose symbol is bounded above and below at a fixed nonzero momentum \cite[(1.67), p.~29]{BalabanPRT1}. The Gaussian factor of $k$ steps is then written exactly as a difference $F_k - F_0$ of two functionals, each of which is one half of the logarithm of the determinant of a gauge-invariant quadratic form, of the action at the level $k$ for $F_k$ and of the action at the level $0$ for $F_0$, divided by the determinant of the corresponding gauge operator, plus a sum of terms which come from the current $J$ of \cite[(1.5), (1.8), p.~261]{BalabanRG1}. The identity is linear algebra (a Schur complement and the matrix determinant lemma), and it telescopes over the steps. Three estimates of symbols then give the theorem: the polarisation of the terms of the current is bounded uniformly in $k$, because gauge covariance makes it vanish to second order at zero momentum; the polarisation of $F_k$ is bounded uniformly in $k$ at a fixed momentum, because it is the one-loop polarisation of a theory on a lattice of unit spacing with a cutoff at unit scale; and the polarisation of $-F_0$, the one-loop polarisation of the bare lattice theory seen from the coarse lattice, is $c\ln(1/\eta_k^2)$ times the action of the background up to a bounded error, with $\eta_k = L^{-k}$. The last estimate is the computation of the one-loop coefficient of the lattice theory in a background field, in which the logarithm comes from the region of momenta between the two cutoffs, and the coefficient $11/48\pi^2$ is the one of the continuum computation, $\kappa$ entering through Lemma~\ref{lem:invariant-forms}.
\end{proof}

We now show how the coupling flow follows.

\subsection{Proof of the Coupling Flow}\label{sec:flow-proof}

\begin{proof}[Proof of Theorem~\ref{thm:coupling-flow}]
\emph{The constants.} Let $C_W$ be the constant of Theorem~\ref{thm:one-loop-window}. We fix $\gamma \in (0, \gamma_1]$ so small that the conditions of Section~\ref{sec:construction} hold with \eqref{eq:tube-class-constants} and that
\begin{equation}\label{eq:flow-gamma}
    \delta' \leq \tfrac{2}{5}c\ln L
    \qquad\text{and}\qquad
    \gamma^{-2} \geq \Bigl(\frac{4}{\beta_0} - 1\Bigr)C_W\;,
\end{equation}
and we put $b \coloneqq 2c\ln L - \delta'$, $b' \coloneqq 2c\ln L + \delta'$, and $\gamma' \coloneqq (\gamma^{-2} + C_W)^{-1/2}$. Then $b \geq (8/5)c\ln L \geq b_*$, $b' \leq (12/5)c\ln L \leq b^*$, and $b'/b \leq 3/2$, and for $\bar g \leq \gamma'$ the number $t = \bar g^{-2}$ satisfies $t \geq \gamma^{-2} + C_W \geq 4C_W/\beta_0$. So $(K, b, b', C_W, \bar g)$ is a tube in the sense of Definition~\ref{def:tube}, for every $K \geq 1$ and $\bar g \in (0, \gamma']$. The choice is not circular, since the constants of Proposition~\ref{prop:coefficients}, and hence $\delta'$, are the same for all tubes.

\emph{The window of the coefficients.} Let $0 \leq m < n \leq K$ and let $h_j \in H_j$ be arbitrary histories, $m \leq j < n$. By Proposition~\ref{prop:coefficients}(i) each $\beta_{j+1}(h_j)$ differs from $b_{j+1}$ by at most $\delta'$, and by \eqref{eq:one-loop-window} the sum of the $b_{j+1}$ differs from $2c(n - m)\ln L$ by at most $C_W$. Hence
\begin{equation}\label{eq:beta-window}
    b(n - m) - C_W \leq \sum_{j=m}^{n-1}\beta_{j+1}(h_j) \leq b'(n - m) + C_W\;.
\end{equation}

\emph{The projected flow.} The set $H_k$ is closed, convex, and nonempty, since it contains $(u_0, \dots, u_k)$. Let $P_k$ be the Euclidean projection of $\mathbb R^{k+1}$ onto it, which does not increase distances and is therefore continuous. For $x_0 \in \mathbb R$ we define $y_0 \coloneqq x_0$ and
\begin{equation}\label{eq:projected-flow}
    y_{k+1} \coloneqq y_k - \beta_{k+1}\bigl(P_k(y_0, \dots, y_k)\bigr)
\end{equation}
for $0 \leq k < K$. This sequence is defined for every $x_0$, whether or not it stays in the tube, and each $y_k$ is a continuous function of $x_0$, by induction on $k$ and Proposition~\ref{prop:coefficients}(iii). Moreover the arguments of the coefficients in \eqref{eq:projected-flow} are histories, so that \eqref{eq:beta-window} applies and gives, for $0 \leq m < n \leq K$,
\begin{equation}\label{eq:projected-window}
    b(n - m) - C_W \leq y_m - y_n \leq b'(n - m) + C_W\;.
\end{equation}

\emph{The bare coupling.} With $m = 0$ and $n = K$, \eqref{eq:projected-window} says that $y_K$ lies between $x_0 - b'K - C_W$ and $x_0 - bK + C_W$. At $x_0 = \ell_0 = t + bK - C_W$ this gives $y_K \leq t$, and at $x_0 = u_0 = t + b'K + C_W$ it gives $y_K \geq t$. By the intermediate value theorem the continuous function $x_0 \mapsto y_K$ takes the value $t$ on $[\ell_0, u_0]$; its set of solutions there is closed, and we let $g_0(K, \bar g)^{-2}$ be the smallest of them.

\emph{The projection is inactive.} For this $x_0$, \eqref{eq:projected-window} with $n = K$ and $y_K = t$ gives $\ell_m \leq y_m \leq u_m$ for every $m < K$, by the definition \eqref{eq:tube-couplings} of $\ell_m$ and $u_m$, and \eqref{eq:projected-window} itself is the condition \eqref{eq:tube-window}. So $(y_0, \dots, y_k) \in H_k$ for every $k \leq K$, the projection $P_k$ leaves it fixed, and \eqref{eq:projected-flow} reads $y_{k+1} = y_k - \beta_{k+1}(y_0, \dots, y_k)$: the sequence is a flow, the densities $\varrho^h_k$ of its histories are those of the construction for the tube started at $g_0$, and its running couplings are $g_k = y_k^{-1/2}$, with $g_K = t^{-1/2} = \bar g$. The first part of \eqref{eq:coupling-window} is $\ell_k \leq y_k \leq u_k$, and the second follows from the lower bound of \eqref{eq:projected-window}, since $y_k - y_l \geq b(l - k) - C_W \geq -C_W$ for $k < l$.
\end{proof}

Three remarks are in order. First, the construction depends on the target: the parameters, the densities, the coefficients, and the bare coupling are those of the tube of $(K, \bar g)$, and objects of different targets are not identified. What is uniform are the constants. Secondly, by Lemma~\ref{lem:tube-couplings} the running couplings of the theorem satisfy Ba{\l}aban's inequalities (2.6), which he derives from the renormalisation group equations and uses in their place \cite[p.~255]{BalabanCRE}; the window \eqref{eq:coupling-window} is the inequality (0.31) of \cite{BalabanRG1} up to the additive constant $C_W$, with $\beta = b/\ln L$ and $\beta' = b'/\ln L$. Thirdly, the bare coupling window \eqref{eq:bare-coupling-window} of the introduction is \eqref{eq:coupling-window} at $k = 0$.

\section{Dressed Loops and the First Step}\label{sec:first-step}

Ba{\l}aban's construction acts on densities, and an observable has to be carried along with them. This section says in what form the loops of an ensemble are carried, and proves that the first step puts them into that form.

\subsection{Two Integrals}\label{sec:two-integrals}

Let $\mathcal C = ((C_1, \chi_1), \dots, (C_n, \chi_n))$ be a loop ensemble at level $q$, let $K > j_q$, and let $g_0 = g_0(K, \bar g)$ be the bare coupling of Theorem~\ref{thm:coupling-flow}. With $\varrho_0 \coloneqq \exp[-g_0^{-2}A - E]$ the bare density of Section~\ref{sec:density}, the expectation of the observable \eqref{eq:ensemble-observable} is the quotient
\begin{equation}\label{eq:two-integrals}
    \mathbb E_K[W_{\mathcal C}] = \frac{Z_K(\mathcal C)}{Z_K}\;,
    \qquad
    Z_K(\mathcal C) \coloneqq \int W_{\mathcal C}(U)\varrho_0(U)\,\dd U\;,
    \qquad
    Z_K \coloneqq \int\varrho_0(U)\,\dd U\;,
\end{equation}
in which the constant $E$ cancels. We apply the construction to the numerator and to the denominator separately. Each operation of the construction preserves the integral of a density: the block averaging $T$ does, because $\int\delta(\bar U(c)V(c)^{-1})\,\dd V(c) = 1$ for the Haar measure, and the $\mathbf R$-operation is defined so that ``integrals of the densities are unchanged'' \cite[p.~176]{BalabanLFR1}; the density it produces ``may not be equal to $\varrho_k$, but is equivalent to it, in the sense that they have equal integrals'' \cite[p.~200]{BalabanLFR1}. Hence $Z_K = \int\varrho_k(V_k)\,\dd V_k$ for every $k \leq K$, and for the numerator we shall construct densities $N_k$ with
\begin{equation}\label{eq:numerator-integral}
    Z_K(\mathcal C) = \int N_k(V_k)\,\dd V_k\;,
    \qquad
    N_0 \coloneqq W_{\mathcal C}\,\varrho_0\;,
\end{equation}
for $0 \leq k \leq K' = K - j_q$. No law of the block variables enters: the quotient of the two integrals at the step $K'$ is the expectation at the cutoff $K$. The point of the construction is the form of $N_k$, which is that of $\varrho_k$ with one more factor in each term.

\subsection{The First Step of the Construction}\label{sec:balaban-first-step}

The first step integrates the fine field with its block averages $V$ on the lattice of the step $1$ held fixed. Ba{\l}aban decomposes the integrand by test functions into terms labelled by two domains, the domain $\Omega_1$ on which the test functions of the plaquettes have passed and the domain $\Lambda_1 \subset \Omega_1$ on which the fluctuation field is small as well, and after the fluctuation integral on $\Lambda_1$ and the renormalisation of the coupling he obtains
\begin{equation}\label{eq:first-step-density}
    \varrho_1(V) = \sum_{\{ \Omega_1, \Lambda_1 \}}\chi_1(\Omega_1)\,\mathbf T_1(\{ \Omega_1, \Lambda_1 \})\exp A_1\Bigl(\frac{1}{g_1^2}, U_1\Bigr)
\end{equation}
\cite[(1.28), p.~253]{BalabanCRE}. Here $\mathbf T_1$ is the operation which holds the integrals over the bare field $V_0$ on $\Omega_1^c$ and over the fluctuation field $A$ on $\Omega_1 \cap \Lambda_1^c$, with the Gaussian weight of the step and the characteristic functions $\chi(\Omega_1 \cap \Lambda_1^c, S_1)$ of the large fluctuations \cite[(1.29), p.~254]{BalabanCRE}, and the effective action is
\begin{equation}\label{eq:first-step-action}
\begin{split}
    A_1\Bigl(\frac{1}{g_1^2}, U_1\Bigr) &= -A\Bigl(\frac{1}{g_1^2}, U_1\Bigr) + \bigl\{ \mathbf E^{(1)}(\Lambda_1, g_0, U_1) - \mathbf E^{(1)}(\Lambda_1, g_0, 1) - \beta_1(g_0)A(\phi_1, U_1) \bigr\}\\
    &\quad + \mathbf B^{(1)}(U_1, A, S_1) - E_1
\end{split}
\end{equation}
\cite[(1.30), p.~254]{BalabanCRE}, with the coupling of \cite[(1.27), p.~253]{BalabanCRE}. These integrals stay pending inside $\varrho_1$: Ba{\l}aban leaves the operations $\mathbf T_k$ of the later steps ``undefined, only their basic general properties are formulated'' \cite[p.~254]{BalabanCRE}. A \emph{labelled term} of $\varrho_1$ is a term of \eqref{eq:first-step-density} with given labels. For the numerator we use the refinement of the labels of \cite[Section~9]{WilsonFirstStep}, in which the domain $\Lambda_1$ is further restricted near the loop by a retreat of the carrier of the loop's integral; the refined labels are again labels of Ba{\l}aban's form, and on each of them the operation $\mathbf T_1$ and the action $A_1$ are those of \eqref{eq:first-step-density} (Supplement, Proposition~S.7).

In the construction of this paper the thresholds of the first step are those of the tube at $\hat g_0$, and the test functions of the plaquettes are soft (Definition~\ref{def:construction}). The first step with the loop inserted is the subject of the companion manuscript \cite{WilsonFirstStep}, for one square loop, Ba{\l}aban's thresholds at $g_0$, and characteristic functions as test functions, and for two disjoint squares in \cite[Section~10]{WilsonFirstStep}. We use its results for a general ensemble and for the construction of this paper, and the changes this requires are proved in Section~\ref{sec:first-step-proof}.

\subsection{Dressed Loops}\label{sec:dressed-loops}

Let us describe the form in which the ensemble is carried, beginning with a heuristic. After $k$ steps the fine field has been integrated out down to the scale $k$, and a loop $C_i$ is a path of $P_k(C_i)$ bonds of that scale. If the block field were exactly flat, the loop would contribute to each term of the density a number, the product $\kappa^{(k)}(C_i, \chi_i)$ of the numbers which the steps $1, \dots, k$ extract from it; since it is not, the loop also contributes the Wilson loop of the background field, corrected by insertions along the loop which are small where the background is close to flat and decay away from the regions of large fields. The insertions interact with each other through the excluded volume of the polymer expansions which produced them, and that interaction is resummed into a factor attached to each family of compatible insertions. The precise form is as follows.

Fix $k \geq 1$ and a labelled term of $\varrho_k$, with the labels of the construction at every scale. A \emph{traversal} at the step $k$ is an oriented bond of the lattice of the scale $k$ which is contained in one of the loops $C_i$, traversed in the orientation of $C_i$; by Lemma~\ref{lem:stopping-index} each loop is a union of $P_k(C_i) = P_q(C_i)a_k$ traversals. For a field $U$ on the fine lattice and an arc $J$ of $C_i$, that is, a set of consecutive traversals, $U(J)$ denotes the product of the transports along the fine bonds of $J$. The objects which make up a dressed loop are the following.

\begin{enumerate}[label=(\alph*)]
\item \emph{Groups.} A group $W$ is made at a step $j \leq k$. It consists of a nonempty finite set of unions of cubes of $\pi_j$, its \emph{localisation domains}, and of a set $\mathcal I(W)$ of pairwise disjoint arcs of the loops, its \emph{intervals}; the set $\mathbb L(W)$ of the loops which carry an interval of $W$ is nonempty. Its \emph{domain} $X_W$ is the union of its localisation domains and of the cubes of $\pi_j$ which meet a fine bond of its intervals. Its \emph{size} $\sigma(W) \geq 1$ is a number which is fixed at the step $j$, and we write $\sigma_j(W)$ when the step is to be shown: for $j \geq 2$ it is $\sigma_j(W) \coloneqq \#_j(X_W) + n_j(W)/M$, where $\#_j(X_W)$ is the number of the cubes of $\pi_j$ in $X_W$ and $n_j(W)$ the number of the traversals of the step $j$ which meet its intervals, and for $j = 1$ it is the size which \cite{WilsonFirstStep} gives to a group. A group is a \emph{cross group} if $\abs{\mathbb L(W)} \geq 2$. A group which is not a cross group has one interval, its \emph{arc}; a cross group may have several intervals on one loop, and its intervals are also called its arcs. The groups of a term carry a symmetric relation, their \emph{compatibility}, for which two groups whose intervals share a bond are incompatible, and a \emph{family} is a finite set $\mathfrak B$ of pairwise compatible groups.
\item \emph{Insertions.} Each group $W$ carries an \emph{insertion} $Y_W$, an element of $\bigotimes_{I \in \mathcal I(W)}\operatorname{End}(V_{\pi_{i(I)}})$, where $C_{i(I)}$ is the loop which contains the interval $I$ and $V_{\pi_i}$ is the space of the representation $\pi_i$; for a group with one interval it is a linear map of $V_{\pi_i}$. It is a function of the block field $V_k$ and of the outer variables of the term, the variables which $\mathbf T_k$ holds pending. Its norm is the projective norm
\begin{equation}\label{eq:projective-norm}
    \norm{Y} \coloneqq \inf\Bigl\{ \sum_\alpha\prod_{I \in \mathcal I(W)}\norm{A_{\alpha,I}} \,\Big|\, Y = \sum_\alpha\bigotimes_{I \in \mathcal I(W)}A_{\alpha,I} \Bigr\}\;,
\end{equation}
the infimum over the finite sums of elementary tensors, with the operator norms of the factors. For a group with one interval this is the operator norm of $Y$.
\item \emph{Words.} Let $\mathfrak B$ be a family whose insertions are elementary tensors, $Y_W = \bigotimes_IA_{W,I}$. Its word $w_k(\mathfrak B)$ is the product, over the loops $C_i$, of the normalised traces on $V_{\pi_i}$ of the ordered products along $C_i$ in which the transport $\pi_i(U_k(J))$ along every arc $J$ outside the intervals of $\mathfrak B$ alternates with the operators $A_{W,I}$ at the intervals $I \subset C_i$. This expression is linear in each factor $A_{W,I}$, and the word of a family with arbitrary insertions is defined from it by multilinearity. Here $U_k$ is the background field of the term, and on a large-field region crossed by a loop the transports also contain outer variables. For the empty family $w_k(\emptyset) = \prod_if^{(i)}_k$ with the \emph{leading words} $f^{(i)}_k \coloneqq \chi_i(U_k(C_i))$.
\item \emph{Interactions.} $Q_k(\mathfrak B) = \exp(-\sum_{\Gamma}\varphi_k(\Gamma))$, the sum over the clusters $\Gamma$ of a hard-core gas of groups, the \emph{reference gas} of the step, which are incompatible with a member of $\mathfrak B$, with numbers $\varphi_k(\Gamma)$.
\item \emph{The scalar exponent.} $\mathcal E_k = \sum_{j \leq k}\sum_Xe^{(j)}_X + E_{\mathrm{const},k}$, a sum of local terms $e^{(j)}_X$ indexed by connected unions $X$ of boxes of the lattice integrated at the step $j$, and a number $E_{\mathrm{const},k}$.
\end{enumerate}

The norm \eqref{eq:projective-norm} is the one in which the words are bounded, whatever the order in which the loops meet the intervals of the cross groups.

\begin{lemma}[Words and Norms]\label{lem:word-norm}
Let $\mathfrak B$ be a family, and let the transports along the arcs outside its intervals be unitary. Then $\abs{w_k(\mathfrak B)} \leq \prod_{W \in \mathfrak B}\norm{Y_W}$.
\end{lemma}

\begin{proof}
If the insertions are elementary tensors, the word is a product over the loops of normalised traces of ordered products of operators on the spaces $V_{\pi_i}$. Each of these traces is at most the product of the operator norms of its factors in modulus, and the transports have the norm one, so that $\abs{w_k(\mathfrak B)} \leq \prod_{W \in \mathfrak B}\prod_{I \in \mathcal I(W)}\norm{A_{W,I}}$. In general we write each $Y_W$ as a finite sum of elementary tensors, expand the word by multilinearity, apply this bound to every term, and take the infimum over the representations in \eqref{eq:projective-norm}.
\end{proof}

The operator norm of $Y_W$ on the tensor product of the spaces of its loops would not serve here. Two loops need not meet the intervals of several cross groups in the same cyclic order, and the word is then not the trace of a product of operators on a tensor product: for three transpositions of the two factors of $V \otimes V$, met by the two loops in opposite cyclic orders, it equals $\dim V$, although a transposition has the operator norm one. The bounds on the insertions of cross groups in the sequel are bounds on sums and integrals of elementary tensors by the sums and integrals of the products of the norms of their factors, and therefore bounds in the norm \eqref{eq:projective-norm}; a product of insertions and transports along the intervals is bounded by the product of the norms, since $\norm{AB} \leq \norm{A}\norm{B}$ on each factor of an elementary tensor.

The bounds of the class are stated with the following quantities. The small parameters of the step $k$ are, with the radius $r_k \coloneqq A_1(\log g_k^{-2})^{p_0}$ of the fluctuation cutoff in units of the coupling,
\begin{equation}\label{eq:class-parameters}
    \tau_k \coloneqq C_\tau\bigl[g_k(1 + r_k)^3 + e^{-c_\tau r_k^2}\bigr]\;,
    \qquad
    \rho_k \coloneqq C_\rho g_k^2(1 + r_k)^4\;,
    \qquad
    \rho_{t,k} \coloneqq C_tg_k(1 + 2r_k)\;,
\end{equation}
with constants fixed in Appendix~\ref{app:constants}. For a traversal $\mathfrak n$ and $t \geq 1$ let $\mathcal P_t(\mathfrak n)$ be the set of the fine plaquettes of $\Omega_k$ within distance $C_{\mathrm{box}}M(t + 1)$ of $\mathfrak n$, the distance measured in units of the scale $k$. The \emph{deviation} of the background near $\mathfrak n$ is
\begin{equation}\label{eq:deviation-quantity}
    \delta^{*}_{\mathfrak n}(U) \coloneqq L^{2k}\sum_{t \geq 1}(t + 1)^3e^{-\eta_wt}\max\bigl\{ \abs{U(\partial p) - 1} \,\big|\, p \in \mathcal P_t(\mathfrak n) \bigr\}\;,
\end{equation}
the maximum over the empty set being $0$; the factor $L^{2k}$ converts the deviation of a fine plaquette into that of a plaquette of the scale $k$. Finally $d_k(\mathfrak n)$ is the distance, in units of the scale $k$, from $\mathfrak n$ to the large-field data of the scale $k$ (the large-field region $Z_k$ of the term), and $N^{(k)}_{\min}$ is the smallest period of the torus in units of the scale $k$.

The deviation $\delta^*$ bounds a sum of insertions, and this is how the class enters most estimates. The inductive step, however, has to evaluate an insertion at a background other than the one at which it was bounded, and to sort the insertions according to the region which their bounds consult; a bound on a sum cannot be used for either, and the class therefore also records a bound for each insertion of the deviation species, with the region to which it refers. The quantities are the following; for a group $W$ they are formed at the step $j$ at which $W$ was made and in the units of the scale $j$, and they are not changed afterwards. The \emph{dependence set} $X^+_W$ of $W$ is the closed $\ell^\infty$-neighbourhood of its domain $X_W$ of radius $8M$ in the units of the fine lattice, which is less than the side of a cube of any of the partitions $\pi_j$, $j \geq 1$; and $\ell_W \coloneqq \#_j(X_W)$, the number of the cubes of $\pi_j$ in the domain, for $j \geq 2$, while $\ell_W \coloneqq 1$ for $j = 1$. Put $R_t \coloneqq C_{\mathrm{box}}M(t + 1)$ for $t \geq 0$. For a traversal $\mathfrak n$ of the step $j$ in an interval of $W$ let $t_{\mathfrak n}(W)$ be the least $t \geq 1$ for which $X_W$ lies in the $\ell^\infty$-ball $B_\infty(\mathfrak n, R_t)$ of radius $R_t$ about $\mathfrak n$. The \emph{comparison region} of $W$ at $\mathfrak n$ is the ball $B_\infty(\mathfrak n, R_{t_{\mathfrak n}(W)+1})$, which contains the dependence set $X^+_W$, and the \emph{deviation of the background at $W$} is
\begin{equation}\label{eq:insertion-deviation}
    \operatorname{Dev}_{\mathfrak n}(W; U) \coloneqq \ell_W\,L^{2j}\max_p\abs{U(\partial p) - 1}\;,
\end{equation}
the maximum over the fine plaquettes $p$ in the comparison region of $W$ at $\mathfrak n$. The \emph{extent} $\operatorname{ext}(W)$ is the least $R \geq 0$ such that $X_W$ and the comparison regions of $W$ at all the traversals of its intervals lie in the ball of radius $RM$ about each of these traversals. For a traversal $\mathfrak n$ of the step $k \geq j$ in an interval of $W$, let $t^{(k)}_{\mathfrak n}(W)$ be the least $t \geq 1$ such that the comparison regions of $W$ at the traversals of the step $j$ contained in $\mathfrak n$ lie within the distance $R_t$ of $\mathfrak n$ in the units of the scale $k$; since the units grow with the step, $t^{(k+1)}_{\mathfrak n'}(W) \leq t^{(k)}_{\mathfrak n}(W)$ for $\mathfrak n \subset \mathfrak n'$. Finally the weights of the tails are
\begin{equation}\label{eq:tail-weights}
    w_t \coloneqq (t + 1)^3e^{-\eta_wt}\;,
    \qquad
    W_t \coloneqq \sum_{t' \geq t}w_{t'}\;,
\end{equation}
so that $\delta^*_{\mathfrak n}(U) = L^{2k}\sum_{t \geq 1}w_t\max\{ \abs{U(\partial p) - 1} \,|\, p \in \mathcal P_t(\mathfrak n) \}$ in \eqref{eq:deviation-quantity}.

\begin{definition}[Dressed Loops]\label{def:dressed-loop}
Let $k \geq 1$. A \emph{dressed loop of the step $k$} for the ensemble $\mathcal C$ on a labelled term of $\varrho_k$ is a function of $V_k$ and of the outer variables of the term of the form
\begin{equation}\label{eq:dressed-loop}
    \mathcal W_k = \prod_{i=1}^{n}\kappa^{(k)}(C_i, \chi_i)\cdot e^{\mathcal E_k}\Bigl[\prod_{i=1}^{n}f^{(i)}_k + \sum_{\mathfrak B \neq \emptyset}Q_k(\mathfrak B)\,w_k(\mathfrak B)\Bigr]\;,
    \qquad
    \kappa^{(k)}(C, \chi) \coloneqq \prod_{j=1}^{k}\kappa_j(C, \chi)\;,
\end{equation}
the sum over the nonempty families, absolutely convergent, with the objects (a)--(e) and the step factors $\kappa_j$ of Definition~\ref{def:step-factor}, such that the following hold.
\begin{enumerate}[label=(\roman*)]
\item \emph{Interactions.} $\abs{Q_k(\mathfrak B)} \leq \prod_{W \in \mathfrak B}e^{\theta(W)}$ with $0 \leq \theta(W) \leq \nu\sigma(W)/16$, where $\sigma(W)$ is the size of $W$ at the step at which it was made; more precisely, the sum of $\abs{\varphi_k(\Gamma)}$ over the clusters $\Gamma$ of the reference gas which are incompatible with a group $W$ is at most $\theta(W)$.
\item \emph{Insertions.} For every traversal $\mathfrak n$, the sum $\nu^{(k)}_{\mathfrak n} \coloneqq \sum_{W \ni \mathfrak n}e^{\theta(W)}\norm{Y_W}$ over the groups of the term made at the step $k$, and over the groups of earlier steps which are not attached to a large-field region of their step, the groups of the cross species defined below being excepted in both cases, satisfies
\begin{equation}\label{eq:insertion-bound}
    \nu^{(k)}_{\mathfrak n} \leq C_\chi\tilde\rho_k\delta^{*}_{\mathfrak n}(U_k) + C\hat\rho_ke^{-c_{\mathrm{att}}d_k(\mathfrak n)/M} + C\rho_{t,k}\,\mathfrak m_k(\mathfrak n) + \sum_{j=1}^{k}CL^{k-j}\hat\rho_je^{-c_{\mathrm{att}}N^{(j)}_{\min}/M}\;,
\end{equation}
where $\tilde\rho_k \coloneqq \rho_k/\min(\alpha_{0,k}, \alpha_{1,k})$ for $k \geq 2$ and $\tilde\rho_1 \coloneqq \rho_1$, $\hat\rho_k \coloneqq C\rho_k$, and $\mathfrak m_k(\mathfrak n)$ is the sum of the marks of the large-field data of the step $k$ near the boundary of $Z_k$, weighted by $e^{-c'_*\dist(\cdot, \mathfrak n)/M}$. The \emph{marks} are the mark weights of \cite[Lemma~9.24]{WilsonFirstStep} with the parameters of the step $k$: a bond $b$ of the region of integration carries the mark $\mathfrak m_b$, which bounds the effect of the conditional mean on the cutoff of $b$ and satisfies $\mathfrak m_b \leq 2C_\varpi r_k^{\dim G}\sum_ee^{-\mu_Hd(b,e)}$, the sum over the bonds $e$ of the boundary of that region, and a \emph{host} $Y$ of the step with an input outside the region carries a mark $\mathfrak m_Y$, which is bounded by a constant of the construction. The hosts are not the localisation domains of the expansions, which are connected unions of arbitrarily many cubes: a host is the elementary box in which one local primitive of the step is evaluated, a box of cubes of side $M$ with at most four cubes along each side \cite[(3.36), (3.75), Lemma~B.10]{WilsonFirstStep}, and only hosts carry marks. A mark of a bond is assigned to the cube $\Delta$ of side $LM$ of the scale $k$ which contains the bond, and a mark of a host to the cube $\Delta$ which contains the lowest cube of the host; $\hat{\mathfrak m}_\Delta$ is the sum of the marks assigned to $\Delta$, and
\begin{equation}\label{eq:marks}
    \mathfrak m_k(\mathfrak n) \coloneqq \sum_\Delta\hat{\mathfrak m}_\Delta e^{-c'_*\dist(\Delta, \mathfrak n)/M}\;,
    \qquad
    \hat{\mathfrak m}_\Delta \leq C_{\mathfrak m}(1 + r_k)^{\dim G}\;.
\end{equation}
The second holds because a cube of side $LM$ has at most $8(LM + 1)^4$ bonds, each with a mark of at most $2C_\varpi r_k^{\dim G}$ times a convergent sum, and because the hosts assigned to it are at most $4^4L^4$ in number, a host being determined by its lowest cube, one of the $L^4$ cubes of side $M$ of $\Delta$, and by its four side lengths. For general localisation domains no such count holds, and none is used: the localisation domains of the step enter the bounds of the insertions through the sums of their activities, with their decay, and not through marks. This bound per cube is the one used when the attached parts are paid (Section~\ref{sec:large-field-step}). The groups attached at an earlier step $j < k$ to a component $Z$ of the large-field region of the step $j$ satisfy, summed over all traversals, the bound $C\hat\rho_j\sum_{\mathfrak n_j}e^{-c_{\mathrm{att}}d^Z_j(\mathfrak n_j)/M}$, the sum over the traversals of the step $j$, and the corresponding bound with $\rho_{t,j}$ for those attached to the boundary layer of $Z$. The groups of the \emph{deviation species}, whose insertions vanish at pure-gauge data on their arcs, satisfy the first term of \eqref{eq:insertion-bound} with weights: $\sum_{W \ni \mathfrak n}e^{\theta(W) + \nu\sigma(W)}\norm{Y_W} \leq C_\chi\tilde\rho_k\delta^{*}_{\mathfrak n}(U_k)$, the sum over these groups, with $\sigma(W)$ the size of $W$ at the step at which it was made. They satisfy it in the following form, of which it is a consequence (Lemma~\ref{lem:tail-to-sum}). The comparison regions of a group $W$ of the deviation species lie in $\Omega_k$. For each such group there is a number $b_W \geq 0$, fixed at the step $j$ at which $W$ was made, such that for every traversal $\mathfrak n$ of the step $j$ in an interval of $W$
\begin{equation}\label{eq:per-insertion}
    e^{\theta(W)}\norm{Y_W} \leq b_W\operatorname{Dev}_{\mathfrak n}(W; U_k)\;,
\end{equation}
and, with $\beta_W \coloneqq b_W\ell_WL^{2j}$, for every traversal $\mathfrak n$ of the step $k$ and every $t \geq 1$,
\begin{equation}\label{eq:tail-clause}
    \sum_{W \ni \mathfrak n,\ t^{(k)}_{\mathfrak n}(W) \geq t}e^{\nu\sigma(W) + \nu_e\operatorname{ext}(W)}\beta_W \leq C_\chi\tilde\rho_kL^{2k}W_t\;,
\end{equation}
the sum over the groups of the deviation species with $\mathfrak n$ in an interval; $\sigma(W)$ and $\operatorname{ext}(W)$ are those of the step at which $W$ was made. The \emph{cross species} consists of the cross groups which the step map or the $\mathbf R$-operation makes from groups none of which is attached to a large-field region or wraps around the torus, and of their descendants. Its insertions carry no deviation bound, they need not vanish at the reference data, and they are not contracted. For every $j \leq k$, the groups of the cross species present at the step $k$ which descend from those made at the step $j$ satisfy
\begin{equation}\label{eq:cross-bound}
\begin{gathered}
    \sum_We^{\theta(W) + \lambda_k(W)\nu\sigma(W)}\norm{Y_W} \leq 2^{k-j+1}\mathsf A^\times_j\;,\\
    \mathsf A^\times_j \coloneqq P_j(\mathcal C)\Bigl[C_{\mathrm{cr}}\mathfrak a_je^{-\nu_\times d_j/M} + \mathbf 1_{\mathbf R}(j)\,C^{\mathbf R}_{\mathrm{cr}}e^{-p_0(g_j)/4}e^{-\kappa d_j/(6MR_j)}\Bigr]\;,
\end{gathered}
\end{equation}
with rates $\lambda_k(W) \in [\frac{15}{16}, 1]$ and the decay rate $\kappa$ of Ba{\l}aban's polymers at the $\mathbf R$-steps (Supplement, Proposition~S.10), where $d_j$ is the least $\ell^\infty$-distance between two loops of the ensemble in units of the scale $j$, so that $d_j \geq a_j$ by Lemma~\ref{lem:stopping-index}, $\mathfrak a_j$ is the amplitude of Definition~\ref{def:amplitude}, and $\mathbf 1_{\mathbf R}(j)$ is $1$ at the steps at which the $\mathbf R$-operation acts and $0$ at the others; for a single loop the cross species is empty. A cross group one of whose parents is attached or wraps around the torus belongs to the species of the first such parent and satisfies the bounds of that species.
\item \emph{Analyticity.} The insertion of a group $W$ depends on the background field only through its restriction to the dependence set $X^+_W$. For real $U_k$ in the support of the indicators of the term, the factors which depend on the background extend analytically to $\mathbf U = e^{iL^{-k}A'}U_k$, for complex $A'$ in the polydisc of the radii $(1 - \beta(1 - 2^{-(k-m)}))\alpha_{1,m}L^{-m}$ per fine bond on $\Omega_m \setminus \Omega_{m+1}$ (and on $\Omega_k$ for $m = k$), and for $k \geq 2$ with the gradient condition of \cite[(2.39), p.~261]{BalabanCRE} and in the variable $\mathbf J$ of \cite[(2.36)]{BalabanCRE}; they are covariant under $G$-valued gauge transformations, and \eqref{eq:insertion-bound} holds there with the deviation of $\mathbf U$ from $U_k$ added to $\delta^{*}_{\mathfrak n}$. The same factors are defined, and analytic in the same polydiscs, at the \emph{reference data} of the term, the trivial field $1$ in place of $U_k$, with $\mathbf J = 0$ and the outer variables trivial, whether or not the reference data lie in the support of the indicators. On the term without large-field regions the domain and the functions are the following, in the units of the fine lattice. For a configuration $U$ with values in $G$ and a field $\alpha$ on the fine bonds with values in the complexified Lie algebra $\mathfrak g^c$ we write $e^{i\alpha}U$ for the configuration $b \mapsto e^{i\alpha(b)}U(b)$. For a bond $b = \langle x, x + e_\mu\rangle$ and a unit vector $e$, one of the $\pm e_\nu$, let $b + e \coloneqq \langle x + e, x + e + e_\mu\rangle$, and let the \emph{covariant difference} of $\alpha$ be
\begin{equation}\label{eq:covariant-difference}
    (\delta_e\alpha)(b) \coloneqq U(x, x + e)\,\alpha(b + e)\,U(x, x + e)^{-1} - \alpha(b)\;,
\end{equation}
where $U(x, x + e)$ is the value of $U$ on the bond from $x$ to $x + e$. For $n \geq 2$ let $\mathcal S_n(U)$ be the set of the pairs $(e^{i\alpha}U, \mathbf J)$, with fields $\alpha$ and $\mathbf J$ on the fine bonds with values in $\mathfrak g^c$, such that for all bonds $b$, unit vectors $e$, and fine plaquettes $p$
\begin{equation}\label{eq:space-of-step}
\begin{gathered}
    \abs{\alpha(b)} < \alpha_{1,n}L^{-n}\;,
    \qquad
    \abs{(\delta_e\alpha)(b)} < \alpha_{1,n}L^{-2n}\;,\\
    \abs{(e^{i\alpha}U)(\partial p) - 1} < c_F\,\alpha_{0,n}L^{-2n}\;,
    \qquad
    \abs{\mathbf J(b)} < \alpha_{0,n}L^{-3n}\;,
\end{gathered}
\end{equation}
with a constant $c_F \in (0, 1]$; $\mathcal S_1(U)$ is the set of the configurations $e^{i\alpha}U$ with $\abs{\alpha(b)} < \alpha_{1,1}L^{-1}$, without a variable $\mathbf J$. In the units of the scale $n$ the fine lattice has the spacing $\xi = L^{-n}$ and the complex part is written $e^{i\xi A'}$, so that $\alpha = \xi A'$, and \eqref{eq:space-of-step} are the conditions $\abs{A'}, \abs{\nabla^\xi_UA'} < \alpha_{1,n}$, $\abs{\partial\mathbf U - 1} < c_F\alpha_{0,n}\xi^2$, and $\abs{\mathbf J} < \alpha_{0,n}$ of \cite[(2.39), (2.34), (2.36), p.~261]{BalabanCRE} on the top domain, the second with the constant $c_F$; a current has the dimension of the third power of an inverse length \cite[(1.8), p.~261]{BalabanRG1}, whence the factor $L^{-3n}$. The constant $c_F$ is so small that these conditions imply those on the averaged configurations, \cite[(2.35), (2.37)]{BalabanCRE}: Ba{\l}aban states that the conditions (i)--(iii) of his spaces with smaller constants imply the condition (iv), by the bounds \cite[(1.17), p.~263]{BalabanRG1}, which he derives in one line from \cite[Proposition~9]{BalabanVAR}. The \emph{space of the step $k$} is the union of the sets $\mathcal S_k(U)$ over the configurations $U$ which are the trivial configuration, the background $U_k(V_k)$ of a block field $V_k$ in the support of the test functions of the term, or a transform of one of these by a $G$-valued gauge transformation. For a set $X$ of bonds let $\mathcal S_n(U)|_X$ be the set defined by the conditions \eqref{eq:space-of-step} on the bonds of $X$ and on the plaquettes and differences which they contain. The insertion of a group $W$ made at a step $j \leq k$ is a holomorphic function of the restriction of $(\mathbf U, \mathbf J)$ to $X^+_W$, of $\mathbf U$ alone for $j = 1$, on the open set
\begin{equation}\label{eq:production-space}
    \mathcal D_W \coloneqq \bigcup_U\mathcal S_j(U)|_{X^+_W}\;,
\end{equation}
the union over the same configurations $U$ at the step $j$, and at the later steps the insertion of $W$ is the same function of the configuration and of the current, both in fine units. For a group with one interval, from the site $x_W$ to the site $x'_W$, and for $G$-valued gauge transformations $v$,
\begin{equation}\label{eq:covariance-law}
    Y_W(\mathbf U^v, R(v)\mathbf J) = \pi(v(x_W))\,Y_W(\mathbf U, \mathbf J)\,\pi(v(x'_W))^{-1}\;,
\end{equation}
where $\mathbf U^v(b) = v(b_-)\mathbf U(b)v(b_+)^{-1}$ and $(R(v)\mathbf J)(b) = v(b_-)\mathbf J(b)v(b_-)^{-1}$ \cite[(1.10), p.~262]{BalabanRG1}, and $\pi$ is the representation of the loop of $W$; for a group with several intervals each factor of the tensor is transformed in this way.
\item \emph{Scalar parts.} $Q_k(\mathfrak B)$ and $\mathcal E_k$ are numbers on each labelled term, independent of $V_k$ and of the outer variables, and
\begin{equation}\label{eq:scalar-bounds}
    \sum_{X \ni x}\abs{e^{(j)}_X}e^{\kappa_e\abs{X}} \leq C\tau_j^2e^{-c_{\mathrm{th}}d^{(j)}(x)}\;,
    \qquad
    \abs{E_{\mathrm{const},k}} \leq \sum_{j \leq k}\epsilon_j(\mathcal C, T)\;,
\end{equation}
where $d^{(j)}(x)$ is the distance of the box $x$ of the scale $j - 1$ from the complement of the carrier of the loop's integral at the step $j$, and $\epsilon_j(\mathcal C, T)$ is the remainder bound of the step $j$ defined in \eqref{eq:step-remainder}.
\item \emph{Uniformity.} The constants depend on the constants of the construction, the characters, and the shape constants $\Lambda(C_i)$, and not on $k$, $K$, the torus, the labelled term, or the number of loops; for an ensemble with admissible loops they may be taken to be nondecreasing functions of $c_{\mathcal C}$ which depend on $\Lambda_*$ in place of the $\Lambda(C_i)$, because $\Lambda(C_i)$ enters only through the counts of Lemma~\ref{lem:square-properties}, which increase with it, and a character only through the number $c_{\pi_i}$ of \eqref{eq:representation-constant} (Section~\ref{sec:results}).
\item \emph{Torus consistency.} On the trivial term and at its reference data, let $T \subset T'$ be two tori of the family $\mathfrak T$ of Section~\ref{sec:reference-gas}. For every insertion and every cluster weight $\varphi_k(\Gamma)$ of an object whose dependence set has diameter less than $\iota_{T,k}/4$, the values on $T$ and on $T'$, at fields which coincide under the covering on the $\iota_{T,k}/4$-neighbourhood of a lift of that set, differ by amounts whose sums rooted at a traversal are at most $C_{\mathrm{rel}}e^{-\iota_{T,k}/(32M)}$ times the right side of \eqref{eq:insertion-bound} for the insertions, and at most $C_{\mathrm{rel}}e^{-\iota_{T,k}/(32M)}\mathfrak a_k$ for the cluster weights.
\item \emph{Geometry and sizes.} For every group $W$ of the term, made at a step $j \leq k$, the domain $X_W$ is connected, in the sense that any two of its cubes are joined by a chain of its cubes in which consecutive ones have a point in common, and $\#_j(X_W) + n_j(W)/M \leq u_\star\sigma(W)$ with a constant $u_\star \geq 1$; and the domains of two incompatible groups have a point in common. On the term without large-field regions the insertions satisfy in addition two bounds in which their sizes on the whole domain of analyticity enter. Let $\norm{Y_W}_{(k)}$ be the supremum of $\norm{Y_W}$ over the \emph{space of the step $k$}, the set of the configurations $(\mathbf U, \mathbf J)$ of (iii) around the real backgrounds in the support of the indicators of the term and around its reference data, and say that a traversal of the step $k$ is \emph{in} an interval of a group if it contains a fine bond of that interval. Then, for every traversal $\mathfrak n$ of the step $k$,
\begin{equation}\label{eq:sup-clause}
    \sum_{W \ni \mathfrak n}e^{\theta(W) + \lambda_c\nu\sigma(W) + \nu_e\operatorname{ext}(W)}\norm{Y_W}_{(k)} \leq C_c\,\rho_k\;,
    \qquad
    \lambda_c \coloneqq \frac{15}{16} - \frac{1}{8192}\;,
\end{equation}
the sum over the groups of the deviation species with $\mathfrak n$ in an interval. Moreover, for a group $W$ made at the step $k$ let $c_W$ be the number subtracted from its insertion at its production (Section~\ref{sec:small-field-map}), which is $0$ if nothing was subtracted, and $S_W \coloneqq \norm{Y_W}_{(k)} + \abs{c_W}$; then
\begin{equation}\label{eq:production-clause}
    \sum_{W \text{ at } p}e^{\theta(W) + \nu\sigma(W) + \nu_e\operatorname{ext}(W)}S_W\,\ell_W \leq 2C_{S'}\rho_k
\end{equation}
for every traversal $p$ of the step $k$ and every cube $p$ of $\pi_k$, the sum over the groups made at the step $k$ which are \emph{at} $p$: those with $p$ in an interval, or with $p \subset X_W$.
\end{enumerate}
The ensemble is \emph{dressed at the step $k$} if $N_k = \sum\chi_k(\Omega_k)\mathbf T_k[\mathcal W_k\exp A_k]$, the sum over the labelled terms of $\varrho_k$, with a dressed loop $\mathcal W_k$ of the step $k$ on each of them.
\end{definition}

\begin{lemma}[From the Tails to the Sum]\label{lem:tail-to-sum}
Let $\mathfrak n$ be a traversal of the step $k$, and let $U$ be a configuration at which the groups $W$ of the deviation species with $\mathfrak n$ in an interval satisfy \eqref{eq:per-insertion}, with $U$ in the place of $U_k$, and \eqref{eq:tail-clause}. Then
\begin{equation}\label{eq:weighted-deviation-sum}
    \sum_{W \ni \mathfrak n}e^{\theta(W) + \nu\sigma(W) + \nu_e\operatorname{ext}(W)}\norm{Y_W} \leq C_\chi\tilde\rho_k\delta^{*}_{\mathfrak n}(U)\;.
\end{equation}
\end{lemma}

\begin{proof}
Let $m_t$ be the maximum of $\abs{U(\partial p) - 1}$ over $p \in \mathcal P_t(\mathfrak n)$ for $t \geq 1$, and $m_0 \coloneqq 0$; the sets $\mathcal P_t(\mathfrak n)$ increase with $t$, so that $m_t$ is nondecreasing. Let $W$ be one of the groups, made at the step $j$, and $\mathfrak n_j \subset \mathfrak n$ a traversal of the step $j$ in its interval. Its comparison region at $\mathfrak n_j$ lies in $\Omega_k$ and within the distance $R_t$ of $\mathfrak n$ in the units of the scale $k$, for $t = t^{(k)}_{\mathfrak n}(W)$, so that its fine plaquettes belong to $\mathcal P_t(\mathfrak n)$, and by \eqref{eq:per-insertion} and \eqref{eq:insertion-deviation}
\[
    e^{\theta(W)}\norm{Y_W} \leq b_W\ell_WL^{2j}\,m_t = \beta_W\,m_{t^{(k)}_{\mathfrak n}(W)}\;.
\]
We write $m_t = \sum_{t' \leq t}(m_{t'} - m_{t'-1})$, sum over the groups, and use \eqref{eq:tail-clause} for each $t'$:
\[
    \sum_{W \ni \mathfrak n}e^{\theta(W) + \nu\sigma(W) + \nu_e\operatorname{ext}(W)}\norm{Y_W} \leq \sum_{t' \geq 1}(m_{t'} - m_{t'-1})\,C_\chi\tilde\rho_kL^{2k}W_{t'} = C_\chi\tilde\rho_kL^{2k}\sum_{t \geq 1}m_tw_t\;,
\]
where the last step is a summation by parts, with $W_{t} - W_{t+1} = w_t$ and $m_tW_{t+1} \to 0$, the numbers $m_t$ being bounded. The right side is $C_\chi\tilde\rho_k\delta^{*}_{\mathfrak n}(U)$ by \eqref{eq:deviation-quantity} and \eqref{eq:tail-weights}.
\end{proof}

Three comments explain the definition. First, the bound \eqref{eq:insertion-bound} separates two kinds of insertions. The first term, the \emph{deviation part}, is small because the background is close to flat; it vanishes when $U_k$ is a pure gauge, by Schur's lemma, which is why it is proportional to $\delta^{*}$. The other terms are \emph{attached} parts, which carry no deviation and are instead small through the decay away from the large-field regions, the boundary layers, and the wrapping of the torus. For several loops there is a third kind, the insertions of the cross species: the value of such an insertion at flat data is invariant under the simultaneous conjugation of its factors by $G$ only, and it is not a scalar in general, so that nothing can be subtracted from it, and it is small because a cross group joins two loops which are far apart in the units of every step, which is the content of \eqref{eq:cross-bound}. Secondly, no smallness of $\sum_W\norm{Y_W}$ over a whole loop is asserted: the sum over families is controlled by the product $\prod_W(1 + e^{\theta(W)}\norm{Y_W})$, whose logarithm is a sum over the traversals of \eqref{eq:insertion-bound}, so that it costs a factor which is exponential in the number of traversals with a small rate. Thirdly, the step factors $\kappa_j$ are defined in Section~\ref{sec:flat-part} by an explicit formula, and since the normalisation in \eqref{eq:dressed-loop} is required to be exactly $\prod_j\kappa_j$, the class is an induction for the normalisation of Definition~\ref{def:renormalised-expectation}.

\subsection{The First Step with a Loop Ensemble}\label{sec:first-step-proof}

\begin{theorem}[The First Step]\label{thm:first-step}
Let $L$ and $\Lambda_*$ be as in Theorem~\ref{thm:main}, let $\mathcal C$ be a loop ensemble at level $q$ whose loops are admissible, and let $\bar g$ and $q$ be as in Theorem~\ref{thm:main} for $\mathcal C$. For every $K$ larger than an index $K_{\mathrm{tor}}(\mathcal C)$, which depends on $\mathcal C$ only through $\max_iP_q(C_i)$ and on $q$ and $T/u$, the density $N_1$ obtained from $N_0 = W_{\mathcal C}\varrho_0$ by the first step of the construction for the target $(K, \bar g)$ satisfies \eqref{eq:numerator-integral}, and the ensemble is dressed at the step $1$.
\end{theorem}

The proof uses the first-step manuscript \cite{WilsonFirstStep} for the analysis of one square loop and of two squares, and consists in checking that each of its statements holds for a general ensemble and for the construction of this paper. We first collect the properties of a general loop that replace those of a square in that analysis; they are consequences of Section~\ref{sec:loops}. Throughout, $C$ is a loop of $\mathcal C$, read in fine units on the lattice $\mathbb T_K$, with $P_0(C)$ bonds, and $\ell$ is a length which is at most a quarter of the separation $a_0 \geq L^3MR_{K'}L^{K'}$ of Lemma~\ref{lem:stopping-index}.

\begin{lemma}[Square Properties of a General Loop]\label{lem:square-properties}
Let $C$ be a loop with chord--arc constant $\Lambda(C)$.
\begin{enumerate}[label=(\roman*)]
\item The bonds of $C$ which meet a closed $\ell^\infty$-cube of side $s$ are at most $\Lambda(C)\cdot 2(s + 1)$, and the sites of $C$ in a set of $\ell^\infty$-diameter $D$ are at most $2\Lambda(C)(D + 1) + 1$.
\item For a set $S$ of $\ell^\infty$-diameter $D$ which meets $C$, the smallest arc of $C$ which contains $S \cap C$ has at most $4\Lambda(C)D + 2$ bonds, or else contains more than half of $C$.
\item A set of $\ell^\infty$-diameter less than $r_0(C)$ which meets $C$ meets at most two consecutive segments of $C$, and on two consecutive segments the chord--arc constant is $2$.
\end{enumerate}
\end{lemma}

\begin{proof}
(i) Two sites $x, y$ of $C$ in the cube have $\abs{x - y}_\infty \leq s$, hence, by Definition~\ref{def:chord-arc}, arc distance at most $\Lambda(C)s$. The sites of $C$ in the cube therefore lie on an arc of at most $\Lambda(C)s$ bonds measured from any one of them in either direction, that is, on at most $2\Lambda(C)s + 1$ sites, and a bond which meets the cube has an end among them or is one of the two bonds at each end; this gives at most $2\Lambda(C)(s + 1)$ bonds for $s \geq 1$, and the count of sites in a set of diameter $D$ is the same argument. (ii) is Lemma~\ref{lem:arcs-around-set}, and (iii) is Lemma~\ref{lem:corner-isolation} together with Lemma~\ref{lem:chord-arc} applied to the path formed by two consecutive segments, which is a bent path with two perpendicular arms.
\end{proof}

\begin{proof}[Proof of Theorem~\ref{thm:first-step}]
The identity \eqref{eq:numerator-integral} at $k = 1$ holds because the first step preserves integrals (Section~\ref{sec:two-integrals}). The rest of the proof goes through the analysis of \cite{WilsonFirstStep} in the order of that paper and states, for each part, the change which a general ensemble or the construction of this paper requires.

\emph{The insertion at a fixed cutoff and the extraction.} The sourced effective interaction of \cite[Sections~2--3]{WilsonFirstStep} is built from the fine bonds of the loop and the local operators of the step, whose ranges are fixed and do not depend on the loop. The loop enters through three counts: the hull of a cluster on the loop, the number of labels which meet a bond, and the number of traversals within a given distance of a point. For a square these are bounded by the elementary geometry of the square; for $C$ they are bounded by Lemma~\ref{lem:square-properties}: the hull constant $32(1 + R_*)$ of \cite[(5.24)]{WilsonFirstStep} becomes $16\Lambda(C)(1 + R_*)$ by (ii), and the number of source labels of one traversal which meet a bond, which is at most a constant $q^{\mathrm{src}}$ for each of the traversals whose initial sites lie in a set of diameter $2R_* + 2$, is at most $q^{\mathrm{src}}(4\Lambda(C)(R_* + 1) + 1)$ by (i). The constants which depend on these counts are fixed before the coupling, and the conditions on the coupling which contain them depend on $\Lambda(C) \leq \Lambda_*$. Every condition on the coupling in this proof which involves the loops is a condition on $g_0$ or on $g_1$, both of which are at most $\gamma_q$ by \eqref{eq:stopping-coupling}, so that it is a requirement on $\gamma_q$, depending on $c_{\mathcal C}$.

\emph{The coefficients.} The perimeter and corner coefficients of \cite[Sections~5--6]{WilsonFirstStep} are defined on the straight line and on the bent path of two perpendicular rays through the coarse sites, in the orientations of the square. We define them for every orientation $v \in \mathcal V \coloneqq \{ \pm e_0, \dots, \pm e_3 \}$ on the line, and for every ordered pair $(v, w)$ of perpendicular orientations on the bent path \eqref{eq:bent-path}; the proofs of their convergence and of their tail bounds use only that the path is a line, or two perpendicular rays, through coarse sites, and not which coordinate directions it uses, so they hold with the same constants. The extraction of the regional scalar of the loop is then the following: allocate each cluster of the reference gas to the coarse segment of its source, give each corner $c$ the clusters allocated to the second half of its incoming segment and to the first half of its outgoing segment, and compare the clusters rooted there with those of the bent path $B_c$. By Lemma~\ref{lem:corner-isolation} a cluster which distinguishes $C$ from $B_c$ near $c$ reaches a distance at least $r_0(C) - L/2 - 2R_*$ from its root, and by the tail bounds of the cluster expansion such clusters cost at most $C\tau_0^2e^{-cr_0(C)/L}$ per root. With $\kappa_C$ the exponential of the sum of the line coefficients over the segments and of the bent-path coefficients over the corners, this gives
\begin{equation}\label{eq:first-step-extraction}
    \Bigl\lvert F_{C,\infty} - \log\kappa_1(C, \chi)\Bigr\rvert \leq C\tau_0^2\Bigl(\frac{P_0(C)}{L} + \abs{\mathcal K(C)}\Bigr)e^{-cr_0(C)/L}\;,
\end{equation}
in place of the square's bound, where $F_{C,\infty}$ is the bulk limit of the logarithm of the reference partition function of the loop and $\kappa_1$ is the step factor of Definition~\ref{def:step-factor} at $j = 1$. The boundary and periodic comparisons of \cite[Section~7 and Appendix~F]{WilsonFirstStep} are made per coarse segment and use the torus only through its smallest period; they hold for $C$ with the sums over the $4n$ coarse segments of a square replaced by the sums over the $P_0(C)/L$ coarse segments of $C$. The bounds of the leading term then hold with $e^{4C\tau_0^2P_0(C)/L}$ in place of $e^{C\tau^2(9n + 1)}$, and the constant $1 + 1/c$, which bounds $(n + 1)e^{-cn}$ for a square, is replaced by $1 + 2P(C)/(e\,c\,r(C))$, since $(P_0(C)/L + \abs{\mathcal K(C)})e^{-cr_0(C)/L} \leq 2(P(C)/r(C))\,(r_0(C)/L)e^{-cr_0(C)/L}$ by Lemma~\ref{lem:chord-arc}.

\emph{The retreated carrier and the labelled terms.} The analysis of \cite[Section~9]{WilsonFirstStep} on the labelled terms uses the loop through the envelopes of the groups, which are arcs of fewer than half of the traversals or the whole loop read from a corner, through the pins of the groups, and through the lifting of the groups to the universal cover. By Lemma~\ref{lem:square-properties}(ii) the number of cubes of side $LM$ which meet an envelope $\mathcal H(W)$ is at most $4 + (1 + 2/L)\abs{\mathcal H(W)}/(LM)$, in place of $8 + \abs{\mathcal H(W)}/(LM)$, so that the constants of the envelopes are unchanged. In the pins, the hull of a component $\gamma$ of a group has at most $(8\Lambda(C) + 1)LMn_\gamma$ traversals, in place of $17LMn_\gamma$, and the condition on the contraction rate which V4 imposes becomes $1000\,\eta\,(8\Lambda(C) + 2) < 1/128$, which holds for $\eta \leq \eta(\mathcal C) \coloneqq 1/(1\,200\,000\,\Lambda(\mathcal C))$. A small interior group has its contacts within a set of diameter less than $r_0(C)$, so that by Lemma~\ref{lem:square-properties}(iii) they lie on two consecutive segments, and its lifting is the one of the square. The torus conditions of \cite{WilsonFirstStep} involve the loop through its diameter and through $4ne^{-c\iota_T} \leq 1$, which reads $(P_0(C)/L)e^{-c\iota_T} \leq 1$; both hold for $K \geq K_{\mathrm{tor}}(\mathcal C)$, because $3\diam(C) < T$ and $\log(P_0(C)/L) \leq \log P_q(C) + K\log L$ while $\iota_T$ grows like $L^K$.

\emph{Several loops.} For two squares, \cite[Section~10]{WilsonFirstStep} groups the clusters of the joint expansion into groups of one loop and cross groups, proves Schur's lemma for each loop separately, and bounds the cross terms by the distance between the squares. Its proofs use the geometry of the pair only through the distance between the two loops and through the property that a set of diameter less than that distance meets one loop. For an ensemble, by Lemma~\ref{lem:stopping-index} the loops are at distance at least $a_0$ on the fine lattice, so that a set of diameter less than $a_0$ meets one loop, and the arguments of \cite[Section~10]{WilsonFirstStep} apply to the loops which a cluster meets. The cross exponent is the sum over the clusters which meet at least two loops; a cluster may meet three or more, and it is in the bound of the exponent, not in the exponent itself, that such a cluster is charged to one pair of the loops which it meets. In the integrands of that section every matrix factor belongs to one loop: at each value of the integration variables the source is the product, over the loops, of the normalised traces of ordered products of unitary matrices \cite[Lemma~10.4]{WilsonFirstStep}, so that the insertion of a cross group is an integral of elementary tensors against the measures of its roots and clusters. The majorants $M(W)$ by which \cite[Section~10]{WilsonFirstStep} bounds the cross groups are the integrals of the products of the norms of the factors, hence bounds in the norm \eqref{eq:projective-norm}, and Lemma~\ref{lem:word-norm} takes the place of the inequality for one loop. The cross groups of the first step have the reference scalar $0$, and their flat values need not vanish \cite[Section~10.3]{WilsonFirstStep}. Their total is at most $C_\times\rho\,\Delta_\times$, with $\Delta_\times$ the sum of $e^{-c_\times\delta_a/(LM)}$ over the fine traversals $a$ and $\delta_a$ the distance of $a$ from the other loops \cite[Lemma~10.7(a)]{WilsonFirstStep}; since a traversal of the scale $1$ consists of $L$ fine traversals at $L$ times its distance in fine units, and $\nu_\times \leq c_\times$, this is \eqref{eq:cross-bound} at $k = j = 1$.

\emph{The construction of this paper.} Two features of the construction differ from those assumed in \cite{WilsonFirstStep}. The thresholds are those of the tube at $\hat g_0$: by Lemma~\ref{lem:tube-thresholds} the radius of the fluctuation cutoff is $t_0 \coloneqq \vartheta_0/g_0 \in [\rho^{-1/2}r_0, r_0]$, and every statement of \cite{WilsonFirstStep} which uses the radius as an upper bound on a passed field holds a fortiori, while each one which uses it as a lower bound holds with the constant in its exponent divided by $\rho \leq 3/2$; the only condition which changes is the one behind $e^{-c_{\mathrm{ph}}r^2} \leq g_0$ in \cite[(9.29)(iii)]{WilsonFirstStep}, with $c_{\mathrm{ph}} = 9c_n\alpha_G/128$, which at the radius $t_0$ reads $(9c_n\alpha_G/(64\rho))A_1^2\log g_0^{-2} \geq 1$ and is a smallness condition on the coupling (Supplement, Proposition~S.7). The test functions of the plaquettes are soft: by Lemma~\ref{lem:soft-decomposition} a term of the decomposition is supported where the passed plaquettes are below $(1 + \kappa_s)$ times the threshold, and every bound of \cite{WilsonFirstStep} drawn from a passed test function holds at the threshold multiplied by $1 + \kappa_s$ (Supplement, Proposition~S.2).

\emph{Assembly.} With these changes, the main theorem of \cite{WilsonFirstStep} and its form on the labelled terms give, on every labelled term, the contribution of the ensemble as $\prod_i\kappa_1(C_i, \chi_i)e^{\mathcal E_1}[\prod_if^{(i)}_1 + \sum_{\mathfrak B \neq \emptyset}Q_1(\mathfrak B)w_1(\mathfrak B)]$, with the insertions bounded per fine traversal by the deviation of the background and the attached parts, the interactions resummed by the cluster expansion against the reference gas, and the scalar exponent the sum of the carrier part, which is a sum of local terms, and of the torus and corner part, which is a constant. The bounds (i)--(iv) of Definition~\ref{def:dressed-loop} at $k = 1$ follow by summing the bounds of \cite{WilsonFirstStep} over the $L$ fine traversals of each traversal of the scale $1$, which multiplies the constants by $L$ and converts the decay lengths $LM$ in fine units into $M$ in units of the scale $1$, and by the inequality $g_0 \leq (1 + \beta_0)g_1$ of Lemma~\ref{lem:tube-couplings}, which converts the parameters at $g_0$ into those at $g_1$. Clause (vi) holds at $k = 1$ with vanishing differences: on the trivial term every insertion and every cluster weight of the first step is a finite sum of absolutely convergent walk series and parameter integrals in local data, each term of which is the same expression on a lift \cite[Sections~E.9.2 and~E.10.4]{WilsonFirstStep}. The constants depend on the loops only through $\Lambda(C_i)$, $P(C_i)/r(C_i)$, and the characters, and on the number of loops not at all, because every count above is a count per traversal or per cube, and a cube of side less than $a_0$ meets one loop.
\end{proof}

\section{The Inductive Step}\label{sec:inductive-class}

We now show that a renormalisation step maps a dressed loop of the step $k$ to one of the step $k + 1$, with the same constants. The difficulty is one of counting. A loop has $P_k(C)$ traversals at the step $k$, and $P_k(C) = P_q(C)a_k$ grows like $L^{K-k}$, so that every quantity which a step contributes in proportion to the perimeter is, summed over the loop, unbounded as $K \to \infty$, even when it is of order $g_k^2$ per traversal. Such quantities can therefore be neither estimated nor neglected; each of them has to be either absorbed exactly into the step factor $\kappa_{k+1}$, or carried along without ever being summed in absolute value over the loop. The deviation parts of the insertions are of the second kind, and what keeps them under control is a contraction: a traversal of the step $k + 1$ consists of $L$ traversals of the step $k$, while the deviation of the background, measured in units of the scale $k + 1$, is $L^2$ times the same deviation measured in units of the scale $k$, so that with the normalisation $\rho_k\delta^*$ of \eqref{eq:insertion-bound} an old insertion contributes to the new bound with a factor of order $L \cdot L^{-2} = L^{-1}$.

\begin{theorem}[The Inductive Step]\label{thm:inductive-step}
Let $L$, $\Lambda_*$, $\bar g$, $\mathcal C$, and $K \geq K_{\mathrm{tor}}(\mathcal C)$ be as in Theorem~\ref{thm:first-step}, and let $1 \leq k < K'$. If the ensemble is dressed at the step $k$, then it is dressed at the step $k + 1$, with the same constants in Definition~\ref{def:dressed-loop}, and \eqref{eq:numerator-integral} holds at $k + 1$.
\end{theorem}

\begin{corollary}[Dressed Loops at Every Step]\label{cor:dressed-every-step}
Under the hypotheses of Theorem~\ref{thm:first-step}, the ensemble is dressed at every step $1 \leq k \leq K'$, with constants independent of $k$, $K$, the torus, and the number of loops.
\end{corollary}

\begin{proof}
Induction on $k$, with Theorem~\ref{thm:first-step} as its base and Theorem~\ref{thm:inductive-step} as its step.
\end{proof}

The proof of Theorem~\ref{thm:inductive-step} occupies the rest of this section. It follows Ba{\l}aban's step \cite[Sect.~3]{BalabanCRE} on the terms with small fields, with the dressed loop carried through each of its operations, his treatment of large fields \cite[Sect.~1]{BalabanLFR2} for the terms whose large-field regions grow, and the $\mathbf R$-operation of \cite{BalabanLFR1,BalabanLFR2}. We fix a labelled term of $\varrho_k$ and write $U_k$ for its background field, $U_{k+1}$ for the background field of the step $k + 1$, and $V^{(k)} \coloneqq M^k(U_{k+1})$ for the block field of the step $k$ determined by $U_{k+1}$.

\subsection{The Evaluation Point}\label{sec:evaluation-point}

The fluctuation field of the step $k + 1$ is the deviation of $V_k$ from $V^{(k)}$, and every factor of the dressed loop has to be expanded around the configuration at which this deviation vanishes. Two facts make this possible. First, $V^{(k)} = M^k(U_{k+1})$ \cite[(2.3), p.~265]{BalabanRG1}, and since the minimal orbit with given averages is unique \cite[p.~256]{BalabanCRE}, $U_k(V^{(k)}) = U_{k+1}$ in the gauge of the step. Secondly, $V^{(k)}$ lies in the support of the test functions of the step $k$ on $\Omega_{k+1}^{\sim}$: by \cite[(3.19), p.~268]{BalabanCRE} the block field $V^{(k)}$ is within $\vartheta_k$ of the local reference fields on the cubes of \cite[(3.16)]{BalabanCRE}, so that it lies in the support of the test function \cite[(3.3), p.~265]{BalabanCRE} with threshold $2\vartheta_k$, and Ba{\l}aban concludes that ``$\chi_k(\square) = 1$ for $\square \subset \Omega_{k+1}^{\sim}$'' \cite[p.~266]{BalabanCRE}. For the soft test functions the same holds with the margin of Lemma~\ref{lem:soft-decomposition}. On the cubes of the band $\Omega^{\sim}_{k+1} \setminus \Omega^{\sim-1}_{k+1}$, which lie in no cube of \cite[(3.16)]{BalabanCRE}, the localisation of the minimiser with constraints of two scales is needed; it is not printed, and it is proved in Supplement, Proposition~S.8, from the uniqueness of the minimal orbit \cite[Theorem~1]{BalabanVAR}.

The regularity of the new background is \cite[p.~268]{BalabanCRE}: $\abs{\partial U_{k+1} - 1} < 2B_3\varepsilon_{k+1}(L^{-1}\eta)^2$ on $\Omega_{k+1}$. By the lasso identity, which writes the holonomy around a plaquette of the scale $j$ as an ordered product of conjugates of the holonomies of the $L^{2j}$ fine plaquettes inside it, and by $\abs{\prod_iA_i - 1} \leq \sum_i\abs{A_i - 1}$ for unitary $A_i$, a plaquette of the scale $k$ in $\Omega_{k+1}$ has deviation at most $2B_3\varepsilon_{k+1}L^{-2}$ under $U_{k+1}$. Consequently, for the deviation \eqref{eq:deviation-quantity},
\begin{equation}\label{eq:deviation-scaling}
    \delta^{*,(k+1)}_{\mathfrak n'}(U) = L^2\,\delta^{*,(k)}_{\mathfrak n'}(U)
    \qquad\text{and}\qquad
    \delta^{*,(k+1)}_{\mathfrak n'}(U_{k+1}) \leq 2c_\sigma B_3\varepsilon_{k+1}\sum_{t \geq 1}(t + 1)^3e^{-\eta_wt}\;,
\end{equation}
where the superscript indicates the scale whose units are used, the first identity holds for the same plaquettes and radii, and $c_\sigma$ converts Ba{\l}aban's matrix norm into the norm of the insertions.

\subsection{The Identity of the Step}\label{sec:step-identity}

The first thing to establish is that the step can be made with the dressed loop inside the integrand. Ba{\l}aban passes from the density $\varrho_k$ to the expansion of $T\varrho_k$ by a sequence of operations on the integrand of a fixed term \cite[pp.~264--270]{BalabanCRE}, and for the density only the integral of the result matters. The dressed loop is a further factor, which depends on the integration variables, and each operation has to be an identity for the product. Most of them are identities for every integrand. Two kinds are not: the gauge fixing requires that the integrand be gauge invariant, and the changes of variables have to be changes of variables in the proper sense, bijections between the domains of the old and of the new variables, with which the further factor can be composed.

We fix a labelled term of $\varrho_k$ and sets $\Omega_{k+1}$, $\Lambda_{k+1}$, $S_{k+1}$ of the step. On the bonds of $\Omega_{k+1}$ the step writes the block field as $V_k = e^{iB'}V^{(k)}$ with $B' = g_kCA_k - h\tilde D(g_kCA_k)$, where $A_k$ is the fluctuation field \cite[(3.11), p.~266; (3.15), pp.~267--268]{BalabanCRE}, and on the bonds of $S_{k+1}$ it changes the variables by the map $\Psi_{k+1}$ of \cite[p.~269]{BalabanCRE} (Supplement, Lemma~17.23). Let
\begin{equation}\label{eq:carried-loop}
    F \coloneqq \mathcal W_k\bigl(e^{iB'}V^{(k)}\bigr) \circ \Psi_{k+1}
\end{equation}
be the dressed loop of the term at this block field, as a function of the new fluctuation variables, of $V_{k+1}$, and of the outer variables. The field $A_k$ is integrated on $\Lambda_{k+1}$, where we write $A$ for it, and it is an outer variable of the step on $\Omega_{k+1} \cap \Lambda_{k+1}^c$ \cite[p.~270]{BalabanCRE}. Let $\mathcal Z_{k+1}$ be the last integral of \cite[(3.25), p.~270]{BalabanCRE}, the integral over $A$ with its factor $z^{(k)}$, and let $\langle\,\cdot\,\rangle$ be the expectation in the probability measure \cite[(3.27), p.~271]{BalabanCRE} at $s = t = 1$, which is the integrand of that integral divided by $\mathcal Z_{k+1}$.

\begin{proposition}[The Identity of the Step]\label{prop:step-identity}
Let $1 \leq k < K'$, and let the ensemble be dressed at the step $k$. Then
\begin{equation}\label{eq:carried-step}
    (TN_k)(V_{k+1}) = \sum\chi_{k+1}(\Omega_{k+1})\,\mathbf T_{k+1}\Bigl[\exp\Bigl(A_k\Bigl(\frac{1}{g_k^2}, U_{k+1}\Bigr) + E_0^{(k)}\Bigr)\,\mathcal Z_{k+1}\,\langle F\rangle\Bigr]\;,
\end{equation}
the sum over the sequences of domains of \cite[(3.25)]{BalabanCRE}; this is Ba{\l}aban's expansion of $T\varrho_k$ with the factor $\langle F\rangle$ in each term. The factor $\langle F\rangle$ is a function of $V_{k+1}$ and of the outer variables of $\mathbf T_{k+1}$, which are those of $\mathbf T_k$, the field $A_k$ on $\Omega_{k+1} \cap \Lambda_{k+1}^c$, and the field $V_k$ on $\Omega_{k+1}^c$.
\end{proposition}

\begin{proof}
By Definition~\ref{def:dressed-loop}, $N_k = \sum\chi_k(\Omega_k)\mathbf T_k[\mathcal W_k\exp A_k]$, and $TN_k$ is the expression \cite[(3.1), p.~264]{BalabanCRE} with the factor $\mathcal W_k$ in each term. We follow the operations by which Ba{\l}aban transforms a term of (3.1) into the terms of (3.25), and we show that each of them is an identity when the integrand is multiplied by a bounded Borel function of the integration variables, composed with the changes of variables where there are any. The function $\mathcal W_k$ is of this kind. Indeed, the transports of a word are unitary at real fields, so that by Lemma~\ref{lem:word-norm} and Definition~\ref{def:dressed-loop}(i)
\[
    \Bigl\lvert\sum_{\mathfrak B \neq \emptyset}Q_k(\mathfrak B)\,w_k(\mathfrak B)\Bigr\rvert \leq \prod_W\bigl(1 + e^{\theta(W)}\norm{Y_W}\bigr) \leq \exp\Bigl(\sum_We^{\theta(W)}\norm{Y_W}\Bigr)\;,
\]
and the last sum is at most the sum, over the finitely many traversals of the torus, of the bounds of Definition~\ref{def:dressed-loop}(ii), which are bounded on the support of the test functions of the term.

\emph{Decompositions of unity.} The decompositions \cite[(3.2), (3.3), p.~265; (3.16), p.~268]{BalabanCRE} multiply the integrand by a sum which equals one, for the soft test functions by Lemma~\ref{lem:soft-decomposition}(i). They hold for every integrand.

\emph{Test functions which equal one.} On $\Omega_{k+1}^{\sim}$ the test functions $\chi_k$ and $\chi_{\mathrm{Ax}}$ of the step $k$ are removed, because they equal one on the support of the new ones \cite[(3.8), (3.9), p.~266]{BalabanCRE}, and on the cubes of $R_{k+1}^c$ the test function of \cite[(3.3)]{BalabanCRE} is removed for the same reason \cite[(3.17)--(3.19), pp.~268--269]{BalabanCRE}; for the soft test functions the removals hold on the plateaus (Supplement, Proposition~S.2(a)). These are identities between factors of the integrand at each point, and they hold for every integrand.

\emph{Regrouping.} At fixed $\Omega_{k+1}$ the terms are summed over the sets $P_{k+1}$ and $Q_{k+1}$ which determine it, and ``the resummation applies only to the characteristic functions and the gauge fixing terms localized in $\Omega_k \backslash \Omega_{k+1}$'' \cite[p.~267]{BalabanCRE}. The sets $R_{k+1}$ are then replaced by the pairs $(\Lambda_{k+1}, S_{k+1})$ of \cite[(3.20), p.~269]{BalabanCRE}. The dressed loop is a function of $V_k$ and of the outer variables which does not depend on these sets, so that it is a common factor of the terms which are regrouped.

\emph{The gauge fixing.} In the blocks of a set $P^1_{k+1}$ of points of the scale $k + 1$ Ba{\l}aban introduces ``the axial gauge fixing as in (1.5), with $V_k$, $\varepsilon_k$, $g_k$ instead of $U$, $\varepsilon_0$, $g_0$'', and obtains ``an expansion of the form (1.6)'' \cite[p.~265]{BalabanCRE}. The identity \cite[(1.5), p.~247]{BalabanCRE} is an integral over gauge transformations $u(x)$ at the points $x \neq y$ of the blocks $B(y)$, $y \in P^1_{k+1}$, and the Faddeev--Popov procedure which follows it replaces $V_k$ by a gauge transform under the integral. It uses that the other factors of the integrand are invariant under the gauge transformations $u$ of $V_k$ which are supported at these points, in Ba{\l}aban's words ``the invariance of all the expressions determining the density with respect to the gauge transformations of the new field'' \cite[p.~250]{BalabanCRE}. We show that $\mathcal W_k(V_k^u) = \mathcal W_k(V_k)$ for these $u$. Let $\tilde u(x') \coloneqq u(x)$ for the fine points $x'$ of the block $B^k(x)$ of a point $x$ of the scale $k$, which is the extension that \cite[p.~250]{BalabanCRE} uses at the first step. The averages are covariant, $M(\{ u\mathbf U_jv \}) = uM(\{ \mathbf U_j \})v$ \cite[(0.6), p.~253]{BalabanRG1}, and the action is invariant, so that $U \mapsto U^{\tilde u}$ maps the minimal orbit of $V_k$ onto that of $V_k^u$. The representative $U_k(V)$ is singled out in the minimal orbit by axial conditions inside the blocks, at every level of their hierarchy: conditions on the averaged contour variables $\mathbf U(y, x)$ of \cite[(0.11), p.~253]{BalabanRG1}, or on the transports along fixed contours \cite[(1.15), p.~78]{BalabanREG}. In a block $B^k(x)$ the transformation $\tilde u$ is constant, so that it conjugates these variables by $u(x)$ and preserves the conditions, and outside these blocks it equals one. Since the conditions ``determine uniquely an element in each orbit'' \cite[p.~78]{BalabanREG}, it follows that $U_k(V_k^u) = U_k(V_k)^{\tilde u}$. By Definition~\ref{def:dressed-loop}(iii) the factors of the words are covariant, so that a closed word is invariant under a gauge transformation of the background which leaves the outer variables fixed; and $\tilde u$ leaves them fixed, because the blocks of $P^1_{k+1}$ lie in $(Z_k'^{\sim5})^c$ \cite[p.~265]{BalabanCRE}, whereas the outer variables of the term are fields on bonds of the large-field region $Z_k \subset Z_k'$. This proves the invariance.

\emph{The fluctuation field.} On $B(\Gamma_{k+1})$ the block field is written $V_k = V_k'V^{(k)}$ and $A_k' = (1/i)\log V_k'$; the $\delta$-functions of the averaging are linearised and removed ``using the operator $C$'', and the field is scaled, $A_k' = g_kA_k$ \cite[(3.11), p.~266; p.~267]{BalabanCRE}. Here ``all the applied formulas are the same as in Sect.~2 [I]'', where the $\delta$-functions serve to ``eliminate the variables $B'(b_0(c))$'' \cite[p.~268]{BalabanRG1}. These operations express the integration variables on the support of the $\delta$-functions through the variables $A_k$, with a Jacobian which Ba{\l}aban includes in the action, and every other factor of the integrand is composed with this parametrisation. For the dressed loop the result is the function $\mathcal W_k(e^{iB'}V^{(k)})$ of \eqref{eq:carried-loop}.

\emph{The change of variables on $S_{k+1}$.} By Supplement, Lemma~17.23, whose two hypotheses hold when $\gamma$ is small, the map $\Psi_{k+1}$ is a diffeomorphism between open sets which contain the supports of the test functions in the new and in the old variables. Therefore the integral of Ba{\l}aban's integrand times a bounded Borel function of the old variables is the integral of his integrand in the new variables times the composition of that function with $\Psi_{k+1}$, which for the dressed loop is $F$.

\emph{The conditional integral.} Finally the integral over $A_k$ is split into the integral over its restriction $A$ to $\Lambda_{k+1}$ and the integral over its restriction to $\Omega_{k+1} \cap \Lambda_{k+1}^c$, which becomes part of the operation $\mathbf T^{(k+1)}$, the quadratic form being completed as in \cite[(3.23)--(3.25), p.~270]{BalabanCRE}. This is Fubini's theorem on the finite torus for an integrable function, and it holds with the bounded factor $F$.

After these operations the term of (3.1) with the factor $\mathcal W_k$ is the sum of the terms of (3.25) with the factor $F$ under the last integral. We have $\mathcal Z_{k+1} \neq 0$, since Ba{\l}aban takes the logarithm of this integral in \cite[(3.26), p.~271]{BalabanCRE}, and by the definition of the expectation the last integral with the factor $F$ equals $\mathcal Z_{k+1}\langle F\rangle$. This is \eqref{eq:carried-step}. The variables on which $\langle F\rangle$ depends are those of $F$ which are not integrated.
\end{proof}

Since the averaging preserves integrals (Section~\ref{sec:two-integrals}), the proposition gives $\int(TN_k)\,\dd V_{k+1} = \int N_k\,\dd V_k$, and the task of the step is to write $\langle F\rangle$ in the form of a dressed loop of the step $k + 1$. At the steps at which the $\mathbf R$-operation acts it is applied afterwards (Section~\ref{sec:large-field-step}).

\subsection{The Expansion of the Words in the Localisation Parameters}\label{sec:localisation-expansion}

To write $\langle F\rangle$ as a dressed loop we first have to see on which part of the fluctuation field each word of $F$ depends, and the difficulty is that this dependence is not local. By \cite[p.~265 and (2.6), p.~266]{BalabanRG1}, whose formulae \cite[p.~267]{BalabanCRE} transfers to the operators of a sequence of domains, the background of the step $k$ at the block field $e^{iB'}V^{(k)}$ is a gauge transform of $e^{i\eta\mathbf H_k(B')}U_{k+1}$, where $\mathbf H_k$ is the function \cite[(1.17), p.~6]{BalabanRG2} of the fluctuation variable $B'$ of \cite[(1.19), p.~6]{BalabanRG2}, built from the propagators $H$, $\tilde G$, and $H_0$ of the step. These propagators have kernels which decay exponentially and do not vanish at any distance, so that the transport of the new configuration along a single bond of a loop depends on the fluctuation field on the whole lattice. Ba{\l}aban meets the same difficulty for the terms of the effective action, and he resolves it by parameters which switch off the propagators across the walls of a partition \cite[pp.~3--7]{BalabanRG2}. We apply his device to the words of the dressed loop.

\emph{The parameters.} Let $\sigma_k$ be the partition of the torus into the cubes $\Delta$ of side $M$ of the scale $k$, the family which \cite[p.~3]{BalabanRG2} takes in its example, here with no cube exempt. To each cube we assign a parameter $s(\Delta)$, and we define the propagators $H(s)$, $\tilde G(s)$, $H_0(s)$ by the rule of \cite[p.~3]{BalabanRG2}: the term of a walk $\omega$ in the expansion \cite[(1.6)]{BalabanRG2} of a propagator is multiplied by $s(\Delta_1) \cdots s(\Delta_m)$, where $\Delta_1, \dots, \Delta_m$ are the cubes of $\sigma_k$ which meet the localisation domain of $\omega$. Let $\mathbf H(s, B')$ be the function \cite[(1.17)]{BalabanRG2} with these propagators; by \cite[(1.5), p.~3]{BalabanRG2} the solutions which enter it are functionals of the propagators with the properties of the original ones. For a set $\sigma \subset \sigma_k$ let $Y(\sigma)$ be the union of its cubes, and call a union of cubes \emph{connected} if any two of its cubes are joined by a chain of its cubes in which consecutive ones share a wall of dimension $d - 1$ \cite[p.~4]{BalabanRG2}; the \emph{components} of $Y(\sigma)$ are its maximal connected parts. We use three properties.
\begin{enumerate}[label=(\alph*)]
\item $\mathbf H(1, B') = \mathbf H_k(B')$, since the propagators at $s = 1$ are the original ones \cite[p.~3]{BalabanRG2}, and $\mathbf H(0, B') = 0$, since the localisation domain of every walk meets a cube, so that all three propagators vanish at $s = 0$.
\item If $s(\Delta) = 0$ for every $\Delta \notin \sigma$, then $\mathbf H(s, B')$ vanishes on a neighbourhood of $Y(\sigma)^c$, and on a component $Y$ of $Y(\sigma)$ it depends only on the restrictions of $s$, of $B'$, and of the configuration to $Y$; we write $\mathbf H_Y(s, B')$ for it. This is the argument of \cite[p.~4]{BalabanRG2}: the kernels of the propagators vanish unless both arguments lie in the interior of one component, the equations \cite[(1.3), (1.4)]{BalabanRG2} separate into equations in the components, and $B'$ is a local function of the fluctuation field.
\item For complex parameters with $\abs{s(\Delta)} \leq e^{\kappa_1}$, where $\kappa_1$ is the constant of \cite[p.~5]{BalabanRG2} with $\delta_1M \geq \kappa_1$, and for $\abs{B'} < \varepsilon_3$, the function $\mathbf H_Y(s, B')$ is analytic in $s$ and $B'$ and
\begin{equation}\label{eq:localised-h-bound}
    \abs{\mathbf H_Y(s, B')} \leq 4B_0e^{16\kappa_1}\abs{B'}_Y\;,
\end{equation}
with $\abs{B'}_Y$ the supremum over $Y$, by \cite[(1.18), p.~6]{BalabanRG2} applied to $B'$ restricted to $Y$; the constants $B_0$, $\varepsilon_2$, $\varepsilon_3$ are those of \cite[pp.~3, 5, 6]{BalabanRG2}, with $4C_2B_0e^{16\kappa_1}\varepsilon_2 \leq 1$, $4C_4B_0^2e^{32\kappa_1}\varepsilon_3 \leq 1$, and $2B_0e^{16\kappa_1}\varepsilon_3 \leq \varepsilon_2$, and ``all the bounds above are uniform on the domain'' of the configurations \cite[p.~6]{BalabanRG2}. None of the estimates \cite[(1.11)--(1.18)]{BalabanRG2} uses the cube which Ba{\l}aban exempts, so that they hold with every cube carrying a parameter.
\end{enumerate}

\emph{The term and the domain of a member.} In this subsection the labelled term of $\varrho_k$ is the term without large-field regions, for which $\Omega_j = \Lambda_j$ is the whole torus for every $j \leq k$ and there are no outer variables; the sets $\Omega_{k+1}$, $\Lambda_{k+1}$, $S_{k+1}$ of the step are arbitrary. On this term the operators of the step $k$ are those of one domain, to which \cite{BalabanRG2} applies as it stands. The terms of $\varrho_k$ with large-field regions, on which the propagators are those of a sequence of domains and the factors of the words depend on outer variables, are the subject of Section~\ref{sec:large-field-step}. For a group $W$ of a family of the step $k$, made at a step $j \leq k$, let $X_W$ be its domain and $X^+_W$ its dependence set (Section~\ref{sec:dressed-loops}); the domain contains the cubes of $\pi_j$ that meet a fine bond of its intervals, and the insertion $Y_W$ depends on the background field only through its restriction to $X^+_W$ (Definition~\ref{def:dressed-loop}(iii)).

\emph{The domain of the factors.} The words have to be evaluated at the configurations $e^{i\eta\mathbf H(s, B')}U_{k+1}$, which for $s \neq 1$ are not backgrounds of block fields, and we have to make sure that their factors are defined there. On the support of the test functions of the step, $\abs{g_kA_k} \leq r_y\vartheta_k$ on every bond; this is the bound \cite[(3.14), p.~267]{BalabanCRE} in the form which the soft test functions give it, with the constant $r_y$ of Supplement, Lemma~17.23, in whose proof this form is derived. For $\gamma$ small the right side is below the constant $\varepsilon_1$ of \cite[p.~7]{BalabanRG2}, and therefore, by \cite[(1.21), p.~7]{BalabanRG2},
\begin{equation}\label{eq:localised-h-size}
    \abs{\mathbf H(s, B')} \leq h_k \coloneqq 4B_0C_1e^{16\kappa_1}r_y\vartheta_k
\end{equation}
for $\abs{s(\Delta)} \leq e^{\kappa_1}$, ``and the same for all admissible norms'', with the absolute constant $C_1$ of \cite[(1.20), p.~6]{BalabanRG2}. Among these norms are those of the condition \cite[(3.14), p.~272]{BalabanRG1}, under which Ba{\l}aban evaluates the terms of the action of the earlier steps at a configuration $e^{i\xi\mathbf A}\mathbf U$, with the argument \cite[(3.13)]{BalabanRG1} in the place of $\mathbf J$: he requires of $\mathbf H_k(s, B')$ the bounds of that condition \cite[p.~7]{BalabanRG2}, and he obtains them from \cite[(1.21)]{BalabanRG2}. The factors of the dressed loop are analytic on the same spaces as those terms (Definition~\ref{def:dressed-loop}(iii)). In the units of the step $k$ the radii of the space of the step $k + 1$ are smaller than those of the step $k$ by at least $\frac12\min(\beta, 1)\min(\alpha_{0,k}, \alpha_{1,k})$: on $\Omega_k \setminus \Omega_{k+1}$ by the factor $1 - \frac12\beta$ of \cite[(2.34)--(2.39), p.~261]{BalabanCRE}, and on $\Omega_{k+1}$ because $\alpha_{i,k+1}/L \leq \frac12\alpha_{i,k}$, by \eqref{eq:tube-two-eight} and $(1 + \beta_0)^2 \leq L/2$. Let $c_{\mathbf J} \geq 1$ be an absolute constant such that, for a point $(\mathbf U, \mathbf J)$ of the space of the step $k + 1$ and a field $\mathbf A$ whose norms in \cite[(3.14)]{BalabanRG1} are at most $a \leq 1$, the configuration $e^{i\xi\mathbf A}\mathbf U$ with the argument \cite[(3.13)]{BalabanRG1} lies in the space of the step $k$ whenever $c_{\mathbf J}a \leq \frac12\min(\beta, 1)\min(\alpha_{0,k}, \alpha_{1,k})$; it exists because that argument differs from the current of $\mathbf U$ by $D^{\xi*}_{\mathbf U}D^\xi_{\mathbf U}\mathbf A + \mathbf F(\mathbf U, \mathbf A)$ with a local operator $\mathbf F$ of the second order in $\mathbf A$ and $\nabla^\xi_{\mathbf U}\mathbf A$ \cite[(3.11), p.~272]{BalabanRG1}. We assume that
\begin{equation}\label{eq:expansion-condition}
    c_{\mathbf J}h_k \leq \tfrac14\min(\beta, 1)\min(\alpha_{0,k}, \alpha_{1,k})\;,
\end{equation}
which holds for $\gamma$ small, since $\vartheta_k/\alpha_{i,k} = (A_1/C^\alpha_i)(\log\hat g_k^{-2})^{p_0 - q_i}$ and $q_0 = q_1 > p_0$ (Appendix~\ref{app:constants}). Then, for every $s$ with $\abs{s(\Delta)} \leq e^{\kappa_1}$, the factors of the words of the step $k$ are defined and analytic at $e^{i\eta\mathbf H(s, B')}U_{k+1}$, with half of the margin between the spaces of the two steps left.

\emph{The expansion of a word.} By the covariance of its factors (Definition~\ref{def:dressed-loop}(iii)) the word of a family $\mathfrak B$ is invariant under gauge transformations of the background, so that at the block field $e^{iB'}V^{(k)}$ it is the same word evaluated at the configuration $e^{i\eta\mathbf H_k(B')}U_{k+1}$. For a family $\mathfrak B$, the empty family included, let $\Phi_{\mathfrak B}(s)$ be its word evaluated at $e^{i\eta\mathbf H(s, B')}U_{k+1}$. For $k \geq 2$ the factors are functions of a configuration and of the variable $\mathbf J$ of \cite[(2.36), p.~261]{BalabanCRE}, and they are evaluated at the argument which \cite[(3.10)--(3.13), p.~272]{BalabanRG1} substitutes for $\mathbf J$ when the configuration is $e^{i\xi\mathbf A}U_{k+1}$: by \cite[(3.11), (3.12)]{BalabanRG1} the current of that configuration is $(L^{j-1}\eta)^3J_{k+1} + D^{\xi*}_{U_{k+1}}D^\xi_{U_{k+1}}\mathbf A + \mathbf F(U_{k+1}, \mathbf A)$, in the normalisation of the step $j$ at which the factor was made, and we take this expression at $\mathbf A = \mathbf H(s, B')$. By (a), $\Phi_{\mathfrak B}(1)$ is the word at the configuration of integration and $\Phi_{\mathfrak B}(0)$ is the word at $U_{k+1}$. The torus has finitely many cubes, and the fundamental theorem of calculus in each parameter gives, as in \cite[(1.9), p.~4]{BalabanRG2},
\begin{equation}\label{eq:parameter-expansion}
    \Phi_{\mathfrak B}(1) = \sum_{\sigma \subset \sigma_k}\prod_{\Delta \in \sigma}\int_0^1\partial_{s(\Delta)}\Phi_{\mathfrak B}(s)\big|_{s(\sigma^c) = 0}\,\dd s(\Delta)\;.
\end{equation}
The term $\sigma = \emptyset$ is $\Phi_{\mathfrak B}(0)$. In the term of $\sigma$, by (b), the configuration on a fine bond is $e^{i\eta\mathbf H_Y}U_{k+1}$ if the bond lies in a component $Y$ of $Y(\sigma)$, and $U_{k+1}$ otherwise. A factor of the word therefore has its value at $U_{k+1}$ unless the interior of a component meets the set on which it depends: for a transport, the bonds of its arc, and for the insertion of a member $W$ of $\mathfrak B$, the dependence set $X^+_W$. We say that a component \emph{touches} a member $W$ if its interior meets $X^+_W$, and we call a component \emph{active} for $\mathfrak B$ if it contains a fine bond of a loop or touches a member of $\mathfrak B$. If a component of $Y(\sigma)$ is not active, the word does not depend on its parameters and the term vanishes, so that only the sets $\sigma$ all of whose components are active remain.

\emph{Arcs and groups.} Let $Y$ be an active component and $C_i$ a loop. Let $S'_i(Y)$ be the set of the fine bonds of $C_i$ in $Y$ together with the intervals on $C_i$ of the members of $\mathfrak B$ which $Y$ touches. An interval $J$ of a member $W$ of $\mathfrak B$ on $C_i$ which contains a bond $b$ of $S'_i(Y)$ lies in $S'_i(Y)$. Indeed, if $b$ lies in an interval of a member touched by $Y$, that interval is $J$, since the intervals of a family are pairwise disjoint. Otherwise $b$ is a bond of $C_i$ in $Y$; the cube of the partition $\pi_j$ of the step at which $W$ was made which contains an end site of $b$ belongs to $X_W$, and this site is an interior point of $Y$, the sites being interior points of the cubes of the partitions (Section~\ref{sec:reference-gas}), so that $Y$ touches $W$ and again $J \subset S'_i(Y)$. We put no arcs between the parts of $S'_i(Y)$: to join them is left to the regrouping of Section~\ref{sec:groups-of-step}, where the arcs which are added are bounded by the sizes of the objects which they join. Two active components are \emph{linked} if they touch a common member of $\mathfrak B$, and the \emph{groups} $G$ of the term are the classes of the equivalence relation which this generates. The \emph{arc set} of a group is $I(G) \coloneqq \bigcup_i\bigcup_{Y \in G}S'_i(Y)$, a union of arcs of the loops. The arc sets of distinct groups are disjoint. Indeed, a fine bond of a loop lies in at most one component, two components having no cube in common; a bond of an interval of a member lies in the sets $S'_i(Y)$ of the components which touch that member, and these are linked; and if a bond of a component $Y$ lies in an interval of a member, then $Y$ touches that member, by what was shown. A group whose arc set meets two loops contains a component with bonds or touched members on both, because the intervals of a member which two linked components touch belong to the sets $S'_i$ of both.

A member touched by a component of $G$ has all its intervals in $I(G)$, and it is touched by no component of another group, by the definition of the groups. A member of $\mathfrak B$ which no component touches is \emph{untouched}; its insertion has its value at $U_{k+1}$, by (b), and its intervals are disjoint from the arc sets of all the groups, by what was shown. Let $\mathfrak B(G)$ be the set of the members touched by the components of $G$, and call the maximal arcs of $I(G)$ the \emph{pieces} of $G$. We define the tensor $Y_G$ by applying the operations $\prod_{Y \in G}\prod_{\Delta \subset Y}\int_0^1\dd s(\Delta)\,\partial_{s(\Delta)}$ to the ordered products, along the pieces of $G$, of the insertions of the members of $\mathfrak B(G)$ and of the transports of $e^{i\eta\mathbf H(s, B')}U_{k+1}$ along the bonds of the pieces between them, at $s(\Delta) = 0$ for every cube outside the components of $G$; for insertions which are not elementary tensors the ordered products are defined by multilinearity, as in Section~\ref{sec:dressed-loops}. So $Y_G$ has one factor for each of its pieces, a linear map between the fibres at the two ends of the piece. By what was said, the parameters of the components of another group do not enter $Y_G$.

\begin{lemma}[The Expansion of a Word]\label{lem:word-expansion}
Let $\mathfrak B$ be a family of the step $k$ on the term without large-field regions, and suppose \eqref{eq:expansion-condition}. For real fluctuation fields in the support of the test functions of the step, the word of $\mathfrak B$ at the configuration of integration equals
\begin{equation}\label{eq:word-expansion}
    \sum_{\mathfrak Y}w(\mathfrak B, \mathfrak Y)\;,
\end{equation}
where the sum runs over the finite sets $\mathfrak Y$ of connected unions of cubes of $\sigma_k$ which are active for $\mathfrak B$ and no two of which have a cube or a wall in common, the empty set included, and $w(\mathfrak B, \mathfrak Y)$ is the word formed, in the cyclic order of each loop, from the tensors $Y_G$ of the groups $G$ of $\mathfrak Y$, the insertions at $U_{k+1}$ of the untouched members of $\mathfrak B$, and the transports of $U_{k+1}$ between them. The term $\mathfrak Y = \emptyset$ is the word of $\mathfrak B$ at $U_{k+1}$; for $\mathfrak B = \emptyset$ it is the product of the leading words of the step $k + 1$. Each $Y_G$ depends on the fluctuation field and on the parameters only on the components of $G$, and on $U_{k+1}$ only along the pieces of $G$, on its components, and on the dependence sets of the members of $\mathfrak B(G)$.
\end{lemma}

\begin{proof}
By \eqref{eq:parameter-expansion} and the remarks which follow it, the word is the sum over the sets $\sigma$ all of whose components are active. The map which sends $\sigma$ to the set $\mathfrak Y$ of the components of $Y(\sigma)$ is a bijection onto the sets of the statement: the components of $Y(\sigma)$ are connected, active, and no two of them share a cube or a wall, and conversely the components of the union of the members of such a set $\mathfrak Y$ are these members. In the term of $\sigma$ the integrand is the word in which the factors on the pieces of the groups are evaluated at the configurations $e^{i\eta\mathbf H_Y}U_{k+1}$, and all other factors, the insertions of the untouched members among them, at $U_{k+1}$. The operations $\int_0^1\dd s(\Delta)\,\partial_{s(\Delta)}$ of the cubes of a component act only on the factors of its group, by (b) and the two remarks before the lemma, and the word is multilinear in its factors, so that the term of $\sigma$ is $w(\mathfrak B, \mathfrak Y)$. The locality of $Y_G$ is (b) with the locality of the factors of the word.
\end{proof}

The expansion is useful because its terms are small: each component of a group contributes a factor which is proportional to the fluctuation field on it and decays exponentially in the number of its cubes, at the rate $\kappa_1$ of the parameters. We prove this here for the groups which touch no member, whose tensors are built from transports alone; the groups with members are estimated in Lemma~\ref{lem:tensor-members} and Proposition~\ref{prop:marked-size} from the bounds of the class. For complex $X$ in the complexified Lie algebra, $\norm{\pi(e^{iX})} \leq e^{c'_\pi\abs X}$ with $c'_\pi \coloneqq 2\sqrt Nc_\pi$: writing $X = X_1 + iX_2$ with $X_1, X_2$ in the Lie algebra, we have $\abs{X_1}, \abs{X_2} \leq \abs X$ for the norm of the test functions, which is invariant under the adjoint, and $\norm{\dd\pi(X)} \leq c_\pi(\norm{X_1} + \norm{X_2})$ by \eqref{eq:representation-constant}, with $\norm{X_j} \leq \sqrt N\abs{X_j}$. Let $\bar\Lambda \coloneqq \max_i\Lambda(C_i)$, write $\#(Y)$ for the number of the cubes of a component and $\#(G) \coloneqq \sum_{Y \in G}\#(Y)$, and fix a constant $\varepsilon_H < \varepsilon_3$ with
\begin{equation}\label{eq:epsilon-h}
    16\bar\Lambda c'_\pi MB_0e^{16\kappa_1}\varepsilon_H \leq 1\;,
\end{equation}
where $c'_\pi$ is taken for the representation with the largest $c_\pi$.

\begin{lemma}[The Size of a Tensor without Members]\label{lem:tensor-size}
In the setting of Lemma~\ref{lem:word-expansion} let $\kappa_1 \geq 1$, and let $G$ be a group of a set $\mathfrak Y$ with $\mathfrak B(G) = \emptyset$. Let the background be $\mathbf U = e^{i\eta a'}U_{k+1}$ with a field $a'$ with values in the complexified Lie algebra for which $\mathbf U$ lies in the space of the step $k + 1$, and let $\norm{a'}_\infty$ be the supremum of $\abs{a'(b)}$ over the fine bonds of $I(G)$. If $\abs{B'}_Y \leq \varepsilon_H$ for every component $Y$ of $G$, then the tensor $Y_G$ formed with $\mathbf U$ in the place of $U_{k+1}$ satisfies
\begin{equation}\label{eq:tensor-size}
    \norm{Y_G} \leq \exp\bigl(c'_\pi\norm{a'}_\infty\abs{I(G)}\bigr)\prod_{Y \in G}\frac{\abs{B'}_Y}{\varepsilon_H}\,e^{-(\kappa_1 - 2)\#(Y)}\;,
\end{equation}
where $\abs{I(G)}$ is the length of the arc set in the units of the step $k$ and the norm is the projective norm \eqref{eq:projective-norm} over the pieces of $G$. Moreover $\abs{I(G)} \leq 4\bar\Lambda M\#(G)$.
\end{lemma}

\begin{proof}
For a family $t = (t_Y)_{Y \in G}$ of complex numbers let $\Phi(s, t)$ be the tensor product, over the pieces of $G$, of the ordered products of the transports $\pi_i\bigl(e^{i\eta\mathbf H_Y(s, t_YB')(b)}e^{i\eta a'(b)}U_{k+1}(b)\bigr)$ along the fine bonds $b$ of the piece, where $Y$ is the component which contains $b$ and the first factor is absent if there is none. By property (c) of the parameters, with the uniformity of Ba{\l}aban's bounds on the space of the configurations, $\Phi$ is analytic in $(s, t)$ for $\abs{s(\Delta)} \leq e^{\kappa_1}$ and $\abs{t_Y} \leq T_Y \coloneqq \varepsilon_H/\abs{B'}_Y$, since $\abs{t_YB'}_Y \leq \varepsilon_H < \varepsilon_3$ there; and $T_Y \geq 1$. By its definition,
\[
    Y_G = \prod_{Y \in G}\prod_{\Delta \subset Y}\int_0^1\partial_{s(\Delta)}\Phi(s, t)\,\dd s(\Delta)\Big|_{t = 1}\;,
\]
at $s(\Delta) = 0$ for the cubes outside the components of $G$. We bound the integrand $g(s, t) \coloneqq \prod_{Y}\prod_{\Delta \subset Y}\partial_{s(\Delta)}\Phi(s, t)$ at $t = 1$ and real $s$ in three steps.

First, $g$ vanishes whenever one $t_Y$ vanishes: then $\mathbf H_Y(s, t_YB') = 0$ for every $s$, by \eqref{eq:localised-h-bound}, so that $\Phi$ does not depend on the parameters of the cubes of $Y$ and its derivatives in them vanish. Schwarz's lemma, applied in each variable $t_Y$ on the disc of radius $T_Y$, therefore gives $\norm{g(s, 1)} \leq \prod_YT_Y^{-1}\sup_t\norm{g(s, t)}$, the supremum over the polydisc. Secondly, each derivative in a parameter $s(\Delta)$ at a point of $[0, 1]$ is bounded by Cauchy's estimate on the circle of radius $e^{\kappa_1} - 1$ about that point, which lies in the disc $\abs{s(\Delta)} \leq e^{\kappa_1}$; since $e^{\kappa_1} - 1 \geq e^{\kappa_1 - 1}$ for $\kappa_1 \geq 1$, this gives $\norm{g(s, t)} \leq e^{-(\kappa_1 - 1)\#(G)}\sup\norm\Phi$, the supremum over the polydisc in $s$ and $t$. Thirdly, $\Phi$ is an elementary tensor, and its norm is the product of the norms of its factors, each of which is a product of transports along single bonds. The transport along a bond $b$ has norm at most $\exp(c'_\pi\eta(\abs{\mathbf H_Y(s, t_YB')(b)} + \abs{a'(b)}))$, because $U_{k+1}(b)$ is unitary, and $\abs{\mathbf H_Y(s, t_YB')} \leq 4B_0e^{16\kappa_1}\varepsilon_H$ on the polydisc by \eqref{eq:localised-h-bound}. The fine bonds of the loops in a cube of $\sigma_k$ have total length at most $4\bar\Lambda M$ in the units of the step $k$, by Lemma~\ref{lem:square-properties}(i) at the scale $k$ and because a cube of side $M < a_k$ meets one loop (Lemma~\ref{lem:stopping-index}). Hence
\[
    \sup\norm\Phi \leq \exp\bigl(c'_\pi \cdot 4\bar\Lambda M\#(G) \cdot 4B_0e^{16\kappa_1}\varepsilon_H\bigr)\exp\bigl(c'_\pi\norm{a'}_\infty\abs{I(G)}\bigr) \leq e^{\#(G)}\exp\bigl(c'_\pi\norm{a'}_\infty\abs{I(G)}\bigr)\;,
\]
by \eqref{eq:epsilon-h}. Multiplying the three bounds, and integrating over the parameters in $[0, 1]$, we obtain \eqref{eq:tensor-size}; the projective norm of an integral of elementary tensors is at most the integral of the products of the norms of their factors.

For the length, a group which touches no member consists of one component, by the definition of the groups, and its arc set is the set of the fine bonds of the loops in that component, of total length at most $4\bar\Lambda M\#(G)$ by the count of the bonds in a cube which was used above.
\end{proof}

For a group of one component the bound \eqref{eq:tensor-size} is of the first order in the fluctuation field, hence of the order of $g_k$ up to logarithms, whereas the insertions of the class are of the order of $g_k^2$. The first order is there: it is the linear response of the transport of a loop to the fluctuation field. It is harmless because it is odd in the fluctuation field and the measure of the step is even up to its potentials (Supplement, Lemma~17.28), and what has to be shown here is that the even part of the tensor is of the second order. On the support of the test functions, $\abs{g_kA_k} \leq r_y\vartheta_k$, as was used for \eqref{eq:localised-h-size}, and therefore $\abs{B'} \leq b_k \coloneqq C_1r_y\vartheta_k$ by \cite[(1.20), p.~6]{BalabanRG2}. The middle member of that display, $O(1)g_k\abs B + 4C_2(O(1)g_k\abs B)^2$, bounds the two terms of $B' = g_kCA_k - h\tilde D(g_kCA_k)$ separately: the function $\tilde D$ ``has an expansion beginning with quadratic terms'' \cite[p.~267]{BalabanRG1}, and $\abs{h\tilde D(g_kCA_k)} \leq 4C_2(c_0g_k\abs{A_k})^2$, where $c_0$ denotes the absolute constant $O(1)$ of that member.

\begin{lemma}[The Even Part of a Tensor of One Component]\label{lem:tensor-parity}
In the setting of Lemma~\ref{lem:tensor-size}, let the group $G$ consist of one component $Y$, let $b_k \leq \frac12\varepsilon_H$, and put $S_Y \coloneqq \exp(c'_\pi\norm{a'}_\infty\abs{I(G)})\,e^{-(\kappa_1 - 2)\#(Y)}$. Let $Y_G = Y^{\mathrm{ev}}_G + Y^{\mathrm{od}}_G$ be the decomposition of the tensor, as a function of the fluctuation field $A_k$, into its even and odd parts under $A_k \mapsto -A_k$. Then, on the support of the test functions of the step,
\begin{equation}\label{eq:tensor-parity}
    \norm{Y^{\mathrm{ev}}_G} \leq c_{\mathrm{ev}}\vartheta_k^2S_Y\;,
    \qquad
    \norm{Y^{\mathrm{od}}_G} \leq \Bigl(\frac{C_1r_y}{\varepsilon_H}\vartheta_k + c_{\mathrm{ev}}\vartheta_k^2\Bigr)S_Y\;,
    \qquad
    c_{\mathrm{ev}} \coloneqq \frac{4C_2c_0^2r_y^2}{\varepsilon_H} + \frac{2C_1^2r_y^2}{\varepsilon_H^2}\;.
\end{equation}
\end{lemma}

\begin{proof}
By property (b) of the parameters, $Y_G$ depends on the fluctuation field through the restriction $X$ of $B'$ to $Y$; we write $Y_G(X)$. The proof of Lemma~\ref{lem:tensor-size}, without its first step, shows that $Y_G$ is analytic in $X$ and $\norm{Y_G(X)} \leq S_Y$ for $\abs X_Y \leq \varepsilon_H$, and $Y_G(0) = 0$. Let $\Lambda$ be its differential at $0$, a linear map, and $R(X) \coloneqq Y_G(X) - \Lambda(X)$. Cauchy's estimates for the function $t \mapsto Y_G(tX)$ on the disc of radius $\varepsilon_H/\abs X_Y$ give
\[
    \norm{\Lambda(X)} \leq S_Y\frac{\abs X_Y}{\varepsilon_H}\;,
    \qquad
    \norm{R(X)} \leq S_Y\sum_{n \geq 2}\Bigl(\frac{\abs X_Y}{\varepsilon_H}\Bigr)^n \leq 2S_Y\Bigl(\frac{\abs X_Y}{\varepsilon_H}\Bigr)^2
\]
for $\abs X_Y \leq \frac12\varepsilon_H$, the first for every $X$ by linearity. Now $B' = g_kCA_k - h\tilde D(g_kCA_k)$, where $C$ is linear \cite[p.~268]{BalabanRG1}, so that the even part of $B'$ under $A_k \mapsto -A_k$ is the even part of $-h\tilde D(g_kCA_k)$, and its restriction to $Y$ is at most $4C_2c_0^2r_y^2\vartheta_k^2$ on the support, by the bound recalled before the lemma and $\abs{g_kA_k} \leq r_y\vartheta_k$. The function $\Lambda$ of the odd part of $B'$ is odd. Hence $Y^{\mathrm{ev}}_G$ is the sum of $\Lambda$ of the even part of $B'$ and of the even part of $R(B')$, and with $\abs{B'}_Y \leq b_k = C_1r_y\vartheta_k \leq \frac12\varepsilon_H$ this gives the first bound of \eqref{eq:tensor-parity}. The second follows from $\norm{Y_G} \leq S_Yb_k/\varepsilon_H$, which is \eqref{eq:tensor-size}, and the first.
\end{proof}

A group which touches members of the family is estimated in the same way, with one difference. The insertions of the members are not entire functions of the configuration, as the transports are, so that the fluctuation field on a component may be scaled only as far as the interpolated configurations stay in the domain on which these insertions are analytic. By the argument which gave \eqref{eq:expansion-condition}, this is the case as long as $c_{\mathbf J} \cdot 4B_0e^{16\kappa_1}$ times the scaled field is at most $\frac14\min(\beta, 1)\min(\alpha_{0,k}, \alpha_{1,k})$; we put
\begin{equation}\label{eq:member-radius}
    r_\alpha \coloneqq \frac{\min(\beta, 1)\min(\alpha_{0,k}, \alpha_{1,k})}{16c_{\mathbf J}B_0e^{16\kappa_1}}\;,
\end{equation}
so that \eqref{eq:expansion-condition} reads $b_k \leq r_\alpha$. The factor which a component gains is then the ratio of its fluctuation field to $\min(\varepsilon_H, r_\alpha)$, a negative power of $\log g_k^{-2}$ in place of a power of the coupling, and the smallness of the tensor comes from the insertions of its members.

\begin{lemma}[The Size of a Tensor with Touched Members]\label{lem:tensor-members}
In the setting of Lemma~\ref{lem:tensor-size}, with a field $a'$ for which $e^{i\eta a'}U_{k+1}$ lies in the space of the step $k + 1$, let $G$ be a group of a set $\mathfrak Y$ with the set $\mathfrak B(G)$ of touched members, and suppose \eqref{eq:expansion-condition} and $b_k \leq \varepsilon_H$. Then, on the support of the test functions of the step,
\begin{equation}\label{eq:tensor-members}
    \norm{Y_G} \leq \exp\bigl(4\bar\Lambda Mc'_\pi\norm{a'}_\infty\#(G)\bigr)\prod_{Y \in G}\frac{\abs{B'}_Y}{\min(\varepsilon_H, r_\alpha)}\,e^{-(\kappa_1 - 2)\#(Y)}\prod_{W \in \mathfrak B(G)}\sup\norm{Y_W}\;,
\end{equation}
where $\norm{a'}_\infty$ is the supremum of $\abs{a'(b)}$ over the fine bonds of the loops in the components of $G$ and each supremum is taken over the configurations $e^{i\eta\mathbf H(s, tB')}e^{i\eta a'}U_{k+1}$, with the corresponding argument in the place of $\mathbf J$, for $\abs{s(\Delta)} \leq e^{\kappa_1}$ and $\abs{t_Y}\abs{B'}_Y \leq \min(\varepsilon_H, r_\alpha)$, which lie in the domain of analyticity of the factors of the step $k$.
\end{lemma}

\begin{proof}
We repeat the proof of Lemma~\ref{lem:tensor-size} with the function $\Phi(s, t)$ which is the tensor product, over the pieces of $G$, of the ordered products of the transports and of the insertions $Y_W$, $W \in \mathfrak B(G)$, at the configuration $e^{i\eta\mathbf H(s, tB')}e^{i\eta a'}U_{k+1}$, where $tB'$ denotes the field $t_YB'$ on each component $Y$ of $G$; the products are formed by multilinearity. For $\abs{t_Y}\abs{B'}_Y \leq \min(\varepsilon_H, r_\alpha)$ the field $\mathbf H(s, tB')$ is at most $4B_0e^{16\kappa_1}r_\alpha$ in the admissible norms, by \eqref{eq:localised-h-bound}, and by the definitions of $c_{\mathbf J}$ and $r_\alpha$ the configuration lies in the space of the step $k$, on which the insertions are analytic (Definition~\ref{def:dressed-loop}(iii)). So $\Phi$ is analytic on the polydisc $\abs{s(\Delta)} \leq e^{\kappa_1}$, $\abs{t_Y} \leq T_Y \coloneqq \min(\varepsilon_H, r_\alpha)/\abs{B'}_Y$, and $T_Y \geq 1$ because $\abs{B'}_Y \leq b_k \leq \min(\varepsilon_H, r_\alpha)$. The three steps are as before. If one $t_Y$ vanishes, $\mathbf H_Y = 0$, no factor depends on the parameters of the cubes of $Y$, and the derivatives vanish, so that Schwarz's lemma gives $\prod_YT_Y^{-1}$. Cauchy's estimate in the parameters gives $e^{-(\kappa_1 - 1)\#(G)}$. And the norm of $\Phi$ is at most the product of the norms of the transports, which run along the fine bonds of the loops in the components of $G$ outside the intervals of the members and are bounded as before by $e^{\#(G)}\exp(4\bar\Lambda Mc'_\pi\norm{a'}_\infty\#(G))$, and of the projective norms of the insertions of the members, since the contraction of tensors along a piece does not increase the product of the norms (Section~\ref{sec:dressed-loops}).
\end{proof}

\subsection{The Fluctuation Integral with Marked Polymers}\label{sec:marked-polymers}

By Lemma~\ref{lem:word-expansion} the function $F$ is, on the term without large-field regions, a finite sum of numbers times words, and in each word only the tensors $Y_G$ depend on the field $A$ which the step integrates, each of them through the fluctuation field on the components of its group. The composition with $\Psi_{k+1}$ in \eqref{eq:carried-loop} does not interfere, because the bonds of $S_{k+1}^{(k)*}$ are separated from $\Lambda_{k+1}$ by a layer of cubes \cite[(3.20), p.~269]{BalabanCRE}, so that $\Psi_{k+1}$ changes outer variables of the step only. The expectation of a word is therefore a contraction of the expectation of the tensor product of its tensors $Y_G$, and what we need is an expansion of such an expectation into local pieces. Ba{\l}aban makes this expansion for one localised factor under the integral: the factor marks a domain $Z_0$ of the polymer expansion, whose activity $H'(Z_0)$ contains it, while the other polymers have the activities of the expansion of the normalisation, and the sum over these is exponentiated \cite[(3.44)--(3.46), pp.~277--278]{BalabanCRE}. We need the same with several marked domains and with factors which are tensors.

\emph{The setting.} We fix a family $\mathfrak B$ and a set $\mathfrak Y$ as in Lemma~\ref{lem:word-expansion}, with the set $\mathfrak G$ of its groups, and put
\begin{equation}\label{eq:tensor-product}
    \mathcal X \coloneqq \bigotimes_{G \in \mathfrak G}Y_G\;,
\end{equation}
a function of $A$ and of the other variables with values in the tensor product of the spaces of the tensors $Y_G$. The word $w(\mathfrak B, \mathfrak Y)$ is multilinear in its factors, so that $\langle w(\mathfrak B, \mathfrak Y)\rangle$ is the contraction of $\langle\mathcal X\rangle$ with the factors which do not depend on $A$, the insertions of the untouched members and the transports of $U_{k+1}$; and $\langle\mathcal X\rangle = I(\mathcal X)/I(1)$, where $I(f)$ denotes the last integral of \cite[(3.25)]{BalabanCRE} with the further factor $f$ under it. For $G \in \mathfrak G$ let $D_G$ be the union of the components of $G$ and of the cubes of $\pi_k$ which share a point with them, and let $D^{\#} \coloneqq \bigcup_{G \in \mathfrak G}D_G$. The bonds of the scale $k$ on whose variables the tensor $Y_G$ depends lie in the interior of $D_G$. Indeed $Y_G$ depends on $A$ through the restriction of $B'$ to its components (Lemma~\ref{lem:word-expansion}), and the variable $B'(b)$ is $g_kA_k(b)$ unless $b$ is a bond $b_0(c)$, for which it is a function of the variables in the blocks of the ends of $c$ and between them, since the $\delta$-function which eliminates it is that of the average \cite[(0.4), p.~253]{BalabanRG1} on the bond $c$ \cite[pp.~267--268]{BalabanRG1}; these blocks have the side $L \leq M$ and are contained in cubes of $\pi_k$ at a distance less than $M$ from the component, and two cubes of $\pi_k$ at a distance less than $M$ share a point.

\emph{The polymers.} The expansion of \cite[Sect.~2]{BalabanRG2} writes the fluctuation integral as a sum over finite sets of \emph{polymers}, localisation domains $Z \in \mathbf D_{k+1}$, with activities $H(Z)$ and the hard core $\zeta$: $\zeta(Z, Z') = 0$ if $Z \cap Z'$ contains a cube or a wall of a cube, and $\zeta(Z, Z') = 1$ otherwise \cite[(2.8)--(2.11), p.~14]{BalabanRG2}. We call two polymers \emph{compatible} if $\zeta = 1$, and write $Z \not\sim Z'$ if they are not. For the measure \cite[(3.27)]{BalabanCRE} of the step the expansion is that of \cite[pp.~275--279]{BalabanCRE}, where ``we repeat all the considerations of Sect.~3.7 [I]'' with the cover and the localised action of \cite[pp.~275--277]{BalabanCRE}; the section of [I] which is meant is \cite[Sect.~2]{BalabanRG2}, whose formulae (2.1)--(2.13) are cited there as (I.7.1)--(I.7.13) (Section~\ref{sec:notes-on-papers}). We call the polymers of the expansion of the normalisation $I(1)$ \emph{vacuum polymers}. A \emph{cluster} is a finite sequence $\mathcal C = (Z_1, \dots, Z_n)$ of vacuum polymers whose incompatibility graph is connected, with the weight
\begin{equation}\label{eq:vacuum-cluster-weight}
    \varphi^T(\mathcal C) \coloneqq \frac{1}{n!}\rho^T(Z_1, \dots, Z_n)H(Z_1) \cdots H(Z_n)\;,
\end{equation}
where $\rho^T$ is the function of \cite[(2.12), p.~14]{BalabanRG2}, which vanishes unless that graph is connected; a cluster is incompatible with a union $X$ of cubes, $\mathcal C \not\sim X$, if one of its polymers is. Ba{\l}aban exponentiates the expansion of the normalisation, ``the factor is represented as $\exp\tilde{\mathbf E}^{(k+1)}(\Lambda_{k+1})$'', and the sum over the polymers compatible with a marked domain as well \cite[pp.~277--278]{BalabanCRE}, and ``the other expectation values in (3.26), (3.28) are treated in exactly the same way'' \cite[p.~279]{BalabanCRE}, the first of which is the expectation in the measure (3.27). We use this in the following form: the series $\sum_{\mathcal C \not\sim X}\abs{\varphi^T(\mathcal C)}$ converges for every union $X$ of cubes of the torus, and the sum of the products of the activities over the finite sets of pairwise compatible vacuum polymers which are compatible with $X$ equals $\exp\sum_{\mathcal C}\varphi^T(\mathcal C)$, the sum over the clusters none of whose polymers is incompatible with $X$. Both statements, with the bound $\sum_{\mathcal C \not\sim X}\abs{\varphi^T(\mathcal C)}e^{d(\mathcal C)} \leq cN(X)$, where $N(X)$ is the number of the cubes of $\pi_{k+1}$ in $X$, $c = \frac{9}{64}e^{-4\kappa}$, and $d(\mathcal C)$ is $9\log 2$ times the total number of the cubes of the polymers of $\mathcal C$, are derived from Ba{\l}aban's bound \cite[(2.38), p.~20]{BalabanRG2} of the activities in Supplement, Lemma~17.24, for $\kappa \geq 2^9(1 + 16\log 2)$.

\begin{proposition}[Marked Polymers]\label{prop:marked-polymers}
In the setting above, call two connected components of a union $Z$ of cubes of $\pi_{k+1}$ \emph{linked} if a set $D_G$, $G \in \mathfrak G$, meets both, and call $Z$ linked if this relation connects all its components. Then
\begin{equation}\label{eq:marked-expansion}
    \langle\mathcal X\rangle = \sum_{\mathfrak M}Q_{k+1}(\mathfrak M)\bigotimes_{M \in \mathfrak M}\hat H(M)\;,
\end{equation}
where:
\begin{enumerate}[label=(\roman*)]
\item the sum runs over the finite sets $\mathfrak M$ of pairwise compatible \emph{marked polymers}: linked unions $M$ of cubes of $\pi_{k+1}$, each with a nonempty set $\mathfrak G(M) \subset \mathfrak G$ of groups whose sets $D_G$ it contains, such that the sets $\mathfrak G(M)$, $M \in \mathfrak M$, form a partition of $\mathfrak G$;
\item $\hat H(M) \coloneqq e^{-\Psi(M)}\tilde H(M)$, where the \emph{marked activity} $\tilde H(M)$ is the activity of $M$ in the expansion \cite[(2.1)--(2.11)]{BalabanRG2} of the integral with the factor $\bigotimes_{G \in \mathfrak G(M)}Y_G$ under it, a tensor with one factor for each piece of each group of $\mathfrak G(M)$, and $\Psi(M) \coloneqq \sum_{\mathcal C \not\sim M}\varphi^T(\mathcal C)$;
\item with $\Psi_{\geq}(T) \coloneqq \sum\varphi^T(\mathcal C)$, the sum over the clusters which are incompatible with every member of a set $T$ of marked polymers,
\begin{equation}\label{eq:vacuum-interaction}
    Q_{k+1}(\mathfrak M) \coloneqq \exp\Bigl(\sum_{T \subset \mathfrak M,\ \abs T \geq 2}(-1)^{\abs T}\Psi_{\geq}(T)\Bigr)\;.
\end{equation}
\end{enumerate}
The tensor $\hat H(M)$ depends on $M$ and on the groups of $\mathfrak G(M)$ alone, and $\Psi_{\geq}(T)$ on $T$ alone.
\end{proposition}

\begin{proof}
We go through the construction \cite[(2.1)--(2.11), pp.~11--14]{BalabanRG2} for the integral $I(\mathcal X)$ and note at each step what the factor $\mathcal X$ changes.

\emph{The Mayer expansion} \cite[(2.1)]{BalabanRG2} of the exponential of the localised action is an identity for the exponential alone, with $\mathcal X$ as a further factor of the integrand. In place of \cite[(2.2)]{BalabanRG2} we put $Y_0 \coloneqq \bigcup_{Y \in \mathbf D}Y \cup D^{\#}$.

\emph{The test functions.} The test function $\chi^{(k)}$ is a product of functions of single bond variables \cite[(3.16), p.~268]{BalabanCRE}, and the decomposition \cite[(2.3)]{BalabanRG2} of the factors of the bonds outside $Y_0$ applies to it, with the thresholds of the construction; $Z_0$ is the smallest localisation domain which contains $Y_0$ and the set $P$ of that decomposition, as in \cite[p.~12]{BalabanRG2}.

\emph{The conditioning} \cite[(2.4)--(2.6)]{BalabanRG2} on the variables outside $Z_0$ uses of the integrand only that it is localised in $Z_0$ with respect to the integration variable. Since $D^{\#} \subset Z_0$, the factor $\mathcal X$ is a function of the variable which is conditioned, so that \cite[(2.5), (2.6)]{BalabanRG2} hold with $\mathcal X$ inside the function $G$ of \cite[(2.6)]{BalabanRG2}.

\emph{The localisation} \cite[(2.8), (2.9)]{BalabanRG2} introduces parameters in the operators $C^{(k)}(Z_0)$, $(C^{(k)})^{1/2}$, and $\Delta_k$, none of which occurs in $\mathcal X$; it gives $I(\mathcal X) = \sum_ZH(Z)$ with activities which now contain $\mathcal X$.

\emph{The factorisation.} Ba{\l}aban's activities factorise over the connected components of $Z$ \cite[(2.10)]{BalabanRG2}, because every factor of the integrand, the Mayer factors, the factors of the test functions, and the conditioned Gaussian integrals at vanishing parameters outside $Z$, splits over them. The factor $\mathcal X$ splits over the groups, and the tensor $Y_G$ of a group depends on the integration variable in $D_G$ only, which may meet several components. Hence the activity factorises over the classes of linked components: a class which contains no set $D_G$ has the activity of a vacuum polymer, and a class $M$ which contains the sets $D_G$ of the groups $G \in \mathfrak G(M) \neq \emptyset$ has an activity $\tilde H(M)$ with values in the tensor product of the spaces of these $Y_G$. A set $D_G$ which meets a class is contained in it, because $D_G \subset D^{\#} \subset Z_0 \subset Z$ and the components of $Z$ which $D_G$ meets are linked. As in \cite[(2.11)]{BalabanRG2} we obtain
\begin{equation}\label{eq:marked-polymer-sum}
    I(\mathcal X) = \sum_{\mathfrak M}\Bigl(\bigotimes_{M \in \mathfrak M}\tilde H(M)\Bigr)\,\Xi(\mathfrak M)\;,
\end{equation}
the sum over the sets $\mathfrak M$ of (i), where $\Xi(\mathfrak M)$ is the sum of the products of the activities over the finite sets of pairwise compatible vacuum polymers which are compatible with every $M \in \mathfrak M$; and $I(1) = \Xi(\emptyset)$.

\emph{The ratio.} By the exponentiation in the form stated before the proposition, applied to the union of the members of $\mathfrak M$ and to the empty set,
\begin{equation}\label{eq:vacuum-ratio}
    \frac{\Xi(\mathfrak M)}{\Xi(\emptyset)} = \exp\Bigl(-\sum_{\mathcal C \not\sim \mathfrak M}\varphi^T(\mathcal C)\Bigr)\;,
\end{equation}
the sum over the clusters which are incompatible with at least one member of $\mathfrak M$. For such a cluster let $s(\mathcal C) \geq 1$ be the number of the members of $\mathfrak M$ with which it is incompatible. It occurs $s(\mathcal C)$ times in $\sum_{M \in \mathfrak M}\Psi(M)$ and once in the exponent of \eqref{eq:vacuum-ratio}, and
\[
    \sum_{T \subset \{ M \,|\, \mathcal C \not\sim M \},\ \abs T \geq 2}(-1)^{\abs T} = (1 - 1)^{s(\mathcal C)} - 1 + s(\mathcal C) = s(\mathcal C) - 1\;,
\]
so that it occurs with the coefficient $s(\mathcal C) - 1$ in the exponent of \eqref{eq:vacuum-interaction}. The rearrangement is permitted, because a cluster lies in at most $2^{s(\mathcal C)}$ of the sums $\Psi_{\geq}(T)$, the torus has finitely many cubes, so that $s(\mathcal C) \leq \abs{\mathfrak M}$, and $\sum_{\mathcal C \not\sim \mathfrak M}\abs{\varphi^T(\mathcal C)}$ converges. Hence
\begin{equation}\label{eq:vacuum-ratio-split}
    \frac{\Xi(\mathfrak M)}{\Xi(\emptyset)} = \prod_{M \in \mathfrak M}e^{-\Psi(M)} \cdot Q_{k+1}(\mathfrak M)\;,
\end{equation}
and dividing \eqref{eq:marked-polymer-sum} by $I(1) = \Xi(\emptyset)$ gives \eqref{eq:marked-expansion}. The last statement holds because the activity of a class of linked components is defined by the integral over the variables of that class.
\end{proof}

A marked polymer may carry several groups, and a group several pieces, on one loop or on several, so that the tensors $\hat H(M)$ are the first objects of the step with several intervals. The marked polymers of a set $\mathfrak M$ have a hard core among themselves and interact softly through $Q_{k+1}$, which is a product of factors, one for each set $T$ of at least two of them; they have no hard core with the domains of the untouched members of $\mathfrak B$, which took no part in the expansion.

\emph{The size of a marked activity.} How small is $\tilde H(M)$? Ba{\l}aban's bound for the activities of the vacuum polymers cannot hold for it, since the factor $\mathcal X_M \coloneqq \bigotimes_{G \in \mathfrak G(M)}Y_G$ has a prescribed support in which nothing decays, and since one of the terms of the activity is the conditional expectation of this factor, which carries no small constant. What holds instead is a bound by the supremum of the factor, with a decay in the part of $M$ outside its support and a price for each cube of the support (Supplement, Lemma~17.27). The supremum is small by the lemmata of Section~\ref{sec:localisation-expansion}, with one exception: a factor which consists of the tensor of a single component without members is of the first order in the fluctuation field. Its odd part enters the activity only through the odd part of the potentials of the fluctuation integral (Supplement, Lemma~17.28), and this gives the second order. For a marked polymer $M$ let $\mathfrak B(M) \coloneqq \bigcup_{G \in \mathfrak G(M)}\mathfrak B(G)$ be the set of its touched members, $n(M)$ the number of the components of its groups, $\#(G_M) \coloneqq \sum_{G \in \mathfrak G(M)}\#(G)$ the number of their cubes, which are cubes of $\pi_k$, and $\#(M)$ the number of the cubes of $\pi_{k+1}$ in $M$. With the constant $A$ of Supplement, Lemma~17.27, and the number $y_k$ of Supplement, Lemma~17.28, put
\begin{equation}\label{eq:marked-size-constants}
    c_\kappa \coloneqq 1 + 81\bigl(5 + \tfrac32L\bigr)\kappa + \tfrac{6561}{16}(L + 2)\kappa + 1594323\,A\;,
    \qquad
    a^{(1)}_k \coloneqq c_{\mathrm{ev}}\vartheta_k^2 + \frac{L^4y_k}{2A}\Bigl(\frac{C_1r_y}{\varepsilon_H}\vartheta_k + c_{\mathrm{ev}}\vartheta_k^2\Bigr)\;.
\end{equation}

\begin{proposition}[The Size of a Marked Activity]\label{prop:marked-size}
On the terms without large-field regions of $\varrho_k$ and of $\varrho_{k+1}$, suppose that the hypotheses of Supplement, Lemmata~17.27 and~17.28, hold, together with \eqref{eq:expansion-condition}, $b_k \leq \frac12\min(\varepsilon_H, r_\alpha)$, and $4\bar\Lambda Mc'_\pi\alpha_{1,k+1} \leq L$. Then, for every marked polymer $M$ and at every configuration of the space of the step $k + 1$,
\begin{equation}\label{eq:marked-size}
    \norm{\tilde H(M)} \leq \mathfrak a(M)\,e^{-(\kappa_1 - 2 - c_\kappa)\#(G_M)}\,e^{-(L + 2)\kappa\#(M)/1296}\prod_{W \in \mathfrak B(M)}\norm{Y_W}_{(k)}\;,
\end{equation}
where $\mathfrak a(M) \coloneqq a^{(1)}_k$ if $\mathfrak B(M) = \emptyset$ and $n(M) = 1$, $\mathfrak a(M) \coloneqq (b_k/\varepsilon_H)^{n(M)}$ if $\mathfrak B(M) = \emptyset$ and $n(M) \geq 2$, and $\mathfrak a(M) \coloneqq (b_k/\min(\varepsilon_H, r_\alpha))^{n(M)}$ if $\mathfrak B(M) \neq \emptyset$. Moreover $a^{(1)}_k \leq c_a\rho_k$ and $(b_k/\varepsilon_H)^2 \leq c_a\rho_k$, with a constant $c_a$ which depends on the constants of the construction and of \eqref{eq:tensor-parity} and not on $k$.
\end{proposition}

\begin{proof}
Let $D \coloneqq \bigcup_{G \in \mathfrak G(M)}D_G$, with $m$ connected components, and let $D'$ be the smallest union of cubes of $\pi_{k+1}$ which contains $D$ and the cubes of $\pi_k$ that touch it, as in Supplement, Section~17. We first count. The set $D_G$ consists of cubes of the components of $G$ and of cubes which share a point with them, so that $D$ has at most $81\,\#(G_M)$ cubes of $\pi_k$, and at most as many connected components; the cubes of $D$ and those which touch it are at most $81^2\,\#(G_M)$, and each cube of $D'$ contains one of them. Hence
\begin{equation}\label{eq:support-counts}
    m \leq 81\,\#(G_M)\;,
    \qquad
    \#(D') \leq 6561\,\#(G_M)\;.
\end{equation}

By Supplement, Lemma~17.27, applied to $M$ with the factor $\mathcal X_M$, and by the remark which follows that lemma,
\[
    \norm{\tilde H(M)} \leq \norm{\mathcal X_M}_{\sup}\,e^{81A\#(D')}\,e^{(5 + \frac32L)\kappa(m - 1)}\,e^{-(L + 2)\kappa(\#(M) - 81\#(D'))/1296}\;,
\]
where $\norm{\mathcal X_M}_{\sup}$ is the supremum of the projective norm of $\mathcal X_M$ over the fluctuation fields of the support of the test functions and over the space of the step $k + 1$. By \eqref{eq:support-counts} the three exponential factors other than $e^{-(L + 2)\kappa\#(M)/1296}$ are together at most $\exp((c_\kappa - 1)\#(G_M))$, where we have bounded $81A$ by $243A$ and used $81 \cdot 6561/1296 = 6561/16$ and $243 \cdot 6561 = 1594323$.

The norm of a tensor product is at most the product of the norms, so that $\norm{\mathcal X_M}_{\sup} \leq \prod_{G \in \mathfrak G(M)}\sup\norm{Y_G}$. On the support of the test functions $\abs{B'}_Y \leq b_k$ for every component $Y$ (Section~\ref{sec:localisation-expansion}), and a configuration $e^{i\eta a'}U_{k+1}$ of the space of the step $k + 1$ has $\abs{a'(b)} \leq \alpha_{1,k+1}/L$ on the fine bonds, in the units of the step $k$, by the radii of Definition~\ref{def:dressed-loop}(iii), so that $\exp(4\bar\Lambda Mc'_\pi\norm{a'}_\infty\#(G)) \leq e^{\#(G)}$ by the last hypothesis. For a group $G$ without members, which is a single component, Lemma~\ref{lem:tensor-size} therefore gives $\norm{Y_G} \leq (b_k/\varepsilon_H)e^{-(\kappa_1 - 3)\#(G)}$, and for a group with members Lemma~\ref{lem:tensor-members} gives
\[
    \norm{Y_G} \leq \prod_{Y \in G}\frac{b_k}{\min(\varepsilon_H, r_\alpha)}\cdot e^{-(\kappa_1 - 3)\#(G)}\prod_{W \in \mathfrak B(G)}\norm{Y_W}_{(k)}\;,
\]
since the configurations over which that lemma takes the suprema lie in the space of the step $k$. Multiplying over the groups, and using $b_k/\varepsilon_H \leq b_k/\min(\varepsilon_H, r_\alpha) \leq 1$, we obtain \eqref{eq:marked-size} in the second and the third case.

In the first case $\mathcal X_M = Y_G$ for a group which consists of one component $Y$ without members. By Supplement, Lemma~17.28, the bound of Lemma~17.27 holds with $\norm{\mathcal X^{\mathrm{ev}}}_{\sup} + (L^4y_k/(2A))\norm{\mathcal X^{\mathrm{od}}}_{\sup}$ in the place of $\norm{\mathcal X_M}_{\sup}$ and $243A$ in the place of $81A$, where $\mathcal X^{\mathrm{ev}}$ and $\mathcal X^{\mathrm{od}}$ are the even and the odd part of the factor under the reflection of the fluctuation field. By Lemma~\ref{lem:tensor-parity}, whose hypothesis $b_k \leq \frac12\varepsilon_H$ holds, these parts are at most $c_{\mathrm{ev}}\vartheta_k^2S_Y$ and $(C_1r_y\vartheta_k/\varepsilon_H + c_{\mathrm{ev}}\vartheta_k^2)S_Y$ with $S_Y \leq e^{-(\kappa_1 - 3)\#(Y)}$, which gives \eqref{eq:marked-size} with $\mathfrak a(M) = a^{(1)}_k$.

For the last assertion, $y_k = 4\alpha_4(\vartheta_k/\varepsilon_1)(8M^4t_k^2 + K_{\mathbf D})$ with $t_k = \vartheta_k/g_k$ (Supplement, Section~17), and $\vartheta_k \leq 1$, so that $a^{(1)}_k \leq c\,\vartheta_k^2(1 + t_k^2)$ with $c \coloneqq c_{\mathrm{ev}} + (2L^4\alpha_4/(A\varepsilon_1))\max(8M^4, K_{\mathbf D})(C_1r_y/\varepsilon_H + c_{\mathrm{ev}})$. By \eqref{eq:tube-threshold-ratio}, $\vartheta_k \leq g_kr_k$, hence $t_k \leq r_k$ and $\vartheta_k^2(1 + t_k^2) \leq g_k^2(1 + r_k)^4 = \rho_k/C_\rho$ by \eqref{eq:class-parameters}. Likewise $(b_k/\varepsilon_H)^2 = (C_1r_y/\varepsilon_H)^2\vartheta_k^2 \leq (C_1r_y/\varepsilon_H)^2\rho_k/C_\rho$, and we take $c_a \coloneqq \max(c, (C_1r_y/\varepsilon_H)^2)/C_\rho$.
\end{proof}

The bound \eqref{eq:marked-size} shows where the smallness of a new group comes from. A marked polymer without members is of the order of $\rho_k$, the square of the coupling up to logarithms. One with members is small through the insertions of the members, which are of that order themselves, and gains for each of its components a factor $b_k/\min(\varepsilon_H, r_\alpha)$, which tends to zero with the coupling as a power of its logarithm, because $q_0 > p_0$. The exponent $c_\kappa$ is the price of the support, and the rate $\kappa_1$ of the localisation parameters, which is chosen after $\kappa$ and $L$ and before $M$ \cite[p.~5]{BalabanRG2}, is taken so large that it pays for it and for the counting of the next section.

\subsection{The Groups of the Step}\label{sec:groups-of-step}

In this subsection the labelled term of $\varrho_{k+1}$ is the term without large-field regions. By Lemma~\ref{lem:word-expansion} and Proposition~\ref{prop:marked-polymers}, applied to every word of $F$,
\begin{equation}\label{eq:pre-grouping}
    \langle F\rangle = \prod_{i=1}^{n}\kappa^{(k)}(C_i, \chi_i)\cdot e^{\mathcal E_k}\sum_{\mathfrak B}Q_k(\mathfrak B)\sum_{\mathfrak Y}\sum_{\mathfrak M}Q_{k+1}(\mathfrak M)\,w(\mathfrak B, \mathfrak Y, \mathfrak M)\;,
\end{equation}
where $Q_k(\emptyset) \coloneqq 1$, the sums run over the families of the step $k$, the sets $\mathfrak Y$ of Lemma~\ref{lem:word-expansion}, and the sets $\mathfrak M$ of Proposition~\ref{prop:marked-polymers}, and $w(\mathfrak B, \mathfrak Y, \mathfrak M)$ is the word in which the tensors $\hat H(M)$ are contracted, in the cyclic order of each loop, with the insertions at $U_{k+1}$ of the untouched members of $\mathfrak B$ and the transports of $U_{k+1}$. This is not yet the form of a dressed loop, for three reasons. The tensors $\hat H(M)$ have several pieces, between which other factors of the word may lie. The terms are tied together by two soft interactions, $Q_k$ and $Q_{k+1}$, the second of which depends on the fields. And the marked polymers have no hard core with the untouched members, whose domains they may meet. All three defects are removed by a regrouping of the terms of \eqref{eq:pre-grouping} which changes none of them: the soft interactions are expanded, and whatever an interaction, an overlap, or the order along a loop ties together is collected into one group of the step $k + 1$. We call two unions of cubes of the partitions of the construction \emph{compatible} if their intersection has dimension at most $d - 2$, so that it contains neither a cube nor a wall of a cube; for two polymers this is the condition $\zeta = 1$.

\emph{The flat parts of the vacuum weights and the links.} The activities of the vacuum polymers are analytic functions of a configuration $(\mathbf U, \mathbf J)$ on the spaces of the step $k + 1$, with bounds which are uniform on these spaces \cite[pp.~6, 15]{BalabanRG2}, and the reference data $(1, 0)$ belong to them. Let $\varphi^T_\circ(\mathcal C)$ be the weight \eqref{eq:vacuum-cluster-weight} of a cluster at the reference data and $y_{\mathcal C} \coloneqq \varphi^T(\mathcal C) - \varphi^T_\circ(\mathcal C)$, a function of the background field on the union $\operatorname{supp}\mathcal C$ of the polymers of $\mathcal C$, since each activity ``is localized in the interior of $Z$ with respect to the external gauge fields'' \cite[p.~14]{BalabanRG2}; the statements on the cluster series used in Proposition~\ref{prop:marked-polymers} hold at the reference data as well. We write $\Psi_\circ$, $\Psi_{\geq,\circ}$, and $Q_\circ$ for the quantities of that proposition with $\varphi^T_\circ$ in the place of $\varphi^T$. For a set $\mathfrak M$ of marked polymers, a \emph{link} is a pair $\lambda = (T, \mathcal C)$ of a nonempty subset $T \subset \mathfrak M$ and a cluster $\mathcal C$ which is incompatible with every member of $T$; its \emph{domain} is $\operatorname{supp}\mathcal C$, its \emph{factor} is
\begin{equation}\label{eq:link-factor}
    u_\lambda \coloneqq \exp\bigl((-1)^{\abs T}y_{\mathcal C}\bigr) - 1\;,
\end{equation}
which vanishes at the reference data, and it is \emph{carried} by the first member of $T$ in a fixed order of the unions of cubes. Since $\prod_Me^{-\Psi(M)} \cdot Q_{k+1}(\mathfrak M)$ is the exponential of $\sum_{T \neq \emptyset}(-1)^{\abs T}\Psi_{\geq}(T)$, where $\Psi_{\geq}(\{ M \}) = \Psi(M)$, and $\Psi_{\geq}(T) - \Psi_{\geq,\circ}(T)$ is the sum of $y_{\mathcal C}$ over the clusters incompatible with every member of $T$,
\begin{equation}\label{eq:link-expansion}
    \prod_{M \in \mathfrak M}e^{-\Psi(M)} \cdot Q_{k+1}(\mathfrak M) = \prod_{M \in \mathfrak M}e^{-\Psi_\circ(M)} \cdot Q_\circ(\mathfrak M)\prod_\lambda(1 + u_\lambda) = \prod_{M \in \mathfrak M}e^{-\Psi_\circ(M)} \cdot Q_\circ(\mathfrak M)\sum_{\mathcal L}\prod_{\lambda \in \mathcal L}u_\lambda\;,
\end{equation}
the sum over the finite sets $\mathcal L$ of links. The product converges absolutely, because a cluster belongs to at most $2^{\abs{\mathfrak M}}$ links, $\abs{u_\lambda} \leq \abs{y_{\mathcal C}}e^{\abs{y_{\mathcal C}}}$, and $\sum_{\mathcal C \not\sim \mathfrak M}(\abs{\varphi^T(\mathcal C)} + \abs{\varphi^T_\circ(\mathcal C)})$ converges. The flat factor is a number, and as in the proof of Proposition~\ref{prop:marked-polymers} it equals $\exp\sum_{\mathcal C}(s(\mathcal C) - 1)\varphi^T_\circ(\mathcal C)$, the sum over the clusters which are incompatible with $s(\mathcal C) \geq 2$ members of $\mathfrak M$; hence
\begin{equation}\label{eq:vacuum-occurrences}
    Q_\circ(\mathfrak M) = \sum_{\mathcal O}\prod_{\mathcal C \in \mathcal O}v_{\mathcal C}\;,
    \qquad
    v_{\mathcal C} \coloneqq \exp\bigl((s(\mathcal C) - 1)\varphi^T_\circ(\mathcal C)\bigr) - 1\;,
\end{equation}
the sum over the finite sets $\mathcal O$ of such clusters, the \emph{vacuum occurrences}, again absolutely convergent.

\emph{The old occurrences.} By Definition~\ref{def:dressed-loop}, $Q_k(\mathfrak B)$ is the exponential of $-\sum_\Gamma\varphi_k(\Gamma)$ over the clusters $\Gamma$ of the reference gas of the step $k$ which are incompatible with a member of $\mathfrak B$, and
\begin{equation}\label{eq:old-occurrences}
    Q_k(\mathfrak B) = \sum_{m \geq 0}\frac{(-1)^m}{m!}\sum_{(\Gamma_1, \dots, \Gamma_m)}\prod_{j=1}^{m}\varphi_k(\Gamma_j)\;,
\end{equation}
the inner sum over the sequences of such clusters, the \emph{old occurrences}. The series converges absolutely, because the sum of $\abs{\varphi_k(\Gamma)}$ over the clusters incompatible with a given group is finite (for $k \geq 2$ by Lemma~\ref{lem:root-bound} at the step $k$, and for $k = 1$ by the corresponding bound of the first step).

Inserting \eqref{eq:link-expansion}, \eqref{eq:vacuum-occurrences}, and \eqref{eq:old-occurrences} into \eqref{eq:pre-grouping}, we obtain a sum over \emph{terms} $t = (\mathfrak B, \mathfrak Y, \mathfrak M, \mathcal L, \mathcal O, (\Gamma_j))$. The \emph{constituents} of a term are its marked polymers, the members of $\mathfrak B$, and the groups of its old occurrences. A vacuum occurrence, like a link, is \emph{carried} by the first, in the fixed order, of the marked polymers with which it is incompatible. A constituent has a \emph{domain}: a marked polymer $M$ the set $M$ together with the domains of the links which it carries, and a group $W$ its dependence set $X^+_W$; two such sets are compatible, as before, if their intersection has dimension at most $d - 2$. It also has \emph{pieces} on the loops. The pieces of a group are its intervals. If the arc sets of the groups of $\mathfrak G(M)$ lie on one loop $C_i$, the marked polymer $M$ has one piece, the arc $I_i(M)$ of $C_i$ obtained by removing the longest maximal run of bonds which are not in these arc sets, the first such run from the base point of $C_i$ if several are longest, and $I_i(M) \coloneqq C_i$ if they cover $C_i$; otherwise its pieces are the pieces of the groups of $\mathfrak G(M)$. A constituent \emph{lies on the loop} $C_i$ if all its pieces are arcs of $C_i$, and it then has a single piece, since a group which is not a cross group has one interval (Section~\ref{sec:dressed-loops}).

\begin{definition}[The Groups of the Step]\label{def:groups-of-step}
Two constituents of a term are \emph{tied} in each of the following cases.
\begin{enumerate}[label=(\roman*)]
\item A member of $\mathfrak B$ is touched by a component of a group of $\mathfrak G(M)$: it is tied to $M$.
\item A marked polymer belongs to the set $T$ of a link of $\mathcal L$, or is incompatible with a vacuum occurrence of $\mathcal O$: it is tied to the marked polymer which carries the link or the occurrence.
\item The domains of a marked polymer and of a group, a member of $\mathfrak B$ or a group of an old occurrence, are incompatible.
\item The domain of a link carried by a marked polymer is incompatible with another marked polymer or with the domain of a link which that one carries.
\item Two groups of one old occurrence are incompatible, or a group of an old occurrence is incompatible with a member of $\mathfrak B$; here the relation is the compatibility of the groups of the step $k$, and two constituents which are the same group of the step $k$ count as incompatible.
\end{enumerate}
The \emph{arc set} of a set of constituents is the union of their pieces and, for each tied pair among them both members of which lie on the same loop $C_i$ and have disjoint pieces, of the shorter of the two arcs of $C_i$ between these pieces; if the two arcs have the same length, the one which follows the earlier piece in the orientation of $C_i$ is taken. No arc is added for a tied pair one of whose members has pieces on several loops. The \emph{groups of the step $k + 1$} in the term $t$ are the classes of the smallest equivalence relation on its constituents which contains the tied pairs and for which, in addition, two classes whose arc sets share a bond coincide, and a member of $\mathfrak B$ one of whose intervals lies in the arc set of a class belongs to that class. For a group $\mathcal G$ we write $I(\mathcal G)$ for its arc set, whose maximal arcs are its \emph{intervals}; its localisation domains are the unions of the cubes of $\pi_{k+1}$ which meet the domain of one of its constituents or a polymer of one of its vacuum occurrences; and its \emph{insertion} $Y_{\mathcal G}$ is the contraction, along its intervals and in the order of each loop, of the tensors $e^{-\Psi_\circ(M)}\tilde H(M)$ of its marked polymers, the insertions at $U_{k+1}$ of its untouched members, and the transports of $U_{k+1}$ between them, multiplied by the factors $u_\lambda$ of the links which its marked polymers carry, the factors $v_{\mathcal C}$ of its vacuum occurrences, and $((-1)^m/m!)\prod_j\varphi_k(\Gamma_j)$ for the sequence of its $m$ old occurrences. Two groups, each given with its constituents, links, and occurrences, are \emph{compatible} if their arc sets are disjoint; the domain of every marked polymer of the one is compatible with the domain of every constituent of the other; every member of $\mathfrak B$ of the one is compatible, as a group of the step $k$, with every member of $\mathfrak B$ of the other; no vacuum occurrence of the one is incompatible with a marked polymer of the other; and no group of an old occurrence of the one is incompatible, as a group of the step $k$, with a member of $\mathfrak B$ of the other. Two old occurrences of different groups are thus not required to be compatible with each other, as the series \eqref{eq:old-occurrences} imposes no condition between its clusters.
\end{definition}

The equivalence relation exists and is unique: the relation in which all constituents are equivalent satisfies the conditions, and the conditions are preserved by intersections of equivalence relations, since the arc set of a set of constituents increases with the set. A link or an occurrence belongs to one group: the marked polymers of the set $T$ of a link are tied to its carrier by (ii), and so are the marked polymers with which a vacuum occurrence is incompatible; and any two groups of an old occurrence are joined by a chain of groups of the occurrence in which consecutive ones are incompatible, the incompatibility graph of a cluster being connected, so that (v) ties them. An interval $J$ of a member $W$ of $\mathfrak B$ which is not in a group $\mathcal G$ is disjoint from $I(\mathcal G)$. Indeed, $J$ does not lie in $I(\mathcal G)$, by the definition of the groups. If it met $I(\mathcal G)$, it would contain an end bond of a maximal arc of $I(\mathcal G)$ and a bond outside $I(\mathcal G)$. That end bond belongs to a piece of a constituent of $\mathcal G$, because an arc which is added between two pieces has a piece at each of its ends; so $J$ would contain a bond of a piece without lying in that piece. For a piece of a marked polymer $M$ this is excluded by the first statement on the sets $S'_i(Y)$ in Section~\ref{sec:localisation-expansion}, by which an interval of a member that contains a bond of such a set lies in it: if $M$ has pieces on several loops, a bond of a piece is a bond of a set $S'_i(Y)$, which is contained in that piece, and if $M$ lies on one loop, $J$ contains an end bond of $I_i(M)$, which is a bond of a set $S'_i(Y) \subset I_i(M)$. It is excluded for an interval of a member of $\mathfrak B$, because the intervals of a family are disjoint; and for an interval of a group of an old occurrence, because that group would then be incompatible with $W$, two groups whose intervals share a bond being incompatible, and $W$ would belong to $\mathcal G$ by (v). So the arc sets of the groups of a term are pairwise disjoint, by the definition; every interval of a constituent lies in the arc set of its group; and no interval of a member meets the arc set of another group. Consequently the word of the term is the word, in the sense of Section~\ref{sec:dressed-loops}, of the insertions $Y_{\mathcal G}$ at the intervals of the groups, with the transports of $U_{k+1}$ between them; this uses only that the word is an ordered product along each loop which is multilinear in its factors, and that the numerical factors of a term are distributed over the groups as in the definition, the factor $1/m!$ of \eqref{eq:old-occurrences} because the sequences $(\Gamma_j)$ which induce given sequences in the groups are $m!/\prod_{\mathcal G}m_{\mathcal G}!$ in number.

\begin{proposition}[The Regrouping]\label{prop:regrouping}
On the terms without large-field regions of $\varrho_k$ and of $\varrho_{k+1}$, suppose that the sum over the terms $t$ of the products of the absolute values of their numerical factors and of the norms of their tensors converges. Then
\begin{equation}\label{eq:regrouped}
    \langle F\rangle = \prod_{i=1}^{n}\kappa^{(k)}(C_i, \chi_i)\cdot e^{\mathcal E_k}\Bigl[\prod_{i=1}^{n}f^{(i)}_{k+1} + \sum_{\mathfrak F \neq \emptyset}w_{k+1}(\mathfrak F)\Bigr]\;,
\end{equation}
the sum over the finite sets $\mathfrak F$ of pairwise compatible groups of the step $k + 1$, each with its constituents, links, and occurrences, where $w_{k+1}(\mathfrak F)$ is the word of the insertions $Y_{\mathcal G}$, $\mathcal G \in \mathfrak F$. Each $Y_{\mathcal G}$ depends on the background field $U_{k+1}$ only on the domains of the constituents of $\mathcal G$ and along its intervals.
\end{proposition}

\begin{proof}
By the convergence assumed, the series obtained by inserting \eqref{eq:link-expansion}, \eqref{eq:vacuum-occurrences}, and \eqref{eq:old-occurrences} into \eqref{eq:pre-grouping} may be summed in any order, and by the remarks after Definition~\ref{def:groups-of-step} the contribution of a term $t$ is $w_{k+1}(\mathfrak F)$ for the set $\mathfrak F$ of its groups; the term without constituents gives the product of the leading words. It remains to see that $t \mapsto \mathfrak F$ is a bijection onto the sets of pairwise compatible groups, the sequences of old occurrences being recorded group by group. The groups of a term are pairwise compatible: their arc sets are disjoint, as was shown; the marked polymers of $\mathfrak M$ are pairwise compatible and the members of $\mathfrak B$ are pairwise compatible as groups of the step $k$, by the definitions of these sets; two constituents in different groups are not tied, which gives the compatibility of the domain of a marked polymer with the domain of a group by (iii), and with that of another marked polymer by (iv); and the conditions on the vacuum and on the old occurrences hold by (ii) and (v). Conversely, let $\mathfrak F$ be a set of pairwise compatible groups, and let $t$ be the union of their data. Its marked polymers are pairwise compatible, so that their components, which lie in the sets $D_G$ of their groups, form a set $\mathfrak Y$ of Lemma~\ref{lem:word-expansion}; the members form a family; and the number $s(\mathcal C)$ of a vacuum occurrence, and with it the marked polymer which carries it, is the same in $t$ as in its group, since it is incompatible with no marked polymer of another group. No constituent of one group is tied in $t$ to one of another: for (i), a component which touches a member $W$ lies in a marked polymer, and the intersection of that polymer with $X^+_W$ then contains an open set, so that the two domains are incompatible; for (ii), a link belongs to one group and a vacuum occurrence is compatible with the marked polymers of the other groups; for (iii) and (iv) by the compatibility of the domains of the marked polymers with those of the other constituents; for (v) by the condition on the old occurrences, each of which is contained in one group. Their arc sets are disjoint, so that no member has an interval in the arc set of another group, and no further identification occurs. Hence the groups of $t$ are the members of $\mathfrak F$. The locality of $Y_{\mathcal G}$ follows from that of the tensors $\tilde H(M)$ (Proposition~\ref{prop:marked-polymers}), of the factors $u_\lambda$, and of the insertions of the members, the other factors being numbers.
\end{proof}

Distinct groups of a set $\mathfrak F$ interact only through the hard core of the compatibility: no soft interaction is left. The hypothesis of the proposition is verified in Section~\ref{sec:count} together with the bounds of the insertions.

\subsection{The Count of the Groups of the Step}\label{sec:count}

Proposition~\ref{prop:regrouping} is so far an identity between a sum and a rearrangement of it, and what it asserts depends on the convergence of that sum. We now prove the convergence, and with it the bound on the groups of the step $k + 1$ from which their place in the class will follow. What makes this a problem of counting is that a group is not the sum of its parts. Its size $\sigma_{k+1}$, in which the bounds of the class are weighted, counts the cubes and the traversals of its domain, and the domain contains arcs which belong to none of the constituents: the arcs which the regrouping adds between them. If the weights of the constituents are to pay for the weight of the group, these arcs have to be charged to the constituents which they join, and each constituent may be charged only in proportion to its own size, however many others it is tied to. We arrange this by giving every constituent a set of cubes which is large enough to contain all that the constituent can add to a group, and whose number is bounded by its own size. Tied constituents then have a cube in common, a group becomes a connected set in the graph of this relation, and the connected sets of a graph are counted by a tree bound, which is the criterion of Koteck\'y and Preiss in the form we need.

Throughout this subsection the labelled terms of $\varrho_k$ and of $\varrho_{k+1}$ are the terms without large-field regions, lengths and distances are measured in the units of the scale $k + 1$ unless the contrary is said, distances are $\ell^\infty$-distances, and a \emph{cube} is a cube of $\pi_{k+1}$, of side $M$. By Lemma~\ref{lem:stopping-index} at the step $k + 1$ two loops are at distance at least $a_{k+1} \geq L^2M$ and the segments of the loops have at least $L^2M$ bonds, so that a cube meets at most one loop, and by Lemma~\ref{lem:square-properties}(i) it meets at most $2\bar\Lambda(M + 1) \leq 4\bar\Lambda M$ traversals of the step $k + 1$.

\begin{lemma}[Coarse Cubes]\label{lem:coarse-cubes}
Let $j \leq k$, and let $X$ be a union of $n$ cubes of $\pi_j$ which is connected in the sense of Definition~\ref{def:dressed-loop}(vii). Then $\diam_\infty(X) \leq nM/L$, and $X$ meets at most $625(1 + n/L)$ cubes.
\end{lemma}

\begin{proof}
A cube of $\pi_j$ has the side $ML^{-(k+1-j)} \leq M/L$ in the units of the scale $k + 1$, and two points of $X$ are joined by a chain of at most $n$ of its cubes in which consecutive ones have a common point; this gives the diameter. For the count let $\mathcal Q$ be the set of the cubes which $X$ meets, and let $\mathcal Q' \subset \mathcal Q$ be maximal among the subsets whose members have centres at distance at least $3M$ from each other. Every member of $\mathcal Q$ has its centre within $2M$ of the centre of a member of $\mathcal Q'$, so that $\abs{\mathcal Q} \leq 5^4\abs{\mathcal Q'}$. For $Q \in \mathcal Q'$ let $Q^*$ be the union of $Q$ and of the $80$ cubes which touch it. The sets $Q^*$, $Q \in \mathcal Q'$, have no cube in common, and every cube of $\pi_j$ lies in one cube, the partitions being nested. If $\mathcal Q'$ has at least two members, then $X$ contains, besides a point of $Q$, a point of another member of $\mathcal Q'$, whose distance from $Q$ is at least $2M$. A chain of cubes of $\pi_j$ in $X$ from the first point to the second begins with a cube contained in $Q^*$ and leaves $Q^*$, and its cubes up to the first one which is not contained in $Q^*$ form a connected set which reaches from a point of $Q$ to a point of the boundary of $Q^*$, at distance $M$ from $Q$; since each of them has diameter at most $M/L$, at least $L$ of them are contained in $Q^*$. Hence $n \geq L\abs{\mathcal Q'}$ in this case, and $\abs{\mathcal Q} \leq 625\max(1, n/L)$ in both.
\end{proof}

\emph{Bodies, windows, and cube sets.} Fix a term $t$ of Section~\ref{sec:groups-of-step}. For a marked polymer $M$ of $t$ let $\mathcal K(M)$ be the set of the clusters of the links and of the vacuum occurrences which $M$ carries. The \emph{body} of a constituent, a group $W$ or a marked polymer $M$, is
\begin{equation}\label{eq:bodies}
    B(W) \coloneqq X^+_W\;,
    \qquad
    B(M) \coloneqq M \cup \bigcup_{\mathcal C \in \mathcal K(M)}\operatorname{supp}\mathcal C \cup \bigcup_{W \in \mathfrak B(M)}X^+_W\;,
\end{equation}
and $d(c) \coloneqq \diam_\infty B(c)$. If the constituent $c$ lies on a loop $C_i$, with the piece $p(c)$, its \emph{window} $\operatorname{win}(c)$ is the arc of $C_i$ which consists of $p(c)$ and of the points of $C_i$ whose distance from an end of $p(c)$, measured along $C_i$, is at most $2\bar\Lambda d(c)$. The \emph{cube set} $K(c)$ is the set of the cubes which meet $B(c)$ or, if $c$ lies on a loop, $\operatorname{win}(c)$; and $n(c)$ is the number of the traversals of the step $k + 1$ which meet $\operatorname{win}(c)$ if $c$ lies on a loop, and which meet a piece of $c$ otherwise.

\begin{lemma}[The Covering of a Group]\label{lem:group-covering}
Let $\mathcal G$ be a group of the step $k + 1$ in the term $t$.
\begin{enumerate}[label=(\roman*)]
\item The body of a constituent is connected, and it contains the domain of the constituent, the polymers of the vacuum occurrences which it carries, and the end sites of its pieces. A piece of a constituent which does not lie on a loop is contained in its body.
\item The bodies of two tied constituents $c$, $c'$ have a point in common. If both lie on the loop $C_i$ and their pieces are disjoint, the arc which Definition~\ref{def:groups-of-step} adds between these pieces has length at most $\bar\Lambda(d(c) + d(c'))$, and it lies in $\operatorname{win}(c)$ if $d(c) \geq d(c')$ and in $\operatorname{win}(c')$ otherwise.
\item The arc set $I(\mathcal G)$ is contained in the union of the windows of the constituents of $\mathcal G$ which lie on a loop and of the pieces of the others, $X_{\mathcal G}$ is contained in the union of the cubes of the sets $K(c)$, and any two constituents of $\mathcal G$ are joined by a chain of constituents of $\mathcal G$ in which the cube sets of consecutive ones have a cube in common.
\item The domain $X_{\mathcal G}$ is connected in the sense of Definition~\ref{def:dressed-loop}(vii). Its number of cubes satisfies $\#(X_{\mathcal G}) \leq \abs{\bigcup_cK(c)}$, the number of the traversals of the step $k + 1$ which meet $I(\mathcal G)$ satisfies $n_{k+1}(\mathcal G) \leq \sum_cn(c)$, and $\operatorname{ext}(\mathcal G) \leq 2\,\#(X_{\mathcal G}) + 3C_{\mathrm{box}} + 1$; the union and the sum run over the constituents of $\mathcal G$.
\item The insertion $Y_{\mathcal G}$ depends on the background field only through its restriction to $X_{\mathcal G}$, and the domains of two incompatible groups of the step $k + 1$ have a point in common.
\end{enumerate}
\end{lemma}

\begin{proof}
(i) For a group $W$ the body is a closed neighbourhood of the connected set $X_W$, which contains the intervals of $W$. Let $M$ be a marked polymer. Each of its connected components, in the sense of the walls, is a connected set, and so is, for each component $Y$ of a group of $\mathfrak G(M)$, the union $Y^\sim$ of $Y$ and of the cubes of $\pi_k$ which share a point with it; the sets $Y^\sim$ lie in $M$, because $D_G$ is their union over the components of $G$. By the definition of a linked union in Proposition~\ref{prop:marked-polymers}, any two connected components of $M$ are joined by a chain of connected components in which consecutive ones meet one set $D_G$, hence meet the sets $Y^\sim$ and $Y'^\sim$ of two components of one group $G$; and two components of a group are joined by a chain of components of that group in which consecutive ones touch a common member $W \in \mathfrak B(M)$, whose dependence set is connected and meets both. So $M \cup \bigcup_{W \in \mathfrak B(M)}X^+_W$ is connected. The support of a cluster is connected, its polymers being connected unions of cubes with a connected graph of incompatibility, and it shares a cube or a wall with $M$ if $M$ carries a link or an occurrence of that cluster. Hence $B(M)$ is connected, and it contains the domain of $M$ and the polymers of its vacuum occurrences by definition. The pieces of the groups of $\mathfrak G(M)$ consist of bonds of the loops in their components, which lie in $M$, and of intervals of members of $\mathfrak B(M)$, which lie in the domains of these members; if $M$ lies on a loop, its piece $I_i(M)$ begins and ends with such bonds.

(ii) We go through the cases of Definition~\ref{def:groups-of-step}. In (i) the member belongs to $\mathfrak B(M)$, and its body is a part of $B(M)$. In (ii) the marked polymer is incompatible with a cluster of $\mathcal K(M)$, for the carrier $M$, so that it shares a cube or a wall with the support of that cluster. In (iii) and (iv) the two domains share a cube or a wall, and they are contained in the bodies. In (v) the two groups are the same group or incompatible groups of the step $k$, and their domains have a common point by Definition~\ref{def:dressed-loop}(vii). Now let $c$ and $c'$ lie on $C_i$ with disjoint pieces $p$ and $p'$, and let $x$ and $x'$ be end sites of $p$ and of $p'$. They lie in the two bodies, which have a common point, so that $\abs{x - x'}_\infty \leq d(c) + d(c')$, and by Definition~\ref{def:chord-arc} the shorter of the two arcs of $C_i$ between $x$ and $x'$ has length at most $\Lambda(C_i)(d(c) + d(c'))$. This arc begins at an end of $p$. It may first run through $p$; it then leaves $p$ at an end and runs outside $p$ until it meets $p'$, which it does at an end of $p'$. Its part between these two ends is one of the two arcs between the pieces, and the shorter of the two, which is the arc that the definition adds, is not longer. That arc begins at an end of $p$, ends at an end of $p'$, and has length at most $2\bar\Lambda\max(d(c), d(c'))$, which gives the last statement.

(iii) By its definition $I(\mathcal G)$ is the union of the pieces of the constituents of $\mathcal G$ and of the arcs of (ii), and the piece of a constituent which lies on a loop is a part of its window; this is the first statement. A localisation domain of $\mathcal G$ consists of cubes which meet the domain of a constituent or a polymer of a vacuum occurrence, hence a body, by (i); and a cube which meets a fine bond of $I(\mathcal G)$ meets a window, or a piece which lies in a body by (i). So $X_{\mathcal G}$ consists of cubes of the sets $K(c)$. For the third statement, call two constituents of the term \emph{joined} if they are connected by a chain of constituents in which the cube sets of consecutive ones have a cube in common. This is an equivalence relation, and it has the three properties which define the groups. Tied constituents are joined, since every cube which contains a common point of their bodies belongs to both cube sets. If the arc sets of two classes share a bond, that bond lies in a window or in a piece of a constituent of each of them, by (ii), and a cube which meets it belongs to the cube sets of both constituents, so that the classes coincide. And a member of $\mathfrak B$ one of whose intervals lies in the arc set of a class has the bonds of that interval in its own body and in a window or a piece of a constituent of the class. Since the groups are the classes of the smallest equivalence relation with these properties, the constituents of a group are joined.

(iv) Let $F$ be the union of the bodies of the constituents of $\mathcal G$ and of $I(\mathcal G)$. The bodies of the constituents of $\mathcal G$ cover the same set as their domains and the polymers of their vacuum occurrences, since the members of $\mathfrak B(M)$ and the clusters of the links are constituents and parts of domains of constituents of $\mathcal G$; so $X_{\mathcal G}$ is the union of the cubes which meet $F$. The set $F$ is connected: the bodies are connected, those of tied constituents meet, every piece and every arc added between two pieces is a connected set which meets the bodies of the constituents concerned at their end sites, and the classes which the definition merges because their arc sets share a bond, or because an interval of a member lies in an arc set, have a common point there. A finite family of closed cubes each of which meets a connected set, and which together cover it, is connected in the sense of Definition~\ref{def:dressed-loop}(vii), since otherwise that set would be the union of two disjoint nonempty closed sets. The bounds on $\#(X_{\mathcal G})$ and on $n_{k+1}(\mathcal G)$ follow from (iii). For the extent put $N \coloneqq \#(X_{\mathcal G})$, and let $\mathfrak n$ be a traversal of the step $k + 1$ in an interval of $\mathcal G$. It meets a cube of $X_{\mathcal G}$, and $X_{\mathcal G}$ is a connected union of $N$ cubes of side $M$, so that $X_{\mathcal G}$ lies in the ball of radius $NM + 1$ about $\mathfrak n$. The radii $R_t = C_{\mathrm{box}}M(t + 1)$ increase in steps of $C_{\mathrm{box}}M \geq 1$, so that $R_{t_{\mathfrak n}(\mathcal G)} \leq \max(2C_{\mathrm{box}}M, NM + 1 + C_{\mathrm{box}}M)$, and the comparison region of $\mathcal G$ at $\mathfrak n$ is a ball of radius at most $(N + 3C_{\mathrm{box}})M$ about $\mathfrak n$. Another traversal of the intervals of $\mathcal G$ lies within $NM + 1$ of $\mathfrak n$, and therefore $X_{\mathcal G}$ and all the comparison regions lie in the ball of radius $(2N + 3C_{\mathrm{box}} + 1)M$ about each of these traversals.

(v) The first statement is the last statement of Proposition~\ref{prop:regrouping}, the domains of the constituents and the intervals of $\mathcal G$ being contained in $X_{\mathcal G}$. For the second we go through the conditions of the compatibility in Definition~\ref{def:groups-of-step}. If the arc sets share a bond, the cubes which meet it lie in both domains. If the domain of a marked polymer of the one is incompatible with the domain of a constituent of the other, or a vacuum occurrence of the one with a marked polymer of the other, the two sets share a cube or a wall, and both are covered by the localisation domains of their groups. If a member of the one, or a group of an old occurrence of the one, is incompatible with a member of the other as a group of the step $k$, their domains have a common point by Definition~\ref{def:dressed-loop}(vii) at the step $k$, and these domains lie in the dependence sets, which the localisation domains cover.
\end{proof}

The lemma reduces the geometry of a group to the cube sets of its constituents, and we now bound these sets. For a polymer the bound has to contain the members which it touches, since the arcs of the polymer contain their intervals.

\begin{lemma}[The Sizes of the Cube Sets]\label{lem:cube-set-sizes}
Let $L \geq 81\bar\Lambda$.
\begin{enumerate}[label=(\roman*)]
\item For a group $W$ of the step $k$ put $s(W) \coloneqq \abs{K(W)} + n(W)/M$. Then
\begin{equation}\label{eq:member-cube-set}
    d(W) \leq \frac{(u_\star\sigma(W) + 1)M}{L}\;,
    \qquad
    s(W) \leq 3000\,\bar\Lambda + 3 \cdot 10^4\,\bar\Lambda u_\star\frac{\sigma(W)}{L}\;,
\end{equation}
and for every cube of $K(W)$ there is a traversal of the step $k$ in an interval of $W$ whose distance from that cube, in the units of the scale $k$, is at most $\max(\operatorname{ext}(W)M + 1,\, 2\bar\Lambda(u_\star\sigma(W) + 1)M)$.
\item For a marked polymer $M$ of the term let $K(v) \coloneqq K(M) \cup \bigcup_{W \in \mathfrak B(M)}K(W)$, let $n(v)$ be the number of the traversals of the step $k + 1$ counted by $n(M)$ or by one of the numbers $n(W)$, $W \in \mathfrak B(M)$, and put $s(v) \coloneqq \abs{K(v)} + n(v)/M$ and
\begin{equation}\label{eq:polymer-reach}
    \hat d(v) \coloneqq \#(M) + \sum_{\mathcal C \in \mathcal K(M)}N(\mathcal C) + \sum_{W \in \mathfrak B(M)}\frac{u_\star\sigma(W) + 1}{L}\;,
\end{equation}
with $N(\mathcal C)$ the number of the cubes of the polymers of $\mathcal C$. Then $d(M) \leq M\hat d(v)$,
\begin{equation}\label{eq:polymer-cube-set}
    s(v) \leq 33 + 121\,\bar\Lambda\Bigl(\#(M) + \sum_{\mathcal C \in \mathcal K(M)}N(\mathcal C)\Bigr) + \sum_{W \in \mathfrak B(M)}\Bigl(3001\,\bar\Lambda + 3.01 \cdot 10^4\,\bar\Lambda u_\star\frac{\sigma(W)}{L}\Bigr)\;,
\end{equation}
and every cube of $K(v)$ lies within the distance $9\bar\Lambda M\hat d(v)$ of every cube of $M$.
\end{enumerate}
\end{lemma}

\begin{proof}
(i) Let $W$ be made at the step $j \leq k$, and put $N_W \coloneqq \#_j(X_W)$, so that $N_W \leq u_\star\sigma(W)$ and $n_j(W) \leq u_\star M\sigma(W)$ by Definition~\ref{def:dressed-loop}(vii). The dependence set lies in the union of $X_W$ and of the cubes of $\pi_j$ which touch it, a connected union of at most $81N_W$ cubes of $\pi_j$, since $8M$ fine bonds are less than the side of a cube of $\pi_j$. By Lemma~\ref{lem:coarse-cubes}, $B(W)$ meets at most $625(1 + 81N_W/L)$ cubes, $X_W$ meets at most $625(1 + N_W/L)$, and $d(W) \leq N_WM/L + 16ML^{-(k+1)} \leq (N_W + 1)M/L$. If $W$ is a cross group, its intervals lie in $X_W$, and each of the cubes which $X_W$ meets is met by at most $4\bar\Lambda M$ traversals, so that $n(W)/M \leq 2500\bar\Lambda(1 + N_W/L)$. If $W$ lies on a loop, its interval has the length $n_j(W)L^{-(k+1-j)} \leq u_\star M\sigma(W)/L$ and lies in $X_W$, and the window adds to it two arcs of length $2\bar\Lambda d(W)$. By Lemma~\ref{lem:cubes-along-arc}, applied on the fine lattice, an arc of length $\ell$ meets at most $8\ell/M + 16\ell/(L^2M) + 32 \leq 9\ell/M + 32$ cubes, the segments having at least $L^2M$ bonds of the scale $k + 1$. Hence the two arcs meet at most $36\bar\Lambda(N_W + 1)/L + 64$ cubes, and $n(W) \leq u_\star M\sigma(W)/L + 4\bar\Lambda d(W) + 2$. Adding up, and using $\bar\Lambda \geq 2$ and $L \geq 81\bar\Lambda$, we obtain the bound on $s(W)$ in both cases, with room to spare. For the last statement, a cube of $K(W)$ meets $B(W)$ or one of the two arcs. In the first case its distance from every traversal of the step $j$ in an interval of $W$ is at most $\operatorname{ext}(W)M$ in the units of the scale $j$, by the definition of the extent, plus the $8M$ fine bonds of the dependence set; in the second its distance from the traversal at the corresponding end of the interval is at most $2\bar\Lambda d(W)$, which is $2\bar\Lambda(N_W + 1)M$ in the units of the scale $k$. A length in the units of the scale $j$ is not larger in those of the scale $k$.

(ii) The body $B(M)$ is a connected union of the cubes of $M$, of the supports of the clusters of $\mathcal K(M)$, and of the bodies of the members of $\mathfrak B(M)$, by Lemma~\ref{lem:group-covering}(i), and the diameter of a connected finite union of closed sets is at most the sum of their diameters; with (i) this gives $d(M) \leq M\hat d(v)$. The cubes which meet $M$ or one of these supports are at most $81(\#(M) + \sum_{\mathcal C}N(\mathcal C))$, and those which meet the body of a member $W$ belong to $K(W)$. If $M$ lies on the loop $C_i$, the sites of the arc sets of its groups lie in $B(M)$, and $I_i(M)$ is the shortest arc of $C_i$ which contains them, so that it has length at most $4\Lambda(C_i)d(M)$ by Lemma~\ref{lem:arcs-around-set} and the remark which follows it. The window then has length at most $8\bar\Lambda d(M)$, it meets at most $72\bar\Lambda\hat d(v) + 32$ cubes, by the count of (i), and $n(M) \leq 8\bar\Lambda M\hat d(v) + 2$. If $M$ does not lie on a loop, its pieces consist of bonds of the loops in $M$, which are met by at most $4\bar\Lambda M\#(M)$ traversals, and of intervals of the members of $\mathfrak B(M)$, which are counted by the numbers $n(W)$. In both cases
\[
    s(v) \leq 33 + 81\Bigl(\#(M) + \sum_{\mathcal C}N(\mathcal C)\Bigr) + 80\,\bar\Lambda\hat d(v) + \sum_{W \in \mathfrak B(M)}s(W)\;,
\]
which is \eqref{eq:polymer-cube-set} by (i), $81 \leq 41\bar\Lambda$, and $L \geq 81\bar\Lambda$. Finally, a point of $B(M)$ is within $d(M)$ of every cube of $M$; a point of the window of $M$ is within $8\bar\Lambda d(M)$ of an end site of $I_i(M)$, which lies in $B(M)$; and a point of the window of a member $W \in \mathfrak B(M)$ is within $2\bar\Lambda d(W) \leq 2\bar\Lambda M\hat d(v)$ of its interval, which lies in $B(M)$. So every cube of $K(v)$ is within $(8\bar\Lambda + 1)M\hat d(v)$ of every cube of $M$.
\end{proof}

\emph{The tree bound.} The sum over the connected sets of a graph is controlled by one condition on the weights of the neighbours of a vertex. We state it with a second weight, which marks one vertex of the set; this is what will show that a group which contains an insertion of an earlier step is small in proportion to that insertion.

\begin{lemma}[Connected Sets of a Graph]\label{lem:tree-bound}
Let $\mathcal V$ be a countable set with a symmetric relation $\sim$, and let $w, b\colon \mathcal V \to [0, \infty)$ be such that, for every $v \in \mathcal V$,
\begin{equation}\label{eq:tree-condition}
    \sum_{v' \sim v}w(v')\,e^{2b(v')} \leq \tfrac12b(v)\;.
\end{equation}
For $v_0 \in \mathcal V$ let $\mathcal S(v_0)$ be the set of the finite subsets of $\mathcal V$ which contain $v_0$ and are connected in the graph of the relation.
\begin{enumerate}[label=(\roman*)]
\item $\sum_{S \in \mathcal S(v_0)}\prod_{v \in S \setminus \{ v_0 \}}w(v) \leq e^{b(v_0)/2}$.
\item Let $m\colon \mathcal V \to [0, \infty)$ and $\epsilon \geq 0$ be such that $\sum_{v' \sim v}m(v')e^{2b(v')} \leq \frac12\epsilon\,b(v)$ for every $v$. Then
\[
    \sum_{S \in \mathcal S(v_0)}\ \sum_{v_1 \in S \setminus \{ v_0 \}}m(v_1)\prod_{v \in S \setminus \{ v_0, v_1 \}}w(v) \leq \epsilon\,b(v_0)\,e^{b(v_0)/2}\;.
\]
\end{enumerate}
\end{lemma}

\begin{proof}
Put $T_0(v) \coloneqq 1$ and $F_0(v) \coloneqq 0$, and for $n \geq 0$
\begin{gather*}
    T_{n+1}(v) \coloneqq \sum_A\prod_{c \in A}w(c)T_n(c)\;,\\
    F_{n+1}(v) \coloneqq \sum_A\sum_{c \in A}\bigl[m(c)T_n(c) + w(c)F_n(c)\bigr]\prod_{c' \in A \setminus \{ c \}}w(c')T_n(c')\;,
\end{gather*}
the sums over the finite sets $A$ of vertices $c \sim v$. Then $T_{n+1}(v) \leq \exp\sum_{c \sim v}w(c)T_n(c)$, and by induction $T_n(v) \leq e^{b(v)/2}$, because $\sum_{c \sim v}w(c)e^{b(c)/2} \leq \frac12b(v)$ by \eqref{eq:tree-condition}. Likewise $F_{n+1}(v) \leq e^{b(v)/2}\sum_{c \sim v}[m(c)T_n(c) + w(c)F_n(c)]$, and by induction $F_n(v) \leq \epsilon\,b(v)e^{b(v)/2}$: the first terms sum to at most $\frac12\epsilon\,b(v)$ by the hypothesis on $m$, and the second to at most $\epsilon\sum_{c \sim v}w(c)b(c)e^{b(c)/2} \leq \frac12\epsilon\,b(v)$, by \eqref{eq:tree-condition} and $xe^{x/2} \leq e^{2x}$.

It remains to compare the sums of the lemma with these numbers. For $S \in \mathcal S(v_0)$ let the \emph{depth} of $S$ be the largest distance from $v_0$ in the graph which the relation induces on $S$. We claim that the sum in (i) over the sets of depth at most $n$ is at most $T_n(v_0)$, and the sum in (ii) at most $F_n(v_0)$; the lemma follows by letting $n$ tend to infinity. For $n = 0$ the only set is $\{ v_0 \}$. Let $S$ have depth at most $n + 1$, and fix an enumeration of $\mathcal V$. Every $u \in S \setminus \{ v_0 \}$ has a parent, the first vertex in the enumeration among its neighbours in $S$ which are closer to $v_0$ than $u$; this defines a tree on $S$ with the root $v_0$. Let $A$ be the set of the neighbours of $v_0$ in $S$, and for $c \in A$ let $S_c$ be the set of $c$ and of its descendants. Then $S_c \in \mathcal S(c)$ with depth at most $n$, the set $S$ is the disjoint union of $\{ v_0 \}$ and of the sets $S_c$, and the product of the weights over $S \setminus \{ v_0 \}$ is the product over $c \in A$ of $w(c)$ times the product over $S_c \setminus \{ c \}$. Since $S$ is determined by $A$ and the sets $S_c$, the sum over $S$ is at most the sum over $A$ of the product over $c \in A$ of $w(c)$ times the sum over $S_c$, which is the claim for (i) by induction. In (ii) the marked vertex $v_1$ belongs to exactly one of the sets $S_c$, where it is $c$ itself or a marked vertex of $S_c \setminus \{ c \}$, and the same argument gives the recursion of $F_{n+1}$.
\end{proof}

\emph{The vertices.} We describe a group of the step $k + 1$ by a set of vertices of three kinds, each with a cube set $K$, a number $s$, and an activity. A \emph{member vertex} is a group $W$ of the step $k$, with $K(W)$ and $s(W)$ of Lemma~\ref{lem:cube-set-sizes} and the activity $\norm{Y_W}_{(k)}$. An \emph{occurrence vertex} is a pair $(\Gamma, n)$ of a cluster $\Gamma$ of the reference gas of the step $k$ and an integer $n \geq 1$, with $K(\Gamma) \coloneqq \bigcup_{W' \in \Gamma}K(W')$, $s(\Gamma) \coloneqq \sum_{W' \in \Gamma}s(W')$, the sum over the distinct groups of $\Gamma$, and the activity $\abs{\varphi_k(\Gamma)}^n/n!$. A \emph{polymer vertex} is a quadruple $v = (M, \mathfrak Y_v, \mathfrak B_v, \mathcal K_v)$ of a union $M$ of cubes, a nonempty finite set $\mathfrak Y_v$ of connected unions of cubes of $\pi_k$ in $M$, a finite set $\mathfrak B_v$ of groups of the step $k$, and a finite set $\mathcal K_v$ of clusters of vacuum polymers which are incompatible with $M$, such that the dependence set of every $W \in \mathfrak B_v$ meets $M$. Its cube set, its number $s(v)$, and its reach $\hat d(v)$ are those of Lemma~\ref{lem:cube-set-sizes}(ii), formed with $\mathfrak B_v$ and $\mathcal K_v$ in the places of $\mathfrak B(M)$ and $\mathcal K(M)$. To define its activity let $\bar\varphi(\mathcal C)$ be the bound of $\abs{\varphi^T(\mathcal C)}$ on the space of the step $k + 1$ of Supplement, Lemma~17.24(iv), let $d(\mathcal C) = 9\log 2 \cdot N(\mathcal C)$, and put
\begin{equation}\label{eq:vertex-constants}
\begin{gathered}
    \xi(\mathcal C) \coloneqq \bar\xi(\mathcal C)e^{\bar\xi(\mathcal C)}\;,
    \qquad
    \bar\xi(\mathcal C) \coloneqq 2e^{d(\mathcal C)}\bar\varphi(\mathcal C)\;,\\
    \kappa' \coloneqq \kappa_1 - 2 - c_\kappa\;,
    \qquad
    r \coloneqq \frac{(L + 2)\kappa}{1296}\;,
    \qquad
    \varpi_k \coloneqq \frac{b_k}{\min(\varepsilon_H, r_\alpha)}\;.
\end{gathered}
\end{equation}
The activity of the polymer vertex $v$ is
\begin{equation}\label{eq:polymer-activity}
    \operatorname{act}(v) \coloneqq \mathfrak a_v\,e^{c\#(M)}\,e^{-\kappa'\#(\mathfrak Y_v)}\,e^{-r\#(M)}\prod_{W \in \mathfrak B_v}\norm{Y_W}_{(k)}\prod_{\mathcal C \in \mathcal K_v}\xi(\mathcal C)\;,
\end{equation}
where $c = \frac{9}{64}e^{-4\kappa}$, $\#(\mathfrak Y_v)$ is the total number of the cubes of the members of $\mathfrak Y_v$, and $\mathfrak a_v \coloneqq c_a\rho_k$ if $\mathfrak B_v = \emptyset$ and $\mathfrak a_v \coloneqq \varpi_k$ otherwise, with the constant $c_a$ of Proposition~\ref{prop:marked-size}. For every vertex $v$ we put
\begin{equation}\label{eq:vertex-weights}
    w(v) \coloneqq \operatorname{act}(v)\,e^{4\nu s(v)}\;,
    \qquad
    b(v) \coloneqq \nu\abs{K(v)}\;,
\end{equation}
and we write $v \sim v'$ if $K(v)$ and $K(v')$ have a cube in common. The \emph{vertex set} $S(\mathcal G)$ of a group $\mathcal G$ of the step $k + 1$ consists of the polymer vertices $(M, \mathfrak Y_M, \mathfrak B(M), \mathcal K(M))$ of its marked polymers, where $\mathfrak Y_M$ is the set of the components of the groups of $\mathfrak G(M)$; of its members which no component touches; and of the pairs $(\Gamma, n)$ of the distinct clusters $\Gamma$ of its old occurrences with their multiplicities $n$ in the sequence of these occurrences. Let finally
\begin{equation}\label{eq:group-product}
    P(\mathcal G) \coloneqq e^{c'_\pi\alpha_{1,k+1}n_{k+1}(\mathcal G)}\prod_M\sup\norm{e^{-\Psi_\circ(M)}\tilde H(M)}\prod_W\norm{Y_W}_{(k)}\prod_\lambda\sup\abs{u_\lambda}\prod_{\mathcal C}\abs{v_{\mathcal C}}\cdot\frac{1}{m!}\prod_{j=1}^{m}\abs{\varphi_k(\Gamma_j)}
\end{equation}
be the product of the norms of the factors of $Y_{\mathcal G}$ in Definition~\ref{def:groups-of-step}, the suprema over the space of the step $k + 1$ and the products over the marked polymers, the untouched members, the links, the vacuum occurrences, and the old occurrences of $\mathcal G$. The first factor bounds the transports: on the space of the step $k + 1$ the transport along one traversal of the step $k + 1$ has norm at most $e^{c'_\pi\alpha_{1,k+1}}$, by the radii of Definition~\ref{def:dressed-loop}(iii), and the insertion of an untouched member is evaluated at a configuration of that space, which belongs to the space of the step $k$ (Section~\ref{sec:localisation-expansion}). Since a contraction of tensors does not increase the product of the norms (Section~\ref{sec:dressed-loops}),
\begin{equation}\label{eq:group-sup}
    \sup\norm{Y_{\mathcal G}} \leq P(\mathcal G)\;,
\end{equation}
the supremum over the space of the step $k + 1$.

\begin{lemma}[The Weight of a Group]\label{lem:group-weight}
Suppose that the hypotheses of Proposition~\ref{prop:marked-size} hold, that $\kappa \geq 2^9(2 + 16\log 2)$, and that $4\bar\Lambda Mc'_\pi\alpha_{1,k+1} \leq \nu \leq 1$. Then the vertex set of every group of the step $k + 1$ is connected in the graph of the relation $\sim$, and for every finite set $S$ of vertices
\begin{equation}\label{eq:group-weight}
    \sum_{\mathcal G\colon S(\mathcal G) = S}e^{\frac52\nu\sigma_{k+1}(\mathcal G) + \nu_e\operatorname{ext}(\mathcal G)}P(\mathcal G)\,\ell_{\mathcal G} \leq \frac{e^{2\nu}}{\nu}\prod_{v \in S}w(v)\;,
\end{equation}
the sum over the groups of the step $k + 1$, in all the terms $t$, with the vertex set $S$.
\end{lemma}

\begin{proof}
The cube set of a polymer vertex of $S(\mathcal G)$ contains those of the marked polymer and of the members which it touches, and the cube set of an occurrence vertex those of the groups of the occurrence, so that the vertex set is connected by Lemma~\ref{lem:group-covering}(iii).

A group is determined by its vertex set up to the links and the vacuum occurrences which its marked polymers carry: the vertex set gives the marked polymers with their components and touched members, hence their groups, since two components are linked if they touch a common member; the untouched members; the old occurrences with their multiplicities, to which there correspond $m!/\prod n!$ sequences; and, for each marked polymer $M$, the set $\mathcal K(M)$ of the clusters of its links and vacuum occurrences. We bound the factors of \eqref{eq:group-product} in turn.

\emph{The clusters.} Fix $M$ and $\mathcal C \in \mathcal K(M)$, and let $s$ be the number of the marked polymers of the group which are incompatible with $\mathcal C$. They are pairwise compatible, and each of them contains a cube which belongs to a polymer of $\mathcal C$ or shares a wall with such a cube, so that $s \leq 9N(\mathcal C)$. The links of $\mathcal C$ which $M$ carries are pairs $(T, \mathcal C)$ with $M$ the first member of $T$, at most $2^{s-1}$ in number, and at most one vacuum occurrence is $\mathcal C$. By \eqref{eq:link-factor} and \eqref{eq:vacuum-occurrences}, $\abs{u_\lambda} \leq 2\bar\varphi e^{2\bar\varphi}$ and $\abs{v_{\mathcal C}} \leq (s - 1)\bar\varphi e^{(s-1)\bar\varphi}$ with $\bar\varphi = \bar\varphi(\mathcal C)$. By Supplement, Lemma~17.24(iv), applied to a polymer of $\mathcal C$, $\bar\varphi e^{d(\mathcal C) + N(\mathcal C)} \leq cN(\mathcal C)$, whence $9N(\mathcal C)\bar\varphi \leq c$ and $\bar\xi(\mathcal C) \leq 1$. The sum, over the nonempty sets of links and occurrences of $\mathcal C$ carried by $M$, of the products of these numbers is therefore at most
\[
    \bigl(1 + 2\bar\varphi e^{2\bar\varphi}\bigr)^{2^{s-1}}\bigl(1 + (s - 1)\bar\varphi e^{(s-1)\bar\varphi}\bigr) - 1 \leq xe^x\;,
    \qquad
    x \coloneqq 2^s\bar\varphi e^c + 2^{s-1}\bar\varphi e^c \leq \bar\xi(\mathcal C)\;,
\]
because $s - 1 \leq 2^{s-1}$, $\frac32e^c \leq 2$, and $2^s \leq e^{d(\mathcal C)}$. This is the factor $\xi(\mathcal C)$ of \eqref{eq:polymer-activity}.

\emph{The other factors.} By Supplement, Lemma~17.24(ii), at the reference data, $\abs{\Psi_\circ(M)} \leq c\,\#(M)$. By Proposition~\ref{prop:marked-size}, and since $\mathfrak a(M) \leq \mathfrak a_v$ in each of its three cases, $\sup\norm{\tilde H(M)}$ is at most the product of the factors of \eqref{eq:polymer-activity} other than $e^{c\#(M)}$ and those of the clusters. The factors of the untouched members are the activities of the member vertices, and the sum of $(1/m!)\prod_j\abs{\varphi_k(\Gamma_j)}$ over the sequences with given multiplicities is the product of the activities of the occurrence vertices. So the sum of $P(\mathcal G)e^{-c'_\pi\alpha_{1,k+1}n_{k+1}(\mathcal G)}$ over the groups with the vertex set $S$ is at most $\prod_{v \in S}\operatorname{act}(v)$.

\emph{The exponent.} The arc set of a group depends on its links, and not on its vertex set alone, but the following bounds hold for every group $\mathcal G$ with the vertex set $S$. By Lemma~\ref{lem:group-covering}(iv), with $N \coloneqq \#(X_{\mathcal G}) \leq \sum_{v \in S}\abs{K(v)}$ and $n_{k+1}(\mathcal G) \leq \sum_{v \in S}n(v)$, we have $\sigma_{k+1}(\mathcal G) = N + n_{k+1}(\mathcal G)/M \leq \sum_{v \in S}s(v)$, $\nu_e\operatorname{ext}(\mathcal G) \leq 2\nu_eN + 2\nu$ because $\nu_e = \nu/(2C_{\mathrm{box}})$, and $\ell_{\mathcal G} = N \leq \nu^{-1}e^{\nu N}$. The hypothesis on $\alpha_{1,k+1}$ gives $c'_\pi\alpha_{1,k+1}n_{k+1}(\mathcal G) \leq \nu n_{k+1}(\mathcal G)/M$. Altogether the exponent is at most $2\nu + (\frac52 + \frac32)\nu\,\sigma_{k+1}(\mathcal G) \leq 2\nu + 4\nu\sum_{v \in S}s(v)$, which with \eqref{eq:vertex-weights} gives \eqref{eq:group-weight}.
\end{proof}

\emph{The rooted sums.} By Lemma~\ref{lem:tree-bound} the sum over the groups is now reduced to the sums of the weights of the vertices which contain a given cube, and these are what the class at the step $k$ and the bounds of the polymers control. The class is stated per traversal, and the first thing to do is to pass from traversals to cubes; this is where the weight of the extent in Definition~\ref{def:dressed-loop}(vii) is spent. Let
\begin{equation}\label{eq:old-sum}
    \mathsf E_k \coloneqq \sup_{\mathfrak n}\sum_{W \ni \mathfrak n}e^{\theta(W) + \lambda_c\nu\sigma(W) + \nu_e\operatorname{ext}(W)}\norm{Y_W}_{(k)}
\end{equation}
be the supremum over the traversals of the step $k$ of the sum over all the groups of the term of $\varrho_k$ with $\mathfrak n$ in an interval, of whatever species; by \eqref{eq:sup-clause} the groups of the deviation species contribute at most $C_c\rho_k$ to it. We use the constants
\begin{equation}\label{eq:count-constants}
\begin{gathered}
    c_r \coloneqq \frac{\nu}{11664\,\bar\Lambda u_\star}\;,
    \qquad
    K_{\mathrm{box}} \coloneqq 6\bar\Lambda LM\sum_{t \geq 0}(t + 3)^5e^{-c_r(t - 2\bar\Lambda)_+}\;,\\
    c_p \coloneqq \min\Bigl(\frac{1}{28800\,\bar\Lambda},\, \frac{\nu L}{1152\,\bar\Lambda u_\star}\Bigr)\;,
    \qquad
    C_{\mathrm{root}} \coloneqq \sum_{t \geq 0}8(2t + 1)^3e^{-c_p(t - 1)_+}\;,
\end{gathered}
\end{equation}
which depend on $\bar\Lambda$, $u_\star$, $\nu$, $L$, and $M$ only, and we say that a vertex $v$ is \emph{at} a cube $\Delta$ if $\Delta \in K(v)$.

\begin{proposition}[The Rooted Sums]\label{prop:rooted-sums}
Suppose that the hypotheses of Lemma~\ref{lem:group-weight} hold, and moreover
\begin{equation}\label{eq:count-conditions}
    \nu \leq \frac{1}{1.2 \cdot 10^6\,\bar\Lambda}\;,
    \qquad
    L \geq 6 \cdot 10^6\,\bar\Lambda u_\star\;,
    \qquad
    r \geq 10\;,
    \qquad
    L^4e^{-\kappa'} \leq \frac{\nu}{1.6 \cdot 10^5\,u_\star}\;.
\end{equation}
Let $\Delta$ be a cube.
\begin{enumerate}[label=(\roman*)]
\item $\sum_{W \text{ at } \Delta}\norm{Y_W}_{(k)}e^{(\nu/16)\sigma(W)} \leq K_{\mathrm{box}}\mathsf E_k$, the sum over the groups of the term of $\varrho_k$, and $\sum_{W \text{ at } \Delta}S_We^{(\nu/16)\sigma(W)} \leq 2K_{\mathrm{box}}C_{S'}\rho_k$, the sum over the groups made at the step $k$.
\item The member vertices satisfy $\sum_{W \text{ at } \Delta}w(W)e^{2b(W)} \leq e^{1/60}K_{\mathrm{box}}\mathsf E_k$.
\item If $10368\,u_\star e^{1/60}C_{S'}\rho_k \leq \nu$, the occurrence vertices satisfy $\sum_{(\Gamma, n) \text{ at } \Delta}w(\Gamma, n)e^{2b(\Gamma)} \leq 6K_{\mathrm{box}}C_{S'}\rho_k$.
\item The polymer vertices with $\mathfrak B_v = \emptyset$ and $M \in \mathbf D_{k+1}$ satisfy $\sum_{v \text{ at } \Delta}w(v)e^{2b(v)} \leq C_\alpha\rho_k$ with $C_\alpha \coloneqq 2C_{\mathrm{root}}c_aL^4e^{-\kappa'}$.
\item If $e^{1/50}K_{\mathrm{box}}\mathsf E_k \leq \frac12$, the polymer vertices of the vertex sets of the groups of the step $k + 1$ which have $\mathfrak B_v \neq \emptyset$ satisfy $\sum_{v \text{ at } \Delta}w(v)e^{2b(v)} \leq C_\beta\varpi_k\mathsf E_k$ with $C_\beta \coloneqq 3C_{\mathrm{root}}K_{\mathrm{box}}L^4e^{-\kappa'}$.
\end{enumerate}
\end{proposition}

\begin{proof}
(i) Let $W$ be at $\Delta$. By Lemma~\ref{lem:cube-set-sizes}(i) there is a traversal $\mathfrak n$ of the step $k$ in an interval of $W$ whose distance $D$ from $\Delta$, in the units of the scale $k$, satisfies $D \leq \operatorname{ext}(W)M + 1$ or $D \leq 2\bar\Lambda(u_\star\sigma(W) + 1)M$. In the first case $\frac12\nu_e\operatorname{ext}(W) \geq \frac12\nu_e(D - 1)/M$, in the second $\frac{1}{64}\nu\sigma(W) \geq (\nu/(128\bar\Lambda u_\star))(D/M - 2\bar\Lambda)$, and in both
\[
    e^{\frac{1}{16}\nu\sigma(W)} \leq e^{-c_r(D/M - 2\bar\Lambda)_+}\,e^{\lambda_c\nu\sigma(W) + \nu_e\operatorname{ext}(W)}\;,
\]
since $\frac{1}{16} + \frac{1}{64} \leq \lambda_c$ and $\nu_e = \nu/5832$. The traversals of the step $k$ at a distance less than $(t + 1)M$ from $\Delta$ meet a cube of side at most $LM + 2(t + 1)M + 2$ in the units of the scale $k$. By Lemma~\ref{lem:stopping-index} the loops which meet that cube are at distance at least $L^2M$ from each other, so that they are at most $(t + 3)^4$ in number, and by Lemma~\ref{lem:square-properties}(i) each of them has at most $2\bar\Lambda(LM + 2(t + 1)M + 3) \leq 6\bar\Lambda LM(t + 3)$ traversals which meet it. Summing over $t \geq 0$ the traversals with $tM \leq D < (t + 1)M$, and for each of them the groups $W$ with that traversal in an interval, we obtain the first bound from \eqref{eq:old-sum}, as $\theta(W) \geq 0$. The second is the same computation with \eqref{eq:production-clause} at the traversals of the step $k$, where $\ell_W \geq 1$ and the weight $\nu\sigma(W)$ takes the place of $\lambda_c\nu\sigma(W)$.

(ii) By \eqref{eq:member-cube-set} and the first two conditions of \eqref{eq:count-conditions},
\begin{equation}\label{eq:member-share}
    4\nu s(W) + 2b(W) \leq 6\nu s(W) \leq \frac{1}{60} + \frac{\nu}{32}\sigma(W)\;,
\end{equation}
because $18000\bar\Lambda\nu \leq \frac{1}{60}$ and $1.8 \cdot 10^5\bar\Lambda u_\star/L \leq \frac{1}{32}$. So (ii) follows from (i).

(iii) The reference gas of the step $k$ is the hard-core gas of the groups made at the step $k$ with the activities $c_W$, $\abs{c_W} \leq S_W$, and the compatibility of the groups. Put $a(W) \coloneqq (\nu/64)\sigma(W)$ and $d(W) \coloneqq \frac{1}{60} + (\nu/32)\sigma(W)$. A group which is incompatible with $W$ has a point of its domain in common with $X_W$, by Definition~\ref{def:dressed-loop}(vii), so that it is at one of the at most $81u_\star\sigma(W)$ cubes of $\pi_k$ which belong to $X_W$ or touch it. By \eqref{eq:production-clause} at the cubes of $\pi_k$,
\[
    \sum_{W' \not\sim W}\abs{c_{W'}}e^{a(W') + d(W')} \leq 81u_\star\sigma(W)\,e^{1/60}\cdot 2C_{S'}\rho_k \leq a(W)\;,
\]
by the hypothesis of (iii). This is the condition of \cite[Theorem~5.4]{FriedliVelenik} with the functions $a$ and $d$, in the form in which Lemma~\ref{lem:root-bound} uses that theorem: the sum of $\abs{\varphi_k(\Gamma)}\prod_{W' \in \Gamma}e^{d(W')}$ over the clusters which contain a group incompatible with an external polymer is at most the sum of $\abs{c_{W'}}e^{a(W') + d(W')}$ over the groups incompatible with it. With a cube of $\pi_k$ in the domain of a group of $\Gamma$ as the external polymer this gives $\abs{\varphi_k(\Gamma)} \leq 2e^{1/60}C_{S'}\rho_k \leq 1$, so that $\sum_{n \geq 1}\abs{\varphi_k(\Gamma)}^n/n! \leq e\abs{\varphi_k(\Gamma)}$. With the cube $\Delta$ as the external polymer, incompatible with the groups which are at $\Delta$, and with \eqref{eq:member-share} for each group of an occurrence, it gives
\[
    \sum_{(\Gamma, n) \text{ at } \Delta}w(\Gamma, n)e^{2b(\Gamma)} \leq e\sum_{W' \text{ at } \Delta}\abs{c_{W'}}e^{\frac{1}{60} + \frac{3}{64}\nu\sigma(W')} \leq 2e^{1 + 1/60}K_{\mathrm{box}}C_{S'}\rho_k\;,
\]
by (i).

(iv) and (v). Let $v$ be a polymer vertex. By \eqref{eq:polymer-cube-set} and the first two conditions of \eqref{eq:count-conditions}, as in (ii),
\[
    4\nu s(v) + 2b(v) \leq 198\nu + \frac{1}{1600}\Bigl(\#(M) + \sum_{\mathcal C \in \mathcal K_v}N(\mathcal C)\Bigr) + \sum_{W \in \mathfrak B_v}\Bigl(\frac{1}{60} + \frac{\nu}{32}\sigma(W)\Bigr)\;.
\]
We first reduce the sum over the vertices at $\Delta$ to a sum over the vertices whose polymer contains a given cube. By Lemma~\ref{lem:cube-set-sizes}(ii) a vertex at $\Delta$ has $9\bar\Lambda\hat d(v) \geq t - 1$ for every cube $\Delta_0$ of $M$, where $tM$ is the distance of the centres of $\Delta_0$ and $\Delta$; then $\#(M) + \sum_{\mathcal C}N(\mathcal C) \geq (t - 1)/(18\bar\Lambda)$ or $\sum_{W \in \mathfrak B_v}(u_\star\sigma(W) + 1)/L \geq (t - 1)/(18\bar\Lambda)$, and in both cases the number
\[
    g(v) \coloneqq \frac{1}{1600}\Bigl(\#(M) + \sum_{\mathcal C \in \mathcal K_v}N(\mathcal C)\Bigr) + \frac{\nu}{64}\sum_{W \in \mathfrak B_v}(\sigma(W) + 1)
\]
is at least $c_p(t - 1)_+$. There are at most $8(2t + 1)^3$ cubes $\Delta_0$ for each $t$, and therefore
\begin{equation}\label{eq:root-reduction}
    \sum_{v \text{ at } \Delta}w(v)e^{2b(v)} \leq C_{\mathrm{root}}\sup_{\Delta_0}\sum_{v\colon \Delta_0 \subset M}w(v)\,e^{2b(v) + g(v)}\;.
\end{equation}
In the sum on the right we sum first over $\mathcal K_v$ and $\mathfrak Y_v$. By the two displays and \eqref{eq:polymer-activity},
\[
    w(v)e^{2b(v) + g(v)} \leq e^{198\nu}\,\mathfrak a_v\,e^{-(r - c - \frac{1}{800})\#(M)}\,e^{-\kappa'\#(\mathfrak Y_v)}\prod_{\mathcal C \in \mathcal K_v}\xi(\mathcal C)e^{N(\mathcal C)/800}\prod_{W \in \mathfrak B_v}\omega(W)\;,
\]
where $\omega(W) \coloneqq \norm{Y_W}_{(k)}\exp(\frac{1}{60} + \frac{1}{64}\nu + \frac{3}{64}\nu\sigma(W))$.
The sum over the sets $\mathcal K_v$ of clusters incompatible with $M$ of the product of the numbers $\xi(\mathcal C)e^{N(\mathcal C)/800} \leq 2e\,\bar\varphi(\mathcal C)e^{d(\mathcal C) + N(\mathcal C)}$ is at most $\exp(2e\,c\,\#(M))$, by Supplement, Lemma~17.24(iv). The sets $\mathfrak Y_v$ are determined by the sets of their cubes, nonempty subsets of the $L^4\#(M)$ cubes of $\pi_k$ in $M$, and
\[
    \sum_{\mathfrak Y_v}e^{-\kappa'\#(\mathfrak Y_v)} \leq \bigl(1 + e^{-\kappa'}\bigr)^{L^4\#(M)} - 1 \leq L^4e^{-\kappa'}e^{2\#(M)}\;,
\]
by the last condition of \eqref{eq:count-conditions}; if $M$ has several connected components and each of them is to contain a member of $\mathfrak Y_v$, the same bound holds with one factor $L^4e^{-\kappa'}$ for each component. Since $7c + \frac{1}{800} \leq 1$, each connected component $P$ of $M$ is thus left with the factor
\[
    \omega(P) \coloneqq L^4e^{-\kappa'}e^{-(r - 3)\#(P)}\;,
    \qquad
    \sum_{P \ni \Delta_0}\omega(P)e^{2\#(P)} \leq L^4e^{-\kappa'}\sum_{n \geq 1}64^{n-1}e^{-(r - 5)n} \leq L^4e^{-\kappa'}\;,
\]
the sum over the connected unions $P \in \mathbf D_{k+1}$ of cubes which contain $\Delta_0$, of which there are at most $64^{n-1}$ with $n$ cubes (Supplement, proof of Lemma~17.24); the last step uses $r \geq 10$.

For (iv), $M$ is connected and $\mathfrak a_v = c_a\rho_k$, and the three bounds give the claim with $e^{198\nu} \leq 2$.

For (v), let $v$ belong to the vertex set of a group, so that $M$ is a marked polymer with $\mathfrak B(M) \neq \emptyset$. Consider the graph whose vertices are the connected components $P$ of $M$ and the members $W \in \mathfrak B(M)$, a component and a member being adjacent if a cube of the component meets the dependence set of the member. It is connected, by the proof of Lemma~\ref{lem:group-covering}(i), since a component of a group which touches a member lies in a connected component of $M$. So the sum over these vertices $v$ with $\Delta_0 \subset M$ is at most $e^{198\nu}\varpi_k$ times the sum, over the connected sets of this kind which contain the component $P_0 \ni \Delta_0$ and at least one member, of the products of the numbers $\omega$. We bound it by Lemma~\ref{lem:tree-bound}(ii), applied to the graph of all the connected unions $P$ of cubes and all the groups $W$ of the term, with the weights $\omega$, the charges $b'(P) \coloneqq \#(P)$ and $b'(W) \coloneqq (\nu/128)(\sigma(W) + 1)$, and the marks $m(W) \coloneqq \omega(W)$, $m(P) \coloneqq 0$. The condition \eqref{eq:tree-condition} holds at a member $W$, because its dependence set meets at most $625u_\star(\sigma(W) + 1)$ cubes, by Lemma~\ref{lem:coarse-cubes} and $L \geq 81$, and
\[
    \sum_{P \sim W}\omega(P)e^{2\#(P)} \leq 625\,u_\star(\sigma(W) + 1)L^4e^{-\kappa'} \leq \tfrac12b'(W)
\]
by the last condition of \eqref{eq:count-conditions}; and at a component $P$, with $\epsilon \coloneqq 2e^{1/50}K_{\mathrm{box}}\mathsf E_k \leq 1$, because by (i)
\[
    \sum_{W \sim P}\omega(W)e^{2b'(W)} \leq \#(P)\sup_{\Delta'}\sum_{W \text{ at } \Delta'}\norm{Y_W}_{(k)}e^{\frac{1}{50} + \frac{1}{16}\nu\sigma(W)} \leq \tfrac12\epsilon\,b'(P)\;,
\]
where we used $\omega(W)e^{2b'(W)} = \norm{Y_W}_{(k)}e^{\frac{1}{60} + \frac{1}{32}\nu + \frac{1}{16}\nu\sigma(W)}$ and $\nu/32 \leq \frac{1}{50} - \frac{1}{60}$, and that a member adjacent to $P$ is at a cube of $P$. The same display is the hypothesis on the marks. Lemma~\ref{lem:tree-bound}(ii) therefore bounds the sum over the connected sets by $\epsilon\,\#(P_0)e^{\#(P_0)/2} \leq \epsilon\,e^{2\#(P_0)}$, and the sum over $P_0 \ni \Delta_0$ gives $\epsilon L^4e^{-\kappa'}$. With \eqref{eq:root-reduction} and $2e^{198\nu + 1/50} \leq 3$ this is (v).
\end{proof}

\begin{theorem}[The Size Bound]\label{thm:size-bound}
Suppose that the hypotheses of Proposition~\ref{prop:rooted-sums} hold, with those of its parts (iii) and (v), and put
\begin{equation}\label{eq:rooted-totals}
    E_\alpha \coloneqq C_\alpha\rho_k\;,
    \qquad
    E_\beta \coloneqq C_\beta\varpi_k\mathsf E_k\;,
    \qquad
    E_o \coloneqq e^{1/60}K_{\mathrm{box}}\mathsf E_k + 6K_{\mathrm{box}}C_{S'}\rho_k\;.
\end{equation}
If $E_\alpha + E_\beta + E_o \leq \frac12\nu$, then the hypothesis of Proposition~\ref{prop:regrouping} holds at every real background of the support of the term, and for every cube $\Delta$
\begin{equation}\label{eq:size-bound-new}
    \sum_{\mathcal G\colon \Delta \subset X_{\mathcal G}}e^{\frac52\nu\sigma_{k+1}(\mathcal G) + \nu_e\operatorname{ext}(\mathcal G)}P(\mathcal G)\,\ell_{\mathcal G} \leq \frac{e^{2\nu}}{\nu}\Bigl[E_\alpha + E_\beta + \frac{2}{\nu}\bigl(E_\alpha + E_\beta + E_o\bigr)E_o + \frac{2}{\nu}\bigl(E_\beta + E_o\bigr)E_\alpha\Bigr]\;,
\end{equation}
the sum over the groups of the step $k + 1$, in all the terms, other than those which consist of a single member of a family and nothing else.
\end{theorem}

\begin{proof}
Let $E \coloneqq E_\alpha + E_\beta + E_o$. The vertex set of a group consists of member vertices, of occurrence vertices, and of polymer vertices which have $\mathfrak B_v = \emptyset$ and a connected polymer, as was shown in the proof of Proposition~\ref{prop:marked-size} for a marked polymer without members, or which have $\mathfrak B_v \neq \emptyset$. Let $\mathcal V$ be the set of these vertices, with the weights and the relation of \eqref{eq:vertex-weights}. A vertex $v' \sim v$ is at one of the cubes of $K(v)$, so that by Proposition~\ref{prop:rooted-sums}
\[
    \sum_{v' \sim v}w(v')e^{2b(v')} \leq \abs{K(v)}\,E \leq \tfrac12b(v)\;,
\]
which is \eqref{eq:tree-condition}. Now let $\mathcal G$ be a group with $\Delta \subset X_{\mathcal G}$. By Lemma~\ref{lem:group-covering}(iii) one of its vertices is at $\Delta$, and we choose a root $v_0$ among these: a polymer vertex with $\mathfrak B_v \neq \emptyset$ if there is one, otherwise a member or an occurrence vertex if there is one, and otherwise a polymer vertex without members. By Lemma~\ref{lem:group-weight} the left side of \eqref{eq:size-bound-new} is at most $e^{2\nu}/\nu$ times the sum over the roots $v_0$ at $\Delta$ and over the connected sets $S \ni v_0$ which occur, of the product of the weights, and we bound this sum for the three kinds of roots. For the roots of the first kind Lemma~\ref{lem:tree-bound}(i) gives at most $\sum_{v_0}w(v_0)e^{b(v_0)/2} \leq E_\beta$. For a root of the second kind the set $S$ has a further element, since a group which consists of a single member is excluded and an old occurrence is tied to a member; we apply Lemma~\ref{lem:tree-bound}(ii) with the marks $m \coloneqq w$ and $\epsilon \coloneqq 2E/\nu$ and obtain at most $\sum_{v_0}w(v_0)\,\epsilon\,b(v_0)e^{b(v_0)/2} \leq (2E/\nu)E_o$. For a root of the third kind, either all the elements of $S$ are polymer vertices without members, and Lemma~\ref{lem:tree-bound}(i) gives at most $E_\alpha$, or $S$ contains a vertex of another kind, and Lemma~\ref{lem:tree-bound}(ii), with the marks $m \coloneqq w$ on the vertices of the other kinds and $m \coloneqq 0$ on the polymer vertices without members, for which $\epsilon = 2(E_\beta + E_o)/\nu$, gives at most $(2(E_\beta + E_o)/\nu)E_\alpha$. In each case we used $b\,e^{b/2} \leq e^{2b}$ and Proposition~\ref{prop:rooted-sums} at $\Delta$.

For the convergence, the terms $t$ correspond to the sets of pairwise compatible groups (Proposition~\ref{prop:regrouping}), and the product of the norms of the tensors and of the numerical factors of a term is the product of the numbers $P(\mathcal G)$ of its groups, without the factors of the transports. The sum over the terms is therefore at most the exponential of the sum of $P(\mathcal G)$ over all the groups. The groups which consist of a single member contribute the sum of $\norm{Y_W}_{(k)}$ over the groups of the term of $\varrho_k$, which is finite by \eqref{eq:old-sum}, the torus having finitely many traversals, and the others contribute at most the right side of \eqref{eq:size-bound-new} times the number of the cubes of the torus.
\end{proof}

The four terms of \eqref{eq:size-bound-new} have different origins, and the difference matters for the induction. The first is the contribution of the groups which the step makes out of nothing, and it is of the order of $\rho_k$ with a constant which does not involve the constants of the class. The others are proportional to the sums of the step $k$, with coefficients which tend to zero with the coupling: $\varpi_k$ is a negative power of $\log g_k^{-2}$, and $E_\alpha$, $E_\beta$, $E_o$ are small. This is the form which a recursion with a fixed point requires.

\subsection{The Small-Field Map and the Flat Parts}\label{sec:small-field-map}

\emph{The groups and the interactions.} By Proposition~\ref{prop:regrouping} the expectation $\langle F\rangle$ is, on the terms without large-field regions, the sum of the words of the families of pairwise compatible groups of Definition~\ref{def:groups-of-step}. A group whose intervals lie on one loop has one interval: each of its constituents then has a single piece, and every identification of Definition~\ref{def:groups-of-step} joins two sets of constituents whose arc sets meet or adds the arc between two pieces, so that the arc set stays an arc. A group with intervals on several loops is a cross group. The polymers which belong to no group form Ba{\l}aban's expansion of the effective action. The families of compatible groups are resummed below against the gas of the groups which they exclude, which gives the interactions $Q_{k+1}(\mathfrak B)$. The identity $N_{k+1} = \sum\chi_{k+1}(\Omega_{k+1})\mathbf T_{k+1}[\mathcal W_{k+1}\exp A_{k+1}]$ is Proposition~\ref{prop:step-identity}, with $\mathcal W_{k+1} = \langle F\rangle$ at the steps at which no $\mathbf R$-operation acts, and \eqref{eq:numerator-integral} holds at $k + 1$ because the step preserves integrals.

\emph{The flat parts.} Every group of the step has a \emph{flat value}, its insertion at the reference data of the term, the trivial background with the labels of the term. By gauge covariance under constant gauge transformations, which fix the trivial background, the flat value of a group on one loop commutes with $\pi_i(G)$, and by Schur's lemma it is a multiple $c_W\,\mathrm{id}$ of the identity. If the group is not attached to a large-field region and contains no part of an attached insertion, the insertion is split as $Y_W = c_W\,\pi_i(U_{k+1}(I_W)) + Y^{\mathrm d}_W$, where $I_W$ is the interval of the group and $\pi_i(U_{k+1}(I_W))$ the transport of the background along it. A word with the first part at $I_W$ is $c_W$ times the word without the group, since that part restores the transport which the group had replaced; the scalars $c_W$ form the reference gas of the step, and the second part $Y^{\mathrm d}_W$ vanishes at flat data and is a deviation. An attached group is not split: its insertion carries no deviation bound, with or without the subtraction, and it is kept whole (Section~\ref{sec:reference-gas}). Section~\ref{sec:flat-part} resums the reference gas along each loop and extracts from it, exactly, the step factor $\kappa_{k+1}$ and the remainder $E^{\circ}_{k+1}$, which joins $\mathcal E_{k+1}$. A cross group carries no flat value. Its insertion at the reference data is invariant under the simultaneous conjugation of its factors by $\pi_{i(I)}(u)$, $u \in G$, and not under the conjugation of a single factor, so that it is not a scalar in general and need not vanish; nothing is subtracted from it, and it is kept whole as an insertion of the cross species (Section~\ref{sec:assembly}).

\emph{The production bound.} The class records, for the groups made at a step, a bound on their sizes on the whole space of that step, \eqref{eq:production-clause}. For the step $k + 1$ it is a consequence of Theorem~\ref{thm:size-bound}, in which the old sums enter only with coefficients that vanish with the coupling; this is what allows the constant to be the same at every step.

\begin{corollary}[The Production Bound]\label{cor:production-bound}
In the setting of Theorem~\ref{thm:size-bound}, let $\mathcal G$ be a group made at the step $k + 1$, that is, a group of the step $k + 1$ which does not consist of a single member. Let $c_{\mathcal G}$ be its flat value if its insertion is split and $0$ otherwise, $Y^{\mathrm d}_{\mathcal G}$ its insertion after the subtraction, and $S_{\mathcal G} \coloneqq \norm{Y^{\mathrm d}_{\mathcal G}}_{(k+1)} + \abs{c_{\mathcal G}}$. Then $S_{\mathcal G} \leq 3e^{\nu\sigma_{k+1}(\mathcal G)}P(\mathcal G)$, and for all numbers $0 \leq \theta(\mathcal G) \leq \nu\sigma_{k+1}(\mathcal G)/16$, every cube $p$ of $\pi_{k+1}$, and every traversal $p$ of the step $k + 1$,
\begin{equation}\label{eq:size-bound}
    \sum_{\mathcal G \text{ at } p}e^{\theta(\mathcal G) + \nu\sigma_{k+1}(\mathcal G) + \nu_e\operatorname{ext}(\mathcal G)}S_{\mathcal G}\,\ell_{\mathcal G} \leq \frac{3e^{2\nu}}{\nu}\Bigl[E_\alpha + E_\beta + \frac{2}{\nu}\bigl(E_\alpha + E_\beta + E_o\bigr)E_o + \frac{2}{\nu}\bigl(E_\beta + E_o\bigr)E_\alpha\Bigr]\;,
\end{equation}
the sum over the groups made at the step $k + 1$ which are at $p$: those with the cube $p$ in $X_{\mathcal G}$, or with the traversal $p$ in an interval; and the same bound holds with the further factor $e^{\frac{7}{16}\nu\sigma_{k+1}(\mathcal G)}$ in each term of the sum. If moreover $\mathsf E_k \leq C_{\mathsf E}\rho_k$ with a constant $C_{\mathsf E}$, and
\begin{equation}\label{eq:size-constant}
    C_{S'} \coloneqq \frac{3e^{2\nu}}{\nu}(1 + \beta_0)^6C_\alpha\;,
\end{equation}
then the right side of \eqref{eq:size-bound} is at most $2C_{S'}\rho_{k+1}$ for $\gamma_q$ sufficiently small, depending on $C_{\mathsf E}$ and on the constants of Proposition~\ref{prop:rooted-sums}.
\end{corollary}

\begin{proof}
The number $c_{\mathcal G}$ is the value of $Y_{\mathcal G}$ at the reference data, which belong to the space of the step $k + 1$, so that $\abs{c_{\mathcal G}} \leq P(\mathcal G)$ by \eqref{eq:group-sup}. The transport along the interval of $\mathcal G$ has norm at most $e^{c'_\pi\alpha_{1,k+1}n_{k+1}(\mathcal G)} \leq e^{\nu\sigma_{k+1}(\mathcal G)}$ on that space, by the hypothesis $4\bar\Lambda Mc'_\pi\alpha_{1,k+1} \leq \nu$ of Lemma~\ref{lem:group-weight} and $n_{k+1}(\mathcal G) \leq M\sigma_{k+1}(\mathcal G)$. Hence $\norm{Y^{\mathrm d}_{\mathcal G}}_{(k+1)} \leq (1 + e^{\nu\sigma_{k+1}(\mathcal G)})P(\mathcal G)$ and $S_{\mathcal G} \leq 3e^{\nu\sigma_{k+1}(\mathcal G)}P(\mathcal G)$. The exponent on the left of \eqref{eq:size-bound} is then at most $(\frac{1}{16} + 2)\nu\sigma_{k+1}(\mathcal G) + \nu_e\operatorname{ext}(\mathcal G)$, which is less than the exponent of \eqref{eq:size-bound-new} by $\frac{7}{16}\nu\sigma_{k+1}(\mathcal G)$, and a traversal in an interval of $\mathcal G$ meets a cube of $X_{\mathcal G}$; so \eqref{eq:size-bound}, with or without the further factor, follows from Theorem~\ref{thm:size-bound}.

For the last statement, $\rho_k \leq (1 + \beta_0)^6\rho_{k+1}$ by \eqref{eq:rho-ratio}, so that the first term on the right of \eqref{eq:size-bound} is $(3e^{2\nu}/\nu)C_\alpha\rho_k \leq C_{S'}\rho_{k+1}$. By \eqref{eq:rooted-totals}, $E_\beta \leq C_\beta C_{\mathsf E}\varpi_k\rho_k$ and $E_o \leq K_{\mathrm{box}}(e^{1/60}C_{\mathsf E} + 6C_{S'})\rho_k$, so that each of the other terms is at most $\rho_k$ times a constant times $\varpi_k$ or $\rho_k/\nu$. Both tend to zero with $\gamma_q$, uniformly in $k \leq K'$: $\rho_k \leq C_\rho\gamma_q^2(1 + A_1(\log\gamma_q^{-2})^{p_0})^4$ because $g_k \leq \gamma_q$ by \eqref{eq:stopping-coupling}, and $\varpi_k$ is at most a constant times $\vartheta_k/\min(\alpha_{0,k}, \alpha_{1,k})$, a negative power of $\log\hat g_k^{-2}$ as in \eqref{eq:expansion-condition}. The other terms together are therefore at most $C_{S'}\rho_{k+1}$ for $\gamma_q$ small.
\end{proof}

By \eqref{eq:size-bound} with $\theta = 0$ the activities $c_{\mathcal G}$ of the reference gas of the step $k + 1$ satisfy the hypothesis of Lemma~\ref{lem:root-bound}, whose proof verifies the condition of Koteck\'y and Preiss with the weight $a^\circ(\mathcal G) = (\nu/32)\sigma_{k+1}(\mathcal G)$; the theorem cited there bounds the sum of the absolute values of the weights of the clusters which are incompatible with a group $\mathcal G$ by $a^\circ(\mathcal G)$. With $\theta(\mathcal G)$ defined as that sum we have $\theta(\mathcal G) \leq \nu\sigma_{k+1}(\mathcal G)/16$, which is clause (i) of Definition~\ref{def:dressed-loop} at the step $k + 1$ for the groups made there, and \eqref{eq:size-bound} with these numbers is the production bound \eqref{eq:production-clause} at the step $k + 1$. The definition of $\theta$ uses the bound with $\theta = 0$ only, so that the argument is not circular.

\subsection{The Deviation Bound of the New Groups}\label{sec:new-deviation}

The production bound controls the size of a new group, and the class asks for more: at a real background the deviation part of an insertion has to be small in proportion to the curvature of the background near the group, as in \eqref{eq:per-insertion}. That it is so comes from the subtraction of the flat value. The deviation part $Y^{\mathrm d}_{\mathcal G}$ is a holomorphic function on the space of the step $k + 1$ which vanishes at the trivial configuration, and near the group the background is, in a suitable gauge, the exponential of a field which is bounded by the local curvature times the linear size of the region (Supplement, Lemma~17.22). On the complex line through the trivial configuration and the background in that gauge, Schwarz's lemma converts the bound of the function on the space into a bound proportional to that field.

We keep the setting of Section~\ref{sec:count}, with the units of the scale $k + 1$, in which the fine lattice has the spacing $\xi = L^{-(k+1)}$. Let $\mathcal G$ be a group made at the step $k + 1$ with one interval, from the site $x_{\mathcal G}$ to the site $x'_{\mathcal G}$, let $\mathfrak n$ be a traversal of the step $k + 1$ in that interval, $t = t_{\mathfrak n}(\mathcal G)$, and let $\mathcal Q = \mathcal Q_{\mathfrak n}(\mathcal G)$ be the smallest cube which is a union of cubes of $\pi_{k+1}$ and contains the ball $B_\infty(\mathfrak n, R_t)$, hence $X_{\mathcal G}$. We say that $\mathcal G$ is of the \emph{deviation species} if for each of these traversals the side of the cube $\mathcal Q^{\sim 5}$, which is obtained from $\mathcal Q$ by adding five layers of cubes, is less than the smallest period of the torus; the other groups with one interval are those of the wrapping species (Section~\ref{sec:assembly}). Put
\begin{equation}\label{eq:cauchy-constants}
    C_g \coloneqq 3B_3(2C_{\mathrm{box}} + 5)M\;,
    \qquad
    c_{\mathrm{an}} \coloneqq \min\Bigl(\frac34,\ \frac{c_F}{2c_G}\Bigr)\;,
    \qquad
    \bar\alpha_{k+1} \coloneqq \min(\alpha_{0,k+1}, \alpha_{1,k+1})\;,
\end{equation}
with the constant $B_3$ of \cite[Theorem~1, p.~279]{BalabanVAR}, the constant $c_F$ of \eqref{eq:space-of-step}, and the constant $c_G$ of \eqref{eq:complex-curvature}.

\begin{proposition}[The Deviation Bound of a New Group]\label{prop:new-deviation}
Let the standing hypotheses of Supplement, Lemma~17.22, hold, let $\gamma$ be sufficiently small, and suppose that the insertions of the groups made at the step $k + 1$ are holomorphic on the space of the step $k + 1$, as functions of the restriction of $(\mathbf U, \mathbf J)$ to $X_{\mathcal G}$, and satisfy \eqref{eq:covariance-law}. Let $V_{k+1}$ be a block field in the support of the term. Then, for every group $\mathcal G$ of the deviation species made at the step $k + 1$ and every traversal $\mathfrak n$ of the step $k + 1$ in its interval,
\begin{equation}\label{eq:new-per-insertion}
    \norm{Y^{\mathrm d}_{\mathcal G}(U_{k+1}, J_{k+1})} \leq \frac{C_g}{c_{\mathrm{an}}\bar\alpha_{k+1}}\,S_{\mathcal G}\operatorname{Dev}_{\mathfrak n}(\mathcal G; U_{k+1})\;,
\end{equation}
so that \eqref{eq:per-insertion} holds at the step $k + 1$ with $b_{\mathcal G} \coloneqq (C_g/(c_{\mathrm{an}}\bar\alpha_{k+1}))e^{\theta(\mathcal G)}S_{\mathcal G}$. Moreover, under the hypotheses of Corollary~\ref{cor:production-bound}, with $\beta_{\mathcal G} = b_{\mathcal G}\ell_{\mathcal G}L^{2(k+1)}$, for every traversal $\mathfrak n$ of the step $k + 1$ and every $t \geq 1$,
\begin{equation}\label{eq:new-tail}
    \sum_{\mathcal G \ni \mathfrak n,\ t_{\mathfrak n}(\mathcal G) \geq t}e^{\nu\sigma_{k+1}(\mathcal G) + \nu_e\operatorname{ext}(\mathcal G)}\beta_{\mathcal G} \leq \frac{4C_gC_{S'}}{c_{\mathrm{an}}}\,\tilde\rho_{k+1}L^{2(k+1)}W_t\;,
\end{equation}
the sum over the groups of the deviation species made at the step $k + 1$ with $\mathfrak n$ in the interval.
\end{proposition}

\begin{proof}
\emph{The region.} The cube $\mathcal Q$ has side at most $2R_t + 1 + 2M$, and $\mathcal Q^{\sim 5}$ lies in the ball $B_\infty(\mathfrak n, R_t + 6M)$, hence in the comparison region $B_\infty(\mathfrak n, R_{t+1})$ of $\mathcal G$ at $\mathfrak n$, because $R_{t+1} - R_t = C_{\mathrm{box}}M \geq 6M$. Let $E$ be the largest value of $\xi^{-2}\abs{U_{k+1}(\partial p) - 1}$ over the fine plaquettes $p$ in $\mathcal Q^{\sim 5}$, as in Supplement, Lemma~17.22. Since $\xi^{-2} = L^{2(k+1)}$, the definition \eqref{eq:insertion-deviation} gives $\ell_{\mathcal G}E \leq \operatorname{Dev}_{\mathfrak n}(\mathcal G; U_{k+1})$. The domain $X_{\mathcal G}$ is a connected union of $\ell_{\mathcal G}$ cubes of side $M$, one of which meets $\mathfrak n$ (Lemma~\ref{lem:group-covering}(iv)), so that it lies in the ball of radius $M\ell_{\mathcal G} + 1$ about $\mathfrak n$, and therefore $R_t \leq \max(2C_{\mathrm{box}}M, M\ell_{\mathcal G} + 1 + C_{\mathrm{box}}M) \leq 2C_{\mathrm{box}}M\ell_{\mathcal G}$ by the definition of $t$. The number $M_{\mathcal Q}$ of that lemma is half the side of $\mathcal Q^{\sim 2}$, rounded up to a multiple of a number which is at most $M$; hence $M_{\mathcal Q} \leq R_t + 5M \leq (2C_{\mathrm{box}} + 5)M\ell_{\mathcal G}$, and
\begin{equation}\label{eq:local-field-size}
    x \coloneqq 3B_3M_{\mathcal Q}E \leq C_g\operatorname{Dev}_{\mathfrak n}(\mathcal G; U_{k+1})\;.
\end{equation}

\emph{The trivial case.} If $x \geq c_{\mathrm{an}}\bar\alpha_{k+1}$, the right side of \eqref{eq:new-per-insertion} is at least $S_{\mathcal G}$, which bounds the left side, since $(U_{k+1}, J_{k+1})$ is a point of the space of the step $k + 1$. We assume from now on that $x < c_{\mathrm{an}}\bar\alpha_{k+1}$. Then $2E \leq 2x/(3B_3M)$ and $2M_{\mathcal Q}E \leq 2x/(3B_3)$ are less than $\bar\alpha_{k+1}$, which tends to zero with $\gamma$, so that for $\gamma$ small the two hypotheses of Supplement, Lemma~17.22, on $E$ and on $M_{\mathcal Q}E$ hold, their right sides being constants; here we used $M_{\mathcal Q} \geq M$.

\emph{The gauge and the function.} By Supplement, Lemma~17.22, applied to the cube $\mathcal Q$, which is embedded in the torus with its five layers by the definition of the species, there is a $G$-valued gauge transformation $u$ on a neighbourhood of $\mathcal Q^{\sim 2}$ such that $U^u_{k+1} = e^{i\xi A}$ on $\mathcal Q^{\sim 2}$ with $\abs A, \abs{\nabla^\xi A} \leq x$ and $\abs{R(u)J_{k+1}} \leq \frac43x$. For a complex number $z$ with $\abs zx < c_{\mathrm{an}}\bar\alpha_{k+1}$ the pair $(e^{iz\xi A}, zR(u)J_{k+1})$ satisfies the conditions \eqref{eq:space-of-step} at $n = k + 1$ around the trivial configuration, on the bonds of $\mathcal Q^{\sim 2}$. Indeed, in the units of the fine lattice its complex part is $\alpha = z\xi A$, of norm at most $\abs zx\,\xi < \alpha_{1,k+1}L^{-(k+1)}$; its covariant differences at the trivial configuration are its differences, of norm at most $\abs zx\,\xi^2 < \alpha_{1,k+1}L^{-2(k+1)}$; by \eqref{eq:complex-curvature} with $U = 1$, and since $\abs zx \leq 1$, its curvature is at most $c_G\xi^2(\abs zx + (\abs zx)^2) \leq 2c_G\abs zx\,\xi^2 < c_F\alpha_{0,k+1}L^{-2(k+1)}$; and the current has norm at most $\frac43\abs zx < \alpha_{0,k+1}$ in the units of the scale $k + 1$. By hypothesis the function
\[
    f(z) \coloneqq Y^{\mathrm d}_{\mathcal G}\bigl(e^{iz\xi A}, zR(u)J_{k+1}\bigr) = Y_{\mathcal G}\bigl(e^{iz\xi A}, zR(u)J_{k+1}\bigr) - c_{\mathcal G}\,\pi\bigl(e^{iz\xi A}(I_{\mathcal G})\bigr)
\]
is therefore holomorphic on the disc of radius $R_X \coloneqq c_{\mathrm{an}}\bar\alpha_{k+1}/x > 1$, where $\norm{f(z)} \leq \norm{Y^{\mathrm d}_{\mathcal G}}_{(k+1)} \leq S_{\mathcal G}$; here we used that $Y_{\mathcal G}$ and the transport along $I_{\mathcal G}$ depend on the configuration only on $X_{\mathcal G} \subset \mathcal Q$ (Lemma~\ref{lem:group-covering}(v)). At $z = 0$ the configuration is the trivial one with the current $0$, and $f(0) = 0$ by the definition of the flat value $c_{\mathcal G}$. At $z = 1$ it is the gauge transform of $(U_{k+1}, J_{k+1})$ by $u$ on $X_{\mathcal G}$, and by \eqref{eq:covariance-law}, which the transport along $I_{\mathcal G}$ satisfies as well,
\[
    f(1) = \pi(u(x_{\mathcal G}))\,Y^{\mathrm d}_{\mathcal G}(U_{k+1}, J_{k+1})\,\pi(u(x'_{\mathcal G}))^{-1}\;,
\]
with unitary factors. Schwarz's lemma, applied to the matrix elements of $f$, gives $\norm{Y^{\mathrm d}_{\mathcal G}(U_{k+1}, J_{k+1})} = \norm{f(1)} \leq S_{\mathcal G}/R_X = S_{\mathcal G}\,x/(c_{\mathrm{an}}\bar\alpha_{k+1})$, and \eqref{eq:new-per-insertion} follows from \eqref{eq:local-field-size}.

\emph{The tails.} By the definitions, $e^{\nu\sigma_{k+1}(\mathcal G) + \nu_e\operatorname{ext}(\mathcal G)}\beta_{\mathcal G}$ is $(C_g/(c_{\mathrm{an}}\bar\alpha_{k+1}))L^{2(k+1)}$ times the term of $\mathcal G$ in \eqref{eq:size-bound}. If $t_{\mathfrak n}(\mathcal G) \geq t \geq 2$, then $X_{\mathcal G}$ is not contained in the ball of radius $R_{t-1} = C_{\mathrm{box}}Mt$ about $\mathfrak n$, so that $M\ell_{\mathcal G} + 1 > C_{\mathrm{box}}Mt$ and $\sigma_{k+1}(\mathcal G) \geq \ell_{\mathcal G} \geq C_{\mathrm{box}}t - 1$. The further factor $e^{\frac{7}{16}\nu\sigma_{k+1}(\mathcal G)}$ of Corollary~\ref{cor:production-bound} is then at least $e^{-\nu}e^{\eta_wt}$, because the rate of the weights \eqref{eq:tail-weights} is $\eta_w = \frac{5}{16}\nu \leq \frac{7}{16}\nu C_{\mathrm{box}}$. So the sum of the terms of \eqref{eq:size-bound} over the groups with $t_{\mathfrak n}(\mathcal G) \geq t$ is at most $2e^\nu C_{S'}\rho_{k+1}e^{-\eta_wt} \leq 2e^\nu C_{S'}\rho_{k+1}W_t$, since $W_t \geq w_t \geq e^{-\eta_wt}$; and for $t = 1$ the sum over all the groups is at most $2C_{S'}\rho_{k+1} \leq 2C_{S'}\rho_{k+1}W_1$, since $W_1 \geq w_1 = 8e^{-\eta_w} \geq 1$. With $2e^\nu \leq 4$ and $\tilde\rho_{k+1} = \rho_{k+1}/\bar\alpha_{k+1}$ this is \eqref{eq:new-tail}.
\end{proof}

The factor $1/\bar\alpha_{k+1}$ in \eqref{eq:new-per-insertion} is the price of Schwarz's lemma: the function is bounded on a space whose radii are $\alpha_{0,k+1}$ and $\alpha_{1,k+1}$, and it is evaluated at a point whose distance from the centre is of the order of the deviation. This is why the first term of \eqref{eq:insertion-bound} has the coefficient $\tilde\rho_{k+1} = \rho_{k+1}/\bar\alpha_{k+1}$, and it is the reason for the two radii having exponents $q_0 = q_1$ larger than that of the thresholds: the deviation of a background of the support is of the order of $\varepsilon_{k+1}$, and $\varepsilon_{k+1}/\bar\alpha_{k+1}$ tends to zero with the coupling.

\subsection{The Contraction of the Deviation Parts}\label{sec:contraction}

It remains to carry the groups of the step $k$ which the step leaves alone. A member of a family which no component touches and which is tied to nothing forms a group of the step $k + 1$ by itself, and its insertion is the old function, evaluated at the new background $U_{k+1}$ (Proposition~\ref{prop:regrouping}). Nothing has to be estimated anew for it. The new background is a background of the step $k$, so that the bounds of the class hold at it as they stand, and what changes is only the unit in which the traversals are counted: a traversal of the step $k + 1$ consists of $L$ traversals of the step $k$, each with its carried groups, while the deviation of the background in the units of the scale $k + 1$ is $L^2$ times the same deviation in the units of the scale $k$. This is the contraction by $L \cdot L^{-2} = L^{-1}$ which was announced at the beginning of the section. The parts of an old insertion which the fluctuation of the step changes do not appear here: an insertion which a component touches lies in a marked polymer and thereby in a group made at the step $k + 1$, whose deviation part was bounded in Proposition~\ref{prop:new-deviation}.

\begin{proposition}[The Carried Groups at the Real Backgrounds]\label{prop:carried-real}
Let the ensemble be dressed at the step $k$, let $\gamma$ be so small that Supplement, Proposition~S.8, applies, and consider the terms without large-field regions of $\varrho_k$ and of $\varrho_{k+1}$, with a block field $V_{k+1}$ in the support of the second.
\begin{enumerate}[label=(\roman*)]
\item For every group $W$ of the deviation species of the step $k$, made at a step $j \leq k$, and every traversal $\mathfrak n$ of the step $j$ in its interval, $e^{\theta(W)}\norm{Y_W(U_{k+1}, J_{k+1})} \leq b_W\operatorname{Dev}_{\mathfrak n}(W; U_{k+1})$, and $Y_W$ vanishes at the reference data.
\item If \eqref{eq:tail-clause} holds at the step $k$ with the constant $C^{(k)}_\chi$, then for every traversal $\mathfrak n'$ of the step $k + 1$ and every $t \geq 1$
\begin{equation}\label{eq:contraction-first}
    \sum_{W \ni \mathfrak n',\ t^{(k+1)}_{\mathfrak n'}(W) \geq t}e^{\nu\sigma(W) + \nu_e\operatorname{ext}(W)}\beta_W \leq L^{-1}\frac{\tilde\rho_k}{\tilde\rho_{k+1}}\,C^{(k)}_\chi\tilde\rho_{k+1}L^{2(k+1)}W_t\;,
\end{equation}
the sum over the groups of the deviation species of the step $k$ with $\mathfrak n'$ in an interval.
\end{enumerate}
\end{proposition}

\begin{proof}
(i) The block field $V^*_k \coloneqq M^k(U_{k+1})$ lies in the support of the term without large-field regions of $\varrho_k$, by Supplement, Proposition~S.8, and $U_k(V^*_k) = U_{k+1}$ by the uniqueness of the minimal orbit (Section~\ref{sec:evaluation-point}). The bound \eqref{eq:per-insertion} of the step $k$ holds at every block field of the support, in particular at $V^*_k$, where the background is $U_{k+1}$ and the argument of the insertion in the place of $\mathbf J$ is the current of that background, which is $J_{k+1}$; the insertion is the same function at both steps (Definition~\ref{def:dressed-loop}(iii)). The insertions of the deviation species vanish at the reference data at the step at which they are made, by the subtraction of the flat value, and the reference data of the two steps are the same configuration.

(ii) A group $W$ of the sum has a traversal $\mathfrak n \subset \mathfrak n'$ of the step $k$ in its interval, one of the $L$ traversals of the step $k$ in $\mathfrak n'$, and $t^{(k)}_{\mathfrak n}(W) \geq t^{(k+1)}_{\mathfrak n'}(W)$ (Section~\ref{sec:dressed-loops}). The left side is therefore at most the sum over these $L$ traversals of the left side of \eqref{eq:tail-clause} at the step $k$, that is, at most $L \cdot C^{(k)}_\chi\tilde\rho_kL^{2k}W_t$, which is the right side of \eqref{eq:contraction-first}.
\end{proof}

The ratio of the amplitudes of the two steps is controlled by Lemma~\ref{lem:tube-couplings}: $g_k \leq (1 + \beta_0)g_{k+1}$ by \eqref{eq:tube-two-six}, $(\log g_k^{-2})^p \leq (1 + \beta_0)(\log g_{k+1}^{-2})^p$ by \eqref{eq:tube-two-seven}, and $\alpha_{i,k+1} \leq (1 + \beta_0)^2\alpha_{i,k}$ by \eqref{eq:tube-two-eight}, so that
\begin{equation}\label{eq:rho-ratio}
    \frac{\rho_k}{\rho_{k+1}} \leq (1 + \beta_0)^6\;,
    \qquad
    \frac{\tilde\rho_k}{\tilde\rho_{k+1}} \leq (1 + \beta_0)^8\;.
\end{equation}

For $k = 1$ the second ratio is $(\rho_1/\rho_2)\min(\alpha_{0,2}, \alpha_{1,2})$, since $\tilde\rho_1 = \rho_1$, and it is at most $(1 + \beta_0)^6$.

\emph{The recursion.} At the step $k + 1$ the groups of the deviation species on the term without large-field regions are the carried ones and those made at the step, and both satisfy \eqref{eq:per-insertion}, by Propositions~\ref{prop:carried-real}(i) and~\ref{prop:new-deviation}. Adding \eqref{eq:contraction-first} and \eqref{eq:new-tail}, we see that the tail clause \eqref{eq:tail-clause} holds at the step $k + 1$ with a constant which satisfies
\begin{equation}\label{eq:deviation-recursion}
    C^{(k+1)}_\chi \leq C_{\mathrm{new}} + aC^{(k)}_\chi\;,
    \qquad
    C_{\mathrm{new}} \coloneqq \frac{4C_gC_{S'}}{c_{\mathrm{an}}}\;,
    \qquad
    a \coloneqq L^{-1}(1 + \beta_0)^8\;,
\end{equation}
by \eqref{eq:rho-ratio}. No term of the second order in the old constants occurs: the contractions of the old insertions with each other and with the fluctuation field are parts of the groups made at the step, and they are contained in $C_{\mathrm{new}}$, which by \eqref{eq:size-constant} involves no constant of the class. At the steps at which the $\mathbf R$-operation acts, $C_{\mathrm{new}}$ also contains the source $C^{\mathbf R}(g_k) \leq C^{\mathbf R}(\gamma_q)$ of the deviation insertions which that operation makes, and $a$ contains the further term $C_aMR_ke^{-p_0(g_k)/4}$ (Section~\ref{sec:large-field-step}). Since $(1 + \beta_0)^8 \leq L/4$ (Appendix~\ref{app:constants}), and for $\gamma_q$ so small that $C_aMR_ke^{-p_0(g_k)/4} \leq 1/L$ for every $k \leq K'$, we have $a \leq \frac12$ at every step. Let $C^* \coloneqq \max(C^{(1)}_\chi, 2C_{\mathrm{new}})$, with the constant $C^{(1)}_\chi$ of the first step. If $C^{(k)}_\chi \leq C^*$, then
\begin{equation}\label{eq:recursion-fixed-point}
    C^{(k+1)}_\chi \leq C_{\mathrm{new}} + aC^* \leq \tfrac12C^* + \tfrac12C^* = C^*\;.
\end{equation}
So the tail clause holds at every step with the constant $C_\chi \coloneqq C^*$, and by Lemma~\ref{lem:tail-to-sum} so does the first term of \eqref{eq:insertion-bound}, with its weights. The factor $\tilde\rho_k = \rho_k/\min(\alpha_{0,k}, \alpha_{1,k})$ in these bounds is the price of Schwarz's lemma in Proposition~\ref{prop:new-deviation}. The supremum form \eqref{eq:sup-clause} of the same groups at the step $k + 1$ does not follow from a recursion of this kind, as we now explain.

\subsection{The Carried Insertions at Complex Data}\label{sec:carried-complex}

The clause \eqref{eq:sup-clause} of Definition~\ref{def:dressed-loop} bounds the insertions on the whole space of the step, and not only at its real backgrounds. For the groups which a step makes this is part of what the count of that step delivers, in the form \eqref{eq:production-clause}. For a group $W$ which was made at an earlier step $j$ and is carried to the step $k + 1$ the question is of another kind, since its insertion is the function which was bounded at the step $j$, and it has now to be bounded on the space of the step $k + 1$ and summed over the groups at a traversal of that step. The obvious bound fails. In the units of the fine lattice the space of the step $k + 1$ lies in that of the step $j$, so that the supremum $\norm{Y_W}_{(j)}$ of the step $j$ bounds the insertion; but a traversal of the step $k + 1$ consists of $L^m$ traversals of the step $j$, $m \coloneqq k + 1 - j$, each with carried groups of its own, and the sum of these suprema is larger than the bound of the step $j$ by the factor $L^m$. At the real backgrounds this factor is compensated by the contraction of the deviation (Section~\ref{sec:contraction}). At complex data there is no deviation to contract, and the compensation has to come from the size of the complex part: the radii of the space of the step $k + 1$ are smaller than those of the step $j$ by the factor $L^{-m}$, up to the ratio of the couplings, so that Schwarz's lemma on the larger space gains $L^{-m}$ for each insertion. This cancels the number of the traversals and no more, and the sum over the generations $m = 1, \dots, k$ would still grow with $k$.

What removes the growth is gauge covariance at complex gauge transformations. The complex part of a configuration of the space of the step $k + 1$ varies slowly, by the condition on its covariant differences, and on a set which is small in the units of the scale $k + 1$ it is therefore, to first order, a complex gauge transformation with a gauge function linear in the position. The dependence set of a carried group is such a set, of diameter $\delta_W$ of the order of $ML^{-m}\sigma(W)$ in these units. Removing the linear part by a gauge transformation which is the identity at the first site of the interval of $W$ leaves a complex part smaller by the factor $\delta_W$, and Schwarz's lemma then gains $L^{-2m}\sigma(W)$, which is summable over the generations; the factor $\sigma(W)$, and the groups whose dependence sets are not small, are paid by a sixteenth of the weight $e^{\nu\sigma(W)}$. This is the reason for the rate $\lambda_c < 1$ in \eqref{eq:sup-clause}. The insertions are covariant under gauge transformations with values in the compact group only, and the first step is to see that the covariance extends.

Throughout this subsection we work on the labelled term without large-field regions, we fix steps $1 \leq j \leq k$ and put $m \coloneqq k + 1 - j$, and we measure lengths in the units of the fine lattice, so that a traversal of the step $n$ consists of $L^n$ fine bonds and a cube of $\pi_n$ has the side $ML^n$. We write $\mathfrak g$ for the Lie algebra of $G$, $\mathfrak g^c \coloneqq \mathfrak g \otimes \mathbb C$, and $G^c$ for the complexification of $G$, the connected complex Lie group $G\exp(i\mathfrak g)$ of matrices with the Lie algebra $\mathfrak g^c$; a representation $\pi$ of $G$ extends to a holomorphic representation of $G^c$. The norm $\abs{\,\cdot\,}$ is that of Definition~\ref{def:dressed-loop}, on the matrices of the size of $G$, and it is invariant under the conjugation by the elements of $G$. The sets $\mathcal S_n(U)$ are those of \eqref{eq:space-of-step}; each of them is open in the space of the pairs of a $G^c$-valued configuration and a $\mathfrak g^c$-valued field, because $\alpha \mapsto e^{i\alpha}U(b)$ maps a neighbourhood of $0$ in $\mathfrak g^c$ biholomorphically onto a neighbourhood of $U(b)$ in $G^c$, and a gauge transformation $v$ with values in $G$ maps $\mathcal S_n(U)|_X$ onto $\mathcal S_n(U^v)|_X$, by the invariance of the norm.

\emph{What is used of the class.} Let $W$ be a group of the deviation species made at the step $j$; it has one interval $I_W$, from a site $x_W$ to a site $x'_W$, which lies in its domain $X_W$. We use the following.
\begin{enumerate}[label=(\alph*)]
\item By Definition~\ref{def:dressed-loop}(iii) the insertion $Y_W$ is holomorphic on the open set $\mathcal D_W$ of \eqref{eq:production-space}, it is bounded by $\norm{Y_W}_{(j)}$ there, it satisfies the covariance law \eqref{eq:covariance-law}, and at the later steps it is the same function of the configuration and of the current.
\item For $V_{k+1}$ in the support of the step $k + 1$ the background $U_{k+1}$ is a background of the step $j$: with $V^*_i \coloneqq M^i(U_{k+1})$ for $j \leq i \leq k$, the block field $V^*_i$ lies in the support of the step $i$ (Supplement, Proposition~S.8, by descending induction on $i$), and $U_i(V^*_i) = U_{k+1}$ by the uniqueness of the minimal orbit (Section~\ref{sec:evaluation-point}). Moreover $\abs{U_{k+1}(\partial p) - 1} < 2B_3\varepsilon_{k+1}L^{-2(k+1)}$ for every fine plaquette $p$ \cite[p.~268]{BalabanCRE}, and for $\gamma$ small the pair of $U_{k+1}$ and its current $J_{k+1}$ is a point of $\mathcal S_{k+1}(U_{k+1})$.
\item The domain $X_W$ is a connected union of at most $u_\star\sigma(W)$ cubes of $\pi_j$, and $I_W$ has at most $u_\star M\sigma(W)L^j$ fine bonds (Definition~\ref{def:dressed-loop}(vii)).
\item The radii of two steps $n' < n$ satisfy $\alpha_{i,n} \leq \varkappa_{n-n'}\,\alpha_{i,n'}$ for $i = 0, 1$, with $\varkappa_m \coloneqq (1 + \beta_0)^2m^{1/2}$, by \eqref{eq:tube-two-eight}, and $\alpha_{1,n} = \lambda_{\mathrm r}\alpha_{0,n}$ with $\lambda_{\mathrm r} \coloneqq C^\alpha_1/C^\alpha_0$, since $q_1 = q_0$ (Appendix~\ref{app:constants}).
\end{enumerate}
We assume
\begin{equation}\label{eq:carried-blocking-condition}
    L \geq 20(1 + \beta_0)^2\;,
\end{equation}
so that $\varkappa_mL^{-m} \leq (1 + \beta_0)^2L^{-1} \leq \frac{1}{20}$ for every $m \geq 1$, because $m^{1/2}L^{-(m-1)} \leq 1$ for $L \geq 2$.

\begin{lemma}[The Space of a Later Step Lies in That of an Earlier One]\label{lem:space-inclusion}
Let $U$ be a $G$-valued configuration. Then $\mathcal S_{k+1}(U) \subset \mathcal S_j(U)$.
\end{lemma}

\begin{proof}
By (d) and \eqref{eq:carried-blocking-condition}, $\alpha_{i,k+1}L^{-(k+1)} \leq \varkappa_mL^{-m}\alpha_{i,j}L^{-j} \leq \alpha_{i,j}L^{-j}$, and the same holds with the powers $L^{-2(k+1)}$ and $L^{-3(k+1)}$, which are smaller. Each of the four radii of \eqref{eq:space-of-step} at $n = k + 1$ is therefore at most the corresponding radius at $n = j$, and for $j = 1$ only the first condition is posed.
\end{proof}

\begin{lemma}[Complexified Covariance]\label{lem:complexified-covariance}
Let $X$ be a finite set of bonds with the set $S$ of its sites, $x, x' \in S$, $\mathcal D$ an open set of pairs $q = (\mathbf U, \mathbf J)$ of a $G^c$-valued configuration and a $\mathfrak g^c$-valued field on $X$ which is invariant under the $G$-valued gauge transformations, and $F$ a holomorphic function on $\mathcal D$ with values in the linear maps of the space of a representation $\pi$ of $G$, such that $F(q^v) = \pi(v(x))F(q)\pi(v(x'))^{-1}$ for every $q \in \mathcal D$ and every $v \in G^S$, where $q^v \coloneqq (\mathbf U^v, R(v)\mathbf J)$. Let $q \in \mathcal D$ and let $t \mapsto u_t$, $t \in [0, 1]$, be a continuous path in $(G^c)^S$ with $u_0 = 1$ and $q^{u_t} \in \mathcal D$ for every $t$. Then
\begin{equation}\label{eq:complexified-covariance}
    F(q^{u_1}) = \pi(u_1(x))\,F(q)\,\pi(u_1(x'))^{-1}\;.
\end{equation}
\end{lemma}

\begin{proof}
The map $u \mapsto q^u$ is holomorphic on $(G^c)^S$, so that $\mathcal U \coloneqq \{ u \in (G^c)^S \,|\, q^u \in \mathcal D \}$ is open, and it contains $G^S$ by the invariance of $\mathcal D$. Since $\pi$ is holomorphic on $G^c$, the function
\[
    \Psi(u) \coloneqq \pi(u(x))^{-1}F(q^u)\pi(u(x')) - F(q)
\]
is holomorphic on $\mathcal U$, and it vanishes on $G^S$ by hypothesis. In the chart $Z \mapsto \exp(Z)$ of $(G^c)^S$ at the identity, with $Z \in (\mathfrak g^c)^S = \mathfrak g^S \oplus i\mathfrak g^S$, the function $\Psi \circ \exp$ is holomorphic near $0$ and vanishes on the real subspace $\mathfrak g^S$, whose complex span is the whole space. All its derivatives at $0$ in the real directions vanish, hence, by the complex linearity of the derivatives of a holomorphic function, all its derivatives at $0$, and $\Psi$ vanishes on a neighbourhood of the identity. By the identity theorem it vanishes on the connected component of $\mathcal U$ which contains the identity, and the path $u_t$ lies in this component, because it lies in $\mathcal U$ and starts at the identity. Hence $\Psi(u_1) = 0$, which is \eqref{eq:complexified-covariance}.
\end{proof}

We use the lemma for the insertion of $W$ with the transport along its interval removed. Let $\mathbf U(I_W)$ be the ordered product of $\mathbf U$ along the fine bonds of $I_W$, from $x_W$ to $x'_W$, an element of $G^c$, and
\begin{equation}\label{eq:base-point-form}
    \hat Y_W(\mathbf U, \mathbf J) \coloneqq Y_W(\mathbf U, \mathbf J)\,\pi\bigl(\mathbf U(I_W)\bigr)^{-1}\;.
\end{equation}
Since $\mathbf U^v(I_W) = v(x_W)\mathbf U(I_W)v(x'_W)^{-1}$, the function $\hat Y_W$ is holomorphic on $\mathcal D_W$ and satisfies \eqref{eq:covariance-law} with $x_W$ in the place of $x'_W$. By Lemma~\ref{lem:complexified-covariance}, for a path $u_t$ as there with $u_1(x_W) = 1$,
\begin{equation}\label{eq:base-point-invariance}
    \hat Y_W(\mathbf U^{u_1}, R(u_1)\mathbf J) = \hat Y_W(\mathbf U, \mathbf J)\;.
\end{equation}

The gauge transformation which we shall use is built from the following function. Three properties of the norm enter, each with a constant which depends only on $G$ and on the norm; we fix $\rho_0 \in (0, 1]$ and $c_G \geq 1$ such that all three hold.
\begin{enumerate}[label=(N\arabic*)]
\item For $X, Q \in \mathfrak g^c$ with $\abs X \leq \rho_0$, $\abs{e^{iX}Qe^{-iX} - Q} \leq c_G\abs X\abs Q$.
\item For $X, Y, Z, A \in \mathfrak g^c$ of norm less than $\rho_0$ the matrix $e^{iX}e^{iY}e^{iZ}e^{-i(X + A)}e^{-iZ}$ is within the radius of convergence of the logarithm about $1$, and
\begin{equation}\label{eq:reduction-function}
    \Phi(X, Y, Z, A) \coloneqq -i\log\bigl(e^{iX}e^{iY}e^{iZ}e^{-i(X + A)}e^{-iZ}\bigr)
\end{equation}
is a holomorphic function with values in $\mathfrak g^c$; the function $\Psi \coloneqq \Phi - (Y - A)$ has all its first and second partial derivatives bounded by $c_G$ in norm where the four arguments have norm at most $\rho_0/2$.
\item Let $U$ be a $G$-valued configuration and $\alpha$ a field with $\abs{\alpha(b)} \leq a_\circ \leq \rho_0$ and $\abs{(\delta_e\alpha)(b)} \leq d_\circ \leq a_\circ$ on the bonds of a fine plaquette $p$. Then
\begin{equation}\label{eq:complex-curvature}
    \abs{(e^{i\alpha}U)(\partial p) - 1} \leq c_G\bigl(\abs{U(\partial p) - 1} + d_\circ + a_\circ^2\bigr)\;.
\end{equation}
\end{enumerate}
Such constants exist. (N1) holds because $Q \mapsto e^{iX}Qe^{-iX} - Q$ is holomorphic in $X$ and vanishes at $X = 0$. In (N2) the product is a holomorphic function of the arguments which equals $1$ at the origin, the logarithm of a product of exponentials of small elements of $\mathfrak g^c$ is given by the series of Baker, Campbell, and Hausdorff and lies in $\mathfrak g^c$, and the derivatives are bounded by Cauchy's estimates. For (N3), let $b_1 = \langle x, x + e_\mu\rangle$, $b_2$, $b_3$, $b_4 = \langle x, x + e_\nu\rangle$ be the bonds of $p$ in the order $(\partial p) = b_1b_2b_3^{-1}b_4^{-1}$, $U_i \coloneqq U(b_i)$, $\alpha_i \coloneqq \alpha(b_i)$, and $F \coloneqq U(\partial p)$. Moving the factors $U_i$ to the right we find
\begin{gather*}
    (e^{i\alpha}U)(\partial p) = e^{i\alpha_1}e^{i\beta_2}e^{-i\beta_3}e^{-i\beta_4}F\;,\\
    \beta_2 \coloneqq U_1\alpha_2U_1^{-1}\;,
    \qquad
    \beta_3 \coloneqq FU_4\alpha_3U_4^{-1}F^{-1}\;,
    \qquad
    \beta_4 \coloneqq F\alpha_4F^{-1}\;.
\end{gather*}
Here $U_1\alpha_2U_1^{-1} = \alpha_4 + (\delta_{e_\mu}\alpha)(b_4)$ and $U_4\alpha_3U_4^{-1} = \alpha_1 + (\delta_{e_\nu}\alpha)(b_1)$ by \eqref{eq:covariant-difference}, and the conjugation by $F$ changes an element $Q$ by at most a constant times $\abs{F - 1}\abs Q$. The sum $\alpha_1 + \beta_2 - \beta_3 - \beta_4$ of the exponents is therefore at most $2d_\circ$ plus a constant times $\abs{F - 1}a_\circ$ in norm, the product of the four exponentials differs from $1 + i(\alpha_1 + \beta_2 - \beta_3 - \beta_4)$ by at most a constant times $a_\circ^2$, since its arguments are bounded by $2a_\circ$, and the multiplication by $F$ changes norms by at most a constant factor. This gives \eqref{eq:complex-curvature}.

Two consequences of the definition \eqref{eq:reduction-function} will be used: $\Phi(X, 0, 0, 0) = 0$ for every $X$ and $\Phi(0, 0, Z, 0) = 0$ for every $Z$, and the linear part of $\Phi$ at the origin is $Y - A$. Hence, where the four arguments have norm at most $\rho_0/2$,
\begin{equation}\label{eq:reduction-remainder}
\begin{gathered}
    \abs{\Psi(X, Y, Z, A)} \leq c_G\bigl[(\abs Y + \abs A)(\abs X + \abs Y + \abs Z + \abs A) + \abs X\abs Z\bigr]\;,\\
    \abs{\partial_X\Psi} \leq c_G(\abs Y + \abs Z + \abs A)\;,
    \qquad
    \abs{\partial_Z\Psi} \leq c_G(\abs X + \abs Y + \abs A)\;,\\
    \abs{\partial_Y\Psi} \leq c_G(\abs X + \abs Y + \abs Z + \abs A)\;.
\end{gathered}
\end{equation}
For the bounds of the derivatives we note that $\partial_X\Psi$ vanishes at $(X, 0, 0, 0)$, $\partial_Z\Psi$ at $(0, 0, Z, 0)$, and $\partial_Y\Psi$ at the origin, and we integrate the second derivatives along the segment from these points. For the first line, $\Psi(X, 0, Z, 0)$ vanishes if $X = 0$ or $Z = 0$, so that it equals the integral of its mixed second derivative over $[0, 1]^2$ along $(sX, tZ)$ and is at most $c_G\abs X\abs Z$; and $\Psi(X, Y, Z, A) - \Psi(X, 0, Z, 0)$ is the integral of the derivative in $(Y, A)$ along $(X, sY, Z, sA)$, which vanishes at the origin and is therefore at most $c_G(\abs X + \abs Y + \abs Z + \abs A)$.

\begin{lemma}[The Gauge Reduction of a Slowly Varying Complex Part]\label{lem:gauge-reduction}
Let $Q$ be a cube of the fine lattice of side $D \geq 1$ and $x_0$ a site of $Q$ with $x_0 + e_\mu \in Q$ for every $\mu$. Let $U = e^{i\zeta}$ on the bonds of $Q$, with a field $\zeta$ with values in $\mathfrak g$, and let $\alpha$ be a field on the bonds of $Q$ with values in $\mathfrak g^c$, such that for all bonds $b$ and unit vectors $e$ with $b, b + e$ in $Q$
\begin{equation}\label{eq:reduction-hypotheses}
    \abs{\alpha(b)} \leq a\;,
    \qquad
    \abs{(\delta_e\alpha)(b)} \leq a'\;,
    \qquad
    \abs{\zeta(b)} \leq z\;,
    \qquad
    \abs{\zeta(b + e) - \zeta(b)} \leq z'\;,
\end{equation}
with $a' \leq a$ and $4aD, a, z \leq \rho_0/2$. Put
\begin{equation}\label{eq:gauge-function}
    \omega(x) \coloneqq \sum_\mu a_\mu\,(x - x_0)_\mu\;,
    \qquad
    a_\mu \coloneqq \alpha(\langle x_0, x_0 + e_\mu\rangle)\;,
    \qquad
    u_t \coloneqq e^{it\omega}\;,
\end{equation}
for the sites $x$ of $Q$ and $t \in [0, 1]$. Then $u_t(x_0) = 1$, $\abs{\omega} \leq 4aD$, and $(e^{i\alpha}U)^{u_t} = e^{i\alpha_t}U$ on the bonds of $Q$, with fields $\alpha_t$ with values in $\mathfrak g^c$ which satisfy, with $a'' \coloneqq a' + c_Gza$ and $r_\Phi \coloneqq c_G[2a(4aD + 2a + z) + 4aDz]$,
\begin{equation}\label{eq:reduced-phase}
    \abs{\alpha_1(b)} \leq 4Da'' + r_\Phi\;,
    \qquad
    \abs{\alpha_t(b)} \leq 2a + r_\Phi\;,
\end{equation}
\begin{equation}\label{eq:reduced-difference}
    \abs{(\delta_e\alpha_t)(b)} \leq a'' + c_Gz(2a + r_\Phi) + c_G\bigl[(2a + z)a + (4aD + 2a + z)a'' + (4aD + 2a)z'\bigr]\;.
\end{equation}
\end{lemma}

\begin{proof}
The value $\omega(x_0) = 0$ and the bound on $\omega$ are immediate from \eqref{eq:gauge-function}, since $\abs{a_\mu} \leq a$ and $\sum_\mu\abs{(x - x_0)_\mu} \leq 4D$. For a bond $b = \langle x, x + e_\mu\rangle$ of $Q$ we have $\omega(x + e_\mu) = \omega(x) + a_\mu$, and therefore
\begin{gather*}
    (e^{i\alpha}U)^{u_t}(b) = e^{it\omega(x)}e^{i\alpha(b)}e^{i\zeta(b)}e^{-i(t\omega(x) + ta_\mu)} = e^{i\alpha_t(b)}e^{i\zeta(b)}\;,\\
    \alpha_t(b) \coloneqq \Phi\bigl(t\omega(x), \alpha(b), \zeta(b), ta_\mu\bigr)\;,
\end{gather*}
by the definition \eqref{eq:reduction-function}; the arguments have norm at most $\rho_0/2$ by hypothesis, and $\alpha_t$ has values in $\mathfrak g^c$ by (N2).

\emph{The plain differences of $\alpha$.} For $b, b + e$ in $Q$, with $b = \langle x, x + e_\mu\rangle$, the value of $U$ on the bond from $x$ to $x + e$ is $e^{\pm i\zeta}$ with the value of $\zeta$ on that bond, so that by \eqref{eq:covariant-difference}, (N1), and \eqref{eq:reduction-hypotheses}
\[
    \abs{\alpha(b + e) - \alpha(b)} \leq \abs{(\delta_e\alpha)(b)} + c_Gz\,\abs{\alpha(b + e)} \leq a' + c_Gza = a''\;.
\]
Two bonds of the direction $\mu$ in $Q$ are joined by a chain of at most $4D$ translations by unit vectors which stays in $Q$, one coordinate after the other, because $Q$ is a cube. Hence $\abs{\alpha(b) - a_\mu} \leq 4Da''$ for every bond $b$ of the direction $\mu$ in $Q$.

\emph{The bounds \eqref{eq:reduced-phase}.} By (N2), $\alpha_t(b) = \alpha(b) - ta_\mu + \Psi(t\omega(x), \alpha(b), \zeta(b), ta_\mu)$, and by the first line of \eqref{eq:reduction-remainder}, with $\abs{t\omega} \leq 4aD$, the last term is at most $r_\Phi$ in norm. For $t = 1$ the first two terms are at most $4Da''$ together, and for general $t$ at most $2a$.

\emph{The bound \eqref{eq:reduced-difference}.} By \eqref{eq:covariant-difference} and (N1), $\abs{(\delta_e\alpha_t)(b)} \leq \abs{\alpha_t(b + e) - \alpha_t(b)} + c_Gz(2a + r_\Phi)$. The arguments of $\Phi$ at $b + e$ are $t\omega(x) + ta_e$, with $a_e \coloneqq \pm a_\nu$ for $e = \pm e_\nu$, then $\alpha(b + e)$, $\zeta(b + e)$, and the same $ta_\mu$. The difference of the linear parts is $\alpha(b + e) - \alpha(b)$, at most $a''$, and by the second and third lines of \eqref{eq:reduction-remainder} the difference of the values of $\Psi$ is at most $c_G(2a + z)\abs{a_e} + c_G(4aD + 2a + z)a'' + c_G(4aD + 2a)z'$, the three terms being the changes of the first, the second, and the third argument, each multiplied by the bound of the corresponding derivative.
\end{proof}

\emph{One carried insertion.} Let $W$ be a group of the deviation species made at the step $j$. We fix a point of the space of the step $k + 1$,
\begin{equation}\label{eq:carried-datum}
    (\mathbf U, \mathbf J) = (e^{i\alpha}U, \mathbf J) \in \mathcal S_{k+1}(U)\;,
\end{equation}
around a real configuration $U$ which is either $U_{k+1}$, for a block field $V_{k+1}$ in the support, with the current $J \coloneqq J_{k+1}$, or the trivial configuration, with $J \coloneqq 0$; in both cases $(U, J) \in \mathcal S_{k+1}(U)$ and, by (b), $U$ is one of the configurations of \eqref{eq:production-space}. By Lemma~\ref{lem:space-inclusion} and (a), both $(\mathbf U, \mathbf J)$ and $(U, J)$ lie in $\mathcal D_W$. We write $a_{(n)} \coloneqq \alpha_{1,n}L^{-n}$ for the radius of the complex part at the step $n$, and
\begin{equation}\label{eq:carried-sizes}
    a \coloneqq a_{(k+1)}\;,
    \qquad
    \mathsf d_W \coloneqq 2u_\star M\sigma(W)L^j\;,
    \qquad
    \delta_W \coloneqq \mathsf d_WL^{-(k+1)} = 2u_\star ML^{-m}\sigma(W)\;.
\end{equation}
By (c) the $\ell^\infty$-diameter of $X_W$ is at most $u_\star\sigma(W)ML^j$, and that of the dependence set $X^+_W$ exceeds it by at most $16M \leq ML^j$; hence $X^+_W$ lies in a cube $Q_W$ of the fine lattice of side $\mathsf d_W$, and $\delta_W$ is its side in the units of the scale $k + 1$. The interval has at most $\mathsf d_W/2$ fine bonds, and $\delta_W \geq 2ML^{-m}$ because $\sigma(W) \geq 1$.

\begin{proposition}[One Carried Insertion]\label{prop:one-carried}
Let $\gamma$ be sufficiently small. With the notation above the following hold.
\begin{enumerate}[label=(\roman*)]
\item For every $W$,
\begin{equation}\label{eq:excess-plain}
    \norm{Y_W(\mathbf U, \mathbf J) - Y_W(U, J)} \leq c_{\mathrm{pl}}\,\varkappa_mL^{-m}\,\norm{Y_W}_{(j)}\;,
    \qquad
    c_{\mathrm{pl}} \coloneqq \max\Bigl(2,\ \frac{8c_G\lambda_{\mathrm r}}{c_FL}\Bigr)\;.
\end{equation}
\item Let $\mu_c \coloneqq \min\bigl(\frac12, c_FM/(c_G\lambda_{\mathrm r})\bigr)$. If $\delta_W \leq \mu_c$, then
\begin{equation}\label{eq:excess-covariant}
    \norm{Y_W(\mathbf U, \mathbf J)} \leq 2\norm{Y_W(U, J)} + \varkappa_mL^{-m}\Bigl(\frac{40\,\delta_W}{\mu_c} + 14L^{-2m}\Bigr)e^{\nu\sigma(W)/8192}\,\norm{Y_W}_{(j)}\;.
\end{equation}
\end{enumerate}
\end{proposition}

The first bound is the gain of the radii alone and holds for every group; the second has the further factor $\delta_W \leq 2u_\star ML^{-m}\sigma(W)$ and is the one which is summable over the generations.

\begin{proof}
We write $\mathsf r_n \coloneqq \alpha_{0,n}L^{-3n}$ for the radius of the current in \eqref{eq:space-of-step} and $a_2 \coloneqq \alpha_{1,k+1}L^{-2(k+1)} = aL^{-(k+1)}$ for that of the covariant differences at the step $k + 1$. By (d) and \eqref{eq:carried-blocking-condition},
\begin{equation}\label{eq:carried-ratios}
    \frac{a}{a_{(j)}} \leq \varkappa_mL^{-m} \leq \frac{1}{20}\;,
    \qquad
    \frac{a_2}{\alpha_{1,j}L^{-2j}} \leq \varkappa_mL^{-2m}\;,
    \qquad
    \frac{\mathsf r_{k+1}}{\mathsf r_j} \leq \varkappa_mL^{-3m} \leq \frac{1}{20L^2}\;.
\end{equation}
The curvature of the real configuration is $f_U \coloneqq \max_p\abs{U(\partial p) - 1} \leq 2B_3\varepsilon_{k+1}L^{-2(k+1)}$ by (b), and $f_U = 0$ for the trivial configuration. Since $\varepsilon_{k+1}/\alpha_{0,k+1}$ and $\alpha_{1,n}$ tend to zero with $\gamma$, uniformly in the steps (Appendix~\ref{app:constants}, $q_0 > p_0$), we may and do assume that
\begin{equation}\label{eq:carried-smallness}
    c_Gf_U \leq \tfrac14c_F\,\alpha_{0,j}L^{-2j}\;,
    \qquad
    c_G\,\alpha_{1,n}^2 \leq \tfrac14c_F\,\alpha_{0,n}\;,
    \qquad
    \alpha_{1,n} \leq \rho_0
\end{equation}
for all steps $n$; the first uses $\alpha_{0,k+1} \leq \varkappa_m\alpha_{0,j}$ and $\varkappa_mL^{-2m} \leq 1$.

\emph{(i).} For $j \geq 2$ put
\[
    R' \coloneqq \min\Bigl(\frac{a_{(j)}}{2a},\ \frac{c_F\,\alpha_{0,j}L^{-2j}}{8c_G\,a_2},\ \frac14\Bigl(\frac{\mathsf r_j}{\mathsf r_{k+1}} - 1\Bigr)\Bigr)\;,
\]
and for $j = 1$ let $R'$ be the first of the three numbers. If $2R' \leq 1$, the left side of \eqref{eq:excess-plain} is at most $2\norm{Y_W}_{(j)} \leq \norm{Y_W}_{(j)}/R'$, because both points lie in $\mathcal D_W$. Let $2R' > 1$. For $\abs z < 2R'$ the pair $(e^{iz\alpha}U, J + z(\mathbf J - J))$ lies in $\mathcal S_j(U)$: the complex part $z\alpha$ has norm less than $2R'a \leq a_{(j)}$ and covariant differences less than $2R'a_2 \leq a_{(j)}L^{-(k+1)} \leq \alpha_{1,j}L^{-2j}$; by \eqref{eq:complex-curvature} with $a_\circ = a_{(j)} \leq \rho_0$ and $d_\circ = 2R'a_2$, and by \eqref{eq:carried-smallness} and the second entry of $R'$, its curvature is at most $c_Gf_U + 2c_GR'a_2 + c_Ga_{(j)}^2 \leq \frac34c_F\alpha_{0,j}L^{-2j}$; and the current has norm less than $\mathsf r_{k+1} + 4R'\mathsf r_{k+1} \leq \mathsf r_j$, by the third entry. The function $g(z) \coloneqq Y_W(e^{iz\alpha}U, J + z(\mathbf J - J)) - Y_W(U, J)$ is therefore holomorphic on the disc of radius $2R' > 1$, it is at most $2\norm{Y_W}_{(j)}$ in norm there, and $g(0) = 0$, so that by Schwarz's lemma, applied to the matrix elements, $\norm{g(1)} \leq \norm{Y_W}_{(j)}/R'$. In both cases the left side of \eqref{eq:excess-plain} is at most $\norm{Y_W}_{(j)}/R'$, and by \eqref{eq:carried-ratios} the reciprocals of the three entries of $R'$ are at most $2\varkappa_mL^{-m}$, $(8c_G\lambda_{\mathrm r}/c_F)\varkappa_mL^{-2m}$, and $4 \cdot \frac{20}{19}\varkappa_mL^{-3m}$, each of which is at most $c_{\mathrm{pl}}\varkappa_mL^{-m}$.

\emph{(ii), the gauge.} The conditions \eqref{eq:space-of-step} are invariant under the $G$-valued gauge transformations, which conjugate $\alpha$ and $\mathbf J$, and so are the norms in \eqref{eq:excess-covariant}, by \eqref{eq:covariance-law} and because $\pi$ is unitary on $G$. For $U = U_{k+1}$ we apply Supplement, Lemma~17.22, to a cube $\mathcal Q$ of $2^4$ cubes of $\pi_{k+1}$ which contains $Q_W$; such a cube exists because $\delta_W \leq \frac12$ is less than the side $M$ of these cubes. Its hypotheses hold for $\gamma$ small, because the local curvature is $E \leq 2B_3\varepsilon_{k+1}$ by (b) and $M_{\mathcal Q} \leq 4M$. It gives a $G$-valued gauge transformation after which $U = e^{i\zeta}$ on $Q_W$ with, in fine units,
\[
    \abs{\zeta(b)} \leq z \coloneqq \hat bL^{-(k+1)}\;,
    \qquad
    \abs{\zeta(b + e) - \zeta(b)} \leq z' \coloneqq \hat bL^{-2(k+1)}\;,
    \qquad
    \hat b \coloneqq 24B_3^2M\varepsilon_{k+1}\;,
\]
and for the trivial configuration $\zeta = 0$. We work in this gauge and apply Lemma~\ref{lem:gauge-reduction} to the cube $Q_W$, with $D = \mathsf d_W$, $x_0 = x_W$, the radii $a$ and $a' = a_2$ of \eqref{eq:carried-datum}, and these $z$, $z'$; its hypotheses hold because $4a\mathsf d_W = 4\alpha_{1,k+1}\delta_W \leq 2\alpha_{1,k+1}$ and $z$ are small. In the bounds of that lemma, $a_2\mathsf d_W = a\delta_W$, $za\mathsf d_W = \hat b\,a\delta_W$, $a^2 = \alpha_{1,k+1}L^{-(k+1)}a \leq \alpha_{1,k+1}a\delta_W$, and $za = \hat b\,a_2$, so that every term is a multiple of $a\delta_W$, or of $a$, or of $a_2$, with a coefficient which is $4$, respectively $2$, respectively $1$, up to terms which vanish with $\alpha_{1,k+1}$ and $\hat b$. For $\gamma$ small we therefore have, for $t \in [0, 1]$ and on the bonds of $Q_W$,
\begin{equation}\label{eq:reduced-datum}
    \abs{\alpha_1} \leq 5a\delta_W\;,
    \qquad
    \abs{\alpha_t} \leq 3a\;,
    \qquad
    \abs{\delta_e\alpha_t} \leq 2a_2\;,
    \qquad
    \abs\omega \leq 4\alpha_{1,k+1}\delta_W\;.
\end{equation}

\emph{(ii), the path.} We claim that $(\mathbf U^{u_t}, R(u_t)\mathbf J) \in \mathcal S_j(U)|_{X^+_W}$ for $t \in [0, 1]$. By Lemma~\ref{lem:gauge-reduction}, $\mathbf U^{u_t} = e^{i\alpha_t}U$, and by \eqref{eq:reduced-datum} and \eqref{eq:carried-ratios} the complex part satisfies $\abs{\alpha_t} \leq 3a < a_{(j)}$ and $\abs{\delta_e\alpha_t} \leq 2a_2 \leq 2\varkappa_mL^{-2m}\alpha_{1,j}L^{-2j} < \alpha_{1,j}L^{-2j}$. For the curvature we use that $\delta_W \leq \mu_c$ and $\delta_W \geq 2ML^{-m}$ give $L^{-m} \leq \mu_c/(2M)$, so that $2a_2 \leq 2\varkappa_mL^{-m} \cdot L^{-m}\lambda_{\mathrm r}\alpha_{0,j}L^{-2j} \leq (\mu_c\lambda_{\mathrm r}/(20M))\alpha_{0,j}L^{-2j} \leq (c_F/(20c_G))\alpha_{0,j}L^{-2j}$ by the definition of $\mu_c$; with \eqref{eq:complex-curvature} for $a_\circ = 3a \leq a_{(j)}$ and \eqref{eq:carried-smallness} the curvature is less than $c_F\alpha_{0,j}L^{-2j}$. Finally $\abs{R(u_t)\mathbf J} \leq (1 + c_G\abs\omega)\abs{\mathbf J} \leq 2\mathsf r_{k+1} < \mathsf r_j$ by (N1). For $j = 1$ only the first of these conditions is posed. Since $u_1(x_W) = 1$, \eqref{eq:base-point-invariance} and \eqref{eq:base-point-form} give
\begin{equation}\label{eq:carried-identity}
    Y_W(\mathbf U, \mathbf J) = \hat Y_W(\mathbf U^{u_1}, R(u_1)\mathbf J)\,\pi\bigl(\mathbf U(I_W)\bigr)\;.
\end{equation}

\emph{(ii), the disc.} Put
\[
    R_A \coloneqq \frac{\mu_c\,a_{(j)}}{10\,a\delta_W}\;,
    \qquad
    R_J \coloneqq \frac13\Bigl(\frac{\mathsf r_j}{\mathsf r_{k+1}} - 1\Bigr)\;,
    \qquad
    R \coloneqq \min(R_A, R_J)\;,
\]
with $R \coloneqq R_A$ for $j = 1$, and $f(z) \coloneqq \hat Y_W(e^{iz\alpha_1}U, J + z(R(u_1)\mathbf J - J))$ for $\abs z < R$, so that $f(1)$ is the first factor of \eqref{eq:carried-identity} and $f(0) = \hat Y_W(U, J)$. Since $\delta_W \leq \mu_c$, we have $R_A \geq a_{(j)}/(10a) \geq 2$ and $R_J \geq 2$ by \eqref{eq:carried-ratios}. For $\abs z < R$ the pair in the argument of $f$ lies in $\mathcal S_j(U)|_{X^+_W}$: by \eqref{eq:reduced-datum} the complex part $z\alpha_1$ has norm less than $5R_Aa\delta_W = \mu_ca_{(j)}/2$; its covariant differences are less than $2R_Aa_2 = \mu_ca_{(j)}L^{-(k+1)}/(5\delta_W) \leq \mu_c\,\alpha_{1,j}L^{-2j}/(10M)$, because $\delta_W \geq 2ML^{-m}$; by \eqref{eq:complex-curvature}, \eqref{eq:carried-smallness}, and the definition of $\mu_c$ its curvature is at most
\[
    c_Gf_U + \frac{c_G\mu_c\lambda_{\mathrm r}}{10M}\alpha_{0,j}L^{-2j} + c_Ga_{(j)}^2 \leq \Bigl(\frac14 + \frac{1}{10} + \frac14\Bigr)c_F\,\alpha_{0,j}L^{-2j}\;;
\]
and the current has norm less than $\mathsf r_{k+1} + 3R_J\mathsf r_{k+1} = \mathsf r_j$, because $\abs{R(u_1)\mathbf J} \leq 2\mathsf r_{k+1}$. Hence $f$ is holomorphic on the disc of radius $R$. Along the interval the transport of $e^{iz\alpha_1}U$ is a product of at most $\mathsf d_W/2$ factors $e^{iz\alpha_1(b)}U(b)$, with $\pi(U(b))$ unitary and $\norm{\pi(e^{iX})} \leq e^{c'_\pi\abs X}$ (Section~\ref{sec:localisation-expansion}), so that by \eqref{eq:base-point-form} and (a)
\begin{align*}
    \norm{f(z)} &\leq \norm{Y_W}_{(j)}\exp\bigl(c'_\pi\tfrac12\mu_ca_{(j)} \cdot \tfrac12\mathsf d_W\bigr) = \norm{Y_W}_{(j)}\exp\bigl(\tfrac12c'_\pi\mu_cu_\star M\alpha_{1,j}\,\sigma(W)\bigr)\\
    &\leq e^{\nu\sigma(W)/8192}\,\norm{Y_W}_{(j)}\;,
\end{align*}
the last for $\gamma$ so small that $4096\,c'_\pi u_\star M\alpha_{1,n} \leq \nu$ for all $n$. By Schwarz's lemma for $f - f(0)$, which is at most twice this bound on the disc, $\norm{f(1) - f(0)} \leq (2/R)e^{\nu\sigma(W)/8192}\norm{Y_W}_{(j)}$. In the same way $\norm{\pi(\mathbf U(I_W))} \leq \exp(c'_\pi a\mathsf d_W/2) = \exp(\frac12c'_\pi\alpha_{1,k+1}\delta_W) \leq 2$, and $\norm{f(0)} = \norm{Y_W(U, J)}$ because $\pi(U(I_W))$ is unitary. By \eqref{eq:carried-identity},
\[
    \norm{Y_W(\mathbf U, \mathbf J)} \leq 2\bigl(\norm{f(0)} + \norm{f(1) - f(0)}\bigr) \leq 2\norm{Y_W(U, J)} + \frac4R\,e^{\nu\sigma(W)/8192}\norm{Y_W}_{(j)}\;,
\]
and $4/R \leq 4/R_A + 4/R_J \leq 40\varkappa_mL^{-m}\delta_W/\mu_c + 12 \cdot \frac{20}{19}\varkappa_mL^{-3m}$ by \eqref{eq:carried-ratios}, which is \eqref{eq:excess-covariant}.
\end{proof}

\emph{The sum over the generations.} It remains to sum the bounds of the proposition over the groups at a traversal.

\begin{proposition}[The Carried Groups in the Supremum Clause]\label{prop:carried-sup}
Let the ensemble be dressed at the steps $1, \dots, k$, and suppose that on the term without large-field regions of the step $k + 1$ the groups carried from these steps satisfy \eqref{eq:per-insertion} and \eqref{eq:tail-clause} at the step $k + 1$, at every real background of the term, and vanish at its reference data. Let $\gamma$ be sufficiently small and $L > (1 + \beta_0)^6$. Then, for every traversal $\mathfrak n'$ of the step $k + 1$,
\begin{equation}\label{eq:carried-sup}
    \sum_{W \ni \mathfrak n'}e^{\theta(W) + \lambda_c\nu\sigma(W) + \nu_e\operatorname{ext}(W)}\norm{Y_W}_{(k+1)} \leq 2C_\chi\tilde\rho_{k+1}\bar\delta^*_{\mathfrak n'} + 2\hat C\,C_{S'}\,\rho_{k+1}\;,
\end{equation}
the sum over the groups of the deviation species made at the steps $j \leq k$ with $\mathfrak n'$ in the interval, where $\bar\delta^*_{\mathfrak n'}$ is the expression \eqref{eq:deviation-quantity} of the step $k + 1$ with each maximum replaced by its supremum over the real backgrounds of the term, and
\begin{equation}\label{eq:generation-constant}
    \hat C \coloneqq \sum_{m \geq 1}(1 + \beta_0)^{6m}\,\varkappa_m\Bigl[L^{-m}\Bigl(\frac{1280\,u_\star M}{e\,\nu\mu_c} + 14L^{-m}\Bigr) + c_{\mathrm{pl}}\exp\Bigl(-\frac{\nu\mu_cL^m}{32u_\star M}\Bigr)\Bigr]\;.
\end{equation}
Here $\theta(W)$, $\sigma(W)$, and $\operatorname{ext}(W)$ are the numbers fixed when $W$ was made, as in \eqref{eq:production-clause}.
\end{proposition}

The constant $\hat C$ is finite because $L > (1 + \beta_0)^6$, and it depends on $L$, $M$, $\nu$, $\beta_0$, $u_\star$, $c_F$, $c_G$, and $\lambda_{\mathrm r}$, and not on $k$ or on the coupling. The first term of \eqref{eq:carried-sup} is at most $2C_\chi\tilde\rho_{k+1} \cdot 2c_\sigma B_3\varepsilon_{k+1}W_1$, because the bound \eqref{eq:deviation-scaling} holds for each maximum and at every real background, and this is less than $\rho_{k+1}$ for $\gamma$ small, since $\tilde\rho_{k+1}\varepsilon_{k+1} = \rho_{k+1}\varepsilon_{k+1}/\min(\alpha_{0,k+1}, \alpha_{1,k+1})$ and $q_0 > p_0$. With \eqref{eq:production-clause} at the step $k + 1$, in which $\ell_W \geq 1$, $S_W \geq \norm{Y_W}_{(k+1)}$, and $\lambda_c \leq 1$, the proposition therefore gives \eqref{eq:sup-clause} at the step $k + 1$ with
\begin{equation}\label{eq:sup-constant}
    C_c \coloneqq 1 + 2(1 + \hat C)\,C_{S'}\;.
\end{equation}
The constant does not come from a recursion: the bound of the carried groups uses the production bounds of the earlier steps and no supremum clause.

\begin{proof}
\emph{The real points.} Let $W$ be a group of the sum, made at the step $j$, and \eqref{eq:carried-datum} a point of the space of the step $k + 1$, around $U$. We apply Proposition~\ref{prop:one-carried}(ii) if $\delta_W \leq \mu_c$ and (i) otherwise, and in both cases
\[
    \norm{Y_W(\mathbf U, \mathbf J)} \leq 2\norm{Y_W(U, J)} + \epsilon_W\;,
\]
with $\epsilon_W$ the second term of \eqref{eq:excess-covariant}, respectively the right side of \eqref{eq:excess-plain}; it does not depend on the point. If $U$ is the trivial configuration, the first term vanishes by hypothesis. If $U = U_{k+1}$, let $\mathfrak n_j \subset \mathfrak n'$ be a traversal of the step $j$ in the interval of $W$ and $t \coloneqq t^{(k+1)}_{\mathfrak n'}(W)$. As in the proof of Lemma~\ref{lem:tail-to-sum}, the comparison region of $W$ at $\mathfrak n_j$ lies within the distance $R_t$ of $\mathfrak n'$ in the units of the scale $k + 1$, so that by \eqref{eq:per-insertion} and \eqref{eq:insertion-deviation} $e^{\theta(W)}\norm{Y_W(U, J)} \leq \beta_W\bar m_t$, where $\bar m_t$ is the supremum over the real backgrounds of the maximum of $\abs{U(\partial p) - 1}$ over $p \in \mathcal P_t(\mathfrak n')$. Taking the supremum over the points we obtain
\begin{equation}\label{eq:carried-pointwise}
    e^{\theta(W)}\norm{Y_W}_{(k+1)} \leq 2\beta_W\,\bar m_{t^{(k+1)}_{\mathfrak n'}(W)} + e^{\theta(W)}\epsilon_W\;.
\end{equation}
The numbers $\bar m_t$ are nondecreasing in $t$, and the summation by parts of the proof of Lemma~\ref{lem:tail-to-sum}, with \eqref{eq:tail-clause} at the step $k + 1$ and $\lambda_c \leq 1$, bounds the weighted sum of the first terms of \eqref{eq:carried-pointwise} over the groups of the sum by $2C_\chi\tilde\rho_{k+1}\bar\delta^*_{\mathfrak n'}$.

\emph{The excess of one generation.} Fix $j \leq k$ and consider the groups of the sum which were made at the step $j$. Each of them has a traversal of the step $j$ in its interval which is contained in $\mathfrak n'$, namely the one which contains a fine bond common to $\mathfrak n'$ and the interval, and $\mathfrak n'$ contains $L^m$ traversals of the step $j$. By \eqref{eq:production-clause} at the step $j$, with $\ell_W \geq 1$ and $S_W \geq \norm{Y_W}_{(j)}$,
\begin{equation}\label{eq:generation-production}
    \sum_We^{\theta(W) + \nu\sigma(W) + \nu_e\operatorname{ext}(W)}\norm{Y_W}_{(j)} \leq 2L^mC_{S'}\,\rho_j\;,
\end{equation}
the sum over these groups. For those with $\delta_W \leq \mu_c$ we have $\lambda_c\nu + \nu/8192 = \frac{15}{16}\nu$ and, by \eqref{eq:carried-sizes} and $te^{-\nu t/16} \leq 16/(e\nu)$,
\begin{align*}
    e^{\lambda_c\nu\sigma(W)}\epsilon_W &\leq e^{\nu\sigma(W)}\varkappa_mL^{-m}\Bigl(\frac{80u_\star ML^{-m}}{\mu_c}\sigma(W)e^{-\nu\sigma(W)/16} + 14L^{-2m}\Bigr)\norm{Y_W}_{(j)}\\
    &\leq e^{\nu\sigma(W)}\varkappa_mL^{-2m}\Bigl(\frac{1280\,u_\star M}{e\,\nu\mu_c} + 14L^{-m}\Bigr)\norm{Y_W}_{(j)}\;.
\end{align*}
For those with $\delta_W > \mu_c$ we have $\sigma(W) > \mu_cL^m/(2u_\star M)$, so that $e^{\lambda_c\nu\sigma(W)} \leq e^{\nu\sigma(W)}e^{-\nu\sigma(W)/16} \leq e^{\nu\sigma(W)}\exp(-\nu\mu_cL^m/(32u_\star M))$ and
\[
    e^{\lambda_c\nu\sigma(W)}\epsilon_W \leq e^{\nu\sigma(W)}c_{\mathrm{pl}}\,\varkappa_mL^{-m}\exp\Bigl(-\frac{\nu\mu_cL^m}{32u_\star M}\Bigr)\norm{Y_W}_{(j)}\;.
\]
Multiplying by $e^{\theta(W) + \nu_e\operatorname{ext}(W)}$ and summing with \eqref{eq:generation-production}, we find that the weighted sum of the $\epsilon_W$ over the groups made at the step $j$ is at most the $m$-th term of \eqref{eq:generation-constant} without the factor $(1 + \beta_0)^{6m}$, times $2C_{S'}\rho_j$, and $\rho_j \leq (1 + \beta_0)^{6m}\rho_{k+1}$ by \eqref{eq:rho-ratio}.

\emph{The sum.} Summing the excesses over $j \leq k$, that is, over $m \geq 1$, and adding the first terms of \eqref{eq:carried-pointwise}, we obtain \eqref{eq:carried-sup}.
\end{proof}

\subsection{Large Fields and the \texorpdfstring{$\mathbf R$}{R}-Operation}\label{sec:large-field-step}

The insertions attached to a large-field region carry no deviation, and they cannot be contracted: a segment of the step $n$ which runs parallel to the large-field data of the step $j < n$ contains $L^{n-j}$ traversals of the step $j$, each carrying an attached part of size $C\hat\rho_j$. They are therefore not kept in the per-traversal bound beyond the step after the one at which they were made. They are carried, with the species of their production, as functions which are never summed over the loop, and their total over each component $Z$ of the large-field region of the step $j$ is paid once, by the small factor which Ba{\l}aban's large-field estimates attach to $Z$.

\emph{The payment.} Ba{\l}aban attaches to every large-field region $Z$ created at the step $j$ the factor $\exp(-\frac12\gamma_0A_1^2p_0^2(g_j)(d'_j(Z) + 1) - 2p_0(g_j))$, with $d'_j(Z)$ the linear size of $Z$, ``the fundamental large field factors, which control convergence of the expansion of the effective densities'' \cite[p.~381]{BalabanLFR2}, and he shows that its part $\exp(-\kappa_j(Z))$ ``controls $K$ renormalization steps, under the assumption that no large fields are created in these steps'', in the sense of the inequality \cite[(1.80), p.~384]{BalabanLFR2}. Two costs have to be paid from it. The first is the total of the attached insertions,
\begin{equation}\label{eq:attached-total}
    A_j(Z) \leq C\hat\rho_jN_j(Z) + C\rho_{t,j}\mathfrak m_j(Z) \leq C_A(d'_j(Z) + 1)\;,
\end{equation}
where $N_j(Z) = \sum_{\mathfrak n_j}e^{-c_{\mathrm{att}}d^Z_j(\mathfrak n_j)/M}$ counts the traversals of the step $j$ near $Z$, which by Lemma~\ref{lem:square-properties}(i) are at most $C\Lambda(C)$ per cube of $Z$, and $\mathfrak m_j(Z)$ is the sum of the marks of the data assigned to $Z$, at most a multiple of the number of cubes of $Z$. The second cost recurs at every step $n$ of the lifetime of $Z$: the traversals near $Z$ at the step $n$ are exposed, the expansion of the step is not available on them, and each costs $C_e\hat\tau_n^2$, $\hat\tau_n^2 \coloneqq \max(\tau_n^2, C\hat\rho_n)$, with a constant $C_e$ of the construction. By Lemma~\ref{lem:square-properties}(i) a cube of side $MR_n$ meets at most $4\Lambda(C)MR_n$ traversals, and along Ba{\l}aban's chain for the regions $S^{n-j}(Z)$ which descend from $Z$ \cite[pp.~384--385]{BalabanLFR2}, whose numbers of cubes are at most $C2^{-(n-j)}d'_j(Z)$ until the region fits into a cube of size $64MR_n$, after which at most $R_j$ steps remain, each with a region of at most $64^d$ cubes,
\begin{equation}\label{eq:lifetime-cost}
    D_j(Z) \coloneqq \sum_n4C_e\Lambda(C)MR_n\hat\tau_n^2\,\#\mathrm{cubes}_n\bigl(S^{n-j}(Z)\bigr) \leq C_d\Lambda(C)MLR_j^2\,\bar\tau^2\,(d'_j(Z) + 1)\;,
\end{equation}
with $\bar\tau^2 \coloneqq \sup_{n \leq K'}\hat\tau_n^2 \leq C\gamma_q^2(\log\gamma_q^{-2})^{6p_0}$, by $g_n \leq \gamma_q$ for $n \leq K'$ from \eqref{eq:stopping-coupling}, the sum running over the steps $n \leq K'$ at which the loops are carried, and $R_n \leq LR_j$ for $n \geq j$ by Lemma~\ref{lem:tube-cubes}. Both costs are at most a polylogarithm of $g_j^{-2}$ of fixed degree times $d'_j(Z) + 1$, while $\kappa_j(Z)$ is of order $(\log g_j^{-2})^{2p_0}(d'_j(Z) + 1)$, so that for $p_0$ large and $\gamma_q$ small, depending on $c_{\mathcal C}$, the factor of $Z$ pays both and leaves a reserve, since $j \leq K'$ and $\log g_j^{-2} \geq \log\gamma_q^{-2}$; Ba{\l}aban's induction over the merging of regions \cite[pp.~384--389]{BalabanLFR2} runs with the summand of \cite[(1.80)]{BalabanLFR2} enlarged by these costs, since both are subadditive under the union of regions up to the slack $2d$ of \cite[p.~386]{BalabanLFR2}. That the attached parts are not increased by the later evaluations at $U_{n+1}$, by the fluctuation parts, and by the merging of several parents into one insertion, and that the parts of an attached insertion which reach the region of small fields are paid by the reserve of their component, are proved in Supplement, Proposition~S.9.

\emph{The $\mathbf R$-operation.} From the step at which the depth condition (N1) first admits a component, the $\mathbf R$-operation acts on the numerator. It integrates a term over the block field on a domain inside a small and old large-field component and replaces it by the normalised fraction of \cite[(1.100), p.~201]{BalabanLFR1}; since the fraction integrates to one over the variables which it replaces, the operation has the normalisation property \cite[(1.102), p.~201]{BalabanLFR1}, $\int(\mathbf R'f)\,\dd V_k = \int f\,\dd V_k$, for every integrable $f$ and not only for the density; the removal of the $\delta$-functions on the components of the first class \cite[p.~378]{BalabanLFR2} is an equivalence which needs the invariance of the integrand under gauge transformations, and the dressed loop is gauge invariant. On the components of the first class the term is expanded by Ba{\l}aban's localisation expansion \cite[pp.~372--376, (1.66)--(1.67)]{BalabanLFR2}, which he makes for ``any term in the sum'' and which applies to the words of the loop, entire functions of their arguments. The part of the resulting polymer gas which is linked to the loop is written exactly as a ratio of two polymer gases, and at flat reference data each group of polymers which interleaves with the loop contracts to a scalar by Schur's lemma, as in Section~\ref{sec:small-field-map}. The flat part of the ratio is therefore the cluster series of a one-dimensional gas along each loop; its perimeter and corner parts are added to the step factor $\kappa_k$ (Definition~\ref{def:step-factor}), and its remainder $E^{\circ,R}_k$, bounded in Proposition~\ref{prop:r-flat-part}, stays in $\mathcal E_k$. The rest is a sum of deviation insertions of the class, each vanishing at pure-gauge data on its hull, with a source of size at most $C^{\mathbf R}(g_k) \coloneqq \exp(-A_0(\log g_k^{-2})^{p_0}/4 + \log g_k^{-2} + C\log\log g_k^{-2})$, which is added to $C_{\mathrm{new}}$ in \eqref{eq:deviation-recursion}, and the groups of polymers which contain insertions of the term add $C_aMR_ke^{-p_0(g_k)/4}$ to the coefficient $a$ there. The details, including the depth $N_k$ of Definition~\ref{def:depth} and the commensurability of the loops with the period $MR_k$ of the gas, are Supplement, Proposition~S.10.

\subsection{Several Loops, the Torus, and the Assembly}\label{sec:assembly}

\emph{Cross groups.} A group whose intervals lie on two or more loops is a cross group. By Lemma~\ref{lem:stopping-index} the loops are at distance at least $a_{k+1} \geq L^2MR_{k+1}$ in units of the scale $k + 1$, so that a set of diameter less than $a_{k+1}$ meets at most one loop, and a cross group contains a connected chain of domains joining two loops, of total size at least $a_{k+1}/M$. Since the flat value of a cross group is not a scalar (Section~\ref{sec:small-field-map}), its insertion is neither split nor contracted, and three things have to be shown for the cross species of Definition~\ref{def:dressed-loop}(ii): that little of it is made at each step, that what has been made grows slowly, and that the total stays bounded.

\emph{The cross species: production.} Let $W$ be a group of the cross species made at the step $k + 1$, and $\mathfrak n$ a traversal in one of its intervals. Its network is connected and contains $\mathfrak n$ and a traversal of another loop, so that it extends from $\mathfrak n$ over at least $\delta_{\mathfrak n}/M$ cubes of $\pi_{k+1}$, where $\delta_{\mathfrak n} \geq d_{k+1}$ is the distance from $\mathfrak n$ to the other loops. The insertion of $W$ is a sum of integrals of ordered products, along each of its intervals, of the tensors $Y_G$ of Lemma~\ref{lem:word-expansion}, old insertions, and transports, and the size bound \eqref{eq:size-bound} is the bound on the corresponding sums and integrals of the products of the norms; it therefore holds for the cross groups in the norm \eqref{eq:projective-norm}. Beyond the weight $e^{\nu\sigma}$ the activities of the polymers of the step carry the factor $e^{-\nu_e}$ for every cube of $\pi_{k+1}$ of their extent, set aside from the rate of the polymers (Supplement, proof of Proposition~S.9(ii)), and $\nu_\times \leq \nu_e$. Rooting each group at the first traversal of its intervals and spending the factor of the extent on the distance, we obtain
\begin{equation}\label{eq:cross-production}
    \sum_We^{\theta(W) + \nu\sigma(W)}\norm{Y_W} \leq C_{\mathrm{cr}}\mathfrak a_{k+1}\sum_{\mathfrak n}e^{-\nu_\times\delta_{\mathfrak n}/M} \leq C_{\mathrm{cr}}\mathfrak a_{k+1}P_{k+1}(\mathcal C)e^{-\nu_\times d_{k+1}/M}\;,
\end{equation}
the first sum over the groups of the cross species which the step map makes at the step $k + 1$ and the second over the traversals of that step. This is the first term of $\mathsf A^\times_{k+1}$ in \eqref{eq:cross-bound}, and at the step $1$ it is the bound of Section~\ref{sec:first-step-proof}. At an $\mathbf R$-step the cross groups which the operation makes add the second term (Supplement, Proposition~S.10(v)).

\emph{The cross species: carrying.} At a later step $n + 1$ an insertion of the cross species is evaluated at the points of the expansion of the step, where its bound, a supremum over the polydisc of its own step, does not grow, and neither do its transports. The parts of it which the step integrates or merges into new groups have weighted mass at most $2\bar c$ times its own, with the constant $\bar c$ of Supplement, Proposition~S.9, and $2\bar c \leq 1$ for $\gamma_q$ small. No distance is gained here, in contrast to the attached insertions: the loops lie in the small-field region on most terms, and the fluctuation field of every later step reaches every insertion. The mass of the cross species of the step $j$ is therefore multiplied by at most $1 + 2\bar c \leq 2$ at each later step, and at an $\mathbf R$-step by a further factor $1 + y^{\mathbf R}_n$, where $\prod_n(1 + y^{\mathbf R}_n) \leq 2$ (Supplement, Propositions~S.9(v) and~S.10(v)). With \eqref{eq:cross-production} this gives \eqref{eq:cross-bound} at the step $k + 1$. The growth by $2^{k-j}$ is paid at the stopping step by the decay of $\mathsf A^\times_j$ in the distance $d_j = L^{K'-j}d_{K'}$, which is exponential in $L^{K'-j}$ (Section~\ref{sec:uniform-bound}).

\emph{The cross exponent.} The cross exponent $E^\times_{k+1}$, the difference between the logarithm of the reference gas of the ensemble and the sum of those of its loops, is a sum over the traversals of terms which vanish unless a cluster reaches from one loop to another, and
\begin{equation}\label{eq:cross-exponent}
    \abs{E^\times_{k+1}} \leq C_{\mathrm{cr}}\mathfrak a_{k+1}\sum_{\mathfrak n}e^{-\nu_\times\delta_{\mathfrak n}/M} \leq C_{\mathrm{cr}}\mathfrak a_{k+1}P_k(\mathcal C)e^{-\nu_\times a_{k+1}/M}\;,
\end{equation}
the sum over the traversals of the step $k$, with the amplitude $\mathfrak a_{k+1}$ of \eqref{eq:amplitude} and the constants $C_{\mathrm{cr}}$, $\nu_\times$ of Appendix~\ref{app:constants}. The right side is summable over $k$, because $a_{k+1} = L^{j_q-q}L^{K'-k-1}$ grows geometrically as $k$ decreases. The flat values are taken loop by loop, and a cross group keeps its insertion; no count depends on the number of loops, since every count is per traversal or per cube, and a cube meets one loop.

\emph{The torus and the wrapping.} The remainder of the flat part contains a term from the periodicity of the torus, of size $C\mathfrak a_{k+1}P_k(\mathcal C)e^{-c\iota_{T,k+1}}$ with $\iota_{T,k} = L\iota_{T,k+1}$, and the insertions whose domains do not lift to the universal cover give the last term of \eqref{eq:insertion-bound}, which grows by at most a factor $L$ per step. Both are summable over the steps: with $x \coloneqq L^{K-j} \geq 1$ and $a \coloneqq c'T/(uM) \geq 1$, which is a lower bound on the period, each is at most $Cx\,e^{-ax}$ and
\begin{equation}\label{eq:wrap-sum}
    \sum_{i \geq 0}L^ie^{-aL^i} \leq e^{-a}\sum_{i \geq 0}L^ie^{-(L^i - 1)} \leq 2e^{-a}\;,
\end{equation}
because $L^ie^{-aL^i} = e^{-a}L^ie^{-a(L^i - 1)} \leq e^{-a}L^ie^{-(L^i - 1)}$ for $a \geq 1$, and the terms of the last series, which decrease in $L$, are at most $1$, $0.74$, $0.20$, $0.01$, and so on at $L = 2$.

\emph{Assembly.} The output of the step on every labelled term of $\varrho_{k+1}$ is a dressed loop of the step $k + 1$: the normalisation is $\prod_i\kappa^{(k)}(C_i, \chi_i)\kappa_{k+1}(C_i, \chi_i)$ by Section~\ref{sec:flat-part}; the leading words are the terms $\mathfrak Y = \emptyset$ of \eqref{eq:word-expansion} for the empty family; the insertions are the old ones carried as above, the fresh ones, the deviation parts of the new groups, the attached ones, and those of the cross species; the interactions are the resummation of the families against the reference gas of the step $k + 1$; and the scalar exponent is $\mathcal E_k$ plus the remainder $E^\circ_{k+1}$ of the flat part, the cross exponent, and the local terms of the step. The clauses (i)--(vi) of Definition~\ref{def:dressed-loop} hold with the same constants by \eqref{eq:recursion-fixed-point}, by the payment of the attached parts, and by Section~\ref{sec:flat-part}, which bounds the remainder and carries the torus consistency (vi) from the step $k$ to the step $k + 1$. This proves Theorem~\ref{thm:inductive-step}, given Supplement, Propositions~S.8--S.10 and Section~\ref{sec:flat-part}.

\section{The Flat Part and the Normalisation}\label{sec:flat-part}

We now determine the number which a step contributes to each loop, and this is the place at which the length of the loops has to be faced directly. At the trivial background the insertions of the step $k + 1$ on one loop are scalars, by Schur's lemma, and the contribution of a loop $C$ becomes the partition function of a hard-core gas of groups on $C$, whose logarithm is a sum over clusters. Every traversal carries clusters of total weight of order $\mathfrak a_{k+1}$, so that the logarithm is of order $P_k(C)\mathfrak a_{k+1}$, while $P_k(C)$ grows like $L^{K-k}$. This number can be neither bounded nor neglected, and it has to be extracted exactly, by a formula whose coefficients do not depend on the loop. Two properties of the gas make this possible: a cluster sees the loop only through the arcs of its groups, and the gas is invariant under the translations which preserve the partitions of the construction. Consequently a cluster which is small compared with the corner isolation cannot tell a loop from an infinite straight line away from the corners, nor from two perpendicular rays near a corner, and the density of the logarithm along the loop is the periodic density of the line or of the bent path, up to errors which decay exponentially in the corner isolation. The coefficients of the normalisation are the sums of these densities over one period, and the excess of the bent path over its two rays. Since the line and the bent path are infinite, these sums are defined on $\mathbb Z^4$, as limits over larger and larger tori, and the third property which we need is that the weights on a torus approach their limits exponentially fast in the period.

\subsection{The Reference Gas}\label{sec:reference-gas}

We fix a step $k + 1 \leq K'$ with $k \geq 1$, the step $1$ being that of Section~\ref{sec:first-step}, and a labelled term $\ell$ of $\varrho_{k+1}$. The \emph{trivial term} $\ell_0$ is the term with $\Omega_j = \Lambda_j = \mathbb T$ for every $j \leq k + 1$. It has no large-field regions and no outer variables, and its reference data, the trivial field, lie in the support of all its test functions, so that the identities of Section~\ref{sec:inductive-class} hold there. The traversals in this section are those of the step $k$, and lengths and distances are measured in units of the scale $k$ unless the contrary is said.

\emph{The partitions and the period.} The cubes of the partitions used by the construction have sides which are powers of $L$, and these partitions are compatible with each other \cite[p.~245]{BalabanCRE}. That $R_j$, $M$, and $M_2$ are powers of $L$ is Ba{\l}aban's choice \cite[p.~245]{BalabanCRE}, and that $M_1$ is one is a standing condition of \cite[(3.1)]{WilsonFirstStep}. Since every side $s$ is odd, the cubes of a partition can be centred at sites, and we take every partition to consist of the cubes $sz + \{-\frac{s-1}{2}, \dots, \frac{s-1}{2}\}^4$, $z \in \mathbb Z^4$, of sites of the scale at which it is used, so that the origin is a site of every lattice and the centre of a cube of every partition; this is compatible with the blocks of the averages, which are centred at the sites of the coarser lattice \cite[(0.3), p.~252]{BalabanRG1}. For the families $\pi_j$ it is Ba{\l}aban's own convention, whose cubes have their centres at the points of a lattice of spacing $M$ \cite[p.~257]{BalabanRG1}, and for the other partitions, which \cite[p.~245]{BalabanCRE} requires only to be compatible, it is our choice. Let $\Pi_{k+1}$ be the largest side, in units of the scale $k$, of the cubes used at the steps $1, \dots, k + 1$: the cubes of the families $\pi_j$, the cubes of side $MR_j$ of which the domains of the step $j$ are unions, the cubes of side $L^2M_2R_j$ of the test functions \cite[(3.2), (3.3), (3.16)]{BalabanCRE}, and the cubes of side $R_1M_1$ of the decoupling \cite[p.~3]{BalabanRG2}. Every translation by a vector of $\Pi_{k+1}\mathbb Z^4$ maps each of these partitions onto itself.

\begin{lemma}[Period Points]\label{lem:period-points}
Let $1 \leq k < K'$. Then $\Pi_{k+1}$ is a power of $L$ with $\Pi_{k+1} \leq L^2MR_k \leq a_k$, so that $\Pi_{k+1}$ divides $a_k$. The corners of the loops of an ensemble at level $q$ lie in $\Pi_{k+1}\mathbb Z^4$, the length of every segment is a multiple of $\Pi_{k+1}$, and so are the periods of the torus.
\end{lemma}

\begin{proof}
At the step $j + 1$ the cubes listed above have, in units of the scale $j$, the sides $M$ for $\pi_j$, $LM \leq MR_j$ for $\pi_{j+1}$, $MR_j$ and $LMR_{j+1} \leq LMR_j$ for the domains of the steps $j$ and $j + 1$, $L^2M_2R_{j+1} \leq LMR_j$ for the test functions, and $R_1M_1 \leq M$ \cite[p.~3]{BalabanRG2}; here we used $R_{j+1} \leq R_j$ from Lemma~\ref{lem:tube-cubes}(ii), $LM_2 < M$ from \cite[(3.1)]{WilsonFirstStep}, and $R_j \geq L$, which holds for $\gamma$ small by Lemma~\ref{lem:tube-cubes}(i). All of them are therefore at most $L^2MR_j$ in units of the scale $j$, that is, at most $L^{j-k}L^2MR_j$ in units of the scale $k$. For $j < k$ the second bound of \eqref{eq:tube-two-nine} gives
\begin{equation}\label{eq:period-scale-comparison}
    L^{j-k}R_j \leq (L + 1)(k - j)^{\beta_0}L^{j-k}R_k \leq (1 + L^{-1})R_k\;,
\end{equation}
because $m^{\beta_0} \leq m \leq L^{m-1}$ for every integer $m \geq 1$. A side which is a power of $L$ and at most $(1 + L^{-1})L^2MR_k < L^3MR_k$ is at most $L^2MR_k$, since $L^2MR_k$ is a power of $L$ as well. Hence $\Pi_{k+1} \leq L^2MR_k$, and $L^2MR_k \leq a_k$ by \eqref{eq:lengths-at-step}. Both $\Pi_{k+1}$ and $a_k = L^{j_q-q}L^{K'-k}$ are powers of $L$, so the first divides the second.

A loop at level $q$ is a path of bonds of $\mathbb T_q$, so that its corners lie in $uL^{-q}\mathbb Z^4$, which is $a_k\mathbb Z^4$ in units of the scale $k$ by the proof of Lemma~\ref{lem:stopping-index}, and its segments, each of which joins two corners along one axis, have lengths which are multiples of $a_k$. The period $T = 2L^{m_T}u$ is $2L^{m_T+q}a_k$ in units of the scale $k$.
\end{proof}

\emph{The reference gas.} By Section~\ref{sec:small-field-map}, the insertion $Y_{\mathcal G}$ of a group $\mathcal G$ of the step $k + 1$ whose intervals lie on one loop takes at the reference data of $\ell$ the value $c_{\mathcal G}\,\mathrm{id}$, with a number $c_{\mathcal G} = c_{\mathcal G}(\ell)$ which is defined by the reference clause of Definition~\ref{def:dressed-loop}(iii); for a cross group we put $c_{\mathcal G} \coloneqq 0$. We also put $c_{\mathcal G} \coloneqq 0$ for an \emph{attached} group, by which we mean a group which is attached to a large-field region of its step or which contains a part of an insertion attached at an earlier step. The insertion of an attached group is thus carried whole, as that of a cross group is. Nothing is lost by this, since no bound of Definition~\ref{def:dressed-loop} for the attached insertions rests on the subtraction of a value at the reference data, and something is gained: the activities of the reference gas, and with them the interactions and the local terms of every later step, involve no attached insertion (Supplement, Lemma~18.2), so that the groups which exist on $\ell$ alone because they contain a part of an attached insertion do not enter the comparison of the reference gas of $\ell$ with that of the trivial term. The groups are those of Definition~\ref{def:groups-of-step}, in which the interactions of the step $k$ and of the vacuum polymers have been expanded and their clusters merged into the groups, the factor of each cluster becoming a factor of the insertion of its group; two groups interact only through the hard core of the compatibility defined there. The \emph{reference gas} of $\ell$ is the hard-core gas of these groups with the activities $c_{\mathcal G}$, and its partition functions are
\begin{equation}\label{eq:reference-partition}
    S^\circ(\ell) \coloneqq \sum_{\mathfrak E}\prod_{\mathcal G \in \mathfrak E}c_{\mathcal G}\;,
    \qquad
    S^{\circ,[D]}(\ell) \coloneqq \sum_{\mathfrak E \perp D}\prod_{\mathcal G \in \mathfrak E}c_{\mathcal G}\;,
\end{equation}
the first sum over all families of pairwise compatible groups, the second over the families $\mathfrak E$ such that $D \cup \mathfrak E$ is a family and $D \cap \mathfrak E = \emptyset$, for a given family $D$. A \emph{cluster} $\Gamma$ is a finite multiset of groups whose incompatibility graph is connected, and its weight is $\Phi^\circ(\Gamma; \ell) \coloneqq \phi^{\mathrm U}(\Gamma)\prod_{\mathcal G \in \Gamma}c_{\mathcal G}$, with the Ursell function $\phi^{\mathrm U}$ of that graph \cite[Chapter~5]{FriedliVelenik}. The set $J_\Gamma$ of the traversals in the arcs of the groups of $\Gamma$ is nonempty, since every group has an arc, and the \emph{network} $N(\Gamma)$ is the union of the domains, the arcs, and the merged clusters of its groups. On the trivial term we write $\Phi^\circ(\Gamma; T)$ for the weights on the torus $T$. The tori of the family $\mathfrak T$ are those whose periods are multiples of the period $T$ of $\mathbb T$; on each of them the construction is defined by the same rules, with the partitions of $\mathbb T$ lifted, and every bound of the preceding sections holds on them with the same constants, since none of these constants depends on the torus.

\begin{definition}[Amplitude]\label{def:amplitude}
For $j \geq 1$ the \emph{amplitude} of the step $j$ is
\begin{equation}\label{eq:amplitude}
    \mathfrak a_j \coloneqq \max\bigl(\tau_{j-1}^2,\, 2\rho_j\bigr)\;,
\end{equation}
and the \emph{weight} of a group of the step $k + 1$ is $a^\circ(\mathcal G) \coloneqq (\nu/32)\sigma_{k+1}(\mathcal G)$, extended additively to clusters.
\end{definition}

The factor $2$ is the value at $q_1 = q_0$ of the factor $1 + (\log g_k^{-2})^{(q_1 - q_0)_+}$ which the exponents of the radii $\alpha_{0,k}$, $\alpha_{1,k}$ contribute to the sizes of the groups, and the construction takes $q_1 = q_0$ (Appendix~\ref{app:constants}). Since $\tau(g)$ and $\rho(g)$ increase for small $g$ and $g_j \leq (1 + \beta_0)\bar g$ for $j \leq K$ by \eqref{eq:tube-two-six} and the window \eqref{eq:coupling-window}, the amplitudes are bounded by
\begin{equation}\label{eq:amplitude-bound}
    \bar{\mathfrak a} \coloneqq \sup_{j \leq K}\mathfrak a_j \leq C\bar g^2(\log\bar g^{-2})^{6p_0}\;,
    \qquad
    \sup_{j \leq K'}\mathfrak a_j \leq C\gamma_q^2(\log\gamma_q^{-2})^{6p_0}\;,
\end{equation}
the second because $g_j \leq \gamma_q$ for $j \leq K'$ by \eqref{eq:stopping-coupling}.

\begin{lemma}[Root Bound and Tail]\label{lem:root-bound}
Let $\gamma_q$ be sufficiently small, depending on $c_{\mathcal C}$. On every labelled term and every torus of $\mathfrak T$, for every traversal $b$ and every $R \geq 0$,
\begin{equation}\label{eq:root-bound}
    \sum_{\mathcal G\colon b \in J_{\mathcal G}}\abs{c_{\mathcal G}}e^{2a^\circ(\mathcal G)} \leq C_{S'}\mathfrak a_{k+1}\;,
    \qquad
    \sum_{\Gamma\colon b \in J_\Gamma}\abs{\Phi^\circ(\Gamma)}e^{a^\circ(\Gamma)} \leq C_{S'}\mathfrak a_{k+1}\;,
\end{equation}
where $J_{\mathcal G}$ is the set of the traversals in the arc of $\mathcal G$, and
\begin{equation}\label{eq:cluster-tail}
    \sum_{\Gamma\colon b \in J_\Gamma,\ \diam N(\Gamma) \geq R}\abs{\Phi^\circ(\Gamma)} \leq C_{S'}\mathfrak a_{k+1}e^{-2c^T_ER}\;,
    \qquad
    c^T_E \coloneqq \frac{\nu}{64A^*_{\mathrm{env}}ML}\;,
\end{equation}
with $A^*_{\mathrm{env}} \coloneqq \max(2, A_{\mathrm{env}} + 1)$ and $A_{\mathrm{env}}$ the constant of the envelopes of \cite[Lemma~E.8]{WilsonFirstStep}.
\end{lemma}

\begin{proof}
For the first bound we note that $\abs{c_{\mathcal G}} \leq S_{\mathcal G}$ in \eqref{eq:size-bound}, and that $e^{2a^\circ} = e^{(\nu/16)\sigma_{k+1}} \leq e^{\nu\sigma_{k+1}}$. Hence the first sum is at most $2C_{S'}\rho_{k+1} \leq C_{S'}\mathfrak a_{k+1}$ by Corollary~\ref{cor:production-bound} with $\theta = 0$, applied at the traversal of the step $k + 1$ which contains $b$, and by \eqref{eq:amplitude}; the same holds for the sums rooted at a traversal of the step $k + 1$ or at a cube of $\pi_{k+1}$. On the terms with large-field regions, and at the steps at which the $\mathbf R$-operation acts, the bound \eqref{eq:size-bound} is that of Section~\ref{sec:large-field-step}.

For the second bound we verify the condition of Koteck\'y and Preiss with the weight $a^\circ$. The domain of a group $\mathcal G'$ which is incompatible with $\mathcal G$ has a point in common with $X_{\mathcal G}$, by Lemma~\ref{lem:group-covering}(v), so that $\mathcal G'$ is at a cube of $\pi_{k+1}$ which belongs to $X_{\mathcal G}$ or touches it. There are at most $81\sigma_{k+1}(\mathcal G) \leq 81M\sigma_{k+1}(\mathcal G)$ such cubes, since $\sigma_{k+1}$ counts the cubes of $X_{\mathcal G}$ and a cube has $81$ cubes which touch it or are equal to it. By the first bound, rooted at these cubes,
\begin{equation}\label{eq:kp-condition}
    \sum_{\mathcal G' \not\sim \mathcal G}\abs{c_{\mathcal G'}}e^{2a^\circ(\mathcal G')} \leq 81M\sigma_{k+1}(\mathcal G)C_{S'}\mathfrak a_{k+1} \leq \frac{\nu}{32}\sigma_{k+1}(\mathcal G) = a^\circ(\mathcal G)\;,
\end{equation}
where the second inequality is $2592\,MC_{S'}\mathfrak a_{k+1} \leq \nu$, which holds for $\gamma_q$ small, depending on $c_{\mathcal C}$, by the second bound of \eqref{eq:amplitude-bound}, because $k + 1 \leq K'$. This is the condition of \cite[Theorem~5.4]{FriedliVelenik} with the functions $a = a^\circ$ and $d = a^\circ$. That theorem gives the convergence of the cluster expansion of the reference gas and of the gases with the activities $c_{\mathcal G}\mathbf 1[\mathcal G \perp D]$, and it bounds the clusters which are incompatible with an external polymer by the sum of $\abs{c_{\mathcal G}}e^{a^\circ(\mathcal G) + d(\mathcal G)}$ over the groups incompatible with it. With the traversal $b$ as the external polymer, incompatible with the groups whose arcs contain it, this is the second bound of \eqref{eq:root-bound}.

For the tail we bound the diameter of the network by the weight. The groups of a cluster are connected under incompatibility, and two incompatible groups have networks at distance at most one cube side $M$, in units of the scale $k + 1$. The network of one group is connected and has diameter at most $M\#_{k+1}(\mathcal G) + n_{k+1}(\mathcal G) = M\sigma_{k+1}(\mathcal G)$ in units of the scale $k + 1$. Hence, in units of the scale $k$, and since $\sigma_{k+1} \geq 1$ for a group with a domain,
\begin{equation}\label{eq:network-diameter}
    \diam N(\Gamma) \leq L\sum_{\mathcal G \in \Gamma}\bigl(M\sigma_{k+1}(\mathcal G) + M\bigr) \leq 2ML\sum_{\mathcal G \in \Gamma}\sigma_{k+1}(\mathcal G) = \frac{64ML}{\nu}\,a^\circ(\Gamma)\;,
\end{equation}
so that $a^\circ(\Gamma) \geq (\nu/(64ML))\diam N(\Gamma) \geq 2c^T_E\diam N(\Gamma)$, because $A^*_{\mathrm{env}} \geq 2$. The clusters in \eqref{eq:cluster-tail} therefore have $e^{a^\circ(\Gamma)} \geq e^{2c^T_ER}$, and the second bound of \eqref{eq:root-bound} gives \eqref{eq:cluster-tail}.
\end{proof}

The rate $2c^T_E$ is the one which the first step allows: there the network of a cluster has diameter at most $A^*_{\mathrm{env}}LM\sum\sigma$ fine bonds, by the envelopes of \cite[Lemma~E.8]{WilsonFirstStep}, and the same argument gives \eqref{eq:cluster-tail} at the step $1$ with the same constant.

\begin{lemma}[Locality and Periodicity]\label{lem:locality}
On the trivial term the following hold.
\begin{enumerate}[label=(\roman*)]
\item Let $\Gamma$ be a cluster of the gas of a loop $C$ on a torus $T$ all of whose groups have arcs of fewer than $P_k(C)/2$ traversals. Let $C'$ be a loop on a torus $T' \in \mathfrak T$ which contains the arcs of the groups of $\Gamma$, with their orientation, each of them as an arc of fewer than $P_k(C')/2$ traversals of $C'$, and let the dependence set of $\Gamma$ embed isometrically in $T'$. Then $\Gamma$ is a cluster of the gas of $C'$ on $T'$ with the same incompatibility graph, and its weight on $T'$ is the weight on $T$ evaluated on the configuration of $T'$ on the image of the dependence set.
\item A translation by a vector of $\Pi_{k+1}\mathbb Z^4$ maps the reference gas of a loop $C$ onto the reference gas of the translated loop, with the same activities and the same weights.
\end{enumerate}
\end{lemma}

\begin{proof}
(i) The insertion of a group of the step $k + 1$ is built from the activities of the polymers of its domains, the tensors $Y_G$ of Lemma~\ref{lem:word-expansion} and the old insertions in its arc, the transports along its arc, and the clusters of the step $k$ merged into it (Section~\ref{sec:small-field-map} and the merging above). None of these depends on the loop outside the arc: the activities are functions of the fields on their polymers, the tensors $Y_G$ and the transports are attached to the arc, and the old insertions and old cluster weights have the same property by induction, starting from the first step, where the groups are products along their intervals \cite[Definition~3.24]{WilsonFirstStep}. The length of the loop enters only through the rule which takes the smallest arc containing a set, and this rule is applied to sets $S$ of traversals inside the arc $A$ of a group of $\Gamma$, at this step or at an earlier one. If $A$ has fewer than half of the traversals of a loop, its complement in that loop lies in one gap of $S$ and is longer than half of the loop, so that this gap is the longest one; the smallest arc containing $S$ is its complement, the part of $A$ between the first and the last traversal of $S$. By the hypotheses on the lengths this description applies on $C$ and on $C'$, and it gives the same arc on both. The hard core depends only on the arcs, the domains, and the merged clusters, so that the incompatibility graph and the Ursell function are the same.

(ii) Every operation of Ba{\l}aban's step, the averages, the minimisers, the propagators and their random walk expansions, the test functions, the localisation, and the activities, is defined from the lattice, the partitions, and the fields alone. A translation by a vector of $\Pi_{k+1}\mathbb Z^4$ preserves the lattices and the partitions, and it maps the trivial field to itself, so it maps each object to the object of the translated data. For the sums of the regular part and for the $\mathbf R$-part this covariance is printed \cite[(2.29), (2.32), p.~260]{BalabanCRE}; for the localised activities it follows from their definitions, provided the choices which the papers leave free are made by translation covariant rules, which we do (Section~\ref{sec:notes-on-papers}).
\end{proof}

The comparison of two tori needs a condition which excludes that a small cluster wraps around the torus. Let $\rho_{\mathrm{dep}} < LMR_1/4$ be the dependence range of the objects of the first step in fine bonds \cite[Sections~E.9.2 and~E.10.4]{WilsonFirstStep}, and put, for $1 \leq j \leq K'$,
\begin{equation}\label{eq:torus-margin}
    \iota_{T,j} \coloneqq \max\Bigl(0,\, \frac{N^{(j)}_{\min}}{4} - \rho_{\mathrm{dep}}L^{-j}\Bigr)\;,
\end{equation}
in units of the scale $j$, so that $\iota_{T,j} = L\iota_{T,j+1}$ whenever both are positive. The dependence set of an object of the step $k + 1$ lies within the distance
\begin{equation}\label{eq:dependence-range}
    R^{(k)}_D \leq 5M + LM + L^2MR_k + \max_{j<k}L^{j-k}R^{(j)}_D \leq 6L^2MR_k
\end{equation}
of its network: the walk pieces of the propagators depend on the field in the fifth enlargement of their domains by cubes of side $M_1$ \cite[Theorem~3.10, p.~416]{BalabanPBG}, at most $5M$; the localisation of the step uses the cubes of $\pi_{k+1}$ \cite[p.~13]{BalabanRG2}, of side $LM$; the test functions use cubes of side $L^2M_2R_{k+1} \leq L^2MR_k$; and the old insertions have their own ranges. The last inequality follows by iteration from $5M + LM + L^2MR_k \leq 2L^2MR_k$, from $R_{k-m} \leq (L + 1)m^{\beta_0}R_k$ by \eqref{eq:tube-two-nine}, and from $(L + 1)\sum_{m \geq 1}mL^{-m} = (L + 1)L/(L - 1)^2 \leq 2$ for $L \geq 5$, the range $\rho_{\mathrm{dep}}L^{-k}$ of the first step being smaller than the term $m = k - 1$.

\begin{lemma}[The Torus Condition]\label{lem:torus-condition}
Let $L \geq 195$. For every loop ensemble at level $q$ and every $1 \leq k \leq K'$,
\begin{equation}\label{eq:torus-condition}
    \iota_{T,k} \geq 16L\cdot 6L^2MR_k \geq 16L\,R^{(k)}_D\;.
\end{equation}
\end{lemma}

\begin{proof}
A loop at level $q$ has a bond of $\mathbb T_q$, so that $uL^{-q} \leq \diam(C)$, and $3\diam(C) < T = 2L^{m_T}u$ by Definition~\ref{def:loop-ensemble}; hence $L^{m_T+q} > 3/2$ and $m_T + q \geq 1$. By Lemma~\ref{lem:period-points} the smallest period in units of the scale $K'$ is $N^{(K')}_{\min} = 2L^{m_T+q}a_{K'} \geq 2L\cdot L^3MR_{K'-1} \geq 2L^4MR_{K'}$, by \eqref{eq:stopping-identity-index} and Lemma~\ref{lem:tube-cubes}(ii). The first-step range is $\rho_{\mathrm{dep}}L^{-K'} \leq L^{1-K'}MR_1/4 \leq MR_{K'}/2$, by \eqref{eq:tube-two-nine} with $m^{\beta_0}L^{-m} \leq L^{-1}$ for $m = K' - 1 \geq 1$, and trivially for $K' = 1$. Hence $\iota_{T,K'} \geq (L^4 - 1)MR_{K'}/2$. For $k < K'$ we have $\iota_{T,k} = L^{K'-k}\iota_{T,K'}$ and $R_k \leq (L + 1)(K' - k)^{\beta_0}R_{K'} \leq (1 + L^{-1})L^{K'-k}R_{K'}$, so that in all cases
\begin{equation}\label{eq:torus-margin-bound}
    \iota_{T,k} \geq \frac{L^4 - 1}{2(1 + L^{-1})}\,MR_k \geq 96L^3MR_k\;,
\end{equation}
where the last inequality is $L^4 - 1 \geq 192L^3 + 192L^2$, which holds for $L \geq 195$. The second inequality of \eqref{eq:torus-condition} is \eqref{eq:dependence-range}.
\end{proof}

A cluster $\Gamma$ \emph{embeds} in $T$ if its dependence set, the $R^{(k)}_D$-neighbourhood of its network, has diameter less than $\iota_{T,k}/4$; the covering map of a torus $T' \in \mathfrak T$ over $T$ is then injective on its lifts, and $\tilde\Gamma$ denotes a lift.

\begin{proposition}[Torus Consistency]\label{prop:torus-consistency}
Let $L$ satisfy \eqref{eq:blocking-factor-condition}, let $\gamma_q$ be sufficiently small, depending on $c_{\mathcal C}$, and assume that the ensemble is dressed at the step $k$. On the trivial term, for a loop $C$ of the ensemble, at the reference data:
\begin{enumerate}[label=(\roman*)]
\item for tori $T \subset T'$ of $\mathfrak T$ and every traversal $b$,
\begin{equation}\label{eq:torus-comparison}
    \sum_{\Gamma\colon b \in J_\Gamma}\abs{\Phi^\circ(\Gamma; T) - \Phi^\circ(\tilde\Gamma; T')} \leq \bar C_Z\mathfrak a_{k+1}e^{-c^Z\iota_{T,k+1}}\;,
    \qquad
    c^Z \coloneqq \frac{\nu}{256A^*_{\mathrm{env}}M}\;,
\end{equation}
where $\Phi^\circ(\Gamma; T) \coloneqq 0$ for a cluster of $T'$ which is not the lift of a cluster embedding in $T$;
\item along the tori $T_m \in \mathfrak T$ of the periods $L^mT$, the weight of every cluster which embeds in some $T_m$ converges to a limit $\Phi^{\circ,\infty}(\Gamma)$, which does not depend on the sequence of tori, and
\begin{equation}\label{eq:torus-limit}
    \sum_{\Gamma\colon b \in J_\Gamma}\abs{\Phi^\circ(\Gamma; T) - \Phi^{\circ,\infty}(\Gamma)} \leq 3\bar C_Z\mathfrak a_{k+1}e^{-c^Z\iota_{T,k+1}}\;;
\end{equation}
\item the clause (vi) of Definition~\ref{def:dressed-loop} holds at the step $k + 1$.
\end{enumerate}
\end{proposition}

\begin{proof}
We split the clusters at $b$ into those which embed in $T$ and those which do not.

\emph{The embedding clusters.} Let $\Gamma$ embed in $T$. The activities of its groups are composed of two kinds of factors. The factors made at the step $k + 1$ are the same functions on $T$ and on $T'$ of the configuration on the dependence set: the walk pieces of the propagators with boundary conditions depend on the field in their domains only \cite[p.~394; Theorem~3.10, p.~416]{BalabanPBG}; a term of the localised expansion ``depends on $\mathbf U$, $\mathbf J$ restricted to $Y_0 \cup \square_0$'' \cite[p.~5]{BalabanRG2}, and the activity $H(Z)$ of the fluctuation integral ``is localized in the interior of $Z$ with respect to the external gauge fields'' \cite[p.~14]{BalabanRG2}; the test functions are local in their cubes; and the couplings are the same numbers on every torus, because the coefficients $\beta_{j+1}$ are defined through the limit of the vacuum polarisation ``as $T^{(j+1)} \nearrow Z^d$'' \cite[p.~264]{BalabanRG1} and the couplings of a tube are prescribed. The factors made at earlier steps, the insertions of the step $k$ inside the fluctuation integral and the merged clusters of the step $k$, differ on the two tori by at most the amounts of the clause (vi) at the step $k$, at fields which coincide on the $\iota_{T,k}/4$-neighbourhood of the dependence set; the deformation of the background by the fluctuation field vanishes outside one lift of the domain of the integral \cite[p.~4]{BalabanRG2}, so this hypothesis holds at every field of the integral. Replacing these factors by their differences one at a time, and summing over the clusters at $b$ with one marked factor by the count which gives \eqref{eq:root-bound}, we obtain for the embedding clusters the bound $2C^{(k+1)}_{\mathrm{rel}}\mathfrak a_{k+1}e^{-\iota_{T,k+1}/(32M)}$, with
\begin{equation}\label{eq:relative-recursion}
    C^{(k+1)}_{\mathrm{rel}} \leq \bigl(L(1 + \beta_0)^7 + \tfrac18\bigr)e^{-(L-1)\iota_{T,k+1}/(32M)}C^{(k)}_{\mathrm{rel}} + C_*\;.
\end{equation}
Here the factor $L$ counts the traversals of the step $k$ in one of the step $k + 1$, $(1 + \beta_0)^7$ compares the amplitudes of the two steps, the exponential is the gain $e^{-\iota_{T,k}/(32M)} = e^{-L\iota_{T,k+1}/(32M)}$ of the clause (vi) at the step $k$ against the rate of the step $k + 1$, and $C_*$ bounds the clusters of Ba{\l}aban's gas of polymers inside the activities which do not embed, whose sets differ between the two tori. By \eqref{eq:torus-margin-bound}, $(L - 1)\iota_{T,k+1}/(32M) \geq 3L^2(L - 1)R_{k+1}$, so the prefactor in \eqref{eq:relative-recursion} is at most $1/2$, and since $C^{(1)}_{\mathrm{rel}} = 0$ by Theorem~\ref{thm:first-step}, we get $C^{(k)}_{\mathrm{rel}} \leq 2C_*$ for every $k$. This is also the clause (vi) at the step $k + 1$ for the cluster weights; for the insertions the same telescoping applies, with the deviation parts, which vanish at the reference data, bounded relative to the deviation by the Cauchy estimate of Section~\ref{sec:contraction}. This proves (iii).

\emph{The other clusters.} If $\Gamma$ does not embed in $T$, its dependence set has diameter at least $\iota_{T,k}/4$, so that its network has diameter at least $\iota_{T,k}/4 - 2R^{(k)}_D \geq \iota_{T,k}/8$ in units of the scale $k$, by Lemma~\ref{lem:torus-condition}. By \eqref{eq:network-diameter} its weight is $a^\circ(\Gamma) \geq (\nu/(64ML))\iota_{T,k}/8 = \nu\iota_{T,k+1}/(512M) \geq c^Z\iota_{T,k+1}$, and the second bound of \eqref{eq:root-bound} shows that these clusters contribute at most $C_{S'}\mathfrak a_{k+1}e^{-c^Z\iota_{T,k+1}}$ on $T$, and the same on $T'$, where the clusters at a lift of $b$ which are not lifts of embedding clusters of $T$ are exactly those whose dependence sets have diameter at least $\iota_{T,k}/4$.

Adding the two parts gives (i), with $\bar C_Z \coloneqq 2C_{S'} + 4C_*$, since $c^Z \leq 1/(32M)$. For (ii), any two tori of $\mathfrak T$ are contained in a third one, whose periods are common multiples of theirs, and $\iota_{T_m,k+1} = L^m\iota_{T,k+1}$. By (i) the weights of a cluster along $(T_m)$ form a Cauchy sequence, whose limit does not depend on the sequence by the comparison through a common torus, and the telescoping sum gives the bound $\bar C_Z\mathfrak a_{k+1}\sum_{m \geq 0}e^{-c^ZL^m\iota_{T,k+1}} \leq 2\bar C_Z\mathfrak a_{k+1}e^{-c^Z\iota_{T,k+1}}$ for the clusters which embed in $T$, because $L^m - 1 \geq m(L - 1)$ and $c^Z\iota_{T,k+1} \geq 1$; the latter follows from $\iota_{T,k+1} \geq 96L^3MR_{k+1}$, by \eqref{eq:torus-margin-bound}, and from $R_{k+1} \geq (\log\hat g_{k+1}^{-2})^r \geq (\log\gamma_q^{-2})^r$, for $\gamma_q$ small, since $k + 1 \leq K'$. The clusters which do not embed in $T$ contribute at most $C_{S'}\mathfrak a_{k+1}e^{-c^Z\iota_{T,k+1}}$ on $T$, and the same in the limit, by Fatou's lemma applied to the bound on each $T_m$, which holds with $\iota_{T,k+1}$ in place of $\iota_{T_m,k+1}$ because the condition is on the diameter of the dependence set. This gives \eqref{eq:torus-limit}.
\end{proof}

\emph{The exact factorisation.} With the reference gas we can now carry out the resummation of Section~\ref{sec:small-field-map}. On a labelled term $\ell$ we split every insertion as $Y_{\mathcal G} = c_{\mathcal G}T(\mathcal G) + Y^{\mathrm d}_{\mathcal G}$, where $T(\mathcal G)$ is the product of the representations of the transports of the background field along the arc of $\mathcal G$. For a group on one loop, $Y^{\mathrm d}_{\mathcal G}$ vanishes at the reference data and is covariant with the endpoints of the arc. For a cross group $c_{\mathcal G} = 0$, so that $Y^{\mathrm d}_{\mathcal G} = Y_{\mathcal G}$ is the whole insertion, which does not vanish at the reference data in general and is bounded by \eqref{eq:cross-bound}. Since every word is multilinear in its insertions, and the transports of the scalar parts join the transports of the gaps,
\begin{equation}\label{eq:multilinear-split}
    w(\mathfrak F) = \sum_{D \subseteq \mathfrak F}\Bigl(\prod_{\mathcal G \in \mathfrak F \setminus D}c_{\mathcal G}\Bigr)w^{\mathrm d}(D)\;,
\end{equation}
where $w^{\mathrm d}(D)$ is the word of the family $D$ with the insertions $Y^{\mathrm d}_{\mathcal G}$, and $w^{\mathrm d}(\emptyset) = \prod_if^{(i)}_{k+1}$. After the merging the families of groups interact only through the hard core, so that the sum over the families $\mathfrak F$ of \eqref{eq:multilinear-split} is a sum over the pairs $(D, \mathfrak E)$ with $\mathfrak F = D \cup \mathfrak E$, and
\begin{equation}\label{eq:exact-factorisation}
    \sum_{\mathfrak F}w(\mathfrak F) = S^\circ(\ell)\Bigl[\prod_if^{(i)}_{k+1} + \sum_{D \neq \emptyset}Q_{k+1}(D)\,w^{\mathrm d}(D)\Bigr]\;,
\end{equation}
with
\begin{equation}\label{eq:interaction-ratio}
    Q_{k+1}(D) \coloneqq \frac{S^{\circ,[D]}(\ell)}{S^\circ(\ell)} = \exp\Bigl(-\sum_{\Gamma \not\perp D}\Phi^\circ(\Gamma; \ell)\Bigr)\;,
\end{equation}
the last sum over the clusters which contain a group incompatible with a member of $D$. The rearrangements are justified by the absolute convergence of $\sum_{\mathfrak F}\prod_{\mathcal G \in \mathfrak F}(\abs{c_{\mathcal G}} + \norm{Y^{\mathrm d}_{\mathcal G}})$, which follows from \eqref{eq:size-bound}; the gas is a gas of groups together with their merged clusters, and the sum over the latter is infinite even on a finite torus. The cluster series of $S^\circ(\ell)$ and of $S^{\circ,[D]}(\ell)$ converge by Lemma~\ref{lem:root-bound}, which gives $S^\circ(\ell) \neq 0$ and the second expression for $Q_{k+1}(D)$, and $\sum_{\Gamma \not\perp \mathcal G}\abs{\Phi^\circ(\Gamma)} \leq a^\circ(\mathcal G) \leq \nu\sigma_{k+1}(\mathcal G)/16$ by \cite[Theorem~5.4]{FriedliVelenik} and \eqref{eq:kp-condition}, which is the clause (i) at the step $k + 1$ with $\theta \coloneqq a^\circ$. For an ensemble the groups of all loops and the cross groups form one gas; since cross groups have the activity $0$, on the trivial term
\begin{equation}\label{eq:cross-exponent-definition}
    \log S^\circ(\ell_0) = \sum_{i=1}^{n}F^\circ_{k+1}(C_i; T) + E^\times_{k+1}\;,
    \qquad
    F^\circ_{k+1}(C; T) \coloneqq \sum_{\Gamma}\Phi^\circ(\Gamma; T)\;,
\end{equation}
the second sum over the clusters of the gas of the single loop $C$, and this defines the cross exponent $E^\times_{k+1}$ of \eqref{eq:cross-exponent}. On a general term the difference $E^\partial_{k+1}(\ell) \coloneqq \log S^\circ(\ell) - \log S^\circ(\ell_0)$ consists of the clusters which see the large-field data of $\ell$; its representation as a sum of local terms $e^{(k+1)}_X$ attached to these data, with the first bound of \eqref{eq:scalar-bounds}, is part of Supplement, Proposition~S.9.

\subsection{The Coefficients}\label{sec:coefficients}

For $v \in \mathcal V$ let $\ell_v \coloneqq \{ tv \, | \, t \in \mathbb Z \}$ be the line through the origin in the orientation $v$, with its traversals numbered by $b \in \mathbb Z$, the traversal $b$ running from $bv$ to $(b + 1)v$. For a perpendicular ordered pair $(v, w)$ let $B_{v,w}$ be the bent path \eqref{eq:bent-path} with $c = 0$, $v_c = v$, and $w_c = w$, with its traversals numbered by $b \in \mathbb Z$: for $b < 0$ the traversal from $bv$ to $(b + 1)v$, and for $b \geq 0$ the traversal from $bw$ to $(b + 1)w$. The \emph{line gas} of $\ell_v$ consists of the clusters $\Gamma$ on $\mathbb Z^4$ whose groups have finite arcs on $\ell_v$, with the weights $\Phi^{\circ,\infty}_v(\Gamma)$ obtained as follows: we take a torus $T' \in \mathfrak T$ in which $\Gamma$ embeds and on it a square $C'$ at level $q$, in a coordinate plane containing $v$ and with its corners in $\Pi_{k+1}\mathbb Z^4$, one side of which contains the arcs of $\Gamma$ in the orientation $v$ and is longer than twice the diameter of the network of $\Gamma$; then $\Gamma$ is a cluster of the gas of $C'$, and its weight does not depend on the choice of $C'$ by Lemma~\ref{lem:locality}(i), whose hypothesis on the lengths holds because an arc of $\Gamma$ lies in the network and therefore has fewer traversals than a side of either square, a quarter of its perimeter. Further, $\Phi^{\circ,\infty}_v(\Gamma)$ is the limit of these weights along $(T_m)$ given by Proposition~\ref{prop:torus-consistency}(ii). The \emph{bent-path gas} of $B_{v,w}$ is defined in the same way, with a square in the plane spanned by $v$ and $w$ which has a corner at the origin with the incoming orientation $v$ and the outgoing orientation $w$, and whose two sides at this corner contain the arcs of $\Gamma$ and are longer than twice the diameter of its network; an arc of $\Gamma$ on these two sides, whose ends are within that diameter of each other, then has fewer traversals than one side. The squares used here need not be admissible. The proofs of Theorems~\ref{thm:first-step} and~\ref{thm:inductive-step} use the loops only through $\Lambda(C) = 2$, $P(C)/r(C) = 8$, and the torus conditions, and the torus conditions, which compare the loop with the periods, hold at the given $K$ once the periods of $T'$ are large compared with the square. The densities and the coefficients of the step $k + 1$ are
\begin{equation}\label{eq:line-coefficient}
    a_v(b) \coloneqq \sum_{\Gamma\colon b \in J_\Gamma}\frac{\Phi^{\circ,\infty}_v(\Gamma)}{\abs{J_\Gamma}}\;,
    \qquad
    \kappa^{(k+1)}_v \coloneqq \sum_{b=0}^{\Pi_{k+1}-1}a_v(b)
\end{equation}
on the line, and, with $a_{v,w}(b)$ defined in the same way on $B_{v,w}$ and with $k_{v,w}(b) \coloneqq a_v(b)$ for $b < 0$ and $k_{v,w}(b) \coloneqq a_w(b)$ for $b \geq 0$,
\begin{equation}\label{eq:corner-coefficient}
    \zeta^{(k+1)}_{v,w} \coloneqq \sum_{b \in \mathbb Z}\bigl(a_{v,w}(b) - k_{v,w}(b)\bigr)
\end{equation}
on the bent path.

\begin{lemma}[The Coefficients]\label{lem:coefficients}
Let $L$ and $\bar g$ be as in Proposition~\ref{prop:torus-consistency}, and put $C_\zeta \coloneqq 2C_{S'}(1 + 2/(e^{2c^T_E} - 1))$.
\begin{enumerate}[label=(\roman*)]
\item The density $a_v$ is $\Pi_{k+1}$-periodic, $\abs{a_v(b)} \leq C_{S'}\mathfrak a_{k+1}$, and $\abs{\kappa^{(k+1)}_v} \leq \Pi_{k+1}C_{S'}\mathfrak a_{k+1}$.
\item $\abs{a_{v,w}(b) - k_{v,w}(b)} \leq 2C_{S'}\mathfrak a_{k+1}e^{-2c^T_E\abs{b}}$ for every $b \in \mathbb Z$, so that the series \eqref{eq:corner-coefficient} converges absolutely and $\abs{\zeta^{(k+1)}_{v,w}} \leq C_\zeta\mathfrak a_{k+1}$.
\item The numbers $\kappa^{(k+1)}_v$ and $\zeta^{(k+1)}_{v,w}$ do not depend on the loops, on the torus, or on the labelled term. They depend on $k$, $K$, and $\bar g$ through the couplings and the parameters of the tube, on the character, and on the constants of the construction.
\end{enumerate}
\end{lemma}

\begin{proof}
(i) By Lemma~\ref{lem:locality}(ii), applied on each torus $T_m$, the weights of a cluster and of its translate by $\Pi_{k+1}v$ agree, and so do their limits. The bound on $a_v(b)$ is \eqref{eq:root-bound} at $R = 0$ with $\abs{J_\Gamma} \geq 1$, valid on each $T_m$ and passed to the limit by Fatou's lemma; summing over one period gives the bound on $\kappa^{(k+1)}_v$.

(ii) Let $b \geq 0$. A cluster of the bent-path gas all of whose arcs lie on the outgoing ray is, by Lemma~\ref{lem:locality}(i), a cluster of the line gas of $\ell_w$ with the same weight, and conversely a cluster of the line gas at $b$ all of whose arcs lie in $\{ b' \geq 0 \}$ is one of the bent-path gas. The terms which survive in $a_{v,w}(b) - a_w(b)$ are therefore the clusters of the bent path at $b$ with an arc on the incoming ray, and the clusters of the line at $b$ with an arc on $\{ b' < 0 \}$. The network of each of them contains the position $bw$ of $b$ and a site $-tv$ or $-tw$ with $t \geq 1$, at distance at least $b$, so that \eqref{eq:cluster-tail} with $R = b$, in each of the two gases, gives the bound. The case $b < 0$ is the same with the two rays exchanged, the distance being at least $\abs{b}$. Summing, with $y \coloneqq e^{-2c^T_E}$,
\begin{equation}\label{eq:corner-coefficient-bound}
    \abs{\zeta^{(k+1)}_{v,w}} \leq 2C_{S'}\mathfrak a_{k+1}\Bigl(1 + 2\sum_{b \geq 1}y^b\Bigr) = 2C_{S'}\mathfrak a_{k+1}\frac{1 + y}{1 - y} = C_\zeta\mathfrak a_{k+1}\;.
\end{equation}

(iii) The coefficients are sums of weights of the line and bent-path gases on $\mathbb Z^4$, which by construction involve neither the loops of the ensemble, nor the torus $\mathbb T$, nor any term but the trivial one. The dependence on the step is through the objects of the steps $1, \dots, k + 1$ on the trivial term.
\end{proof}

The coefficients are kept separately for each orientation and each ordered pair of perpendicular orientations, since the covariance of the construction under the hypercubic group is not established; no relation between them is used. With them we define the factor of a step. At the step $1$ the coefficients are those of the first step: $\Pi_1 \coloneqq L$, and $\kappa^{(1)}_v$ and $\zeta^{(1)}_{v,w}$ are the perimeter and corner coefficients of \cite[(5.33), (6.4)]{WilsonFirstStep} in the orientation $v$ and at the corner $(v, w)$ (Section~\ref{sec:first-step-proof}). A step $j$ is an \emph{$\mathbf R$-step} if $j > N_j$, so that the $\mathbf R$-operation acts on the numerator at that step.

\begin{definition}[Step Factor]\label{def:step-factor}
For a loop $C$ with the character $\chi$ and $1 \leq j \leq K'$ let $N^{(j-1)}_v(C)$ be the number of traversals of $C$ of the orientation $v$ at the step $j - 1$. The \emph{step factor} is $\kappa_j(C, \chi) \coloneqq \exp\log\kappa_j(C, \chi)$ with
\begin{equation}\label{eq:step-factor}
    \log\kappa_j(C, \chi) \coloneqq \sum_{v \in \mathcal V}\frac{N^{(j-1)}_v(C)}{\Pi_j}\,\kappa^{(j)}_v + \sum_{c \in \mathcal K(C)}\zeta^{(j)}_{v_c,w_c} + \mathbf 1_{\mathbf R}(j)\Bigl[\sum_{v \in \mathcal V}N^{(j-1)}_v(C)\,\mu^{\mathbf R}_{j,v} + \sum_{c \in \mathcal K(C)}\zeta^{\mathbf R}_{j,v_c,w_c}\Bigr]\;,
\end{equation}
where $\mathbf 1_{\mathbf R}(j)$ is $1$ at an $\mathbf R$-step and $0$ otherwise, and $\mu^{\mathbf R}_{j,v}$, $\zeta^{\mathbf R}_{j,v,w}$ are the coefficients of Proposition~\ref{prop:r-flat-part}. The coefficients depend on $\chi$ through the representation $\pi$ in the insertions.
\end{definition}

The factor is not zero, being an exponential, and $\log\kappa_j$ has exactly the structure announced in Section~\ref{sec:results}: a term proportional to the length of the segments of each orientation and a term for each corner, with coefficients which are numbers of the step.

\subsection{The Remainder of a Step}\label{sec:flat-remainder}

\begin{theorem}[The Flat Part of a Step]\label{thm:flat-part}
Let $L$ and $\bar g$ be as in Proposition~\ref{prop:torus-consistency}, let $1 \leq k < K'$, and assume that the ensemble is dressed at the step $k$. For every loop $C$ of the ensemble,
\begin{equation}\label{eq:flat-part}
    F^\circ_{k+1}(C; T) = \sum_{v \in \mathcal V}\frac{N^{(k)}_v(C)}{\Pi_{k+1}}\,\kappa^{(k+1)}_v + \sum_{c \in \mathcal K(C)}\zeta^{(k+1)}_{v_c,w_c} + E^\circ_{k+1}(C, T)\;,
\end{equation}
with a number $E^\circ_{k+1}(C, T)$ which satisfies, with $C^{\mathcal K}_E \coloneqq 2C_{S'} + 2C_\zeta$,
\begin{equation}\label{eq:flat-remainder}
    \abs{E^\circ_{k+1}(C, T)} \leq C^{\mathcal K}_E\mathfrak a_{k+1}\bigl(P_k(C) + \abs{\mathcal K(C)}\bigr)e^{-2c^T_Er_k(C)} + 3\bar C_Z\mathfrak a_{k+1}P_k(C)e^{-c^Z\iota_{T,k+1}}\;.
\end{equation}
\end{theorem}

\begin{proof}
We allocate every cluster equally to the traversals of $J_\Gamma$,
\begin{equation}\label{eq:allocation}
    F^\circ_{k+1}(C; T) = \sum_{b \in C}a_C(b; T)\;,
    \qquad
    a_C(b; T) \coloneqq \sum_{\Gamma\colon b \in J_\Gamma}\frac{\Phi^\circ(\Gamma; T)}{\abs{J_\Gamma}}\;,
\end{equation}
the sum over the $P_k(C)$ traversals, which converges absolutely by \eqref{eq:root-bound}. Every traversal is assigned to one corner (Section~\ref{sec:corner-neighbourhood}), and we compare the density at a traversal $b$ assigned to the corner $c$ with the density of the bent path $B_c = c + B_{v_c,w_c}$ at the corresponding traversal.

\emph{Small clusters.} Let $\Gamma$ be a cluster at $b$ whose network has diameter less than $r_k(C)$. By Lemma~\ref{lem:corner-isolation} every site of $C$ or of $B_c$ outside $S_{\mathrm{in}}(c) \cup S_{\mathrm{out}}(c)$ is at distance at least $r_k(C)$ from the position of $b$, while the arcs of the groups of $\Gamma$ lie in its network; so these arcs lie in $S_{\mathrm{in}}(c) \cup S_{\mathrm{out}}(c)$, within distance $r_k(C)$ of the position of $b$, hence in an arc of fewer than $4r_k(C) \leq P_k(C)/2$ traversals, since the chord--arc constant of two consecutive segments is $2$ by Lemma~\ref{lem:square-properties}(iii). The cluster does not wrap around the torus, because $r_k(C) \leq \diam_k(C) < N^{(k)}_{\min}/3$. By Lemma~\ref{lem:locality}(i), whose hypothesis on the lengths holds on $C$ by this count and on the squares of the bent-path gas by their definition, it is a cluster of the bent-path gas of $B_c$, and conversely a cluster of that gas at $b$ whose network has diameter less than $r_k(C)$ is a cluster of the gas of $C$. By Lemma~\ref{lem:period-points} the corner $c$ lies in $\Pi_{k+1}\mathbb Z^4$, so by Lemma~\ref{lem:locality}(ii) the weights of the gas of $B_c$ are those of $B_{v_c,w_c}$ at the translated traversal. Their weights on $T$ and on $\mathbb Z^4$ differ by at most $3\bar C_Z\mathfrak a_{k+1}e^{-c^Z\iota_{T,k+1}}$ in the sum at $b$, by \eqref{eq:torus-limit}.

\emph{Large clusters.} The clusters at $b$ whose networks have diameter at least $r_k(C)$ contribute at most $C_{S'}\mathfrak a_{k+1}e^{-2c^T_Er_k(C)}$ in the gas of $C$ and the same in the gas of $B_c$, by \eqref{eq:cluster-tail}. This includes the clusters with a group whose arc has at least $P_k(C)/2$ traversals, whose networks have diameter at least $P_k(C)/(2\Lambda(C)) \geq r_k(C)$ by Lemma~\ref{lem:long-arcs} and Lemma~\ref{lem:chord-arc}.

Together, writing $b$ also for the corresponding traversal of $B_{v_c,w_c}$,
\begin{equation}\label{eq:density-comparison}
    \bigl\lvert a_C(b; T) - a_{v_c,w_c}(b)\bigr\rvert \leq 2C_{S'}\mathfrak a_{k+1}e^{-2c^T_Er_k(C)} + 3\bar C_Z\mathfrak a_{k+1}e^{-c^Z\iota_{T,k+1}}\;.
\end{equation}

\emph{The corner sums.} Let $\ell$ and $\ell'$ be the lengths of $S_{\mathrm{in}}(c)$ and $S_{\mathrm{out}}(c)$. The traversals assigned to $c$ are the traversals $b$ of $B_{v_c,w_c}$ with $-\floor{\ell/2} \leq b \leq \ceil{\ell'/2} - 1$. By \eqref{eq:corner-coefficient} the sum of $a_{v_c,w_c}(b)$ over them is the sum of $k_{v_c,w_c}(b)$ over them, plus $\zeta^{(k+1)}_{v_c,w_c}$, minus the sum of $a_{v_c,w_c}(b) - k_{v_c,w_c}(b)$ over the other $b$, which have $b \geq \ceil{\ell'/2} \geq s_k(C)/2$ or $\abs{b} \geq \floor{\ell/2} + 1 \geq (s_k(C) + 1)/2$. By Lemma~\ref{lem:coefficients}(ii), with $y = e^{-2c^T_E}$, this last sum is at most
\begin{equation}\label{eq:corner-tail-sum}
    2C_{S'}\mathfrak a_{k+1}\Bigl(\sum_{b \geq s_k(C)/2}y^b + \sum_{b \geq (s_k(C)+1)/2}y^b\Bigr) \leq \frac{4C_{S'}\mathfrak a_{k+1}}{1 - y}e^{-c^T_Es_k(C)} \leq 2C_\zeta\mathfrak a_{k+1}e^{-2c^T_Er_k(C)}\;,
\end{equation}
because $C_\zeta \geq 2C_{S'}/(1 - y)$ and $s_k(C) \geq 2r_k(C)$.

\emph{The segment sums.} A segment $\sigma$ from the corner $c$ to the corner $c''$, of orientation $v$ and length $\ell_\sigma$, receives the terms $k_{v_c,w_c}(b) = a_v(b)$, $0 \leq b < \ceil{\ell_\sigma/2}$, from $c$, and the terms $k_{v_{c''},w_{c''}}(b) = a_v(b)$, $-\floor{\ell_\sigma/2} \leq b < 0$, from $c''$, read on the line $c + \ell_v$, since $c'' = c + \ell_\sigma v$ and both corners lie in $\Pi_{k+1}\mathbb Z^4$. Together these are the terms $a_v(b)$, $0 \leq b < \ell_\sigma$, of the line $\ell_v$, by Lemma~\ref{lem:locality}(ii), and since $\ell_\sigma$ is a multiple of $\Pi_{k+1}$ (Lemma~\ref{lem:period-points}) and $a_v$ is $\Pi_{k+1}$-periodic, their sum is exactly $(\ell_\sigma/\Pi_{k+1})\kappa^{(k+1)}_v$. Summing over the segments gives the first term of \eqref{eq:flat-part}, since $\sum_{\sigma \text{ of orientation } v}\ell_\sigma = N^{(k)}_v(C)$.

\emph{Conclusion.} Adding \eqref{eq:density-comparison} over the $P_k(C)$ traversals and \eqref{eq:corner-tail-sum} over the $\abs{\mathcal K(C)}$ corners gives \eqref{eq:flat-remainder}, with $e^{-2c^T_Er_k(C)}$ for both corner terms.
\end{proof}

At the steps at which the $\mathbf R$-operation acts, its polymer gas linked to the loop has a flat part of the same structure, with the period $MR_j$ of the cubes of the $\mathbf R$-operation in place of $\Pi_{j}$, the decay $e^{-\kappa d_j}$ of its activities in place of \eqref{eq:cluster-tail}, and a group's hull in place of its network.

\begin{proposition}[The Flat Part of an $\mathbf R$-Step]\label{prop:r-flat-part}
Let $j \leq K'$ be an $\mathbf R$-step. There are numbers $\mu^{\mathbf R}_{j,v}$, $v \in \mathcal V$, and $\zeta^{\mathbf R}_{j,v,w}$, $(v, w)$ perpendicular, independent of the loops, the torus, and the labelled term, with $\abs{\mu^{\mathbf R}_{j,v}}, \abs{\zeta^{\mathbf R}_{j,v,w}} \leq C^R_Ee^{-p_0(g_j)/4}$, such that for every loop $C$ of a dressed ensemble the flat part of the exponent which the $\mathbf R$-operation of the step $j$ contributes along $C$ is $\sum_vN^{(j-1)}_v(C)\mu^{\mathbf R}_{j,v} + \sum_{c \in \mathcal K(C)}\zeta^{\mathbf R}_{j,v_c,w_c} + E^{\circ,\mathbf R}_j(C, T)$, with
\begin{equation}\label{eq:r-flat-remainder}
    \abs{E^{\circ,\mathbf R}_j(C, T)} \leq C^R_EMR_je^{-p_0(g_j)/4}\Bigl[\bigl(P_j(C) + \abs{\mathcal K(C)}\bigr)e^{-c^R_Er_j(C)/(MR_j)} + P_j(C)e^{-c^R_E\iota_{T,j}/(MR_j)}\Bigr]\;.
\end{equation}
\end{proposition}

\begin{proof}
Supplement, Proposition~S.10. The proof is that of Theorem~\ref{thm:flat-part} for the gas of the $\mathbf R$-operation linked to the loop, whose flat activities are periodic along each segment with the period $MR_j$; every segment is a union of such periods starting at a corner, because $MR_jL^j$ divides $MR_{K'}L^{K'}$, which divides the fine length $L^{j_q-q+K'}$ of a bond of $\mathbb T_q$ by \eqref{eq:stopping-identity-index}.
\end{proof}

\emph{The identity of the step.} We can now complete the step of Section~\ref{sec:inductive-class}. On every labelled term, by \eqref{eq:exact-factorisation}, \eqref{eq:cross-exponent-definition}, Theorem~\ref{thm:flat-part}, and Proposition~\ref{prop:r-flat-part},
\begin{equation}\label{eq:step-identity}
    \prod_i\kappa^{(k)}(C_i, \chi_i)\,e^{\mathcal E_k}S^\circ(\ell) = \prod_i\kappa^{(k+1)}(C_i, \chi_i)\,e^{\mathcal E_{k+1}}\;,
\end{equation}
with $\mathcal E_{k+1} \coloneqq \mathcal E_k + E^\partial_{k+1}(\ell) + \sum_iE^\circ_{k+1}(C_i, T) + E^\times_{k+1} + \mathbf 1_{\mathbf R}(k + 1)\bigl(\sum_iE^{\circ,\mathbf R}_{k+1}(C_i, T) + E^{\mathbf R\times}_{k+1}\bigr)$, the last terms entering at an $\mathbf R$-step together with the flat part of the $\mathbf R$-gas, with the cross constant $E^{\mathbf R\times}_{k+1}$ of Supplement, Proposition~S.10(iv)(c), and $E_{\mathrm{const},k+1} \coloneqq E_{\mathrm{const},k} + \sum_iE^\circ_{k+1}(C_i, T) + E^\times_{k+1} + \mathbf 1_{\mathbf R}(k + 1)\bigl(\sum_iE^{\circ,\mathbf R}_{k+1}(C_i, T) + E^{\mathbf R\times}_{k+1}\bigr)$, which are numbers on each labelled term. The identity is exact, and it is the normalisation clause of the step.

\subsection{The Normalisation}\label{sec:normalisation}

The remainders of the steps are collected in the following bound, in which $P_{j-1}(\mathcal C) \coloneqq \sum_iP_{j-1}(C_i)$, $\abs{\mathcal K(\mathcal C)} \coloneqq \sum_i\abs{\mathcal K(C_i)}$, and $r_{j-1}(\mathcal C) \coloneqq \min_ir_{j-1}(C_i)$: for $1 \leq j \leq K'$,
\begin{equation}\label{eq:step-remainder}
\begin{split}
    \epsilon_j(\mathcal C, T) &\coloneqq \mathfrak a_j\Bigl[C^{\mathcal K}_E\bigl(P_{j-1}(\mathcal C) + \abs{\mathcal K(\mathcal C)}\bigr)e^{-2c^T_Er_{j-1}(\mathcal C)} + 3\bar C_ZP_{j-1}(\mathcal C)e^{-c^Z\iota_{T,j}} + C_{\mathrm{cr}}P_{j-1}(\mathcal C)e^{-\nu_\times a_j/M}\Bigr]\\
    &\quad + \mathbf 1_{\mathbf R}(j)\,C^R_EMR_je^{-p_0(g_j)/4}\Bigl[\bigl(P_j(\mathcal C) + \abs{\mathcal K(\mathcal C)}\bigr)e^{-c^R_Er_j(\mathcal C)/(MR_j)} + P_j(\mathcal C)e^{-c^R_E\iota_{T,j}/(MR_j)}\Bigr]\\
    &\quad + \mathbf 1_{\mathbf R}(j)\,C^{\mathbf R\times}_EP_j(\mathcal C)e^{-p_0(g_j)/4}e^{-\kappa a_j/(6MR_j)}\;.
\end{split}
\end{equation}
At $j = 1$ the first line bounds the remainder of the first step, by \eqref{eq:first-step-extraction}, by the periodic comparison of \cite[Appendix~F]{WilsonFirstStep}, and by the cross terms of \cite[Section~10]{WilsonFirstStep}. By \eqref{eq:flat-remainder}, \eqref{eq:cross-exponent}, \eqref{eq:r-flat-remainder}, and Supplement, Proposition~S.10(iv)(c), the constant part of the scalar exponent satisfies the bound of Definition~\ref{def:dressed-loop}(iv) at every step.

\begin{corollary}[Sum of the Remainders]\label{cor:remainder-sum}
Let $L$ be as in Proposition~\ref{prop:torus-consistency}, and let $\mathcal C$ be a loop ensemble at level $q$ whose loops are admissible, with $\gamma_q$ as in that proposition. For every $K > j_q$,
\begin{equation}\label{eq:remainder-sum}
    \sum_{j=1}^{K'}\epsilon_j(\mathcal C, T) \leq B_E(\mathcal C) \coloneqq \sum_{i=1}^{n}\Bigl[\Bigl(\frac{P(C_i)}{r(C_i)} + P_q(C_i)\Bigr)\bigl(C'_E\bar{\mathfrak a} + C''_ES^{\mathbf R}(\bar g)\bigr)\Bigr]\;,
\end{equation}
where $C'_E$ and $C''_E$ depend on the constants of the construction, on $\Lambda_*$, and on $c_{\mathcal C}$, nondecreasingly in $c_{\mathcal C}$, and $S^{\mathbf R}(\bar g) \coloneqq \sup_K\sum_{j \leq K}(MR_j)^2e^{-p_0(g_j)/4}$ is finite and tends to $0$ as $\bar g \to 0$. In particular $B_E(\mathcal C)$ does not depend on $K$.
\end{corollary}

\begin{proof}
We sum the five parts of \eqref{eq:step-remainder} over $j$, loop by loop, with $m \coloneqq K' - j + 1$, so that the lengths of the step $j - 1$ are those of the stopping step multiplied by $L^m$.

\emph{The corner part.} The exponent of this part contains the length $r_{j-1}(\mathcal C) = \min_ir_{j-1}(C_i)$ of the ensemble, which for a loop $C$ of the ensemble may be smaller than its own $r_{j-1}(C)$, so that we sum with $r(\mathcal C)$ throughout. The ratio $P_{j-1}(C)/r_{j-1}(\mathcal C) = P(C)/r(\mathcal C)$ does not depend on $j$, and $x \coloneqq r_{K'}(\mathcal C) \geq a_{K'}/2 \geq 1$ by \eqref{eq:lengths-at-step} for each loop, since $s_{K'}, \delta_{K'} \geq a_{K'}$ and $P_{K'} \geq 4a_{K'}$. Lemma~\ref{lem:scale-sums} with $c = 2c^T_E$, together with $xe^{-2c^T_Ex} \leq 1/(2ec^T_E)$, $e^{-2c^T_Ex} \leq 1$, and $\abs{\mathcal K(C)} \leq P(C)/(2r(C)) \leq P(C)/(2r(\mathcal C))$ from Lemma~\ref{lem:chord-arc}, gives
\begin{equation}\label{eq:corner-part-sum}
    \sum_{m \geq 1}\Bigl(\frac{P(C)}{r(\mathcal C)}xL^m + \abs{\mathcal K(C)}\Bigr)e^{-2c^T_ExL^m} \leq \frac{P(C)}{r(\mathcal C)}\Bigl(\frac{1}{2ec^T_E} + \frac{1}{c^T_E} + \frac12 + \frac{1}{2c^T_E}\Bigr)\;.
\end{equation}
Since every loop of an ensemble at level $q$ has $r(C_i) \geq \frac12uL^{-q}$, the same holds for $r(\mathcal C)$, and $P(C)/r(\mathcal C) \leq 2P_q(C)$, a term of $B_E(\mathcal C)$.

\emph{The torus part.} By the proof of Lemma~\ref{lem:torus-condition}, $\iota_{T,K'} \geq a_{K'}(L^{m_T+q} - L^{-3})/2 \geq a_{K'}$, and $\iota_{T,j} = L^{m-1}\iota_{T,K'}$, while $P_{j-1}(C) = P_q(C)a_{K'}L^m \leq P_q(C)L\,L^{m-1}\iota_{T,K'}$. Lemma~\ref{lem:scale-sums} with $x = \iota_{T,K'}$ and $c = c^Z$, and $c^Zx \geq 1$, give $\sum_mP_{j-1}(C)e^{-c^Z\iota_{T,j}} \leq 3P_q(C)Lxe^{-c^Zx} \leq 3P_q(C)L/(ec^Z)$.

\emph{The cross part.} With $a_j = a_{K'}L^{m-1}$ and $P_{j-1}(C) = P_q(C)La_j$, Lemma~\ref{lem:scale-sums} with $x = a_{K'}$ and $c = \nu_\times/M$, and $\nu_\times a_{K'}/M \geq 1$, which holds for $\gamma_q$ small since $a_{K'} \geq L^3MR_{K'}$ and $R_{K'} \geq (\log\hat g_{K'}^{-2})^r \geq (\log\gamma_q^{-2})^r$, give at most $3P_q(C)La_{K'}e^{-\nu_\times a_{K'}/M} \leq 3P_q(C)LM/(e\nu_\times)$.

\emph{The $\mathbf R$-part.} Since $P_j(C) = (P(C)/r(\mathcal C))r_j(\mathcal C)$ and $te^{-ct/(MR_j)} \leq MR_j/(ec)$ for $t \geq 0$, the first bracket is at most $(P(C)/r(\mathcal C))MR_j(1/(ec^R_E) + 1/2)$; since $\iota_{T,j} \geq a_j$ and $P_j(C) = P_q(C)a_j$, the second is at most $P_q(C)MR_j/(ec^R_E)$. The sum over $j$ is therefore at most $C^R_E(P(C)/r(\mathcal C) + P_q(C))(1/2 + 2/(ec^R_E))S^{\mathbf R}(\bar g)$, and $P(C)/r(\mathcal C) \leq 2P_q(C)$ as in the corner part. In the same way the cross part of the $\mathbf R$-steps, with $P_j(\mathcal C) = P_q(\mathcal C)a_j$ and $a_je^{-\kappa a_j/(6MR_j)} \leq 6MR_j/(e\kappa)$, sums to at most $6C^{\mathbf R\times}_EP_q(\mathcal C)S^{\mathbf R}(\bar g)/(e\kappa)$. The number $S^{\mathbf R}(\bar g)$ is finite and tends to $0$ with $\bar g$: by Lemma~\ref{lem:tube-couplings} and the definition of a tube, $g_j^{-2} \geq \hat g_j^{-2}/\rho \geq (2/3)(\bar g^{-2} + b'(K - j))$, with $n = K - j$, and $(MR_j)^2 \leq 4L^2M^2(\log g_j^{-2})^{2r}$ by Lemma~\ref{lem:tube-cubes}(i), so that $(MR_j)^2e^{-p_0(g_j)/4} \leq 4L^2M^2y^{2r}e^{-(A_0/4)y^{p_0}}$ at $y = \log g_j^{-2}$. The two factors move in opposite directions as $y$ grows, and we bound their product: the function $y \mapsto y^{2r}e^{-(A_0/4)y^{p_0}}$ is decreasing for $y^{p_0} \geq 8r/(A_0p_0)$, which for $\bar g$ small contains every $y \geq \log((2/3)\bar g^{-2})$, so the product is at most its value at $y = \log((2/3)(\bar g^{-2} + b'n))$. The resulting series over $n \geq 0$ converges because $p_0 > 1$, uniformly in $K$, with terms which tend to $0$ as $\bar g \to 0$.

The amplitudes are at most $\bar{\mathfrak a}$. Adding the five parts and using $P(C_i)/r(\mathcal C) \leq 2P_q(C_i)$ gives \eqref{eq:remainder-sum}, with the factor $3$ which this costs contained in $C'_E$ and $C''_E$.
\end{proof}

\begin{definition}[Normalisation]\label{def:normalisation}
The normalisation $\kappa_K^{(q)}(C, \chi)$ of a loop $C$ at level $q$ with the character $\chi$ is the product \eqref{eq:normalisation-product} of the step factors $\kappa_k(C, \chi)$ of Definition~\ref{def:step-factor}, $1 \leq k \leq K'$.
\end{definition}

The normalisation removes from the expectation a perimeter term and a term for each corner, and nothing else. Since the lengths of the segments scale with the steps, $N^{(j-1)}_v(C) = L^{-(j-1)}N^{(0)}_v(C)$, the definitions give the following description, which shows which divergences are removed.

\begin{proposition}[Structure of the Normalisation]\label{prop:normalisation-structure}
For a loop $C$ at level $q$ with the character $\chi$ and $K > j_q$,
\begin{equation}\label{eq:normalisation-structure}
    \log\kappa^{(q)}_K(C, \chi) = \sum_{v \in \mathcal V}N^{(0)}_v(C)\,m_{v,K} + \sum_{c \in \mathcal K(C)}z_{v_c,w_c,K}\;,
\end{equation}
with the numbers
\begin{equation}\label{eq:mass-and-corner}
    m_{v,K} \coloneqq \sum_{j=1}^{K'}L^{-(j-1)}\Bigl(\frac{\kappa^{(j)}_v}{\Pi_j} + \mathbf 1_{\mathbf R}(j)\mu^{\mathbf R}_{j,v}\Bigr)\;,
    \qquad
    z_{v,w,K} \coloneqq \sum_{j=1}^{K'}\bigl(\zeta^{(j)}_{v,w} + \mathbf 1_{\mathbf R}(j)\zeta^{\mathbf R}_{j,v,w}\bigr)\;,
\end{equation}
which depend on $K$, $\bar g$, $q$, the character, and the constants of the construction, but not on $C$ or on the torus, and satisfy $\abs{m_{v,K}} \leq \sum_{j \leq K'}L^{-(j-1)}(C_{S'}\mathfrak a_j + C^R_Ee^{-p_0(g_j)/4})$ and $\abs{z_{v,w,K}} \leq \sum_{j \leq K'}(C_\zeta\mathfrak a_j + C^R_Ee^{-p_0(g_j)/4})$.
\end{proposition}

\begin{proof}
Insert $N^{(j-1)}_v(C) = L^{-(j-1)}N^{(0)}_v(C)$ into \eqref{eq:step-factor} and sum over $j$. The bounds are Lemma~\ref{lem:coefficients}, its counterpart at the step $1$ \cite[(5.31), (6.5)]{WilsonFirstStep}, and Proposition~\ref{prop:r-flat-part}. The dependence on $q$ is through $K' = K - j_q$ only.
\end{proof}

The number $N^{(0)}_v(C)$ is the length of the segments of orientation $v$ in fine units, of order $L^K$, so that the first sum in \eqref{eq:normalisation-structure} is the self-energy of the loop, proportional to the physical length of each orientation divided by the lattice spacing, and the second sum, which grows at most like $\sum_jg_j^2(\log g_j^{-2})^{6p_0}$ in $K$, is the renormalisation of the corners.

\section{The Identity at the Stopping Scale and the Uniform Bound}\label{sec:stopping-scale}

We now stop the induction and read off what it says about the renormalised expectation. Two features of the problem seem to stand in the way of a bound which is uniform in $K$. For the normalisation $\prod_i\kappa^{(q)}_K(C_i, \chi_i)$ we have no bound, from above or from below, which is uniform in $K$. By Proposition~\ref{prop:normalisation-structure} its logarithm is a sum of the counts $N^{(0)}_v(C)$, which are proportional to the lengths of the loops in the units of the fine lattice and grow like $L^K$, multiplied by the numbers $m_{v,K}$, and the bounds of that proposition on $m_{v,K}$ do not make these products small or even bounded. Whether the normalisation in fact tends to zero or to infinity is not decided by these estimates; for a trivial character, for which the observable is the constant one, it is constant. What matters is that an estimate of $\mathbb E_K[W_{\mathcal C}]$ which is not exact in the normalisation is of no use once it is divided by a number about which only this much is known. The terms of the density with large fields, on the other hand, are not small by counting: in four dimensions the factor which a large field contributes ``does not give any positive power of $\varepsilon$'' \cite[p.~244]{BalabanCRE}, while the number of places at which a large field may occur grows like $\varepsilon_K^{-4}$. Both difficulties disappear at the stopping step $K'$. The normalisation is the same number in every term of the dressed loops, so that it can be divided out exactly; and the density at the step $K'$ is a function of the block field on the lattice $\mathbb T_{j_q}$, whose spacing $uL^{-j_q}$ does not depend on $K$, so that the terms with large fields can be compared with the term without them as whole integrals over a fixed space, at a price which is exponential in the number of sites of $\mathbb T_{j_q}$ and in nothing else. The outcome is an exact identity, Theorem~\ref{thm:stopping-identity}, and from it the bound, Theorem~\ref{thm:uniform-bound}.

Throughout this section $\mathcal C = ((C_1, \chi_1), \dots, (C_n, \chi_n))$ is an ensemble as in Theorem~\ref{thm:main}, and
\begin{equation}\label{eq:index-k-zero}
    K \geq K_0(\mathcal C) \coloneqq \max\bigl(j_q + 1,\, K_{\mathrm{tor}}(\mathcal C)\bigr)\;,
\end{equation}
so that $K' = K - j_q \geq 1$ and, by Corollary~\ref{cor:dressed-every-step}, the ensemble is dressed at the step $K'$. We use the units of the step $K'$, in which the lattice of the block field $V \coloneqq V_{K'}$ is $\mathbb T_{j_q}$ with the spacing $1$ and
\begin{equation}\label{eq:stopping-sites}
    N_q \coloneqq \frac{T}{uL^{-j_q}} = 2L^{m_T + j_q}
\end{equation}
sites per side, and the fine lattice $\mathbb T_K$ has the spacing $L^{-K'}$. Every loop of $\mathcal C$ is a path of bonds of $\mathbb T_q \subset \mathbb T_{j_q}$, with $P_{K'}(C_i) = P_q(C_i)L^3MR_{K'-1}$ bonds of $\mathbb T_{j_q}$ by Lemma~\ref{lem:stopping-index}, and we write $P_{K'}(\mathcal C) \coloneqq \sum_iP_{K'}(C_i)$. We write $g_{K'}$ for the running coupling of the step $K'$; it depends on $K$, but by \eqref{eq:tube-two-six} with $n - m = j_q$ it satisfies $g_{K'} \leq (1 + \beta_0)g_K = (1 + \beta_0)\bar g$, and by Lemma~\ref{lem:tube-couplings}(i) $\log g_{K'}^{-2} \leq \log\hat g_{K'}^{-2}$, where $\hat g_{K'} = \hat g_{K-j_q}$ does not depend on $K$. The parameters $\varepsilon_{K'}$, $R_{K'}$, $\alpha_{0,K'}$, $\alpha_{1,K'}$ of the step $K'$ are functions of $\hat g_{K'}$ (Definition~\ref{def:construction}), and they do not depend on $K$ either.

\subsection{The Identity}\label{sec:identity}

A \emph{labelled term} $\ell$ of $\varrho_{K'}$ is a term of the expansion \eqref{eq:density-expansion} at $k = K'$ with given labels: the sequences of domains, with the refined labels of Section~\ref{sec:balaban-first-step} at the first step, together with the labels of the $\mathbf R$-operations, a renormalised component being removed from the sequences \cite[(0.3), p.~176]{BalabanLFR1}. We write
\begin{equation}\label{eq:term-integral}
    \varrho_\ell(V) \coloneqq \chi_{K'}(\Omega_{K'})\,\mathbf T_{K'}(\ell)\bigl[\exp A_{K'}\bigr]\;,
    \qquad
    Z_\ell \coloneqq \int\varrho_\ell(V)\,\dd V\;,
\end{equation}
for the density of $\ell$ and its integral, $\dd V$ being the normalised Haar measure on $G^{\mathcal B}$ with $\mathcal B$ the set of the positive bonds of $\mathbb T_{j_q}$. For a function $F$ of $V$ and of the outer variables of $\ell$, and for $Z_\ell > 0$, we put
\begin{equation}\label{eq:term-expectation}
    \langle F\rangle_\ell \coloneqq \frac{1}{Z_\ell}\int\chi_{K'}(\Omega_{K'})\,\mathbf T_{K'}(\ell)\bigl[F\exp A_{K'}\bigr]\,\dd V\;.
\end{equation}
The term $\ell_0$ is the one whose large-field region $Z_{K'} = \Lambda_{K'}^c$ is empty; by the nesting \eqref{eq:domain-sequences} it has $\Omega_j = \Lambda_j = \mathbb T$ for every $j \leq K'$, and it collects, by \cite[(0.6), p.~176]{BalabanLFR1}, all the histories whose large-field components have been renormalised. For a labelled term $\ell$ we write $Z(\ell)$ for the set of the connected components of its large-field region at the step $K'$, so that $Z(\ell) = \emptyset$ if and only if $\ell = \ell_0$.

\begin{theorem}[Identity at the Stopping Scale]\label{thm:stopping-identity}
Under the hypotheses of Theorem~\ref{thm:main}, let $K \geq K_0(\mathcal C)$. The numbers $w_\ell \coloneqq Z_\ell/\sum_{\ell'}Z_{\ell'}$ are nonnegative and $\sum_\ell w_\ell = 1$. With $\widehat{\mathcal W}_\ell$ the dressed loop of the step $K'$ on $\ell$ divided by $\prod_i\kappa^{(q)}_K(C_i, \chi_i)$,
\begin{equation}\label{eq:stopping-identity}
    W_K(\mathcal C) = \sum_{\ell \,|\, Z_\ell > 0}w_\ell\,\bigl\langle\widehat{\mathcal W}_\ell\bigr\rangle_\ell\;,
    \qquad
    \widehat{\mathcal W}_\ell = e^{\mathcal E_{K'}(\ell)}\Bigl[\prod_{i=1}^{n}f^{(i)}_{K'} + \sum_{\mathfrak B \neq \emptyset}Q_{K'}(\mathfrak B)\,w_{K'}(\mathfrak B)\Bigr]\;.
\end{equation}
On the term $\ell_0$ the exponent $\mathcal E_{K'}(\ell_0) = E_{\mathrm{const},K'}$ is a number, and $\varrho_{\ell_0}$ is a function of the field $V$ on $\mathcal B$ alone. Consequently, with the measure $\mu_K$ on $G^{\mathcal B}$ of total mass $w_{\ell_0}$ defined by
\begin{equation}\label{eq:stopping-measure}
    \dd\mu_K(V) \coloneqq \frac{\varrho_{\ell_0}(V)\,\dd V}{\sum_\ell Z_\ell}\;,
\end{equation}
and with $W_{\mathcal C}(V)$ the observable \eqref{eq:ensemble-observable} of the field $V$ on $\mathbb T_{j_q}$,
\begin{equation}\label{eq:stopping-identity-leading}
    W_K(\mathcal C) = e^{E_{\mathrm{const},K'}}\int W_{\mathcal C}(V)\,\dd\mu_K(V) + r_K(\mathcal C)\;,
\end{equation}
where
\begin{equation}\label{eq:stopping-remainder}
    r_K(\mathcal C) \coloneqq e^{E_{\mathrm{const},K'}}\int\Bigl[\prod_{i=1}^{n}f^{(i)}_{K'} - W_{\mathcal C}(V) + \sum_{\mathfrak B \neq \emptyset}Q_{K'}(\mathfrak B)\,w_{K'}(\mathfrak B)\Bigr]\dd\mu_K(V) + \sum_{\ell \neq \ell_0 \,|\, Z_\ell > 0}w_\ell\,\bigl\langle\widehat{\mathcal W}_\ell\bigr\rangle_\ell\;.
\end{equation}
\end{theorem}

The identity is exact, and it is a convex combination: $W_K(\mathcal C)$ is an average, with the weights $w_\ell$, of expectations of the normalised dressed loops. The measure $\mu_K$ lives on the space $G^{\mathcal B}$, which does not depend on $K$, and the function $W_{\mathcal C}(V)$ which it integrates in the leading term does not depend on $K$ either; everything that moves with $K$ is in the measure, in the number $E_{\mathrm{const},K'}$, and in the remainder.

\begin{proof}
We first show that every term is nonnegative. At real configurations the test functions of the construction take values in $[0, 1]$ (Lemma~\ref{lem:soft-decomposition}(i)), the operation $\mathbf T_{K'}(\ell)$ consists of integrals against nonnegative measures (Haar measures, $\delta$-functions of the block averaging, Gaussian weights, and characteristic functions), and the effective action is real; Ba{\l}aban proves that the densities are positive, ``hence the denominators are positive, and the operation $\mathbb R$ is well defined'' \cite[p.~176]{BalabanLFR1}, and the terms which the $\mathbf R$-operation produces are such a term multiplied by the exponential of a real sum \cite[(0.6), p.~176]{BalabanLFR1}. Hence $Z_\ell \geq 0$, and \eqref{eq:term-expectation} is the expectation with respect to a probability measure on the variables of $\ell$. Since every operation of the construction preserves the integral of a density (Section~\ref{sec:two-integrals}), $\sum_\ell Z_\ell = \int\varrho_{K'}\,\dd V = Z_K > 0$, and the $w_\ell$ are nonnegative numbers with sum one.

Next we use the numerator. By \eqref{eq:numerator-integral} at $k = K'$ and Definition~\ref{def:dressed-loop},
\begin{equation}\label{eq:numerator-terms}
    Z_K(\mathcal C) = \int N_{K'}(V)\,\dd V = \sum_\ell\int\chi_{K'}(\Omega_{K'})\,\mathbf T_{K'}(\ell)\bigl[\mathcal W_{K'}(\ell)\exp A_{K'}\bigr]\,\dd V = \sum_{\ell \,|\, Z_\ell > 0}Z_\ell\,\bigl\langle\mathcal W_{K'}(\ell)\bigr\rangle_\ell\;,
\end{equation}
where a term with $Z_\ell = 0$ contributes nothing, since it is the integral of the function $\mathcal W_{K'}(\ell)$, which is bounded on the support of the term by Lemma~\ref{lem:term-bound} below, against a nonnegative measure of total mass zero. The normalisation in front of the dressed loop is the same number on every term, namely $\prod_i\kappa^{(K')}(C_i, \chi_i) = \prod_i\kappa^{(q)}_K(C_i, \chi_i)$ by \eqref{eq:dressed-loop} and \eqref{eq:normalisation-product}, and it is not zero. Dividing \eqref{eq:numerator-terms} by $Z_K$ and by this number, and using \eqref{eq:two-integrals} and Definition~\ref{def:renormalised-expectation}, we obtain \eqref{eq:stopping-identity}.

On the term $\ell_0$ there are no large-field data at any step. The carrier of the loop's integral is therefore the whole torus at every step, the distances $d^{(j)}(x)$ in \eqref{eq:scalar-bounds} are infinite, and every local term $e^{(j)}_X$ vanishes, so that $\mathcal E_{K'}(\ell_0) = E_{\mathrm{const},K'}$; for $j \geq 2$ this is also the statement $E^\partial_j(\ell_0) = 0$ of the definition of $E^\partial_j$ in Section~\ref{sec:reference-gas}. The operation $\mathbf T_{K'}(\ell_0)$ holds no integrations, so that $\varrho_{\ell_0}(V) = \chi_{K'}(\mathbb T)\exp A_{K'}(U_{K'}(V))$, where the $\mathbf R$-part and the boundary terms left by renormalised components are functions of $V$ \cite[(0.6), p.~176]{BalabanLFR1}, \cite[p.~392]{BalabanLFR2}. The insertions on $\ell_0$ are functions of $V$ as well, since there are no outer variables. Separating the term $\ell_0$ in \eqref{eq:stopping-identity} and adding and subtracting $e^{E_{\mathrm{const},K'}}\int W_{\mathcal C}\,\dd\mu_K$ gives \eqref{eq:stopping-identity-leading} with the remainder \eqref{eq:stopping-remainder}.
\end{proof}

\subsection{The Coarse Loop}\label{sec:coarse-loop}

On the term $\ell_0$ the leading words $f^{(i)}_{K'} = \chi_i(U_{K'}(V)(C_i))$ are the Wilson loops of the background field, a field on the fine lattice which depends on $K$, while the leading term of \eqref{eq:stopping-identity-leading} contains the Wilson loops of $V$ itself. The two differ by little, because $V$ is the $K'$-fold block average of $U_{K'}(V)$ and $U_{K'}(V)$ is close to flat on the scale of a plaquette of $\mathbb T_{j_q}$. To quantify this we need a comparison between an iterated average and a transport, which we now prove.

The average of \cite[(0.4), p.~253]{BalabanRG1} maps a field $U$ on a lattice to the field $\bar U = M(U)$ on the lattice whose spacing is $L$ times larger. With $d = 4$, let $B(y)$ be the set of the $L^d$ points of the finer lattice in the cube of side $L$ centred at $y$ \cite[(0.3), p.~252]{BalabanRG1}, $\mathbf G(y, x)$ the family of the shortest paths from $y$ to $x$ which run through the coordinate directions in some order, one path for each order, $[x, x']$ the translate of $c$ to $x$, and $U(c)$ the transport along the straight path of $L$ bonds of $c$ \cite[p.~252]{BalabanRG1}. For $x \in B(c_-)$, $\Gamma \in \mathbf G(c_-, x)$, and $\Gamma' \in \mathbf G(c_+, x')$ put $\ell(x, \Gamma, \Gamma') \coloneqq \Gamma \cup [x, x'] \cup (-\Gamma') \cup (-c)$, a closed path at $c_-$, and $v(x, \Gamma, \Gamma') \coloneqq L^{-d}\abs{\mathbf G(c_-, x)}^{-1}\abs{\mathbf G(c_+, x')}^{-1}$. Then
\begin{equation}\label{eq:average}
    \bar U(c) = \exp\Bigl[i\sum_{x, \Gamma, \Gamma'}v(x, \Gamma, \Gamma')\,\frac1i\log U\bigl(\ell(x, \Gamma, \Gamma')\bigr)\Bigr]U(c)\;,
\end{equation}
and the weights $v(x, \Gamma, \Gamma')$ are nonnegative with sum one. We write $M^j$ for the $j$-fold iterate of $M$, the average of order $j$, as the averages of higher order are defined recursively in \cite[(43), p.~24]{BalabanAVG}, and $T_c(U)$ for the transport of $U$ along the straight path of the fine bonds of a bond $c$ of a coarser lattice. Here and below $\norm{\cdot}$ is the operator norm, $\mathfrak g$ is the Lie algebra of $G$ realised as the Hermitian matrices $X$ with $e^{iX} \in G$, $r_G \in (0, 1/10]$ is a radius such that the principal logarithm of every $W \in G$ with $\norm{W - 1} \leq r_G$ lies in $i\mathfrak g$, which exists because the exponential map is a local diffeomorphism at $0$, and for a unitary representation $\pi$ of $G$ and $Y \in \mathfrak g$ we write $\pi(e^{iY}) = e^{i\pi_*(Y)}$, so that $\pi_*(Y) = -i\,\dd\pi(iY)$ and, with the number $c_\pi$ of \eqref{eq:representation-constant}, $\norm{\pi(e^{iY}) - 1} \leq \norm{\pi_*(Y)} \leq c_\pi\norm{Y}$.

Three elementary facts about unitary matrices will be used. For unitary $W$ and $g$, $\norm{gWg^{-1} - 1} = \norm{W - 1} = \norm{W^{-1} - 1}$, and for any matrices $a_1, \dots, a_m$,
\begin{equation}\label{eq:product-deviation}
    \Bigl\lVert\prod_{l=1}^{m}(1 + a_l) - 1\Bigr\rVert \leq \prod_{l=1}^{m}\bigl(1 + \norm{a_l}\bigr) - 1\;.
\end{equation}
Secondly, if $W$ is unitary and $\norm{W - 1} \leq e^s - 1$ with $0 \leq s \leq 1/10$, then the principal logarithm $W = e^{iX}$ has
\begin{equation}\label{eq:log-bound}
    \norm{X} \leq 1.06\,s\;:
\end{equation}
an eigenvalue $e^{i\phi}$ of $W$ with $\phi \in (-\pi, \pi]$ has $2\abs{\sin(\phi/2)} = \abs{e^{i\phi} - 1} \leq e^s - 1 \leq 1.052\,s$, because $(e^s - 1)/s$ is increasing and equals $1.0517\dots$ at $s = 1/10$, and $2\arcsin(y/2) \leq y(1 - y^2/4)^{-1/2} \leq 1.002\,y$ for $0 \leq y \leq 0.11$. Thirdly, $\norm{e^{iX} - 1} \leq \norm{X}$ for Hermitian $X$, since $\abs{e^{i\phi} - 1} \leq \abs{\phi}$; consequently, for $X, F \in \mathfrak g$ with $s \coloneqq \norm{X} + \norm{F} \leq \min(1/10, \log(1 + r_G))$, \eqref{eq:product-deviation} gives $\norm{e^{iX}e^{iF} - 1} \leq e^s - 1 \leq r_G$, so that $e^{iX}e^{iF} = e^{iE}$ with $E \in \mathfrak g$ and $\norm{E} \leq 1.06\,s$ by \eqref{eq:log-bound}.

\begin{lemma}[The Coarse Loop]\label{lem:coarse-loop}
Let $L \geq 21$ be odd and $k \geq 1$, let $a_4 \coloneqq 7.2$, and let $U$ be a $G$-valued field on a periodic lattice of spacing $\eta$, with periods which are multiples of $L^k\eta$, such that $\norm{U(\partial p) - 1} \leq t$ for every plaquette $p$ and
\begin{equation}\label{eq:chart-condition}
    a_4L^{2k}t \leq \min\bigl(1/10,\ \log(1 + r_G)\bigr)\;.
\end{equation}
\begin{enumerate}[label=(\roman*)]
\item For every $0 \leq j \leq k$ and every bond $c$ of the lattice of spacing $L^j\eta$ there is $E_c \in \mathfrak g$ with $(M^jU)(c) = e^{iE_c}T_c(U)$ and $\norm{E_c} \leq a_4L^{2j}t$.
\item For every closed path $C$ of $P$ bonds of the lattice of spacing $L^k\eta$ and every character $\chi = (\dim\pi)^{-1}\operatorname{Tr}\pi$ of an irreducible unitary representation $\pi$,
\begin{equation}\label{eq:coarse-loop}
    \bigl\lvert\chi(U(C)) - \chi\bigl((M^kU)(C)\bigr)\bigr\rvert \leq e^{c_\pi a_4PL^{2k}t} - 1\;.
\end{equation}
\end{enumerate}
\end{lemma}

The constant $a_4 = 0.3(d - 1)(d + 4)$ at $d = 4$ does not grow with $j$: each level of averaging loses only the factor $(d + 3)/L$ on the deviation of the previous level, and the deviation of the plaquettes of the new level, which is $L^2$ times larger, is what sets the scale.

\begin{proof}
(i) We argue by induction on $j$, the case $j = 0$ being $E_c = 0$. Let $0 \leq j < k$, assume (i) at the level $j$, and write $W \coloneqq M^jU$, a $G$-valued field on the lattice of spacing $L^j\eta$, which we use as the unit of length in this step. Let $c$ be a bond of the lattice of spacing $L^{j+1}\eta$, from $c_-$ to $c_+ = c_- + Le_\mu$, so that $(M^{j+1}U)(c) = e^{iX_c}W(c)$ by \eqref{eq:average}, with $X_c$ the convex combination in the exponent and $W(c)$ the product of $W$ over the $L$ bonds $b$ of $c$.

\emph{The transport.} By the hypothesis of the induction, $W(c) = \prod_{b \subset c}e^{iE_b}T_b(U) = R_c\,T_c(U)$ with $R_c \coloneqq \prod_{b \subset c}g_be^{iE_b}g_b^{-1} \in G$, where $g_b$ is the transport of $U$ along the part of $c$ before $b$. By \eqref{eq:product-deviation} and the third fact, $\norm{R_c - 1} \leq e^{s_c} - 1$ with $s_c \coloneqq a_4L^{2j+1}t$, and $s_c \leq a_4L^{2k-1}t \leq \min(1/10, \log(1 + r_G))/L$ by \eqref{eq:chart-condition}, since $2j + 1 \leq 2k - 1$. Hence $R_c = e^{iF_c}$ with $F_c \in \mathfrak g$ and $\norm{F_c} \leq 1.06\,a_4L^{2j+1}t$ by \eqref{eq:log-bound}.

\emph{The loops of the average.} Let $x = c_- + n$ with $n \in \mathbb Z^d$ and $\abs{n_\nu} \leq (L - 1)/2$, which is the condition $x \in B(c_-)$ for odd $L$, let $\Gamma \in \mathbf G(c_-, x)$ and $\Gamma' \in \mathbf G(c_+, x')$, and let $\ell \coloneqq \Gamma \cup [x, x'] \cup (-\Gamma') \cup (-c)$. The loop $\ell$ has at most $(d + 2)L$ bonds: $\Gamma$ and $\Gamma'$ have $\abs{n}_1 \leq d(L - 1)/2$ each, and $[x, x']$ and $c$ have $L$ each. As above, $W(\ell) = R_\ell\,U(\ell)$, where $U(\ell)$ is the transport of $U$ around the fine loop obtained by replacing each bond of $\ell$ by its straight fine path, and $\norm{R_\ell - 1} \leq e^{(d+2)a_4L^{2j+1}t} - 1$. To bound $U(\ell)$ we contract $\ell$ across plaquettes of the lattice of spacing $L^j\eta$. Let $\Gamma_+$ be the translate of $\Gamma$ by $Le_\mu$, which lies in $\mathbf G(c_+, x')$. The paths $\Gamma'$ and $\Gamma_+$ have the same displacement $n$, and passing from the order of the directions in $\Gamma'$ to that in $\Gamma_+$ by transpositions of adjacent straight segments, each pair of directions $\alpha \neq \beta$ is transposed at most once and the transposition sweeps $\abs{n_\alpha}\abs{n_\beta}$ plaquettes. This replaces $\Gamma'$ by $\Gamma_+$; we then move $[x, x']$ back along $\Gamma$ one bond at a time, which sweeps $L$ plaquettes for a bond of direction $\pm e_\nu$ with $\nu \neq \mu$, and none for a bond of direction $\pm e_\mu$, after which the loop is trivial. The number of plaquettes swept is therefore at most
\begin{equation}\label{eq:loop-area}
\begin{split}
    A_L &\coloneqq \sum_{\alpha < \beta}\abs{n_\alpha}\abs{n_\beta} + L\bigl(\abs{n}_1 - \abs{n_\mu}\bigr)\\
    &\leq \frac{d(d - 1)}{2}\Bigl(\frac{L - 1}{2}\Bigr)^2 + \frac{(d - 1)L^2}{2} \leq \frac{(d - 1)(d + 4)}{8}L^2 = \frac{a_4}{2.4}L^2\;.
\end{split}
\end{equation}
Each sweep multiplies the transport around the loop, at its base point, by a conjugate of the transport of $U$ around the boundary of a plaquette of the lattice of spacing $L^j\eta$, and the latter is an ordered product of conjugates of the transports around the $L^{2j}$ fine plaquettes inside it. By \eqref{eq:product-deviation} and the first fact, $\norm{U(\ell) - 1} \leq (1 + t)^{A_LL^{2j}} - 1 \leq e^{A_LL^{2j}t} - 1$, and therefore $\norm{W(\ell) - 1} \leq e^{s_\ell} - 1$ with
\begin{equation}\label{eq:loop-deviation}
    s_\ell \coloneqq A_LL^{2j}t + (d + 2)a_4L^{2j+1}t \leq \Bigl(\frac{1}{2.4} + \frac{d + 2}{L}\Bigr)a_4L^{2j+2}t \leq 0.71\,a_4L^{2(j+1)}t \leq a_4L^{2k}t\;,
\end{equation}
where we used \eqref{eq:loop-area} and $(d + 2)/L \leq 6/21$. By \eqref{eq:chart-condition}, $s_\ell \leq \min(1/10, \log(1 + r_G))$, so that $\norm{W(\ell) - 1} \leq r_G$; since $W(\ell) \in G$, its principal logarithm lies in $i\mathfrak g$, and $\norm{(1/i)\log W(\ell)} \leq 1.06\,s_\ell$ by \eqref{eq:log-bound}. The convex combination $X_c$ therefore lies in $\mathfrak g$ and $\norm{X_c} \leq 1.06\,(A_LL^{2j} + (d + 2)a_4L^{2j+1})t$.

\emph{The new level.} We have $(M^{j+1}U)(c) = e^{iX_c}e^{iF_c}T_c(U)$, and
\begin{equation}\label{eq:new-level}
\begin{split}
    \norm{X_c} + \norm{F_c} &\leq 1.06\bigl(A_LL^{2j} + (d + 3)a_4L^{2j+1}\bigr)t\\
    &\leq 1.06\Bigl(\frac{1}{2.4} + \frac{d + 3}{L}\Bigr)a_4L^{2(j+1)}t \leq 0.795\,a_4L^{2(j+1)}t\;,
\end{split}
\end{equation}
since $(d + 3)/L \leq 1/3$ for $L \geq 21$. The right side is at most $0.795\min(1/10, \log(1 + r_G))$ by \eqref{eq:chart-condition}, so that by the third fact $e^{iX_c}e^{iF_c} = e^{iE_c}$ with $E_c \in \mathfrak g$ and $\norm{E_c} \leq 1.06 \cdot 0.795\,a_4L^{2(j+1)}t \leq a_4L^{2(j+1)}t$. This is (i) at the level $j + 1$; in particular $M^{j+1}U$ is $G$-valued, as the induction requires.

(ii) Let $c_1, \dots, c_P$ be the bonds of $C$ in their order, put $T_l \coloneqq T_{c_l}(U)$ and $H_l \coloneqq T_1 \cdots T_{l-1}$, and use (i) at $j = k$, with $E_{c^{-1}} = -T_c^{-1}E_cT_c$ for a bond traversed against its orientation, which has the same norm. Then
\begin{equation}\label{eq:coarse-holonomy}
    (M^kU)(C) = \prod_{l=1}^{P}e^{iE_{c_l}}T_l = \Bigl(\prod_{l=1}^{P}H_le^{iE_{c_l}}H_l^{-1}\Bigr)\,U(C)\;,
\end{equation}
and since $\abs{(\dim\pi)^{-1}\operatorname{Tr}((\pi(R) - 1)\pi(H))} \leq \norm{\pi(R) - 1}$ for unitary $\pi(H)$, \eqref{eq:product-deviation} and $\norm{\pi(e^{iE}) - 1} \leq c_\pi\norm{E}$ give $\abs{\chi(U(C)) - \chi((M^kU)(C))} \leq \prod_l(1 + c_\pi a_4L^{2k}t) - 1 \leq e^{c_\pi a_4PL^{2k}t} - 1$.
\end{proof}

We apply the lemma at the stopping step. Ba{\l}aban states the regularity of the background in a norm $\abs{\cdot}$ on matrices; we let $c_h \geq 1$ be a constant with $c_h^{-1}\abs{Y} \leq \norm{Y} \leq c_h\abs{Y}$ for all $N \times N$ matrices $Y$, and $B_3$ is the constant of \cite[Theorem~1, p.~279]{BalabanVAR}, which depends on $d$ and $L$ only.

\begin{corollary}[The Background at the Stopping Step]\label{cor:background-stopping}
Let $\bar g$ be so small that
\begin{equation}\label{eq:chart-stopping}
    2a_4c_h(1 + 2\kappa_s)B_3\,\varepsilon_{K'} \leq \min\bigl(1/10,\ \log(1 + r_G)\bigr)\;.
\end{equation}
Then on the term $\ell_0$, for every real $V$ in the support of $\varrho_{\ell_0}$ and every $i$,
\begin{equation}\label{eq:background-stopping}
    \bigl\lvert f^{(i)}_{K'} - W_{\chi_i}(C_i; V)\bigr\rvert \leq e^{x(C_i)} - 1\;,
    \qquad
    x(C) \coloneqq 2c_\pi a_4c_h(1 + 2\kappa_s)B_3\,P_{K'}(C)\,\varepsilon_{K'}\;,
\end{equation}
and $\abs{\prod_if^{(i)}_{K'} - W_{\mathcal C}(V)} \leq \sum_i(e^{x(C_i)} - 1)$. The numbers $x(C_i)$ do not depend on $K$.
\end{corollary}

\begin{proof}
On $\ell_0$ the determining set is the whole lattice $\mathbb T_{j_q}$, so that the background field $U_{K'} = U_{K'}(V)$ is a minimiser of the action on the fine fields $U$ with $M^{K'}(U) = V$ on every bond of $\mathbb T_{j_q}$ \cite[(2.11), (2.12), p.~256]{BalabanCRE}. Its plaquettes satisfy $\abs{\partial U_{K'} - 1} < 2B_3\varepsilon_{K'}(L^{-1}\eta)^2$ on $\Omega_{K'} = \mathbb T$ \cite[p.~268]{BalabanCRE}, where $\eta = L^{-(K'-1)}$ is the fine spacing in the units of the step $K' - 1$, so that $(L^{-1}\eta)^2 = L^{-2K'}$ is the area of a fine plaquette in the units of the step $K'$. This bound is drawn from passed test functions; for the soft test functions of the construction it holds with $\varepsilon_{K'}$ replaced by $(1 + 2\kappa_s)\varepsilon_{K'}$ (Supplement, Proposition~S.2). Hence the hypotheses of Lemma~\ref{lem:coarse-loop} hold with $k = K'$, $\eta = L^{-K'}$, and $t \coloneqq 2c_h(1 + 2\kappa_s)B_3\varepsilon_{K'}L^{-2K'}$, for which $a_4L^{2K'}t$ is the left side of \eqref{eq:chart-stopping}, and the periods of $\mathbb T_K$ are multiples of the spacing of $\mathbb T_{j_q}$. On $\ell_0$ the leading word is $f^{(i)}_{K'} = \chi_i(U_{K'}(C_i))$, and \eqref{eq:coarse-loop} with $P = P_{K'}(C_i)$ gives \eqref{eq:background-stopping}, because $c_\pi a_4P_{K'}(C_i)L^{2K'}t = x(C_i)$. For the product we telescope: $\abs{\prod_ia_i - \prod_ib_i} \leq \sum_i\abs{a_i - b_i}$ for complex numbers of modulus at most one, and $\abs{\chi_i(g)} \leq 1$ for $g \in G$. Finally $P_{K'}(C_i)$ and $\varepsilon_{K'}$ do not depend on $K$.
\end{proof}

\subsection{The Weight of the Large-Field Terms}\label{sec:large-field-weight}

We next show that the terms $\ell \neq \ell_0$ carry a small total weight, even after their integrals are multiplied by the factors which the dressed loop acquires near the large-field regions. Since $\mathbb T_{j_q}$ is fixed, this can be done by comparing whole integrals: the integral of $\ell_0$ is bounded below by the Haar volume of a ball of configurations, and the integral of every other term above by the supremum of its small-field exponent, each at a price exponential in $N_q^4$. The large-field factors of Ba{\l}aban's construction then pay for the history of each component, with half of their gain left over.

We need the following quantities, for a component $X \in Z(\ell)$. The \emph{history} of $X$ consists of the components of the large-field regions $Z_n$, $n \leq K'$, of the term which lie in $X$, together with the new large-field regions $Z'$ created at the steps $j' \leq K'$ which are merged into $X$; we write $\abs{Y}_n$ for the number of points of the scale $n$ in a set $Y$, $\#_n(Y) \coloneqq \abs{Y}_n(MR_n)^{-4}$ for the number of cubes of side $MR_n$ of the scale $n$ in a union $Y$ of such cubes, and $d'_n(Y)$ for its linear size in terms of these cubes \cite[p.~380]{BalabanLFR2}. Ba{\l}aban's factor \cite[(1.79), p.~383]{BalabanLFR2} contains for each new region $Z'$ created at the step $j'$ the gain $\exp(-\frac12\gamma_0A_1^2p_0^2(g_{j'})(d'_{j'}(Z') + 1))$, which holds for the construction of this paper with $\frac12\gamma_0$ replaced by $\frac12\gamma_0/\rho \geq \frac13\gamma_0$ (Proposition~\ref{prop:construction-admissible}); here $\gamma_0$, $A_1$, and $p_0(g)$ are those of \cite{BalabanLFR2}. The regions created at the first step carry instead the gain $\exp(-\frac12c_4\vartheta_V^2A_1^2(\log g_1^{-2})^{2p_0}(d'_1(Z') + 1))$ of the first step at the retreated region, with two constants $c_4$, $\vartheta_V$ of that step (Supplement, Proposition~S.7(iii)). We put
\begin{equation}\label{eq:history-quantities}
    P(X) \coloneqq \sum_{n \leq K'}L^4\abs{Z_n \cap X}_n\;,
    \qquad
    G(X) \coloneqq \sum_{Z'}\Upsilon_{j'}\bigl(d'_{j'}(Z') + 1\bigr)\;,
\end{equation}
with $\Upsilon_{j'} \coloneqq \frac16\gamma_0A_1^2p_0^2(g_{j'})$ for $j' \geq 2$ and $\Upsilon_1 \coloneqq \frac14c_4\vartheta_V^2A_1^2(\log g_1^{-2})^{2p_0}$, the \emph{history volume} of $X$ and half of its gains, and we let $C_P \geq 2$ be the constant of Appendix~\ref{app:constants}.

\begin{proposition}[The Weight of the Large-Field Terms]\label{prop:large-field-weight}
Let $\bar g$ be sufficiently small, depending on the constants of the construction, and let $K > j_q$. Then $Z_{\ell_0} > 0$, and for every choice of numbers $\omega(X)$, one for each component $X$ of each term $\ell \neq \ell_0$, which depend on $X$ and its history only and satisfy $1 \leq \omega(X) \leq \exp(C_PP(X) + \frac12G(X))$,
\begin{equation}\label{eq:large-field-weight}
    \sum_{\ell \neq \ell_0}Z_\ell\prod_{X \in Z(\ell)}\omega(X) \leq \delta_{\mathrm{LF}}Z_{\ell_0}\;,
    \qquad
    \delta_{\mathrm{LF}} \coloneqq e^{2C_K}\bigl(e^{s_{\mathrm{LF}}} - 1\bigr)\;,
\end{equation}
with
\begin{equation}\label{eq:large-field-constants}
    C_K \coloneqq C_{\mathrm{fs}}N_q^4\bigl(1 + \log g_{K'}^{-1}\bigr)\;,
    \qquad
    s_{\mathrm{LF}} \coloneqq 2N_q^4(MR_{K'})^{-4}\exp\bigl(-c_{\mathrm{LF}}\,p_0(g_{K'})\bigr)\;,
\end{equation}
$c_{\mathrm{LF}} \coloneqq (\rho(1 + \beta_0))^{-1}$, and $C_{\mathrm{fs}}$ the constant of \eqref{eq:final-scale-constant} below. In particular $1 - w_{\ell_0} \leq \delta_{\mathrm{LF}}$, and $\delta_{\mathrm{LF}} \leq \bar\delta_{\mathrm{LF}}$, where $\bar\delta_{\mathrm{LF}}$ is the number obtained by replacing $\log g_{K'}^{-1}$ by $\frac12\log\hat g_{K'}^{-2}$ and $p_0(g_{K'})$ by $p_0((1 + \beta_0)\bar g)$, which depends on $q$, $T/u$, $\bar g$, and the constants of the construction, and not on $K$.
\end{proposition}

\begin{proof}
We follow the terms through four steps: a common normalisation, an upper bound for each $Z_\ell$, a lower bound for $Z_{\ell_0}$, and the sum over the terms. We write $g_a \coloneqq g_{K'}$, and assume $g_a \leq 1$.

\emph{The normalisation.} Every term is written in Ba{\l}aban's completed form: for its components renormalised at the step $K'$ the missing terms of the effective action ``with localization domains intersecting the domain $Z$'' are added ``to complete the effective action'' and subtracted from the remaining expression \cite[p.~377, (1.70)]{BalabanLFR2}, and its surviving components are written with the normalised operations of \cite[(1.79), p.~383]{BalabanLFR2}, which holds for all large-field regions and not only for those which satisfy the conditions of the $\mathbf R$-operation \cite[p.~383--384]{BalabanLFR2}. A term then reads $\chi_{K'}\prod_{X \in Z(\ell)}\mathbf T'_{K'}(X)\exp A'_{K'}$, where the constants subtracted in the completion, ``the sum of normalization terms and vacuum energy terms with localizations intersecting the large field domain $Z$'' \cite[p.~378]{BalabanLFR2}, stand inside $\mathbf T'_{K'}(X)$. We divide every term by $e^{-E^{\ell_0}}$, with $E^{\ell_0}$ the constant of the effective action of $\ell_0$, a single number for all terms, which changes neither the $w_\ell$ nor \eqref{eq:large-field-weight}. The constant $E^\ell$ of a term is obtained by subtracting localised one-step vacuum energies from the initial constant \cite[p.~262]{BalabanCRE}, so that $E^{\ell_0} - E^\ell$ is the sum of the localised normalisation and vacuum-energy terms whose localisations meet the history of a component of $\ell$; a normalisation constant is $O(\log g_n^{-2})$ per point of the scale $n$ and the vacuum energies are $O(1)$ per point \cite[p.~380]{BalabanLFR2}, \cite[p.~264]{BalabanCRE}, and with a constant $C_\lambda$ of $d$, $G$, and these bounds which absorbs the enlargement of the localisations by a layer of cubes and their exponential tails,
\begin{equation}\label{eq:constant-difference}
    \abs{E^{\ell_0} - E^\ell} \leq \sum_{X \in Z(\ell)}\sum_{n \leq K'}C_\lambda\bigl(1 + \log g_n^{-2}\bigr)L^4\abs{Z_n \cap X}_n\;.
\end{equation}

\emph{The upper bound.} The action is nonnegative, so that $-A(1/g_{K'}^2, U_{K'}) \leq 0$. On the support of the term, the regular part, the $\mathbf R$-part, and the boundary part are bounded by \cite[(2.45), (2.46), (2.48), pp.~263--264]{BalabanCRE} by $C_{\mathcal E}\sum_n\abs{\Gamma_n}$, with $C_{\mathcal E} \coloneqq E_1(1 - L^{-\beta})^{-1} + 1 + 2B_\partial$ and $B_\partial$ the constant printed $B_1$ in \cite[(2.47)]{BalabanCRE}. The couplings of the completed action to the exterior of a renormalised component are at most $\alpha e^{-(1+2\beta)\kappa d_{k,Z}(Y)}$, with $\alpha$ arbitrarily small for $\gamma$ small \cite[(1.68), p.~377]{BalabanLFR2}, and summed over the localisations which meet a cube of side $MR_{K'}$ they are at most a constant $C_\partial \leq 1$, so that they contribute at most $C_\partial N_q^4$. The last layer $\Gamma_{K'}$ of the determining set lies in $\mathbb T_{j_q}$ and has at most $N_q^4$ points, and every earlier layer $\Gamma_n \subset \Omega_{n+1}^c \subset Z_{n+1}$ lies in a component of $Z(\ell)$, with $\abs{\Gamma_n \cap X}_n \leq L^4\abs{Z_{n+1} \cap X}_{n+1}$, since the large-field regions grow and merge until they are renormalised, and a renormalised component leaves the sequence \cite[(0.3), p.~176]{BalabanLFR1}. Together with \eqref{eq:constant-difference}, the exponent of the term, less the common constant, is at most $(C_{\mathcal E} + C_\partial)N_q^4 + \sum_{X \in Z(\ell)}P^c(X)$ with the \emph{history cost}
\begin{equation}\label{eq:history-cost}
    P^c(X) \coloneqq \sum_{n \leq K'}c_nL^4\abs{Z_n \cap X}_n\;,
    \qquad
    c_n \coloneqq C_{\mathcal E} + C_\lambda\bigl(1 + \log g_n^{-2}\bigr)\;.
\end{equation}
Since $\mathbf T'_{K'}(X)$ is a positive operation, $\mathbf T'_{K'}(X)F \leq (\mathbf T'_{K'}(X)1)\sup F$ for $F \geq 0$ \cite[(1.73), p.~380]{BalabanLFR2}, and integrating over $V$ with the normalised Haar measure,
\begin{equation}\label{eq:term-majorant}
    Z_\ell \leq M_\ell \coloneqq e^{C_K}\prod_{X \in Z(\ell)}\Bigl[\sup\mathbf T'_{K'}(X)1\cdot e^{P^c(X)}\Bigr]\;,
\end{equation}
where we used $(C_{\mathcal E} + C_\partial)N_q^4 \leq C_K$, which holds for the constant \eqref{eq:final-scale-constant} below.

\emph{The lower bound.} Let $\mathfrak B \coloneqq \{ V \,|\, \norm{V(b) - 1} < r_{\mathfrak B} \text{ for every } b \in \mathcal B \}$ with $r_{\mathfrak B} \coloneqq g_a/(10c_hB_3)$. The averaged extension $M^{\cdot}(Q^{s*}_{K'}V)$ of \cite[(2.16), p.~257]{BalabanCRE}, from which the minimisers are built, takes on each bond either the value $1$ or a value $V(c)^{\pm1}$: the field $Q^{s*}_{K'}V$ equals $1$ on the bonds inside the blocks and $V(c)$ on every bond which joins the two blocks of $c$ \cite[(1.3), p.~246]{BalabanCRE}, and in \eqref{eq:average} the paths $\Gamma$, $\Gamma'$ stay inside blocks, while $[x, x']$ and the straight path of a bond each cross the face between two blocks at most once. Hence on $\mathfrak B$ every plaquette of the block data is a product of at most four factors within $r_{\mathfrak B}$ of $1$, so that by \eqref{eq:product-deviation} its deviation is at most $c_h((1 + r_{\mathfrak B})^4 - 1) \leq 5c_hr_{\mathfrak B} = g_a/(2B_3)$ in the norm $\abs{\cdot}$. Then \cite[Theorem~1, p.~279]{BalabanVAR}, with its $\varepsilon_1 = g_a/(2B_3) \leq a_1$ for $g_a$ small, puts the local minimisers and the minimiser $U_{K'}(V)$ in the space of the configurations with $\abs{U(\partial p) - 1} < B_3\varepsilon_1\eta^2 = (g_a/2)\eta^2$ for every fine plaquette, $\eta = L^{-K'}$. So every test function of the term $\ell_0$ passes, because $g_a/2 \leq (\rho^{1/2}/2)\hat g_{K'} < \hat g_{K'}p_0(\hat g_{K'}) = \varepsilon_{K'}$ by Lemma~\ref{lem:tube-couplings}(i), and $\chi_{K'}(\mathbb T) = 1$ on $\mathfrak B$. For unitary $W$, $1 - \operatorname{Re}\tr W = \frac{1}{2N}\operatorname{Tr}((W - 1)^*(W - 1)) \leq \frac12\norm{W - 1}^2$, and there are $6N_q^4L^{4K'}$ fine plaquettes, so that on $\mathfrak B$
\begin{equation}\label{eq:trivial-term-action}
    \frac{1}{g_a^2}A\bigl(U_{K'}(V)\bigr) \leq \frac{1}{g_a^2}\cdot6N_q^4L^{4K'}\cdot\frac{c_h^2}{2}\Bigl(\frac{g_a}{2}\Bigr)^2L^{-4K'} = \frac34c_h^2N_q^4\;;
\end{equation}
and $\mathbf E_{K'} + \mathbf R_{K'} \geq -C_{\mathcal E}N_q^4$ by \cite[(2.45), (2.46)]{BalabanCRE} with $\sum_n\abs{\Gamma_n} = N_q^4$. The Haar measure of the ball of radius $r \leq 1$ about $1$ in $G$ is at least $(c_Gr)^{\dim G}$, with a constant $c_G$ of $G$, by the exponential coordinates, and $\mathcal B$ has $4N_q^4$ bonds. Therefore
\begin{equation}\label{eq:trivial-term-lower}
    Z_{\ell_0} \geq \exp\Bigl(-\frac34c_h^2N_q^4 - C_{\mathcal E}N_q^4\Bigr)\bigl(c_Gr_{\mathfrak B}\bigr)^{4N_q^4\dim G} \geq e^{-C_K}\;,
\end{equation}
with
\begin{equation}\label{eq:final-scale-constant}
    C_{\mathrm{fs}} \coloneqq 1 + c_h^2 + C_\partial + C_{\mathcal E} + 4\dim G\bigl(1 + \log_+(10c_hB_3/c_G)\bigr)\;,
\end{equation}
a constant of $d$, $L$, $G$, and the constants of Ba{\l}aban's theorems, independent of $K$; for the last inequality of \eqref{eq:trivial-term-lower} we used $\log(c_Gr_{\mathfrak B}) \geq -\log g_a^{-1} - \log_+(10c_hB_3/c_G)$ and $g_a \leq 1$.

\emph{The component factors.} Let $X \in Z(\ell)$, and let $\kappa^{\mathrm{tr}}(X)$ be the exact exponent of the factor of $X$ in \cite[(1.79)]{BalabanLFR2}, so that $\sup\mathbf T'_{K'}(X)1 \leq \exp(-\kappa^{\mathrm{tr}}(X) - (2/\rho)p_0(g_{j(X)}))$, with $j(X)$ the step of the first large-field region in $X$ \cite[p.~384]{BalabanLFR2}. By Supplement, Proposition~S.11, which is Ba{\l}aban's induction for \cite[(1.80), p.~384]{BalabanLFR2} with half of the gains set aside and the cost of each step raised,
\begin{equation}\label{eq:modified-kappa}
    \kappa^{\mathrm{tr}}(X) = \hat\kappa(X) + G(X) + H(X)\;,
    \qquad
    H(X) \coloneqq \sum_{n \leq K'}\lambda_n\sum_{Z''}M^4R_n^5d'_n(Z'')\;,
\end{equation}
with $\lambda_n \coloneqq 1 + 3 \cdot 2^3L^4C_\lambda(1 + \log g_n^{-2})$,
the inner sum over the components $Z''$ of $Z_n$ in $X$, where $\hat\kappa(X) \geq 0$, and $\hat\kappa(X) \geq \hat c_0\#_{K'}(X)$ if $X$ survives the stopping step, that is, if it fails one of the conditions of the $\mathbf R$-operation, with $\hat c_0 \coloneqq M^4R_{K'+1}^5(3 \cdot 2^3)^{-1}(2L^2)^{-4}$.

The history volume and the history cost are paid by $H(X)$ and half of $G(X)$:
\begin{equation}\label{eq:history-payment}
    P^c(X) + C_PP(X) \leq H(X) + \frac12G(X)\;,
\end{equation}
provided that, for every $n \leq K'$,
\begin{equation}\label{eq:history-condition}
    3 \cdot 2^3L^4\bigl(C_{\mathcal E} + C_P\bigr) \leq R_n
    \qquad\text{and}\qquad
    2^4L^4M^4R_n^4\bigl(c_n + C_P\bigr) \leq \frac12\Upsilon_n\;.
\end{equation}
Indeed, a component $Z''$ of $Z_n$ in $X$ with $d'_n(Z'') \neq 0$ has $\#_n(Z'') \leq 3 \cdot 2^3d'_n(Z'')$ \cite[p.~384]{BalabanLFR2}, so that its points cost $(c_n + C_P)L^4\abs{Z''}_n \leq 3 \cdot 2^3L^4(c_n + C_P)M^4R_n^4d'_n(Z'') \leq \lambda_nM^4R_n^5d'_n(Z'')$, by the first condition of \eqref{eq:history-condition} and $R_n \geq 1$. A component with $d'_n(Z'') = 0$ has at most $2^4$ cubes, because a tree of length zero is a point and at most $2^4$ cubes of a partition contain a given point. It contains no image $S(\cdot)$ of an earlier region, since $S$ adds ten layers of cubes \cite[p.~384]{BalabanLFR2}, so it consists of new regions created at the step $n$, whose gain enters $G(X)$ once; its points cost at most $2^4L^4M^4R_n^4(c_n + C_P) \leq \frac12\Upsilon_n$, which is half of the contribution of one new region to $G(X)$, by the second condition. The conditions \eqref{eq:history-condition} hold for $\bar g$ small, since $R_n \geq (\log g_n^{-2})^r$ and $R_n \leq 2L(\log g_n^{-2})^r$ by Lemma~\ref{lem:tube-cubes}, while $\Upsilon_n$ is a multiple of $(\log g_n^{-2})^{2p_0}$ and $2p_0 > 4r + 1$ (Appendix~\ref{app:constants}).

By \eqref{eq:modified-kappa}, \eqref{eq:history-payment}, and $p_0(g_{j}) \geq (1 + \beta_0)^{-1}p_0(g_a)$ for $j \leq K'$, which follows from $g_j \leq (1 + \beta_0)g_a$ by \eqref{eq:tube-two-six} for $\bar g$ small,
\begin{equation}\label{eq:component-factor}
    \sup\mathbf T'_{K'}(X)1\cdot e^{P^c(X)} \leq \exp\Bigl(-\hat\kappa(X) - C_PP(X) - \frac12G(X) - 2c_{\mathrm{LF}}\,p_0(g_a)\Bigr)\;.
\end{equation}
Let $C_{\mathrm{ent}}$ be a constant such that at most $C_{\mathrm{ent}}^m$ connected unions of $m$ cubes contain a given cube. If $X$ survives the stopping step, then $\hat\kappa(X) \geq \hat c_0\#_{K'}(X) \geq (\log C_{\mathrm{ent}} + 1)\#_{K'}(X)$ for $\bar g$ small, since $\hat c_0$ grows with $R_{K'+1}$. If $X$ is renormalised at the stopping step, then $\#_{K'}(X) \leq 100^4$ by the first condition of the $\mathbf R$-operation, and $c_{\mathrm{LF}}p_0(g_a) \geq (\log C_{\mathrm{ent}} + 1)100^4$ for $\bar g$ small. In both cases, for every $\omega(X) \leq \exp(C_PP(X) + \frac12G(X))$,
\begin{equation}\label{eq:component-weight}
    \omega(X)\sup\mathbf T'_{K'}(X)1\cdot e^{P^c(X)} \leq \exp\Bigl(-c_{\mathrm{LF}}\,p_0(g_a) - (\log C_{\mathrm{ent}} + 1)\#_{K'}(X)\Bigr)\;.
\end{equation}

\emph{The sum.} The components of $Z(\ell)$ are disjoint connected unions of the $N_q^4(MR_{K'})^{-4}$ cubes of side $MR_{K'}$ of $\mathbb T_{j_q}$, each of one of the two classes. The sum over the labels with given components and classes, the sequences of domains and the labels of the $\mathbf R$-operations, is contained in Ba{\l}aban's bound, since ``the summations over the admissible sequences can be replaced by the factors $\exp O(1)(MR_j)^{-d}|Z_j|$'' in the first exponential of \cite[(1.79)]{BalabanLFR2} \cite[p.~383]{BalabanLFR2}, and \eqref{eq:component-weight} is a bound on the resulting sum for each component. Dropping the disjointness of the components, and with the factor $2$ for the two classes,
\begin{equation}\label{eq:pattern-sum}
\begin{split}
    \sum_{\ell \neq \ell_0}M_\ell\prod_{X \in Z(\ell)}\omega(X) &\leq e^{C_K}\sum_{m \geq 1}\frac{1}{m!}\Bigl(\frac{2N_q^4}{(MR_{K'})^4}e^{-c_{\mathrm{LF}}p_0(g_a)}\sum_{m' \geq 1}C_{\mathrm{ent}}^{m'}e^{-(\log C_{\mathrm{ent}} + 1)m'}\Bigr)^m\\
    &\leq e^{C_K}\bigl(e^{s_{\mathrm{LF}}} - 1\bigr)\;,
\end{split}
\end{equation}
because $\sum_{m' \geq 1}e^{-m'} < 1$. With $Z_\ell \leq M_\ell$ from \eqref{eq:term-majorant} and $e^{C_K} \leq e^{2C_K}Z_{\ell_0}$ from \eqref{eq:trivial-term-lower}, this is \eqref{eq:large-field-weight}. Taking $\omega = 1$ gives $\sum_{\ell \neq \ell_0}Z_\ell \leq \delta_{\mathrm{LF}}Z_{\ell_0} \leq \delta_{\mathrm{LF}}\sum_\ell Z_\ell$, that is, $1 - w_{\ell_0} \leq \delta_{\mathrm{LF}}$. Finally $C_K$ increases with $\log g_a^{-1} \leq \frac12\log\hat g_{K'}^{-2}$, and $s_{\mathrm{LF}}$ decreases with $p_0(g_a) \geq p_0((1 + \beta_0)\bar g)$, while $N_q$ and $R_{K'}$ do not depend on $K$; so $\delta_{\mathrm{LF}} \leq \bar\delta_{\mathrm{LF}}$.
\end{proof}

At a fixed level the number $\bar\delta_{\mathrm{LF}}$ tends to zero with $\bar g$, although it is not needed for the bound. Indeed, $L^{j_q} = L^{q+3}MR_{K'-1}$ grows like $(\log\bar g^{-2})^r$, so that $C_K$ is of the order of $(\log\bar g^{-2})^{4r+1}$, while $c_{\mathrm{LF}}p_0((1 + \beta_0)\bar g)$ is of the order of $(\log\bar g^{-2})^{p_0}$ with $p_0 > 4r + 1$. This smallness is used in Section~\ref{sec:reflection-positivity}.

\subsection{The Uniform Bound}\label{sec:uniform-bound}

It remains to bound the normalised dressed loops on each term. On $\ell_0$ the bound is a number which does not depend on $K$, because the number of traversals at the stopping step is fixed; on the other terms the dressed loop acquires, in addition, the costs of the insertions attached to the large-field regions and of the traversals near them, and these are paid as in Proposition~\ref{prop:large-field-weight}. We put
\begin{equation}\label{eq:sup-amplitudes}
    \bar\rho \coloneqq C_\rho(1 + \beta_0)^2\bar g^2\Bigl(1 + A_1\bigl(\log((1 + \beta_0)^{-2}\bar g^{-2})\bigr)^{p_0}\Bigr)^4\;,
\end{equation}
$S_w \coloneqq \sum_{t \geq 1}(t + 1)^3e^{-\eta_wt}$, and $S_{\mathrm{wr}}(b) \coloneqq \sum_{i \geq 0}L^ie^{-bL^i}$ for $b > 0$, so that $\rho_j \leq \bar\rho$ for every $j \leq K$ when $\bar g$ is small, because $g_j \leq (1 + \beta_0)\bar g$ by \eqref{eq:tube-two-six} and $x \mapsto x(1 + A_1(\log x^{-1})^{p_0})^4$ is increasing for small $x$. With the constants of \eqref{eq:insertion-bound}, let
\begin{equation}\label{eq:insertion-total}
    I_*(\mathcal C) \coloneqq P_{K'}(\mathcal C)\Bigl[\frac{2c_\sigma C_\chi(1 + 2\kappa_s)B_3S_w\,\bar\rho\,\varepsilon_{K'}}{\min(\alpha_{0,K'}, \alpha_{1,K'})} + C^2\bar\rho\,S_{\mathrm{wr}}\Bigl(\frac{c_{\mathrm{att}}N_q}{M}\Bigr)\Bigr] + A^\times(\mathcal C)\;,
\end{equation}
where the last term is the total of the cross species,
\begin{equation}\label{eq:cross-total}
    A^\times(\mathcal C) \coloneqq P_{K'}(\mathcal C)\Bigl[10\,C_{\mathrm{cr}}\bar{\mathfrak a}\,e^{-\nu_\times d_\times/M} + 2C^{\mathbf R}_{\mathrm{cr}}S^{\mathbf R}(\bar g)\,e^{-\kappa d_\times/(24MR_{K'})}\Bigr]\;,
\end{equation}
with $d_\times = d_{K'} \geq a_{K'}$ the least $\ell^\infty$-distance between two loops of $\mathcal C$ in units of the scale $K'$, and $A^\times(\mathcal C) \coloneqq 0$ for a single loop. These numbers do not depend on $K$. The first two terms of \eqref{eq:insertion-total} are sums over the loops, and $A^\times$ does not decrease when a loop is added to an ensemble.

\begin{lemma}[The Total of the Cross Species]\label{lem:cross-total}
Let $\gamma_q$ be sufficiently small, depending on $c_{\mathcal C}$. On every labelled term and at every step $k \leq K'$, the groups of the cross species present at the step $k$ satisfy $\sum_We^{\theta(W)}\norm{Y_W} \leq A^\times(\mathcal C)$.
\end{lemma}

\begin{proof}
By \eqref{eq:cross-bound}, summed over $j \leq k$, the left side is at most $\sum_{j \leq K'}2^{K'-j+1}\mathsf A^\times_j$. In units of the scale $j$ the lengths of the stopping step are multiplied by $L^m$, $m \coloneqq K' - j$, so that $P_j(\mathcal C) = L^mP_{K'}(\mathcal C)$ and $d_j = L^md_\times$. We use the following consequence of Lemma~\ref{lem:scale-sums}: for $c > 0$ and $x > 0$ with $cx \geq 1$,
\begin{equation}\label{eq:doubling-sum}
    \sum_{m \geq 0}(2L)^me^{-cxL^m} \leq e^{-cx}\Bigl(1 + \frac{4}{cx}\Bigr) \leq 5e^{-cx}\;,
    \qquad
    (2L)^me^{-cxL^m} \leq e^{-cx/2}\;.
\end{equation}
Indeed $L^m - 1 \geq 2m$ and $cx \geq \log 2$ give $2^m \leq e^{cx(L^m - 1)/2}$, hence $(2L)^me^{-cxL^m} \leq e^{-cx/2}L^me^{-(c/2)xL^m}$. The first bound follows by the lemma with $c/2$ in place of $c$, after division by $x$, and the second from $te^{-(c/2)t} \leq 2/(ec)$ at $t = xL^m$ and $2/(ecx) \leq 1$. With $x = d_\times$ and $c = \nu_\times/M$, where $cx \geq \nu_\times a_{K'}/M \geq 1$ as in the proof of Corollary~\ref{cor:remainder-sum}, the first terms of the $\mathsf A^\times_j$ contribute at most $2C_{\mathrm{cr}}\bar{\mathfrak a}P_{K'}(\mathcal C) \cdot 5e^{-\nu_\times d_\times/M}$. For the second terms put $y \coloneqq \kappa d_\times/(12MR_{K'})$, so that $y \geq \kappa L^3/12 \geq 2\log(2L)$, because $d_\times \geq a_{K'} \geq L^3MR_{K'}$, $L \geq 195$, and $\kappa \geq 4^d$ as in Supplement, Proposition~S.10. For $m \geq 1$, \eqref{eq:tube-two-nine} gives $R_j \leq (L + 1)m^{\beta_0}R_{K'} \leq 2LmR_{K'}$, hence $\kappa d_j/(6MR_j) \geq yL^{m-1}/m \geq y(L/2)^{m-1}$, and for $m = 0$ the exponent is $2y$. Since $(L/2)^{m-1} - \frac12 \geq m$ for $m \geq 2$ and $L \geq 8$, and $y/2 \geq \log(2L)$, we have $(2L)^me^{-y(L/2)^{m-1}} \leq e^{-y/2}$ for every $m \geq 1$. The second terms therefore contribute at most $2C^{\mathbf R}_{\mathrm{cr}}P_{K'}(\mathcal C)e^{-y/2}\sum_j\mathbf 1_{\mathbf R}(j)e^{-p_0(g_j)/4}$, and the last sum is at most $S^{\mathbf R}(\bar g)$ of Corollary~\ref{cor:remainder-sum}, because $MR_j \geq 1$.
\end{proof}

For a component $X \in Z(\ell)$ we collect the costs of its history in
\begin{equation}\label{eq:attached-weight}
    \omega_A(X) \coloneqq \exp\Bigl(\sum_{n \leq K'}C_{A,n}\#_n(Z_n \cap X)\Bigr)\;,
    \qquad
    C_{A,n} \coloneqq 2C_A + 8C_e\Lambda_*MR_n\hat\tau_n^2\;,
\end{equation}
with $C_A$ the constant of \eqref{eq:attached-total}, $\hat\tau_n$ as in \eqref{eq:lifetime-cost}, and $C_e$ the constant of the cost of an exposed traversal in Section~\ref{sec:large-field-step}.

\begin{lemma}[The Terms at the Stopping Step]\label{lem:term-bound}
Let $\bar g$ satisfy \eqref{eq:chart-stopping}. For every labelled term $\ell$ of $\varrho_{K'}$ and every real configuration of its variables in the support of its indicators,
\begin{equation}\label{eq:term-bound}
    \bigl\lvert\widehat{\mathcal W}_\ell\bigr\rvert \leq e^{B_E(\mathcal C) + I_*(\mathcal C)}\prod_{X \in Z(\ell)}\omega_A(X)\;,
\end{equation}
with $B_E(\mathcal C)$ of Corollary~\ref{cor:remainder-sum}. On $\ell_0$, moreover, $\abs{\sum_{\mathfrak B \neq \emptyset}Q_{K'}(\mathfrak B)w_{K'}(\mathfrak B)} \leq e^{I_*(\mathcal C)} - 1$. If, for every $n \leq K'$,
\begin{equation}\label{eq:attached-condition}
    C_{A,n} \leq L^4(MR_n)^4\;,
\end{equation}
which holds for $\gamma_q$ small, depending on $c_{\mathcal C}$, then $1 \leq \omega_A(X) \leq \exp(C_PP(X))$.
\end{lemma}

\begin{proof}
We first bound the sum over the families. For a family $\mathfrak B$ the word $w_{K'}(\mathfrak B)$ is the normalised trace of a product of unitary transports, which contain the background and, on large-field regions, outer variables of the term, and of the insertions of $\mathfrak B$, so that $\abs{w_{K'}(\mathfrak B)} \leq \prod_{W \in \mathfrak B}\norm{Y_W}$ by Lemma~\ref{lem:word-norm}; with Definition~\ref{def:dressed-loop}(i), the sum over the finite families of groups of the term is bounded by a product over the groups. The groups of a term may be infinite in number, since their labels are not bounded; the product is then the limit of the products over finite sets of groups, which converge because $\sum_We^{\theta(W)}\norm{Y_W}$ is finite by the bounds below, and the sum over the families converges absolutely to a number bounded by it:
\begin{equation}\label{eq:family-sum}
    \Bigl\lvert\prod_if^{(i)}_{K'} + \sum_{\mathfrak B \neq \emptyset}Q_{K'}(\mathfrak B)\,w_{K'}(\mathfrak B)\Bigr\rvert \leq \prod_W\bigl(1 + e^{\theta(W)}\norm{Y_W}\bigr) \leq \exp\Bigl(\sum_We^{\theta(W)}\norm{Y_W}\Bigr)\;,
\end{equation}
where we used $\abs{f^{(i)}_{K'}} \leq 1$, and where the sum without the leading term is at most the right side minus one. Every group contains a traversal, so that the sum over the groups which are neither attached nor of the cross species is at most $\sum_{\mathfrak n}\nu^{(K')}_{\mathfrak n}$, the sum over the $P_{K'}(\mathcal C)$ traversals of the step $K'$, and the groups of the cross species contribute at most $A^\times(\mathcal C)$ on every term, by Lemma~\ref{lem:cross-total}.

\emph{The term $\ell_0$.} On $\ell_0$ there are no large-field data at the step $K'$, so that $d_{K'}(\mathfrak n) = \infty$ and $\mathfrak m_{K'}(\mathfrak n) = 0$ in \eqref{eq:insertion-bound}, and there are no attached groups, since the attached parts of the renormalised components have become deviation insertions of the class (Section~\ref{sec:large-field-step} and Supplement, Proposition~S.10). By \eqref{eq:deviation-scaling} at $k + 1 = K'$, with $\varepsilon_{K'}$ enlarged by the factor $1 + 2\kappa_s$ of the soft test functions as in the proof of Corollary~\ref{cor:background-stopping}, $\delta^*_{\mathfrak n}(U_{K'}) \leq 2c_\sigma(1 + 2\kappa_s)B_3\varepsilon_{K'}S_w$, and $\tilde\rho_{K'} \leq \bar\rho/\min(\alpha_{0,K'}, \alpha_{1,K'})$. The smallest period of the torus in units of the scale $j$ is $N^{(j)}_{\min} = N_qL^{K'-j}$, so that with $i = K' - j$ the last term of \eqref{eq:insertion-bound} is at most $C^2\bar\rho\sum_{i \geq 0}L^ie^{-c_{\mathrm{att}}N_qL^i/M} = C^2\bar\rho\,S_{\mathrm{wr}}(c_{\mathrm{att}}N_q/M)$. Summing over the traversals and adding the cross species, $\sum_{\mathfrak n}\nu^{(K')}_{\mathfrak n} + A^\times(\mathcal C) \leq I_*(\mathcal C)$, and \eqref{eq:family-sum} gives both assertions on $\ell_0$, since $\abs{e^{\mathcal E_{K'}(\ell_0)}} = e^{\operatorname{Re}E_{\mathrm{const},K'}} \leq e^{B_E(\mathcal C)}$ by Definition~\ref{def:dressed-loop}(iv) and Corollary~\ref{cor:remainder-sum}.

\emph{The other terms.} On a term $\ell \neq \ell_0$ three further contributions enter, each attached to a component of a large-field region $Z_n$, $n \leq K'$, which lies in a component $X \in Z(\ell)$. The groups attached at the step $n$ to a component $Z$ of $Z_n$, including the second and third terms of \eqref{eq:insertion-bound} at $n = K'$, have total weight at most $A_n(Z) \leq C_A(d'_n(Z) + 1) \leq 2C_A\#_n(Z)$ by \eqref{eq:attached-total}, where $d'_n(Z) + 1 \leq 2\#_n(Z)$ for a connected union of cubes, since a tree through the centres of adjacent cubes meets all of them and has edges of length at most twice the side. At the step $n$ the traversals within a cube of side $MR_n$ of $Z_n$ are exposed, at most $4\Lambda_*MR_n$ per cube by Lemma~\ref{lem:square-properties}(i), since such a cube meets one loop by Lemma~\ref{lem:stopping-index}; each costs at most $C_e\hat\tau_n^2$, and so does each traversal in the local terms $e^{(n)}_X$ of $E^\partial_n(\ell)$ attached to $Z_n$ (Supplement, Proposition~S.9), which gives at most $8C_e\Lambda_*MR_n\hat\tau_n^2$ per cube. The constant part of $\mathcal E_{K'}(\ell)$ is bounded by $B_E(\mathcal C)$ as on $\ell_0$. These contributions add to the exponent in \eqref{eq:family-sum} and to $\operatorname{Re}\mathcal E_{K'}(\ell)$ at most $\sum_{X \in Z(\ell)}\sum_{n \leq K'}C_{A,n}\#_n(Z_n \cap X)$, which gives \eqref{eq:term-bound}.

Finally, $\omega_A(X) \geq 1$ trivially, and under \eqref{eq:attached-condition} $\sum_nC_{A,n}\#_n(Z_n \cap X) \leq \sum_nL^4\abs{Z_n \cap X}_n = P(X) \leq C_PP(X)$. The condition \eqref{eq:attached-condition} holds for $\gamma_q$ small, depending on $c_{\mathcal C}$, because both terms of $C_{A,n}$ are at most $\frac12L^4(MR_n)^4$ for $n \leq K'$. Indeed, $R_n \geq (\log\hat g_n^{-2})^r \geq (\log\gamma_q^{-2})^r$ by Definition~\ref{def:tube-parameters} and \eqref{eq:stopping-coupling}, so that $4C_A \leq L^4(MR_n)^4$ once $L^4M^4(\log\gamma_q^{-2})^{4r} \geq 4C_A$. Moreover $R_n \leq 2L(\log g_n^{-2})^r$ by Lemma~\ref{lem:tube-cubes}(i) and $\hat\tau_n^2 \leq Cg_n^2(\log g_n^{-2})^{6p_0}$, so that $16C_e\Lambda_*MR_n\hat\tau_n^2 \leq C'\gamma_q^2(\log\gamma_q^{-2})^{6p_0+r}$, since $x \mapsto x(\log x^{-1})^{6p_0+r}$ is increasing for small $x$, and the right side is at most $1 \leq L^4(MR_n)^4$ for $\gamma_q$ small.
\end{proof}

\begin{theorem}[Uniform Bound]\label{thm:uniform-bound}
Under the hypotheses of Theorem~\ref{thm:main}, let $\bar g$ be so small that \eqref{eq:chart-stopping} and the hypothesis of Proposition~\ref{prop:large-field-weight} hold, and $\gamma_q$ so small that \eqref{eq:attached-condition} holds. For every $K \geq K_0(\mathcal C)$,
\begin{equation}\label{eq:uniform-bound}
    \abs{W_K(\mathcal C)} \leq B(\mathcal C) \coloneqq e^{B_E(\mathcal C) + I_*(\mathcal C)}\bigl(1 + \bar\delta_{\mathrm{LF}}\bigr)\;,
\end{equation}
and the remainder \eqref{eq:stopping-remainder} satisfies
\begin{equation}\label{eq:remainder-bound}
    \abs{r_K(\mathcal C)} \leq e^{B_E(\mathcal C)}\Bigl[\sum_{i=1}^{n}\bigl(e^{x(C_i)} - 1\bigr) + e^{I_*(\mathcal C)} - 1\Bigr] + e^{B_E(\mathcal C) + I_*(\mathcal C)}\,\bar\delta_{\mathrm{LF}}\;.
\end{equation}
The number $B(\mathcal C)$ does not depend on $K$; it depends on $\mathcal C$ through $\sum_i(P(C_i)/r(C_i) + P_q(C_i))$, through $P_q(\mathcal C)$, and through the characters, and on $q$, $T/u$, $\bar g$, and the constants of the construction.
\end{theorem}

\begin{proof}
By \eqref{eq:stopping-identity}, the triangle inequality, and Lemma~\ref{lem:term-bound},
\begin{equation}\label{eq:uniform-bound-split}
    \abs{W_K(\mathcal C)} \leq e^{B_E(\mathcal C) + I_*(\mathcal C)}\Bigl(w_{\ell_0} + \sum_{\ell \neq \ell_0}w_\ell\prod_{X \in Z(\ell)}\omega_A(X)\Bigr)\;,
\end{equation}
where we used that $\omega_A(X)$ is a number of the label, so that $\langle\prod_X\omega_A(X)\rangle_\ell = \prod_X\omega_A(X)$. By Lemma~\ref{lem:term-bound} the numbers $\omega(X) \coloneqq \omega_A(X)$ satisfy the hypothesis of Proposition~\ref{prop:large-field-weight}, which gives $\sum_{\ell \neq \ell_0}w_\ell\prod_X\omega_A(X) \leq \delta_{\mathrm{LF}}Z_{\ell_0}/\sum_{\ell}Z_\ell \leq \bar\delta_{\mathrm{LF}}$; together with $w_{\ell_0} \leq 1$ this is \eqref{eq:uniform-bound}. For the remainder, the integrand of the first term of \eqref{eq:stopping-remainder} is at most $\sum_i(e^{x(C_i)} - 1) + e^{I_*(\mathcal C)} - 1$ by Corollary~\ref{cor:background-stopping} and Lemma~\ref{lem:term-bound}, the measure $\mu_K$ has mass $w_{\ell_0} \leq 1$, and $\abs{e^{E_{\mathrm{const},K'}}} \leq e^{B_E(\mathcal C)}$; the second term is bounded as above. The dependence of $B(\mathcal C)$ is read from Corollary~\ref{cor:remainder-sum}, \eqref{eq:insertion-total}, and Proposition~\ref{prop:large-field-weight}: $P_{K'}(\mathcal C) = P_q(\mathcal C)L^3MR_{K'-1}$, and $N_q$, $\varepsilon_{K'}$, $R_{K'}$, $\alpha_{i,K'}$ are functions of $q$, $T/u$, $\bar g$, and the constants of the construction.
\end{proof}

The bound \eqref{eq:uniform-bound} is the first half of \eqref{eq:main-bound}. Two features of it deserve a comment. First, no condition on the coupling in this section involves the size of the loops or their number: \eqref{eq:chart-stopping} and the conditions of Proposition~\ref{prop:large-field-weight} concern the construction alone, and \eqref{eq:attached-condition} is a condition on $\gamma_q$ which involves the shape constant $\Lambda_*$ and the representations through $c_{\mathcal C}$. The size and the number of the loops enter only the value of $B(\mathcal C)$, through $B_E(\mathcal C)$ and $I_*(\mathcal C)$, which are sums over the loops. Secondly, the bound does not say that the large-field terms are small: $\bar\delta_{\mathrm{LF}}$ is a fixed number at a fixed coupling. What makes the terms with large fields harmless is that their weights, multiplied by the costs of the dressed loop, are summable against the term without large fields on the fixed lattice $\mathbb T_{j_q}$.

\section{The Difference of the Walk Factors}\label{sec:walk-factor-difference}

The limit $K \to \infty$ is obtained by comparing the construction for the target $(K, \bar g)$ with the construction for $(K + 1, \bar g)$. The two share every lattice but the finest, and on a shared lattice they perform the same operations, with the same parameters, on data which differ by little; what differs is the fine lattice, on which every background field, every propagator, and every term of the effective action is computed. This section and the next prove the two estimates of the comparison which are not contained in Ba{\l}aban's analysis of a single run. The present one concerns the propagators. The effective actions are built from the inverses of the quadratic forms of the fluctuation integrals, the propagators of \cite{BalabanPBG}, through their expansions into generalised random walks, and the comparison needs a bound on the difference between the terms of one walk in the two constructions which is small in the ratio of the fine spacing to the spacing of the scale on which the walk lives. The difficulty is that the two terms do not act on the same space: one maps sources on the lattice of spacing $\varepsilon$ to fields on that lattice, the other does the same on the lattice of spacing $\varepsilon/L$. We compare them through the one-step average and its right inverses. The comparison is then exact up to a bilinear defect, and this defect is small only between fields which are regular at the scale of their shell; it is the decay of the propagators across the shells which turns this into a bound of the order of a power of $\varepsilon/a_j$.

\subsection{Two Runs on Common Lattices}\label{sec:two-runs}

Throughout this section and the next two, $\mathcal C$ is an ensemble as in Theorem~\ref{thm:main}, $K \geq K_0(\mathcal C)$, and $K' = K - j_q$. \emph{Run~A} is the construction for the target $(K, \bar g)$, started at the bare coupling $g_0(K, \bar g)$ of Theorem~\ref{thm:coupling-flow} and stopped after $K'$ steps; \emph{run~B} is the construction for $(K + 1, \bar g)$, started at $g_0(K + 1, \bar g)$ and stopped after $K' + 1$ steps. Both are constructions for tubes with the same numbers $b$, $b'$, and $C_W$, which Theorem~\ref{thm:coupling-flow} provides independently of $K$. We write $\varepsilon \coloneqq \varepsilon_K = uL^{-K}$ for the fine spacing of run~A and
\begin{equation}\label{eq:shared-spacing}
    a_k \coloneqq \varepsilon L^k\;,
    \qquad\text{so that}\qquad
    \frac{\varepsilon}{a_k} = L^{-k}\;.
\end{equation}
The step $k$ of run~A maps a density of the field on the lattice of spacing $a_{k-1}$ to one on the lattice of spacing $a_k$, and so does the step $k + 1$ of run~B, because $(k + 1) - (K + 1) = k - K$. We call $k$ the \emph{shared scale} of these two steps, $0 \leq k \leq K'$; the step $1$ of run~B, which maps the lattice of spacing $\varepsilon/L$ to that of spacing $\varepsilon$, is its \emph{extra step}. The couplings at the shared scale $k$ are $g^A_k$ and $g^B_{k+1}$, and we put
\begin{equation}\label{eq:coupling-difference}
    \delta_k \coloneqq (g^B_{k+1})^{-2} - (g^A_k)^{-2}\;.
\end{equation}

\begin{lemma}[Shared Parameters]\label{lem:shared-parameters}
For every integer $k \leq K$ the numbers $u_k$ and $\ell_k$ of \eqref{eq:tube-couplings} of run~A are those of run~B at $k + 1$. Consequently:
\begin{enumerate}[label=(\roman*)]
\item every parameter of Definition~\ref{def:tube-parameters}, the depth $N_k$ of Definition~\ref{def:depth}, and every threshold, radius, cube size, and soft width of the construction at the shared scale $k$ are the same in both runs;
\item the stopping index $j_q$, the stopping lattice $\mathbb T_{j_q}$, and the number $N_q$ of \eqref{eq:stopping-sites} are the same in both runs, and run~B's stopping step is $K' + 1$;
\item for $0 \leq k \leq K'$, $\abs{\delta_k} \leq u_k - \ell_k \leq \frac12\min\bigl((g^A_k)^{-2}, (g^B_{k+1})^{-2}\bigr)$.
\end{enumerate}
\end{lemma}

\begin{proof}
By \eqref{eq:tube-couplings}, $u_k$ for run~A is $\bar g^{-2} + b'(K - k) + C_W$, and $u_{k+1}$ for run~B is $\bar g^{-2} + b'((K + 1) - (k + 1)) + C_W$, the same number; the same holds for $\ell$. (i) The parameters of Definition~\ref{def:tube-parameters} are functions of $\hat g_k = u_k^{-1/2}$, and $N_0(k)$ and $N'_k$ in Definition~\ref{def:depth} are defined through $\hat g_k$ and through the cube sizes $R_{k'}$ at the integers $k' \leq k$, which correspond to each other in the same way; the remaining parameters are, by Definition~\ref{def:construction}(i), functions of $\hat g_k$ as well. (ii) By Definition~\ref{def:stopping-index}, $j_q$ is defined through the numbers $R^{[n]} = R_{K-n}$, which depend on $n$ and not on $K$. (iii) The squared inverse couplings $(g^A_k)^{-2}$ and $(g^B_{k+1})^{-2}$ are coordinates of histories of the two tubes (Theorem~\ref{thm:coupling-flow}), so both lie in $[\ell_k, u_k]$, which gives the first inequality, and $u_k - \ell_k = (b' - b)(K - k) + 2C_W$. Since $b' \leq \frac32b$ and $C_W \leq \beta_0t/4 \leq t/5$ by \eqref{eq:tube-class}, this is at most $\frac12(t + b(K - k) - C_W) = \frac12\ell_k$, and both squared inverse couplings are at least $\ell_k$.
\end{proof}

The lemma is the reason for the tubes. With Ba{\l}aban's own parameters the cube size $R_k$ of the two runs at a shared scale would differ wherever a power of $L$ lies between $(\log(g^A_k)^{-2})^r$ and $(\log(g^B_{k+1})^{-2})^r$, and then the two expansions would have different terms, which could not be compared one by one. With the parameters of the tubes the two runs perform the same decompositions on every shared lattice, and they differ only through the fine lattice and through the couplings, the latter by the numbers $\delta_k$, whose size is determined in Section~\ref{sec:cauchy}.

We need one more object, which carries a field of run~A to the fine lattice of run~B. For a field $U$ on the lattice of spacing $\varepsilon$ whose plaquettes are small, let $\iota(U)$ be the minimiser of the Wilson action on the lattice of spacing $\varepsilon/L$ among the fields $W$ with $M(W) = U$, the minimal configuration of \cite[(2.11), (2.12), p.~256]{BalabanCRE} for one averaging step. By \cite[Theorem~1, p.~279]{BalabanVAR} it exists, it is unique up to gauge transformations, it is regular, and it is an analytic and gauge covariant function of $U$; moreover $M\iota(U) = U$ and $\iota(1) = 1$. For a domain $X$ we write $\iota_X$ for the minimiser localised to $X$ in the sense of \cite[(2.14)--(2.16), p.~257]{BalabanCRE}, which depends on $U$ on a neighbourhood $X^\sim$ of $X$ and satisfies $M\iota_X(U) = U$ there and $\iota_X(1) = 1$. Finally, for a field $U^B$ on the fine lattice of run~B we call
\begin{equation}\label{eq:aligned-background}
    \tilde U \coloneqq MU^B
\end{equation}
the \emph{aligned background}: a field on the fine lattice of run~A, which satisfies run~A's constraints whenever $U^B$ satisfies run~B's, since the iterated averages of run~B are those of run~A composed with $M$.

\subsection{Propagators and Their Walk Expansions}\label{sec:walk-expansions}

We recall the objects to be compared. The fluctuation integral of a step is Gaussian to leading order, with a quadratic form built from the covariant Laplacian at the background, the block averages, and the gauge fixing, and its inverse is expressed through three operators of \cite{BalabanPBG}: $G'$, the inverse of the scalar operator $\Delta'_a$; $(Q'G'^2Q'^*)^{-1}$, an operator on the block data; and $G$, the inverse of the one-form operator $\Delta_a$. Each of them is approximated by local inverses: for a cube $\square$ of the scale $j$, Ba{\l}aban takes a sequence of cubes $\Omega_{j+1}(\square) \subset \Omega_j(\square) = \tilde\square^4 \subset \Omega_{j-1}(\square) \subset \dots \subset \Omega_0(\square) \subset \tilde\square^5$ with the same centre and
\begin{equation}\label{eq:local-sequence}
    \dist\bigl(\Omega_n(\square)^c, \Omega_{n+1}(\square)\bigr) = 2R_0M_0L^n\eta\;,
\end{equation}
on which the operators are taken with Dirichlet conditions outside $\Omega_0(\square)$ and with the averaging of the scale $n$ on the shell $\Omega_n(\square) \setminus \Omega_{n+1}(\square)$, and he calls $G'_\square$, $C_\square = (Q'G'^2_\square Q'^*)^{-1}$, and $G_\square$ the resulting operators \cite[pp.~408--409]{BalabanPBG}. Here $\eta$ is the fine spacing and $R_0M_0$ the smallest size for which the results of \cite{BalabanPRT2} hold, with $M = KR_0M_0$ \cite[p.~408]{BalabanPBG}. With a partition of unity $\sum_\square h_\square^2 = 1$ he puts $G'_0 = \sum_\square h_\square G'_\square h_\square$, and similarly $C_0$ and $G_0$ \cite[(3.87), p.~409]{BalabanPBG}, writes $\Delta'_aG'_0 = I - R'$ with the small remainder $R' = \sum_\square K(h_\square)G'_\square h_\square$ \cite[(3.88)--(3.89), p.~409]{BalabanPBG}, and expands $G' = G'_0\sum_nR'^n$ into a sum over walks of cubes \cite[Theorem~3.7, (3.90), p.~409]{BalabanPBG}. The same construction, with more kinds of factors, gives the expansion of $G$,
\begin{equation}\label{eq:walk-expansion}
    G = \sum_\omega R_0(X_0)R_{\alpha_1}(X_1)\cdots R_{\alpha_n}(X_n)\;,
\end{equation}
over walks $\omega = ((0, X_0), (\alpha_1, X_1), \dots, (\alpha_n, X_n))$ with $X_{i-1} \cap X_i \neq \emptyset$, whose term depends on the configuration on $\tilde X_0^5 \cup \dots \cup \tilde X_n^5$ and satisfies
\begin{equation}\label{eq:walk-bound}
    \bigl\lvert(R_0(X_0)\cdots R_{\alpha_n}(X_n)J)(x)\bigr\rvert \leq O(1)(L^j\eta)^2O(M^{-1/2})^{\abs{\omega}}M^{-\frac12\abs{\omega}}e^{-\frac12\delta_0d(\omega, y, y')}\abs{J}
\end{equation}
for $x \in \Delta(y)$, $y \in \Lambda_j$, and $\operatorname{supp}J \subset \Delta(y')$, together with the corresponding bounds in the norms of \cite[(3.42)--(3.47)]{BalabanPBG} \cite[Theorem~3.10, (3.107), (3.108), p.~416]{BalabanPBG}. Here $d$ is the weighted distance on the set $\mathfrak B$ of the blocks of all scales \cite[(2.46), p.~231]{BalabanPRT2}, which \cite[p.~397]{BalabanPBG} cites as (2.36) (Section~\ref{sec:notes-on-papers}), and $d(\omega, y, y')$ is the infimum of $d(y, y_0) + d(y_0, y_1) + \dots + d(y_{n-1}, y_n) + d(y_n, y')$ over the points $y_i \in X_i \cap \mathfrak B$ \cite[(3.93), p.~410, and p.~413]{BalabanPBG}. The operators $H$, $H_0$, and $\tilde G$ through which the fixed point of the fluctuation integral is written \cite[(1.2)--(1.4), p.~2]{BalabanRG2} have expansions of the same form, ``a propagator is represented by the sum (3.107) [13]'', and the term of a walk is bounded ``as (3.108) [13]'' in ``all operator norms in formulations of theorems in [13]'' \cite[(1.6), (1.7), p.~3]{BalabanRG2}. Ba{\l}aban states these bounds for whole terms of walks and not for their factors, of which he says only that each ``satisfies a bound of the type (3.89), possibly with an additional power of $L^j\eta$'' \cite[p.~413]{BalabanPBG}; the comparison below is accordingly a comparison of whole terms.

What the two runs share is read off from these definitions. The cubes of the scale $j$ are cells of the lattice $ML^j\eta\mathbb Z^4$, and $ML^j\varepsilon = ML^{j+1}(\varepsilon/L)$, so that for $\square$ at the shared scale $j$ the cubes $\square$, $\tilde\square^n$, the partitions of unity as functions on the torus, the localisation domains, and the walks are the same in both runs. The two local sequences are related by
\begin{equation}\label{eq:shared-shells}
    \Omega^B_{n+1}(\square) = \Omega^A_n(\square)\;,
    \qquad 0 \leq n \leq j + 1\;,
\end{equation}
because, by \eqref{eq:local-sequence}, $\dist(\Omega_n(\square)^c, \tilde\square^4) = 2R_0M_0L^n\eta(L^{j-n} - 1)/(L - 1)$ \cite[p.~408]{BalabanPBG}, and the value for run~B at $n + 1$, with $\eta = \varepsilon/L$ and $j + 1$ in place of $j$, is the value for run~A at $n$. So the shells of the two runs coincide at every common spacing $L^n\varepsilon$, and run~B has one more shell $\Omega^B_0 \setminus \Omega^B_1$, of width $2R_0M_0\varepsilon/L$ at the spacing $\varepsilon/L$, where run~A's operators have their Dirichlet exterior. What the runs do not share is the space on which the operators act: on $\Omega_0(\square)$ every factor of run~A acts on fields and sources on the lattice of spacing $\varepsilon$, every factor of run~B on the lattice of spacing $\varepsilon/L$.

\subsection{The Galerkin Identities}\label{sec:galerkin}

Fix $\square$ at the shared scale $j$ and one of the local operators: $\mathcal O^X$ is $\Omega_0\Delta_a\Omega_0$ or $\Omega_0\Delta'_a\Omega_0$ of run $X \in \{ A, B \}$, acting on the space $\mathfrak F^X$ of fields (or functions) on the bonds (or sites) of the fine lattice of run $X$ in $\Omega_0^X(\square)$, at run~B's background $U^B$ for $X = B$ and at the aligned background $\tilde U$ for $X = A$, and $G^X \coloneqq (\mathcal O^X)^{-1}$ on $\mathfrak F^X$. The pairings $\langle\cdot, \cdot\rangle_X$ carry the weights $\eta_X^4$ of \cite[p.~393]{BalabanPBG} and are bilinear, so that $\mathcal O^X$ is symmetric on the complex domain of \cite[(3.37)]{BalabanPBG}, as the analytic extension of a symmetric form. Let $Q_1$ be the linearised one-step average at $U^B$, which maps fields of run~B to fields of run~A. Two averages have to be kept apart here. The operators of \cite{BalabanPBG} are built with the average of the fluctuation relative to the background \cite[(3.13), p.~392; (3.113), p.~418]{BalabanPBG}, and $Q_1$ is always its linear part; the average $M$ of the construction has the linear part $Q^{\mathrm{pl}}_1$, the plain one of \eqref{eq:lifted-term}. The two differ by a gauge transformation of the sites of run~A, $Q_1 = Q^{\mathrm{pl}}_1 + \bar D\Lambda_1$ with $\Lambda_1A$ the block means of the contour variables of $A$, and this term does not vanish at a flat background; but every gauge-invariant function of the averaged configuration is the same for the two averages, and the minimisers of the two constraints are gauge transforms of each other (Supplement, Proposition~S.12(j)). Let $\chi_A$ be the characteristic function of $\Omega^A_0(\square)$, and let $\mathcal L$ be a right inverse of $Q_1$ with
\begin{equation}\label{eq:lift-support}
    \mathcal L\mathfrak F^A \subset \mathfrak F^B\;.
\end{equation}
The \emph{bilinear defect} of $\square$ is
\begin{equation}\label{eq:bilinear-defect}
    \beta_\square(\mathcal F, \mathcal A) \coloneqq \langle\chi_AQ_1\mathcal F, \mathcal O^A\chi_AQ_1\mathcal A\rangle_A - \langle\mathcal F, \mathcal O^B\mathcal A\rangle_B\;,
    \qquad \mathcal F, \mathcal A \in \mathfrak F^B\;,
\end{equation}
and $\beta_\square(\cdot, \mathcal A)^\sharp$ denotes the source of run~B which represents the linear form $\mathcal F \mapsto \beta_\square(\mathcal F, \mathcal A)$ on $\mathfrak F^B$. The defect measures how far the quadratic form of run~A, pulled back by the average, is from that of run~B.

\begin{lemma}[Galerkin Identities]\label{lem:galerkin}
Under \eqref{eq:lift-support}, for every source $g$ of run~B and every source $f$ of run~A,
\begin{align}
    \chi_AQ_1G^Bg - G^A\chi_A\mathcal L^*g &= G^A\chi_A\mathcal L^*\beta_\square(\cdot, G^Bg)^\sharp\;, \label{eq:galerkin-coarse}\\
    \mathcal LG^Af - G^BQ_1^*\chi_Af &= -G^B\beta_\square(\cdot, \mathcal LG^Af)^\sharp\;. \label{eq:galerkin-fine}
\end{align}
\end{lemma}

\begin{proof}
For \eqref{eq:galerkin-coarse} put $u_B \coloneqq G^Bg$ and $u_A \coloneqq G^A\chi_A\mathcal L^*g$. For $X \in \mathfrak F^A$ we have $\mathcal LX \in \mathfrak F^B$ by \eqref{eq:lift-support} and $\chi_AQ_1\mathcal LX = X$, so that
\begin{equation}\label{eq:galerkin-coarse-proof}
\begin{gathered}
    \langle X, \mathcal O^Au_A\rangle_A = \langle\mathcal LX, g\rangle_B = \langle\mathcal LX, \mathcal O^Bu_B\rangle_B\;,\\
    \langle X, \mathcal O^A\chi_AQ_1u_B\rangle_A = \langle\chi_AQ_1\mathcal LX, \mathcal O^A\chi_AQ_1u_B\rangle_A\;.
\end{gathered}
\end{equation}
The difference of the two right sides is $\beta_\square(\mathcal LX, u_B) = \langle X, \chi_A\mathcal L^*\beta_\square(\cdot, u_B)^\sharp\rangle_A$. Both $\chi_AQ_1u_B - u_A$ and $G^A\chi_A\mathcal L^*\beta_\square(\cdot, u_B)^\sharp$ lie in $\mathfrak F^A$, on which $X$ is arbitrary and $\mathcal O^A$ is invertible, and \eqref{eq:galerkin-coarse} follows. For \eqref{eq:galerkin-fine} put $u_A \coloneqq G^Af$ and $u_B \coloneqq G^BQ_1^*\chi_Af$. For $Y \in \mathfrak F^B$, since $\chi_AQ_1\mathcal Lu_A = u_A$,
\begin{equation}\label{eq:galerkin-fine-proof}
    \langle Y, \mathcal O^Bu_B\rangle_B = \langle\chi_AQ_1Y, f\rangle_A = \langle\chi_AQ_1Y, \mathcal O^A\chi_AQ_1\mathcal Lu_A\rangle_A = \langle Y, \mathcal O^B\mathcal Lu_A\rangle_B + \beta_\square(Y, \mathcal Lu_A)\;,
\end{equation}
so that $\langle Y, \mathcal O^B(\mathcal Lu_A - u_B)\rangle_B = -\beta_\square(Y, \mathcal Lu_A)$ for every $Y \in \mathfrak F^B$, and $\mathcal Lu_A - u_B \in \mathfrak F^B$ by \eqref{eq:lift-support}.
\end{proof}

The identities are exact, and the whole difficulty is in the size of the defect. This is not small as a bilinear form: for a field $\mathcal A$ which alternates at the spacing $\varepsilon/L$ inside one block, $Q_1\mathcal A = 0$ and $\beta_\square(\mathcal A, \mathcal A) = -\langle\mathcal A, \mathcal O^B\mathcal A\rangle_B$, which is of the order $(\varepsilon/L)^{-2}$ times the square of the norm of $\mathcal A$. In the identities, however, the defect is paired with $G^Bg$ or $\mathcal LG^Af$, which are regular at the scale of their shell by \cite[(3.43)]{BalabanPBG}, and between regular fields it is small: on a shell of spacing $L^n\eta$ the lattices of the two runs are averages of each other, and the defect of two regular fields is an oscillation of relative order $L^{-n}$. To make this precise we use, for a block $y$ of the shell $n$, the norm
\begin{equation}\label{eq:shell-norm}
    \norm{\mathcal A}_{1,\beta;y} \coloneqq \sup\bigl((L^n\eta)^{-1}\abs{\mathcal A} + \abs{\nabla\mathcal A}\bigr) + (L^n\eta)^\beta\norm{\nabla\mathcal A}_\beta\;,
\end{equation}
the supremum and the H\"older quotient of \cite[(3.40)]{BalabanPBG} taken over the blocks within weighted distance $2$ of $y$. We index the shells by run~A's index $n$, so that run~A's shell $n$ is run~B's shell $n + 1$ by \eqref{eq:shared-shells}, and run~B's extra shell has the index $-1$.

\begin{proposition}[The Defect of the Local Operators]\label{prop:local-defect}
Let $\beta \in [\frac34, 1)$, and let $R_0M_0$ satisfy the conditions of Appendix~\ref{app:conditions}. There is a constant $C_{\mathrm{df}}(\beta)$, independent of the shell, of $j$, $k$, $K$, the torus, the block, and the field, such that for each of the operators of $G'_\square$, $C_\square$, and $G_\square$, at aligned backgrounds on the complex domain of \cite[(3.37)]{BalabanPBG}, and for each $\mathcal A \in \mathfrak F^B$, there are fields $s_0$, $s_1$, $s_2$ of run~B with
\begin{equation}\label{eq:defect-form}
    \beta_\square(\mathcal F, \mathcal A) = \langle D^*\mathcal F, s_2\rangle_B + \langle\nabla\mathcal F, s_1\rangle_B + \langle\mathcal F, s_0\rangle_B
    \qquad\text{for every } \mathcal F \in \mathfrak F^B\;,
\end{equation}
and, on every block $y$ of the shell $n$,
\begin{equation}\label{eq:defect-sources}
    \sup_{\Delta(y)}\bigl(\abs{s_1} + \abs{s_2} + L^n\eta\abs{s_0}\bigr) \leq C_{\mathrm{df}}(\beta)\varpi_n\norm{\mathcal A}_{1,\beta;y}\;,
    \qquad
    \varpi_n \coloneqq \begin{cases} L^{-\beta n}, & n \geq 1,\\ 1, & n \in \{ 0, -1 \}, \end{cases}
\end{equation}
together with a H\"older bound of $s_1$ and $s_2$ at the scale of their shell; the same holds with the two arguments of $\beta_\square$ exchanged.
\end{proposition}

\begin{proof}
We show the mechanism on the averaging terms, which can be computed exactly, and refer to Supplement, Proposition~S.12, for the other terms. The quadratic form of \cite[(3.16)]{BalabanPBG} contains $Q^*aQ$ with the weight $(L^n\eta)^{d-2}$ on the bonds of the shell $n$. Run~B's shell $n + 1$ is run~A's shell $n$, with the weight $(L^{n+1}\varepsilon/L)^{d-2} = (L^n\varepsilon)^{d-2}$, and run~B's averaging of the scale $n + 1$ is run~A's averaging of the scale $n$ composed with the one-step average, $Q^B_{n+1}(U^B) = Q^A_n(\tilde U)Q_1$ (Supplement, Proposition~S.12(a)), while $\chi_A = 1$ on run~A's shells. Hence the averaging part of $\beta_\square$ is exactly $-a\sum_b(\varepsilon/L)^{d-2}\langle\mathcal F(b), \mathcal A(b)\rangle$, the sum over the bonds of run~B's extra shell, which has the form \eqref{eq:defect-form} with $s_0$ supported there and $\varpi_{-1} = 1$. For the scalar operator the coefficients of the averaging on the shells are $a_n$, with $1/a_n = a^{-1}\sum_{i=0}^{n-1}L^{-2i}$ for $n \geq 1$ \cite[p.~394]{BalabanPBG}; run~A's shell $n$ carries $a_n$ and run~B's the same shell $a_{n+1}$, and $\abs{a_{n+1} - a_n}/a_n \leq L^{-2n}/(1 - L^{-2}) \leq 2L^{-\beta n}$, which is again of the form \eqref{eq:defect-sources}. The covariant Laplacian, the curvature term, and the gauge-fixing term are not compositions of averages; on run~A's shells $n \geq 1$ their defect between regular fields is an oscillation of relative order $L^{-n}$, of which the H\"older exponent retains $L^{-\beta n}$, and on the two shells of the finest spacing nothing relates the two runs and $\varpi = 1$. The derivation of these bounds, for $G'_\square$, for $C_\square$, which needs an identity of its own because the data spaces of the two runs differ by run~B's extra layer of data, and for $G_\square$, whose gauge-fixing part is nonlocal through $P = G'Q'^*(Q'G'^2Q'^*)^{-1}Q'G'$, is Supplement, Proposition~S.12(b)--(d).
\end{proof}

The following lemma shows why the shells of the finest spacing, on which the defect is of order one, cost only a factor $(\varepsilon/a_j)^\beta$: a propagator of the scale $j$ reaches them only after crossing $j$ shells, each of width $2R_0M_0$ in its own units.

\begin{lemma}[Reduction over the Shells]\label{lem:shell-reduction}
Let $\mathcal L$ satisfy \eqref{eq:lift-support}, have the range of one block of run~A, and satisfy $\sum_{z \in \Delta(y)}\eta_B^4(\abs{\nabla\mathcal Lv(z)} + (L^n\eta)^{-1}\abs{\mathcal Lv(z)}) \leq C_{\mathcal L}\sum_{y'''}\sum_{w \in \Delta(y''')}\eta_A^4(\abs{\nabla v(w)} + (L^n\eta)^{-1}\abs{v(w)})$, the sum over the $3^4$ blocks $y'''$ next to $y$. Assume $\frac16\delta_0R_0M_0 > 8\log c_0(\frac16) + 1$ and $4\delta_0R_0M_0 \geq 3(1 + \beta)\log L$, with $c_0$ the constant of \cite[(2.58)]{BalabanPRT2}. Then for blocks $y$, $y'$ of $\square$ in the layers $j$ and $j + 1$ of the local sequence, $x \in \Delta(y)$, and a source $g$ of run~B with $\operatorname{supp}g \subset \Delta(y')$,
\begin{equation}\label{eq:shell-reduction}
    \bigl\lvert(\chi_AQ_1G^B_\square - G^A_\square\chi_A\mathcal L^*)g(x)\bigr\rvert \leq C_\square(\beta)\Bigl(\frac{\varepsilon}{a_j}\Bigr)^\beta(L^j\eta)^2e^{-\frac13\delta_0d(y, y')}\abs{g}\;,
\end{equation}
with $C_\square(\beta) \coloneqq 6c_0^4(\frac16)c_\beta C_3$, $C_3 \coloneqq 9\cdot3^4e^{4\delta_0}(1 + L)C_{\mathcal L}C_{\mathrm{df}}(\beta)B_0B_0(\beta)$, and $c_\beta \coloneqq L^{1+\beta}(3 + L^\beta + L^{1+\beta}/(L - 1))$, where $B_0$ and $B_0(\beta)$ are the constants of \cite[(3.42), (3.43)]{BalabanPBG}.
\end{lemma}

\begin{proof}
By \eqref{eq:galerkin-coarse} and Proposition~\ref{prop:local-defect} with $\mathcal A = u_B \coloneqq G^B_\square g$, the left side is $\abs{(G^A\chi_A\mathcal L^*(Ds_2 + \nabla^*s_1 + s_0))(x)}$, and we split the sources over the blocks $y''$ of the shells. We bound the contribution of a block $y''$ of the shell $n$ in two legs.

\emph{The leg from $y'$ to $y''$.} The bounds \cite[(3.42), (3.43)]{BalabanPBG} for $G^B_\square$, which hold for the local sequences by \cite[Corollary~3.6, p.~408]{BalabanPBG}, give $\norm{u_B}_{1,\beta;y''} \leq 3e^{2\delta_0}B_0(\beta)(L^n\eta)e^{-\delta_0d(y'', y')}\abs{g}$, where $e^{2\delta_0}$ accounts for the blocks within distance $2$ and $B_0 \leq B_0(\beta)$. By \eqref{eq:defect-sources},
\begin{equation}\label{eq:leg-two}
    \sup_{\Delta(y'')}\bigl(\abs{s_1} + \abs{s_2} + L^n\eta\abs{s_0}\bigr) \leq 3e^{2\delta_0}C_{\mathrm{df}}(\beta)\varpi_nB_0(\beta)(L^n\eta)e^{-\delta_0d(y'', y')}\abs{g}\;.
\end{equation}

\emph{The leg from $y''$ to $y$.} For a source $\sigma$ on $\Delta(y'')$, the symmetry of $G^A$ gives $(G^A\chi_A\mathcal L^*\nabla^*\sigma)(x) = \langle\nabla\mathcal L[G^A(\cdot, x)], \sigma\rangle_B$, which is at most $\sup\abs{\sigma}$ times the block sum of $\nabla\mathcal LG^A(\cdot, x)$. By the hypothesis on $\mathcal L$ and \cite[(3.42)]{BalabanPBG} for $G^A$, this is at most
\begin{equation}\label{eq:leg-one}
    \sup\abs{\sigma}\cdot3^4e^{2\delta_0}C_{\mathcal L}B_0\bigl[(L^j\eta) + (L^j\eta)^2(L^n\eta)^{-1}\bigr]e^{-\delta_0d(y, y'')}\;,
\end{equation}
and the same holds with $D$ in place of $\nabla^*$, while for $s_0$ the factor is $3^4e^{2\delta_0}C_{\mathcal L}B_0(L^j\eta)^2e^{-\delta_0d(y, y'')}$, that is $(L^j\eta)^2(L^n\eta)^{-1}$ times $L^n\eta\abs{s_0}$. Since $n \leq j + 1$, $L^j\eta \leq L(L^j\eta)^2(L^n\eta)^{-1}$.

\emph{The product and the sum.} The factor $(L^n\eta)^{-1}$ of the second leg cancels the factor $L^n\eta$ of the first, so that the shell $n$ contributes at most $C_3\varpi_n(L^j\eta)^2\sum_{y''}e^{-\delta_0[d(y, y'') + d(y'', y')]}\abs{g}$. One third of the decay gives $e^{-\frac13\delta_0d(y, y')}$ by the triangle inequality \cite[(2.54)]{BalabanPRT2}; one third gives the count $\sum_{y''}e^{-\frac13\delta_0d(y, y'')} \leq 6c_0^4(\frac16)$ over one shell, by \cite[(2.58)]{BalabanPRT2} with $\alpha = \frac13$ under the first condition; and the last third is at most $e^{-\frac43\delta_0R_0M_0(j-n-1)_+} \leq L^{-(1+\beta)(j-n-1)_+}$, because the weighted distance from a block of $\square$ to the shell $n$ is at least $2R_0M_0(j - n - 1)$ \cite[(2.60)]{BalabanPRT2}, both legs cross the shells in between, and the second condition holds. Summing over the shells with \eqref{eq:defect-sources}: the shells $j$ and $j + 1$ give at most $2L^{-\beta j}$; the shells $1 \leq n \leq j - 1$ give $\sum_nL^{-\beta n}L^{-(1+\beta)(j-n-1)} \leq L^{1+\beta}L^{-\beta j}/(L - 1)$; the shell $0$ gives $L^{-(1+\beta)(j-1)} \leq L^\beta L^{-\beta j}$; and run~B's extra shell, at distance at least $2R_0M_0j$, gives at most $L^{-\beta j}$. Hence $\sum_n\varpi_nL^{-(1+\beta)(j-n-1)_+} \leq c_\beta L^{-\beta j}$, with room for the intersected sequences of the cubes next to a coarser layer \cite[p.~408]{BalabanPBG}, and $L^{-j} = \varepsilon/a_j$ by \eqref{eq:shared-spacing}.
\end{proof}

\subsection{The Walk Terms}\label{sec:walk-terms}

We can now state the comparison of the walk terms. Let $W^X_\omega$ be the term of the walk $\omega$ in run $X$'s expansion of one of the operators $H$, $H_0$, or $\tilde G$ of \cite[(1.2)--(1.4), p.~2]{BalabanRG2}, at run $X$'s background, on a set of the $\mathbf R$-operation whose layers lie at shared scales; for these operators the walks of run~A and of run~B are the same, because the cubes are. Let $Y$ be a localisation domain of the walk, and let $\mathcal L_\iota$ be the operator $H_1$ of \cite[(3.129), p.~421]{BalabanPBG} of the one-step problem on the neighbourhood $Y^\sim$ at the background $U^B$, which assigns to a field of run~A the minimiser of the linearised action among the fields of run~B with this average. It is a right inverse of $Q_1$, and when $U^B$ is the localised lift $\iota_Y(\tilde U)$ of its own average it is the derivative of $\iota_Y$ at $\tilde U$ up to a linear gauge transformation, the part of this derivative in the Landau gauge (Supplement, Proposition~S.12(e), (j)). The operator $\mathcal J'$ is the identity on data, for $H$ and $H_0$, and $\mathcal L_\iota^*$ on sources, for $\tilde G$; $m(\omega)$ is the finest scale of a layer of the set which meets $\bigcup_i\tilde X_i^5$.

\begin{theorem}[Walk-Factor Difference]\label{thm:walk-difference}
Let the constants satisfy the conditions of Appendix~\ref{app:conditions}, let $\bar g$ be sufficiently small, and let $K' \geq 2m_{\mathrm{rc}} + 2$, where $m_{\mathrm{rc}}$ is the least integer with $L^{m_{\mathrm{rc}}} \geq (8\log L + \delta_0)/\delta'_0$ and $\delta'_0$ is the rate of Supplement, Proposition~S.12. For every walk $\omega$ of the expansions of $H$, $H_0$, and $\tilde G$ at the backgrounds of the two runs on a set whose layers lie at the shared scales $j \geq m_{\mathrm{rc}}$, among them all the scales $j \geq K'/2 - 1$,
\begin{equation}\label{eq:walk-difference}
    \bigl\lvert(W^B_\omega - \mathcal L_\iota W^A_\omega\mathcal J')\lambda(x)\bigr\rvert \leq \delta_w(\omega)\cdot O(1)(L^j\eta)^pO(M^{-1/2})^{\abs{\omega}}M^{-\frac12\abs{\omega}}e^{-\frac16\delta_0d(\omega, y, y')}\abs{\lambda}\;,
\end{equation}
in the supremum norm, where $(L^j\eta)^p$ is the power of \eqref{eq:walk-bound} for the operator in question, and
\begin{equation}\label{eq:walk-defect}
    \delta_w(\omega) \leq C_w\Bigl(\frac{\varepsilon}{a_{m(\omega)}}\Bigr)^{\gamma_w}\;,
    \qquad
    \gamma_w \coloneqq \min\Bigl(\beta,\, \frac{1 - \varepsilon_0}{3}\Bigr)\;,
\end{equation}
with a constant $C_w$ of the constants of the construction and of $\varepsilon_0$, $\varepsilon_1$, and $\beta$, which does not depend on $K$, $k$, $\varepsilon$, or the torus. If run~A's term is evaluated at the aligned background $\tilde U$ in place of its own, the same holds with the exponent $\beta$ in place of $\gamma_w$. The averaged difference $Q_1W^B_\omega - W^A_\omega\mathcal J'$ satisfies the same bound, with the same map $\mathcal J'$, the identity on the data for $H$ and $H_0$ and $\mathcal L_\iota^*$ on the sources for $\tilde G$. The first differences of the two differences satisfy these bounds with the power of the scale lowered by one and with $\delta_w(\omega)^{\beta/(1+\beta)}$ in place of $\delta_w(\omega)$.
\end{theorem}

\begin{proof}
The proof has three parts: the factors at the aligned background, the change from the aligned background to run~A's own, and the nonlinear functions through which $H$, $H_0$, and $\tilde G$ are built from the propagators.

\emph{The factors.} A term of a walk is a product of factors, each of which is a local operator $G'_\square$, $C_\square$, $G_\square$, multiplied by partition functions, commutators with them, and averaging operators \cite[pp.~409--416]{BalabanPBG}. For two products $a_1\cdots a_n$ and $b_1\cdots b_n$,
\begin{equation}\label{eq:telescoping}
    a_1\cdots a_n - b_1\cdots b_n = \sum_{i=1}^{n}a_1\cdots a_{i-1}(a_i - b_i)b_{i+1}\cdots b_n\;,
\end{equation}
and we apply this first to the averaged difference $Q_1W^B_\omega - W^A_\omega\mathcal J'$. Between two consecutive factors of a walk term stands a fine field, a fine source, or data, and the two runs are compared with $Q_1$ on the fine fields, with the adjoint $\mathcal L^*$ of the cut-off lift $\mathcal L \coloneqq \chi_B\mathcal L_\iota$ on the fine sources, and with the identity on the data. The difference is then the sum, over the factors $c$ of the term, of the product of run~A's factors to the left of $c$, the \emph{intertwining defect} of $c$, and run~B's factors to the right. For a local propagator the intertwining defect is $Q_1G^B - G^A\mathcal L^*$, the left side of \eqref{eq:galerkin-coarse}. For a factor $\mathcal K$ which maps fine fields to fine sources, a commutator with a partition function, it is $\mathcal L^*\mathcal K^B - \mathcal K^AQ_1$, and $\langle X, (\mathcal K^AQ_1 - \mathcal L^*\mathcal K^B)\mathcal A\rangle_A$ is the bilinear defect of $\mathcal K$ at the pair $(\mathcal LX, \mathcal A)$, because $\chi_AQ_1\mathcal LX = X$. For the averages it vanishes, by Supplement, Proposition~S.12(a). The only property of the lift which this uses is that it is an exact right inverse on $\mathfrak F^A$ (Supplement, Proposition~S.12(e)); the product $\mathcal LQ_1$ is not the identity on the fields of run~B and does not occur. Each defect of a local propagator is bounded by Lemma~\ref{lem:shell-reduction}, which costs the factor $C_\square(\beta)(\varepsilon/a_j)^\beta$ and lowers the rate $\frac12\delta_0$ to $\frac16\delta_0$; the explicit factors, multiplications by the partition functions and commutators with them, are functions on the torus which both runs take to be the same, and their defect between the two lattices is of relative order $\varepsilon/a_j$ or $(\varepsilon/a_j)^\beta$, by the oscillation of these functions over the reach of the kernel of $Q_1$ (Supplement, Proposition~S.12(f)). The sum \eqref{eq:telescoping} has $\abs{\omega} + 1$ terms, which the factor $O(M^{-1/2})^{\abs{\omega}}$ absorbs after a change of the constant. Since a cube of the walk lies at a scale $j \geq m(\omega)$, $(\varepsilon/a_j)^\beta \leq (\varepsilon/a_{m(\omega)})^\beta$, and the shells at the scales $n < j$ inside $\tilde\square^5$, which are not layers of the set, enter only through the constant $c_\beta$ of Lemma~\ref{lem:shell-reduction}. This gives the bound of the averaged difference at the aligned background, with the exponent $\beta$.

The fine form \eqref{eq:walk-difference} follows from the averaged one, and it is here that the product $\mathcal L_\iota Q_1$ enters, applied to a field of run~B which is regular. We write
\begin{equation}\label{eq:fine-from-averaged}
    (W^B_\omega - \mathcal L_\iota W^A_\omega\mathcal J')\lambda = (1 - \mathcal L_\iota Q_1)W^B_\omega\lambda + \mathcal L_\iota(Q_1W^B_\omega - W^A_\omega\mathcal J')\lambda\;.
\end{equation}
The second term is the lift of the averaged difference, and the kernel of $\mathcal L_\iota$ is bounded and decays exponentially at the scale $\varepsilon$ \cite[(3.133), p.~422]{BalabanPBG}. The first term is the difference between a field of run~B and the lift of its average. For an arbitrary field this difference is as large as the field, but for the lift $\mathcal L_\iota = G_1Q^*(QG_1Q^*)^{-1}$ it has an exact expression,
\begin{equation}\label{eq:reconstruction-identity}
    (1 - \mathcal L_\iota Q_1)u = \tilde C_1\bigl(D^*D + DRD^* + \Delta' - \Delta'_\pi - \Delta^{(2)}_\pi\bigr)u\;,
    \qquad
    \tilde C_1 \coloneqq G_1 - \mathcal L_\iota QG_1\;,
\end{equation}
in which $\tilde C_1$ is the propagator of the fluctuations of one step, of the size $\varepsilon^2$ and of the range $\varepsilon$, and the operator in the bracket is of the second order with small terms of lower order. Since $u = W^B_\omega\lambda$ satisfies the bound of its walk also in the norm of the first differences, with one power of $a_j$ less \cite[(3.108)]{BalabanPBG}, the first term of \eqref{eq:fine-from-averaged} is of the relative order $\varepsilon/a_j$ (Supplement, Proposition~S.12(i)). The weights of both estimates decay at the scale $\varepsilon$, while the bounds of the walk vary at the scales $a_j \geq \varepsilon L^{m_{\mathrm{rc}}}$ of the layers, which ensures that the former are summed against the latter without a loss in the rate (Supplement, Proposition~S.12(h)); the condition $K' \geq 2m_{\mathrm{rc}} + 2$ places the scales $j \geq K'/2 - 1$ of the large-field terms in this range. This gives \eqref{eq:walk-difference} at the aligned background, with the exponent $\beta$.

The first differences are bounded by interpolation. The averaged difference is the difference of two fields of run~A, each of which satisfies the bound of its walk in the H\"older norm of the first differences, and a field $f$ at the scale $a$ satisfies $\sup\abs{\nabla f} \leq 4(\sup\abs f)^{\beta/(1+\beta)}([\nabla f]_\beta)^{1/(1+\beta)} + 4a^{-1}\sup\abs f$; this turns the factor $\delta_w(\omega)$ of the supremum into $\delta_w(\omega)^{\beta/(1+\beta)}$ for the first differences. For the fine difference the first differences of the two terms of \eqref{eq:fine-from-averaged} are bounded separately: those of the first carry the factor $(\varepsilon/a_j)^\beta$, because the two leading terms of \eqref{eq:reconstruction-identity} vanish on covariantly constant fields up to the curvature, and those of the second are bounded by the first differences of the averaged difference, because $\mathcal L_\iota$ maps regular fields of run~A to regular fields of run~B (Supplement, Proposition~S.12(e), (h), (i)).

\emph{The background.} Run~A's own background is $U^A = e^{i\varepsilon A}\tilde U$ with $a_{j(x)}\abs{A(x)} + a_{j(x)}^2\abs{\nabla A(x)} \leq \bar d(x) \coloneqq C_U(\varepsilon/a_{j(x)})^{(1-\varepsilon_0)/3}\varepsilon_1$, where $j(x)$ is the layer of $x$ and $C_U$ a constant of $M$ (Proposition~\ref{prop:minimiser-distance}). Put $\bar d(\omega) \coloneqq \sup\bar d(x)$ over $\bigcup_i\tilde X_i^5$ and $U(t) \coloneqq e^{it\varepsilon A}\tilde U$. For $\abs{t} \leq r \coloneqq \alpha'_*/(2\bar d(\omega))$, with $\alpha'_* \coloneqq \min(\alpha_1, \alpha_J/\hat x_J)$ the radius of the class of pairs of a configuration and a current of Supplement, Proposition~S.12(g), the configuration $U(t)$ lies in the domain of \cite[(3.37)]{BalabanPBG} layer by layer, where the local propagators are analytic with the bounds of \cite[Theorems~3.1--3.3]{BalabanPBG} \cite[Theorem~3.4, p.~400]{BalabanPBG}, and so is the term of a walk, since the proof of \eqref{eq:walk-bound} uses only these bounds \cite[p.~410]{BalabanPBG}. Under $4\bar d(\omega) \leq \alpha'_*$ we have $r \geq 2$, and Cauchy's estimate gives
\begin{equation}\label{eq:background-change}
    \bigl\lvert W^A_\omega(U^A) - W^A_\omega(\tilde U)\bigr\rvert \leq \sup_{\abs{t} \leq 1}\abs{\partial_tW^A_\omega(U(t))} \leq \frac{\sup_{\abs{t} = r}\abs{W^A_\omega(U(t))}}{r - 1} \leq \frac{4\bar d(\omega)}{\alpha'_*}
\end{equation}
times the bound \eqref{eq:walk-bound}, in each of its norms. With $\bar d(\omega) \leq C_U\varepsilon_1(\varepsilon/a_{m(\omega)})^{(1-\varepsilon_0)/3}$ this is of the form \eqref{eq:walk-defect}.

\emph{The nonlinear functions.} The operators $H$, $H_0$, and $\tilde G$ are built from the propagators through the fixed point $\mathbf A_0$ of \cite[(1.4), p.~2]{BalabanRG2}, the operator $D$ of \cite[(1.3)]{BalabanRG2}, and the solution of a Neumann series; their walk expansions are composite walks with vertices at which three chains meet \cite[p.~427]{BalabanPBG}. The comparison of these functions between the two runs, at complex values of the decoupling parameters $s(\Delta)$ on the circles $\abs{s(\Delta)} = e^{\kappa_1}$ on which they enter Cauchy integrals \cite[(1.23), p.~7]{BalabanRG2}, is Supplement, Proposition~S.12(h); it costs the further constants collected in $C_w$ and no further power of $\varepsilon/a_j$.
\end{proof}

The proof shows where the exponent comes from. The local operators are compared at the rate $(\varepsilon/a_j)^\beta$ for every $\beta < 1$, and the loss is in the change of the background, whose rate $(1 - \varepsilon_0)/3$ is that of the distance between the minimisers of the two runs in the norm of the analyticity domains. This distance is the subject of the next section.

\section{The Large-Field Terms of Two Cutoffs}\label{sec:large-field-comparison}

We now turn to the terms of the two densities which carry large fields. Here the obvious approach, a comparison of the two runs term by term, meets two obstacles. The first is that most of these terms cannot be compared at all with an error that is small relative to the term: a large-field region which was created at a fine scale has been integrated with the test functions, cutoffs, and Gaussian weights of the fine lattice of its run, and the two fine lattices are different. The second is that the test functions are evaluated on fields of the two runs which differ, so that a term of one run may be supported where the corresponding term of the other vanishes. The resolution of the first obstacle is that we do not need to compare the terms with an old large-field region at all: such a term carries the gain of every scale through which the region has lived, and the weight of Proposition~\ref{prop:large-field-weight} shows that the terms of this kind have a total weight which tends to zero with $K$, in each run separately. The terms that remain have their large fields at scales coarser than $K'/2$, where the fine lattice is invisible up to a power of $\varepsilon/a_{K'/2} = L^{-K'/2}$. The second obstacle is removed by the soft test functions, whose change is proportional to that of their arguments.

\subsection{The Matched Decomposition}\label{sec:matched-decomposition}

Both densities at the stopping step have the form \eqref{eq:density-expansion}, a sum over the admissible sequences of domains, with one pair $(\Omega_j, \Lambda_j)$ for each step, and with the components which the $\mathbf R$-operation has renormalised removed from the sequences \cite[(0.3), p.~176]{BalabanLFR1}. The sequences of run~B have one entry more, that of the extra step.

\begin{lemma}[Matched Decomposition]\label{lem:matched-decomposition}
Call a labelled term of either run, at its stopping step, \emph{early} if its pair of domains differs from $(\mathbb T, \mathbb T)$ at a step whose shared scale is at most $1$: for run~B the extra step or the step $2$, for run~A the first step. Then every labelled term is of exactly one of two kinds:
\begin{enumerate}[label=(\Alph*)]
\item[(M)] it is not early; deleting the pair of the extra step from a term of run~B of this kind gives a term of run~A of this kind, and this is a bijection $\tau^+ \mapsto \tau$ between the terms of kind (M) of the two runs;
\item[(U)] it is early; then it has a component of its large-field region at its stopping step whose history contains a large-field region created at a shared scale at most $1$.
\end{enumerate}
\end{lemma}

\begin{proof}
A component which is renormalised at some step leaves the sequences, so a term whose pair at some step is not $(\mathbb T, \mathbb T)$ has a large-field region created at that step which belongs to a component alive at the stopping step; for an early term this is the kind (U). The shared scale $1$ has to be excluded together with the extra step because the two runs do not perform the same operation there: it is the first step of run~A, which is carried out at the retreated region, with labels refined by that region (Supplement, Proposition~S.7), whereas for run~B it is the second step, an ordinary one. A term of run~A whose pair at the first step is $(\mathbb T, \mathbb T)$ has the empty retreat, because the retreat removes cubes near the boundary of the region of integration and the torus has none, so that its label at this step is the pair alone; and a term of run~B whose pairs at the extra step and at the step $2$ are $(\mathbb T, \mathbb T)$ has no further label at these steps either. For the later steps, the admissibility conditions on the sequences \cite[pp.~255--256]{BalabanCRE} and the conditions of the $\mathbf R$-operation \cite[p.~177, (i), (ii)]{BalabanLFR1} are conditions on the domains, stated with the cube sizes, the distances $MR_j$, and the depth $N_j$ of the scale; by Lemma~\ref{lem:shared-parameters} these are the same at a shared scale, so that the conditions of run~B at the steps $3, \dots, K' + 1$ are those of run~A at the steps $2, \dots, K'$. This gives the bijection.
\end{proof}

For a matched pair $(\tau^+, \tau)$ the two terms are the same operations applied to ingredients which depend on the run. The pair is the \emph{all-small-field pair} if $\tau = \ell_0$. Otherwise let $j_0$ be the least scale at which one of its terms has a nonempty large-field region; since every nonempty large-field region of a term belongs, after the $\mathbf R$-operations, to the history of a component which is alive at the stopping step, $j_0$ is the scale at which the oldest of these components was born, and $\Lambda_i = \Omega_i = \mathbb T$ for $i < j_0$ in both terms. With
\begin{equation}\label{eq:young-old}
    m \coloneqq \floor{K'/2}\;,
\end{equation}
the pair is \emph{old} if $K' - j_0 \geq m$ and \emph{young} otherwise. A young pair has all its large fields at the shared scales $i \geq j_0 - 1 \geq K' - m \geq K'/2$; this is the case in which a comparison is possible.

The terms which are not compared are disposed of by the following lemma, which uses only the weight of Section~\ref{sec:large-field-weight}. We write
\begin{equation}\label{eq:loop-majorant}
    \mathfrak b(\mathcal C) \coloneqq e^{B_E(\mathcal C) + I_*(\mathcal C)}\;,
\end{equation}
so that $\mathfrak b(\mathcal C) \leq B(\mathcal C)$, and we use the notation of Theorem~\ref{thm:stopping-identity} for either run, with $Z$ the integral of the density and $N \coloneqq \sum_\ell Z_\ell\langle\widehat{\mathcal W}_\ell\rangle_\ell$, so that $W = N/Z$.

\begin{lemma}[Terms with an Old Component]\label{lem:old-component}
Let $\bar g$ be as in Theorem~\ref{thm:uniform-bound}, let $p \geq 0$, and let $\mathcal S$ be a set of labelled terms of one run, none of them $\ell_0$, each of which has a component $X$ of its large-field region at the stopping step with $P(X) \geq p$. Then
\begin{equation}\label{eq:old-component}
    \sum_{\ell \in \mathcal S}Z_\ell \leq e^{-C_Pp}\delta_{\mathrm{LF}}Z_{\ell_0}\;,
    \qquad
    \sum_{\ell \in \mathcal S}Z_\ell\abs{\langle\widehat{\mathcal W}_\ell\rangle_\ell} \leq \mathfrak b(\mathcal C)e^{-(C_P-1)p}\delta_{\mathrm{LF}}Z_{\ell_0}\;,
\end{equation}
and removing the terms of $\mathcal S$ from $N$ and $Z$ changes $W = N/Z$ by at most $2B(\mathcal C)e^{-(C_P-1)p}\bar\delta_{\mathrm{LF}}$.
\end{lemma}

\begin{proof}
For each $\ell \in \mathcal S$ choose a component $X_\ell$ with $P(X_\ell) \geq p$. For the first bound let $\omega(X_\ell) \coloneqq e^{C_PP(X_\ell)}$ and $\omega(X) \coloneqq 1$ for the other components; these numbers satisfy the hypothesis of Proposition~\ref{prop:large-field-weight}, so that
\begin{equation}\label{eq:old-component-proof}
    \sum_{\ell \in \mathcal S}Z_\ell \leq e^{-C_Pp}\sum_{\ell \in \mathcal S}Z_\ell\,\omega(X_\ell) \leq e^{-C_Pp}\sum_{\ell \neq \ell_0}Z_\ell\prod_{X \in Z(\ell)}\omega(X) \leq e^{-C_Pp}\delta_{\mathrm{LF}}Z_{\ell_0}\;.
\end{equation}
For the second we use Lemma~\ref{lem:term-bound}, $\abs{\langle\widehat{\mathcal W}_\ell\rangle_\ell} \leq \mathfrak b(\mathcal C)\prod_X\omega_A(X)$, and the numbers $\omega(X_\ell) \coloneqq \omega_A(X_\ell)e^{(C_P-1)P(X_\ell)}$ and $\omega(X) \coloneqq \omega_A(X)$ otherwise. Since $1 \leq \omega_A(X) \leq e^{P(X)}$ under \eqref{eq:attached-condition}, as the proof of Lemma~\ref{lem:term-bound} shows, they satisfy $1 \leq \omega(X) \leq e^{C_PP(X)}$, and the same computation gives the second bound. Finally, if $N_{\mathcal S}$ and $Z_{\mathcal S}$ are the contributions of $\mathcal S$, then $Z - Z_{\mathcal S} \geq Z_{\ell_0}$ and
\begin{equation}\label{eq:removal}
    \Bigl\lvert\frac{N - N_{\mathcal S}}{Z - Z_{\mathcal S}} - \frac{N}{Z}\Bigr\rvert = \Bigl\lvert\frac{N_{\mathcal S} - WZ_{\mathcal S}}{Z - Z_{\mathcal S}}\Bigr\rvert \leq \frac{\abs{N_{\mathcal S}} + \abs{W}Z_{\mathcal S}}{Z_{\ell_0}} \leq \bigl(\mathfrak b(\mathcal C) + B(\mathcal C)\bigr)e^{-(C_P-1)p}\delta_{\mathrm{LF}}\;,
\end{equation}
by \eqref{eq:old-component} and Theorem~\ref{thm:uniform-bound}, and $\delta_{\mathrm{LF}} \leq \bar\delta_{\mathrm{LF}}$.
\end{proof}

\begin{corollary}[The Unmatched Terms and the Old Pairs]\label{cor:unmatched-old}
Under the hypotheses of Lemma~\ref{lem:old-component}, removing the terms of kind (U) from both runs changes $W_{K+1}(\mathcal C) - W_K(\mathcal C)$ by at most $4B(\mathcal C)e^{-(C_P-1)K'}\bar\delta_{\mathrm{LF}}$, and removing the old pairs from both runs changes $W_{K+1}(\mathcal C) - W_K(\mathcal C)$ by at most $4B(\mathcal C)e^{-(C_P-1)(m+1)}\bar\delta_{\mathrm{LF}}$.
\end{corollary}

\begin{proof}
The large-field regions grow and merge until they are renormalised \cite[(0.3), p.~176]{BalabanLFR1}, so a component alive at the stopping step whose history begins at the scale $j$ meets the large-field region at every scale from $j$ to the stopping step, and its history volume \eqref{eq:history-quantities}, in which every scale with a nonempty intersection contributes at least $L^4$, is at least the number of these scales. For a term of kind (U) this number is at least $K'$, the number of the steps from the shared scale $1$ to the stopping step, and for each term of an old pair it is at least $m + 1$. Lemma~\ref{lem:old-component} gives the assertions.
\end{proof}

\subsection{The Shells of the Test Functions}\label{sec:shells}

On a young pair the two terms are compared ingredient by ingredient, and the first ingredients are the test functions. Let $\Pi_X$ be the product of the test functions and characteristic functions of run $X$'s term and $\varrho_X$ the rest of its integrand. Since all factors of $\Pi_X$ take values in $[0, 1]$,
\begin{equation}\label{eq:three-piece}
    \Pi_A\varrho_A - \Pi_B\varrho_B = \min(\Pi_A, \Pi_B)(\varrho_A - \varrho_B) + (\Pi_A - \Pi_B)_+\varrho_A - (\Pi_B - \Pi_A)_+\varrho_B\;,
\end{equation}
which puts each of the last two pieces on the support of one run's own term, where that run's bounds apply.

\begin{lemma}[Shells of the Soft Test Functions]\label{lem:soft-shells}
Let a soft test function $\bar\chi(f_\square/t)$ of a cube $\square$ have the same threshold $t$ in both runs, and let $\abs{f^A_\square - f^B_\square} \leq \delta_\square t$ wherever $\min(f^A_\square, f^B_\square) < (1 + 2\kappa_s)t$. Then $\abs{\bar\chi(f^A_\square/t) - \bar\chi(f^B_\square/t)} \leq \delta_\square/\kappa_s$, and the same holds for $1 - \bar\chi$. Consequently, if $\Pi_A$ and $\Pi_B$ consist of such soft test functions and of characteristic functions which are the same in both runs,
\begin{equation}\label{eq:soft-shell-bound}
    (\Pi_A - \Pi_B)_+\varrho_A \leq \Bigl(\sum_\square\frac{\delta_\square}{\kappa_s}\Bigr)\,\Pi^{\sharp}_A\varrho_A\;,
\end{equation}
where $\Pi^\sharp_A$ is the product of the sharp test functions of run~A at the thresholds $(1 + \kappa_s)t$ for the passed and $t$ for the failed test functions; for a test function of the cascade of \cite[(1.22)--(1.24), p.~181]{BalabanLFR1} the factor $\delta_\square/\kappa_s$ is replaced by $(N + 1)\delta_\square/\kappa_s$ and the thresholds by those of Supplement, Proposition~S.2.
\end{lemma}

\begin{proof}
Where $\min(f^A_\square, f^B_\square) \geq (1 + 2\kappa_s)t$ both values of $\bar\chi$ vanish. Elsewhere $\abs{f^A_\square/t - f^B_\square/t} \leq \delta_\square$, and $\bar\chi$ is $\kappa_s^{-1}$-Lipschitz by Lemma~\ref{lem:soft-decomposition}(iii). Replacing the factors of $\Pi_A$ by those of $\Pi_B$ one at a time, with all factors in $[0, 1]$, gives $(\Pi_A - \Pi_B)_+ \leq \mathbf 1\{\Pi_A > 0\}\sum_\square\delta_\square/\kappa_s$, and on $\{\Pi_A > 0\}$ every passed test function of run~A has $f^A_\square < (1 + \kappa_s)t$ and every failed one $f^A_\square > t$, by Lemma~\ref{lem:soft-decomposition}(ii). The cascade test functions are soft with the width $\kappa_s/(N + 1)$ (Supplement, Proposition~S.2), which gives the factor $N + 1$.
\end{proof}

The right side of \eqref{eq:soft-shell-bound} is a term of run~A with sharp test functions at enlarged thresholds, and Ba{\l}aban's bounds hold for it by Proposition~\ref{prop:construction-admissible}: a shell costs, relative to Ba{\l}aban's bound of the term, the factor $\delta_\square/\kappa_s$. The characteristic functions which restrict the fluctuation fields are sharp, and their thresholds differ between the runs through the factor $g_j$ of Definition~\ref{def:construction}(ii). For them we use the density of the Gaussian weight.

\begin{lemma}[Shells of the Coupling]\label{lem:coupling-shells}
Let $\abs{\delta_j} \leq \frac12(g^A_j)^{-2}$, and let the fluctuation field $A_j(b)$ of a bond be restricted in the two runs by $\abs{A_j(b)} < t$ and $\abs{A_j(b)} < t'$ with $t = \vartheta_j/g_j$ and $t' = \vartheta_j/g'_j$, $g_j$ and $g'_j$ the two couplings at the shared scale $j$. Then $\abs{t' - t} \leq 2tg_j^2\abs{\delta_j}$, and the mass of the shell between the two radii under the normalised Gaussian weight of the step, with the other variables fixed, is at most
\begin{equation}\label{eq:coupling-shell}
    C_{\mathrm{sh}}\,p_0(\hat g_j)^{\dim G}g_j^2\abs{\delta_j}\;,
\end{equation}
with a constant $C_{\mathrm{sh}}$ of $L$, $G$, and $A_1/A_0$. The same holds for the field $B_j$ of \cite[(1.55), p.~188]{BalabanLFR1} with $\delta'_j$ in place of $\vartheta_j$.
\end{lemma}

\begin{proof}
With $x \coloneqq g_j^{-2}$ and $x' \coloneqq g_j'^{-2}$ we have $t'/t = g_j/g'_j = (x'/x)^{1/2}$ and $\abs{x' - x} = \abs{\delta_j} \leq \frac12x$, so that $x'/x \in [\frac12, \frac32]$, and $\abs{u^{1/2} - 1} \leq \abs{u - 1}$ for $u > 0$ gives $\abs{t'/t - 1} \leq \abs{\delta_j}/x = g_j^2\abs{\delta_j}$; exchanging the runs gives the bound with either coupling and the factor $2$. The Gaussian weight of the fluctuation field has the precision $Q = C^*\Delta^{(j)}C - C^*\Delta^{(j)}CC^{(j)}C^*\Delta^{(j)}C \leq C^*\Delta^{(j)}C$ of \cite[(2.21), p.~258]{BalabanCRE}, because $C^{(j)} \geq 0$. The conditional variance of $A_j(b)$, the other variables fixed, is $(Q_{bb})^{-1} \geq \norm{C^*\Delta^{(j)}C}^{-1}$, so that the conditional density of $A_j(b) \in \mathfrak g$ is at most $(\norm{C^*\Delta^{(j)}C}/2\pi)^{\dim G/2}$, a constant of the lattice. The shell $\{ \min(t, t') \leq \abs{A_j(b)} < \max(t, t') \}$ has Lebesgue measure at most $C\abs{t' - t}\max(t, t')^{\dim G - 1}$, and $t \leq (A_1/A_0)p_0(\hat g_j)\cdot\hat g_j/g_j \leq (A_1/A_0)p_0(\hat g_j)$ by Lemma~\ref{lem:tube-couplings}(i), which gives \eqref{eq:coupling-shell}. The same argument applies to $B_j$ under the Gaussian form of \cite[p.~188]{BalabanLFR1}, and to the complement of a characteristic function, whose shell is the same set.
\end{proof}

The positivity of the quadratic form is the same fact by which Ba{\l}aban bounds the integrals over the fluctuation fields in the large-field operations \cite[p.~383]{BalabanLFR2}, and so a shell term is at most \eqref{eq:coupling-shell} times his bound of the term. Two features of these shell bounds matter for what follows. The plaquette shells cost $\delta_\square/\kappa_s$, which will be a power of $\varepsilon/a_j$; the coupling shells cost a multiple of $g_j^2\abs{\delta_j}$, which is controlled in Section~\ref{sec:cauchy} together with the other differences of the couplings.

\subsection{The Configurations of the Two Runs}\label{sec:configurations}

The functions $f_\square$ of the plaquette and reference test functions, and all integrands of the large-field operations, are computed from minimisers of the action on the fine lattice of the run. At equal data the minimisers of the two runs differ by a lattice artefact, which is the content of the following proposition. For a field $W$ on the fine lattice of run~B we write $MW$ for its average on the fine lattice of run~A, and for a point $x$ of a determining set $j(x)$ is the layer of $x$.

\begin{proposition}[Minimiser Distance]\label{prop:minimiser-distance}
Let $0 < \varepsilon_0 < 1/125$, let the constants satisfy the conditions of Appendix~\ref{app:conditions}, and let $\bar g$ be sufficiently small. There are constants $C_G(\varepsilon_0)$ and $C_U$, independent of $K$, $k$, $\varepsilon$, and the torus, with the following properties.
\begin{enumerate}[label=(\roman*)]
\item Let $U^A$ and $U^B$ be the minimisers of the two runs for a determining set at the shared scales with the same data $V$, which satisfy the hypothesis (7) of \cite[Theorem~1, p.~279]{BalabanVAR} with the parameter $\varepsilon_1$. Then, in a gauge, $U^A = e^{i\varepsilon A'}MU^B$ with
\begin{equation}\label{eq:minimiser-distance}
    a_{j(x)}\abs{A'(x)} + a_{j(x)}^2\abs{\nabla A'(x)} \leq 2\bar\rho(x)\;,
    \qquad
    \bar\rho(x) \coloneqq C_G(\varepsilon_0)\Bigl(\frac{\varepsilon}{a_{j(x)}}\Bigr)^{1-\varepsilon_0}\varepsilon_1\;.
\end{equation}
\item For each of the configurations $U''_k$, $U^0_k$, $U_{k,\Lambda}$, $U_k$ of the $\mathbf R$-operation of \cite[\S1]{BalabanLFR2} on a young history born at the scale $n$, at the data of the two runs, $c^A = e^{i\varepsilon A_c}Mc^B$ with
\begin{equation}\label{eq:coarse-side}
    a_{j(x)}\abs{A_c(x)} + a_{j(x)}^2\abs{\nabla A_c(x)} \leq \bar d(x) \coloneqq C_U\Bigl(\frac{\varepsilon}{a_{\min(j(x), n)}}\Bigr)^{(1-\varepsilon_0)/3}\varepsilon_1\;.
\end{equation}
\end{enumerate}
Here $\varepsilon_1 \leq \bar\varepsilon_1$, the number of \eqref{eq:data-parameter} below, which depends on $\bar g$ and the constants of the construction and not on $K$.
\end{proposition}

The proof is Supplement, Proposition~S.13. In (i), which is the case of a single run of \cite[Theorem~1]{BalabanVAR} on two lattices, $MU^B$ satisfies run~A's constraint, and the difference of the two actions at $MU^B$ is the blocking artefact $b(W) \coloneqq A_\varepsilon(MW) - A_{\varepsilon/L}(W)$, whose first variation is a tangential current of order $\varepsilon/a_j$; the gap of the action at the minimum and the bounds of \cite{BalabanPBG} for the Green operator of the constrained problem turn it into \eqref{eq:minimiser-distance}, with the loss $\varepsilon_0$ in the H\"older norm. In (ii) the data themselves differ, on the region $\Lambda$ where the $\mathbf R$-operation integrates, through the minimiser $U_0$ of \cite[(1.3), p.~357]{BalabanLFR2}; the minimiser of the nested problem of \cite[Proposition~1, p.~194]{BalabanLFR1} is compared by the growth of its action away from the minimum. If $\Delta$ is the distance between the minimising data of the two runs, the quadratic growth $\frac12\kappa_\Lambda\Delta^2$ is at most the change of the difference of the two runs' actions between the two minimisers, a sum over blocks of the smaller of a bound of order $(\varepsilon/a_n)^{1-\varepsilon_0}$ and a bound linear in $\Delta$; with $\min(u, v) \leq (uv)^{1/2}$ this gives $\Delta^{3/2} \leq C(\varepsilon/a_n)^{(1-\varepsilon_0)/2}$, whence the exponent $(1 - \varepsilon_0)/3$.

The data parameter of the regularity of the history is
\begin{equation}\label{eq:data-parameter}
\begin{aligned}
    \varepsilon_1 &\coloneqq \max\Bigl(C\max_{j \leq K'}(\varepsilon_j + \vartheta_j),\, C_MR_{K'}\varepsilon_{K'}\Bigr)\\
    &\leq \bar\varepsilon_1 \coloneqq \max\bigl(C(1 + A_1/A_0)\bar gp_0(\bar g),\, C_MLA_0\bar g(\log\bar g^{-2})^{p_0+r}\bigr)\;,
\end{aligned}
\end{equation}
where $C$ is the constant with which \cite[p.~178]{BalabanLFR1} bounds the data of a large-field region by its thresholds and $C_M$ a constant of $M$ \cite[(1.83), p.~197]{BalabanLFR1}; the bound holds because $\varepsilon_j + \vartheta_j = (1 + A_1/A_0)\hat g_jp_0(\hat g_j)$ with $\hat g_j \leq \bar g$, because $R_{K'}\varepsilon_{K'} < LA_0\hat g_{K'}(\log\hat g_{K'}^{-2})^{p_0+r}$ by Lemma~\ref{lem:tube-cubes}(i), and because $g \mapsto gp_0(g)$ and $g \mapsto g(\log g^{-2})^{p_0+r}$ increase for small $g$. The point of (i) is the exponent: the tested quantities of a cube at the scale $j$ differ between the runs by a power of $\varepsilon/a_j$, which is small precisely because the young histories live at the scales $j \geq K'/2 - 1$. In terms of the thresholds:

\begin{corollary}[The Tested Quantities]\label{cor:tested-quantities}
Under the hypotheses of Proposition~\ref{prop:minimiser-distance}, let $f^X_\square$ be the quantity which a plaquette or reference test function of a cube $\square$ at the scale $j$ of a young history evaluates in run $X$, at the threshold $t$. Wherever $\min(f^A_\square, f^B_\square) < (1 + 2\kappa_s)t$, at equal data,
\begin{equation}\label{eq:tested-quantities}
    \abs{f^A_\square - f^B_\square} \leq \delta^s_jt\;,
    \qquad
    \delta^s_j \coloneqq \delta^f_j + 7g_j^2\abs{\delta_j}\;,
    \qquad
    \delta^f_j \coloneqq C_f\Bigl(\frac{\varepsilon}{a_j}\Bigr)^{1-2\varepsilon_0}\;,
\end{equation}
with a constant $C_f$ which is independent of $K$, $k$, $\varepsilon$, and the torus, and at most polynomial in $j_q$.
\end{corollary}

\begin{proof}
For the test functions of the local minimisers of the cubes \cite[(1.3), (1.4), (1.7), (1.8), p.~178]{BalabanLFR1}, \cite[(2.17), p.~257]{BalabanCRE}, the plaquettes of $U^A$ and of $MU^B$ differ by the curl of $\varepsilon A'$ with transports, at most $a_j^2\abs{\nabla A'} + C(a_j\abs{A'})^2 + C(a_j\abs{A'})B_3\varepsilon_1 \leq 3\bar\rho$ in the units of the scale for $\bar g$ small, by \eqref{eq:minimiser-distance} and \cite[(8)]{BalabanVAR}; the plaquettes of $MU^B$ and of $U^B$ differ by the blocking artefact, of the same order. Relative to the threshold, which is at least $(1 - \beta)\varepsilon_j \geq \delta'_j$, this is at most a constant times $(\varepsilon/a_j)^{1-\varepsilon_0}\varepsilon_1/\delta'_j$. On a young history the scales satisfy $K' - j + 1 \leq 8(1 + j_q)(j + 1)$, and $\varepsilon_1/\delta'_j$ is at most a polynomial in $j$ and $\log\hat g_j^{-2}$ with coefficients polynomial in $j_q$, by \eqref{eq:data-parameter}, Lemma~\ref{lem:tube-couplings}, and $p_1 < p_0$; this polynomial is absorbed by $(a_j/\varepsilon)^{\varepsilon_0} = L^{\varepsilon_0j}$, which gives $\delta^f_j$. The test functions of the multiscale minimisers, those of the cascade \cite[(1.22)--(1.24), p.~181]{BalabanLFR1}, of the reference configurations \cite[(1.27), p.~182]{BalabanLFR1}, and of \cite[(1.75), p.~193]{BalabanLFR1}, are bounded in the same way, with the layer of the tested plaquette in place of $j$ (Supplement, Proposition~S.13(c)). The coupling part $7g_j^2\abs{\delta_j}$ enters only where a tested quantity contains the scaled field $g_jB_j$, as in \cite[(1.27), (1.59)]{BalabanLFR1}, by the computation of Lemma~\ref{lem:coupling-shells}.
\end{proof}

\subsection{The Large-Field Comparison}\label{sec:lf-comparison}

We can now state the comparison of the large-field terms of a young pair. Let $\tilde\theta \in (0, 1)$ be the rate of \eqref{eq:theta-tilde} in Section~\ref{sec:cauchy}; in this section we use only $\tilde\theta \geq L^{-(1-\varepsilon_0)/4}$ and $\tilde\theta \geq L^{-\theta_V/8}$, with $\theta_V$ of \eqref{eq:age-exponent}. The ingredients of a term are the factors of the integrands of its one-step operations and the terms of its effective action. For a young pair we distinguish three kinds of them at the scales $i \geq j_0 - 1$:
\begin{enumerate}[label=($\alpha$)]
\item[($\alpha$)] the ingredients of the small-field integration: the frozen terms of the effective action, the Gaussian weights and their normalisations, the function \cite[(1.6), p.~178]{BalabanLFR1} of the block variables, and the localised constants;
\item[($\beta$)] the test functions of the plaquettes and of the reference configurations, and the characteristic functions of the fluctuation fields;
\item[($\gamma$)] the large-field integrands: the action and the quadratic form of \cite[(1.70), p.~377]{BalabanLFR2}, the minimiser $U_0$ and the operators $H_{1,k}$ and $\Delta_1(\zeta_0)$ of the denominators of the renormalisation quotients \cite[(1.1)--(1.2), p.~356]{BalabanLFR2}, and the factors $V(Y)$ of \cite[(1.91)]{BalabanLFR2} and the boundary terms of the steps at which a large-field region existed, as functions of the configurations and the propagators of the run, the terms of the effective action from which they are built being held fixed.
\end{enumerate}
The dependence of the factors $V(Y)$ and of the boundary terms on the terms of the effective action from which they are built is of the kind ($\alpha$).
The ingredients ($\alpha$) are compared in Section~\ref{sec:cauchy}, by the recursion which also compares the all-small-field pair. The ingredients ($\beta$) and ($\gamma$) are compared by the following theorem, which says that their run dependence costs, per point of the history of a component, an amount which decays geometrically in the scale at which the component was born. For a component $X$ of a young pair let $n(X)$ be its birth scale and $\abs{Z_i \cap X}_i$ the number of points of the scale $i$ of its history at the scale $i$, as in \eqref{eq:history-quantities}. Write $M^A_\tau$ for the majorant \eqref{eq:term-majorant} of run~A's term and $M^B_{\tau^+}$ for that of run~B's.

\begin{theorem}[Large-Field Comparison]\label{thm:lf-comparison}
Let the constants satisfy the conditions of Appendix~\ref{app:conditions}, and let $\bar g$ be sufficiently small. There are a constant $C_M$ of $L$ and $G$, a constant $C_h$ which is independent of $K$, $\varepsilon$, and the torus and at most polynomial in $j_q$, and an index $k_{\mathrm{LF}}$ with the following property. For $K' \geq k_{\mathrm{LF}}$ and every young matched pair $(\tau^+, \tau)$, let $\varrho^{A,B}[\tau]$ be run~A's term with its ingredients of the kinds ($\beta$) and ($\gamma$) replaced by run~B's, except for the factors \eqref{eq:three-piece} of the shells. Then
\begin{equation}\label{eq:lf-comparison}
    \int\bigl\lvert\varrho^{A,B}[\tau] - \varrho^A[\tau]\bigr\rvert\,\dd V \leq \Bigl(\exp\Bigl(\sum_{X}\sum_{i}C_Ms^{\mathrm{LF}}_{n(X)}L^4\abs{Z_i \cap X}_i\Bigr) - 1 + \sigma_\delta\Bigr)M^A_\tau\;,
    \qquad
    s^{\mathrm{LF}}_n \coloneqq C_h\tilde\theta^n\;,
\end{equation}
the sum over the components $X$ of the pair and the scales $i$ of their histories, where $\sigma_\delta \coloneqq \sum_XP(X)\sup_{i \geq K'-m}\abs{\delta_i}$ collects the coupling parts of the shells; the shell pieces of \eqref{eq:three-piece} satisfy the same bound, with $M^B_{\tau^+}$ for the piece on run~B's support. The same bounds hold for the numerators and the denominators of the renormalisation quotients of \cite[(1.1), (1.2), p.~356]{BalabanLFR2} at the complex arguments of the inductive class, relative to Ba{\l}aban's majorants of their moduli, and linearly in the norms of the classes of the effective actions of the two runs.
\end{theorem}

\begin{proof}
We treat the two kinds of ingredients in turn and then sum.

\emph{The ingredients ($\beta$).} By the redundancy of Ba{\l}aban's test functions \cite[(3.6)--(3.9), p.~265--266; p.~268--269]{BalabanCRE}, a test function of the scale $j$ which is not identically one on the support of the other characteristic functions of the term lies in the large-field region of the term or in its collar, so on a young pair it lies on a cube of a history at a scale $j \geq j_0 - 1$, or on a cube of the stopping scale. By Lemma~\ref{lem:soft-shells} and Corollary~\ref{cor:tested-quantities} each of them contributes the factor $\delta^s_j/\kappa_s$, and $(N + 1)\delta^s_j/\kappa_s$ for a cascade test function; by Lemma~\ref{lem:coupling-shells} each characteristic function of a fluctuation field contributes at most $C_{\mathrm{sh}}p_0(\hat g_j)^{\dim G}g_j^2\abs{\delta_j}$. A cube of the scale $j$ carries the test functions of the step $j$ and of the $\mathbf R$-operations of the steps $k' \geq j$ with $k' - N_{k'} \leq j$, at most $14$ of each step and two of the cascade, and there are at most $s_T(j) \coloneqq 2 + N_j$ such steps, because $N$ is nonincreasing (Lemma~\ref{lem:depth}). Since $N_j \leq \log_LC_{\mathrm{dep}} + \nu_L\log_L\log g_j^{-2}$ by \eqref{eq:depth-bounds}, the first parts of all the test functions of a cube are at most a constant times $(j + 1)s_T(j)L^{-(1-2\varepsilon_0)j}$, and with $L^{-(1-2\varepsilon_0)j} \leq \tilde\theta^jL^{-(3-7\varepsilon_0)j/4}$, which holds because $\tilde\theta \geq L^{-(1-\varepsilon_0)/4}$, and $j \geq n(X) - 1$, this is at most $s^{\mathrm{LF}}_{n(X)}$ once $C_h$ contains $\tilde\theta^{-1}\sup_j(j + 1)s_T(j)L^{-(3-7\varepsilon_0)j/4}$, which is finite since $s_T(j)$ grows like $\log\log j$. The coupling parts of all the test functions and characteristic functions of a cube sum to at most $(7/\kappa_s)16s_T(j)^2g_j^2\abs{\delta_j} + 8s_T(j)C_{\mathrm{sh}}p_0(\hat g_j)^{\dim G}g_j^2\abs{\delta_j} \leq \abs{\delta_j}$ for $\bar g$ small, and the cubes of the histories are counted by $P(X)$.

\emph{The ingredients ($\gamma$).} By Proposition~\ref{prop:minimiser-distance}, the configurations of the $\mathbf R$-operation of the two runs differ, in the units of their layers, by $(\varepsilon/a_n)^{(1-\varepsilon_0)/3}$ times a constant, where $a_n \geq a_{n(X)-1}$ is the spacing of the finest layer of the history. Each large-field integrand is a function of these configurations whose two-run difference is bounded per point of the history by a constant times $g_j^{-2}(\varepsilon/a_n)^{(1-\varepsilon_0)/3}\varepsilon_1^2$, for the actions and quadratic forms, by a constant times $(\varepsilon/a_n)^{(1-\varepsilon_0)/3}N_k$, for the forms of the denominators, and, for the factors $V(Y)$ relative to their bound $\abs{V(Y)} \leq \alpha e^{-(1+2\beta)\kappa d(Y)}$ of \cite[(1.68), p.~377]{BalabanLFR2}, by a sum over the generations $j$ of the terms which they re-sum of powers $\theta_V$ of the configuration distance relative to the radii of analyticity, with the factor $L^{-(\beta'/2)(m-j)}$ of the age, where Theorem~\ref{thm:walk-difference} compares the propagators through which they are built; the boundary terms are compared through the recursion of their generation, under the third condition of \eqref{eq:t-one}. These bounds are Supplement, Proposition~S.14. The factors $g_j^{-2}$, $\varepsilon_1^2/\varepsilon_j^2$, $N_k$, and the polynomial dependence on the scale are absorbed into the geometric factor, as in the first part, since $(1 - \varepsilon_0)/3 > (1 - \varepsilon_0)/4$, and for the factors $V(Y)$ since $\tilde\theta \geq L^{-\theta_V/8}$.

\emph{The sum.} Every ingredient of the two kinds enters per point of the history of a component at a scale $i \geq n(X) - 1$, and its contribution relative to Ba{\l}aban's bound is at most $C_Ms^{\mathrm{LF}}_{n(X)}$ per point plus the coupling parts, with $C_M$ the number of ingredients per point, a constant of $L$ and $G$. Ba{\l}aban's bound of a term is proved by estimating the integrals over the group variables, with their normalised weights, by one, and those over the fluctuation fields by the positivity of their quadratic forms \cite[p.~383]{BalabanLFR2}. A product of factors which each change by a relative amount $x_\iota \geq 0$ changes by at most $\prod_\iota(1 + x_\iota) - 1 \leq \exp(\sum_\iota x_\iota) - 1$ relative to the product of the bounds, which gives \eqref{eq:lf-comparison}. For the renormalisation quotients the same argument applies to the moduli of the numerators and denominators at complex arguments, where the shells are those of the real integration variables, and the dependence on the effective actions of the two runs enters linearly through the norms of their classes (Supplement, Proposition~S.14(e)).
\end{proof}

The theorem is the per-point form in which the large fields enter the recursion of Section~\ref{sec:cauchy}: the factor $s^{\mathrm{LF}}_{n(X)}$ is attached to the birth scale of the component and not to the scale of the point, so that its sum over all points of a history, counted by $P(X)$, is controlled by the weight of Proposition~\ref{prop:large-field-weight}. The stopping-scale test functions of the all-small-field pair are the one place where a shell does not lie on a history, and they are treated directly.

\begin{proposition}[The Stopping-Scale Shells]\label{prop:stopping-shells}
Under the hypotheses of Theorem~\ref{thm:lf-comparison}, replacing the stopping-scale test function $\chi_{K'}(\mathbb T)$ of run~A's all-small-field term by run~B's, in the numerator and in the normalisation, and the same for run~B, changes $W_{K+1}(\mathcal C) - W_K(\mathcal C)$ by at most $b^{(S)}_KB(\mathcal C)$, with
\begin{equation}\label{eq:stopping-shells}
    b^{(S)}_K \coloneqq 4\kappa_s^{-1}e^{2C_K}N_q^4\delta^f\;,
    \qquad
    \delta^f \coloneqq C_fL^{-(1-2\varepsilon_0)K'}\;,
\end{equation}
which is summable in $K$.
\end{proposition}

\begin{proof}
The term $\ell_0$ carries the test function \eqref{eq:test-function} of the stopping step, a product over at most $N_q^4$ cubes, built on each run's fine lattice at the common threshold $\varepsilon_{K'}$, and no other test function which is not identically one (Supplement, Proposition~S.2). By Corollary~\ref{cor:tested-quantities} at $j = K'$, where $\varepsilon/a_{K'} = L^{-K'}$ and $\delta_{K'}$ does not enter, and by Lemma~\ref{lem:soft-shells}, the shell of each run is at most $N_q^4\delta^f/\kappa_s$ times that run's term with sharp test functions at the enlarged thresholds, whose integral is at most $e^{C_K}$ by the proof of Proposition~\ref{prop:large-field-weight}. Relative to $Z_{\ell_0} \geq e^{-C_K}$ this is at most $N_q^4\delta^f\kappa_s^{-1}e^{2C_K}$, and the change of the numerator is at most $\mathfrak b(\mathcal C) \leq B(\mathcal C)$ times this by Lemma~\ref{lem:term-bound}. Arguing as in \eqref{eq:removal} for the numerator and the normalisation of each run gives the factor $4B(\mathcal C)$. The sum over $K$ converges because $C_K$ and $N_q$ do not depend on $K$.
\end{proof}

\section{The Limit along the Full Sequence}\label{sec:cauchy}

We can now compare $W_{K+1}(\mathcal C)$ with $W_K(\mathcal C)$. By Theorem~\ref{thm:stopping-identity} each of them is the average of the normalised dressed loops against the density of its run on the common space $G^{\mathcal B}$, so that the comparison is a comparison of two densities and of two families of functions on one space. The difference between the densities is, at first sight, of order one: the exponents of the two runs contain the effective actions of all the steps, and the terms of the first steps, produced on the fine lattices, differ by quantities which do not tend to zero. What saves the comparison is the decay of Ba{\l}aban's effective actions in the age of their terms, by which a term produced $m$ steps before the stopping step contributes at most $E_1L^{-\beta m}$ per point of the stopping lattice \cite[(2.43), p.~263]{BalabanCRE}. The old terms can therefore be bounded without comparing them, and only the recent ones, produced at scales at which the fine lattice is far away, have to be compared. The comparison of the recent terms is a recursion: the difference of the terms produced by one step is bounded by the differences of the terms which enter it, through a contraction which is Ba{\l}aban's, and by sources which are the lattice artefacts of Sections~\ref{sec:walk-factor-difference} and~\ref{sec:large-field-comparison}. Its solution decays geometrically in the scale, up to a term from the fine scales which is a negative power of $K$.

\begin{theorem}[Comparison of Two Cutoffs]\label{thm:cauchy}
Under the hypotheses of Theorem~\ref{thm:main} there are an index $K_1 \geq K_0(\mathcal C)$ and nonnegative numbers $a_K$, given for $K \geq K_1$ by \eqref{eq:cauchy-sequence} and for $K_0(\mathcal C) \leq K < K_1$ by $a_K \coloneqq 2B(\mathcal C)$, such that $\sum_Ka_K < \infty$ and $\abs{W_{K+1}(\mathcal C) - W_K(\mathcal C)} \leq a_K$ for every $K \geq K_0(\mathcal C)$. The index $K_1$ depends on $\mathcal C$ through the quantities on which $B(\mathcal C)$ depends, and on $q$, $T/u$, $\bar g$, and the constants of the construction.
\end{theorem}

For the finitely many $K < K_1$ there is nothing to prove, since $\abs{W_K(\mathcal C)} \leq B(\mathcal C)$ for every $K \geq K_0(\mathcal C)$ by Theorem~\ref{thm:uniform-bound}. The rest of this section proves the bound for $K \geq K_1$, first for the all-small-field pair and then for the other terms. Throughout, the two runs are those of Section~\ref{sec:two-runs}.

\subsection{The Exponents of the All-Small-Field Pair}\label{sec:exponent-difference}

On the term $\ell_0$ of either run the large-field regions are empty, the coupling \eqref{eq:position-coupling} is constant on the last domain, and the exponent of the density at $V$ is, by \eqref{eq:effective-action}--\eqref{eq:regular-part},
\begin{equation}\label{eq:trivial-exponent}
    -\frac{1}{g_{\mathrm{stop}}^2}A_{\mathrm{fine}}(U(V)) + \sum_j\mathrm{irr}_j(U(V)) + \sum_j\bigl[\mathbf R^{(j)}(U(V)) - \mathbf R^{(j)}(1)\bigr] - E\;,
\end{equation}
with $\mathrm{irr}_j(U) \coloneqq \mathbf E^{(j)}(U) - \mathbf E^{(j)}(1) - \beta_jA(\phi_j, U)$ the quantity bounded by \eqref{eq:balaban-theorem-two}, where for run~A $g_{\mathrm{stop}} = g^A_{K'}$, $A_{\mathrm{fine}} = A_\varepsilon$, and $U = U^A(V)$ is the minimiser on the lattice of spacing $\varepsilon$ with $M^{K'}U = V$, and for run~B $g_{\mathrm{stop}} = g^B_{K'+1}$, $A_{\mathrm{fine}} = A_{\varepsilon/L}$, and $U = U^B(V)$ with $M^{K'+1}U = V$. To compare the terms of the same shared scale we carry run~A's arguments to run~B's fine lattice. For a localisation domain $X$ and a term $\mathbf E^{B,(j+1)}(X, \cdot, z)$ of run~B let
\begin{equation}\label{eq:lifted-term}
    \mathcal G_j(X, (\mathbf U, \mathbf J), z) \coloneqq \mathbf E^{B,(j+1)}\bigl(X, (\iota_X\mathbf U, Q^{\mathrm{pl}}_1(\iota_X\mathbf U)^*\mathbf J), z\bigr)\;,
    \qquad
    \Delta_j \coloneqq \mathcal G_j - \mathbf E^{A,(j)}\;,
\end{equation}
where $Q^{\mathrm{pl}}_1(W)$ is the plain linearised one-step average at $W$ and $\mathbf J$ is the current variable of the spaces of \cite[(1.10)--(1.16), p.~262]{BalabanRG1}; both are evaluated at the current $\mathbf J = J(U)$ of the configuration, so that $\Delta_j$ is a functional of fields on the lattice of spacing $\varepsilon$. We write $\mathrm{irr}(\Delta_j)(U) \coloneqq \Delta_j(U) - \Delta_j(1) - \Pi^\Delta A(h_z, U) - N_j(U)$, with $\Pi^\Delta$ the isotropic coefficient of $\Delta_j$ in the sense of \cite[(3.57)--(3.61), pp.~281--282]{BalabanCRE}, $h_z$ the tent functions of \cite[(3.65), p.~283]{BalabanCRE}, and $N_j$ the part of the second moment of $\mathcal G_j$ off its isotropic coefficient: since $\mathcal G_j$ evaluates every domain at its own lift, the Euclidean covariance by which Ba{\l}aban shows that this part vanishes \cite[(3.60)--(3.61), p.~282]{BalabanCRE} is not available for it, and $N_j$ is kept as a term of its own, a sum of the $36$ plane-pair actions $\tr F_{\kappa\mu}F_{\lambda\nu}$ of \cite[(3.57)]{BalabanCRE}. Its coefficients are small for a reason which is specific to the lift: for affine data the minimisers of both runs are affine fields and the global lift carries the affine field of run~A to that of run~B, so that the second moment of $\mathcal G_j$ would be that of run~B's own family if every domain were evaluated at the global lift; what remains is the difference between the localised and the global lift on a domain, which decays exponentially in $M_1L^j$, and each coefficient of $N_j$, as well as $\Pi^{\mathcal G} - \beta^B_{j+1}$, is at most $C_\Pi E_0e^{-\delta_cM_1L^j}$ (Supplement, Proposition~S.15(d)). Similarly $\mathrm{irr}^{\mathcal G}_j(U) \coloneqq \mathcal G_j(U) - \mathcal G_j(1) - \Pi^{\mathcal G}A(h_z, U) - N_j(U)$, and $\mathcal G_j(1) = \mathbf E^{B,(j+1)}(1)$ because $\iota_X1 = 1$ and $J(1) = 0$.

\begin{lemma}[Exponent Difference]\label{lem:exponent-difference}
At every $V$ in the support of both terms $\ell_0$, the exponent of run~B minus that of run~A is
\begin{multline}\label{eq:exponent-difference}
    D(V) = -\delta_{K'}A_{\varepsilon/L}(U^B) - (g^A_{K'})^{-2}\bigl[A_{\varepsilon/L}(U^B) - A_\varepsilon(U^A)\bigr] + \mathrm{irr}^B_1(U^B) + \sum_{j=1}^{K'}\mathrm{irr}(\Delta_j)(U^A)\\
    + \sum_{j=1}^{K'}\bigl[\mathrm{irr}^B_{j+1}(U^B) - \mathrm{irr}^{\mathcal G}_j(U^A)\bigr] + D^{\mathbf R}(V) - (E^B - E^A)\;,
\end{multline}
where $D^{\mathbf R}$ consists of the same three groups for the $\mathbf R$-terms and $E^B - E^A$ is a number.
\end{lemma}

\begin{proof}
The coupling of run~B at its stopping step is $(g^B_{K'+1})^{-2} = (g^A_{K'})^{-2} + \delta_{K'}$ by \eqref{eq:coupling-difference}, which gives the first two terms. The regular part of run~B consists of its terms of the extra step, $\mathrm{irr}^B_1$, and of those of the steps $j + 1$, $1 \leq j \leq K'$. For the latter, $\mathrm{irr}^B_{j+1}(U^B) - \mathrm{irr}^A_j(U^A) = [\mathrm{irr}^B_{j+1}(U^B) - \mathrm{irr}^{\mathcal G}_j(U^A)] + [\mathrm{irr}^{\mathcal G}_j(U^A) - \mathrm{irr}^A_j(U^A)]$, and the second bracket is $\mathrm{irr}(\Delta_j)(U^A)$, because the subtractions at $U = 1$, of the isotropic part, and of $N_j$ are linear and $\Pi^\Delta = \Pi^{\mathcal G} - \beta^A_j$, the coefficient of run~A being absorbed into the coupling by \eqref{eq:position-coupling}. The same holds for the $\mathbf R$-terms.
\end{proof}

The groups of \eqref{eq:exponent-difference} are of four kinds: the coupling term, which is controlled by the recursion through $\delta_{K'}$; the difference of the two constrained minima of the action, a lattice artefact; the terms of the extra step and the old terms, which are bounded directly; and the differences of the recent terms of the same shared scale, which are the subject of the recursion. Since $MU^B$ satisfies run~A's constraint, the second group is
\begin{equation}\label{eq:minima-difference}
    A_{\varepsilon/L}(U^B) - A_\varepsilon(U^A) = \bigl[A_\varepsilon(MU^B) - A_\varepsilon(U^A)\bigr] - b(U^B)\;,
    \qquad
    b(W) \coloneqq A_\varepsilon(MW) - A_{\varepsilon/L}(W)\;,
\end{equation}
the gap between two configurations of run~A's problem and the blocking artefact.

\begin{proposition}[Classical Depth]\label{prop:classical-depth}
Let $0 < \varepsilon_0 < 1/2$. There are $k_0(\varepsilon_0)$ and $C_b$ such that for $K' \geq k_0(\varepsilon_0)$ and every $V$ in the support of both terms $\ell_0$,
\begin{equation}\label{eq:classical-depth}
    \abs{A_{\varepsilon/L}(U^B(V)) - A_\varepsilon(U^A(V))} \leq 2C_bL^{-K'}\varepsilon_1^2N_q^4\;.
\end{equation}
\end{proposition}

\begin{proof}
The blocking artefact is a sum over the coarse plaquettes of local Taylor remainders, each bounded by the regularity of its cube \cite[Theorem~1, (9), (10), p.~279]{BalabanVAR}, which gives $\abs{b(U^B)} \leq C_bL^{-K'}\varepsilon_1^2N_q^4$, and the gap is at most $C_{\mathrm{gap}}(\varepsilon_0)L^{-(2-2\varepsilon_0)K'}\varepsilon_1^2N_q^4$, which is at most the same bound for $K' \geq k_0(\varepsilon_0)$; both are Supplement, Proposition~S.13(a), (b).
\end{proof}

\subsection{The Direct Bounds}\label{sec:direct-bounds}

The terms of the extra step and the terms of large age are bounded without comparing them.

\begin{lemma}[The Extra Step and the Old Tail]\label{lem:direct-bounds}
Let $\bar g$ be as in Theorem~\ref{thm:main}, and let $1 \leq m \leq K'$. On the support of both terms $\ell_0$,
\begin{equation}\label{eq:extra-step}
    \abs{\mathrm{irr}^B_1(U^B)} + \abs{\mathbf R^{B,(1)}(U^B) - \mathbf R^{B,(1)}(1)} \leq \bigl(E_1L^{-\beta K'} + R_1(bK - C_W)^{-\kappa_0/2}\bigr)N_q^4\;,
\end{equation}
and the differences of the regular and $\mathbf R$-parts of the shared scales $j \leq K' - m$ satisfy
\begin{equation}\label{eq:old-tail}
    \sum_{j \leq K'-m}\abs{\mathrm{irr}^B_{j+1}(U^B) - \mathrm{irr}^A_j(U^A)} \leq \frac{2E_1N_q^4L^{-\beta m}}{1 - L^{-\beta}}\;,
    \qquad
    \sum_{j \leq K'-m}\abs{\Delta\mathbf R_j} \leq C_{\mathbf R}N_q^4(1 + m)^{1-\kappa_0/2}\;,
\end{equation}
with $\Delta\mathbf R_j$ the difference of the $\mathbf R$-parts of the shared scale $j$ and $C_{\mathbf R}$ a constant of $R_1$, $b$, and $\kappa_0$.
\end{lemma}

\begin{proof}
On $\ell_0$ the determining set is the stopping lattice itself, $\Gamma_{K'} = \mathbb T_{j_q}$ with $N_q^4$ points, so that in \eqref{eq:balaban-theorem-two} at the generation $j$ only the term $n = K'$ (respectively $K' + 1$ for run~B) occurs, and it gives $E_1L^{-\beta(K'-j)}N_q^4$ for the regular part and $R_1g_j^{\kappa_0}N_q^4$ for the $\mathbf R$-part. At the extra step $j = 1$ of run~B this is $E_1L^{-\beta K'}N_q^4$, and $(g^B_1)^{-2} \geq \bar g^{-2} + bK - C_W$ by \eqref{eq:coupling-window} for the target $(K + 1, \bar g)$, which gives \eqref{eq:extra-step}. For \eqref{eq:old-tail} the triangle inequality and the same bound for each run give $\sum_{j \leq K'-m}2E_1L^{-\beta(K'-j)}N_q^4$, and the geometric series gives the first bound. For the second, $g_j^{-2} \geq g_{K'}^{-2} + b(K' - j) - C_W$ by \eqref{eq:coupling-window} gives $g_j^{\kappa_0} \leq C(1 + K' - j)^{-\kappa_0/2}$ for $\bar g$ small, and $\sum_{i \geq m}(1 + i)^{-\kappa_0/2} \leq C(1 + m)^{1-\kappa_0/2}$ since $\kappa_0 > 2$.
\end{proof}

We use \eqref{eq:old-tail} at the age $m_K \coloneqq \floor{K^{1/2}}$, where it is summable in $K$ because $\kappa_0 \geq 13$ (Appendix~\ref{app:constants}), and it uses no information on the differences of the terms. The terms of the shared scales $j > K' - m_K$ are compared, and this is where Ba{\l}aban's proof of \eqref{eq:balaban-theorem-two} has to be run for a difference.

\subsection{The Frozen Differences}\label{sec:frozen-differences}

The terms $\mathbf E^{A,(j)}$ and $\mathcal G_j$ are two families of terms of the generation $j$ on the same lattices, with the localisation, analyticity, gauge invariance, and bound of Section~\ref{sec:density}, for $\mathcal G_j$ on a common polydisc of reduced radii $\alpha'_{i,m} \coloneqq \alpha_{i,m}/(4B_3)$. For their difference we use the \emph{class norm} $d_j$, the least number with $\abs{\Delta_j(X, (\mathbf U, J(\mathbf U)), z)} \leq d_jE_0e^{-\kappa d_j(X)}$ on the evaluation sets of this polydisc, the sets of its pairs at which the constructions below read the terms, namely the minimisers of real and of complex data with their currents and the configurations which the steps build from them (Supplement, Proposition~S.15), and $D_k \coloneqq \sum_{j \leq k}L^{-\beta'(k-j)}d_j$ for its carried sum. Here
\begin{equation}\label{eq:age-exponent}
    \beta' \coloneqq \tfrac12\min(\beta_H, 1 - \beta_E) \geq \tfrac{19}{40}\;,
    \qquad
    \theta_V \coloneqq \frac{\beta'}{2(4 + \beta')}\;,
\end{equation}
with $\beta_H \geq 19/20$ the H\"older index in the regularity bounds of \cite[(3.31), p.~276]{BalabanRG1}, which Ba{\l}aban states for every index less than one and which hold up to the index one for the minimisers \cite[Theorem~1, (9), p.~279]{BalabanVAR}, and $\beta_E \leq 1/20$ the positive number of \cite[(3.67), p.~283]{BalabanCRE}, which may be taken arbitrarily small at the cost of a larger constant. We write $d^{\mathbf R}_j$ for the class norm of the difference of the $\mathbf R$-terms of the shared scale $j$ and $S_k \coloneqq \sum_{j \leq k}d^{\mathbf R}_j$.

\begin{proposition}[Theorem~2 for a Difference]\label{prop:theorem-two-difference}
Let the constants satisfy the conditions of Appendix~\ref{app:conditions} and $\bar g$ be small. There are constants $E'_1$, $E_V$, and $C_\alpha$, independent of $j$, $k$, $K$, and the torus, such that for $z \in \Lambda^0_j \cap (\Omega_n \setminus \Omega_{n+1})$ and every regular background $U$ of the class of the step,
\begin{equation}\label{eq:theorem-two-difference}
    \abs{\mathrm{irr}(\Delta_j)(U)} \leq d_j\,C_\alpha E_0(L^{j-n})^{4+\beta'}\;,
\end{equation}
and on the all-small-field pair, at $n = K'$,
\begin{equation}\label{eq:variation-group}
    \abs{\mathrm{irr}^B_{j+1}(U^B) - \mathrm{irr}^{\mathcal G}_j(U^A)} \leq E_V\Bigl[\Bigl(\frac{\hat\sigma_{K',j}}{r_{0,K'}}\Bigr)^{\theta_V} + \bigl(C_\Pi\varepsilon_1^2e^{-\delta_cM_1L^j}\bigr)^{\theta_V}\Bigr](L^{j-K'})^{4+\beta'/2}\;,
\end{equation}
where $\hat\sigma_{K',j}$ is the distance between $U^B$ and the lift $\iota(U^A)$ in the norm of the evaluation sets, bounded by Proposition~\ref{prop:minimiser-distance} and the tail of the localised lift (Supplement, Proposition~S.15(d)), $r_{0,K'}$ the radius of these sets, and $C_\Pi E_0e^{-\delta_cM_1L^j}$ the bound of $\Pi^{\mathcal G} - \beta^B_{j+1}$ and of the coefficients of $N_j$, with $\delta_c$ the rate of Supplement, Proposition~S.12. The same holds for the $\mathbf R$-terms with $d^{\mathbf R}_j$ and without the factor of the age.
\end{proposition}

\begin{proof}
For \eqref{eq:theorem-two-difference} we run Ba{\l}aban's proof of \cite[Theorem~2, (2.43)]{BalabanCRE}, which occupies \cite[pp.~279--283]{BalabanCRE} and rests on \cite[Sects.~3--4]{BalabanRG1}, with the difference $\Delta_j$ in place of a single family; the proof is linear in the family, and it uses of the family only its localisation, its analyticity on the polydiscs of \cite[(3.13)--(3.14)]{BalabanRG1} with the radii of the step, and its bound, all of which $\Delta_j$ has with $d_jE_0$ in place of $E_0$, except the covariance which removes the plane-resolved moment, which is kept as $N_j$. The H\"older regularity of the background, of index $\beta_H$, gives the power $4 + \beta'$ in place of $4$. This is Supplement, Proposition~S.15(a)--(b). For \eqref{eq:variation-group} the two values are the irrelevant parts of one term of run~B at two configurations, $U^B$ and the lift $\iota(U^A)$, each bounded by the age bound for one term; interpolating the configuration along a line of length $\hat\sigma_{K',j}$ in the norm of the evaluation sets and using Cauchy's estimate on a disc of radius $r_{0,K'}$, then interpolating between the bound of the age and the bound of the distance with the weight $\theta_V$, gives \eqref{eq:variation-group}; this is Supplement, Proposition~S.15(c).
\end{proof}

\subsection{The Coefficients of the Two Runs}\label{sec:coefficient-difference}

The recursion of the next subsection contains the differences $\delta_k$ of the couplings of the two runs, which are sums of differences of the coefficients of the flow. The coefficient of a step is the second moment of its regular term at the trivial configuration \cite[(3.57)--(3.64), pp.~281--283]{BalabanCRE}, and one is tempted to bound the difference of two coefficients by the class norm of the difference of the two terms. This fails. The class norm $d_j$ bounds $\Delta_j$ on polydiscs whose radii $\alpha'_{i,j}$ are proportional to $\hat g_j$, up to powers of its logarithm, and Cauchy's estimate for a second derivative on such a polydisc costs the inverse square of the radius,
\begin{equation}\label{eq:radius-cost}
    \abs{\Pi^\Delta_j} \leq C\hat g_j^{-2}\,d_jE_0\;.
\end{equation}
Nothing better follows from $d_j$: the function $z \mapsto z^2$ has the supremum $r^2$ on the disc of radius $r$ and the second coefficient one, and in the same way a multiple of the action $A(h_z, \cdot)$ with a coefficient of the order $\hat g_j^{-2}d_jE_0$ has a class norm of the order $d_j$. With \eqref{eq:radius-cost} the recursion would not close, because the differences of the couplings would feed back into every step with a kernel which is not summable over the later steps. For a single run the coefficient is bounded all the same (Supplement, Proposition~S.4), and the reason is a structure of the regular terms which we now use for the difference of two runs.

For a run and a complex number $s$ let $\mathbf E^{(j)}(s)$ be the family of regular terms which the steps of the run produce when the $\mathbf R$-terms which enter the regular part \cite[p.~271]{BalabanCRE} are multiplied by $s$, the couplings, the cutoffs, and all other terms being those of the run (Supplement, Lemma~3.4). At $s = 1$ these are the regular terms of the run. At $s = 0$ they are the terms of the small-field recursion of \cite{BalabanRG1,BalabanRG2} along the couplings of the run, and by Ba{\l}aban's Theorem~3 \cite[p.~264]{BalabanRG1} they are analytic and bounded on spaces whose radii $\alpha_0$, $\alpha_1$ do not depend on the coupling. For $\abs s \leq \varrho_{j-1}$, where
\begin{equation}\label{eq:family-radius}
    \varrho_k \coloneqq 1 + \hat g_k^{-2}\;,
\end{equation}
they satisfy the bounds of the regular terms on the spaces of the tube, and they differ from the terms at $s = 0$ by at most $C_\rho g_{j-1}^{\kappa_0-4}E_0e^{-\kappa d_j(X)}$, because the $\mathbf R$-terms are of the order $g^{\kappa_0}$ (Supplement, Lemma~3.5). Let $\mathcal G_j(s)$ be the family of run~B at the lift, as in \eqref{eq:lifted-term}, and $\Delta_j(s) \coloneqq \mathcal G_j(s) - \mathbf E^{A,(j)}(s)$, so that $\Delta_j(1) = \Delta_j$. Besides $d_j$ we use two class norms:
\begin{itemize}
\item $d^\sharp_j$, the least number with $\abs{\Delta_j(X, \cdot, z; s)} \leq d^\sharp_jE_0e^{-\kappa d_j(X)}$ on the evaluation sets of the polydisc of the tube, for every $\abs s \leq \varrho_{j-1}$;
\item $d^\bullet_j$, the least number with $\abs{\Delta_j(X, \cdot, z; 0)} \leq d^\bullet_jE_0e^{-\kappa d_j(X)}$ on the evaluation sets of the polydisc of the fixed radii $\alpha^\bullet_i \coloneqq \alpha_i/(4B_3)$.
\end{itemize}
Since $\varrho_{j-1} \geq 1$, $d_j \leq d^\sharp_j$. We write $D^\sharp_k$ and $D^\bullet_k$ for the carried sums, defined as $D_k$, and $u_j \coloneqq \max(d^\bullet_j, d^\sharp_j)$.

\begin{lemma}[The Coefficients of the Two Runs]\label{lem:coefficient-difference}
There is a constant $C_\beta$, which does not depend on $j$, $K$, the couplings, or the torus, such that at every shared scale $j$
\begin{equation}\label{eq:coefficient-difference}
    \abs{\beta^B_{j+1} - \beta^A_j} \leq C_\beta u_j + C_\Pi E_0e^{-\delta_cM_1L^j}\;.
\end{equation}
\end{lemma}

\begin{proof}
By the definition of $\Delta_j$, $\Pi^\Delta_j = \Pi^{\mathcal G}_j - \beta^A_j$, so that $\beta^B_{j+1} - \beta^A_j = \Pi^\Delta_j - (\Pi^{\mathcal G}_j - \beta^B_{j+1})$, and the bracket, the isotropic part of the moment defect of the lifted family, is at most $C_\Pi E_0e^{-\delta_cM_1L^j}$ (Supplement, Proposition~S.15(d)). The coefficient is linear in the family, and we split it as
\begin{equation}\label{eq:coefficient-split}
    \Pi^\Delta_j = \Pi\bigl(\Delta_j(0)\bigr) + \Pi\bigl(\Delta_j(1) - \Delta_j(0)\bigr)\;,
\end{equation}
where $\Pi(F)$ denotes the isotropic coefficient of a family $F$. The family $\Delta_j(0)$ is bounded by $d^\bullet_jE_0e^{-\kappa d_j(X)}$ on a polydisc about the trivial configuration whose radius is a fixed multiple of $\min(\alpha^\bullet_0, \alpha^\bullet_1)$. Cauchy's estimate for its second moment costs therefore a constant which does not depend on the coupling, and $\abs{\Pi(\Delta_j(0))} \leq C^\bullet_\beta d^\bullet_j$. For the second term we use the analyticity in $s$. At every point of the evaluation sets of the tube the function $s \mapsto \Delta_j(X; s)$ is bounded by $d^\sharp_jE_0e^{-\kappa d_j(X)}$ on the disc of radius $\varrho_{j-1}$, and for $s \in [0, 1]$ the disc of the radius $\hat g_{j-1}^{-2}$ about $s$ lies in it, so that Cauchy's estimate for the derivative in $s$ gives
\begin{equation}\label{eq:family-difference}
    \abs{\Delta_j(X; 1) - \Delta_j(X; 0)} \leq \hat g_{j-1}^2\,d^\sharp_jE_0e^{-\kappa d_j(X)}\;.
\end{equation}
The second moment of this family on the polydisc of the tube costs a constant times $\hat g_j^{-2}$, as in \eqref{eq:radius-cost}, and $\hat g_{j-1}^2\hat g_j^{-2} \leq 1$ because $\hat g_j$ increases with $j$; hence $\abs{\Pi(\Delta_j(1) - \Delta_j(0))} \leq C^\sharp_\beta d^\sharp_j$. This is \eqref{eq:coefficient-difference} with $C_\beta \coloneqq 2\max(C^\bullet_\beta, C^\sharp_\beta)$, and the constants are those of Supplement, Proposition~S.15(e).
\end{proof}

The lemma replaces the radius of the tube by two radii which the construction provides: that of Ba{\l}aban's small-field spaces, for the part of the regular terms which does not involve the $\mathbf R$-terms, and that of the disc of $s$, for the part which does. The second is available because the $\mathbf R$-terms are small: a disc of the radius $\hat g^{-2}$ costs two of the $\kappa_0$ powers of the coupling which they carry. The price is that the two norms $d^\bullet$ and $d^\sharp$ have to be carried through the recursion in the place of $d$, and that the $\mathbf R$-terms enter the recursion of $d^\sharp$ with the factor $\varrho_k$.

\subsection{The Generated Difference of One Step}\label{sec:recursion}

At a shared scale $k$ the two runs apply the same step to inputs on the same lattice, with the same block data and the same fluctuation variable, and the step's output is the logarithm of a fluctuation integral. The two integrals differ by the interaction, which collects the fluctuation parts of the frozen terms and of the fine actions, by the Gaussian form, which is built from the propagators of the run, and by the cutoff of the fluctuation field. Let $\mathbf P^A$ and $\mathbf P^B$ be the two interactions, $\Delta^{(k)}_A$ and $\Delta^{(k)}_A + \delta\Delta^{(k)}$ the two quadratic forms of the fluctuation field, and for $t \in [0, 1]$ put $\mathbf P_t \coloneqq \mathbf P^A + t(\mathbf P^B - \mathbf P^A)$, $\Delta_t \coloneqq \Delta^{(k)}_A + t\delta\Delta^{(k)}$, and
\begin{equation}\label{eq:interpolated-integral}
    Z_t \coloneqq \int\chi^{A,(k)}\exp\Bigl[-\frac12\langle A, C^*\Delta_tCA\rangle + \mathbf P_t\Bigr]\,\dd A\;,
\end{equation}
with run~A's cutoff in both, so that run~B's integral is $Z_1$ plus the shell term of Lemma~\ref{lem:coupling-shells}.

\begin{lemma}[Interpolation Identity]\label{lem:interpolation}
Let $k$ be a shared scale at which the number $r_k$ of \eqref{eq:gaussian-depth} satisfies $C_{\mathrm{cov}}r_k \leq \alpha^G_4$, with $\alpha^G_4$ the room of the Gaussian integral of Supplement, Proposition~S.16(a). On the class of the scale $k + 1$ and for every $t \in [0, 1]$ the integral $Z_t$ does not vanish, and with $\langle\cdot\rangle_t$ the expectation of the normalised integrand of $Z_t$ and $\log Z_t$ the branch which is continuous in $t$,
\begin{equation}\label{eq:interpolation-identity}
    \log Z_1 - \log Z_0 = \int_0^1\Bigl[\langle\mathbf P^B - \mathbf P^A\rangle_t - \frac12\bigl\langle\langle A, C^*\delta\Delta^{(k)}CA\rangle\bigr\rangle_t\Bigr]\,\dd t\;.
\end{equation}
For every $t$ the normalised integrand is of the kind which Ba{\l}aban's cluster expansion of the step treats: its interaction $\mathbf P_t$ satisfies the inductive bounds, and its quadratic form is run~A's plus the form $t\,\delta\Delta^{(k)}$, which is a small quadratic term of the exponent.
\end{lemma}

\begin{proof}
We begin with the last statement, on which the others rest. The interaction $\mathbf P_t$ is a convex combination of two interactions which satisfy the inductive bounds \cite[(1.18), p.~263]{BalabanRG1}, \cite[(2.31), (2.42)]{BalabanCRE}, and these bounds define convex sets. For the quadratic form we do not argue through its positivity, which says nothing at the complex backgrounds of the class. We write $\Delta_t = \Delta^A + t\,\delta\Delta^{(k)}$ and treat the second term as a part of the exponent: by Supplement, Proposition~S.17, its kernel is at most $r_k$ times the bound of the kernel of the form itself, with exponential decay, so that $\abs{\frac12t\langle A, C^*\delta\Delta^{(k)}CA\rangle} \leq \frac12C_{\mathrm{cov}}r_k\sum_b\abs{A(b)}^2$, and a quadratic term with a coefficient within the room $\alpha^G_4$ is absorbed into the Gaussian integral of run~A by the steps \cite[(2.14)--(2.25), pp.~15--17]{BalabanRG2}, as in Supplement, Proposition~S.16(a). The integrand of $Z_t$ is thus one to which Ba{\l}aban's expansion applies, with run~A's covariance, at real and at complex backgrounds of the class.

In this expansion the integral which normalises the expectation, taken with the Gaussian measure of run~A, ``is represented as $\exp\tilde{\mathbf E}^{(k+1)}(\Lambda_{k+1})$'', with $\tilde{\mathbf E}^{(k+1)}$ the convergent sum of analytic terms of the exponentiated cluster expansion \cite[p.~277]{BalabanCRE}; since the Gaussian normalisation of run~A is the inverse square root of the determinant of an invertible form, $Z_t \neq 0$. The integrand is differentiable in $t$, with the derivative equal to the integrand times $\mathbf P^B - \mathbf P^A - \frac12\langle A, C^*\delta\Delta^{(k)}CA\rangle$, which on the support of the cutoff is the sum of a bounded function and a quadratic form in $A$ and is dominated, uniformly in $t$, by an integrable function, because on the support of the cutoff $\abs A$ is bounded. So $Z_t$ is continuously differentiable, $\partial_t\log Z_t = \partial_tZ_t/Z_t$ is the expectation of the derivative of the exponent, and integrating over $t$ gives \eqref{eq:interpolation-identity}, the device by which Ba{\l}aban separates the terms of a step \cite[(3.26)--(3.27), p.~271]{BalabanCRE}.
\end{proof}

At the steps at which $C_{\mathrm{cov}}r_k > \alpha^G_4$ the lemma is not used. By \eqref{eq:gaussian-depth} below these are the steps $k \leq k_r$ with a number $k_r$ which does not depend on $K$, and there each of the two outputs satisfies the bounds of its class, so that the class norm of their difference is at most $2$; this is the bound \eqref{eq:recursion} below with a depth source $e_{k+1} \geq 2$, and the constant $C_e$ of the depth sources is chosen at least $2\tilde\theta^{-k_r-1}$.

The right side of \eqref{eq:interpolation-identity} is a cluster expansion with one marked activity, the piece of $\mathbf P^B - \mathbf P^A$ or of the quadratic insertion on a domain, and the localised terms of the difference of the step's outputs are the localised terms of this marked expansion. The marked activity consists of the fluctuation parts of the frozen differences of the generations $j \leq k$, of size $D_k + S_k$ in the class norm; of the coupling, through the factors $g_k$ in the backgrounds and the coefficients $g_k^{n-2}$ of the terms of order $n \geq 3$, whose two values differ relatively by at most $g_k^2\abs{\delta_k}$; of the depth sources $e_{k+1}$, which are the differences of the two runs in the objects of the fine lattices; and of the shell of the cutoff. The depth sources need a word. One of them is the difference $\delta\Delta^{(k)}$ of the quadratic forms. The others are the differences of the parts of the third and higher orders in the fluctuation field of the two actions, of the current terms, and of the configurations at which the frozen terms are read. The classical depth of Proposition~\ref{prop:classical-depth} is not among them: it bounds the difference of the two actions at the two minimisers, which enters an exponent with the factor $g_k^{-2}$, and this factor is bounded at the stopping scale only. At the earlier scales the two actions are compared through their expansions in the fluctuation field, in which the terms of the orders zero and one are absent and the term of the second order is the quadratic form; the remaining terms are at most $g_kp_0(\hat g_k)^3$ per point, uniformly in the regularity of the background, and their two-run differences are smaller by a power of $\varepsilon/a_k$ (Supplement, Proposition~S.16(b)). The quadratic forms differ by at most
\begin{equation}\label{eq:gaussian-depth}
    r_k \leq C_{\delta\Delta}\bigl[L^{-\gamma_H(k-1)} + L^{-(1-\varepsilon_0)k}\bigr]\;,
    \qquad
    \gamma_H \coloneqq \frac{\beta\gamma_w}{1 + \beta}\;,
\end{equation}
times the bound of the form itself, by Theorem~\ref{thm:walk-difference} and Supplement, Proposition~S.17, and the shell of the cutoff contributes at most $C\abs{\delta_k}g_k^2e^{-cp_0(\hat g_k)^2}$ per point, by Lemma~\ref{lem:coupling-shells} and the Gaussian tail beyond the radius $\vartheta_k/g_k$. All of these sources are bounded at backgrounds whose regularity is measured by constants and not by the coupling, and they satisfy $e_{k+1} \leq C_e\tilde\theta^{k+1}$ with a constant which does not depend on $K$. The marked-cluster bound of Supplement, Proposition~S.16, then bounds the generated difference per point of the scale $k + 1$ by $\varepsilon^*_k(D_k + S_k) + C_g\varepsilon^*_kg_k^2\abs{\delta_k} + e_{k+1}$. The argument reads the regular terms of the two runs only through their localisation, their analyticity, and their bounds, and it applies in the same way to the two families of Section~\ref{sec:coefficient-difference}. For the families $\Delta(s)$ on the spaces of the tube, uniformly in $\abs s \leq \varrho_k$, the couplings, the cutoffs, and the propagators are those of the two runs, so that $\delta_k$ and $e_{k+1}$ do not depend on $s$, and the only change is that the $\mathbf R$-terms which enter the regular part are multiplied by $s$:
\begin{equation}\label{eq:recursion}
    d^\sharp_{k+1} \leq \varepsilon^*_k(D^\sharp_k + \varrho_kS_k) + C_g\varepsilon^*_kg_k^2\abs{\delta_k} + e_{k+1}\;,
    \qquad
    \varepsilon^*_k \coloneqq C_XC_{\mathrm{pt}}C_{\mathrm{KP}}L^4C^{\mathbf A}_\alpha\bigl[C_T\mu_k + 4a_k\bigr]\;.
\end{equation}
For the families at $s = 0$ on the spaces of the fixed radii no $\mathbf R$-term enters, the expansion of the step is that of \cite[Sect.~2]{BalabanRG2} on these spaces, and the sources are the same differences of propagators, minimisers, and Gaussian forms, which Sections~\ref{sec:walk-factor-difference} and~\ref{sec:large-field-comparison} bound at backgrounds whose regularity is measured by constants and not by the coupling:
\begin{equation}\label{eq:recursion-fixed}
    d^\bullet_{k+1} \leq \varepsilon^*_kD^\bullet_k + C_g\varepsilon^*_kg_k^2\abs{\delta_k} + e_{k+1}\;.
\end{equation}
In both inequalities $\varepsilon^*_k$ and $e_{k+1}$ denote the larger of the coefficients and of the sources of the two settings (Supplement, Proposition~S.16(b)). Here $\mu_k$ is the square of the ratio of the cutoff's supremum to the radius of analyticity of the frozen terms, a negative power of $\log g_k^{-2}$, $a_k$ the bound of the interaction's activities, a positive power of $g_k$ times a power of $\log g_k^{-2}$, and the other constants are those of Supplement, Proposition~S.16, independent of $k$ and $K$. The coefficient of the flow is the marginal part of the step's output, and the difference of the coefficients of the two runs is bounded by Lemma~\ref{lem:coefficient-difference}. From $\delta_K = 0$, which holds because both runs are pinned at $g^A_K = g^B_{K+1} = \bar g$, and from $\delta_j = \sum_{i > j}(\beta^B_{i+1} - \beta^A_i)$, the coupling differences are therefore $\abs{\delta_j} \leq C_\beta\sum_{i > j}u_i + 2C_\Pi E_0e^{-\delta_cM_1L^{j+1}}$, and the second term, which decays faster than any geometric sequence, is counted among the depth sources $e_i$ of \eqref{eq:recursion}; the steps between the stopping step and the unit scale, which both runs perform for the coupling prescription, obey the same recursion. The $\mathbf R$-terms have no decay in their age beyond the volume factor \cite[p.~283]{BalabanCRE}, and their recursion, Supplement, Proposition~S.16(c), together with the renormalisation quotients of the $\mathbf R$-operation, in which Theorem~\ref{thm:lf-comparison} enters through $h_n \coloneqq C_h\tilde\theta^n$, gives
\begin{equation}\label{eq:r-recursion}
    S_k \leq C_S\Bigl[\sum_{j \leq k}c_j\bigl(\mathcal D'_{[j]} + D_{j-1}\bigr) + \sum_{j \leq k}c^B_j\mathcal D^{*\prime}_j + \sum_{j \leq k}c^V_j\mathcal D^{V\prime}_j + C_\nu w_K\Bigr]\;,
\end{equation}
with $c_j$, $c^B_j$, $c^V_j$ constants times $(100MR_j)^4g_j^{\kappa_0}$, whose sum over $j$ is bounded uniformly in $K$ for $\kappa_0 > 2$, $\mathcal D'_{[j]}$, $\mathcal D^{*\prime}_j$, $\mathcal D^{V\prime}_j$ the maxima over the history of the integrands' differences per point, built from the $d_i$, the $\abs{\delta_i}$, the $e_i$, and the $h_i$ (Supplement, Proposition~S.16(d)), and $w_K$ of \eqref{eq:weight-w} below.

We now solve the recursion. Put
\begin{equation}\label{eq:theta-tilde}
    \theta \coloneqq L^{-\beta'} + \varepsilon^*\;,
    \qquad
    \tilde\theta \coloneqq \max\Bigl(\frac{1 + \theta}{2},\, L^{-\gamma_H},\, L^{-(1-2\varepsilon_0)/3},\, L^{-(1-\varepsilon_0)/4},\, L^{-\theta_V/8}\Bigr)\;,
\end{equation}
where $\varepsilon^* \coloneqq \sup_{g \leq (1+\beta_0)\gamma}\varepsilon^*(g)$ is the supremum of the coefficient of \eqref{eq:recursion} as a function of the coupling, which bounds $\varepsilon^*_k$ at every shared scale because $g_k \leq (1 + \beta_0)\bar g$ by \eqref{eq:tube-two-six}; and put
\begin{equation}\label{eq:weight-w}
    \ell_K \coloneqq \bigl(\log(bK + \bar g^{-2})\bigr)^{4r+2p_0}\;,
    \qquad
    \mu_K \coloneqq K^{3-\kappa_0/2}\ell_K\;,
    \qquad
    w_K \coloneqq (K + 1)^2\bigl(\tilde\theta^{K'/8} + 2^{-K'/16}\bigr) + \mu_K\;.
\end{equation}

\begin{proposition}[Solution of the Recursion]\label{prop:recursion-solution}
Let $\gamma$ satisfy the conditions (T1) of \eqref{eq:t-one}. Then $\tilde\theta < 1$, and for every shared scale $i$
\begin{equation}\label{eq:recursion-solution}
    d_i \leq u_i \leq 2C_s(i\tilde\theta^i + w_K)\;,
    \qquad
    S_i \leq C''_Sw_K\;,
    \qquad
    \abs{\delta_i} \leq 2C_sC_\beta\bigl[2(1 - \tilde\theta)^{-2}(i + 1)\tilde\theta^{i+1} + (K + 1 - i)w_K\bigr]\;,
\end{equation}
with constants $C_s$ and $C''_S$ which do not depend on $K$. In particular, at the stopping step,
\begin{equation}\label{eq:stopping-differences}
    D_{K'} + S_{K'} + C\varepsilon_{K'}^2\abs{\delta_{K'}} + e_{K'} \leq C_a\bigl[(K + 1)\tilde\theta^{K'} + (j_q + 2)w_K\bigr]\;,
\end{equation}
with a constant $C_a$ which does not depend on $K$ and depends on $j_q$ at most as $L^{C\varepsilon_0j_q}$ times a polynomial.
\end{proposition}

\begin{proof}
In \eqref{eq:r-recursion} the differences of the regular terms are those of the two runs, which are bounded by $d_i \leq u_i$. Substituting \eqref{eq:r-recursion} and the bound on $\abs{\delta_i}$ into \eqref{eq:recursion} and \eqref{eq:recursion-fixed} gives $u \leq \mathcal Tu + s$ for $u_i = \max(d^\bullet_i, d^\sharp_i)$, with an operator $\mathcal T$ on sequences which is positive, monotone, and positively homogeneous, and a source $s$ built from the $e_i$, the $h_i$, and $C_\nu w_K$, with $s_1 \coloneqq 2$ for the first step, at which both runs' outputs satisfy the bound of their class. Put $m_i \coloneqq i\tilde\theta^i + w_K$ and $\norm{u}_m \coloneqq \sup_iu_i/m_i$. Since $\tilde\theta > L^{-\beta'}$, $\sum_{j \leq i}L^{-\beta'(i-j)}m_j \leq (1 - L^{-\beta'}/\tilde\theta)^{-1}m_i$, and $m_i \leq \tilde\theta^{-1}m_{i+1}$. The factor $\varrho_i$ of \eqref{eq:recursion} is paid by the $\mathbf R$-terms: for $j \leq i$, $\varrho_i \leq 1 + \hat g_j^{-2} \leq \frac52g_j^{-2}$ by Lemma~\ref{lem:shared-parameters}(iii), so that $\varrho_ic_j$ is at most a constant times $(100MR_j)^4g_j^{\kappa_0-2}$, whose sum over $j$ is still bounded uniformly in $K$, while the direct bounds of the fine scales in the source are multiplied by $\varrho_i \leq (1 + \bar g^{-2} + b' + C_W)(K + 1)$, which is the reason for the powers of $K$ in \eqref{eq:weight-w}. The sums over the fine scales are paid by the factors $\varrho_ic_j$, which for $j \leq K/2$ are at most $C(bK/2)^{1-\kappa_0/2}\ell_K$, and the sums over the coarse scales by $\tilde\theta^j$, which together give at most a constant times $\mu_K \leq w_K$; and the coupling differences enter through $g_i^2\sum_{l > i}m_l \leq (2(1 - \tilde\theta)^{-2} + b^{-1})m_{i+1}$, which uses $g_i^2(K - i) \leq b^{-1}$ by \eqref{eq:coupling-window}. Hence $\norm{\mathcal T}_m \leq \varepsilon^*C_{\mathcal T} \leq \frac14$ by the first condition of \eqref{eq:t-one}, and since the system is finite, $\norm{u}_m < \infty$ and $\norm{u}_m \leq \norm{\mathcal T}_m\norm{u}_m + \norm{s}_m \leq \frac14\norm{u}_m + C_s$. This gives the first bound of \eqref{eq:recursion-solution}, and the others follow from \eqref{eq:r-recursion} and the bound on $\abs{\delta_i}$. The first condition implies $\varepsilon^* \leq (1 - L^{-\beta'})/4$, because $C_{\mathcal T} \geq (1 - L^{-\beta'})^{-1}$, so that $(1 + \theta)/2 \leq (3 + L^{-\beta'})/4 < 1$, and the other entries of \eqref{eq:theta-tilde} are less than one. At $i = K'$, $K'\tilde\theta^{K'} \leq (K + 1)\tilde\theta^{K'}$ and $(K + 1 - K')w_K = (j_q + 1)w_K$, which gives \eqref{eq:stopping-differences}; the details, the constants, and the depth sources $e_i \leq C_e\tilde\theta^i$ are Supplement, Proposition~S.16(e).
\end{proof}

The bound \eqref{eq:stopping-differences} is summable in $K$, because $(K + 1)\tilde\theta^{K'}$ and $(j_q + 2)(K + 1)^2(\tilde\theta^{K'/8} + 2^{-K'/16})$ decay exponentially at fixed $j_q$, and $\mu_K = K^{3-\kappa_0/2}\ell_K$ with $\kappa_0 \geq 13$. The second condition of \eqref{eq:t-one} is the one which allows the recent $\mathbf R$-terms, which are carried through $N'_k \leq N'(\hat g_k)$ steps without contraction \cite[pp.~270--271]{BalabanCRE}, to be summed by a window lemma; the memory $N'$ grows as a double logarithm of the inverse coupling, while $\hat\varepsilon^*$ decays as a negative power of its logarithm. The same recursion holds for the deviation insertions of the dressed loops, which are produced by the same fluctuation integral with the loop's words as further marked activities, with their class norms $d^{\mathrm{ins}}_j$ in place of $d_j$ and an additional term $\varepsilon^Q_kD^{\mathrm{ins}}_k$ from the old vertices of the groups, which enter as numbers, with $\varepsilon^Q(g)$ of order $g^2$ times a power of $\log g^{-2}$ (Supplement, Proposition~S.16(f)). The coefficient $\varepsilon^Q$ depends on the representations, but the recursion of the insertions runs over the steps $k \leq K'$ only, at which $g_k \leq \gamma_q$ by \eqref{eq:stopping-coupling}; under the condition
\begin{equation}\label{eq:insertion-recursion-condition}
    \varepsilon^Q(g)C_{\mathcal T} \leq \frac14
\end{equation}
for every $g \leq \gamma_q$, which holds for $\gamma_q$ small, depending on $c_{\mathcal C}$, the operator of this recursion has the norm at most $(\varepsilon^* + \sup_{g \leq \gamma_q}\varepsilon^Q(g))C_{\mathcal T} \leq \frac12$, by the first condition of \eqref{eq:t-one}, and its solution is $d^{\mathrm{ins}}_i \leq 2C^*_{\mathrm{ins}}m_i$.

\subsection{The All-Small-Field Pair}\label{sec:trivial-pair}

On the common support of the two terms $\ell_0$ the groups of \eqref{eq:exponent-difference}, less the number $E^B - E^A$, are bounded at every $V$ by
\begin{equation}\label{eq:x-k}
    X_K \coloneqq X^{\mathrm{dir}}_K + N_q^4C_a\bigl[(K + 1)\tilde\theta^{K'} + (j_q + 2)w_K\bigr] + X^{\mathrm{var}}_K + X^{\mathbf R}_K + X^{\mathrm{cl}}_K\;,
    \qquad
    x_K \coloneqq e^{X_K} - 1\;,
\end{equation}
where $X^{\mathrm{dir}}_K$ is the sum of the right sides of \eqref{eq:extra-step} and \eqref{eq:old-tail} at $m = m_K$, the second term is \eqref{eq:stopping-differences} summed over the $N_q^4$ points of the stopping lattice, $X^{\mathrm{var}}_K$ is the sum over $j$ of \eqref{eq:variation-group}, $X^{\mathbf R}_K$ the bound of the corresponding group of the $\mathbf R$-terms of the ages less than $m_K$, and $X^{\mathrm{cl}}_K \coloneqq 2C_b(g^A_{K'})^{-2}L^{-K'}\varepsilon_1^2N_q^4$ by \eqref{eq:classical-depth}, together with its counterpart on the young pairs (Supplement, Proposition~S.13(d)). Each of these is summable in $K$: the variation group by $\hat\sigma_{K',j} \leq C(\varepsilon/a_{K'})^{(1-\varepsilon_0)/3}\varepsilon_1$ and the factor of the age, and the coupling factor $(g^A_{K'})^{-2} \leq \bar g^{-2} + b'j_q + C_W$ of the classical group is bounded. The loop's terms are compared on the same support in four parts.

\emph{The leading words.} By Proposition~\ref{prop:minimiser-distance}(i), $U^A = e^{i\varepsilon A'}MU^B$ in a gauge, with $\varepsilon\abs{A'(c)} \leq 4C_G(\varepsilon_0)(\varepsilon/a_{K'})^{2-\varepsilon_0}\varepsilon_1$ on every bond $c$ of the fine lattice of run~A, and by \eqref{eq:average} $MU^B(c) = e^{iY_c}W(c)$, with $W(c)$ the transport of $U^B$ along the $L$ fine bonds of $c$ and $\abs{Y_c} \leq C_Y(\varepsilon/a_{K'})^2\varepsilon_1$ by Lemma~\ref{lem:coarse-loop}. So the holonomies of $C_i$ in the two backgrounds are $\prod_cW(c)$ and, up to a gauge transformation which leaves the trace unchanged, $\prod_ce^{i\varepsilon A'(c)}e^{iY_c}W(c)$. Telescoping inside the normalised trace, with $\norm{\pi(e^{iX}) - 1} \leq c_\pi\norm{X} \leq c_\pi c_h\abs{X}$ for the unitary representation $\pi$ (Section~\ref{sec:coarse-loop}), and summing over the $P_{K'}(C_i)L^{K'}$ fine bonds of $C_i$,
\begin{equation}\label{eq:leading-word-difference}
    \Bigl\lvert\prod_if^{(i),B} - \prod_if^{(i),A}\Bigr\rvert \leq \lambda_K \coloneqq \sum_ic_{\pi_i}c_h\bigl(4C_G(\varepsilon_0) + C_Y\bigr)\varepsilon_1P_{K'}(C_i)L^{-(1-\varepsilon_0)K'}\;,
\end{equation}
because every factor has modulus at most one and $L^{K'}(\varepsilon/a_{K'})^{2-\varepsilon_0} = L^{-(1-\varepsilon_0)K'}$. The words with insertions are covariant at the slots of their insertions, and the same telescoping with the insertions held fixed changes the normalised dressed loop by at most $\lambda_K\prod_b(1 + \norm{Y_b}) \leq \lambda_K\mathfrak b(\mathcal C)$.

\emph{The deviation insertions.} By the recursion for the insertions their difference changes the normalised dressed loop by at most $\mathfrak b(\mathcal C)X_K$.

\emph{The wrap insertions and the interaction weights.} The insertions which wrap around the torus have no deviation bound and are not contracted by the recursion, but their mass decays as $e^{-cN^{(j)}_{\min}/M}$ in the smallest period of the torus at their scale, $N^{(j)}_{\min} \geq c(T/u)L^{K'-j+j_q}$. Those of age at least $m^w_K$, the least $m$ with $a_wL^m \geq 3\log(K + 2)$, are not compared and have total mass at most $\tau^w_K \coloneqq 2A_w\max(L, 3L\log(K + 2)/a_w)(K + 2)^{-3}$ in each run; the younger ones are compared at their birth, by the interpolation of the constituents of their step described below for the loop constant, and then carried, with a difference which grows by at most $2L$ per step. This gives a change of the normalised dressed loop of at most $x^w_K\mathfrak b(\mathcal C)$, with $x^w_K$ summable in $K$; the interaction weights $Q_{K'}(\mathfrak B)$ of the families, cluster activities of the reference gas of the stopping step, are compared in the same way, with a change at most $x^Q_K\mathfrak b(\mathcal C)$. Both are Supplement, Proposition~S.16(g).

\emph{The loop constant.} The constant of the dressed ensemble is $E_{\mathrm{const}} = \sum_{j \leq K'}E_j$, where $E_j \coloneqq E_{\mathrm{const},j} - E_{\mathrm{const},j-1}$ collects, by \eqref{eq:step-identity}, the flat remainders of the loops and the cross term of the step $j$, and at an $\mathbf R$-step also the remainders $E^{\circ,\mathbf R}_j(C_i, T)$ and the cross constant $E^{\mathbf R\times}_j$. We index these terms by the age $m \coloneqq K' + 1 - j$ in run~A and $m \coloneqq K' + 2 - j$ in run~B, so that the two terms of the same shared scale have the same $m$. Each $E_j$ is at most $\epsilon_j(\mathcal C, T)$ of \eqref{eq:step-remainder}, and each of the five parts of $\epsilon_j(\mathcal C, T)$ is at most $A_EL^{m/2}e^{-a_EL^{m/2}}$, with constants $A_E$ and $a_E$ which do not depend on $K$. Indeed, in the first three parts the lengths are those of the stopping step multiplied by $L^m$, as in the proof of Corollary~\ref{cor:remainder-sum}, and $te^{-ct} \leq (2/c)e^{-ct/2}$ for $t \geq 0$. In the two parts of the $\mathbf R$-steps the lengths $r_j(C_i)$, $\iota_{T,j}$, and $a_j$ are at least $a_j/2$, by \eqref{eq:lengths-at-step} and the proof of Lemma~\ref{lem:torus-condition}, and $a_j = a_{K'}L^{m-1} \geq L^{m+2}MR_{K'}$; since $R_{K'} \geq (\log\hat g_{K'}^{-2})^r$, $R_j < L(\log\hat g_j^{-2})^r$ by Lemma~\ref{lem:tube-cubes}(i), and $\log\hat g_j^{-2} = \log(\hat g_{K'}^{-2} + b'(m - 1)) \leq (1 + \log(1 + b'm))\log\hat g_{K'}^{-2}$, we get $a_j/(MR_j) \geq L^{m+1}(1 + \log(1 + b^*m))^{-r} \geq c_LL^{m/2}$, with $c_L > 0$ a constant of $L$, $r$, and $b^*$. Since moreover $MR_je^{-p_0(g_j)/4}$ is bounded uniformly in $j$, and $P_j(C_i) = P_q(C_i)a_j$, the bound $te^{-ct} \leq (2/c)e^{-ct/2}$ with $t = a_j/(MR_j)$ gives the claim for these parts as well. The terms of age $m \geq m^E_K$, the least $m$ with $a_EL^{m/2} \geq 2\log(K + 2)$, are not compared; their sum is at most $\tau^E_K \coloneqq 2A_E\max(L^{1/2}, 2L^{1/2}\log(K + 2)/a_E)(K + 2)^{-2}$ in each run, and for $K \geq K_1$ the extra step of run~B is among them. For the younger terms, which are produced at $\mathbf R$-steps of both runs, we interpolate: every constituent $x$ of the step, an activity of the effective action, a marked activity of the loop, a factor $V(Y)$, or the operation of a young component, is replaced by $x^\varsigma \coloneqq x^A + \varsigma(x^B - x^A)$, $\varsigma \in \mathbb C$, after the old components, those born at a scale at most $K' - \floor{K'/2}$, have been removed by multiplying their operations by a parameter $\mathfrak z$ with $\abs{\mathfrak z} \leq e^{C_PK'/4}$. This is possible because the bound \eqref{eq:component-factor} holds at every step and for every component separately: an old component has lived through at least $K'/4$ steps at a young step, so that its operation has a total mass at most $e^{-C_PK'/4}$ times the bound under which the expansion of the step converges, and the multiplication by $\mathfrak z$ keeps each activity within that bound. The relative differences of the constituents are at most the number $\rho^E_K$ defined in the proof of Supplement, Proposition~S.16(g)--(h), which is at most a constant times $(K + 1)w_K + \ell_KL^{-\gamma_HK'/4}$ plus a multiple of $x^w_K$, so that it tends to zero and is summable in $K$ after multiplication by any power of $\log(K + 2)$; for $\abs{\varsigma} \leq 1/(2\rho^E_K)$ the constituents are at most twice their bounds, the cluster expansion of the step converges, and the per-step term is analytic in $\varsigma$ with the bound $A^{E,2}L^{m/2}e^{-a^{E,2}L^{m/2}}$ at the enlarged amplitudes and the lowered decay; for the cross constant this holds for each of the cluster expansions, of the gas linked to $\mathcal C$ and of the gases linked to the single loops $C_i$, whose constants it subtracts. Cauchy's estimate on the disc of radius $1/(2\rho^E_K) \geq 2$ gives
\begin{equation}\label{eq:loop-constant-difference}
    \abs{E_j(1) - E_j(0)} \leq 4\rho^E_KA^{E,2}L^{m/2}e^{-a^{E,2}L^{m/2}}\;,
\end{equation}
and the old components cost at most $2e^{-C_PK'/4}A_EL^{m/2}e^{-a_EL^{m/2}}$ in each run, by Cauchy's estimate in $\mathfrak z$. Summing over $m$ with $\sum_{m \geq 1}\Lambda^me^{-a\Lambda^m} \leq e^{-a}(1 + 6/a)$ for $\Lambda \geq 2$, we obtain $\abs{E^B_{\mathrm{const}} - E^A_{\mathrm{const}}} \leq \Delta E_K$, with $\Delta E_K \leq \min(2\tau^E_K + \Delta^E_K, 2B_E(\mathcal C))$ summable in $K$. The details are Supplement, Proposition~S.16(h).

\begin{lemma}[The All-Small-Field Pair]\label{lem:trivial-pair}
For $K \geq K_1$, with run~A's constant on both runs, the terms $\ell_0$ change $W_{K+1}(\mathcal C) - W_K(\mathcal C)$ by at most $(3x_K + \lambda_K + x^w_K + x^Q_K)\,B(\mathcal C)$ on their common support, and by at most $b^{(S)}_KB(\mathcal C)$ through the shells of the stopping-scale test functions.
\end{lemma}

\begin{proof}
On the common support, run~B's density is $e^{D(V) + E^B - E^A}$ times run~A's, and $\abs{D(V) + E^B - E^A} \leq X_K$, so that $\int\abs{\varrho^B_{\ell_0} - \varrho^A_{\ell_0}}\,\dd V \leq x_KZ^A_{\ell_0}$, where the number $E^B - E^A$ cancels between the numerator and the normalisation. The four parts above give $\int\abs{\widehat{\mathcal W}^B\varrho^B_{\ell_0} - \widehat{\mathcal W}^A\varrho^A_{\ell_0}}\,\dd V \leq \mathfrak b(\mathcal C)(x_K + X_K + \lambda_K + x^w_K + x^Q_K)Z^A_{\ell_0}$. As in \eqref{eq:removal}, $N'/Z' - N/Z = (\Delta N - W'\Delta Z)/Z$ with $\abs{W'} \leq B(\mathcal C)$ and $Z \geq Z^A_{\ell_0}$, and $X_K \leq x_K$. The shells are Proposition~\ref{prop:stopping-shells}.
\end{proof}

\subsection{The Pairs with Large Fields}\label{sec:large-field-pairs}

For a young matched pair Theorem~\ref{thm:lf-comparison} compares the ingredients ($\beta$) and ($\gamma$); the ingredients ($\alpha$) are compared by \eqref{eq:recursion-solution} per point of a history at a scale $i \geq K' - m$: the frozen exponents and the depth terms by $C_a[(i + 1)\tilde\theta^i + w_K]$; the coupling term of the action by $C\varepsilon_i^2\abs{\delta_i} \leq \abs{\delta_i}$; the function \cite[(1.6), p.~178]{BalabanLFR1}, $z^{-1}\exp[-g_i^{-2}(1 - \operatorname{Re}\tr V_i)]$, relatively by at most $e^{4\abs{\delta_i}} - 1 \leq 10\abs{\delta_i}$, because $0 \leq 1 - \operatorname{Re}\tr V \leq 2$ and the normalisation $z$ has the common cutoff $\varepsilon_i$; the Gaussian form by $Cr_i\abs{A_i(b)}^2$ per variable, by \eqref{eq:gaussian-depth}, which on the cutoff is at most $Cr_ip_0(\hat g_i)^2$ and where a large-field function leaves $A_i$ unbounded is absorbed by the positivity of the form; and the localised normalisation constants, of order $\log g_i^{-2}$ per variable \cite[p.~380]{BalabanLFR2}, whose dependence on the coupling is the term $-\frac12\dim G\log g_i^{-2}$ per variable of the normalisation of a Gaussian integral over a variable scaled by $g_i$, by $\frac12\dim G\abs{\log(1 + g_i^2\delta_i)} \leq \abs{\delta_i}$. Collecting, with $C_M$ the number of ingredients per point,
\begin{equation}\label{eq:s-k}
\begin{aligned}
    s'_K \coloneqq{}& 2C_a\bigl[(K' + 1)\tilde\theta^{K'-m} + w_K\bigr] + 10\sup_{i \geq K'-m}\abs{\delta_i}\\
    &+ 4C(1 + C_{\mathrm{cov}})C_{\delta\Delta}\ell_K\bigl[L^{-\gamma_H(K'-m-1)} + L^{-(1-\varepsilon_0)(K'-m)}\bigr]\;,
\end{aligned}
\end{equation}
$s_K \coloneqq s'_K + L^4C_h\tilde\theta^{K'/2}$, and $\omega^s(X) \coloneqq (1 + C_MP(X))e^{C_Ms_KP(X)}$, which is at most $e^{C_PP(X)}$ when $C_P \geq 2C_M$ and $s_K \leq 1$.

\begin{proposition}[The Young Pairs]\label{prop:young-pairs}
Let $k_*$ be the least integer with $k_* \geq k_{\mathrm{LF}}$ such that $s'_K \leq \frac12$ and $L^4C_h\tilde\theta^{K'/2} \leq \frac12$ for $K' \geq k_*$. For $K' \geq k_*$ the young matched pairs change $W_{K+1}(\mathcal C) - W_K(\mathcal C)$ by at most $\bar\delta_{\mathrm{LF}}(5s_K + 3x_K + s^{W}_K)B(\mathcal C)$, where $s^W_K$, summable in $K$, bounds the change of the dressed loops near the large-field regions (Supplement, Proposition~S.14(f)). The threshold $k_*$ exists, and $s_K$ is summable in $K$.
\end{proposition}

\begin{proof}
By \eqref{eq:three-piece}, \eqref{eq:lf-comparison}, and the bounds on the ingredients ($\alpha$), each of whose contributions enters per point of a history and is counted by $P(X)$, with the small-field exponent on the last domain bounded by $X_K$ as for the all-small-field pair,
\begin{equation}\label{eq:young-pair-bound}
\begin{gathered}
    \int\abs{\varrho^{B'}[\tau^+] - \varrho^A[\tau]}\,\dd V \leq (s_K + x_K)\prod_{X}\omega^s(X)\,M^A_\tau\;,\\
    \int(\Pi_B - \Pi_A)_+\varrho_B\,\dd V \leq s_K\prod_X\omega^s(X)\,M^B_{\tau^+}\;,
\end{gathered}
\end{equation}
where $\varrho^{B'} \coloneqq \min(\Pi_A, \Pi_B)\varrho_B$ and we used $(e^a - 1) + (e^b - 1) \leq e^{a+b} - 1$, $e^{\sum\phi_X} - 1 \leq (\sum\phi_X)e^{\sum\phi_X}$, and $\sum_XC_MP(X) \leq \prod_X(1 + C_MP(X))$. Since $Z_\ell \leq M_\ell$ and the weights $\omega^s(X)$ are admissible in Proposition~\ref{prop:large-field-weight}, by the proof of that proposition the sums over the young pairs are at most $\delta_{\mathrm{LF}}(s_K + x_K)Z^A_{\ell_0}$ and $\delta_{\mathrm{LF}}s_KZ^B_{\ell_0}$. Against the normalised dressed loops the first bound holds with $\mathfrak b(\mathcal C)(2s_K + 2x_K + s^W_K)$ in place of $s_K + x_K$, the deviation insertions at the points of the histories being compared by the recursion and the dressings near the large-field regions by Supplement, Proposition~S.14(f). Removing run~B's own-support pieces changes $W_{K+1}$ by at most $2B(\mathcal C)s_K\delta_{\mathrm{LF}}$, as in \eqref{eq:removal}, and the rest changes $W_{K+1} - W_K$ by at most $B(\mathcal C)\delta_{\mathrm{LF}}(3s_K + 3x_K + s^W_K)$ through the numerators and the normalisations. Finally $s'_K \to 0$: $(K' + 1)\tilde\theta^{K'-m} \to 0$ since $K' - m \geq K'/2$, the covariance terms decay geometrically against the polylogarithm $\ell_K$, and for $i \geq K'/2$ the bound \eqref{eq:recursion-solution} gives $\abs{\delta_i} \leq C[K'\tilde\theta^{K'/2} + Kw_K]$, where $Kw_K$ is summable because $K\mu_K = K^{4-\kappa_0/2}\ell_K$ and $\kappa_0 > 10$.
\end{proof}

\subsection{Proof of the Comparison}\label{sec:cauchy-proof}

Let $K_1 \geq K_0(\mathcal C)$ be the least index such that every $K \geq K_1$ satisfies $K' \geq \max(k_0(\varepsilon_0), k_*, 2m_{\mathrm{rc}} + 2)$, with $m_{\mathrm{rc}}$ the integer of Theorem~\ref{thm:walk-difference}, $\max(m^E_K, m^w_K) \leq \floor{K'/4}$, $\rho^E_K \leq \frac14$, and the further conditions of Supplement, Proposition~S.16, each of which holds for $K$ large because $m^E_K$ and $m^w_K$ grow as $\log\log K$ while $\rho^E_K$, $x^*_K$, and $s_K$ tend to zero. For $K \geq K_1$ put
\begin{equation}\label{eq:cauchy-sequence}
    a_K \coloneqq \Bigl[3x_K + \lambda_K + x^w_K + x^Q_K + b^{(S)}_K + \bar\delta_{\mathrm{LF}}\bigl(a^{(U)}_K + b^{(M)}_K\bigr)\Bigr]B(\mathcal C) + \bigl(e^{\Delta E_K} - 1\bigr)B(\mathcal C)\;,
\end{equation}
with $a^{(U)}_K \coloneqq 4e^{-(C_P-1)K'}$ and $b^{(M)}_K \coloneqq 4e^{-(C_P-1)(\floor{K'/2}+1)} + 5s_K + 3x_K + s^W_K$.

\begin{proof}[Proof of Theorem~\ref{thm:cauchy}]
For $K < K_1$ the bound $\abs{W_{K+1} - W_K} \leq 2B(\mathcal C)$ is Theorem~\ref{thm:uniform-bound}. Let $K \geq K_1$. We first replace run~B's constant $E^B_{\mathrm{const}}$ by run~A's: with $W^{B,A} \coloneqq e^{E^A_{\mathrm{const}} - E^B_{\mathrm{const}}}W_{K+1}(\mathcal C)$,
\begin{equation}\label{eq:constant-swap}
    W_{K+1} - W_K = \bigl(W^{B,A} - W_K\bigr) + \bigl(e^{E^B_{\mathrm{const}} - E^A_{\mathrm{const}}} - 1\bigr)W^{B,A}\;,
\end{equation}
and $\abs{W^{B,A}} \leq B(\mathcal C)$ because the bound of Lemma~\ref{lem:term-bound} holds for run~B's dressed loop with run~A's constant, so that the second term is at most $(e^{\Delta E_K} - 1)B(\mathcal C)$. For the first, the terms of kind (U) of the two runs change the difference by at most $\bar\delta_{\mathrm{LF}}a^{(U)}_KB(\mathcal C)$ and the old pairs change the difference by at most $4e^{-(C_P-1)(m+1)}\bar\delta_{\mathrm{LF}}B(\mathcal C)$, by Corollary~\ref{cor:unmatched-old}; the young pairs change it by at most $\bar\delta_{\mathrm{LF}}(5s_K + 3x_K + s^W_K)B(\mathcal C)$, by Proposition~\ref{prop:young-pairs}; and the all-small-field pair by at most $(3x_K + \lambda_K + x^w_K + x^Q_K + b^{(S)}_K)B(\mathcal C)$, by Lemma~\ref{lem:trivial-pair}. This gives $\abs{W_{K+1} - W_K} \leq a_K$. Each term of \eqref{eq:cauchy-sequence} is summable in $K$: $x_K \leq X_Ke^{X_K}$ with $X_K$ summable and bounded; $\lambda_K$, $a^{(U)}_K$, and the first term of $b^{(M)}_K$ decay geometrically; $b^{(S)}_K$ by Proposition~\ref{prop:stopping-shells}; $s_K$ and $s^W_K$ by Proposition~\ref{prop:young-pairs}; $x^w_K$, $x^Q_K$, and $\Delta E_K$ by the bounds above; and $e^{\Delta E_K} - 1 \leq \Delta E_Ke^{2B_E(\mathcal C)}$. Since $B(\mathcal C)$ and $\bar\delta_{\mathrm{LF}}$ do not depend on $K$, $\sum_Ka_K < \infty$.
\end{proof}

\begin{proof}[Proof of Theorem~\ref{thm:main}]
\emph{The thresholds.} We sort the requirements on the coupling of Sections~\ref{sec:construction}--\ref{sec:reflection-positivity} into the two kinds of Section~\ref{sec:results}. Those of the first kind are finitely many upper bounds on $\bar g$ in terms of the constants of the construction and of $\Lambda_*$, among them $\bar g \leq \gamma'$ of Theorem~\ref{thm:coupling-flow}, and we let $\gamma_0 > 0$ be the smallest of these bounds. Each requirement of the second kind, for an ensemble with $c_{\mathcal C} = c$, bounds a nondecreasing function of $\gamma_q$ from above by a positive quantity which does not increase with $c$, because the constants which it contains are nondecreasing functions of $c_{\mathcal C}$ (Definition~\ref{def:dressed-loop}(v)); it therefore holds whenever $\gamma_q$ lies below a positive number which does not increase with $c$. We let $\gamma_\sharp(c) > 0$ be the smallest of these numbers, a nonincreasing function of $c$ which depends only on the constants of the construction and on $\Lambda_*$, and we put
\begin{equation}\label{eq:level-threshold}
    q_\sharp(c, \bar g) \coloneqq \max\Bigl(0,\ \Bigl\lceil\frac{\gamma_\sharp(c)^{-2} - \bar g^{-2} + C_W}{b}\Bigr\rceil\Bigr)\;,
\end{equation}
with the constants $b$ and $C_W$ of Theorem~\ref{thm:coupling-flow}. This function is nondecreasing in $c$, because $\gamma_\sharp$ does not increase, and nondecreasing in $\bar g$; it vanishes as soon as $\bar g^{-2} \geq \gamma_\sharp(c)^{-2} + C_W$, and it depends only on the constants of the construction and on $\Lambda_*$. If $q \geq q_\sharp(c, \bar g)$, then $bq \geq \gamma_\sharp(c)^{-2} - \bar g^{-2} + C_W$, and the first inequality of \eqref{eq:stopping-coupling-bound} gives
\begin{equation}\label{eq:level-threshold-consequence}
    \gamma_q^{-2} \geq \bar g^{-2} + bq - C_W \geq \gamma_\sharp(c)^{-2}\;,
\end{equation}
that is, $\gamma_q \leq \gamma_\sharp(c)$.

\emph{The bound and the limit.} Let $\bar g \in (0, \gamma_0]$, and let $\mathcal C$ be a loop ensemble with admissible loops at a level $q \geq q_\sharp(c_{\mathcal C}, \bar g)$. Then $\bar g$ satisfies the requirements of the first kind, and $\gamma_q$ satisfies those of the second kind for $\mathcal C$ by \eqref{eq:level-threshold-consequence}, so that $\bar g$ and $q$ are as in Theorem~\ref{thm:main} for $\mathcal C$ in the sense of Section~\ref{sec:results}. The bound of \eqref{eq:main-bound} is therefore Theorem~\ref{thm:uniform-bound} with $K_0 \coloneqq K_0(\mathcal C)$ and $B \coloneqq B(\mathcal C)$, and the summability of the differences is Theorem~\ref{thm:cauchy}. A sequence with summable differences is a Cauchy sequence, which gives the limit \eqref{eq:main-limit}. The dependence of $K_0$ and $B$ is read from \eqref{eq:index-k-zero}, Theorem~\ref{thm:first-step}, and Theorem~\ref{thm:uniform-bound}.
\end{proof}

\subsection{Several Loops}\label{sec:several-loops-cauchy}

The proof has been written for an ensemble of $n$ loops throughout, and it remains to say where the number of loops enters. Every comparison above is made per traversal of a bond of the stopping step, per insertion, per cube, or per point of a history, and a cube of side less than the separation $a_k$ of Lemma~\ref{lem:stopping-index} meets one loop. The number of loops therefore enters only through sums over the loops: the leading words through $\lambda_K$, a sum over the loops, because every factor of $\prod_if^{(i)}$ has modulus at most one; the transports and the insertions through $\mathfrak b(\mathcal C)$, the exponential of a sum over the loops; the wraps through the total perimeter $P_{K'}(\mathcal C)$ in their constants $A_w$; and the loop constant through the sum over the loops of the single-loop constants and the constants of the cross groups, whose per-step terms are bounded as those of one loop with the total perimeter in place of the perimeter (Corollary~\ref{cor:remainder-sum}). The insertions of the cross species are not contracted, so that the recursion of the deviation insertions does not apply to them, and they are compared as the wraps are. Those born at the step $K' - m$ have total mass at most $2^{m+1}\mathsf A^\times_{K'-m}$ in each run, by \eqref{eq:cross-bound}, which decays doubly exponentially in $m$; those of age at least $m^\times_K$, an index which grows like $\log\log K$ and is not the age $m_K = \floor{K^{1/2}}$ of the old terms of the effective action, are not compared, and their total $\tau^\times_K$ in each run is summable in $K$. The younger ones are produced by the fluctuation integrals and the $\mathbf R$-operations of the last $m^\times_K$ steps, are compared at their birth by the interpolation of the constituents of their step, and are then carried. Their differences cannot be bounded separately from those of the wraps: a difference of the cross species can pass into a wrap, and the wraps inside the groups of the reference gas change the cluster occurrences which every new group merges, the deviation groups among them. The differences of the young insertions of the two species, of the activities of the reference gas which contain wraps, and of the part of the deviation insertions which comes from these occurrences therefore satisfy one linear recursion over the last $m^\times_K$ steps, with a factor per step which does not depend on $K$ and with sources which are summable in $K$; its solution is at most a power of $\log K$ times these sources. Together with the weights this changes the normalised dressed loop by at most $x^\times_K\mathfrak b(\mathcal C)$ with $x^\times_K$ summable in $K$, and for several loops $x^\times_K$ takes the place of $x^w_K + x^Q_K$ (Supplement, Proposition~S.16(i)). The counts of the positions of a cross group are those of one loop with the total perimeter, because two loops are at distance at least $a_k$. No condition on the coupling or on $L$ involves $n$; the index $K_1$ involves $\mathcal C$ through $B(\mathcal C)$, the total perimeter, and the thresholds of the single loops.

\section{Symmetries, Reflection Positivity, and Nonvanishing}\label{sec:reflection-positivity}

The symmetries and the reflection positivity are properties of the lattice measures at every cutoff, and one expects them to pass to the limits without effort, since a limit of equal numbers is an equality and a limit of positive semidefinite matrices is positive semidefinite. The difficulty lies in the normalisation. The renormalised expectations are the expectations divided by the products $\prod_i\kappa^{(q)}_K(C_i, \chi_i)$, so that a symmetry $h$ of the lattice measures survives the division only if $\kappa^{(q)}_K(hC, \chi) = \kappa^{(q)}_K(C, \chi)$ for every loop; and the positivity of a matrix $(\mathbb E_K[F_a\,\Theta F_b])_{a,b}$ survives it only if the number by which the entry $(a, b)$ is divided factorises as a number of $a$ times the complex conjugate of the same number of $b$, which for the ensembles $\mathcal D_a \cup \theta\bar{\mathcal D}_b$ means that the normalisation of a reflected loop with the conjugate character must be the complex conjugate of the normalisation of the loop. The coefficients of the normalisation are defined by cluster expansions in which these symmetries are not visible, so we obtain them from the exact values of the flat parts, whose symmetries follow from those of the construction, and from the extraction of the coefficients along dilations of the loop. The construction is symmetric under the reflections in coordinate hyperplanes, but it is not known to be symmetric under the exchange of two axes, and this is where our covariance statements stop. The nonvanishing is then a consequence of the identity at the stopping scale: at small coupling the measure $\mu_K$ has almost full mass, the background is almost flat on its support, and all the corrections tend to zero with $\bar g$.

\subsection{Symmetries and Reflection Positivity on the Lattice}\label{sec:lattice-reflection-positivity}

A symmetry $h$ of level $q$ (Section~\ref{sec:results}) acts on the bonds of $\mathbb T_K$, $K \geq q$, by $h(x, y) \coloneqq (hx, hy)$, on paths bond by bond, and on gauge fields by $(hU)(b) \coloneqq U(hb)$, so that $(hU)(\gamma) = U(h\gamma)$ for every path $\gamma$. For a function $F$ of the gauge field we put $(F \circ h)(U) \coloneqq F(hU)$; in particular $W_\chi(C) \circ h = W_\chi(hC)$ for a loop $C$ at level $q$, and $W_{\mathcal C} \circ h = W_{h\mathcal C}$ for an ensemble.

\begin{lemma}[Invariance of the Lattice Measures]\label{lem:lattice-invariance}
Let $K \geq q$ and let $h$ be a symmetry of level $q$. For every bounded measurable function $F$ of the gauge field on $\mathbb T_K$, $\mathbb E_K[F \circ h] = \mathbb E_K[F]$ and $\mathbb E_K[\bar F] = \overline{\mathbb E_K[F]}$.
\end{lemma}

\begin{proof}
We write $s_U(p) \coloneqq 1 - \operatorname{Re}\tr U(\partial p)$. The holonomy $(hU)(\partial p) = U(h\partial p)$ is that of $U$ around the plaquette $hp$, from one of its corners and in one of its two orientations, and $\operatorname{Re}\tr$ is invariant under conjugation and, on unitary matrices, under inversion; hence $s_{hU}(p) = s_U(hp)$. Since $h$ permutes the plaquettes, the action of \eqref{eq:lattice-measure} is invariant under $U \mapsto hU$. The map $h$ permutes the unoriented bonds, and the value of $hU$ on a positive bond $b$ is the value of $U$ on $hb$ if $hb$ is positive and the inverse of its value on $(hb)^{-1}$ otherwise; since the Haar measure is invariant under inversion, the bijection $U \mapsto hU$ preserves the product Haar measure. The change of variables $U \mapsto hU$ in the integral defining $\mathbb E_K[F \circ h]$ therefore gives the first assertion, and the second holds because $\mu_{K,g_0}$ is a positive measure.
\end{proof}

For the reflection positivity we use the reflections of level $q$, for which the positivity of the Wilson action was observed by Osterwalder and Seiler \cite{OsterwalderSeiler}; we give the short proof, since it is the basis of everything in this section. Let $\mu \in \{ 0, \dots, 3 \}$ and $a \in uL^{-q}\mathbb Z$. Since $a$ and $T/2 = L^{m_T}u$ are multiples of $\varepsilon_K$ for $K \geq q$, the reflection $\theta_{\mu,a}$ is an involution of the sites of $\mathbb T_K$, which fixes the two hyperplanes $P_{\mu,a} \coloneqq \{ x_\mu \equiv a \} \cup \{ x_\mu \equiv a + T/2 \}$ pointwise and exchanges the open half-tori $H_{\mu,a} \coloneqq \{ a < x_\mu < a + T/2 \}$ and $\theta_{\mu,a}H_{\mu,a}$. For a bounded measurable function $F$ of the gauge field we put $(\Theta_{\mu,a}F)(U) \coloneqq \overline{F(\theta_{\mu,a}U)}$; the map $\Theta_{\mu,a}$ is antilinear and multiplicative.

Let $\bar\Lambda_+$ be the set of the sites with $x_\mu - a \in [0, T/2]$ modulo $T$, and $\bar\Lambda_- \coloneqq \theta_{\mu,a}\bar\Lambda_+$, so that $\bar\Lambda_+ \cap \bar\Lambda_- = P_{\mu,a}$. Let $E_0$ be the set of the bonds with both ends in $P_{\mu,a}$, $E_+$ the set of the other bonds with both ends in $\bar\Lambda_+$, and $E_- \coloneqq \theta_{\mu,a}E_+$. Every bond belongs to exactly one of these sets: a bond orthogonal to $e_\mu$ has both ends on one hyperplane $\{ x_\mu = \mathrm{const} \}$, and a bond along $e_\mu$ joins two neighbouring hyperplanes, which lie both in $\bar\Lambda_+$ or both in $\bar\Lambda_-$, because the two hyperplanes of $P_{\mu,a}$ are hyperplanes of $\mathbb T_K$ at distance $T/2 \geq 2\varepsilon_K$. The bonds of $E_0$ are orthogonal to $e_\mu$, so $\theta_{\mu,a}$ fixes each of them with its orientation. We write $\mathcal A_+$ for the set of the bounded measurable functions of the restriction of $U$ to $E_+ \cup E_0$; for $F \in \mathcal A_+$ the function $\Theta_{\mu,a}F$ depends on the restriction to $E_- \cup E_0$.

\begin{lemma}[Reflection Positivity on the Lattice]\label{lem:lattice-reflection-positivity}
Let $K \geq q$, $\mu \in \{ 0, \dots, 3 \}$, and $a \in uL^{-q}\mathbb Z$. For every bounded measurable $F$, $\mathbb E_K[\Theta_{\mu,a}F] = \overline{\mathbb E_K[F]}$, and for every $F \in \mathcal A_+$, $\mathbb E_K[F\,\Theta_{\mu,a}F] \geq 0$.
\end{lemma}

\begin{proof}
We write $\theta$ for $\theta_{\mu,a}$, which is a symmetry of level $q$, and $s_U(p) \coloneqq 1 - \operatorname{Re}\tr U(\partial p)$ as in the proof of Lemma~\ref{lem:lattice-invariance}. That lemma, applied to $\bar F$, gives $\mathbb E_K[\Theta F] = \mathbb E_K[\bar F \circ \theta] = \mathbb E_K[\bar F] = \overline{\mathbb E_K[F]}$, which is the first assertion.

For the second, a plaquette lies in $\bar\Lambda_+$ or in $\bar\Lambda_-$, by the argument given for the bonds, and in both exactly when it lies in $P_{\mu,a}$. Put
\begin{equation}\label{eq:half-action}
    B(U) \coloneqq \sum_{p \subset \bar\Lambda_+,\ p \not\subset P_{\mu,a}}s_U(p) + \frac12\sum_{p \subset P_{\mu,a}}s_U(p)\;,
\end{equation}
a real function in $\mathcal A_+$. The reflection maps the plaquettes of $\bar\Lambda_+$ which do not lie in $P_{\mu,a}$ bijectively onto those of $\bar\Lambda_-$, and each plaquette of $P_{\mu,a}$ onto itself, so that, since $s_{\theta U}(p) = s_U(\theta p)$, $B(\theta U)$ is the sum \eqref{eq:half-action} with $\bar\Lambda_-$ in place of $\bar\Lambda_+$, and the action is $B + B \circ \theta$. With $f \coloneqq Fe^{-g_0^{-2}B} \in \mathcal A_+$ and since $B$ is real, the integrand of $\mathbb E_K[F\,\Theta F]$ against the Haar measure is therefore
\begin{equation}\label{eq:reflection-integrand}
    F(U)\overline{F(\theta U)}\,e^{-g_0^{-2}(B(U) + B(\theta U))} = f(U)\,\overline{f(\theta U)}\;.
\end{equation}
Write $U = (U_+, U_0, U_-)$ for the restrictions to $E_+$, $E_0$, and $E_-$. The restriction of $\theta U$ to $E_+ \cup E_0$ is $(\varphi(U_-), U_0)$, where $\varphi(U_-)(b) \coloneqq U(\theta b)$ for $b \in E_+$ defines a map $\varphi \colon G^{E_-} \to G^{E_+}$ which maps the product Haar measure to the product Haar measure, as in the proof of Lemma~\ref{lem:lattice-invariance}. By Fubini's theorem, the integral of \eqref{eq:reflection-integrand} is
\begin{equation}\label{eq:reflection-square}
    \int\Bigl(\int f(U_+, U_0)\,\dd U_+\Bigr)\overline{\Bigl(\int f(\varphi(U_-), U_0)\,\dd U_-\Bigr)}\,\dd U_0 = \int\Bigl\lvert\int f(U_+, U_0)\,\dd U_+\Bigr\rvert^2\dd U_0 \geq 0\;,
\end{equation}
and $\mathbb E_K[F\,\Theta F]$ is this number divided by $Z_{K,g_0} > 0$.
\end{proof}

The lemma makes $\langle F, G\rangle \coloneqq \mathbb E_K[F\,\Theta_{\mu,a}G]$ a positive semidefinite sesquilinear form on $\mathcal A_+$, linear in $F$ and antilinear in $G$; such a form is Hermitian, by polarisation. If $F \in \mathcal A_+$ has the mean $c \coloneqq \mathbb E_K[F]$, then the constant $c$ lies in $\mathcal A_+$, $\mathbb E_K[\Theta_{\mu,a}F] = \bar c$ by the lemma, and
\begin{equation}\label{eq:reflection-covariance}
    0 \leq \langle F - c, F - c\rangle = \mathbb E_K[F\,\Theta_{\mu,a}F] - \abs{\mathbb E_K[F]}^2\;.
\end{equation}

The Wilson loops behave as follows. Let $C$ be a loop whose sites lie in $H_{\mu,a}$, with the character $\chi$ of the representation $\pi$. Its bonds have both ends in $H_{\mu,a}$, hence lie in $E_+$, and $W_\chi(C) \in \mathcal A_+$. Since $\pi$ is unitary, $\overline{\chi(g)} = \bar\chi(g)$ for the character $\bar\chi$ of the conjugate representation, which is irreducible and unitary, and therefore
\begin{equation}\label{eq:reflected-loop}
    (\Theta_{\mu,a}W_\chi(C))(U) = \overline{\chi\bigl(U(\theta_{\mu,a}C)\bigr)} = W_{\bar\chi}(\theta_{\mu,a}C; U)\;,
\end{equation}
where $\theta_{\mu,a}C$ is the image of $C$ with the image orientation, a loop at level $q$ with the same chord--arc constant, ratio $P/r$, and number $P_q$ as $C$ (Section~\ref{sec:results}). For an ensemble $\mathcal D$ all of whose loops lie in $H_{\mu,a}$, the ensemble $\theta_{\mu,a}\bar{\mathcal D}$ consists of the images with the conjugate characters, and by multiplicativity $\Theta_{\mu,a}W_{\mathcal D} = W_{\theta_{\mu,a}\bar{\mathcal D}}$.

\subsection{The Covariance of the Normalisation}\label{sec:reflection-normalisation}

The division by the normalisations keeps the symmetries of the lattice measures only where the normalisations share them, and the following proposition says where they do.

\begin{proposition}[Covariance of the Normalisation]\label{prop:reflection-normalisation}
Let $L$ and $\Lambda_*$ be as in Theorem~\ref{thm:main}, let $C$ be an admissible loop at level $q$ with the character $\chi$ of a representation $\pi$, and let $\bar g$ and $q$ be as in Theorem~\ref{thm:main} for $((C, \chi))$. For every $K > j_q$ and every $g \in \mathfrak G_q$,
\begin{equation}\label{eq:reflection-normalisation}
    \kappa^{(q)}_K(gC, \chi) = \kappa^{(q)}_K(C, \chi)\;,
    \qquad
    \kappa^{(q)}_K(C, \bar\chi) = \overline{\kappa^{(q)}_K(C, \chi)}\;,
    \qquad
    \kappa^{(q)}_K(C^{-1}, \chi) = \kappa^{(q)}_K(C, \bar\chi)\;.
\end{equation}
In particular $\kappa^{(q)}_K(\theta_{\mu,a}C, \bar\chi) = \overline{\kappa^{(q)}_K(C, \chi)}$ for every $\mu$ and every $a$ with $2a \in uL^{-q}\mathbb Z$.
\end{proposition}

For the translations in $\mathfrak G_q$ the first identity can be read off from Proposition~\ref{prop:normalisation-structure}, since a translate of $C$ has the same numbers $N^{(0)}_v$ and the same pairs $(v_c, w_c)$ at its corners as $C$. The structure of the normalisation does not help with the reflections $\theta_\mu$, nor with the other two identities: the loops $\theta_\mu C$ and $C^{-1}$ have different numbers of segments of a given orientation and different pairs at their corners, and the coefficients of $\chi$ and of $\bar\chi$ are different numbers, while nothing relates the coefficients of different orientations, or of conjugate characters, to each other. The coefficients of Definition~\ref{def:step-factor} are sums of cluster weights, and \eqref{eq:reflection-normalisation} would follow from symmetries of the individual weights; but the expansions which define the weights involve many choices whose behaviour under the reflections we do not wish to follow. We use instead the exact values of the flat parts, which do not depend on any choice. For a torus $T' \in \mathfrak T$, a loop $C$ at level $q$ on $T'$, and $0 \leq j \leq K'$, let $\widehat{\mathcal W}_j(C, \chi; T')$ be the value at the reference data of the dressed loop of the step $j$ of the single-loop ensemble $((C, \chi))$ on the trivial term, with $\widehat{\mathcal W}_0 \coloneqq 1$, whenever this ensemble is dressed at the step $j$.

\begin{lemma}[Flat Values under the Symmetries]\label{lem:flat-values}
Let $T' \in \mathfrak T$, let $C$ be a loop at level $q$ on $T'$ with the character $\chi$, let $1 \leq j \leq K'$, and let $\mu \in \{ 0, \dots, 3 \}$. If the ensemble $((C, \chi))$ is dressed at the steps $1, \dots, j$, then
\begin{equation}\label{eq:flat-value}
    \widehat{\mathcal W}_j(C, \chi; T') = \kappa^{(j)}(C, \chi)\,e^{E_{\mathrm{const},j}(C, \chi; T')}\;,
\end{equation}
where $E_{\mathrm{const},j}(C, \chi; T')$ is the constant of the scalar exponent of the step $j$ on the trivial term, and
\begin{equation}\label{eq:flat-value-symmetries}
    \begin{gathered}
        \widehat{\mathcal W}_j(\theta_\mu C, \chi; T') = \widehat{\mathcal W}_j(C, \chi; T')\;,
        \qquad
        \widehat{\mathcal W}_j(C, \bar\chi; T') = \overline{\widehat{\mathcal W}_j(C, \chi; T')}\;,
        \\
        \widehat{\mathcal W}_j(C^{-1}, \chi; T') = \widehat{\mathcal W}_j(C, \bar\chi; T')\;,
    \end{gathered}
\end{equation}
each identity holding whenever the single-loop ensembles which occur in it are dressed at the steps $1, \dots, j$.
\end{lemma}

\begin{proof}
At the reference data the leading word is $\chi(1) = 1$, and the insertions of the nonempty families vanish: after the exact factorisation \eqref{eq:exact-factorisation}, and after its counterpart at the first step in the proof of Theorem~\ref{thm:first-step}, the insertions of the families are the parts $Y^{\mathrm d}_{\mathcal G}$, which vanish at the reference data, because for a single loop every insertion takes a scalar value there, by Schur's lemma. On the trivial term the scalar exponent is the number $E_{\mathrm{const},j}$, as in the proof of Theorem~\ref{thm:stopping-identity}. This gives \eqref{eq:flat-value}.

For \eqref{eq:flat-value-symmetries} we express the flat values through the construction itself. The reference data of the trivial term are those at the block field $V_j = 1$, which lies in the support of all its test functions, and on the trivial term the operation $\mathbf T_j$ holds no integrations. Let $N_j[F]$ be the density obtained from $F\varrho_0$ by the operations of the first $j$ steps on $T'$, and $N_j[F](1; \ell_0)$ the value of its term $\ell_0$ at $V_j = 1$. Since the operations which produce the terms of the numerator are those which produce the terms of the density, applied to a function which carries one more factor, Definition~\ref{def:dressed-loop} gives
\begin{equation}\label{eq:flat-value-ratio}
    \widehat{\mathcal W}_j(C, \chi; T') = \frac{N_j[W_\chi(C)](1; \ell_0)}{N_j[1](1; \ell_0)}\;,
\end{equation}
where $N_j[1](1; \ell_0) = \exp A_j(1) > 0$. The right side depends on the loop and its character only through the function $W_\chi(C)$. The map $F \mapsto N_j[F](1; \ell_0)$ has two properties, which are proved in Supplement, Proposition~S.18. It is \emph{real}: $N_j[\bar F](1; \ell_0) = \overline{N_j[F](1; \ell_0)}$, because at real fields every operation of the construction is a sum, a multiplication by a real function, or an integral against a real kernel (Haar and Gaussian measures, $\delta$-functions of the block averaging, test functions, and exponentials of real actions \cite[(3.27), p.~271]{BalabanCRE}, \cite[p.~380]{BalabanLFR2}). And it is \emph{invariant under each $\theta_\mu$}: $N_j[F \circ \theta_\mu](1; \ell_0) = N_j[F](1; \ell_0)$, because $\theta_\mu$ maps every lattice and every partition of the construction to itself, the partitions being centred at the origin (Section~\ref{sec:reference-gas}), maps the trivial term and the trivial field to themselves, and commutes with every operation: with the averages, which Ba{\l}aban defines so as to be ``symmetric with respect to lattice Euclidean transformations'' \cite[p.~252]{BalabanRG1}, with the gauge fixing, introduced ``to preserve the Euclidean invariance'' \cite[p.~247]{BalabanCRE}, with the fluctuation integral \cite[(2.17), (2.18), p.~269]{BalabanRG1}, whose fluctuation fields are written first in an axial gauge \cite[(1.19), p.~79]{BalabanREG} with contours which start at the centre of a block and change the coordinates in a fixed order \cite[(1.7), p.~18]{BalabanPRT1}, an order which the reflection of one coordinate about the centre does not change, with the soft test functions, which are functions of plaquette deviations in norms invariant under conjugation and inversion, and with the $\mathbf R$-operation, whose gauge-fixing graphs are averaged over ``$2^d$ reflections in subsets of coordinates'' \cite[(1.100), p.~201]{BalabanLFR1}, \cite[p.~379]{BalabanLFR2}.

The three identities now follow from \eqref{eq:flat-value-ratio}. Since $W_\chi(C) \circ \theta_\mu = W_\chi(\theta_\mu C)$, the invariance gives $N_j[W_\chi(\theta_\mu C)](1; \ell_0) = N_j[W_\chi(C)](1; \ell_0)$. Since $\pi$ is unitary, $\overline{\chi(V)} = \chi(V^{-1}) = \bar\chi(V)$ for $V \in G$, so that $W_{\bar\chi}(C) = \overline{W_\chi(C)}$, and the reality gives $N_j[W_{\bar\chi}(C)](1; \ell_0) = \overline{N_j[W_\chi(C)](1; \ell_0)}$. Finally $U(C^{-1}) = U(C)^{-1}$, so that $W_\chi(C^{-1}) = \chi(U(C)^{-1}) = \bar\chi(U(C)) = W_{\bar\chi}(C)$ is the same function as $W_{\bar\chi}(C)$. Dividing by the positive number $N_j[1](1; \ell_0)$ gives \eqref{eq:flat-value-symmetries}.
\end{proof}

\begin{proof}[Proof of Proposition~\ref{prop:reflection-normalisation}]
We first reduce the first identity to the reflections $\theta_\mu$. Every $g \in \mathfrak G_q$ is a composition $g = g_1 \cdots g_r$ of translations by vectors of $uL^{-q}\mathbb Z^4$ and reflections $\theta_\mu$. Each loop $g_{i+1} \cdots g_rC$ is an admissible loop at level $q$ with the character $\chi$ (Section~\ref{sec:results}), so that $\bar g$ and $q$ are as in Theorem~\ref{thm:main} for it, the conditions of that theorem involving the loop only through $c_\pi$ and $\Lambda_*$; and the normalisation of a loop is unchanged by a translation in $\mathfrak G_q$, as noted after the proposition. It therefore suffices to prove the first identity for $g = \theta_\mu$, and the last assertion follows from the first two identities, since $\theta_{\mu,a} \in \mathfrak G_q$ when $2a \in uL^{-q}\mathbb Z$.

The three remaining identities have a common form. Let $(C', \chi', \chi^\circ, \epsilon)$ be one of the three quadruples
\begin{equation}\label{eq:symmetry-cases}
    (\theta_\mu C, \chi, \chi, \mathrm{id})\;,
    \qquad
    (C, \bar\chi, \chi, \mathrm{conj})\;,
    \qquad
    (C^{-1}, \chi, \bar\chi, \mathrm{id})\;,
\end{equation}
where $\mathrm{id}$ is the identity of $\mathbb C$ and $\mathrm{conj}$ the complex conjugation, so that the identity to be proved reads $\kappa^{(q)}_K(C', \chi') = \epsilon(\kappa^{(q)}_K(C, \chi^\circ))$. The pairs $(C, \chi^\circ)$ and $(C', \chi')$ satisfy the hypotheses of the proposition, because $c_{\bar\pi} = c_\pi$ and $\Lambda(C') = \Lambda(C)$. It suffices to prove $\kappa_j(C', \chi') = \epsilon(\kappa_j(C, \chi^\circ))$ for every step $1 \leq j \leq K'$, and since $\kappa_j$ is the exponential of \eqref{eq:step-factor} and $\epsilon$ commutes with the exponential, it suffices to prove that
\begin{equation}\label{eq:log-defect}
    D_j(C) \coloneqq \log\kappa_j(C', \chi') - \epsilon\bigl(\log\kappa_j(C, \chi^\circ)\bigr) \in 2\pi i\mathbb Z\;.
\end{equation}
Fix $j$. Let $c_0$ be a corner of $C$, and for an integer $m \geq 1$ let $C_m$ be the image of $C$ under the dilation $x \mapsto c_0 + m(x - c_0)$, read on a torus $T'_m \in \mathfrak T$ whose period we choose below. The path $C_m$ visits no site twice, its corners lie in $c_0 + uL^{-q}\mathbb Z^4$, its segments have the orientations of those of $C$ and $m$ times their lengths, and its corners have the same pairs $(v_c, w_c)$ as those of $C$. Hence $C_m$ is a loop at level $q$ on $T'_m$ as soon as $3\diam(C_m) < T'_m$, with $\Lambda(C_m) = \Lambda(C)$, $P(C_m)/r(C_m) = P(C)/r(C)$, $P_{j-1}(C_m) = mP_{j-1}(C)$, and $r_{j-1}(C_m) = mr_{j-1}(C)$. We choose the period of $T'_m$ so large that, in addition, $K \geq K_{\mathrm{tor}}$ for $C_m$ on $T'_m$, which is possible because the torus conditions of the first step read $(P_0(C_m)/L)e^{-c\iota_{T'_m}} \leq 1$ in the proof of Theorem~\ref{thm:first-step}, and that $mP_{j-1}(C)e^{-c_2\iota_{T'_m,j}} \leq 1/m$, with the constant $c_2 > 0$ of \eqref{eq:dilation-remainder} below. Let $C'_m$ be obtained from $C_m$ as $C'$ is obtained from $C$, that is, $C'_m$ is $\theta_\mu C_m$, $C_m$, or $C_m^{-1}$; it is the dilation of $C'$ by the factor $m$ about $\theta_\mu c_0$, $c_0$, or $c_0$, since $\theta_\mu$ is linear. Then the ensembles $((C_m, \chi^\circ))$ and $((C'_m, \chi'))$ are dressed at every step up to $K'$, by Theorems~\ref{thm:first-step} and~\ref{thm:inductive-step}, applied on $T'_m$; their hypotheses hold because $\Lambda(C'_m) = \Lambda(C_m) = \Lambda(C)$, because $C'_m$ has the same numbers $P_i$ and $r_i$ as $C_m$ at every step $i$, so that the torus conditions hold for both, and because $c_{\bar\pi} = c_\pi$, so that $\bar g$ and $q$ are as in Theorem~\ref{thm:main} for both ensembles.

By \eqref{eq:step-factor}, $\log\kappa_j(C_m, \chi^\circ)$ is an affine function of $m$, and so is $\log\kappa_j(C'_m, \chi')$, because $C'_m$ is a dilation of $C'$ by the factor $m$. Hence
\begin{equation}\label{eq:affine-defect}
    D_j(C_m) = mA + B
\end{equation}
with numbers $A$ and $B$ which do not depend on $m$, and $D_j(C) = A + B$. On the other hand, by \eqref{eq:flat-value} at the steps $j$ and $j - 1$, for a single loop on the trivial term,
\begin{equation}\label{eq:flat-step-ratio}
    \frac{\widehat{\mathcal W}_j(C_m, \chi^\circ; T'_m)}{\widehat{\mathcal W}_{j-1}(C_m, \chi^\circ; T'_m)} = \exp\bigl(\log\kappa_j(C_m, \chi^\circ) + E_j(C_m, \chi^\circ; T'_m)\bigr)\;,
\end{equation}
where $E_j \coloneqq E_{\mathrm{const},j} - E_{\mathrm{const},j-1}$ is the remainder of the step, which for a single loop is $E^\circ_j + \mathbf 1_{\mathbf R}(j)E^{\circ,\mathbf R}_j$ by \eqref{eq:step-identity} at $j \geq 2$, and the remainder bounded in \eqref{eq:first-step-extraction} and by the periodic comparison of \cite[Appendix~F]{WilsonFirstStep} at $j = 1$. The same identity holds for $(C'_m, \chi')$, and by \eqref{eq:flat-value-symmetries} at the steps $j$ and $j - 1$ the left side for $(C'_m, \chi')$ is the image under $\epsilon$ of the left side for $(C_m, \chi^\circ)$. Since $\epsilon$ commutes with the exponential, taking exponentials apart gives
\begin{equation}\label{eq:defect-modulo}
    D_j(C_m) \in 2\pi i\mathbb Z + \epsilon\bigl(E_j(C_m, \chi^\circ; T'_m)\bigr) - E_j(C'_m, \chi'; T'_m)\;.
\end{equation}
The remainders tend to zero as $m \to \infty$. Indeed, by \eqref{eq:flat-remainder} and \eqref{eq:r-flat-remainder} at $j \geq 2$, and by \eqref{eq:first-step-extraction} and the periodic comparison of \cite[Appendix~F]{WilsonFirstStep} at $j = 1$, and since $P_j \leq P_{j-1}$ and $r_{j-1} \geq r_j$, each of them is at most a constant times
\begin{equation}\label{eq:dilation-remainder}
    \bigl(mP_{j-1}(C) + \abs{\mathcal K(C)}\bigr)e^{-c_1mr_j(C)} + mP_{j-1}(C)\,e^{-c_2\iota_{T'_m,j}}\;,
\end{equation}
with positive constants $c_1$ and $c_2$ which depend on $j$ but not on $m$ (at $j \geq 2$ the smaller of $2c^T_E$ and $c^R_E/(MR_j)$, and of $c^Z$ and $c^R_E/(MR_j)$), and with the constants of $\chi$ and $\bar\chi$; the first term tends to zero because $r_j(C) > 0$, and the second is at most $1/m$ by the choice of $T'_m$. Let $\delta_m$ be the modulus of the difference of the two remainders in \eqref{eq:defect-modulo}, so that $\delta_m \to 0$. By \eqref{eq:affine-defect}, $A = D_j(C_{m+1}) - D_j(C_m)$ lies within the distance $\delta_m + \delta_{m+1}$ of $2\pi i\mathbb Z$ for every $m$, and since $2\pi i\mathbb Z$ is closed, $A \in 2\pi i\mathbb Z$. Then $B = D_j(C_m) - mA$ lies within the distance $\delta_m$ of $2\pi i\mathbb Z$ for every $m$, and $B \in 2\pi i\mathbb Z$. Hence $D_j(C) = A + B \in 2\pi i\mathbb Z$, which is \eqref{eq:log-defect}.
\end{proof}

The argument determines the logarithms of the normalisations only modulo $2\pi i$, which is all that \eqref{eq:reflection-normalisation} asks for; it does not show that the coefficients themselves are real, nor that they are invariant under the reflections, since the dilations of one loop see only the combinations of the coefficients which occur in that loop.

\subsection{Proof of the Symmetries}\label{sec:symmetries-proof}

\begin{proof}[Proof of Corollary~\ref{cor:symmetries}]
Each of the ensembles $\mathcal C'$ of the statement, namely $\mathcal C$ with its loops in another order, $g\mathcal C$, $\bar{\mathcal C}$, $\mathcal C$ with one loop reversed and its character conjugated, and $h\mathcal C$, consists of admissible loops at level $q$ which are pairwise vertex-disjoint, the maps $g$ and $h$ being bijections, and has $c_{\mathcal C'} = c_{\mathcal C}$, because $c_{\bar\pi} = c_\pi$. So $q \geq q_\sharp(c_{\mathcal C'}, \bar g)$, the limit $W(\mathcal C')$ exists by Theorem~\ref{thm:main}, and $\bar g$ and $q$ are as in that theorem for every single loop of $\mathcal C$ with its character, since $c_{\pi_i} \leq c_{\mathcal C}$ and $q_\sharp$ is nondecreasing. It suffices to prove the identities for $W_K$ at a fixed $K > j_q$.

(i) Both the observable \eqref{eq:ensemble-observable} and the product of the normalisations in \eqref{eq:renormalised-expectation} are products over the loops, which do not depend on the order of the factors.

(ii) By Lemma~\ref{lem:lattice-invariance}, $\mathbb E_K[W_{g\mathcal C}] = \mathbb E_K[W_{\mathcal C} \circ g] = \mathbb E_K[W_{\mathcal C}]$, and by Proposition~\ref{prop:reflection-normalisation}, $\kappa^{(q)}_K(gC_i, \chi_i) = \kappa^{(q)}_K(C_i, \chi_i)$ for every $i$. Dividing gives $W_K(g\mathcal C) = W_K(\mathcal C)$.

(iii) Since $W_{\bar\chi_i}(C_i) = \overline{W_{\chi_i}(C_i)}$, as in the proof of Lemma~\ref{lem:flat-values}, Lemma~\ref{lem:lattice-invariance} gives $\mathbb E_K[W_{\bar{\mathcal C}}] = \overline{\mathbb E_K[W_{\mathcal C}]}$, and Proposition~\ref{prop:reflection-normalisation} gives $\kappa^{(q)}_K(C_i, \bar\chi_i) = \overline{\kappa^{(q)}_K(C_i, \chi_i)}$; dividing gives $W_K(\bar{\mathcal C}) = \overline{W_K(\mathcal C)}$. If $C_i$ is replaced by $C_i^{-1}$ and $\chi_i$ by $\bar\chi_i$, the observable does not change, since $W_{\bar\chi_i}(C_i^{-1}) = W_{\chi_i}(C_i)$, and neither does the normalisation, since $\kappa^{(q)}_K(C_i^{-1}, \bar\chi_i) = \kappa^{(q)}_K(C_i, \chi_i)$ by the third identity of \eqref{eq:reflection-normalisation} applied to $(C_i, \bar\chi_i)$.

(iv) By Lemma~\ref{lem:lattice-invariance}, $\mathbb E_K[W_{h\mathcal C}] = \mathbb E_K[W_{\mathcal C}]$, and by Proposition~\ref{prop:reflection-normalisation}, $\kappa^{(q)}_K(hC_i, \chi_i) = \kappa^{(q)}_K(g_iC_i, \chi_i) = \kappa^{(q)}_K(C_i, \chi_i)$ for every $i$.
\end{proof}

The proof of (iv) shows what is missing for a general symmetry $h$ which exchanges two axes. Writing $h(x) = Ox + t$, the image $hC$ has $N^{(0)}_v(hC) = N^{(0)}_{O^{-1}v}(C)$ and the pairs $(Ov_c, Ow_c)$ at its corners, so that Lemma~\ref{lem:lattice-invariance} and Proposition~\ref{prop:normalisation-structure} give
\begin{equation}\label{eq:permutation-defect}
    \frac{W_K(h\mathcal C)}{W_K(\mathcal C)} = \prod_{i=1}^{n}\exp\Bigl(\sum_{v \in \mathcal V}N^{(0)}_v(C_i)\bigl(m_{v,K} - m_{Ov,K}\bigr) + \sum_{c \in \mathcal K(C_i)}\bigl(z_{v_c,w_c,K} - z_{Ov_c,Ow_c,K}\bigr)\Bigr)
\end{equation}
whenever $W_K(\mathcal C) \neq 0$, with the coefficients of the characters $\chi_i$. The covariance of the limits under $h$ therefore follows from the equality of the coefficients of the orientations $v$ and $Ov$ and of the corner pairs $(v, w)$ and $(Ov, Ow)$, and the method of Proposition~\ref{prop:reflection-normalisation} would give it from the invariance $N_j[F \circ h](1; \ell_0) = N_j[F](1; \ell_0)$ of the flat values. We do not know whether the construction has this invariance. The averages, the first gauge fixing, the designated bonds \cite[p.~267, ll.~19--20]{BalabanRG1}, and the $\mathbf R$-operation, whose graphs are averaged over the permutations of the coordinates as well \cite[p.~201, ll.~34--38]{BalabanLFR1}, are symmetric under the exchange of two axes; but the fluctuation fields of the steps $k \geq 2$ are written in the axial gauge of \cite[(1.19), p.~79]{BalabanREG}, whose contours change the coordinates in one fixed order \cite[(1.7), p.~18]{BalabanPRT1}, and the Landau gauge into which they are then transformed is defined relative to this preliminary choice, since the restricted gauge transformations do not form a group and the orbit ``depends on an element $U$ chosen to generate the set'' \cite[pp.~80--81]{BalabanREG}. The test functions of the fluctuation fields \cite[(3.16), p.~268]{BalabanCRE}, and with them the term $\ell_0$, are therefore not known to be invariant under the exchange of two axes. Ba{\l}aban states the covariance of his effective actions under every Euclidean transformation which preserves the lattices \cite[p.~269, ll.~9--26]{BalabanRG1}, \cite[(3.39), p.~274]{BalabanCRE}, but the operations from which we would derive the invariance of the numerator include this axial gauge, and we have not found an argument which removes it. An axial gauge in the averaged variables, as for the first gauge fixing, would give the symmetry, at the price of proving Ba{\l}aban's existence, uniqueness, and bounds for the Landau gauge again from it.

The positivity has no counterpart for the reflections which exchange two axes, for a reason which has nothing to do with the normalisations. Such a reflection is $x \mapsto Ox + t$ with $O$ the reflection in the hyperplane orthogonal to $e_\mu - e_\nu$ or to $e_\mu + e_\nu$. Its fixed points, if there are any, satisfy one equation $x_\mu \mp x_\nu \equiv c$ modulo $T$, so that they form a single hyperplane of $\mathbb T$, whose complement is connected: it is the image of the connected set $(0, T) \times (\mathbb R/T\mathbb Z)^3$ under the continuous map which sends $(s, y)$ to the point whose coordinates other than $x_\mu$ are those of $y$ and for which $x_\mu \mp x_\nu = c + s$. There are therefore no two half-tori which the reflection exchanges and whose union with the fixed points is the torus, and the proof of Lemma~\ref{lem:lattice-reflection-positivity} has nothing to apply to.

\subsection{Proof of the Reflection Positivity}\label{sec:proof-reflection-positivity}

\begin{proof}[Proof of Corollary~\ref{cor:reflection-positivity}]
Let $\theta = \theta_{\mu,s}$ with $s \in uL^{-q}\mathbb Z$. The two half-tori bounded by its fixed points are $H_{\mu,s}$ and $H_{\mu,s+T/2}$, and $\theta_{\mu,s+T/2} = \theta_{\mu,s}$, because $2s + T \equiv 2s$ modulo $T$; since $T/2 = L^{m_T}u \in uL^{-q}\mathbb Z$, we may therefore assume $H_+ = H_{\mu,s}$, and we write $\Theta \coloneqq \Theta_{\mu,s}$ and $H_- \coloneqq \theta H_+$. For $0 \leq a, b \leq m$ the family $\mathcal D_a \cup \theta\bar{\mathcal D}_b$ is a loop ensemble at level $q$ with admissible loops: the loops of $\mathcal D_a$ lie in $H_+$ and those of $\theta\bar{\mathcal D}_b$ in $H_-$, which is disjoint from $H_+$, so that they are pairwise vertex-disjoint, and the images are admissible loops at level $q$ by Section~\ref{sec:results}. Its number $c_{\mathcal D_a \cup \theta\bar{\mathcal D}_b} = \max(c_{\mathcal D_a}, c_{\mathcal D_b})$ is at most $c$, because $c_{\theta\bar{\mathcal D}_b} = c_{\mathcal D_b}$, so that $q \geq q_\sharp(c, \bar g) \geq q_\sharp(c_{\mathcal D_a \cup \theta\bar{\mathcal D}_b}, \bar g)$. By Theorem~\ref{thm:main} the limits $W(\mathcal D_a \cup \theta\bar{\mathcal D}_b)$ exist, with $W_K(\emptyset) = 1$ for every $K$.

Fix $K > j_q$, write $F_a \coloneqq W_{\mathcal D_a}$, with $F_0 \coloneqq 1$, and $\kappa_a$ for the product of the normalisations $\kappa^{(q)}_K$ of the loops of $\mathcal D_a$, with $\kappa_0 \coloneqq 1$. By the last assertion of Proposition~\ref{prop:reflection-normalisation}, which applies to each loop of $\mathcal D_b$ because its number $c_\pi \leq c$ gives $q \geq q_\sharp(c_\pi, \bar g)$, the product of the normalisations of the loops of $\theta\bar{\mathcal D}_b$ is $\bar\kappa_b$. Since $F_a \in \mathcal A_+$ and $\Theta F_b = W_{\theta\bar{\mathcal D}_b}$ by Section~\ref{sec:lattice-reflection-positivity}, the functions $G_a \coloneqq F_a/\kappa_a \in \mathcal A_+$ satisfy, by Definition~\ref{def:renormalised-expectation},
\begin{equation}\label{eq:gram-identity}
    W_K(\mathcal D_a \cup \theta\bar{\mathcal D}_b) = \frac{\mathbb E_K[F_a\,\Theta F_b]}{\kappa_a\bar\kappa_b} = \mathbb E_K[G_a\,\Theta G_b] = \langle G_a, G_b\rangle\;,
\end{equation}
because $\Theta$ is antilinear. For $\xi \in \mathbb C^m$ we obtain $\sum_{a,b}\xi_a\bar\xi_bW_K(\mathcal D_a \cup \theta\bar{\mathcal D}_b) = \langle\sum_a\xi_aG_a, \sum_b\xi_bG_b\rangle \geq 0$ by Lemma~\ref{lem:lattice-reflection-positivity}, and the matrix is Hermitian because the form is. Letting $K \to \infty$ entry by entry, the limit matrix is Hermitian and positive semidefinite, which is (i).

For (ii), $W_K(\theta\bar{\mathcal D}_a) = \mathbb E_K[\Theta F_a]/\bar\kappa_a = \overline{\mathbb E_K[F_a]/\kappa_a} = \overline{W_K(\mathcal D_a)}$ by Lemma~\ref{lem:lattice-reflection-positivity}, and the limit gives $W(\theta\bar{\mathcal D}_a) = \overline{W(\mathcal D_a)}$. The mean of $G_a$ is $w_a \coloneqq W_K(\mathcal D_a)$, and the mean of $\Theta G_b$ is $\bar w_b$, so that
\begin{equation}\label{eq:connected-gram}
    \langle G_a - w_a, G_b - w_b\rangle = W_K(\mathcal D_a \cup \theta\bar{\mathcal D}_b) - w_a\bar w_b\;,
\end{equation}
and the matrix of the right sides is Hermitian and positive semidefinite by the same argument; its limit is the matrix of (ii).
\end{proof}

\subsection{Nonvanishing and the Growth Bound}\label{sec:nonvanishing}

The connected functions of Corollary~\ref{cor:reflection-positivity}(ii), and the logarithms of Theorem~\ref{thm:non-gaussian}, are useful only if the limits do not vanish, and we now show that they do not, on the ensembles whose loops are not too long for the coupling.

\begin{proposition}[Nonvanishing]\label{prop:nonvanishing}
Let $L$ and $\Lambda_*$ be as in Theorem~\ref{thm:main}, and let $q \geq 0$. There are $\bar g_1 \in (0, \gamma_0]$ and a function $\ell \colon [0, \infty) \times (0, \bar g_1] \to [0, \infty]$, nonincreasing in its first argument and with $\ell(c, \bar g) \to \infty$ as $\bar g \to 0$ for every $c$, which depend only on the constants of the construction, $\Lambda_*$, $q$, and $T/u$, such that the following holds for every $\bar g \in (0, \bar g_1]$. If $\mathcal C$ is a loop ensemble at level $q$ whose loops are admissible and whose lengths satisfy
\begin{equation}\label{eq:nonvanishing-length}
    \sum_{i=1}^{n}\Bigl(P_q(C_i)^2 + \frac{P(C_i)}{r(C_i)}\Bigr) \leq \ell(c_{\mathcal C}, \bar g)\;,
\end{equation}
then $q \geq q_\sharp(c_{\mathcal C}, \bar g)$, the limit $W(\mathcal C)$ exists, and, for every $K \geq K_0(\mathcal C)$,
\begin{equation}\label{eq:nonvanishing-bound}
    \abs{W_K(\mathcal C)} \geq \frac12\;,
    \qquad\text{so that}\qquad
    \abs{W(\mathcal C)} \geq \frac12\;.
\end{equation}
\end{proposition}

The condition \eqref{eq:nonvanishing-length} bounds the lengths of the loops by a number which depends on their representations through $c_{\mathcal C}$ and grows without limit as the coupling decreases. Unlike the conditions of Theorem~\ref{thm:main}, it does restrict the representations at a given coupling: since a loop has at least four bonds, it fails whenever $\ell(c_{\mathcal C}, \bar g) < 16$. It involves the number of loops only through the sum, and for the ensembles which are generated from a fixed one by reflections in coordinate hyperplanes it is needed only for the ensemble from which they are generated (Proposition~\ref{prop:reflection-products}).

For the lower bound we need that the Wilson loop of a nearly flat field is close to one, with a bound in terms of the length of the loop and the size of a box containing it.

\begin{lemma}[Holonomy of a Nearly Flat Field]\label{lem:flat-holonomy}
Let $U$ be a $G$-valued field on a periodic lattice of spacing $\eta$, let $B$ be a box of side $D\eta$ in it, with $D \geq 1$ an integer smaller than every period in lattice units, and let $C$ be a closed path of $P$ bonds in $B$. If $\norm{U(\partial p) - 1} \leq t$ for every plaquette $p$ of $B$ and $s \coloneqq 3DPt \leq \min(1/10, \log(1 + r_G))$, then for every irreducible unitary representation $\pi$ of $G$ with the character $\chi$,
\begin{equation}\label{eq:flat-holonomy}
    \norm{\pi(U(C)) - 1} \leq 1.06\,c_\pi s
    \qquad\text{and}\qquad
    \abs{\chi(U(C)) - 1} \leq 1.06\,c_\pi s\;.
\end{equation}
\end{lemma}

\begin{proof}
Let $o$ be the corner of $B$ with the smallest coordinates, and for a site $y$ of $B$ let $\gamma_y$ be the path from $o$ to $y$ which changes first the coordinate $0$, then $1$, $2$, and $3$. For a bond $b = (y, y')$ of $B$ put $\ell_b \coloneqq \gamma_y\,b\,\gamma_{y'}^{-1}$, a closed path at $o$, and $U^\gamma(b) \coloneqq U(\ell_b)$. Then $U^\gamma$ is the gauge transform of $U$ by $y \mapsto U(\gamma_y)^{-1}$, so that $U^\gamma(C)$ is conjugate to $U(C)$ and $\chi(U^\gamma(C)) = \chi(U(C))$.

We first bound $U^\gamma(b)$. It suffices to take $b$ in a positive direction $e_\nu$, since $U^\gamma(b^{-1}) = U^\gamma(b)^{-1}$. Let $p$ be the path formed by the legs of $\gamma_y$ in the directions $\nu' > \nu$, a path of $n \leq 3D$ bonds from a site $z$ to $y$. Then $\gamma_y = \gamma_z\,p$ and $\gamma_{y'} = \gamma_z\,e\,(p + e_\nu)$, with $e$ the bond from $z$ to $z + e_\nu$, so that $\ell_b = \gamma_z\,p\,b\,(p + e_\nu)^{-1}e^{-1}\gamma_z^{-1}$, and $U(\ell_b)$ is conjugate to the holonomy $H$ of the boundary of the strip swept by $p$ in the direction $e_\nu$. With $p = (p_1, \dots, p_n)$ and $q_l$ the plaquette spanned by $p_l$ and $e_\nu$, telescoping along the strip writes $H = \lambda_n\cdots\lambda_1$, where $\lambda_l$ is the holonomy of $q_l$, from one of its corners and in one of its orientations, conjugated by the transport along $p_1 \cdots p_{l-1}$. Each $\lambda_l$ satisfies $\norm{\lambda_l - 1} = \norm{U(\partial q_l) - 1} \leq t$, since the operator norm of $W - 1$ is invariant under unitary conjugation and inversion, and \eqref{eq:product-deviation} gives
\begin{equation}\label{eq:axial-bond}
    \norm{U^\gamma(b) - 1} = \norm{H - 1} \leq (1 + t)^{3D} - 1\;.
\end{equation}
The path $C$ has $P$ bonds, and \eqref{eq:product-deviation} applied to $U^\gamma(C) = \prod_{b \in C}U^\gamma(b)$ gives $\norm{U^\gamma(C) - 1} \leq (1 + t)^{3DP} - 1 \leq e^s - 1$. Since $s \leq \min(1/10, \log(1 + r_G))$, the principal logarithm $U^\gamma(C) = e^{iX}$ has $X \in \mathfrak g$ and $\norm{X} \leq 1.06\,s$ by \eqref{eq:log-bound}, so that $\norm{\pi(U^\gamma(C)) - 1} \leq c_\pi\norm{X}$ by \eqref{eq:representation-constant}. Since $\pi(U^\gamma(C))$ is conjugate to $\pi(U(C))$ by a unitary matrix, $\norm{\pi(U(C)) - 1} = \norm{\pi(U^\gamma(C)) - 1} \leq 1.06\,c_\pi s$. Finally $\abs{\chi(W) - 1} = \abs{(\dim\pi)^{-1}\operatorname{Tr}(\pi(W) - 1)} \leq \norm{\pi(W) - 1}$.
\end{proof}

At the stopping step we apply the lemma to the background field. For a loop $C$ at level $q$ with the character of $\pi$ put
\begin{equation}\label{eq:flat-loop-constant}
    y(C) \coloneqq 3.18\,c_\pi c_h(1 + 2\kappa_s)B_3\,P_{K'}(C)^2\,\varepsilon_{K'}\;,
\end{equation}
which does not depend on $K$, since $P_{K'}(C)$ and $\varepsilon_{K'}$ do not (Section~\ref{sec:stopping-scale}).

\begin{proof}[Proof of Proposition~\ref{prop:nonvanishing}]
We first bound the leading words on $\ell_0$. On the support of $\varrho_{\ell_0}$ every fine plaquette of the background satisfies $\norm{U_{K'}(\partial p) - 1} \leq t \coloneqq 2c_h(1 + 2\kappa_s)B_3\varepsilon_{K'}L^{-2K'}$, as in the proof of Corollary~\ref{cor:background-stopping}. In fine units the loop $C_i$ has $P \coloneqq P_{K'}(C_i)L^{K'}$ bonds, and it lies in a box of side $D \leq P/2$, since a closed path of $P$ bonds has $\ell^\infty$-diameter at most $P/2$; the box embeds in the torus because $3\diam(C_i) < T$. Hence $3DPt \leq 3c_h(1 + 2\kappa_s)B_3P_{K'}(C_i)^2\varepsilon_{K'}$, and Lemma~\ref{lem:flat-holonomy} gives $\abs{f^{(i)}_{K'} - 1} \leq 1.06\,c_{\pi_i} \cdot 3DPt \leq y(C_i)$ whenever $3c_h(1 + 2\kappa_s)B_3P_{K'}(C_i)^2\varepsilon_{K'} \leq \min(1/10, \log(1 + r_G))$. Since $\abs{f^{(i)}_{K'}} \leq 1$, telescoping gives $\abs{\prod_if^{(i)}_{K'} - 1} \leq \sum_iy(C_i)$.

Next we bound the term $\ell_0$ of the identity \eqref{eq:stopping-identity} from below. On $\ell_0$ the normalised dressed loop is $e^{E_{\mathrm{const},K'}}[\prod_if^{(i)}_{K'} + \sum_{\mathfrak B \neq \emptyset}Q_{K'}(\mathfrak B)w_{K'}(\mathfrak B)]$, where $\abs{E_{\mathrm{const},K'}} \leq B_E(\mathcal C)$ by Definition~\ref{def:dressed-loop}(iv) and Corollary~\ref{cor:remainder-sum}, and the sum over the families has modulus at most $e^{I_*(\mathcal C)} - 1$ by Lemma~\ref{lem:term-bound}. Hence
\begin{equation}\label{eq:trivial-term-from-below}
    \bigl\lvert\bigl\langle\widehat{\mathcal W}_{\ell_0}\bigr\rangle_{\ell_0}\bigr\rvert \geq e^{-B_E(\mathcal C)}\Bigl(2 - e^{I_*(\mathcal C)} - \sum_iy(C_i)\Bigr)\;.
\end{equation}
The other terms are bounded as in the proof of Theorem~\ref{thm:uniform-bound}: $\sum_{\ell \neq \ell_0}w_\ell\abs{\langle\widehat{\mathcal W}_\ell\rangle_\ell} \leq e^{B_E(\mathcal C) + I_*(\mathcal C)}\bar\delta_{\mathrm{LF}}$, and $w_{\ell_0} \geq 1 - \bar\delta_{\mathrm{LF}}$ by Proposition~\ref{prop:large-field-weight}. By \eqref{eq:stopping-identity} and the triangle inequality, for every $K \geq K_0(\mathcal C)$,
\begin{equation}\label{eq:nonvanishing-lower}
    \abs{W_K(\mathcal C)} \geq \bigl(1 - \bar\delta_{\mathrm{LF}}\bigr)e^{-B_E(\mathcal C)}\Bigl(2 - e^{I_*(\mathcal C)} - \sum_iy(C_i)\Bigr) - e^{B_E(\mathcal C) + I_*(\mathcal C)}\bar\delta_{\mathrm{LF}}\;.
\end{equation}

It remains to choose $\bar g_1$ and $\ell$ so that the four quantities $B_E(\mathcal C)$, $I_*(\mathcal C)$, $\sum_iy(C_i)$, and $\bar\delta_{\mathrm{LF}}$ are small. Each of them is at most a sum over the loops of a number of the loop, multiplied by a constant which is a nondecreasing function of $c_{\mathcal C}$ and by a function of $\bar g$ which tends to zero, at fixed $q$ and $T/u$, except for $\bar\delta_{\mathrm{LF}}$, which does not depend on the ensemble. For $B_E(\mathcal C)$ this is \eqref{eq:remainder-sum}, with the constants $C'_E$ and $C''_E$, since $\bar{\mathfrak a} \leq C\bar g^2(\log\bar g^{-2})^{6p_0}$ by \eqref{eq:amplitude-bound} and $S^{\mathbf R}(\bar g) \to 0$; for $\bar\delta_{\mathrm{LF}}$ it is the remark after Proposition~\ref{prop:large-field-weight}. For the other two we note that $\hat g_{K'}^{-2} = \bar g^{-2} + b'j_q + C_W$ by \eqref{eq:tube-couplings}, so that $\hat g_{K'} \leq \bar g$, and $\log\hat g_{K'}^{-2} \leq 2\log\bar g^{-2}$ for $\bar g$ small, because $j_q$ grows only like $\log\log\bar g^{-2}$ by \eqref{eq:stopping-identity-index}. The numbers $P_{K'}(C) = P_q(C)L^3MR_{K'-1}$ are at most a constant times $P_q(C)(\log\hat g_{K'}^{-2})^r$ by Lemma~\ref{lem:tube-cubes}; by Definition~\ref{def:construction}, $\varepsilon_{K'} = \hat g_{K'}p_0(\hat g_{K'})$ is $\hat g_{K'}$ times a power of $\log\hat g_{K'}^{-2}$, and $\varepsilon_{K'}/\min(\alpha_{0,K'}, \alpha_{1,K'})$ is a power of $\log\hat g_{K'}^{-2}$; the number $\bar\rho$ of \eqref{eq:sup-amplitudes} is $\bar g^2$ times a power of $\log\bar g^{-2}$; and $S_{\mathrm{wr}}(b) \leq e^{-b}(1 + 2/b)$ by Lemma~\ref{lem:scale-sums} with $x = 1$. The cross total \eqref{eq:cross-total} is at most $\sum_iP_{K'}(C_i)$ times $10\,C_{\mathrm{cr}}\bar{\mathfrak a} + 2C^{\mathbf R}_{\mathrm{cr}}S^{\mathbf R}(\bar g)$, which tends to zero with $\bar g$ as in the bound on $B_E(\mathcal C)$. Hence, apart from the factors $P_q(C_i)$, $c_{\pi_i}$, and the constants $C_\chi$ and $C$ of \eqref{eq:insertion-bound}, the numbers \eqref{eq:flat-loop-constant} and \eqref{eq:insertion-total} are at most a positive power of $\bar g$ times a power of $\log\bar g^{-2}$. Collecting, there are functions $h_E$, $h_I$, and $h_y$ on $(0, \gamma_0]$ which tend to zero as $\bar g \to 0$, and nondecreasing functions $F_E$ and $F_I$ on $[0, \infty)$, all of which depend only on the constants of the construction, $\Lambda_*$, $q$, and $T/u$, such that
\begin{equation}\label{eq:nonvanishing-quantities}
\begin{gathered}
    B_E(\mathcal C) \leq F_E(c_{\mathcal C})h_E(\bar g)\sum_i\Bigl(\frac{P(C_i)}{r(C_i)} + P_q(C_i)\Bigr)\;,\\
    I_*(\mathcal C) \leq F_I(c_{\mathcal C})h_I(\bar g)\sum_iP_q(C_i)\;,
    \qquad
    y(C_i) \leq c_{\pi_i}h_y(\bar g)P_q(C_i)^2\;,
\end{gathered}
\end{equation}
and such that $3c_h(1 + 2\kappa_s)B_3P_{K'}(C_i)^2\varepsilon_{K'} \leq h_y(\bar g)P_q(C_i)^2$. We let $\bar g_1 \in (0, \gamma_0]$ be so small that $\bar\delta_{\mathrm{LF}} \leq 1/20$ for $\bar g \leq \bar g_1$, put $m_G \coloneqq \min(1/10, \log(1 + r_G))$, and define $\ell(c, \bar g) \coloneqq 0$ if $q < q_\sharp(c, \bar g)$ and otherwise
\begin{equation}\label{eq:nonvanishing-function}
    \ell(c, \bar g) \coloneqq \min\Bigl(\frac{1}{10F_E(c)h_E(\bar g)},\ \frac{1}{20F_I(c)h_I(\bar g)},\ \frac{1}{20\,c\,h_y(\bar g)},\ \frac{m_G}{h_y(\bar g)}\Bigr)\;,
\end{equation}
with the convention $1/0 \coloneqq \infty$. Each entry of the minimum does not increase with $c$, and neither does the indicator of $q \geq q_\sharp(c, \bar g)$, since $q_\sharp$ is nondecreasing in $c$; so $\ell$ does not increase with $c$. For fixed $c$ we have $q_\sharp(c, \bar g) = 0 \leq q$ for $\bar g$ small, by Theorem~\ref{thm:main}, and every entry of the minimum tends to infinity as $\bar g \to 0$; so $\ell(c, \bar g) \to \infty$.

Let now $\bar g \leq \bar g_1$ and let $\mathcal C$ satisfy \eqref{eq:nonvanishing-length}. A closed path which visits no site twice has at least four bonds, so that the left side of \eqref{eq:nonvanishing-length} is at least $16$; hence $\ell(c_{\mathcal C}, \bar g) > 0$, which by definition means $q \geq q_\sharp(c_{\mathcal C}, \bar g)$, and Theorem~\ref{thm:main} gives the limit $W(\mathcal C)$. Since $P_q(C_i) \leq P_q(C_i)^2$, the bounds \eqref{eq:nonvanishing-quantities} and \eqref{eq:nonvanishing-function} give $B_E(\mathcal C) \leq 1/10$, $I_*(\mathcal C) \leq 1/20$, $\sum_iy(C_i) \leq c_{\mathcal C}h_y(\bar g)\sum_iP_q(C_i)^2 \leq 1/20$, and the condition of Lemma~\ref{lem:flat-holonomy} for every loop, and $\bar\delta_{\mathrm{LF}} \leq 1/20$ by the choice of $\bar g_1$. Then \eqref{eq:nonvanishing-lower} gives
\begin{equation}\label{eq:nonvanishing-number}
    \abs{W_K(\mathcal C)} \geq 0.95\,e^{-0.1}\bigl(2 - e^{0.05} - 0.05\bigr) - 0.05\,e^{0.15} \geq 0.71 \geq \frac12\;,
\end{equation}
and in the limit $\abs{W(\mathcal C)} \geq 1/2$. For a loop $C$ at level $q$ in $H_+$ for which $\mathcal D \coloneqq ((C, \chi))$ satisfies \eqref{eq:nonvanishing-length}, we have $q \geq q_\sharp(c_\pi, \bar g)$, as just shown, and Corollary~\ref{cor:reflection-positivity} with $m = 1$, $\mathcal D_1 = \mathcal D$, and $c = c_\pi$ gives $W(\mathcal D \cup \theta\bar{\mathcal D}) - \abs{W(\mathcal D)}^2 \geq 0$, so that $W(\mathcal D \cup \theta\bar{\mathcal D}) \geq \abs{W(\mathcal D)}^2 \geq 1/4$.
\end{proof}

The same quantities give an upper bound which factorises over the loops: each loop contributes a factor which is exponential in its length in units of $\mathbb T_q$ and in its ratio $P/r$, with a rate which tends to zero with the coupling, and the number of loops enters only through the number of factors.

\begin{proposition}[Growth Bound]\label{prop:growth-bound}
Let $L$, $\Lambda_*$, and $q \geq 0$ be as in Proposition~\ref{prop:nonvanishing}, and let $\bar g_1$, $F_E$, $F_I$, $h_E$, and $h_I$ be as in its proof. Put $a(c, \bar g) \coloneqq F_E(c)h_E(\bar g) + F_I(c)h_I(\bar g)$, a function which is nondecreasing in $c$ and tends to zero as $\bar g \to 0$ for every $c$. For every $\bar g \in (0, \bar g_1]$ and every loop ensemble $\mathcal C$ at level $q \geq q_\sharp(c_{\mathcal C}, \bar g)$ whose loops are admissible, and for every $K \geq K_0(\mathcal C)$,
\begin{equation}\label{eq:growth-bound}
    \abs{W_K(\mathcal C)} \leq \frac{21}{20}\prod_{i=1}^{n}\exp\Bigl(a(c_{\mathcal C}, \bar g)\Bigl(\frac{P(C_i)}{r(C_i)} + P_q(C_i)\Bigr)\Bigr)\;,
\end{equation}
and the limit $W(\mathcal C)$ satisfies the same bound.
\end{proposition}

\begin{proof}
By Theorem~\ref{thm:uniform-bound}, $\abs{W_K(\mathcal C)} \leq e^{B_E(\mathcal C) + I_*(\mathcal C)}(1 + \bar\delta_{\mathrm{LF}})$, and $\bar\delta_{\mathrm{LF}} \leq 1/20$ for $\bar g \leq \bar g_1$ by the choice of $\bar g_1$. The bounds \eqref{eq:nonvanishing-quantities} on $B_E(\mathcal C)$ and $I_*(\mathcal C)$ were derived in the proof of Proposition~\ref{prop:nonvanishing} from the definitions \eqref{eq:remainder-sum} and \eqref{eq:insertion-total} alone, without the condition \eqref{eq:nonvanishing-length}, and $P_q(C_i) \leq P(C_i)/r(C_i) + P_q(C_i)$ gives $B_E(\mathcal C) + I_*(\mathcal C) \leq a(c_{\mathcal C}, \bar g)\sum_i(P(C_i)/r(C_i) + P_q(C_i))$. The limit exists by Theorem~\ref{thm:main}, and the bound passes to it.
\end{proof}

Since $P_q(C)uL^{-q}$ is the length of $C$, the bound \eqref{eq:growth-bound} is a perimeter bound with the rate $a(c_{\mathcal C}, \bar g)L^q/u$ per unit length, together with a factor for the shape of each loop; with Proposition~\ref{prop:nonvanishing} it brackets the limits between $\frac12$ and this bound on the ensembles of short loops. Continuity of $W$ in the positions of the loops is not asserted, since the limits are defined on the loops of the lattices $\mathbb T_q$ only.

The condition \eqref{eq:nonvanishing-length} involves the number of loops, because the bounds on $B_E(\mathcal C)$, $I_*(\mathcal C)$, and $\sum_iy(C_i)$ are bounds on sums over the loops. For ensembles which are generated from a fixed one by reflections this dependence disappears. We say that an ensemble $\mathcal D$ at level $q$ is \emph{generated by reflections from} an ensemble $\mathcal B$ if $\mathcal D = \mathcal B$, or if $\mathcal D = \mathcal D' \cup \theta_{\mu,a}\bar{\mathcal D}'$ for an ensemble $\mathcal D'$ which is generated by reflections from $\mathcal B$ and all of whose loops lie in $H_{\mu,a}$, for some $\mu$ and some $a \in uL^{-q}\mathbb Z$. The \emph{blocks} of $\mathcal B$ are $\{ \mathcal B \}$, and those of $\mathcal D' \cup \theta_{\mu,a}\bar{\mathcal D}'$ are the blocks of $\mathcal D'$ and their images under $\theta_{\mu,a}$ with the conjugate characters. Every block is the image of $\mathcal B$ under a composition of reflections, with all its characters conjugated or none, so that it has the same chord--arc constants, ratios $P/r$, and numbers $P_q$ as $\mathcal B$, and either the characters of $\mathcal B$ or their conjugates. In particular $c_{\mathcal B_i} = c_{\mathcal B} = c_{\mathcal D}$ for every block $\mathcal B_i$ of $\mathcal D$, because $c_{\bar\pi} = c_\pi$, and every block satisfies \eqref{eq:nonvanishing-length} if $\mathcal B$ does.

\begin{proposition}[Products over Reflections]\label{prop:reflection-products}
Let $L$ and $\Lambda_*$ be as in Theorem~\ref{thm:main}, let $\mathcal B$ be a loop ensemble at level $q$ whose loops are admissible, with $\bar g$ and $q$ as in Theorem~\ref{thm:main} for $\mathcal B$, and let $\mathcal D$ be generated by reflections from $\mathcal B$, with the blocks $\mathcal B_1, \dots, \mathcal B_N$. Then $\abs{W_K(\mathcal D)} \geq \prod_{i=1}^{N}\abs{W_K(\mathcal B_i)}$ for every $K > j_q$. If $\bar g \leq \bar g_1$ and $\mathcal B$ satisfies \eqref{eq:nonvanishing-length}, then $\abs{W(\mathcal D)} \geq 2^{-N}$; the condition does not involve $\mathcal D$, and in particular not its number of loops.
\end{proposition}

\begin{proof}
We argue by induction along the generation. For $\mathcal D = \mathcal B$ there is nothing to prove. Let $\mathcal D = \mathcal D' \cup \theta_{\mu,a}\bar{\mathcal D}'$, with the loops of $\mathcal D'$ in $H_{\mu,a}$, and write $\Theta \coloneqq \Theta_{\mu,a}$ and $F \coloneqq W_{\mathcal D'} \in \mathcal A_+$. Then $W_{\mathcal D} = F\,\Theta F$ by Section~\ref{sec:lattice-reflection-positivity}, and \eqref{eq:reflection-covariance} with Lemma~\ref{lem:lattice-reflection-positivity} gives
\begin{equation}\label{eq:doubling-inequality}
    \abs{\mathbb E_K[W_{\mathcal D}]} = \mathbb E_K[F\,\Theta F] \geq \abs{\mathbb E_K[F]}^2 = \abs{\mathbb E_K[W_{\mathcal D'}]}\,\abs{\mathbb E_K[W_{\theta_{\mu,a}\bar{\mathcal D}'}]}\;.
\end{equation}
The normalisation of $\mathcal D$ is the product of those of $\mathcal D'$ and $\theta_{\mu,a}\bar{\mathcal D}'$, and dividing by its modulus gives $\abs{W_K(\mathcal D)} \geq \abs{W_K(\mathcal D')}\abs{W_K(\theta_{\mu,a}\bar{\mathcal D}')}$. By the induction hypothesis $\abs{W_K(\mathcal D')}$ is at least the product over the blocks of $\mathcal D'$; and since $\theta_{\mu,a}\bar{\mathcal D}'$ is generated by reflections from $\theta_{\mu,a}\bar{\mathcal B}$, with the reflections conjugated by $\theta_{\mu,a}$, and has the images of the blocks of $\mathcal D'$ as its blocks, the same holds for it. This proves the first assertion. If $\bar g \leq \bar g_1$ and $\mathcal B$ satisfies \eqref{eq:nonvanishing-length}, then so does every block, so that $\abs{W_K(\mathcal B_i)} \geq 1/2$ for $K \geq K_0(\mathcal B_i)$ by Proposition~\ref{prop:nonvanishing}, and the limit gives $\abs{W(\mathcal D)} \geq 2^{-N}$; the limit exists by Theorem~\ref{thm:main}, because the loops of $\mathcal D$ are admissible and $c_{\mathcal D} = c_{\mathcal B}$, so that $\bar g$ and $q$ are as in Theorem~\ref{thm:main} for $\mathcal D$.
\end{proof}

The lower bound decays with the number of blocks, but it is positive at a coupling which is fixed once $\mathcal B$ is. The number of loops which can be generated in this way is limited only by the torus: starting from one loop of diameter $h$ and doubling in each direction at a hyperplane at distance $h$ from the current configuration, the extent in the direction $\mu$ after $k_\mu$ doublings is at most $2^{k_\mu}\cdot3h - 2h$, so that of the order of $4\log_2(T/h)$ doublings fit, and the number of loops is of the order of $(T/h)^4$.

\section{Non-Gaussianity}\label{sec:non-gaussianity}

We now prove Theorem~\ref{thm:non-gaussian}. The proof computes the number $\mathfrak Y$ of \eqref{eq:witness} in two ways. For a Gaussian charge model it vanishes identically, by an algebraic cancellation which uses nothing but two facts: $\log\abs{W}$ is a quadratic form in the current of the loop, linear in the covariance, and the logarithm of a normalisation is a linear function of the lengths and the corners (Section~\ref{sec:gaussian-models}). For the limits of Yang--Mills theory we compute $\mathfrak Y$ to the order $\bar g^2$ at the stopping scale, where the identity of Theorem~\ref{thm:stopping-identity} expresses $W_K$ through a measure on the fixed lattice $\mathbb T_{j_q}$ which, at small coupling, is concentrated near the flat fields.

Three things have to be shown, and each meets a difficulty of its own. The normalisations have to drop out of $\mathfrak Y$, which they do because $\mathfrak Y$ was built to annihilate every function of the form of Proposition~\ref{prop:normalisation-structure} (Section~\ref{sec:normalisations-drop-out}). The pair terms have to be of the fourth order. Under the measure of the stopping scale the Wilson loop of each of two separated rectangles depends, up to corrections which decay with their distance, on the field near that rectangle only, and its deviation from one has an $L^2$-norm of the order $\bar g^2$ rather than $\bar g$, because its part of first order is a trace of an element of the Lie algebra and vanishes; the covariance of two such deviations is then of the order $\bar g^4$ (Section~\ref{sec:pair-terms}). Finally the third difference of the rectangles has to be bounded below by a positive multiple of $\bar g^2$, and this is the heart of the matter. After the insertions are shown to cancel in the third difference, by the translation invariance which made the coefficients of the normalisation independent of the loop (Section~\ref{sec:shape-term}), the third difference is a combination of the second moments of the fluxes of the background field through the rectangles $I_a$. An integration by parts against the density of the stopping scale identifies these second moments, to the leading order, with a quadratic form evaluated at a flat field near the given one (Section~\ref{sec:flux-moments}). At a flat field this form is the covariance of the fluxes of the block field of a free field, which differs from the covariance of the free field itself by a term that decays on the scale of the blocks (Section~\ref{sec:flat-covariance}), and the third difference of the free covariances is computed explicitly (Section~\ref{sec:free-field}).

Throughout the section $q$, the plane spanned by $e_\mu$ and $e_\nu$, $t$, $s$, $o$, and $\pi$ are fixed as in Section~\ref{sec:results}, with $\Lambda_* \geq 1 + \max(s/t, 4t/s)$ and $L$ as in Theorem~\ref{thm:main}, and the eleven diagrams are those defined there. Every requirement in this section that $\bar g$ be small is an upper bound on $\bar g$ in terms of $q$, $t/u$, $s/u$, $T/u$, $\pi$, and the constants of the construction, and $\bar g_0$ in Theorem~\ref{thm:non-gaussian} is the least of these bounds.

\subsection{The Witness}\label{sec:witness}

For an ensemble $\mathcal D$ we write $\abs{\mathcal D}$ for the number of its loops, $N_v(\mathcal D) \coloneqq \sum_{C \in \mathcal D}N_v(C)$ for the total length of its segments of orientation $v \in \mathcal V$, and, for a perpendicular pair $(v, w)$, $K_{v,w}(\mathcal D)$ for the number of the corners $c$ of its loops with $(v_c, w_c) = (v, w)$. The number \eqref{eq:witness} is
\begin{equation}\label{eq:witness-weights}
    \mathfrak Y(W) = \sum_{\mathcal D}n_{\mathcal D}\log\abs{W(\mathcal D)} = \mathfrak y_S(W) - \mathfrak y_2(W) + \mathfrak y_3(W)\;,
\end{equation}
the sum over the six diagrams $R_1$, $R_2$, $R_3$, $R_4$, $P_2$, $P_3$ with the weights $n_{\mathcal D} = 1, -3, 3, -1, -1, 1$, where $\mathfrak y_S(W) \coloneqq \sum_{n=1}^{4}n_{R_n}\log\abs{W(R_n)}$ is the \emph{shape term} and $\mathfrak y_k(W) \coloneqq \log\abs{W(P_k)/W(R_1)^2}$ are the pair terms. To each diagram we assign the set $\operatorname{ind}(\mathcal D)$ of the indices $a$ of the rectangles $I_a$ which tile it, or which are its loops: $\operatorname{ind}(I_a) \coloneqq \{ a \}$, $\operatorname{ind}(R_n) \coloneqq \{ 0, \dots, n - 1 \}$, $\operatorname{ind}(R^{(1)}_n) \coloneqq \{ 1, \dots, n \}$, and $\operatorname{ind}(P_k) \coloneqq \{ 0, k \}$. The Gaussian models are tested by a second set of weights on all eleven diagrams,
\begin{equation}\label{eq:symmetric-weights}
    n^{\mathrm{sym}}_{\mathcal D} \coloneqq
    \begin{cases}
        0, 1, 1, -1 & \text{for } \mathcal D = I_0, I_1, I_2, I_3\;,\\
        -1, -2, 2, 1 & \text{for } \mathcal D = R_2, R^{(1)}_2, R_3, R^{(1)}_3\;,\\
        -1, -1, 1 & \text{for } \mathcal D = R_4, P_2, P_3\;,
    \end{cases}
\end{equation}
and we put $\mathfrak Y^{\mathrm{sym}}(W) \coloneqq \sum_{\mathcal D}n^{\mathrm{sym}}_{\mathcal D}\log\abs{W(\mathcal D)}$, the sum over the eleven diagrams.

\begin{lemma}[Counts and Index Sums]\label{lem:witness-counts}
\begin{enumerate}[label=(\roman*)]
\item A rectangle of width $at$ and height $s$ among the nine has $N_{\pm e_\mu} = at$, $N_{\pm e_\nu} = s$, and $N_v = 0$ for the other $v$, and it has exactly one corner of each of the types $(e_\mu, e_\nu)$, $(e_\nu, -e_\mu)$, $(-e_\mu, -e_\nu)$, $(-e_\nu, e_\mu)$ and no other corner.
\item Let $f$ be a function on the eleven diagrams such that $f(\mathcal D)$ depends only on $\mathcal D$ through $\abs{\mathcal D}$, the numbers $N_v(\mathcal D)$, and the numbers $K_{v,w}(\mathcal D)$, and linearly. Then $\sum_{\mathcal D}n_{\mathcal D}f(\mathcal D) = 0$, $\sum_nn_{R_n}f(R_n) = 0$, $f(P_k) = 2f(R_1)$, and $\sum_{\mathcal D}n^{\mathrm{sym}}_{\mathcal D}f(\mathcal D) = 0$.
\item If $f(I_a) = f(I_0)$ for every $a$ and $f(R^{(1)}_n) = f(R_n)$ for $n \in \{ 2, 3 \}$, then $\sum_{\mathcal D}n^{\mathrm{sym}}_{\mathcal D}f(\mathcal D) = \sum_{\mathcal D}n_{\mathcal D}f(\mathcal D)$, the left sum over the eleven diagrams and the right one over the six.
\item For $0 \leq a \leq b \leq 3$, the sum $\mathsf s(a, b)$ of $n^{\mathrm{sym}}_{\mathcal D}$ over the diagrams with $a, b \in \operatorname{ind}(\mathcal D)$ vanishes.
\end{enumerate}
\end{lemma}

\begin{proof}
(i) The rectangle is traversed from its corner of least coordinates in the directions $e_\mu$, $e_\nu$, $-e_\mu$, $-e_\nu$, in this order, which gives the four corners and the lengths.

(iii) The weights $n^{\mathrm{sym}}$ sum to $0 + 1 + 1 - 1 = 1$ over $I_0, \dots, I_3$, to $-1 - 2 = -3$ over $R_2$, $R^{(1)}_2$, and to $2 + 1 = 3$ over $R_3$, $R^{(1)}_3$, which are the weights $n$ of $R_1$, $R_2$, and $R_3$, and they agree with $n$ on $R_4$, $P_2$, $P_3$.

(ii) It suffices to take for $f$ each of the functions $\abs{\mathcal D}$, $N_v(\mathcal D)$, $K_{v,w}(\mathcal D)$. On the six diagrams $R_1, R_2, R_3, R_4, P_2, P_3$, by (i), $\abs{\mathcal D} = 1, 1, 1, 1, 2, 2$, $N_{\pm e_\mu}(\mathcal D) = t, 2t, 3t, 4t, 2t, 2t$, $N_{\pm e_\nu}(\mathcal D) = s, s, s, s, 2s, 2s$, the other $N_v$ vanish, and $K_{v,w}(\mathcal D) = \abs{\mathcal D}$ for the four types of (i) and $0$ for the others. Hence
\begin{equation}\label{eq:count-sums}
    \sum_{\mathcal D}n_{\mathcal D}\abs{\mathcal D} = 1 - 3 + 3 - 1 - 2 + 2 = 0\;,
    \qquad
    \sum_{\mathcal D}n_{\mathcal D}N_{\pm e_\mu}(\mathcal D) = (1 - 6 + 9 - 4 - 2 + 2)t = 0\;,
\end{equation}
and $\sum_{\mathcal D}n_{\mathcal D}N_{\pm e_\nu}(\mathcal D) = (1 - 3 + 3 - 1 - 2 + 2)s = 0$. On the rectangles alone the same sums are $1 - 3 + 3 - 1 = 0$, $(1 - 6 + 9 - 4)t = 0$, and $0$. The loops of $P_k$ have the lengths and the corners of $R_1$, which gives $f(P_k) = 2f(R_1)$. Finally $f$ takes the same value on a rectangle and on its translates, so that (iii) and the first assertion give the last one.

(iv) The diagrams which contain a given pair of indices, with their weights $n^{\mathrm{sym}}$, are: for $(0, 0)$, $I_0$, $R_2$, $R_3$, $R_4$, $P_2$, $P_3$, with $0 - 1 + 2 - 1 - 1 + 1$; for $(1, 1)$, $I_1$, $R_2$, $R^{(1)}_2$, $R_3$, $R^{(1)}_3$, $R_4$, with $1 - 1 - 2 + 2 + 1 - 1$; for $(2, 2)$, $I_2$, $R^{(1)}_2$, $R_3$, $R^{(1)}_3$, $R_4$, $P_2$, with $1 - 2 + 2 + 1 - 1 - 1$; for $(3, 3)$, $I_3$, $R^{(1)}_3$, $R_4$, $P_3$, with $-1 + 1 - 1 + 1$; for $(0, 1)$, $R_2$, $R_3$, $R_4$, with $-1 + 2 - 1$; for $(0, 2)$, $R_3$, $R_4$, $P_2$, with $2 - 1 - 1$; for $(0, 3)$, $R_4$, $P_3$, with $-1 + 1$; for $(1, 2)$, $R^{(1)}_2$, $R_3$, $R^{(1)}_3$, $R_4$, with $-2 + 2 + 1 - 1$; and for $(1, 3)$ and $(2, 3)$, $R^{(1)}_3$, $R_4$, with $1 - 1$. Each sum is zero.
\end{proof}

\subsection{Gaussian Charge Models}\label{sec:gaussian-models}

The Gaussian charge models of Section~\ref{sec:results} fail this criterion for an algebraic reason, which the next proposition records.

\begin{proposition}[Gaussian Charge Models]\label{prop:gaussian-models}
For every Gaussian charge model, its function $W$ does not vanish on the eleven diagrams and $\mathfrak Y^{\mathrm{sym}}(W) = 0$. If $W_m$ are the functions of a sequence of Gaussian charge models, each with its own charge, Gaussian vector, and normalisation, and $W_m(\mathcal D) \to W(\mathcal D) \neq 0$ for each of the eleven diagrams, then $\mathfrak Y^{\mathrm{sym}}(W) = 0$.
\end{proposition}

\begin{proof}
For a real Gaussian variable $Z$ we have $\abs{\mathbb E[e^{iZ}]} = e^{-\operatorname{Var}Z/2}$. Hence $W(\mathcal D) \neq 0$ and
\begin{equation}\label{eq:gaussian-log}
    \log\abs{W(\mathcal D)} = -\frac{\mathsf g^2}{2}\sum_{a,b \in \operatorname{ind}(\mathcal D)}\operatorname{Cov}(\mathsf X_a, \mathsf X_b) - \sum_{C \in \mathcal D}\log\abs{\kappa(C)}\;.
\end{equation}
We multiply by $n^{\mathrm{sym}}_{\mathcal D}$ and sum over the eleven diagrams. In the first term the coefficient of $\operatorname{Cov}(\mathsf X_a, \mathsf X_b)$, $a \leq b$, is $-\frac12\mathsf g^2$ or $-\mathsf g^2$ times $\mathsf s(a, b)$, which vanishes by Lemma~\ref{lem:witness-counts}(iv). The second term is a function of $\mathcal D$ which is linear in the numbers $N_v(\mathcal D)$ and $K_{v,w}(\mathcal D)$, by the definition of a Gaussian charge model, so that its weighted sum vanishes by Lemma~\ref{lem:witness-counts}(ii). Hence $\mathfrak Y^{\mathrm{sym}}(W) = 0$. For the limit, $\mathfrak Y^{\mathrm{sym}}(W) = \lim_m\mathfrak Y^{\mathrm{sym}}(W_m) = 0$, since $\log\abs{\cdot}$ is continuous on $\mathbb C \setminus \{ 0 \}$.
\end{proof}

The cancellation can be read off in the models whose covariance is invariant under the translation by $te_\mu$, $\operatorname{Cov}(\mathsf X_a, \mathsf X_b) = c^{\mathrm G}_{\abs{a-b}}$. There $\operatorname{Var}\mathsf X(R_n) = nc^{\mathrm G}_0 + 2\sum_{1 \leq k < n}(n - k)c^{\mathrm G}_k$ and $\operatorname{Var}\mathsf X(P_k) = 2c^{\mathrm G}_0 + 2c^{\mathrm G}_k$, so that, the normalisations cancelling by Lemma~\ref{lem:witness-counts}(ii),
\begin{equation}\label{eq:stationary-gaussian}
    \mathfrak y_S(W) = \mathsf g^2\bigl(c^{\mathrm G}_3 - c^{\mathrm G}_2\bigr)\;,
    \qquad
    \mathfrak y_k(W) = -\mathsf g^2c^{\mathrm G}_k\;,
    \qquad
    \mathfrak Y(W) = \mathfrak Y^{\mathrm{sym}}(W) = 0\;:
\end{equation}
the covariances of the separated rectangles which make up the shape term also make up the pair terms.

\begin{remark}[Gaussian Fields with Values in the Lie Algebra]\label{rem:lie-gaussian}
The models of Proposition~\ref{prop:gaussian-models} are abelian, and the number $\mathfrak Y$ does not distinguish Yang--Mills theory from Gaussian fields with values in $\mathfrak g$ and path-ordered holonomies. Let $(T_\alpha)$ be a basis of $\mathfrak g$ which is orthonormal for an $\mathrm{Ad}$-invariant inner product, with $\sum_\alpha\pi_*(T_\alpha)^2 = C\cdot1$, let $B = \sum_\alpha B^\alpha T_\alpha$ on the bonds of a lattice containing the rectangles, with $B^\alpha$ real, centred, and jointly Gaussian, $\operatorname{Cov}(B^\alpha(e), B^\beta(f)) = \delta^{\alpha\beta}G(e, f)$, and $G$ invariant under the translation by $te_\mu$, and let $W(\mathcal D) \coloneqq \mathbb E[\prod_{C \in \mathcal D}\chi(\prod_{e \in C}e^{i\mathsf gB(e)})]$. Expanding the ordered product to the second order in $\mathsf g$, the term of first order of $\chi(\prod_ee^{i\mathsf gB(e)})$ is $i\mathsf g(\dim\pi)^{-1}\operatorname{Tr}\pi_*(\sum_eB(e))$, which vanishes because $\operatorname{Tr}\circ\pi_*$ vanishes on $\mathfrak g$ (Lemma~\ref{lem:near-one} below), the term of second order has the mean $-\frac12\mathsf g^2C\operatorname{Var}\mathsf X^1(C)$ with $\mathsf X^1(C) \coloneqq \sum_{e \in C}B^1(e)$, the terms of odd order have mean zero, and the remainder of fourth order has a finite moment. Hence the pair terms are $O(\mathsf g^4)$ and $\mathfrak Y(W) = \mathsf g^2C(c_3 - c_2) + O(\mathsf g^4)$ with $c_k \coloneqq \operatorname{Cov}(\mathsf X^1(I_0), \mathsf X^1(I_k))$, which for the free covariance is positive by Lemma~\ref{lem:free-field} below. These fields share with Yang--Mills theory the feature on which the theorem rests.
\end{remark}

\subsection{The Normalisations Drop Out}\label{sec:normalisations-drop-out}

We next show that the limits exist on the eleven diagrams and that $\mathfrak Y(W)$ is the limit of the same number computed from the expectations $N_K(\mathcal D) \coloneqq \mathbb E_K[W_{\mathcal D}]$, in which no normalisation appears.

\begin{proposition}[The Normalisations Drop Out]\label{prop:normalisations-drop-out}
\begin{enumerate}[label=(\roman*)]
\item Let $K > j_q$. If $N_K(\mathcal D) \neq 0$ on the six diagrams $R_1, \dots, R_4, P_2, P_3$, then each of $\mathfrak Y$, $\mathfrak y_S$, $\mathfrak y_2$, and $\mathfrak y_3$ takes the same value on $W_K$ and on $N_K$.
\item There is $\bar g_1' > 0$ such that for $0 < \bar g \leq \bar g_1'$ the limits $W(\mathcal D)$ exist on the eleven diagrams, $\abs{W_K(\mathcal D)} \geq 1/2$ for $K \geq K_0(\mathcal D)$, $\abs{W(\mathcal D)} \geq 1/2$, $W_K(I_a) = W_K(I_0)$ and $W_K(R^{(1)}_n) = W_K(R_n)$ for $K > j_q$, and $W(I_a) = W(I_0)$ and $W(R^{(1)}_n) = W(R_n)$. For these $\bar g$, $\mathfrak Y^{\mathrm{sym}}(W) = \mathfrak Y(W) = \lim_K\mathfrak Y(N_K)$, $\mathfrak y_S(W) = \lim_K\mathfrak y_S(N_K)$, and $\mathfrak y_k(W) = \lim_K\mathfrak y_k(N_K)$.
\item Let $\bar g \leq \bar g_1'$ and $q' \geq q$. If the limits $W^{(q')}$ of the renormalised expectations with the normalisation of the level $q'$ exist and do not vanish on the six diagrams, then $\mathfrak Y(W^{(q')}) = \mathfrak Y(W)$.
\end{enumerate}
\end{proposition}

\begin{proof}
(i) By Proposition~\ref{prop:normalisation-structure}, $\log\abs{\kappa^{(q)}_K(C, \chi)} = \sum_vN^{(0)}_v(C)\operatorname{Re}m_{v,K} + \sum_{c \in \mathcal K(C)}\operatorname{Re}z_{v_c,w_c,K}$, where $N^{(0)}_v(C) = N_v(C)/\varepsilon_K$ is the number of bonds of $\mathbb T_K$ of orientation $v$ on $C$ and the coefficients do not depend on $C$. Hence $\sum_{C \in \mathcal D}\log\abs{\kappa^{(q)}_K(C, \chi)}$ is a linear function of the numbers $N_v(\mathcal D)$ and $K_{v,w}(\mathcal D)$, and by Lemma~\ref{lem:witness-counts}(ii) its weighted sums with the weights of $\mathfrak Y$ and of $\mathfrak y_S$ vanish, and its value on $P_k$ is twice its value on $R_1$. Since $\log\abs{W_K(\mathcal D)} = \log\abs{N_K(\mathcal D)} - \sum_{C \in \mathcal D}\log\abs{\kappa^{(q)}_K(C, \chi)}$, this is (i).

(ii) Each of the eleven diagrams is an ensemble at level $q$ of at most two rectangles with the representation $\pi$. A rectangle with sides $a \geq b$ has $s(C) = b$, $\delta(C) = b$, and $P(C) = 2(a + b)$, so that $\Lambda(C) \leq \max(2, 1 + a/b)$ by Lemma~\ref{lem:chord-arc}, $r(C) = \min(b/2, b, (a + b)/4) = b/2$, and $P(C)/r(C) = 4(1 + a/b)$; for the nine rectangles $a/b \leq \max(s/t, 4t/s)$, so that they are admissible and $P(C)/r(C) \leq 4\Lambda_*$, and $P_q(C) \leq 2(4t + s)L^q/u$. The left side of \eqref{eq:nonvanishing-length} is therefore at most $\ell_{\mathrm{w}} \coloneqq 8(4t + s)^2L^{2q}/u^2 + 8\Lambda_*$ for each diagram. Let $\bar g_1$ and $\ell$ be those of Proposition~\ref{prop:nonvanishing} at the level $q$, and let $\bar g_1' \leq \bar g_1$ be so small that $\ell(c_\pi, \bar g) \geq \ell_{\mathrm{w}}$ for $\bar g \leq \bar g_1'$, which is possible because $\ell(c_\pi, \bar g) \to \infty$ as $\bar g \to 0$. Proposition~\ref{prop:nonvanishing} then gives $q \geq q_\sharp(c_\pi, \bar g)$, the limits, and the lower bounds. The diagrams $I_a$ and $R^{(1)}_n$ are the images of $I_0$ and $R_n$ under the translations by $ate_\mu$ and $te_\mu$, which belong to $\mathfrak G_q$ because $t \in uL^{-q}\mathbb N$, so that Corollary~\ref{cor:symmetries}(ii) gives the equalities for $W_K$ and for $W$. By the lower bounds $N_K(\mathcal D) \neq 0$ for $K$ large, and (i) together with the continuity of $\log\abs{\cdot}$ on $\mathbb C \setminus \{ 0 \}$ gives the three limits. Finally $\mathfrak Y^{\mathrm{sym}}(W) = \mathfrak Y(W)$ by Lemma~\ref{lem:witness-counts}(iii) applied to $f = \log\abs{W}$.

(iii) By \eqref{eq:level-change} and Proposition~\ref{prop:normalisation-structure} at the level $q'$, $\log\abs{\kappa^{(q')}_K(C, \chi)}$ is again a linear function of the lengths and the corner numbers of $C$ with coefficients independent of $C$, so that (i) holds with $W^{(q')}_K$ in place of $W_K$. Since $W^{(q')}_K(\mathcal D) \to W^{(q')}(\mathcal D) \neq 0$, also $N_K(\mathcal D) \neq 0$ for $K$ large, and $\mathfrak Y(W^{(q')}) = \lim_K\mathfrak Y(N_K) = \mathfrak Y(W)$ by (ii).
\end{proof}

In other words, $\mathfrak Y(W)$ does not depend on the renormalisation prescription within the class of perimeter and corner normalisations, and in particular not on the level at which the diagrams are read; any change of the masses $m_{v,K}$, of the corner constants, or of a constant factor per loop leaves it unchanged.

\subsection{The Free Field}\label{sec:free-field}

The lower bound on the shape term rests on an explicit computation for the free field. At the stopping scale the field is close to a flat field, whose holonomies around the cycles of the torus need not be trivial, and in the components of the Lie algebra on which these holonomies act diagonally the free field acquires twisted periodicity conditions. We therefore compute the covariance of the fluxes for every twist, and we need a lower bound which is uniform in it.

Let $\eta > 0$ and $\mathsf T \in \eta\mathbb N$, let $\theta \in [0, 1)^4$ be a \emph{twist}, and let $\mathbb T^*_{\eta,\theta}$ be the set of the $(\mathsf T/\eta)^4$ momenta $p = (2\pi/\mathsf T)(n + \theta)$, $n \in \mathbb Z^4$, taken modulo $(2\pi/\eta)\mathbb Z^4$, with the momentum $p = 0$ removed if $\theta = 0$. Put $\hat p_\alpha \coloneqq (2/\eta)\sin(\eta p_\alpha/2)$ and $\hat p^2 \coloneqq \sum_\alpha\hat p_\alpha^2$, which is positive on $\mathbb T^*_{\eta,\theta}$, and define for $z \in \eta\mathbb Z^4$ the kernel
\begin{equation}\label{eq:flux-kernel}
    k_\theta(z) \coloneqq \frac{1}{\mathsf T^4}\sum_{p \in \mathbb T^*_{\eta,\theta}}e^{ip\cdot z}\,\frac{\hat p_\mu^2 + \hat p_\nu^2}{\hat p^2}\;.
\end{equation}
For a rectangle $\mathsf I$ of plaquettes of the lattice $\eta\mathbb Z^4$ in the plane of $e_\mu$ and $e_\nu$, which we identify with the set of the base points of its plaquettes, and for $\mathsf w \in \eta\mathbb Z$, put
\begin{equation}\label{eq:flux-covariance}
    c^\theta(\mathsf w) \coloneqq \eta^4\sum_{x \in \mathsf I}\,\sum_{y \in \mathsf I + \mathsf we_\mu}\operatorname{Re}k_\theta(x - y)\;.
\end{equation}
At $\theta = 0$ this is the covariance of the fluxes $\mathsf X(\mathsf I) = \oint_{\partial\mathsf I}\mathsf B$ and $\mathsf X(\mathsf I + \mathsf we_\mu)$ of the free field, the centred Gaussian $1$-form on the periodic lattice of spacing $\eta$ and period $\mathsf T$ with $\mathbb E[\mathsf B_\alpha(x)\mathsf B_\beta(y)] = \delta_{\alpha\beta}\mathsf T^{-4}\sum_{p}e^{ip\cdot(x-y)}/\hat p^2$: by the lattice Stokes formula $\mathsf X(\mathsf I) = \eta^2\sum_{x \in \mathsf I}\mathsf F(x)$ with the field strength $\mathsf F = \partial_\mu\mathsf B_\nu - \partial_\nu\mathsf B_\mu$, and $\mathbb E[\mathsf F(x)\mathsf F(y)] = k_0(x - y)$ because $\abs{e^{i\eta p_\alpha} - 1}^2 = \eta^2\hat p_\alpha^2$.

\begin{lemma}[The Free Field's Third Difference]\label{lem:free-field}
Let $\mathsf I$ have width $t$ along $e_\mu$ and height $s$ along $e_\nu$, with $t, s \in \eta\mathbb N$, $12t \leq \mathsf T$, $s \leq \mathsf T/2$, and $\eta \leq \mathsf T/4$. Then for every twist $\theta$,
\begin{equation}\label{eq:free-field}
    c^\theta(3t) - c^\theta(2t) \geq \frac{s^2t^3}{112\,\mathsf T^5}\;.
\end{equation}
\end{lemma}

\begin{proof}
We split the multiplier of \eqref{eq:flux-kernel} as $(\hat p_\mu^2 + \hat p_\nu^2)/\hat p^2 = 1 - \sum_{\alpha \notin \{ \mu, \nu \}}\hat p_\alpha^2/\hat p^2$ and treat the two parts separately.

\emph{The constant part.} For $z \in \eta\mathbb Z^4$ the sum of $e^{ip\cdot z}$ over all the $(\mathsf T/\eta)^4$ momenta $(2\pi/\mathsf T)(n + \theta)$ factorises over the coordinates, and in each coordinate it equals $(\mathsf T/\eta)e^{2\pi i\theta_\alpha z_\alpha/\mathsf T}$ if $z_\alpha \in \mathsf T\mathbb Z$ and $0$ otherwise. We take lifts in which $\mathsf I + \mathsf we_\mu$ is the translate of $\mathsf I$ by $\mathsf we_\mu$, with $t \leq \mathsf w \leq \mathsf T - t$. For $x \in \mathsf I$ and $y \in \mathsf I + \mathsf we_\mu$ we then have $0 < y_\mu - x_\mu < \mathsf w + t \leq \mathsf T$, so that $x - y \notin \mathsf T\mathbb Z^4$, and the constant part contributes $0$ to $k_\theta(x - y)$ if $\theta \neq 0$, and $-\mathsf T^{-4}$, the removed momentum, if $\theta = 0$. Either way its contribution to $c^\theta(\mathsf w)$ does not depend on $\mathsf w$ in this range, which contains $2t$ and $3t$ because $12t \leq \mathsf T$.

\emph{The resolvent in one direction.} In the second part we fix $p_\perp \coloneqq (p_\alpha)_{\alpha \neq \mu}$ with $\sum_{\alpha \notin \{ \mu, \nu \}}\hat p_\alpha^2 > 0$, the other terms vanishing, put $\mathsf m^2 \coloneqq \sum_{\alpha \neq \mu}\hat p_\alpha^2 > 0$, and let $\vartheta > 0$ with $\cosh\vartheta = 1 + \eta^2\mathsf m^2/2$. The sum over $p_\mu$ is then
\begin{equation}\label{eq:one-direction-resolvent}
    \mathsf h(z) \coloneqq \frac{1}{\mathsf T}\sum_{p_\mu}\frac{e^{ip_\mu z}}{\hat p_\mu^2 + \mathsf m^2} = \frac{\eta}{2\sinh\vartheta}\sum_{n \in \mathbb Z}e^{2\pi i\theta_\mu n}e^{-\vartheta\abs{z - n\mathsf T}/\eta}\;,
    \qquad
    z \in \eta\mathbb Z\;,
\end{equation}
the sum over $p_\mu \in (2\pi/\mathsf T)(\mathbb Z + \theta_\mu)$ modulo $2\pi/\eta$. Indeed, both sides satisfy $h(z + \mathsf T) = e^{2\pi i\theta_\mu}h(z)$ and solve $\eta^{-2}(2h(z) - h(z + \eta) - h(z - \eta)) + \mathsf m^2h(z) = \eta^{-1}\sum_ne^{2\pi i\theta_\mu n}\delta_{z,n\mathsf T}$, whose solution in the space of such $h$ is unique because the operator has the eigenvalues $\hat p_\mu^2 + \mathsf m^2 > 0$ there. For the left side this is the sum of $e^{ip_\mu z}$ over the momenta, computed above in one coordinate. For the right side, $g(z) \coloneqq (\eta/(2\sinh\vartheta))e^{-\vartheta\abs{z}/\eta}$ solves the equation with the source $\eta^{-1}\delta_{z,0}$: at $z \neq 0$ because $2\cosh\vartheta = 2 + \eta^2\mathsf m^2$, and at $z = 0$ because $2 - 2e^{-\vartheta} + \eta^2\mathsf m^2 = 2\cosh\vartheta - 2e^{-\vartheta} = 2\sinh\vartheta$; the sum $\sum_ne^{2\pi i\theta_\mu n}g(z - n\mathsf T)$ converges absolutely and has the periodicity and the sources. Since $g$ is real and even, $\mathsf h(-z) = \overline{\mathsf h(z)}$, and $\operatorname{Re}\mathsf h$ is even.

\emph{Monotonicity.} For $0 \leq z \leq \mathsf T$ put $y \coloneqq \vartheta z/\eta$, $\mathsf N \coloneqq \vartheta\mathsf T/\eta$, $\mathsf q \coloneqq e^{-\mathsf N}$, and $\psi \coloneqq 2\pi\theta_\mu$. Summing the two geometric series over $n \leq 0$ and $n \geq 1$,
\begin{equation}\label{eq:resolvent-real-part}
    \frac{2\sinh\vartheta}{\eta}\operatorname{Re}\mathsf h(z) = \frac{S(y)}{\abs{1 - \mathsf qe^{i\psi}}^2}\;,
    \qquad
    S(y) \coloneqq e^{-y}(1 - \mathsf q\cos\psi) + \mathsf qe^{y}(\cos\psi - \mathsf q)\;.
\end{equation}
We claim that $-S'(y) \geq (1 - \mathsf q)e^{-y}(1 - e^{2y - \mathsf N})$ for $0 \leq y \leq \mathsf N/2$. Indeed $-S'(y) = e^{-y}(1 - \mathsf q\cos\psi) - \mathsf qe^{y}(\cos\psi - \mathsf q)$. If $\cos\psi \leq \mathsf q$, the second term is nonnegative and the claim follows from $1 - \mathsf q\cos\psi \geq 1 - \mathsf q$. Otherwise $\cos\psi - \mathsf q \leq 1 - \mathsf q\cos\psi$, the difference being $-(1 - \cos\psi)(1 + \mathsf q)$, and $\mathsf qe^y = e^{2y - \mathsf N}e^{-y}$, so that $-S'(y) \geq e^{-y}(1 - \mathsf q\cos\psi)(1 - e^{2y - \mathsf N})$, which gives the claim. Since $\abs{1 - \mathsf qe^{i\psi}}^2 \leq 4$, for $0 \leq z < z' < \mathsf T/2$,
\begin{equation}\label{eq:resolvent-decrease}
    \operatorname{Re}\mathsf h(z) - \operatorname{Re}\mathsf h(z') \geq \frac{\eta(1 - \mathsf q)}{8\sinh\vartheta}\int_{\vartheta z/\eta}^{\vartheta z'/\eta}e^{-y}\bigl(1 - e^{2y - \mathsf N}\bigr)\,\dd y > 0\;,
\end{equation}
so that $\operatorname{Re}\mathsf h$ is strictly decreasing on $[0, \mathsf T/2)$, for every twist.

\emph{The difference.} For $x \in \mathsf I$ and $y \in \mathsf I + \mathsf we_\mu$ the coordinates other than $x_\mu$ and $x_\nu$ agree, and the two rectangles have the same rows. The second part of $c^\theta(\mathsf w)$ is therefore $-\mathsf T^{-3}\sum_{p_\perp}\bigl(\sum_{\alpha \notin \{ \mu, \nu \}}\hat p_\alpha^2\bigr)\abs{\eta\sum_{x_\nu}e^{ip_\nu x_\nu}}^2\eta^2\sum_{x_\mu,y_\mu}\operatorname{Re}\mathsf h(x_\mu - y_\mu - \mathsf w)$, with $x_\nu$ over the rows and $x_\mu$, $y_\mu$ over the columns of $\mathsf I$, and by the first step
\begin{equation}\label{eq:third-difference-sum}
    c^\theta(3t) - c^\theta(2t) = \frac{1}{\mathsf T^3}\sum_{p_\perp}\Bigl(\sum_{\alpha \notin \{ \mu, \nu \}}\hat p_\alpha^2\Bigr)\Bigl\lvert\eta\sum_{x_\nu}e^{ip_\nu x_\nu}\Bigr\rvert^2\eta^2\sum_{x_\mu,y_\mu}\bigl[\operatorname{Re}\mathsf h(z + 2t) - \operatorname{Re}\mathsf h(z + 3t)\bigr]\;,
\end{equation}
with $z \coloneqq y_\mu - x_\mu \in (-t, t)$, where we used that $\operatorname{Re}\mathsf h$ is even. Since $t < z + 2t < z + 3t < 4t \leq \mathsf T/3$, every bracket is positive by \eqref{eq:resolvent-decrease}, and every term of \eqref{eq:third-difference-sum} is nonnegative. We keep one of them. Let $\alpha$ and $\alpha'$ be the directions other than $\mu$ and $\nu$, and choose $p_\nu$, $p_\alpha$, $p_{\alpha'}$ in their classes with $\abs{p_\nu}, \abs{p_{\alpha'}} \leq \pi/\mathsf T$ and $\pi/\mathsf T \leq \abs{p_\alpha} \leq 2\pi/\mathsf T$, which is possible because every class $(2\pi/\mathsf T)(\mathbb Z + \theta_\beta)$ meets both $[-\pi/\mathsf T, \pi/\mathsf T]$ and $\{ \pi/\mathsf T \leq \abs{p} \leq 2\pi/\mathsf T \}$. Since $\eta \leq \mathsf T/4$, these momenta have $\abs{\eta p/2} \leq \pi/4$, and $\sin y \geq (2\sqrt2/\pi)y$ on $[0, \pi/4]$ gives $\hat p_\alpha^2 \geq 8/\mathsf T^2$, $\mathsf m^2 \leq 6\pi^2/\mathsf T^2$ because $\hat p_\beta \leq \abs{p_\beta}$, and, with $s \leq \mathsf T/2$,
\begin{equation}\label{eq:row-sum}
    \Bigl\lvert\eta\sum_{x_\nu}e^{ip_\nu x_\nu}\Bigr\rvert = \eta\,\Bigl\lvert\frac{\sin(sp_\nu/2)}{\sin(\eta p_\nu/2)}\Bigr\rvert \geq \frac{2\sqrt2}{\pi}\,s\;,
\end{equation}
the left side being $s$ if $p_\nu = 0$.

\emph{The bracket.} Now $\eta\mathsf m \leq \sqrt6\pi/4 < 2$, so that $\vartheta = 2\operatorname{arsinh}(\eta\mathsf m/2) \in [\eta\mathsf m/2, \eta\mathsf m]$, because $y/2 \leq \operatorname{arsinh}y \leq y$ on $[0, 1]$, and $\sinh\vartheta = \eta\mathsf m(1 + \eta^2\mathsf m^2/4)^{1/2} \leq \sqrt2\,\eta\mathsf m$. In \eqref{eq:resolvent-decrease} with $z + 2t$ and $z + 3t$ the interval of integration starts at most at $3\vartheta t/\eta \leq 3\mathsf mt$ and has length $\vartheta t/\eta \geq \mathsf mt/2$; on it $2y - \mathsf N \leq (\vartheta/\eta)(8t - \mathsf T) \leq -\vartheta\mathsf T/(3\eta) \leq -\mathsf m\mathsf T/6$; and $1 - \mathsf q \geq 1 - e^{-\mathsf m\mathsf T/2}$. With $1 - e^{-x} \geq x(1 - x/2)$ for $x \geq 0$, the bracket is at least
\begin{equation}\label{eq:bracket-lower}
    \frac{1}{8\sqrt2\,\mathsf m}\bigl(1 - e^{-\mathsf m\mathsf T/2}\bigr)\bigl(1 - e^{-\mathsf m\mathsf T/6}\bigr)e^{-3\mathsf mt}\,\frac{\mathsf mt}{2}\Bigl(1 - \frac{\mathsf mt}{4}\Bigr) \geq \frac{t}{512\sqrt2} \geq \frac{t}{725}\;,
\end{equation}
because $2\sqrt2 \leq \mathsf m\mathsf T \leq \sqrt6\pi$ and $\mathsf mt \leq \sqrt6\pi/12$ give $e^{-\mathsf m\mathsf T/2} \leq e^{-\sqrt2} \leq 1/4$, $e^{-\mathsf m\mathsf T/6} \leq e^{-\sqrt2/3} \leq 5/8$, $e^{-3\mathsf mt} \geq e^{-2} \geq 2/15$, and $1 - \mathsf mt/4 \geq 5/6$, and $(3/4)(3/8)(2/15)(5/6) = 1/32$. The kept term of \eqref{eq:third-difference-sum} is therefore at least $\mathsf T^{-3}(8/\mathsf T^2)(8s^2/\pi^2)\,t^2\,t/725 = 64s^2t^3/(725\pi^2\mathsf T^5)$, since the sum over $x_\mu$ and $y_\mu$ has $(t/\eta)^2$ terms, and $64/(725\pi^2) \geq 1/112$.
\end{proof}

For the free field with the charge $\mathsf g$ the fluxes $\mathsf X_a \coloneqq \oint_{I_a}\mathsf B$ form a Gaussian charge model whose covariance is invariant under the translation by $te_\mu$, with $c^{\mathrm G}_k = c^0(kt)$. By \eqref{eq:stationary-gaussian} its number $\mathfrak Y$ vanishes, while its shape term is $\mathsf g^2(c^0(3t) - c^0(2t)) > 0$ by the lemma: in the abelian theory the positive shape term is cancelled exactly by the pair terms.

\subsection{The Pair Terms}\label{sec:pair-terms}

We now show that the pair terms $\mathfrak y_k(W) = \log\abs{W(P_k)/(W(I_0)W(I_k))}$ are of the fourth order in the coupling, up to logarithms. Reflection positivity gives only their sign, and only for some $t$: if $(k + 1)t/2 \in uL^{-q}\mathbb Z$, the reflection $\theta \coloneqq \theta_{\mu,o_\mu+(k+1)t/2}$ is of level $q$ and maps $I_0$ onto $I_k$ with the opposite orientation, so that Corollaries~\ref{cor:symmetries}(iii) and~\ref{cor:reflection-positivity}(ii) give $W(P_k) = W(I_0 \cup \theta\bar I_0) \geq \abs{W(I_0)}^2$ and $\mathfrak y_k(W) \geq 0$. An upper bound needs more, namely that the two rectangles are nearly independent under the measure of the stopping scale.

The argument has two parts. At the stopping step the normalised dressed loop of $P_k$ on the term $\ell_0$ is, up to terms which decay super-polynomially in $\bar g$, the product of those of $I_0$ and $I_k$, because the dressed loop of an ensemble is built from objects which see only the loops that carry them, and the objects which reach from one rectangle to the other are paid by their size; this is Lemma~\ref{lem:pair-stopping}. The pair term is then governed by the covariance, under the probability measure $\langle\cdot\rangle_{\ell_0}$, of the deviations from one of the two normalised dressed loops. Each deviation consists of the deviation of a character of the background holonomy, which is quadratic in the distance of the holonomy from $1$ by Lemma~\ref{lem:near-one} below, and of the insertions, which are of the order $\bar g^2$; its $L^2$-norm is therefore of the order $\bar g^2$, and the covariance of the order $\bar g^4$ (Lemma~\ref{lem:stopping-sizes} and Theorem~\ref{thm:pair-terms}).

\begin{lemma}[Characters near One]\label{lem:near-one}
Let $\pi$ be an irreducible unitary representation of $G$ of dimension $n_\pi$ with the character $\chi$. Then $\operatorname{Tr}\pi_*(Y) = 0$ for every $Y \in \mathfrak g$ and $\det\pi(U) = 1$ for every $U \in G$. If $U \in G$ and $h \coloneqq \norm{\pi(U) - 1} \leq 2/n_\pi$, then
\begin{equation}\label{eq:near-one}
    \abs{\chi(U) - 1} \leq \frac35h^2\;.
\end{equation}
\end{lemma}

\begin{proof}
The differential $\dd\pi$ is a homomorphism of Lie algebras from $\operatorname{Lie}G$ into the Lie algebra of $U(n_\pi)$, and the trace is a homomorphism from the latter into the abelian Lie algebra $i\mathbb R$; so $\operatorname{Tr}\circ\dd\pi$ vanishes on the commutators, which span $\operatorname{Lie}G$ because $G$ is simple. Since $\pi_*(Y) = -i\,\dd\pi(iY)$ (Section~\ref{sec:coarse-loop}), $\operatorname{Tr}\pi_*(Y) = 0$, and $\det\pi(e^{iY}) = e^{i\operatorname{Tr}\pi_*(Y)} = 1$. The set of the $U \in G$ with $\det\pi(U) = 1$ is therefore a subgroup which contains the image of the exponential map, a neighbourhood of $1$, and since a connected group is generated by every neighbourhood of $1$, it is $G$.

If $n_\pi = 1$, then $\pi = \det\pi = 1$ and $\chi = 1$. Let $n_\pi \geq 2$, so that $h \leq 1$, and let $e^{i\omega_j}$, $\omega_j \in (-\pi, \pi]$, be the eigenvalues of the unitary matrix $\pi(U)$. Each satisfies $2\abs{\sin(\omega_j/2)} = \abs{e^{i\omega_j} - 1} \leq h$, and since $\arcsin$ is convex on $[0, \frac12]$ with $\arcsin\frac12 < 0.524$, we have $\arcsin y \leq 1.05\,y$ there, so that $\abs{\omega_j} = 2\arcsin(\abs{e^{i\omega_j} - 1}/2) \leq 1.05\,h$. Hence $\abs{\sum_j\omega_j} \leq 1.05\,n_\pi h \leq 2.1$, while $e^{i\sum_j\omega_j} = \det\pi(U) = 1$, so that $\sum_j\omega_j = 0$. Therefore
\begin{equation}\label{eq:near-one-sum}
    \chi(U) - 1 = \frac{1}{n_\pi}\sum_{j=1}^{n_\pi}\bigl(e^{i\omega_j} - 1 - i\omega_j\bigr)\;,
\end{equation}
and $\abs{e^{i\omega} - 1 - i\omega} \leq \omega^2/2$ gives $\abs{\chi(U) - 1} \leq (1.05\,h)^2/2 \leq \frac35h^2$.
\end{proof}

\emph{The objects at the stopping step.} Let $\mathcal D$ be one of the eleven diagrams. By the proof of Proposition~\ref{prop:normalisations-drop-out}(ii) its loops are admissible, and for $\bar g \leq \bar g_1'$ we have $q \geq q_\sharp(c_\pi, \bar g)$, so that the results of Section~\ref{sec:stopping-scale} apply to it. The index $K_0(\mathcal D)$ of \eqref{eq:index-k-zero} depends on $\mathcal D$ only through the largest $P_q$ of its loops, and we put $K_{\mathrm w} \coloneqq \max(K_0(R_4), j_q + 2)$, which is at least $K_0(\mathcal D)$ for the eleven diagrams. Let $K \geq K_{\mathrm w}$, so that the stopping step $K' = K - j_q \geq 2$ is common to the eleven diagrams. The density $\varrho_{K'}$, its labelled terms, the weights $w_\ell$, and the expectations $\langle\cdot\rangle_\ell$ of Theorem~\ref{thm:stopping-identity} do not involve the ensemble, and we write $w \coloneqq w_{\ell_0}$ and $\langle\cdot\rangle \coloneqq \langle\cdot\rangle_{\ell_0}$, a probability measure on $G^{\mathcal B}$. On $\ell_0$ we write $E(\mathcal D) \coloneqq E_{\mathrm{const},K'}$ for the constant of the dressed loop of $\mathcal D$, and
\begin{equation}\label{eq:stopping-bracket}
    S(\mathcal D) \coloneqq \prod_{C \in \mathcal D}f^{(C)}_{K'} + \sum_{\mathfrak B \neq \emptyset}Q_{K'}(\mathfrak B)\,w_{K'}(\mathfrak B)
\end{equation}
for its bracket, a function of $V$, and we put $\Lambda(\mathcal D) \coloneqq \sum_{\ell \neq \ell_0 \,|\, Z_\ell > 0}w_\ell\langle\widehat{\mathcal W}_\ell\rangle_\ell$, computed with the dressed loops of $\mathcal D$. Then \eqref{eq:stopping-identity} reads
\begin{equation}\label{eq:pair-stopping-form}
    W_K(\mathcal D) = w\,e^{E(\mathcal D)}\bigl\langle S(\mathcal D)\bigr\rangle + \Lambda(\mathcal D)\;.
\end{equation}
For a single rectangle $C$ we write $\xi_C \coloneqq S(C) - 1$ and $I(C) \coloneqq S(C) - f^{(C)}_{K'}$, so that $\abs{I(C)} \leq e^{I_*(C)} - 1$ by Lemma~\ref{lem:term-bound}. Finally let
\begin{equation}\label{eq:typical-threshold}
    \bar C \coloneqq C_{\mathrm{fs}}N_q^4\Bigl(1 + \frac12\log\hat g_{K'}^{-2}\Bigr)\;,
    \qquad
    T_{\mathrm{ng}} \coloneqq 2\bar C + 4\log\bar g^{-2}\;,
\end{equation}
which do not depend on $K$, and, for each of the nine rectangles $C$,
\begin{equation}\label{eq:pair-sizes}
    \bar h(C) \coloneqq 4.8\,c_\pi c_hB_3P_{K'}(C)^2(2NT_{\mathrm{ng}})^{1/2}(1 + \beta_0)\bar g\;,
    \qquad
    \sigma(C) \coloneqq \frac35\bar h(C)^2 + e^{I_*(C)} - 1 + 3\bar g^4\;.
\end{equation}

\begin{lemma}[Sizes at the Stopping Step]\label{lem:stopping-sizes}
For $\bar g$ sufficiently small, depending on $q$, $t/u$, $s/u$, $T/u$, $\pi$, and the constants of the construction, the following hold for every $K \geq K_{\mathrm w}$ and every one of the eleven diagrams $\mathcal D$.
\begin{enumerate}[label=(\roman*)]
\item $\abs{E(\mathcal D)} \leq \bar g^{10}$, $1 - w \leq \bar g^{10}$, $\abs{\Lambda(\mathcal D)} \leq \bar g^{10}$, and $I_*(\mathcal D) \leq \bar g$.
\item The event $\mathcal V \coloneqq \{ V \in \operatorname{supp}\varrho_{\ell_0} \,|\, A(U_{K'}(V)) \leq T_{\mathrm{ng}}g_{K'}^2 \}$, with $A$ the Wilson action, has $\langle\mathbf 1_{\mathcal V^c}\rangle \leq \bar g^8$, and on $\mathcal V$, for each of the nine rectangles $C$, $\norm{\pi(U_{K'}(V)(C)) - 1} \leq \bar h(C)$ and $\abs{f^{(C)}_{K'} - 1} \leq \frac35\bar h(C)^2$.
\item For each of the nine rectangles $C$, $\abs{\xi_C} \leq 3$ on the support of $\varrho_{\ell_0}$ and $\langle\abs{\xi_C}^2\rangle^{1/2} \leq \sigma(C)$.
\item There are $C_\sigma$ and $m_\sigma$, depending on the same quantities, with $\sigma(C) \leq C_\sigma\bar g^2(\log\bar g^{-2})^{m_\sigma}$ for each of the nine rectangles.
\end{enumerate}
\end{lemma}

\begin{proof}
We first collect the sizes of the parameters of the stopping step. By \eqref{eq:tube-couplings}, $\hat g_{K'}^{-2} = \bar g^{-2} + b'j_q + C_W$, and $j_q$ grows like $\log\log\bar g^{-2}$ by \eqref{eq:stopping-identity-index}, so that $\hat g_{K'} \leq \bar g$ and $\log\bar g^{-2} \leq \log\hat g_{K'}^{-2} \leq 2\log\bar g^{-2}$ for $\bar g$ small, as in the proof of Proposition~\ref{prop:nonvanishing}. Hence $(\log\bar g^{-2})^r \leq R_{K'} \leq R_{K'-1} \leq 2L(\log\hat g_{K'-1}^{-2})^r \leq C(\log\bar g^{-2})^r$ by Lemma~\ref{lem:tube-cubes}, with constants $C$ of the kind of the statement here and below. Consequently $P_{K'}(C) = P_q(C)L^3MR_{K'-1} \leq C(\log\bar g^{-2})^r$ for each of the nine rectangles, $N_q = 2L^{m_T+q}L^3MR_{K'-1} \leq C(\log\bar g^{-2})^r$, and $\bar C + T_{\mathrm{ng}} \leq C(\log\bar g^{-2})^{4r+1}$. In units of the step $K'$ the corners of the rectangles lie in $a_{K'}\mathbb Z^4$ with $a_{K'} = L^{j_q-q} \geq L^3MR_{K'}$ by Lemma~\ref{lem:stopping-index}, and their shortest segments, the distances between their opposite sides, and the distance $(k - 1)t/(uL^{-j_q})$ of $I_0$ from $I_k$ are at least $a_{K'}$.

(i) By Definition~\ref{def:dressed-loop}(iv), $\abs{E(\mathcal D)} \leq \sum_{j \leq K'}\epsilon_j(\mathcal D, T)$, and we sum the parts of \eqref{eq:step-remainder} as in the proof of Corollary~\ref{cor:remainder-sum}, keeping the exponentials. For a loop $C$ of $\mathcal D$ put $x \coloneqq Lr_{K'}(C) \geq La_{K'}/2 \geq L^4MR_{K'}/2$. The corner part is at most $C^{\mathcal K}_E\bar{\mathfrak a}(P(C)/r(C))(x + 1)(1 + 1/(c^T_Ex))e^{-2c^T_Ex}$ by Lemma~\ref{lem:scale-sums} with $c = 2c^T_E$ and $\abs{\mathcal K(C)} \leq P(C)/(2r(C))$, and $2c^T_Ex \geq \nu L^3R_{K'}/(64A^*_{\mathrm{env}})$. The torus part is at most $9\bar C_Z\bar{\mathfrak a}P_q(C)L\iota_{T,K'}e^{-c^Z\iota_{T,K'}}$, by the proof of Corollary~\ref{cor:remainder-sum}, and since $y \mapsto ye^{-c^Zy}$ decreases for $c^Zy \geq 1$ and $\iota_{T,K'} \geq 96L^3MR_{K'}$ by \eqref{eq:torus-margin-bound}, it is at most $9\bar C_Z\bar{\mathfrak a}P_q(C)L \cdot 96L^3MR_{K'}e^{-3\nu L^3R_{K'}/(8A^*_{\mathrm{env}})}$. The cross part is at most $3C_{\mathrm{cr}}\bar{\mathfrak a}P_q(C)L \cdot L^3MR_{K'}e^{-\nu_\times L^3R_{K'}}$ in the same way, with $a_{K'} \geq L^3MR_{K'}$. Each of these three is at most a constant times $(\log\bar g^{-2})^re^{-c_\flat(\log\bar g^{-2})^r}$ with $c_\flat \coloneqq \min(\nu/(64A^*_{\mathrm{env}}), \nu_\times)L^3$, and since $r = 2$ this is at most $\bar g^{11}$ for $\bar g$ small. The $\mathbf R$-part is at most $C^R_E(P(C)/r(C) + P_q(C))(\frac12 + 2/(ec^R_E))S^{\mathbf R}(\bar g)$ by the proof of Corollary~\ref{cor:remainder-sum}, and the cross part of the $\mathbf R$-steps, present for the pairs, is at most $6C^{\mathbf R\times}_EP_q(\mathcal D)S^{\mathbf R}(\bar g)/(e\kappa)$ by the same proof. To see that $S^{\mathbf R}(\bar g)$ decays faster than every power of $\bar g$, put $y_n \coloneqq \frac23(\bar g^{-2} + b'n)$ and $a \coloneqq (A_0/4)(\log y_0)^{p_0-1}$, which tends to infinity as $\bar g \to 0$ because $p_0 > 1$. The term $j = K - n$ of $S^{\mathbf R}(\bar g)$ is at most $4L^2M^2(\log g_j^{-2})^{2r}\exp(-(A_0/4)(\log y_n)^{p_0})$, by that proof, and $(\log y_n)^{p_0} \geq (\log y_0)^{p_0-1}\log y_n$, while $g_j^{-2} \leq \hat g_j^{-2} \leq 3y_n$ by Lemma~\ref{lem:tube-couplings}(i) and \eqref{eq:tube-couplings}; for $\bar g$ small, $(\log3y)^{2r} \leq y$ for $y \geq y_0$ and $a \geq 3$, so that
\begin{equation}\label{eq:r-sum-decay}
    S^{\mathbf R}(\bar g) \leq 4L^2M^2\sum_{n \geq 0}y_n^{1-a} \leq 4L^2M^2y_0^{1-a}\Bigl(1 + \frac{3y_0}{2b'(a - 2)}\Bigr) \leq 8L^2M^2y_0^{2-a}\;,
\end{equation}
the sum being bounded by its first term and the integral $\int_0^\infty y_u^{1-a}\,\dd u$, which is at most $\bar g^{11}$ for $\bar g$ small. Adding the parts over the at most two loops of $\mathcal D$ gives $\abs{E(\mathcal D)} \leq \bar g^{10}$. For the large-field terms, $1 - w \leq \bar\delta_{\mathrm{LF}}$ by Proposition~\ref{prop:large-field-weight}, and $\abs{\Lambda(\mathcal D)} \leq e^{B_E(\mathcal D) + I_*(\mathcal D)}\bar\delta_{\mathrm{LF}}$ by the proof of Theorem~\ref{thm:uniform-bound}; here $\bar\delta_{\mathrm{LF}} \leq e^{2\bar C}(e^{\bar s} - 1)$ with $\bar s \coloneqq 2N_q^4\exp(-c_{\mathrm{LF}}A_0(\log((1 + \beta_0)^{-2}\bar g^{-2}))^{p_0})$, whose logarithm is at most $C(\log\bar g^{-2})^{4r+1} - c(\log\bar g^{-2})^{p_0}$ with $p_0 \geq 5r > 4r + 1$, so that $\bar\delta_{\mathrm{LF}} \leq \bar g^{11}$ for $\bar g$ small. Finally $B_E(\mathcal D) \leq 1$ and $I_*(\mathcal D) \leq \bar g$ for $\bar g$ small, by \eqref{eq:nonvanishing-quantities} and by (iv) below, whose proof does not use (i).

(ii) We first bound the probability of $\mathcal V^c$. In the proof of Proposition~\ref{prop:large-field-weight}, every term is divided by the common factor $e^{-E^{\ell_0}}$, after which the density of $\ell_0$ is at most $\exp(-g_{K'}^{-2}A(U_{K'}(V)) + C_{\mathcal E}N_q^4)$, since $\chi_{K'}(\mathbb T) \leq 1$ and the regular part, the $\mathbf R$-part, and the boundary part are bounded by $C_{\mathcal E}\sum_n\abs{\Gamma_n} = C_{\mathcal E}N_q^4$ on $\ell_0$, and its integral is at least $e^{-C_K}$ by \eqref{eq:trivial-term-lower}. Since $C_{\mathcal E}N_q^4 \leq C_K \leq \bar C$, because $C_{\mathrm{fs}} \geq C_{\mathcal E}$ and $\log g_{K'}^{-1} \leq \frac12\log\hat g_{K'}^{-2}$, the normalised Haar measure gives
\begin{equation}\label{eq:typical-event}
    \langle\mathbf 1_{\mathcal V^c}\rangle \leq e^{C_K}\int_{\mathcal V^c}e^{-g_{K'}^{-2}A(U_{K'}(V)) + C_{\mathcal E}N_q^4}\,\dd V \leq e^{2\bar C - T_{\mathrm{ng}}} = \bar g^8\;.
\end{equation}
On $\mathcal V$ we apply Supplement, Lemma~17.6, with $T = T_{\mathrm{ng}}$ and $\varepsilon^{\mathrm{act}} \coloneqq 3(2NT_{\mathrm{ng}})^{1/2}g_{K'}$. Its conditions hold for $\bar g$ small: $\varepsilon_{K'}$ is small; $2B_3\varepsilon_{K'} \leq a_0$; $\varepsilon^{\mathrm{act}} \leq a_1$; and $\varepsilon^{\mathrm{act}} \leq 2\varepsilon_{K'} = 2\hat g_{K'}p_0(\hat g_{K'})$, because $g_{K'} \leq \rho^{1/2}\hat g_{K'}$ by Lemma~\ref{lem:tube-couplings}(i), so that it suffices that $9 \cdot 2NT_{\mathrm{ng}}\rho \leq 4A_0^2(\log\hat g_{K'}^{-2})^{2p_0}$, whose left side is at most $C(\log\bar g^{-2})^{4r+1}$ while $2p_0 > 4r + 1$. The lemma gives $\abs{U(\partial p) - 1} < B_3\varepsilon^{\mathrm{act}}L^{-2K'}$ on the fine plaquettes, $U \coloneqq U_{K'}(V)$, and hence $\norm{U(\partial p) - 1} \leq t \coloneqq c_hB_3\varepsilon^{\mathrm{act}}L^{-2K'}$. In fine units a rectangle $C$ of the nine is a closed path of $P \coloneqq P_{K'}(C)L^{K'}$ bonds in a box of side at most $P/2$, which embeds in the torus because $3\diam(C) < T$. Hence $3(P/2)Pt \leq \frac32c_hB_3P_{K'}(C)^2\varepsilon^{\mathrm{act}}$, which is at most $\min(1/10, \log(1 + r_G))$ for $\bar g$ small, and Lemma~\ref{lem:flat-holonomy} gives $\norm{\pi(U(C)) - 1} \leq 1.06 \cdot \frac32c_\pi c_hB_3P_{K'}(C)^2\varepsilon^{\mathrm{act}} \leq \bar h(C)$, using $g_{K'} \leq (1 + \beta_0)\bar g$. For $\bar g$ small, $\bar h(C) \leq 2/n_\pi$, and Lemma~\ref{lem:near-one} gives $\abs{f^{(C)}_{K'} - 1} = \abs{\chi(U(C)) - 1} \leq \frac35\bar h(C)^2$.

(iii) On the support of $\varrho_{\ell_0}$, $\abs{\xi_C} \leq \abs{f^{(C)}_{K'} - 1} + \abs{I(C)} \leq 2 + e^{I_*(C)} - 1 \leq 3$ by (i), and on $\mathcal V$, $\abs{\xi_C} \leq \frac35\bar h(C)^2 + e^{I_*(C)} - 1$ by (ii). Splitting $\xi_C = \xi_C\mathbf 1_{\mathcal V} + \xi_C\mathbf 1_{\mathcal V^c}$, Minkowski's inequality and \eqref{eq:typical-event} give $\langle\abs{\xi_C}^2\rangle^{1/2} \leq \frac35\bar h(C)^2 + e^{I_*(C)} - 1 + 3\langle\mathbf 1_{\mathcal V^c}\rangle^{1/2} \leq \sigma(C)$.

(iv) By the sizes collected above, $\bar h(C)^2 \leq C\bar g^2P_{K'}(C)^4T_{\mathrm{ng}} \leq C\bar g^2(\log\bar g^{-2})^{8r+1}$. In \eqref{eq:insertion-total}, $\bar\rho \leq C\bar g^2(\log\bar g^{-2})^{4p_0}$ by \eqref{eq:sup-amplitudes}, $\varepsilon_{K'}/\min(\alpha_{0,K'}, \alpha_{1,K'})$ is a power of $\log\hat g_{K'}^{-2}$ by Definition~\ref{def:construction}, $S_{\mathrm{wr}}(b) \leq 3$ for $b \geq 1$ by Lemma~\ref{lem:scale-sums}, and $P_{K'}(C) \leq C(\log\bar g^{-2})^r$, so that $I_*(C) \leq C\bar g^2(\log\bar g^{-2})^{m'}$ with an exponent $m'$ of the constants of stage $0$, and $e^{I_*(C)} - 1 \leq 2I_*(C)$. With $3\bar g^4 \leq 3\bar g^2$ this gives (iv) with $m_\sigma \coloneqq \max(8r + 1, m')$.
\end{proof}

We next make precise the statement that a dressed loop sees only the loops which carry its objects. The \emph{network} of a group is the union of its domains, of all scales, its intervals, and the networks of the clusters merged into it; it is connected, contains the networks of its old members, and has diameter at most $M\sigma(W)$ in units of the step at which it was made, by the proof of Lemma~\ref{lem:root-bound}. For an ensemble $\mathcal C$ with a loop $C$ of the character $\chi_C$ we write $\mathcal C_C \coloneqq ((C, \chi_C))$.

\begin{lemma}[Locality between Ensembles]\label{lem:ensemble-locality}
Let $\mathcal C$ be an ensemble as in Theorem~\ref{thm:main}, $C$ one of its loops, and $K \geq K_0(\mathcal C)$. On the trivial term, at every step $1 \leq k \leq K'$, the following hold.
\begin{enumerate}[label=(\roman*)]
\item Every group of $\mathcal C$ which is not a cross group and whose intervals lie on $C$ is a group of $\mathcal C_C$, with the same domains, intervals, network, and insertion, the insertion being the same function of the block field. Two such groups are compatible for $\mathcal C$ if and only if they are compatible for $\mathcal C_C$.
\item Every group of $\mathcal C_C$ whose network is at distance greater than $2$ from the other loops of $\mathcal C$, in units of the step $k$, is a group of $\mathcal C$ of the kind (i).
\item A cluster of the reference gas of $\mathcal C$ of the step $k$ all of whose groups are of the kind (i) is a cluster of the reference gas of $\mathcal C_C$ with the same weight, and a cluster of $\mathcal C_C$ all of whose groups satisfy the hypothesis of (ii) is a cluster of $\mathcal C$ with the same weight.
\end{enumerate}
\end{lemma}

\begin{proof}
A group, with its insertion, is determined by its constituents. At the step $1$ these are its roots and its cluster occurrences, and a group is the only group of the term formed by them \cite[Definition~10.1, Lemma~10.2(a)]{WilsonFirstStep}. At a step $k + 1 \geq 2$ they are the polymers of the step merged into it, its old members, which are groups of the step $k$, the slots on its arcs, and the clusters of the reference gas of the step $k$ merged into it (Section~\ref{sec:small-field-map} and the merging rule of Section~\ref{sec:reference-gas}), and at an $\mathbf R$-step also the polymers of the $\mathbf R$-operation linked to it (Supplement, Proposition~S.10). Its domains are those of its constituents, its intervals are the smallest arcs of its loops which contain the intervals of its members and the traversals of its slots and merged clusters, and its insertion is the Gaussian integral of the step of the product of the activities of its polymers, the insertions of its old members, its slots, and the weights of its merged clusters, less its flat part \eqref{eq:multilinear-split}. Compatibility is a relation between the domains, the arcs, and the merged clusters of two groups, and a cluster of the reference gas is a multiset of groups whose weight is the Ursell function of its incompatibility graph times the flat values of its groups. None of these rules involves a loop which carries no interval of the group: the polymers and their activities are those of Ba{\l}aban's expansion, which does not involve the ensemble; the slots and the transports are attached to the traversals of the arcs, along the same background field; and the smallest arc of $C$ which contains a set of traversals of $C$ is determined by $C$. We argue by induction on $k$.

\emph{The step $1$.} By \cite[Lemma~10.2(c)]{WilsonFirstStep}, applied to $C$ and to each other loop of $\mathcal C$ as in the proof of Theorem~\ref{thm:first-step}, the groups of $\mathcal C$ with intervals on $C$ only are the groups of $\mathcal C_C$ none of whose roots has a contact on another loop, with the same intervals, weights, and compatibility. A root with a contact on a loop $C'$ contains a bond of the dependence set $\mathcal F_r$ of a link $r$ of $C'$ \cite[(3.87), Lemma~3.5]{WilsonFirstStep}, whose bonds have both ends in the two blocks of the first step at the ends of the bond of $C'$ which contains $r$, so that it lies within distance $1$ of $C'$ in units of the step $1$. This gives (i) and (ii) at the step $1$, and (iii) follows from the description of the clusters.

\emph{The step $k \to k + 1$.} Let $W$ be a group of $\mathcal C$ of the step $k + 1$ which is not a cross group and has its intervals on $C$. Its old members and the groups of its merged clusters have their intervals on $C$, since the arcs of $W$ contain theirs; in particular no merged cluster contains groups of two loops, which would put arcs on both into $W$. By (i) and (iii) at the step $k$ they are objects of $\mathcal C_C$ with the same insertions and weights, and the rules above assemble from the same constituents a group of $\mathcal C_C$ with the same domains, intervals, network, and insertion. Conversely, let $W$ be a group of $\mathcal C_C$ of the step $k + 1$ whose network is at distance greater than $2$ from the other loops, in units of the step $k + 1$. The networks of its old members and merged clusters lie in that of $W$, at distance greater than $2L$ from the other loops in units of the step $k$, so that by (ii) and (iii) at the step $k$ they are objects of $\mathcal C$ of the kind (i); none of its polymers meets a traversal of another loop; and the rules assemble from these constituents a group of $\mathcal C$ which is not a cross group. Compatibility is the same in both ensembles because it depends only on the domains, the arcs, and the merged clusters, and (iii) follows from (i) and (ii) by the description of the clusters.
\end{proof}

The groups of $\mathcal C_C$ which do not satisfy the hypothesis of (ii) have no counterpart in $\mathcal C$, in general: their constituents may meet objects of another loop, with which they form a cross group of $\mathcal C$. These groups, the cross groups, and the clusters which contain groups of two loops are the objects which distinguish the pair from the two rectangles, and the next lemma shows that they are negligible.

\begin{lemma}[The Pair at the Stopping Step]\label{lem:pair-stopping}
Let $k \in \{ 2, 3 \}$, $K \geq K_{\mathrm w}$, and let $D \coloneqq (k - 1)tL^{j_q}/u$ be the distance of $I_0$ from $I_k$ in units of the step $K'$. On the support of $\varrho_{\ell_0}$,
\begin{equation}\label{eq:pair-stopping}
    \bigl\lvert S(P_k) - S(I_0)S(I_k)\bigr\rvert \leq 3e^{I_*(P_k)}\bigl(e^{N_{\mathrm{ng}}} - 1\bigr)\;,
\end{equation}
where, with $I_{\mathrm{wr}}(P_k) \coloneqq P_{K'}(P_k)C^2\bar\rho\,S_{\mathrm{wr}}(c_{\mathrm{att}}N_q/M)$ the second part of \eqref{eq:insertion-total} and $A^\times(P_k)$ its last part, the cross total \eqref{eq:cross-total} of the pair, for which $d_\times = D$,
\begin{equation}\label{eq:near-mass}
    N_{\mathrm{ng}} \coloneqq A^\times(P_k) + e^{-\nu D/(3M)}I_*(P_k) + I_{\mathrm{wr}}(P_k) + LP_{K'}(P_k)C_{S'}\mathfrak a_{K'}e^{-\nu D/(96A^*_{\mathrm{env}}M)}\;.
\end{equation}
\end{lemma}

\begin{proof}
Let $\mathcal D$ be one of $P_k$, $I_0$, $I_k$. We call a group of $\mathcal D$ at the step $K'$ on $\ell_0$ \emph{far} if it is not a cross group and its network lies within distance $D/3$ of the loop $C$ which carries its intervals, in units of the step $K'$, and a cluster of the reference gas of the step $K'$ \emph{far} if all its groups are far and carried by the same loop. The other groups and clusters are \emph{near}. We put
\begin{equation}\label{eq:far-part}
    F(\mathcal D) \coloneqq \sum_{\mathfrak B \text{ far}}Q^{\mathrm{far}}(\mathfrak B)\,w_{K'}(\mathfrak B)\;,
    \qquad
    Q^{\mathrm{far}}(\mathfrak B) \coloneqq \exp\Bigl(-\sum_{\Gamma \text{ far},\ \Gamma \not\perp \mathfrak B}\Phi^\circ(\Gamma)\Bigr)\;,
\end{equation}
the first sum over the families of far groups, the empty family included, and the second over the far clusters which contain a group incompatible with a member of $\mathfrak B$.

\emph{The far parts factorise.} Let $C \in \{ I_0, I_k \}$ and let $C'$ be the other rectangle. A far group of $P_k$ carried by $C$ has its network at distance at least $D - D/3 > 2$ from $C'$, so by Lemma~\ref{lem:ensemble-locality}(i) it is a group of $\mathcal C_C = C$ with the same network, hence far, and with the same insertion; conversely a far group of $C$ is a group of $P_k$ of the kind (i) by Lemma~\ref{lem:ensemble-locality}(ii), with the same network. By Lemma~\ref{lem:ensemble-locality}(iii) the same holds for the far clusters, with the same weights. Two incompatible groups have networks at distance at most $M$ in units of the step $K'$, by the proof of Lemma~\ref{lem:root-bound}, while the networks of far groups carried by $I_0$ and by $I_k$ are at distance at least $D/3 > M$; so these groups are compatible, and so are a far cluster carried by one rectangle and a far group carried by the other. Hence the families of far groups of $P_k$ are the unions $\mathfrak B_0 \cup \mathfrak B_k$ of families of far groups of $I_0$ and of $I_k$, and $Q^{\mathrm{far}}(\mathfrak B_0 \cup \mathfrak B_k)$ is the product of the two factors of $I_0$ and of $I_k$. The insertions of far groups act on the space of one loop, so that the operator of which $w_{K'}(\mathfrak B_0 \cup \mathfrak B_k)$ is the normalised trace is the tensor product of those of $\mathfrak B_0$ and $\mathfrak B_k$, along the same background field, and the normalised trace of a tensor product is the product of the normalised traces. Therefore
\begin{equation}\label{eq:far-factorisation}
    F(P_k) = F(I_0)\,F(I_k)\;.
\end{equation}

\emph{The near parts are small.} Let $\mathfrak M \coloneqq \sum_We^{\theta(W)}\norm{Y_W}$, the sum over all groups of $\mathcal D$, $\mathfrak N_{\mathrm g} \coloneqq \sum_{W \text{ near}}e^{\theta(W)}\norm{Y_W}$, and $\mathfrak N_{\mathrm c} \coloneqq \sum_{\Gamma \text{ near}}\abs{\Phi^\circ(\Gamma)}$. By \eqref{eq:interaction-ratio} and the bound after it, $Q_{K'}(\mathfrak B) = Q^{\mathrm{far}}(\mathfrak B)\exp(-\sum_{\Gamma \text{ near}, \Gamma \not\perp \mathfrak B}\Phi^\circ(\Gamma))$ for a family of far groups, and $\abs{Q^{\mathrm{far}}(\mathfrak B)} \leq \exp(\sum_{\Gamma \not\perp \mathfrak B}\abs{\Phi^\circ(\Gamma)}) \leq \prod_{W \in \mathfrak B}e^{\theta(W)}$, while $\abs{w_{K'}(\mathfrak B)} \leq \prod_{W \in \mathfrak B}\norm{Y_W}$ as in the proof of Lemma~\ref{lem:term-bound}. Hence, separating in $S(\mathcal D) - F(\mathcal D)$ the families of far groups from those with a near member, and bounding sums over families by sums over sets of groups,
\begin{equation}\label{eq:near-difference}
    \abs{S(\mathcal D) - F(\mathcal D)} \leq e^{\mathfrak M}\bigl(e^{\mathfrak N_{\mathrm c}} - 1\bigr) + \prod_{W \text{ far}}\bigl(1 + e^{\theta(W)}\norm{Y_W}\bigr)\Bigl(\prod_{W \text{ near}}\bigl(1 + e^{\theta(W)}\norm{Y_W}\bigr) - 1\Bigr) \leq e^{\mathfrak M}\bigl(e^{\mathfrak N_{\mathrm g} + \mathfrak N_{\mathrm c}} - 1\bigr)\;,
\end{equation}
where we used $(e^a - 1) + (e^b - 1) \leq e^{a+b} - 1$ for $a, b \geq 0$. In the same way $\abs{S(\mathcal D)} \leq e^{\mathfrak M}$ and $\abs{F(\mathcal D)} \leq e^{\mathfrak M}$, and $\mathfrak M \leq I_*(\mathcal D)$ by the proof of Lemma~\ref{lem:term-bound}, since there are no attached groups on $\ell_0$.

We bound the near masses. The groups of the cross species of $P_k$ have the total at most $A^\times(P_k)$ by Lemma~\ref{lem:cross-total}, with $d_\times = D$, and a single rectangle has none. A near group which is not a cross group has a network which contains a traversal of its loop and a point at distance at least $D/3$ from it; if it was made at the step $j \leq K'$, its network has diameter at least $L^{K'-j}D/3$ in units of the step $j$, so that $\sigma(W) \geq D/(3M)$. If it is of the deviation species, the weighted bound of Definition~\ref{def:dressed-loop}(ii) gives for these groups the total $e^{-\nu D/(3M)}\sum_{\mathfrak n}C_\chi\tilde\rho_{K'}\delta^*_{\mathfrak n}(U_{K'}) \leq e^{-\nu D/(3M)}I_*(\mathcal D)$, by the proof of Lemma~\ref{lem:term-bound}. On $\ell_0$ there are no attached groups, so that a group which is neither of the cross species nor of the deviation species is one of the groups whose domains do not lift to the universal cover (Section~\ref{sec:assembly}), and the total of these is at most $I_{\mathrm{wr}}(\mathcal D)$ by the proof of Lemma~\ref{lem:term-bound}. Thus $\mathfrak N_{\mathrm g} \leq A^\times(\mathcal D) + e^{-\nu D/(3M)}I_*(\mathcal D) + I_{\mathrm{wr}}(\mathcal D)$. A near cluster of nonzero weight contains no cross group, whose flat value is $0$; it contains either groups of both rectangles, which are joined by a chain of incompatible groups, or a group whose network reaches distance $D/3$ from its loop, and in both cases its network has diameter at least $D/3$ in units of the step $K'$, that is $LD/3$ in units of the step $K' - 1$, and contains a traversal of its groups' arcs. By \eqref{eq:cluster-tail} at the step $K'$, summed over the $LP_{K'}(\mathcal D)$ traversals of the step $K' - 1$, $\mathfrak N_{\mathrm c} \leq LP_{K'}(\mathcal D)C_{S'}\mathfrak a_{K'}e^{-2c^T_ELD/3}$, and $2c^T_ELD/3 = \nu D/(96A^*_{\mathrm{env}}M)$. Since $A^\times$, $I_*$, $I_{\mathrm{wr}}$, and $P_{K'}$ do not decrease when a loop is added, $\mathfrak N_{\mathrm g} + \mathfrak N_{\mathrm c} \leq N_{\mathrm{ng}}$ for each of the three diagrams.

\emph{Conclusion.} By \eqref{eq:far-factorisation},
\begin{equation}\label{eq:pair-telescope}
    S(P_k) - S(I_0)S(I_k) = \bigl(S(P_k) - F(P_k)\bigr) - F(I_0)\bigl(S(I_k) - F(I_k)\bigr) - S(I_k)\bigl(S(I_0) - F(I_0)\bigr)\;,
\end{equation}
and \eqref{eq:near-difference} with $I_*(I_0) + I_*(I_k) \leq I_*(P_k)$ gives \eqref{eq:pair-stopping}.
\end{proof}

\begin{theorem}[The Pair Terms]\label{thm:pair-terms}
For $\bar g$ sufficiently small, depending on $q$, $t/u$, $s/u$, $T/u$, $\pi$, and the constants of the construction, and for $k \in \{ 2, 3 \}$,
\begin{equation}\label{eq:pair-terms}
    \abs{\mathfrak y_k(W)} \leq 16\,\sigma(I_0)^2 \leq 16C_\sigma^2\bar g^4(\log\bar g^{-2})^{2m_\sigma}\;.
\end{equation}
\end{theorem}

\begin{proof}
Let $K \geq K_{\mathrm w}$, and write $X_{\mathcal D} \coloneqq e^{E(\mathcal D)}S(\mathcal D)$, so that $W_K(\mathcal D) = w\langle X_{\mathcal D}\rangle + \Lambda(\mathcal D)$ by \eqref{eq:pair-stopping-form}, and abbreviate the indices $P_k$, $I_0$, $I_k$ by $P$, $0$, $k$. Expanding the product $W_K(I_0)W_K(I_k)$ and adding and subtracting $w\langle X_0X_k\rangle$ and $w\langle X_0\rangle\langle X_k\rangle$,
\begin{equation}\label{eq:pair-decomposition}
\begin{split}
    \Delta &\coloneqq W_K(P_k) - W_K(I_0)W_K(I_k) = w\bigl(\langle X_0X_k\rangle - \langle X_0\rangle\langle X_k\rangle\bigr) + w\bigl\langle X_P - X_0X_k\bigr\rangle\\
    &\quad + w(1 - w)\langle X_0\rangle\langle X_k\rangle + \Lambda(P_k) - w\Lambda(I_0)\langle X_k\rangle - w\Lambda(I_k)\langle X_0\rangle - \Lambda(I_0)\Lambda(I_k)\;.
\end{split}
\end{equation}

The first term is a covariance, in which the constant parts of $S = 1 + \xi$ drop out: $\langle X_0X_k\rangle - \langle X_0\rangle\langle X_k\rangle = e^{E(I_0) + E(I_k)}\langle(\xi_{I_0} - \langle\xi_{I_0}\rangle)(\xi_{I_k} - \langle\xi_{I_k}\rangle)\rangle$. By the Cauchy--Schwarz inequality, $\langle\abs{\xi - \langle\xi\rangle}^2\rangle \leq \langle\abs{\xi}^2\rangle$, and Lemma~\ref{lem:stopping-sizes}(i), (iii), its modulus is at most $e^{2\bar g^{10}}\sigma(I_0)\sigma(I_k)$, and $\sigma(I_k) = \sigma(I_0)$ because $I_k$ is a translate of $I_0$, with the same $P_{K'}$ and $I_*$.

The other terms are small. On the support of $\varrho_{\ell_0}$, $X_P - X_0X_k = e^{E(P_k)}(S(P_k) - S(I_0)S(I_k)) + (e^{E(P_k)} - e^{E(I_0) + E(I_k)})S(I_0)S(I_k)$. By Lemma~\ref{lem:stopping-sizes}(i), $\abs{E(\mathcal D)} \leq \bar g^{10}$ and $I_*(\mathcal D) \leq \bar g \leq 1/10$, and $N_{\mathrm{ng}} \leq \bar g^{10}$ for $\bar g$ small, because each of its four terms carries one of the factors $e^{-\nu_\times D/M}$ and $S^{\mathbf R}(\bar g)$ (in \eqref{eq:cross-total}, with $d_\times = D$; the second is smaller than every power of $\bar g$ by \eqref{eq:r-sum-decay}), $e^{-\nu D/(3M)}$, $e^{-2c_{\mathrm{att}}L^3R_{K'}}$ (from $S_{\mathrm{wr}}(b) \leq 3e^{-b}$ for $b \geq 1$ and $N_q/M \geq 2L^3R_{K'}$), and $e^{-\nu D/(96A^*_{\mathrm{env}}M)}$, with $D \geq L^3MR_{K'}$ and $R_{K'} \geq (\log\bar g^{-2})^2$, against a power of $\log\bar g^{-2}$. Lemma~\ref{lem:pair-stopping}, $\abs{S(\mathcal D)} \leq e^{I_*(\mathcal D)}$, and $\abs{e^a - e^b} \leq e^{\max(\abs a, \abs b)}\abs{a - b}$ give $\abs{X_P - X_0X_k} \leq 14\bar g^{10}$, and $\abs{\langle X_{\mathcal D}\rangle} \leq 2$. With Lemma~\ref{lem:stopping-sizes}(i) for $1 - w$ and $\Lambda$, the sum of the moduli of the terms of \eqref{eq:pair-decomposition} other than the first is at most $(14 + 4 + 1 + 4 + 1)\bar g^{10} \leq \bar g^9$ for $\bar g \leq 1/24$. Since $\sigma(I_0) \geq 3\bar g^4$, we have $\bar g^9 \leq \frac12\sigma(I_0)^2$, and $e^{2\bar g^{10}} \leq \frac32$, so that
\begin{equation}\label{eq:pair-difference-bound}
    \abs{\Delta} \leq 2\,\sigma(I_0)^2\;.
\end{equation}

Finally we take the logarithm. By Proposition~\ref{prop:normalisations-drop-out}(ii), $W_K(R_1) = W_K(I_0) = W_K(I_k)$ and $\abs{W_K(I_0)} \geq \frac12$, so that $W_K(P_k) = W_K(I_0)W_K(I_k)(1 + z)$ with $z \coloneqq \Delta/(W_K(I_0)W_K(I_k))$, $\abs{z} \leq 8\sigma(I_0)^2 \leq \frac12$ for $\bar g$ small by Lemma~\ref{lem:stopping-sizes}(iv), and
\begin{equation}\label{eq:pair-log}
    \abs{\mathfrak y_k(W_K)} = \bigl\lvert\log\abs{1 + z}\bigr\rvert \leq 2\abs{z} \leq 16\,\sigma(I_0)^2\;,
\end{equation}
because $\log(1 + x) \leq x$ and $\log(1 - x) \geq -2x$ for $0 \leq x \leq \frac12$. The right side does not depend on $K$, and $\mathfrak y_k(W_K) \to \mathfrak y_k(W)$ as $K \to \infty$ by the continuity of $\log\abs{\cdot}$ on $\mathbb C \setminus \{ 0 \}$, since the limits on the eleven diagrams exist and do not vanish (Proposition~\ref{prop:normalisations-drop-out}(ii)). The second inequality of \eqref{eq:pair-terms} is Lemma~\ref{lem:stopping-sizes}(iv).
\end{proof}

\subsection{The Shape Term}\label{sec:shape-term}

The shape term $\mathfrak y_S$ is a combination of the four rectangles $R_1, \dots, R_4$ alone, and we reduce it, in three steps, to the second moments of the fluxes of the background field. At the stopping step its logarithms are, up to terms of the order $\bar g^4$, the real parts of the expectations of the deviations $\xi_{R_n}$, and each deviation is the sum of the deviation of the leading word and of the insertions. The insertions cancel in the shape term. This is the mechanism which made the coefficients of the normalisation independent of the loop: the expectation of an insertion is a number attached to a traversal, which, by the invariance of the construction under the translations that preserve its partitions, depends only on the position of the traversal relative to the nearest corner, so that its sum over a rectangle is a linear function of the lengths of the sides and of the number of corners of each type, and the weights of the shape term annihilate such functions (Proposition~\ref{prop:insertions-cancel}). The leading words, finally, are expanded to the second order in the logarithms of the holonomies of the rectangles $I_a$, and the identity which writes the boundary of $R_n$ as a product of loops around $I_0, \dots, I_{n-1}$ turns the shape term into a second difference of the covariances of these logarithms (Lemma~\ref{lem:lasso}).

\begin{proposition}[The Shape Term at the Stopping Step]\label{prop:shape-stopping}
For $\bar g$ sufficiently small, depending on $q$, $t/u$, $s/u$, $T/u$, $\pi$, and the constants of the construction, and for $K \geq K_{\mathrm w}$,
\begin{equation}\label{eq:shape-stopping}
    \Bigl\lvert\mathfrak y_S(W_K) - \sum_{n=1}^{4}n_{R_n}\operatorname{Re}\bigl\langle\xi_{R_n}\bigr\rangle\Bigr\rvert \leq \bar g^9 + 8\max_n\sigma(R_n)^2\;.
\end{equation}
\end{proposition}

\begin{proof}
By \eqref{eq:pair-stopping-form}, $W_K(R) = we^{E(R)}\langle S(R)\rangle(1 + z_R)$ with $z_R \coloneqq \Lambda(R)/(we^{E(R)}\langle S(R)\rangle)$, for each of the four rectangles $R$. By Lemma~\ref{lem:stopping-sizes}, $\abs{\langle S(R)\rangle - 1} = \abs{\langle\xi_R\rangle} \leq \sigma(R) \leq \frac12$, $w \geq 1 - \bar g^{10}$, $\abs{E(R)} \leq \bar g^{10}$, and $\abs{\Lambda(R)} \leq \bar g^{10}$, so that $\abs{z_R} \leq 3\bar g^{10} \leq \frac12$ and $\abs{\log\abs{1 + z_R}} \leq 6\bar g^{10}$. Taking logarithms,
\begin{equation}\label{eq:shape-logarithms}
    \log\abs{W_K(R)} = \log w + \operatorname{Re}E(R) + \log\abs{\langle S(R)\rangle} + \log\abs{1 + z_R}\;.
\end{equation}
The weights $n_{R_n} = 1, -3, 3, -1$ sum to zero, so that $\log w$ drops out of $\mathfrak y_S(W_K)$, and their moduli sum to $8$, so that the second and the fourth terms contribute at most $8 \cdot 7\bar g^{10} \leq \bar g^9$. For the third term, $\abs{\log(1 + x) - x} \leq \abs{x}^2$ for complex $x$ with $\abs{x} \leq \frac12$, by the series of the logarithm, and $\log\abs{1 + x} = \operatorname{Re}\log(1 + x)$; with $x = \langle\xi_R\rangle$ this gives $\abs{\log\abs{\langle S(R)\rangle} - \operatorname{Re}\langle\xi_R\rangle} \leq \sigma(R)^2$, and \eqref{eq:shape-stopping} follows.
\end{proof}

Since $\xi_R = (f^{(R)}_{K'} - 1) + I(R)$, the sum in \eqref{eq:shape-stopping} splits into the contribution of the leading words and that of the insertions. We treat the insertions first. The objects of the construction which we use are those of the stopping step on the term $\ell_0$; the partitions of the steps $1, \dots, K'$ are invariant under the translations by the vectors of $\Pi\mathbb Z^4$, in units of the step $K'$, where $\Pi \coloneqq L^{-1}\Pi_{K'}$ with the number $\Pi_{K'}$ of Lemma~\ref{lem:period-points}, a power of $L$. By that lemma and \eqref{eq:stopping-identity-index}, $M \leq \Pi \leq LMR_{K'-1} = L^{-2}a_{K'}$, so that $\Pi$ divides $a_{K'}$, and the corners of the four rectangles, their sides $t_{\mathrm b} \coloneqq tL^{j_q}/u$ and $s_{\mathrm b} \coloneqq sL^{j_q}/u$ in units of the step $K'$, and the differences of their corners lie in $\Pi\mathbb Z^4$ and $\Pi\mathbb Z$, with $\min(t_{\mathrm b}, s_{\mathrm b}) \geq a_{K'} \geq L^2\Pi$. We put
\begin{equation}\label{eq:locality-length}
    \mathsf D \coloneqq \Pi\Bigl\lfloor\frac{\min(t_{\mathrm b}, s_{\mathrm b})}{3\Pi}\Bigr\rfloor\;,
    \qquad\text{so that}\qquad
    \frac29\min(t_{\mathrm b}, s_{\mathrm b}) \leq \mathsf D \leq \frac13\min(t_{\mathrm b}, s_{\mathrm b})\;,
\end{equation}
because $\min(t_{\mathrm b}, s_{\mathrm b})/(3\Pi) \geq L^2/3 \geq 3$. Every side of each of the four rectangles has a length $N_v \geq 3\mathsf D$, and $\mathsf D \geq \frac29L^3MR_{K'}$.

\begin{proposition}[The Insertions Cancel]\label{prop:insertions-cancel}
For $\bar g$ sufficiently small, depending on the quantities of Proposition~\ref{prop:shape-stopping}, and for $K \geq K_{\mathrm w}$,
\begin{equation}\label{eq:insertions-cancel}
    \Bigl\lvert\sum_{n=1}^{4}n_{R_n}\bigl\langle I(R_n)\bigr\rangle\Bigr\rvert \leq 8\max_n\mathsf e(R_n)\;,
\end{equation}
where, for a rectangle $C$,
\begin{equation}\label{eq:insertion-error}
    \mathsf e(C) \coloneqq \frac12I_*(C)^2e^{I_*(C)} + I_*(C)\bigl(\bar h(C) + 2\bar g^8\bigr) + 2e^{-\nu(\mathsf D/4 - M)/(64M)}I_*(C) + I_{\mathrm{wr}}(C)\;,
\end{equation}
and $8\max_n\mathsf e(R_n) \leq C\bar g^3(\log\bar g^{-2})^{m'_\sigma}$ with constants of the same kind.
\end{proposition}

\begin{proof}
Let $C$ be one of the four rectangles. On $\ell_0$, $I(C) = \sum_{\mathfrak B \neq \emptyset}Q_{K'}(\mathfrak B)w_{K'}(\mathfrak B)$, the members of a family having pairwise disjoint arcs on $C$. For a traversal $\mathfrak n$ of the step $K'$ and a point $V$ of the support of $\varrho_{\ell_0}$ let $\bar\nu_{\mathfrak n} \coloneqq \sum_{W \ni \mathfrak n}e^{\theta(W)}\norm{Y_W}$, a function of $V$, so that $\sum_{\mathfrak n}\bar\nu_{\mathfrak n} \leq I_*(C)$ at every such $V$, by the proof of Lemma~\ref{lem:term-bound}. The bounds of the class are bounds on these sums at each $V$, and not on the sums of the suprema of their terms; accordingly every estimate below is made at a fixed $V$ and then integrated, the sums over the groups and the integral being exchanged by Tonelli's theorem under these absolute bounds.

\emph{Families of two or more members.} The members of a family have distinct first traversals. Bounding each member through its first traversal, $\abs{Q_{K'}(\mathfrak B)w_{K'}(\mathfrak B)} \leq \prod_{W \in \mathfrak B}e^{\theta(W)}\norm{Y_W}$ shows that the families of $m$ members contribute at most the elementary symmetric sum $e_m((\bar\nu_{\mathfrak n})_{\mathfrak n}) \leq I_*(C)^m/m!$, and those with $m \geq 2$ at most $e^{I_*(C)} - 1 - I_*(C) \leq \frac12I_*(C)^2e^{I_*(C)}$.

\emph{One member.} For $\mathfrak B = \{ W \}$, let $J$ be the arc of $W$, $\mathsf t_W \coloneqq \pi(U_{K'}(J))$ the transport along it, and $\mathsf H_W \coloneqq \pi(U_{K'}(C_J))$, with $C_J$ the loop $C$ read from the initial site of $J$. Then $w_{K'}(\{ W \}) = \tr(Y_W\mathsf t_W^{-1}\mathsf H_W)$, and $\mathsf H_W$ is unitarily conjugate to $\pi(U_{K'}(C))$, so that $\norm{\mathsf H_W - 1} = h_C \coloneqq \norm{\pi(U_{K'}(C)) - 1}$ and $\abs{w_{K'}(\{ W \}) - \tr(Y_W\mathsf t_W^{-1})} \leq \norm{Y_W}h_C$. Since $h_C \leq \bar h(C)$ on $\mathcal V$, $h_C \leq 2$ everywhere, and $\langle\mathbf 1_{\mathcal V^c}\rangle \leq \bar g^8$ (Lemma~\ref{lem:stopping-sizes}(ii)), and $\abs{Q_{K'}(\{ W \})} \leq e^{\theta(W)}$,
\begin{equation}\label{eq:one-member}
    \Bigl\lvert\bigl\langle I(C)\bigr\rangle - \sum_W\mathsf z^0_W\Bigr\rvert \leq \frac12I_*(C)^2e^{I_*(C)} + I_*(C)\bigl(\bar h(C) + 2\bar g^8\bigr)\;,
    \qquad
    \mathsf z^0_W \coloneqq Q_{K'}(\{ W \})\bigl\langle\tr(Y_W\mathsf t_W^{-1})\bigr\rangle\;,
\end{equation}
the sum over all groups of $C$, which converges absolutely, since $\sum_W\abs{\mathsf z^0_W} \leq I_*(C)$.

\emph{Small and large groups.} We call a group $W$ \emph{small} if it is of the deviation species and $M\sigma(W) + 1 \leq \mathsf D/4$, and \emph{large} otherwise. A large group of the deviation species has $\sigma(W) \geq (\mathsf D/4 - 1)/M$, so that by the weighted bound of Definition~\ref{def:dressed-loop}(ii) these groups contribute at most $e^{-\nu(\mathsf D/4 - 1)/M}I_*(C)$ to $\sum_W\abs{\mathsf z^0_W}$, and the groups of the other species at most $I_{\mathrm{wr}}(C)$, as in the proof of Lemma~\ref{lem:pair-stopping}.

\emph{The interaction, localised.} For a small $W$ let $Q^{\mathsf D}(W) \coloneqq \exp(-\sum'\Phi^\circ(\Gamma))$, the sum over the clusters of the reference gas of the step $K'$ which contain a group incompatible with $W$ and whose networks lie within distance $\mathsf D/4$ of the network of $W$, and $\mathsf z_W \coloneqq Q^{\mathsf D}(W)\langle\tr(Y_W\mathsf t_W^{-1})\rangle$. A cluster which contains a group incompatible with $W$ and is left out of this sum has a group within distance $M$ of the network of $W$ and a network of diameter at least $\mathsf D/4 - M$, in units of the step $K'$, so that $a^\circ(\Gamma) \geq \nu(\mathsf D/4 - M)/(64M)$ by \eqref{eq:network-diameter}. The theorem of Koteck\'y and Preiss in the form \cite[Theorem~5.4]{FriedliVelenik}, with the condition \eqref{eq:kp-condition}, gives $\sum_{\Gamma \not\perp W}\abs{\Phi^\circ(\Gamma)}e^{a^\circ(\Gamma)} \leq a^\circ(W) = \theta(W)$, so that the clusters left out have $\sum\abs{\Phi^\circ(\Gamma)} \leq \theta(W)e^{-\nu(\mathsf D/4 - M)/(64M)}$. Since $Q_{K'}(\{ W \}) = Q^{\mathsf D}(W)e^{-z}$ with $z$ the sum of $\Phi^\circ(\Gamma)$ over these clusters, $\abs{Q^{\mathsf D}(W)} \leq e^{\theta(W)}$, $\abs{z} \leq \theta(W)$, and $\abs{e^{-z} - 1} \leq \abs{z}e^{\abs{z}}$,
\begin{equation}\label{eq:localised-interaction}
    \abs{Q_{K'}(\{ W \}) - Q^{\mathsf D}(W)} \leq \theta(W)e^{2\theta(W)}e^{-\nu(\mathsf D/4 - M)/(64M)} \leq e^{\theta(W) + \nu\sigma(W)}e^{-\nu(\mathsf D/4 - M)/(64M)}\;,
\end{equation}
because $\theta(W)e^{\theta(W)} \leq e^{2\theta(W)}$ and $\theta(W) \leq \nu\sigma(W)/16$. By the weighted bound, $\sum_{W \text{ small}}\abs{\mathsf z^0_W - \mathsf z_W} \leq e^{-\nu(\mathsf D/4 - M)/(64M)}I_*(C)$, and $\sum_{W \text{ small}}\abs{\mathsf z_W} \leq I_*(C)$.

\emph{Agreement.} Let $C$ and $C'$ be two of the four rectangles, possibly equal, $\tau$ a translation by a vector of $\Pi\mathbb Z^4$, $W$ a small group of $C$ whose arc has the first traversal $\mathfrak p$, and $\mathbb B$ the open ball of radius $\mathsf D$ about the initial site of $\mathfrak p$. Suppose that $\tau$ maps the bonds of $C$ with initial site in $\mathbb B$ onto the bonds of $C'$ with initial site in $\tau\mathbb B$, with their orientations. We claim that $\tau W$ is a small group of $C'$ with the first traversal $\tau\mathfrak p$ and that $\mathsf z_{\tau W}(C') = \mathsf z_W(C)$. Indeed, the network of $W$ lies within $\mathsf D/4$ of the initial site of $\mathfrak p$, and the networks of the clusters of $\sum'$ within $\mathsf D/2$; so the arcs of $W$ and of the groups of these clusters are paths of consecutive bonds of $C$ with initial sites in $\mathbb B$, hence also of $\tau^{-1}C'$, since both loops visit no site twice, and each has fewer than $2\mathsf D \leq P_{K'}(C)/6$ traversals. By the argument of Lemma~\ref{lem:locality}(i), these groups, their flat values, the hard core between them, the weights $\Phi^\circ$, and the insertion $Y_W$ as a function of $V$, depend on the loop only through these arcs; and by the argument of Lemma~\ref{lem:locality}(ii), which applies to the objects of a step as functions of the field, $\tau$ carries them to groups, clusters, and weights of $C'$ with the same sizes and species, with $\Phi^\circ(\tau\Gamma; C') = \Phi^\circ(\Gamma; C)$, $Y_{\tau W}(C')(\tau V) = Y_W(C)(V)$, and $\mathsf t_{\tau W}(\tau V) = \mathsf t_W(V)$, where $(\tau V)(b) \coloneqq V(\tau^{-1}b)$. The same argument with $\tau^{-1}$ shows that $\Gamma \mapsto \tau\Gamma$ maps the clusters of $\sum'$ for $W$ onto those for $\tau W$, so that $Q^{\mathsf D}(\tau W) = Q^{\mathsf D}(W)$. Finally the density $\varrho_{\ell_0}$ is composed of the same operations, so that $\varrho_{\ell_0}(\tau V) = \varrho_{\ell_0}(V)$, and the Haar measure on $G^{\mathcal B}$ is invariant under $\tau$, the periods of the torus being multiples of $\Pi$; hence $\langle F\circ\tau^{-1}\rangle = \langle F\rangle$ for bounded measurable $F$, and with $F(V) \coloneqq \tr(Y_W(C)(V)\mathsf t_W(V)^{-1})$ this gives the claim.

\emph{Corners and sides.} For a traversal $\mathfrak p$ of $C$ let $\mathsf y_C(\mathfrak p) \coloneqq \sum\mathsf z_W(C)$, the sum over the small groups $W$ with the first traversal $\mathfrak p$. The initial site of $\mathfrak p$ lies on a side of $C$ of orientation $v$, at a distance $\mathsf d(\mathfrak p) \in [0, N_v)$ from the initial corner of that side. We call $\mathfrak p$ a corner traversal of that corner if $\mathsf d(\mathfrak p) < \mathsf D$, a corner traversal of the terminal corner of the side if $\mathsf d(\mathfrak p) \geq N_v - \mathsf D$, and a middle traversal otherwise; since $N_v \geq 3\mathsf D$, these classes partition the traversals. For a middle traversal $\mathfrak p$ the bonds of $C$ with initial site at distance less than $\mathsf D$ from that of $\mathfrak p$ are the bonds of its own side at distance less than $\mathsf D$ along it, because the adjacent sides are at distance at least $\mathsf D$ and the opposite side at distance at least $3\mathsf D$. For $0 \leq \varrho < \Pi$ let $\mathfrak p_v(\varrho)$ be the traversal of the side of $R_1$ of orientation $v$ with $\mathsf d(\mathfrak p_v(\varrho)) = \mathsf D + \varrho$, a middle traversal, since $\varrho < \Pi \leq \mathsf D \leq N_v(R_1) - 2\mathsf D$. If $\mathsf d(\mathfrak p) \equiv \mathsf D + \varrho$ modulo $\Pi$, the translation $\tau$ with $\tau\mathfrak p_v(\varrho) = \mathfrak p$ has its vector in $\Pi\mathbb Z^4$, and it maps the bonds of $R_1$ near $\mathfrak p_v(\varrho)$ onto those of $C$ near $\mathfrak p$; the agreement, applied with $\tau$ and $\tau^{-1}$, gives $\mathsf y_C(\mathfrak p) = \mathsf y_{R_1}(\mathfrak p_v(\varrho))$. For a corner traversal $\mathfrak p$ of a corner $\mathfrak c$, the bonds of $C$ with initial site at distance less than $\mathsf D$ from that of $\mathfrak p$ lie on the two sides at $\mathfrak c$, within $2\mathsf D$ of $\mathfrak c$. The translation which takes the corner of $R_1$ of the type of $\mathfrak c$ to $\mathfrak c$ has its vector in $\Pi\mathbb Z^4$, and it maps the corner traversals of that corner of $R_1$ onto those of $\mathfrak c$, together with the bonds near each; so the corner total $\mathsf Z_{\mathfrak c} \coloneqq \sum_{\mathfrak p}\mathsf y_C(\mathfrak p)$, over the corner traversals of $\mathfrak c$, depends only on the type of $\mathfrak c$.

\emph{The sum.} The middle traversals of the side of orientation $v$ have $\mathsf d(\mathfrak p) = \mathsf D, \dots, N_v - \mathsf D - 1$; there are $N_v - 2\mathsf D \in \Pi\mathbb Z$ of them, and each residue of $\mathsf d(\mathfrak p) - \mathsf D$ modulo $\Pi$ occurs equally often. Hence, the rearrangement being justified by the absolute convergence,
\begin{equation}\label{eq:small-sum}
    \sum_{W \text{ small}}\mathsf z_W(C) = \sum_{\mathfrak c}\mathsf Z_{\mathfrak c} + \sum_{v \in \mathcal V}\frac{N_v(C) - 2\mathsf D}{\Pi}\,\mathsf S_v\;,
    \qquad
    \mathsf S_v \coloneqq \sum_{\varrho=0}^{\Pi-1}\mathsf y_{R_1}(\mathfrak p_v(\varrho))\;,
\end{equation}
the first sum over the four corners of $C$ and the terms with $N_v(C) = 0$ absent. Each rectangle has one corner of each of the four types and one side of each of the orientations $\pm e_\mu$, $\pm e_\nu$ (Lemma~\ref{lem:witness-counts}(i)), and the weights $n_{R_n}$ satisfy $\sum_nn_{R_n} = 0$ and $\sum_nn_{R_n}N_v(R_n) = 0$ for every $v$, by the proof of Lemma~\ref{lem:witness-counts}(ii). Therefore $\sum_nn_{R_n}\sum_{W \text{ small}}\mathsf z_W(R_n) = 0$.

\emph{Conclusion.} By \eqref{eq:one-member}, the bound on the large groups, \eqref{eq:localised-interaction}, and the vanishing of the weighted sum of the small groups, $\abs{\langle I(C)\rangle - \sum_{W \text{ small}}\mathsf z_W(C)} \leq \mathsf e(C)$, where we used $e^{-\nu(\mathsf D/4 - 1)/M} \leq e^{-\nu(\mathsf D/4 - M)/(64M)}$; the moduli of the weights sum to $8$, which gives \eqref{eq:insertions-cancel}. For the order, $I_*(C) \leq C\bar g^2(\log\bar g^{-2})^{m'}$ and $\bar h(C) \leq C\bar g(\log\bar g^{-2})^{4r+1/2}$ by the proof of Lemma~\ref{lem:stopping-sizes}(iv), while $\mathsf D/(4M) \geq L^3(\log\bar g^{-2})^r/18$ and $I_{\mathrm{wr}}(C)$ carries the factor $e^{-2c_{\mathrm{att}}L^3R_{K'}}$, so that the last two terms of \eqref{eq:insertion-error} decay faster than every power of $\bar g$, since $r = 2$.
\end{proof}

It remains to treat the leading words. On $\mathcal V$ we write $U \coloneqq U_{K'}(V)$, $o_a \coloneqq o + ate_\mu$ for the initial corner of $I_a$, and, for $0 \leq a \leq b \leq 3$, $\mathsf L_{a,b}$ for the closed path which runs from $o_a$ along the bottom side of $R_4$ to $o_b$, around $I_b$, and back to $o_a$, so that $\mathsf L_{a,a} = I_a$. By the proof of Lemma~\ref{lem:stopping-sizes}(ii), applied to the defining representation of $G$, for which the number \eqref{eq:representation-constant} is $1$, every $I_b$ satisfies $\norm{U(I_b) - 1} \leq h^G \coloneqq 1.06 \cdot \frac32c_hB_3P_{K'}(I_0)^2\varepsilon^{\mathrm{act}}$ on $\mathcal V$, and $U(\mathsf L_{a,b})$ is unitarily conjugate to $U(I_b)$. If $h^G \leq r_G$, the principal logarithm $X_{a,b} \coloneqq -i\log U(\mathsf L_{a,b})$ therefore lies in $\mathfrak g$, and $\norm{X_{a,b}} \leq 1.05\,h^G$, by the argument of the proof of Lemma~\ref{lem:near-one} applied to the eigenvalues. We put $X_a \coloneqq X_{a,a}$,
\begin{equation}\label{eq:flux-moments}
    \mathsf K_{a,b} \coloneqq \tr\bigl(\pi_*(X_a)\,\pi_*(X_{a,b})\bigr)\;,
    \qquad
    \Psi_k \coloneqq \bigl\langle\mathbf 1_{\mathcal V}\mathsf K_{0,k}\bigr\rangle\;,
\end{equation}
which are real, since both factors are Hermitian, and $\mathsf l \coloneqq 1.05\,c_\pi h^G$.

\begin{lemma}[The Lasso]\label{lem:lasso}
For $\bar g$ sufficiently small, depending on the quantities of Proposition~\ref{prop:shape-stopping}, and for $K \geq K_{\mathrm w}$, $h^G \leq r_G$, $\mathsf l \leq 1/16$, and
\begin{equation}\label{eq:lasso}
    \sum_{n=1}^{4}n_{R_n}\bigl\langle f^{(R_n)}_{K'} - 1\bigr\rangle = \Psi_3 - \Psi_2 + r_H\;,
    \qquad
    \abs{r_H} \leq 512\,\mathsf l^3 + 16\bar g^8\;.
\end{equation}
\end{lemma}

\begin{proof}
The conditions on $h^G$ and $\mathsf l$ hold for $\bar g$ small, since $h^G \leq \bar h(I_0)/c_\pi$, which tends to zero with $\bar g$ by Lemma~\ref{lem:stopping-sizes}(iv).

\emph{The lasso.} As closed paths from $o$, $\partial R_1 = \mathsf L_{0,0}$, and $\partial R_{n+1}$ is $\mathsf L_{0,n}$ followed by $\partial R_n$: in the product the return of $\mathsf L_{0,n}$ along the bottom side cancels against the first part of $\partial R_n$, and the left side of $I_n$ against the right side of $R_n$. Hence $U(\partial R_n) = U(\mathsf L_{0,n-1})\cdots U(\mathsf L_{0,0})$. For $a \leq b$, the paths $\mathsf L_{0,a}$ and $\mathsf L_{0,b}$ both run first from $o$ to $o_a$ and finally back, so that $U(\mathsf L_{0,a}) = \mathsf tU(I_a)\mathsf t^{-1}$ and $U(\mathsf L_{0,b}) = \mathsf tU(\mathsf L_{a,b})\mathsf t^{-1}$ with the same transport $\mathsf t \in G$ from $o$ to $o_a$. The principal logarithm commutes with unitary conjugation, so $X_{0,a} = \mathsf tX_a\mathsf t^{-1}$ and $X_{0,b} = \mathsf tX_{a,b}\mathsf t^{-1}$, and since $\pi_*(gXg^{-1}) = \pi(g)\pi_*(X)\pi(g)^{-1}$, the numbers $\tr(\pi_*(X_{0,a})\pi_*(X_{0,b}))$ equal $\mathsf K_{a,b}$.

\emph{The second order.} Put $Z_a \coloneqq \pi_*(X_{0,a})$, a Hermitian matrix with $\pi(U(\mathsf L_{0,a})) = e^{iZ_a}$ and $\norm{Z_a} \leq c_\pi\norm{X_{0,a}} \leq \mathsf l$, and $S \coloneqq \sum_{a<n}\norm{Z_a} \leq 4\mathsf l \leq \frac14$. Write $e^{iZ_a} = 1 + \mathsf k_a$ with $\mathsf k_a = iZ_a + \mathsf j_a$ and $\mathsf j_a = -\frac12Z_a^2 + \mathsf E_a$; the exponential series gives $\norm{\mathsf E_a} \leq \frac14\norm{Z_a}^3$, $\norm{\mathsf j_a} \leq 0.57\norm{Z_a}^2$, and $\norm{\mathsf k_a} \leq 1.3\norm{Z_a}$ for $\norm{Z_a} \leq \frac14$. In the expansion of the ordered product of the $1 + \mathsf k_a$, the terms with three or more factors $\mathsf k$ have total norm at most $(1.3S)^3e^{1.3S}/6 \leq 0.51S^3$; the terms with two factors are $-\sum Z_aZ_b$, over the ordered pairs of the product, up to $0.66\sum_{a<b}\norm{Z_a}\norm{Z_b}(\norm{Z_a} + \norm{Z_b}) \leq 0.22S^3$, because the last sum is at most $S^3/3$; and $\sum_a\norm{\mathsf E_a} \leq S^3/4$. So
\begin{equation}\label{eq:lasso-expansion}
    \pi(U(\partial R_n)) = 1 + i\sum_{a<n}Z_a - \frac12\sum_{a<n}Z_a^2 - \sum Z_aZ_b + R_n\;,
    \qquad
    \norm{R_n} \leq S^3\;.
\end{equation}
By Lemma~\ref{lem:near-one}, $\tr Z_a = 0$, and $\tr(Z_aZ_b) = \tr(Z_bZ_a)$, so that $f^{(R_n)}_{K'} - 1 = -\frac12\tr(\sum_{a<n}Z_a)^2 + r_n$ with $\abs{r_n} \leq S^3 \leq 64\,\mathsf l^3$. With $\mathsf v(m) \coloneqq \sum_{n>m}n_{R_n}$, so that $\mathsf v(0), \dots, \mathsf v(3) = 0, -1, 2, -1$,
\begin{equation}\label{eq:lasso-weights}
    \sum_{n=1}^{4}n_{R_n}\tr\Bigl(\sum_{a<n}Z_a\Bigr)^2 = \sum_{a,b \leq 3}\mathsf v(\max(a, b))\tr(Z_aZ_b) = \sum_a\mathsf v(a)\mathsf K_{a,a} + 2\sum_{a<b}\mathsf v(b)\mathsf K_{a,b}\;.
\end{equation}

\emph{Stationarity.} Let $\tau$ be the translation by $t_{\mathrm b}e_\mu$, whose vector lies in $\Pi\mathbb Z^4$, acting on $V$ as in the proof of Proposition~\ref{prop:insertions-cancel}. By the argument of Lemma~\ref{lem:locality}(ii), $U_{K'}(\tau V) = \tau U_{K'}(V)$, so that $U_{K'}(\tau V)(\mathsf L_{a,b}) = U_{K'}(V)(\mathsf L_{a-1,b-1})$ for $a \geq 1$, and the event $\mathcal V$ is invariant under $\tau$, since $A(U_{K'}(\cdot))$ and the support of $\varrho_{\ell_0}$ are. With the invariance of $\langle\cdot\rangle$ under $\tau$, applied $a$ times, $\langle\mathbf 1_{\mathcal V}\mathsf K_{a,b}\rangle = \Psi_{b-a}$.

\emph{Conclusion.} Hence $\langle\mathbf 1_{\mathcal V}\sum_nn_{R_n}\tr(\sum_{a<n}Z_a)^2\rangle = \Psi_0\sum_a\mathsf v(a) + 2\sum_{a<b}\mathsf v(b)\Psi_{b-a} = 2(\Psi_2 - \Psi_3)$, because $\sum_a\mathsf v(a) = 0$ and the coefficients of $\Psi_1$, $\Psi_2$, $\Psi_3$ in the second sum are $-1 + 2 - 1 = 0$, $2 - 1 = 1$, and $-1$. On $\mathcal V^c$ we use $\abs{f^{(R_n)}_{K'} - 1} \leq 2$, the moduli of the weights summing to $8$, and $\langle\mathbf 1_{\mathcal V^c}\rangle \leq \bar g^8$. Together, $\sum_nn_{R_n}\langle f^{(R_n)}_{K'} - 1\rangle = \Psi_3 - \Psi_2 + r_H$ with $\abs{r_H} \leq 8 \cdot 64\,\mathsf l^3 + 16\bar g^8$.
\end{proof}

\subsection{The Coarse Covariance at Flat Fields}\label{sec:flat-covariance}

The second moments $\Psi_k$ are expectations under the measure $\langle\cdot\rangle$ of the stopping step, which at small coupling is concentrated near the fields $V$ whose minimal configurations are nearly flat. The integration by parts of Section~\ref{sec:flux-moments} will express them, to the order $\bar g^2$, through a quadratic form evaluated at such fields, and in this section we compute that form at the fields which are exactly flat. There it is the covariance of the fluxes of a Gaussian field on $\mathbb T_{j_q}$, the block field of the free field, and the question is how far it is from the covariance of the free field itself, which Lemma~\ref{lem:free-field} bounds below. The two differ by the covariance of the fluxes given the block averages (Lemma~\ref{lem:flux-splitting}). This is the covariance of a constrained Gaussian field of Ba{\l}aban's, whose kernel decays exponentially on the scale of the blocks, so that it contributes almost nothing to the separated rectangles $I_0$ and $I_k$ (Lemma~\ref{lem:fluctuation-part}). The integration by parts needs, besides, an inverse of the Hessian of the action, which has a kernel at a flat field; we regularise it by a small number $\lambda$, and the price, a term proportional to $\lambda$, is controlled by a coercivity estimate for the block field (Lemmata~\ref{lem:coarse-coercivity} and~\ref{lem:flux-norm}).

\emph{The setting.} We use the units of the step $K'$ of Section~\ref{sec:stopping-scale}: the block field $V$ lives on the bonds $\mathcal B$ of $\mathbb T_{j_q}$, which has the spacing $1$ and $N_q$ sites per side, and the fine lattice $\mathbb T_K$ has the spacing $\eta \coloneqq L^{-K'}$. In these units the rectangles have the sides $t_{\mathrm b} = tL^{j_q}/u$ and $s_{\mathrm b} = sL^{j_q}/u$, integers with $12t_{\mathrm b} < N_q$ and $3s_{\mathrm b} < N_q$, and we write $\mathsf I_a$ for the set of the plaquettes of $\mathbb T_K$ in the plane of $e_\mu$ and $e_\nu$ inside the rectangle bounded by $I_a$, so that $\abs{\mathsf I_a} \coloneqq \eta^2\#\mathsf I_a = s_{\mathrm b}t_{\mathrm b}$. For a field $V$ which satisfies \cite[(7), p.~278]{BalabanVAR}, the minimal configurations $U$ with $M^{K'}U = V$ form one orbit of the gauge transformations of $\mathbb T_K$ which equal $1$ at the sites of $\mathbb T_{j_q}$ \cite[(4), p.~278; Theorem~1, p.~279]{BalabanVAR}, and $U_{K'}(V)$ denotes any of them. Since the corner $o$ of the rectangles is a site of $\mathbb T_q \subset \mathbb T_{j_q}$, the holonomy of a closed path based at $o$ is the same for all of them, so that the logarithms $X_{a,b}$ of Section~\ref{sec:shape-term} are functions of $V$; and $\mathsf A(V) \coloneqq A(U_{K'}(V))$, the Wilson action of the minimal configuration, is a function of $V$ which is invariant under the gauge transformations of $\mathbb T_{j_q}$. On $\mathfrak g$ we use the inner product $X \cdot Y \coloneqq \tr(XY)$, which is invariant under $\mathrm{Ad}$ and is the product of \cite[p.~392]{BalabanPBG}, with $\abs X_2 \coloneqq (X \cdot X)^{1/2}$, so that $\abs X_2 \leq \norm X \leq N^{1/2}\abs X_2$ because $\tr$ is normalised. We put $d_{\mathfrak g} \coloneqq \dim G$ and fix an orthonormal basis $(t_\alpha)_{\alpha \leq d_{\mathfrak g}}$ of $\mathfrak g$.

\emph{The Casimir number.} For the representation $\pi$ of the eleven diagrams, with the normalised trace $\tr$ of its space, we define $C_2(\pi)$ and record a consequence of Schur's lemma:
\begin{equation}\label{eq:casimir-number}
    \sum_{\alpha=1}^{d_{\mathfrak g}}\pi_*(t_\alpha)^2 = C_2(\pi) \cdot 1\;,
    \qquad
    \tr\bigl(\pi_*(Y)\,\pi_*(Z)\bigr) = \frac{C_2(\pi)}{d_{\mathfrak g}}\,Y \cdot Z \quad (Y, Z \in \mathfrak g)\;.
\end{equation}
Indeed, the sum does not depend on the orthonormal basis, since an orthogonal change of basis leaves $\sum_\alpha t_\alpha \otimes t_\alpha$ invariant, and $\pi(g)\pi_*(t_\alpha)\pi(g)^{-1} = \pi_*(\mathrm{Ad}(g)t_\alpha)$ maps it to the sum for the basis $(\mathrm{Ad}(g)t_\alpha)$; so it commutes with $\pi(G)$ and is a scalar, $\pi$ being irreducible. The form $(Y, Z) \mapsto \tr(\pi_*(Y)\pi_*(Z))$ is real, symmetric, and $\mathrm{Ad}$-invariant, so it equals $(SY) \cdot Z$ for a symmetric operator $S$ on $\mathfrak g$ which commutes with $\mathrm{Ad}(G)$; an eigenspace of $S$ is then invariant under $\mathrm{Ad}(G)$, hence an ideal of $\mathfrak g$, hence $\mathfrak g$ itself because $G$ is simple, and the trace of the first identity gives $S = C_2(\pi)/d_{\mathfrak g}$. Finally $C_2(\pi) > 0$: the number $n_\pi C_2(\pi) = \sum_\alpha\operatorname{Tr}(\pi_*(t_\alpha)^2)$ is a sum of squared Hilbert--Schmidt norms of Hermitian matrices, which vanishes only if $\pi_* = 0$, and then $\pi = 1$ on the exponentials, which generate the connected group $G$, while $\pi$ is not trivial.

\emph{The fluxes and their derivatives.} Let $\mathcal V_3 \coloneqq \{ V \in \operatorname{supp}\varrho_{\ell_0} \,|\, \mathsf A(V) < 3T_{\mathrm{ng}}g_{K'}^2 \}$, an open neighbourhood of $\mathcal V$ in the support. The proof of Lemma~\ref{lem:stopping-sizes}(ii) with $3T_{\mathrm{ng}}$ in place of $T_{\mathrm{ng}}$, whose conditions hold for $\bar g$ small by the same argument, multiplies $\varepsilon^{\mathrm{act}}$ by $\sqrt3$ and gives $\norm{U(I_b) - 1} \leq \sqrt3h^G$ on $\mathcal V_3$, so that the logarithms $X_{a,b}$ of Section~\ref{sec:shape-term} are defined on $\mathcal V_3$ and have $\norm{X_{a,b}} \leq 1.05\sqrt3h^G \leq 2h^G$ once $\sqrt3h^G \leq r_G$. We put
\begin{equation}\label{eq:flux-coordinates}
    X^\alpha \coloneqq t_\alpha \cdot X_{0,0}\;,
    \qquad
    \Xi_\alpha \coloneqq t_\alpha \cdot \bigl(X_{0,3} - X_{0,2}\bigr)\;,
    \qquad\text{so that}\qquad
    \mathsf K_{0,3} - \mathsf K_{0,2} = \frac{C_2(\pi)}{d_{\mathfrak g}}\sum_\alpha X^\alpha\,\Xi_\alpha
\end{equation}
by \eqref{eq:casimir-number}. For $u \in \mathfrak g^{\mathcal B}$ we write $e^{iu}V$ for the field $(e^{iu_b}V(b))_b$, $\partial_uF(V) \coloneqq \frac{\dd}{\dd s}F(e^{isu}V)|_{s=0}$ for a function $F$, and $\nabla F(V) \in \mathfrak g^{\mathcal B}$ for its gradient, with $\langle u, \nabla F\rangle = \partial_uF$ for the inner product $\langle u, u'\rangle \coloneqq \sum_bu_b \cdot u'_b$; we put $\norm u \coloneqq \langle u, u\rangle^{1/2}$ and $\norm u_\infty \coloneqq \max_b\norm{u_b} \leq N^{1/2}\norm u$, and we let $\Delta(V)$ be the symmetric operator on $\mathfrak g^{\mathcal B}$ with $\langle u, \Delta(V)u\rangle = \frac{\dd^2}{\dd s^2}\mathsf A(e^{isu}V)|_{s=0}$. For $\lambda > 0$ such that $\Delta(V) + \lambda$ is invertible we define
\begin{equation}\label{eq:c-lambda}
    c_\lambda(V) \coloneqq \sum_{\alpha=1}^{d_{\mathfrak g}}\bigl\langle\nabla X^\alpha, (\Delta(V) + \lambda)^{-1}\nabla\Xi_\alpha\bigr\rangle\;.
\end{equation}
This number does not depend on the orthonormal basis, and it is invariant under the gauge transformations $\omega$ of $\mathbb T_{j_q}$: these act on $\mathfrak g^{\mathcal B}$ by the orthogonal map $u_b \mapsto \mathrm{Ad}(\omega(b_-))u_b$, under which $\nabla$ and $\Delta$ are covariant, and on the logarithms based at $o$ by $\mathrm{Ad}(\omega(o))$, which changes the basis. The functions $X^\alpha$, $\Xi_\alpha$, and $\mathsf A$ are real analytic on $\mathcal V_3$, by \cite[Proposition~9, p.~309]{BalabanVAR}, as Lemma~\ref{lem:moved-minimiser} below makes precise.

\emph{The linearisation at a flat field.} We call $V_0$ \emph{flat} if $V_0(\partial p) = 1$ for every plaquette $p$ of $\mathbb T_{j_q}$. On the $\mathfrak g$-valued $1$-forms and $2$-forms of $\mathbb T_K$ we use the products $\langle A, A'\rangle \coloneqq \eta^4\sum_bA_b \cdot A'_b$ and $\langle F, F'\rangle \coloneqq \eta^4\sum_pF_p \cdot F'_p$, over the bonds and the positively oriented plaquettes, and for a configuration $U$ we write $D$ for the covariant exterior derivative of $1$-forms at $U$, $(DA)(p)$ being the coefficient of $i\eta^2$ in $U'(\partial p)U(\partial p)^{-1}$ to first order for $U' = e^{i\eta A}U$, and $\partial_U$ for the covariant derivative of $0$-forms, as in \cite[(3.8)--(3.10), p.~392]{BalabanPBG}, which writes $D$ for both. Let $Q$ be the linearisation at $U$ of the average of order $K'$, $M^{K'}(e^{i\eta A}U) = e^{i(QA + O(A^2))}M^{K'}(U)$ on each bond of $\mathcal B$, and $Q^{\mathrm{dr}}$ the linear part of the dressed average of Supplement, Lemma~17.7, at $U$. For a $2$-form $\chi$ and $t \in \mathfrak g$ let $\chi \otimes t$ be the $2$-form $p \mapsto \chi(p)t$, and let $\chi_a \coloneqq \eta^{-2}\mathbf 1_{\mathsf I_a}$, the plaquettes being oriented from $e_\mu$ to $e_\nu$, so that $\langle\chi_a \otimes t, F\rangle = \eta^2\sum_{p \in \mathsf I_a}t \cdot F(p)$, the flux of $t \cdot F$ through $I_a$.

\begin{lemma}[Flat Fields]\label{lem:flat-fields}
Let $V_0$ be flat, and let $\sigma_1, \dots, \sigma_4$ be integers such that each hyperplane $\{ x_\beta \equiv \sigma_\beta + \frac12 \}$ has distance at least $3$ from the rectangle bounded by $R_4$.
\begin{enumerate}[label=(\roman*)]
\item There are commuting $h_1, \dots, h_4 \in G$ and a gauge transformation $\omega$ of $\mathbb T_{j_q}$ with $\omega(o) = 1$ such that $V_0^\omega$ equals $h_\beta$ on the bonds $\langle y, y + e_\beta\rangle$ with $y_\beta \equiv \sigma_\beta$ and $1$ on the other bonds, and $U_{K'}(V_0^\omega)$ may be taken to be the flat configuration $U_{\mathrm s}$ which equals $h_\beta$ on the bonds $\langle x, x + \eta e_\beta\rangle$ of $\mathbb T_K$ with $x_\beta \equiv \sigma_\beta + \frac12(1 - \eta)$ and $1$ on the other bonds.
\item Let $(E_r)_{r \leq d_{\mathfrak g}}$ be an orthonormal basis of $\mathfrak g^{\mathbb C}$, for the Hermitian extension of the product, with $\mathrm{Ad}(h_\beta)E_r = e^{2\pi i\theta_{r,\beta}}E_r$ and $\theta_r \in [0, 1)^4$. In the gauge of (i), the components $\tr(E_r^*B)$ of a $\mathfrak g$-valued form $B$ on $\mathbb T_K$ or on $\mathbb T_{j_q}$, read in the box $\prod_\beta(\sigma_\beta + \frac12, \sigma_\beta + \frac12 + N_q)$ and extended by $B(x + N_qe_\beta) = e^{2\pi i\theta_{r,\beta}}B(x)$, are scalar forms with the twisted periodicity $\theta_r$; the operators $D$, $\partial_U$, $Q$, $Q^{\mathrm{dr}}$ at $U_{\mathrm s}$, and the translations of $\mathbb T_{j_q}$ by whole blocks, act on each component as the corresponding operators at $U = 1$ on such forms, $Q^{\mathrm{dr}}$ as the average $Q_{K'}$ of \cite[(1.18), p.~20]{BalabanPRT1} with centred blocks.
\item Let $H$ be a linear map from the forms on $\mathcal B$ to the $1$-forms on $\mathbb T_K$ such that $QHu = u$ and $\norm{DHu} = \min\{ \norm{DA} \,|\, QA = u \}$, and $K \coloneqq (DH)^*DH$. The operator $DH$ does not depend on the choice of $H$ and is the same for $Q^{\mathrm{dr}}$ in place of $Q$, and so is the subspace $D(\ker Q)$.
\item In the gauge of (i), $\Delta(V_0) = K$, $\nabla X^\alpha(V_0) = H^*D^*(\chi_0 \otimes t_\alpha)$, and $\nabla\Xi_\alpha(V_0) = H^*D^*((\chi_3 - \chi_2) \otimes t_\alpha)$.
\end{enumerate}
\end{lemma}

\begin{proof}
(i) The torus cut along the four hyperplanes $\{ x_\beta \equiv \sigma_\beta + \frac12 \}$ is a box, and the transport of $V_0$ from $o$ along paths in the box does not depend on the path, $V_0$ being flat; with $\omega(y)$ the inverse of the transport from $o$ to $y$, $V_0^\omega = 1$ on the bonds of the box, and on a bond which crosses the hyperplane of direction $\beta$ it equals the holonomy $h_\beta$ of a closed path from $o$ which winds once in that direction, independently of the bond by flatness. Two such holonomies commute, since the closed path around the cycles $\beta$ and $\gamma$ and back is contractible. The configuration $U_{\mathrm s}$ is flat, since a plaquette which meets the fine bonds of a hyperplane contains two of them with opposite orientations, or four at the intersection of two hyperplanes, where the $h_\beta$ commute. Near the hyperplane of direction $\beta$ it is the gauge transform of $1$ by the transformation equal to $1$ on one side and to $h_\beta^{-1}$ on the other, which is $1$ or $h_\beta^{-1}$ at the sites of $\mathbb T_{j_q}$, and the average \cite[(0.4), p.~253]{BalabanRG1} is covariant under such a transformation through its values at the endpoints of each bond \cite[(0.6), p.~253]{BalabanRG1}; so all the averages of $U_{\mathrm s}$ are configurations of the same form, and $M^{K'}U_{\mathrm s} = V_0^\omega$. Its action is $0$, so it is a minimum, and it lies on the unique critical orbit \cite[Theorem~1, p.~279]{BalabanVAR}.

(ii) The gauge transformations which cross the hyperplanes act on $\mathfrak g$ by $\mathrm{Ad}(h_\beta^{\pm1})$, which multiplies the component $r$ by $e^{\pm2\pi i\theta_{r,\beta}}$; off the hyperplanes $U_{\mathrm s} = 1$, and the operators are covariant (Supplement, Lemma~17.7(a), and \cite[(3.32), p.~395]{BalabanPBG}), so in the extended components they are the operators at $U = 1$. At $U = 1$ each step of the dressed average is the straight average with centred blocks (Supplement, proof of Lemma~17.7), whose composite is \cite[(1.18), p.~20]{BalabanPRT1}.

(iii) The set $\mathcal F(u) \coloneqq \{ DA \,|\, QA = u \}$ is an affine subspace, so it has a unique point of least norm, which is $DHu$ for every admissible $H$. At $U_{\mathrm s}$ the gauge directions satisfy $D\partial_U = 0$, and differentiating the covariance \cite[(0.6), p.~253]{BalabanRG1} gives $Q\partial_U\tilde\omega = \bar D(\tilde\omega|_{\mathbb T_{j_q}})$, with $\bar D$ the covariant derivative of $0$-forms of $\mathbb T_{j_q}$ at $V_0^\omega$; so $\mathcal F(u) = \mathcal F(u + \bar D\omega')$ for every $0$-form $\omega'$ of $\mathbb T_{j_q}$, the restriction to the sites being onto. The linearisation of \cite[(0.4)]{BalabanRG1} at $1$ averages $A$ over the paths $\Gamma \cup [x, x'] \cup (-\Gamma')$, with $\Gamma$ from $c_-$ to $x$ and $\Gamma'$ from $c_+$ to $x'$ inside the blocks, since the factors $(-c)$ and $\mathbf U(c)$ cancel to first order; the parts along $\Gamma$ and $\Gamma'$ are the values at $c_-$ and $c_+$ of a $0$-form $\psi_A$ of $\mathbb T_{j_q}$, so that one step is the straight average minus $\bar D\psi_A$. Composing over the steps, with the identity $Q_{K'}\partial\tilde\omega = \bar DQ'_{K'}\tilde\omega$ of \cite[(1.20), p.~20]{BalabanPRT1}, $Q - Q^{\mathrm{dr}}$ takes values in the range of $\bar D$ in each component by (ii), so $\mathcal F(u)$ is the same for $Q^{\mathrm{dr}}$, and $DH$ is the same. Finally $D(\ker Q)$ is the orthogonal complement of $D(\operatorname{ran}H)$ in $\operatorname{ran}D$ by Lemma~\ref{lem:flux-splitting}(ii) below, for $Q$ and for $Q^{\mathrm{dr}}$.

(iv) Let $u \in \mathfrak g^{\mathcal B}$. The map $A \mapsto \frac1i\log(M^{K'}(e^{i\eta A}U_{\mathrm s})(V_0^\omega)^{-1})$ is analytic near $0$ with the derivative $Q$, which has the right inverse $H$, so the implicit function theorem gives an analytic family $A_{\mathrm c}(s) = sHu + O(s^2)$ with $M^{K'}(e^{i\eta A_{\mathrm c}(s)}U_{\mathrm s}) = e^{isu}V_0^\omega$. At the flat $U_{\mathrm s}$ the expansion \cite[(3.12), p.~392]{BalabanPBG} has $J = 0$ and $\Delta = D^*D$, since the form $\Delta'$ of \cite[(3.10)]{BalabanPBG} vanishes, so that $\mathsf A(e^{isu}V_0^\omega) \leq A(e^{i\eta A_{\mathrm c}(s)}U_{\mathrm s}) = \frac12s^2\norm{DHu}^2 + O(s^3)$. Conversely, by \cite[Proposition~9, p.~309]{BalabanVAR}, $U_{K'}(e^{isu}V_0^\omega) = (e^{i\eta A(s)}U_{\mathrm s})^{w(s)}$ with $A$ and the gauge transformation $w$ analytic, $A(0) = 0$, $w(0) = 1$, and by the covariance $M^{K'}(e^{i\eta A(s)}U_{\mathrm s}) = (e^{isu}V_0^\omega)^{w(s)^{-1}}$ with $w$ read at the sites, so that $QA'(0) = u + \bar D\omega'$ for a $0$-form $\omega'$; hence $DA'(0) \in \mathcal F(u)$ and $\mathsf A(e^{isu}V_0^\omega) = \frac12s^2\norm{DA'(0)}^2 + O(s^3) \geq \frac12s^2\norm{DHu}^2 + O(s^3)$. Therefore $\langle u, \Delta(V_0)u\rangle = \norm{DHu}^2$ and $DA'(0) = DHu$. The loops $\mathsf L_{0,k}$ lie at distance at least $3$ from the hyperplanes, where $U_{\mathrm s} = 1$, so by the lattice Stokes formula the holonomy of $\mathsf L_{0,k}$ at $e^{i\eta A(s)}U_{\mathrm s}$ is $1 + is\eta^2\sum_{p \in \mathsf I_k}DA'(0)(p) + O(s^2)$, the conjugation by the transport along the bottom side of $R_4$ and by $w(s)(o)$ changing only terms of the second order because the holonomy at $s = 0$ is $1$. Hence $\partial_uX^\alpha = \langle\chi_0 \otimes t_\alpha, DHu\rangle$ and $\partial_u\Xi_\alpha = \langle(\chi_3 - \chi_2) \otimes t_\alpha, DHu\rangle$, which is the claim.
\end{proof}

At a flat field the problem has become linear, and in each component it is the problem of the free field with a twist. We first split the fluxes, in an abstract setting, into the part which the block averages determine and the rest.

\begin{lemma}[Flux Splitting]\label{lem:flux-splitting}
Let $E$, $F$, $E'$ be finite-dimensional real or complex inner product spaces, $D \colon E \to F$ and $Q \colon E \to E'$ linear, $Q$ onto.
\begin{enumerate}[label=(\roman*)]
\item There is a linear $H \colon E' \to E$ with $QH = 1$ such that $HB$ minimises $\norm{DA}$ on $\{ QA = B \}$ for every $B$, and every such $H$ has $D^*DH(E') \subset \operatorname{ran}Q^*$.
\item $\operatorname{ran}D = D(\operatorname{ran}H) \oplus D(\ker Q)$, the two summands being orthogonal.
\item $\Pi_{\operatorname{ran}D} = \Pi_{D\operatorname{ran}H} + \Pi_{D\ker Q}$, and $\Pi_{D\operatorname{ran}H} = DH(H^*D^*DH)^+H^*D^*$, where $\Pi$ denotes the orthogonal projection and ${}^+$ the Moore--Penrose inverse.
\end{enumerate}
\end{lemma}

\begin{proof}
(i) Let $\iota$ be the inclusion of $\ker Q$. Since $Q$ is onto, $QQ^+ = 1$ and $\{ QA = B \} = Q^+B + \ker Q$. Put $H \coloneqq (1 - \iota(D\iota)^+D)Q^+$, so that $QH = 1$ because $Q\iota = 0$. Then $DHB = (1 - D\iota(D\iota)^+)DQ^+B$ is $DQ^+B$ minus its orthogonal projection onto $D(\ker Q) = \operatorname{ran}D\iota$, that is the point of least norm of $DQ^+B + D(\ker Q)$, so $HB$ is a minimiser. At every minimiser $HB$ the first variation gives $\langle DHB, Dz\rangle = 0$ for $z \in \ker Q$, so $D^*DHB \in (\ker Q)^\perp = \operatorname{ran}Q^*$.

(ii) For $z \in \ker Q$ and $B \in E'$, by (i), $\langle DHB, Dz\rangle = \langle Q^*y, z\rangle = \langle y, Qz\rangle = 0$ for some $y$; and every $A$ is $HQA + (A - HQA)$ with $Q(A - HQA) = 0$.

(iii) The first identity follows from (ii), and the second is $T(T^*T)^+T^* = \Pi_{\operatorname{ran}T}$ for $T \coloneqq DH$.
\end{proof}

We apply the lemma with $E$ the $\mathfrak g$-valued $1$-forms on $\mathbb T_K$, $F$ the $2$-forms, $E'$ the forms on $\mathcal B$, and $D$, $Q$, $H$ of Lemma~\ref{lem:flat-fields} at a flat field. In the component $r$ the projection $\Pi_{\operatorname{ran}D}$ is diagonal in the momenta $p = (2\pi/N_q)(n + \theta_r)$, and its entry on the plaquettes of the plane of $e_\mu$ and $e_\nu$ is $(\hat p_\mu^2 + \hat p_\nu^2)/\hat p^2$, the momentum $p = 0$ being absent if $\theta_r = 0$: the range of $D$ at the momentum $p$ consists of the forms $a \wedge b$ with $a_\beta = \eta^{-1}(e^{i\eta p_\beta} - 1)$, $\abs{a_\beta}^2 = \hat p_\beta^2$, and $\abs{\iota_{\bar a}(e_\mu \wedge e_\nu)}^2/\abs a^2$ is that entry. Hence, with the kernel \eqref{eq:flux-kernel} at $\mathsf T = N_q$ and the number \eqref{eq:flux-covariance} for $\mathsf I = \mathsf I_0$,
\begin{equation}\label{eq:free-flux-identification}
    \operatorname{Re}\bigl\langle\chi_0, \Pi_{\operatorname{ran}D,r}\chi_k\bigr\rangle = c^{\theta_r}(kt_{\mathrm b})\;,
    \qquad
    \mathfrak f_r \coloneqq \operatorname{Re}\bigl\langle\chi_0, \Pi_{D\ker Q,r}(\chi_3 - \chi_2)\bigr\rangle\;,
\end{equation}
where the index $r$ marks the operator on the component $r$; $\mathfrak f_r$ is the \emph{fluctuation part} of the third difference. In the language of Gaussian fields, Lemma~\ref{lem:flux-splitting}(iii) says that the fluxes of the free field, with covariance $\Pi_{\operatorname{ran}D}$, are the sum of the fluxes of the block field, with covariance $\Pi_{D\operatorname{ran}H}$, and of an independent field with the covariance given the block averages.

To bound $\mathfrak f_r$ we need pointwise values of a field which is known through an equation and through local $L^2$-norms. The next lemma provides them; it is a discrete interior estimate, uniform in the spacing.

\begin{lemma}[An Interior Bound]\label{lem:interior-bound}
On $\varepsilon\mathbb Z^4$, $\varepsilon \leq 1/16$, let $\nabla_\kappa f(x) \coloneqq \varepsilon^{-1}(f(x + \varepsilon e_\kappa) - f(x))$, $\nabla_\kappa^*$ its adjoint for $\langle f, h\rangle \coloneqq \varepsilon^4\sum_x\bar f(x)h(x)$, $\Delta \coloneqq \sum_\kappa\nabla_\kappa^*\nabla_\kappa$, $B_\rho$ the Euclidean ball of radius $\rho$ about a point $x_0$, and $\norm f_{B_\rho} \coloneqq \norm{f\mathbf 1_{B_\rho}}$. There is $C_{\mathrm{LR}}$, a number, such that if $f$ and $g_1, \dots, g_4$ are complex functions with $\Delta f = \sum_\kappa c_\kappa\nabla_\kappa g_\kappa$ on $B_1$, for constants $\abs{c_\kappa} \leq 1$, then $\abs{f(x)} \leq C_{\mathrm{LR}}(\norm f_{B_1} + \sup_{B_1}\abs g)$ for $x \in B_{1/2}$.
\end{lemma}

\begin{proof}
Let $G$ be the Green's function of the simple random walk on $\mathbb Z^4$, with $\mathcal LG = -\mathbf 1_{\{0\}}$ for the generator $\mathcal L$, and $\mathcal G(z) \coloneqq \frac18\varepsilon^{-2}G(z/\varepsilon)$. Since $\Delta$ acts on functions of $y/\varepsilon$ as $-8\varepsilon^{-2}\mathcal L$, $\Delta_y\mathcal G(x - y) = \varepsilon^{-4}\mathbf 1_{\{ y = x \}}$, and $h(x) = \varepsilon^4\sum_y\mathcal G(x - y)(\Delta h)(y)$ for finitely supported $h$. By \cite[Theorem~4.3.1]{LawlerLimic}, $G(z) = a_4\abs z^{-2} + O(\abs z^{-4})$, so that $\abs{G(z)} \leq c\abs z^{-2}$ and, by the mean value theorem for the leading term, $\abs{G(z + e_\kappa) - G(z)} \leq c\abs z^{-3}$ for $z \neq 0$, while $G(0)$ is finite; hence $\abs{\mathcal G(z)} \leq c(\abs z \vee \varepsilon)^{-2}$ and $\abs{\nabla_\kappa\mathcal G(z)}, \abs{\nabla_\kappa^*\mathcal G(z)} \leq c(\abs z \vee \varepsilon)^{-3}$, with numbers $c$. Extend $f$ and $g$ by $0$ off $B_1$, let $\psi \in C^\infty(\mathbb R^4)$ with $\psi = 1$ on $B_{3/4}$ and $\operatorname{supp}\psi \subset B_{7/8}$, and let $x \in B_{1/2}$. Then $f(x) = (\psi f)(x)$ is the sum of $\varepsilon^4\sum_y\mathcal G(x - y)\psi(y)(\Delta f)(y)$ and $\varepsilon^4\sum_y\mathcal G(x - y)([\Delta, \psi]f)(y)$, in which only values in $B_{7/8+2\varepsilon} \subset B_1$ enter. Summing by parts, the first sum is $\sum_\kappa c_\kappa\varepsilon^4\sum_y\nabla_\kappa^*(\psi\mathcal G(x - \cdot))(y)g_\kappa(y)$, and $\abs{\nabla_\kappa^*(\psi\mathcal G(x - \cdot))(y)} \leq c'(\abs{x - y} \vee \varepsilon)^{-3}$ on $B_1$, whose $\varepsilon^4$-sum over $B_1$ is at most a number $c''$; so the first sum is at most $4c''\sup_{B_1}\abs g$. In the second, $[\Delta, \psi]h(y) = -\varepsilon^{-2}\sum_\kappa\sum_\pm(\psi(y \pm \varepsilon e_\kappa) - \psi(y))h(y \pm \varepsilon e_\kappa)$ vanishes unless $\frac34 - \varepsilon \leq \abs{y - x_0} < \frac78 + \varepsilon$, so that $\abs{x - y} \geq \frac18$ there, and $[\Delta, \psi]^* = -[\Delta, \psi]$. With $\psi(y \pm \varepsilon e_\kappa) - \psi(y) = \pm\varepsilon\partial_\kappa\psi(y) + r_\pm$, $\abs{r_\pm} \leq \frac12\varepsilon^2\norm{\partial^2\psi}_\infty$, we get $\abs{[\Delta, \psi]\mathcal G(x - \cdot)(y)} \leq c'''$, a number, so the second sum is at most $c'''\varepsilon^4\sum_{B_1}\abs f \leq c'''\abs{B_2}^{1/2}\norm f_{B_1}$ by the Cauchy--Schwarz inequality.
\end{proof}

The fluctuation part is a covariance of Ba{\l}aban's constrained Gaussian field, and its decay is inherited from his bounds on the propagators.

\begin{lemma}[The Fluctuation Part]\label{lem:fluctuation-part}
There are $C_7$ and $\delta_2 > 0$, depending on $d$, $L$, and $G$, such that at every flat field and for every component $r$, $\abs{\mathfrak f_r} \leq C_7(s/t)^2e^{-\delta_2t_{\mathrm b}/2}$.
\end{lemma}

\begin{proof}
We work at $U_{\mathrm s}$ in the gauge of Lemma~\ref{lem:flat-fields}(i), with the sequence $\Omega_j = \mathbb T$ of the term $\ell_0$ and the dressed average, so that $Q = Q^{\mathrm{dr}}$ by Lemma~\ref{lem:flat-fields}(iii); distances are in units of the blocks of $\mathbb T_{j_q}$, where the distance of \cite[(2.46), p.~231]{BalabanPRT2} is at least the $\ell^1$-distance of the blocks.

\emph{Ba{\l}aban's constrained covariance.} The cubes of the class \cite[(3.35), p.~396]{BalabanPBG} have side $O(1)M < N_q$, so on each of them $U_{\mathrm s}$ is a gauge transform of $1$, and it satisfies \cite[(3.35)--(3.36)]{BalabanPBG} with $A = 0$. By Supplement, Lemma~17.7(c), Ba{\l}aban's operators and \cite[Theorems~3.1, 3.3, 3.13]{BalabanPBG} are available for it with constants of $d$, $L$, and $G$, and Supplement, Lemma~17.8, gives \cite[(3.132), p.~422]{BalabanPBG}. Since $U_{\mathrm s}$ is flat, the form $\Delta'$ and the current $J$ of \cite[(3.10)--(3.11), p.~392]{BalabanPBG} vanish, hence so do $\Delta'_\pi$, $\Delta^{(2)}$, and $\Delta^{(2)}_\pi$ by \cite[(3.120), (3.134), (3.135), pp.~419--422]{BalabanPBG}; by \cite[(3.119)]{BalabanPBG} and $D\partial_U = 0$, $\Delta_\pi = D^*D$; and by \cite[(3.122), (3.128)]{BalabanPBG}, with $a = 1$, $G_1 = G = (D^*D + \partial_UR\partial_U^* + Q^*Q)^{-1}$, where $R$ is Ba{\l}aban's projection on the $0$-forms. The covariance $\mathfrak G$ of the Gaussian field constrained by $QA = 0$ and $R\partial_U^*A = 0$ \cite[(3.148), p.~425]{BalabanPBG} is then
\begin{equation}\label{eq:constrained-covariance}
    \mathfrak G = G - G\partial_UR\partial_U^*G - GQ^*(QGQ^*)^{-1}QG
\end{equation}
by \cite[(3.153), p.~426]{BalabanPBG}, and it has the bounds of \cite[Theorem~3.3]{BalabanPBG} except the one with the covariant Laplacian \cite[Theorem~3.13, p.~426]{BalabanPBG}.

\emph{The flux covariance.} For a $1$-form $\lambda$ let $v \coloneqq \mathfrak G\lambda$ and $w \coloneqq (QGQ^*)^{-1}QG\lambda$. Then $v = G(\lambda - \partial_UR\partial_U^*G\lambda - Q^*w)$, and $Qv = 0$, $R\partial_U^*v = 0$, because $\mathfrak G = \mathfrak PG_1$ with $\mathfrak P$ the projection onto $S \coloneqq \{ QA = 0, R\partial_U^*A = 0 \}$ \cite[(3.153); p.~425]{BalabanPBG}. Applying $G^{-1}$ we obtain
\begin{equation}\label{eq:constrained-equation}
    D^*Dv = \lambda - \partial_UR\partial_U^*G\lambda - Q^*w\;,
\end{equation}
and for $z \in S$ it gives $\langle D^*Dv - \lambda, z\rangle = -\langle\partial_U^*G\lambda, R\partial_U^*z\rangle - \langle w, Qz\rangle = 0$, $R$ being an orthogonal projection. Moreover $D(S) = D(\ker Q)$: for $z \in \ker Q$, the form $z' \coloneqq z - \partial_UG'R\partial_U^*z$ has $Qz' = 0$, because $Q\partial_UG'R = \bar DQ'G'R = 0$ \cite[p.~426, l.~17]{BalabanPBG}, $R\partial_U^*z' = 0$, because $z'$ is $z$ transformed to the gauge $R\partial_U^*A = 0$ \cite[p.~419, ll.~26--27]{BalabanPBG}, and $Dz' = Dz$. For $\lambda \coloneqq D^*\phi$ we therefore have $Dv \in D(\ker Q)$ and $Dv - \phi \perp D(\ker Q)$, that is,
\begin{equation}\label{eq:constrained-flux}
    D\mathfrak GD^* = \Pi_{D\ker Q}\;,
    \qquad
    \mathfrak f_r = \operatorname{Re}\bigl\langle\chi_0, D\mathfrak GD^*(\chi_3 - \chi_2)\bigr\rangle_r\;,
\end{equation}
all these operators commuting with the decomposition into components.

\emph{The local equation.} Let $x_0$ be the base point of a plaquette of $\mathsf I_0 \cup \mathsf I_2 \cup \mathsf I_3$, $B_\rho$ the ball of radius $\rho$ about it, and $\phi$ a $2$-form which vanishes on $B_2$. On $B_3$ we have $U_{\mathrm s} = 1$, so $D = d$ and $\partial_U = \nabla$ there. For $F \coloneqq D\mathfrak GD^*\phi$, $dF = 0$ on $B_2$, and \eqref{eq:constrained-equation} with $\lambda = D^*\phi = 0$ on $B_{3/2}$ gives $d^*F = -Q^*w + \nabla(\cdot)$ there, hence $dd^*F = -dQ^*w$. On $2$-forms $dd^* + d^*d = \sum_\kappa\nabla_\kappa^*\nabla_\kappa$ componentwise, because the forward and backward differences commute, so $\Delta F_{\mu\nu} = -\nabla_\mu(Q^*w)_\nu + \nabla_\nu(Q^*w)_\mu$ on $B_1$, and Lemma~\ref{lem:interior-bound}, for each of the six components of each component $r$, gives
\begin{equation}\label{eq:local-flux-bound}
    \abs{F(x_0)} \leq \sqrt6\,C_{\mathrm{LR}}\Bigl(\norm F_{B_1} + \sup_{B_1}\abs{Q^*w}\Bigr)\;.
\end{equation}

\emph{The multipliers.} The average $Q$ is a mean of line integrals over one block length \cite[(1.18), p.~20]{BalabanPRT1}, so for a bond $b$ of $\mathcal B$ and the unit form $e_b$ at $b$, $\abs{Q^*e_b} \leq 1$, $Q^*e_b$ is supported on the fine bonds between the blocks of $b_-$ and $b_+$, $\abs{QA(b)} \leq \sup\abs A$, and $\abs{QA(b)} \leq \sqrt2\norm A_{\tilde\Delta(b)}$, with $\tilde\Delta(b)$ the two blocks which it reads. On $1$-forms the components of the covariant gradient $\nabla_U$ of \cite[(3.39), p.~397]{BalabanPBG} are $D_\mu A_\nu$ and $(DA)_{\mu\nu} = D_\mu A_\nu - D_\nu A_\mu$, so $\abs{DA} \leq 2\abs{\nabla_UA}$ pointwise. Let $\delta_2 \coloneqq \frac12\min(\delta_0, \delta_1)$ with $\delta_0$ of \cite[Theorem~3.1]{BalabanPBG} and $\delta_1$ of \eqref{eq:pbg-3132-local} below, and $C_\delta \coloneqq (2/(1 - e^{-\delta_2}))^4 \geq \sum_{y \in \mathbb Z^4}e^{-\delta_2\abs y_1}$. At one scale \cite[(3.132)]{BalabanPBG} reads
\begin{equation}\label{eq:pbg-3132-local}
    \abs{(QGQ^*)^{-1}(b, b')} \leq O(1)e^{-\delta_1d(b,b')}\;.
\end{equation}
For $\phi = \chi_k$ we have $QGD^*\phi(b) = \langle\chi_k, DGQ^*e_b\rangle$, and \eqref{eq:pbg-3132-local} together with \cite[(3.42)]{BalabanPBG} for $G$ \cite[Theorem~3.3]{BalabanPBG} gives $\abs{w(y)} \leq C_2\abs{\mathsf I_k}e^{-\delta_2d(y,\mathsf I_k)}$, with $d(y, \mathsf I_k)$ the distance of $y$ from the blocks which meet $\mathsf I_k$, the convolution of the two exponentials being bounded by $e^{-\delta_2d}C_\delta$. For $\phi$ supported in $B_1(x)$ with $\norm\phi \leq 1$ the same argument with \cite[(3.46)]{BalabanPBG} for $G$ and $\abs{QA(b)} \leq \sqrt2\norm A_{\tilde\Delta(b)}$ gives $\abs{w_\phi(y)} \leq C_4e^{-\delta_2d(y,B_1(x))}$; here and below the $C_i$ depend on $d$, $L$, and $G$.

\emph{The bound.} Fix $k \in \{ 2, 3 \}$, the base point $x$ of a plaquette of $\mathsf I_0$, and $F \coloneqq D\mathfrak GD^*\chi_k$. Then $\norm F_{B_1(x)}$ is the supremum of $\abs{\langle\phi, F\rangle}$ over the $\phi$ supported in $B_1(x)$ with $\norm\phi \leq 1$, and $\langle\phi, F\rangle = \langle D\mathfrak GD^*\phi, \chi_k\rangle = \eta^2\sum_{p \in \mathsf I_k}F_\phi(p)$ with $F_\phi \coloneqq D\mathfrak GD^*\phi$, $\mathfrak G$ being symmetric. The blocks which meet $B_1(x)$ or $B_1(x')$, $x' \in \mathsf I_k$, and their neighbours are at distance at least $(k - 1)t_{\mathrm b} - 5 \geq t_{\mathrm b} - 5$; so \eqref{eq:local-flux-bound} at $x'$, with \cite[(3.46)]{BalabanPBG} for $\mathfrak G$ and the bound on $w_\phi$, gives $\abs{F_\phi(x')} \leq C_5e^{-\delta_2(t_{\mathrm b}-5)}$, hence $\norm F_{B_1(x)} \leq C_5\abs{\mathsf I_k}e^{-\delta_2(t_{\mathrm b}-5)}$. Then \eqref{eq:local-flux-bound} at $x$, with the bound on $w$, gives $\abs{F(x)} \leq C_6\abs{\mathsf I_k}e^{-\delta_2(t_{\mathrm b}-5)}$, and $\abs{\langle\chi_0, F\rangle} \leq C_6\abs{\mathsf I_0}\abs{\mathsf I_k}e^{-\delta_2(t_{\mathrm b}-5)}$. Since $\abs{\mathsf I_a} = s_{\mathrm b}t_{\mathrm b} = (s/t)t_{\mathrm b}^2$ and $y^4e^{-\delta_2y/2} \leq (8/(e\delta_2))^4$, the two values of $k$ give the claim with $C_7 \coloneqq 2C_6e^{5\delta_2}(8/(e\delta_2))^4$.
\end{proof}

The coarse form $K$ has a kernel, and the following estimate controls how far a coarse field orthogonal to the gauge directions can vary without costing action.

\begin{lemma}[Coercivity of the Block Field]\label{lem:coarse-coercivity}
At a flat field, in a component $r$, let $\bar D$ be the covariant derivative of the $0$-forms of $\mathbb T_{j_q}$ and $\nabla_\mu$ the difference of a form on $\mathcal B$ in the direction $\mu$, acting on each of its components.
\begin{enumerate}[label=(\roman*)]
\item $\operatorname{ran}K \subset \ker\bar D^*$.
\item Every form $u$ on $\mathcal B$ with $\bar D^*u = 0$ has $\norm{\nabla_\mu u} \leq 4\norm{DHu}$.
\end{enumerate}
\end{lemma}

\begin{proof}
By Lemma~\ref{lem:flat-fields}(ii) and (iii) we may lift the component to quasi-periodic forms, on which $D = d$, $Q = Q_{K'}$ with centred blocks, and the products are taken over one period. (i) For a $0$-form $\omega$ of $\mathbb T_{j_q}$ let $\tilde\omega$ be constant on each block with $Q'_{K'}\tilde\omega = \omega$. By \cite[(1.20), p.~20]{BalabanPRT1}, $Q\partial\tilde\omega = \bar D\omega$, and $d\partial\tilde\omega = 0$; since the curvature of a minimiser is unique, $DH\bar D\omega = 0$. Hence $K\bar D = 0$ and $\bar D^*K = 0$.

(ii) Put $A \coloneqq Hu$ and $F \coloneqq dA$, so that $u = QA$. By \cite[(1.18)]{BalabanPRT1}, $u_b = \sum_{x \in B(b_-)}\eta^4\int_{[x,x+e_\nu]}A$ for $b = \langle b_-, b_- + e_\nu\rangle$, the line integral being $\eta$ times the sum over the fine bonds. For $\nu \neq \mu$ the lattice Stokes formula on the square $S_x$ spanned by $[x, x + e_\nu]$ and $[x, x + e_\mu]$ gives $\int_{[x,x+e_\nu]+e_\mu}A - \int_{[x,x+e_\nu]}A = \psi(x + e_\nu) - \psi(x) + \Phi_F(S_x)$, with $\psi(z) \coloneqq \int_{[z,z+e_\mu]}A$ and $\Phi_F(S) \coloneqq \eta^2\sum_{q \subset S}F_{\mu\nu}(q)$; for $\nu = \mu$ it holds with $\Phi_F \coloneqq 0$. The translation by $e_\mu$ maps $B(b_-)$ onto $B(b_- + e_\mu)$ and that by $e_\nu$ maps it onto $B(b_+)$, so averaging gives
\begin{equation}\label{eq:difference-stokes}
    \nabla_\mu u = \bar D(Q'_{K'}\psi) + \mathcal SF\;,
    \qquad
    (\mathcal SF)_b \coloneqq \sum_{x \in B(b_-)}\eta^4\Phi_F(S_x)\;.
\end{equation}
If $\bar D^*u = 0$, then $\bar D^*\nabla_\mu u = \nabla_\mu\bar D^*u = 0$, so $\norm{\nabla_\mu u}^2 = \langle\nabla_\mu u, \mathcal SF\rangle$ and $\norm{\nabla_\mu u} \leq \norm{\mathcal SF}$. A fine plaquette of the plane of $e_\mu$ and $e_\nu$ lies in $S_x$ for at most $\eta^{-2}$ points $x$ of $B(b_-)$, and all the $S_x$ lie in $R_b \coloneqq B(b_-) + [0, 1]e_\mu + [0, 1]e_\nu$, which contains at most $4\eta^{-4}$ such plaquettes; so $\abs{(\mathcal SF)_b} \leq \eta^4\sum_{q \subset R_b}\abs{F_{\mu\nu}(q)} \leq 2(\eta^4\sum_{q \subset R_b}\abs{F_{\mu\nu}(q)}^2)^{1/2}$ by the Cauchy--Schwarz inequality. Each plaquette lies in $R_b$ for at most four bonds $b$ of each direction, so $\norm{\mathcal SF}^2 \leq 16\norm F^2$.
\end{proof}

\begin{lemma}[The Flux of the Block Field]\label{lem:flux-norm}
At a flat field, in each component, $\norm{\Pi_{D\operatorname{ran}H}\chi_0}^2 \leq C_V(s_{\mathrm b} + 2)(t_{\mathrm b} + 2)$, with $C_V$ depending on $d$, $L$, and $G$.
\end{lemma}

\begin{proof}
The norm is the supremum of $\abs{\langle\chi_0, F\rangle}$ over the $F \in D(\operatorname{ran}H)$ with $\norm F = 1$. Let $F = DHu$. By Lemma~\ref{lem:flux-splitting}(i), $D^*F = Q^*\sigma$ for some $\sigma$, and applying $H^*$, $\sigma = H^*D^*DHu = Ku$. Near a base point $x_0$ of a plaquette of $\mathsf I_0$ we have $D = d$, so $dF = 0$, $d^*F = Q^*Ku$, and $\Delta F = dQ^*Ku$ on $B_1(x_0)$, and Lemma~\ref{lem:interior-bound} gives $\abs{F(x_0)} \leq \sqrt6\,C_{\mathrm{LR}}(\norm F_{B_1(x_0)} + \sup_{B_1(x_0)}\abs{Q^*Ku})$. By \cite[(3.133), p.~422]{BalabanPBG}, with the factor of \cite[p.~422, ll.~32--33]{BalabanPBG} for the variational $H$, which is $1$ at one scale, and $\abs{DA} \leq 2\abs{\nabla_UA}$, $\abs{DHe_b(x)} \leq C_He^{-\delta_H\abs{y(x) - b_-}_1}$, with $\delta_H \coloneqq \frac12\delta_1$ and $y(x)$ the block of $x$. With the weighted norms $\norm F_{w,y}^2 \coloneqq \eta^4\sum_xe^{-\delta_H\abs{y(x) - y}_1}\abs{F(x)}^2$ and $C'_\delta \coloneqq (2/(1 - e^{-\delta_H}))^4$, the Cauchy--Schwarz inequality gives $\abs{(Ku)_b} = \abs{\langle DHe_b, F\rangle} \leq C_HC'^{1/2}_\delta\norm F_{w,b_-}$. A fine bond lies in the supports of at most two of the $Q^*e_b$, $B_1(x_0)$ meets these supports only for $\abs{b_- - y(x_0)}_1 \leq 5$, and it lies in the blocks with $\abs{y - y(x_0)}_1 \leq 4$, so that $\abs{F(x_0)} \leq C_8\norm F_{w,y(x_0)}$ with $C_8 \coloneqq \sqrt6\,C_{\mathrm{LR}}(e^{2\delta_H} + 2C_HC'^{1/2}_\delta e^{5\delta_H})$. Now $\langle\chi_0, F\rangle = \eta^2\sum_{p \in \mathsf I_0}F_{\mu\nu}(p)$, and $\mathsf I_0$ meets at most $(s_{\mathrm b} + 2)(t_{\mathrm b} + 2)$ blocks, in each of which it has at most $\eta^{-2}$ plaquettes; by the Cauchy--Schwarz inequality over these blocks and $\sum_ye^{-\delta_H\abs{y(x) - y}_1} \leq C'_\delta$, $\abs{\langle\chi_0, F\rangle} \leq C_8\sum_y\norm F_{w,y} \leq C_8C'^{1/2}_\delta((s_{\mathrm b} + 2)(t_{\mathrm b} + 2))^{1/2}\norm F$, which gives the claim with $C_V \coloneqq C_8^2C'_\delta$.
\end{proof}

We can now compute the number \eqref{eq:c-lambda} at a flat field. The regularisation enters through the operator $P$ below, which is positive and is small on the fluxes of the block field because these vary slowly along the translations by whole blocks.

\begin{proposition}[The Coarse Covariance at a Flat Field]\label{prop:flat-covariance}
Let $V_0$ be flat, $0 < \lambda \leq 1$, and, in the gauge of Lemma~\ref{lem:flat-fields}(i), $P \coloneqq \Pi_{D\operatorname{ran}H} - DH(K + \lambda)^{-1}H^*D^*$. Then:
\begin{enumerate}[label=(\roman*)]
\item with $C_7$, $\delta_2$ of Lemma~\ref{lem:fluctuation-part} and $C_V$ of Lemma~\ref{lem:flux-norm},
\begin{equation}\label{eq:flat-covariance}
    c_\lambda(V_0) \geq d_{\mathfrak g}\Bigl(\frac{s^2t^3}{112\,T^5} - C_7\Bigl(\frac st\Bigr)^2e^{-\delta_2t_{\mathrm b}/2} - 10\pi^2C_V\lambda t_{\mathrm b}^2(s_{\mathrm b} + 2)(t_{\mathrm b} + 2)\Bigr)\;;
\end{equation}
\item $DH\lambda(K + \lambda)^{-1}u = PDHu$ for every form $u$ on $\mathcal B$, and in every component $r$ and for every $F \in D(\operatorname{ran}H)$,
\begin{equation}\label{eq:lambda-term-flat}
    \bigl\lvert\bigl\langle\chi_3 - \chi_2, P_rF\bigr\rangle\bigr\rvert \leq 2\pi t_{\mathrm b}\bigl(\lambda C_V(s_{\mathrm b} + 2)(t_{\mathrm b} + 2)\bigr)^{1/2}\norm F\;.
\end{equation}
\end{enumerate}
\end{proposition}

\begin{proof}
\emph{The exact form.} By Lemma~\ref{lem:flat-fields}(iv) and the gauge invariance of $c_\lambda$, $c_\lambda(V_0) = \sum_\alpha\langle\chi_0 \otimes t_\alpha, \Pi_\lambda((\chi_3 - \chi_2) \otimes t_\alpha)\rangle$ with $\Pi_\lambda \coloneqq DH(K + \lambda)^{-1}H^*D^*$. The sum over the orthonormal basis is a partial trace over $\mathfrak g$, which is the same in the basis $(E_r)$, and $D$, $Q$, $H$, $K$ commute with the decomposition into components, so $c_\lambda(V_0) = \sum_r\operatorname{Re}\langle\chi_0, \Pi_{\lambda,r}(\chi_3 - \chi_2)\rangle$, the total being real. By Lemma~\ref{lem:flux-splitting}(iii), $\Pi_{D\operatorname{ran}H} = DHK^+H^*D^*$, and $H^*D^*$ maps into $(\ker DH)^\perp = (\ker K)^\perp$, so that $P = DHg(K)H^*D^*$ with $g(x) \coloneqq \lambda x^{-1}(x + \lambda)^{-1}$ on $(\ker K)^\perp$ and $g \coloneqq 0$ on $\ker K$; in particular $0 \leq P \leq \Pi_{D\operatorname{ran}H}$. By Lemma~\ref{lem:flux-splitting}(iii) and \eqref{eq:free-flux-identification},
\begin{equation}\label{eq:flat-decomposition}
    c_\lambda(V_0) = \sum_r\bigl(c^{\theta_r}(3t_{\mathrm b}) - c^{\theta_r}(2t_{\mathrm b}) - \mathfrak f_r - \mathfrak l_r\bigr)\;,
    \qquad
    \mathfrak l_r \coloneqq \operatorname{Re}\bigl\langle\chi_0, P_r(\chi_3 - \chi_2)\bigr\rangle\;.
\end{equation}

\emph{The translation.} Let $\mathcal T$ be the translation by one block in the direction $\mu$, acting on the forms of the component $r$ with the phase of Lemma~\ref{lem:flat-fields}(ii) across the hyperplane. It is unitary and commutes with $D$ and $Q$, hence with $DH$, whose value at $u$ is the point of least norm of $\{ DA \,|\, QA = u \}$, and with $K$; and $\chi_k = \mathcal T^{kt_{\mathrm b}}\chi_0$, the rectangles not meeting the hyperplanes. With $J \coloneqq H^*D^*\chi_0$ and $v \coloneqq g(K)^{1/2}J$ we get $\langle\chi_0, P_r\mathcal T^{kt_{\mathrm b}}\chi_0\rangle = \langle v, \mathcal T^{kt_{\mathrm b}}v\rangle = \int e^{ikt_{\mathrm b}p}\,\dd\nu(p)$, with $\nu \geq 0$ the spectral measure of $\mathcal T$ in $v$ on $(-\pi, \pi]$, so that
\begin{equation}\label{eq:lambda-spectral}
    \mathfrak l_r = \int\bigl(\cos3t_{\mathrm b}p - \cos2t_{\mathrm b}p\bigr)\,\dd\nu(p)\;.
\end{equation}
Since $\cos3x - \cos2x = -2\sin\frac{5x}2\sin\frac x2$, $\abs{\cos3x - \cos2x} \leq \frac52x^2$; and $p^2 \leq \frac{\pi^2}4\abs{1 - e^{ip}}^2$ on $(-\pi, \pi]$. Hence $\abs{\mathfrak l_r} \leq \frac{5\pi^2}8t_{\mathrm b}^2\norm{(1 - \mathcal T)v}^2 = \frac{5\pi^2}8t_{\mathrm b}^2\norm{\nabla_\mu v}^2$.

\emph{The coercivity.} Since $g(K)$ vanishes on $\ker K$, $v \in \operatorname{ran}K \subset \ker\bar D^*$ by Lemma~\ref{lem:coarse-coercivity}(i), and Lemma~\ref{lem:coarse-coercivity}(ii) gives $\norm{\nabla_\mu v}^2 \leq 16\langle v, Kv\rangle = 16\langle J, Kg(K)J\rangle \leq 16\lambda\langle J, K^+J\rangle = 16\lambda\norm{\Pi_{D\operatorname{ran}H}\chi_0}^2$, because $Kg(K) = \lambda(K + \lambda)^{-1} \leq \lambda K^+$ on $(\ker K)^\perp$, and $\langle J, K^+J\rangle = \langle\chi_0, \Pi_{D\operatorname{ran}H}\chi_0\rangle$. With Lemma~\ref{lem:flux-norm}, $\abs{\mathfrak l_r} \leq 10\pi^2C_V\lambda t_{\mathrm b}^2(s_{\mathrm b} + 2)(t_{\mathrm b} + 2)$.

\emph{Conclusion of (i).} Lemma~\ref{lem:free-field} with $\mathsf T = N_q$, $\eta = L^{-K'}$, $t = t_{\mathrm b}$, and $s = s_{\mathrm b}$, whose conditions hold because $12t_{\mathrm b} < N_q$ and $3s_{\mathrm b} < N_q$, gives $c^{\theta_r}(3t_{\mathrm b}) - c^{\theta_r}(2t_{\mathrm b}) \geq s_{\mathrm b}^2t_{\mathrm b}^3/(112N_q^5) = s^2t^3/(112T^5)$ for every twist, the number being invariant under the common rescaling of the lengths; Lemma~\ref{lem:fluctuation-part} bounds $\mathfrak f_r$, and there are $d_{\mathfrak g}$ components.

(ii) On $(\ker K)^\perp$, $\lambda(K + \lambda)^{-1} = g(K)K$, and $DH$ vanishes on $\ker K$, which $\lambda(K + \lambda)^{-1}$ maps into itself; so $DH\lambda(K + \lambda)^{-1}u = DHg(K)Ku = DHg(K)H^*D^*DHu = PDHu$. Since $0 \leq P \leq 1$, the Cauchy--Schwarz inequality for the form $\langle\cdot, P\cdot\rangle$ gives $\abs{\langle\chi_3 - \chi_2, P_rF\rangle} \leq \langle\chi_3 - \chi_2, P_r(\chi_3 - \chi_2)\rangle^{1/2}\norm F$, and as in the translation step $\langle\chi_3 - \chi_2, P_r(\chi_3 - \chi_2)\rangle = \int\abs{e^{it_{\mathrm b}p} - 1}^2\,\dd\nu(p) \leq t_{\mathrm b}^2\int p^2\,\dd\nu(p) \leq \frac{\pi^2}4t_{\mathrm b}^2\norm{\nabla_\mu v}^2$, which the coercivity step bounds by $4\pi^2t_{\mathrm b}^2\lambda C_V(s_{\mathrm b} + 2)(t_{\mathrm b} + 2)$.
\end{proof}

\subsection{The Second Moments of the Fluxes}\label{sec:flux-moments}

We now compute $\Psi_3 - \Psi_2$ to the order $\bar g^2$. Near the fields of small action the measure $\langle\cdot\rangle$ has the density $\exp(-g_{K'}^{-2}\mathsf A + \mathcal A')$, whose second part is of the order one per block, and an integration by parts against the Haar measure turns $\langle X^\alpha\partial_w\mathsf A\rangle$ into $g_{K'}^2\langle\partial_wX^\alpha\rangle$, up to terms with the derivatives of $w$ and of $\mathcal A'$. If $w$ solves $(\Delta + \lambda)w = \nabla\Xi_\alpha$, then $\sum_\alpha\partial_wX^\alpha$ is the number $c_\lambda$ of Section~\ref{sec:flat-covariance}; but such a $w$ satisfies $\partial_w\mathsf A = \Xi_\alpha$ only approximately, and since $X^\alpha$ is of the order $\bar g$, the defect $\partial_w\mathsf A - \Xi_\alpha$ has to be small compared with $\bar g$ on all the fields of small action. At a flat field both $\nabla\mathsf A$ and $\Xi_\alpha$ vanish, and at the field $e^{iu}W_0$ near a flat $W_0$ the defect is, to first order in $u$, the number $-\lambda\langle(K + \lambda)^{-1}\nabla\Xi_\alpha(W_0), u\rangle$, which the regularisation keeps small. The difficulty is that a field of action of the order $g_{K'}^2$ need not lie within a distance of the order $g_{K'}$ of the flat fields: these form a variety with singular points, the holonomies with large centralisers, near which the Hessian $K = \Delta(W_0)$ has small eigenvalues in the directions which leave the variety. We meet it with a {\L}ojasiewicz inequality for commuting holonomies and a stratification of the variety by the dimension of the centraliser, which give at every field of small action a flat field and a lower bound on the Hessian there that are matched to each other: the field lies close to the flat field on the scale which the bound sets (Proposition~\ref{prop:near-flat}). The defect is then followed along a segment from a field which differs from the flat one only in the directions of the kernel of $K$ (Proposition~\ref{prop:ibp-error}).

\emph{The cutoff and the regularisation.} We keep the units and the notation of Section~\ref{sec:flat-covariance}, and we write $g \coloneqq g_{K'}$. The data parameter $\varepsilon_1$ of stage~$5$, at which \cite[Theorem~1, p.~279]{BalabanVAR} is applied, gives the constants $r_V \coloneqq C_1\varepsilon_1/(2N^{1/2})$ and $\varepsilon_V \coloneqq c_hB_5\varepsilon_1$, with $C_1$ and $B_5$ of \cite[(172)--(173), p.~305]{BalabanVAR}, and we write $\mathcal R$ for the set of the $V \in G^{\mathcal B}$ with $\abs{V(\partial p') - 1} < \varepsilon_1$ for every plaquette $p'$ of $\mathbb T_{j_q}$. Let $\varepsilon^\zeta \coloneqq 3(8NT_{\mathrm{ng}})^{1/2}g$ and $\mathcal O \coloneqq \{ V \in G^{\mathcal B} \,|\, \abs{V(\partial p') - 1} < \varepsilon^\zeta \text{ for every } p' \}$, an open subset of $\mathcal R$ for $\bar g$ small. On $\mathcal O$, \cite[Theorem~1]{BalabanVAR} with its $\varepsilon_1$ equal to $\varepsilon^\zeta$ gives the minimal configurations, with $\abs{U_{K'}(V)(\partial p) - 1} < B_3\varepsilon^\zeta\eta^2$ on the fine plaquettes. Since $B_3\varepsilon^\zeta < \varepsilon_{K'}$ for $\bar g$ small, by $T_{\mathrm{ng}} \leq C(\log\bar g^{-2})^{4r+1}$, $g \leq \rho^{1/2}\hat g_{K'}$, and $p_0 > 2r + \frac12$, every test function of the term $\ell_0$ passes on $\mathcal O$, as in the lower bound of the proof of Proposition~\ref{prop:large-field-weight}; hence $\chi_{K'}(\mathbb T) = 1$ and $\varrho_{\ell_0} = \exp A_{K'}(U_{K'}(V))$ on $\mathcal O$, and $\mathcal O \subset \operatorname{supp}\varrho_{\ell_0}$. We fix $\vartheta \in C^\infty(\mathbb R)$ with $0 \leq \vartheta \leq 1$, $\vartheta = 1$ on $(-\infty, 1]$, and $\vartheta = 0$ on $[2, \infty)$, and we put $\zeta(V) \coloneqq \vartheta(\mathsf A(V)/(T_{\mathrm{ng}}g^2))$ for $V \in \operatorname{supp}\varrho_{\ell_0}$ and $\zeta(V) \coloneqq 0$ otherwise. By Supplement, Lemma~17.6, with $T = 2T_{\mathrm{ng}}$, every $V \in \operatorname{supp}\varrho_{\ell_0}$ with $\mathsf A(V) \leq 2T_{\mathrm{ng}}g^2$ has $\abs{V(\partial p') - 1} \leq 2(4NT_{\mathrm{ng}})^{1/2}g < \varepsilon^\zeta$ for every $p'$. So $\zeta$ vanishes outside the compact set $\mathcal K$ of the fields with $\abs{V(\partial p') - 1} \leq 2(4NT_{\mathrm{ng}})^{1/2}g$ for every $p'$, which lies in $\mathcal O$; it is smooth on $G^{\mathcal B}$, because $\mathsf A$ is real analytic on $\mathcal O$ \cite[Proposition~9, p.~309]{BalabanVAR}; $\zeta = 1$ on $\mathcal V$; and $\operatorname{supp}\zeta \subset \mathcal V_3$. On $\operatorname{supp}\zeta$ the plaquettes satisfy $\norm{V(\partial p') - 1} \leq \varepsilon_\zeta \coloneqq 2c_h(4NT_{\mathrm{ng}})^{1/2}g$, and $\norm{X_{0,b}} \leq 2h^G$ by Section~\ref{sec:flat-covariance}.

We fix an exponent $p_\lambda$ with $14r + 2 < p_\lambda < q_0 - 11r - \frac12$, an interval which is not empty because $q_0 > 25r + \frac52$ (Appendix~\ref{app:constants}), and we put
\begin{equation}\label{eq:regularisation}
    \begin{gathered}
        \lambda \coloneqq N_q^{-2}\bigl(\log\bar g^{-2}\bigr)^{-p_\lambda}\;,
        \qquad
        F_\alpha \coloneqq \zeta X^\alpha\;,
        \\
        w_\alpha \coloneqq (\Delta(V) + \lambda)^{-1}\nabla\Xi_\alpha\;,
        \qquad
        \mathfrak e_\alpha \coloneqq \langle w_\alpha, \nabla\mathsf A\rangle - \Xi_\alpha\;,
    \end{gathered}
\end{equation}
the last two at the fields at which $\Delta(V) + \lambda$ is invertible. The sizes are those collected in the proof of Lemma~\ref{lem:stopping-sizes}: for $\bar g$ small and $K \geq K_{\mathrm w}$, $N_q$ and $P \coloneqq P_{K'}(R_4) = 2(4t_{\mathrm b} + s_{\mathrm b}) \leq 2N_q$ are at most $C(\log\bar g^{-2})^r$, $T_{\mathrm{ng}} \leq C(\log\bar g^{-2})^{4r+1}$, and $h^G \leq C\bar g(\log\bar g^{-2})^{4r+1/2}$; moreover $\bar g^2/2 \leq \hat g_{K'}^2 \leq g^2 \leq (1 + \beta_0)^2\bar g^2$, because $\hat g_{K'}^{-2} = \bar g^{-2} + b'j_q + C_W$, $t_{\mathrm b} \geq a_{K'} \geq (\log\bar g^{-2})^r$, and the reduced radii $r_{i,K'} \coloneqq \alpha_{i,K'}/(8B_3)$ of Supplement, Section~15, satisfy $r_{i,K'} \geq c\,\bar g(\log\bar g^{-2})^{q_0}$ by Definition~\ref{def:tube-parameters}. Here and below $C$, $c$, and $m$ are positive constants which depend on $q$, $t/u$, $s/u$, $T/u$, $\pi$, $p_\lambda$, and the constants of the construction, and we write $X \preceq \bar g^a$ if $\abs X \leq C\bar g^a(\log\bar g^{-2})^m$; if $X \preceq \bar g^{a+b}$ with $b > 0$, then $\abs X \leq \bar g^{a+b/2}$ for $\bar g$ small.

The minimal configuration depends analytically on its data, and we need this in a quantitative form, at complex data and with the gauge transformation which returns it to Ba{\l}aban's representative.

\begin{lemma}[The Moved Minimiser]\label{lem:moved-minimiser}
Let $V_1 \in \mathcal R$, and let $U_1$ be a minimal configuration for $V_1$.
\begin{enumerate}[label=(\roman*)]
\item There are maps $B \mapsto \mathcal H(B)$, into the $\mathfrak g^{\mathbb C}$-valued $1$-forms on $\mathbb T_K$, and $B \mapsto w_B$, into the $G^{\mathbb C}$-valued functions on the sites of $\mathbb T_K$, analytic on the set of the $\mathfrak g^{\mathbb C}$-valued forms $B$ on $\mathcal B$ with $\norm B_\infty < r_V$, with $\mathcal H(0) = 0$, $w_0 = 1$,
\begin{equation}\label{eq:moved-minimiser-bounds}
    \norm{\mathcal H(B)_b} < \varepsilon_V\;,
    \qquad
    \norm{(\nabla^\eta\mathcal H(B))(b)} < \varepsilon_V\;,
    \qquad
    \norm{\log w_B(x)} \leq 72d\,\varepsilon_V\;,
\end{equation}
with $\nabla^\eta$ the covariant lattice derivative at $U_1$ of \cite[(19), p.~281]{BalabanVAR} in units of the blocks, with the same bound $\varepsilon_V$ for the second derivative $D^{\eta*}_{U_1}D^\eta_{U_1}\mathcal H(B)$ of \cite[(19)]{BalabanVAR}, which \cite[Proposition~9, p.~309]{BalabanVAR} bounds together with the first two, and such that, for real $B$, $(e^{i\eta\mathcal H(B)}U_1)^{w_B}$ is a minimal configuration for the data $e^{iB}V_1$. Consequently, for real $B$, $\mathsf A(e^{iB}V_1) = A(e^{i\eta\mathcal H(B)}U_1)$, and the holonomy at a site $y$ of $\mathbb T_{j_q}$ of a closed path of $\mathbb T_K$ based at $y$ is $w_B(y)$ times its holonomy at $e^{i\eta\mathcal H(B)}U_1$ times $w_B(y)^{-1}$.
\item For $\norm B_\infty < r_V$, $\norm{\mathcal H(B)_b}, \norm{(\nabla^\eta\mathcal H(B))(b)}, \norm{(D^{\eta*}_{U_1}D^\eta_{U_1}\mathcal H(B))(b)} \leq \varepsilon_V\norm B_\infty/r_V$ and $\norm{\log w_B(x)} \leq 72d\,\varepsilon_V\norm B_\infty/r_V$. For $\norm B_\infty \leq r_V/2$, every $\mathfrak g^{\mathbb C}$-valued form $v$ on $\mathcal B$, and $m \geq 1$, the derivatives $\partial^m_vf(B) \coloneqq (\dd/\dd z)^mf(B + zv)|_{z=0}$ satisfy $\norm{\partial^m_v\mathcal H(B)_b}, \norm{\partial^m_v(\nabla^\eta\mathcal H)(B)(b)} \leq m!\,\varepsilon_V(2\norm v_\infty/r_V)^m$ and $\norm{\partial^m_v\log w_B(x)} \leq 72d\,m!\,\varepsilon_V(2\norm v_\infty/r_V)^m$.
\end{enumerate}
\end{lemma}

\begin{proof}
(i) The data $V_1$ satisfy \cite[(7), p.~278]{BalabanVAR} with the parameter $\varepsilon_1$, and $\norm B_\infty < r_V$ gives $\abs B \leq N^{1/2}\norm B < C_1\varepsilon_1/2$, hence $\abs{e^{iB} - 1} \leq e^{\abs B} - 1 < C_1\varepsilon_1$, by \cite[(20), p.~21]{BalabanAVG} and the submultiplicativity of $\abs\cdot$; so $V' \coloneqq e^{iB}$ satisfies \cite[(172), p.~305]{BalabanVAR}. By \cite[Proposition~9, p.~309]{BalabanVAR} and \cite[pp.~305--307]{BalabanVAR}, the minimal configuration for $V'V_1$, transformed to the Landau gauge relative to $U_1$, is $e^{i\eta\mathcal H(B)}U_1$ with $\mathcal H$ analytic in the $\mathfrak g^{\mathbb C}$-valued $B$ satisfying \cite[(172)]{BalabanVAR} \cite[p.~307, ll.~1--2]{BalabanVAR}, and $\mathcal H$ satisfies \cite[(19)--(21), p.~281]{BalabanVAR} with $\varepsilon_2 = B_5\varepsilon_1$. On the term $\ell_0$ every set $\Omega_j$ of the determining sequence is the torus, so the conditions \cite[(19)]{BalabanVAR} for $j = K'$, where $L^{K'}\eta = 1$, read $\abs{\mathcal H} < B_5\varepsilon_1$ and $\abs{\nabla^\eta\mathcal H} < B_5\varepsilon_1$, which give the first two bounds after multiplication by $c_h$. The configuration in the Landau gauge is carried back to the axial gauge by a gauge transformation which is analytic in $\mathcal H$, hence in $B$ \cite[p.~307, ll.~2--5]{BalabanVAR}, and which is normalised by the condition that its averaged transformation be $1$ on the sets $\Lambda_j$ \cite[p.~281, ll.~1--3]{BalabanVAR}; by the uniqueness in Supplement, Lemma~17.14, applied at the $G$-valued base $\mathbf U = U_1$, with $A' = 0$ and $X = \mathcal H(B)$, so that $\varepsilon_A = 0$ and $\varepsilon_X = B_5\varepsilon_1$, this is the transformation $w_B$ of that lemma, and $\abs{\log w_B} \leq 72dB_5\varepsilon_1$. For real $B$ the minimal configurations for $e^{iB}V_1$ form one orbit, the action is gauge invariant, and a gauge transformation $w$ changes the holonomy at $y$ of a closed path based at $y$ by the conjugation with $w(y)$.

(ii) The function $z \mapsto \mathcal H(zB/\norm B_\infty)_b$ is analytic on $\abs z < r_V$, bounded by $\varepsilon_V$, and vanishes at $z = 0$, so Schwarz's lemma, applied to its pairings with the functionals of norm one, gives $\norm{\mathcal H(B)_b} \leq \varepsilon_V\norm B_\infty/r_V$; the same argument applies to $(\nabla^\eta\mathcal H(B))(b)$ and to $\log w_B(x)$. For $\norm B_\infty \leq r_V/2$ and $\abs z < r_V/(2\norm v_\infty)$ we have $\norm{B + zv}_\infty < r_V$, and Cauchy's estimate on this disc gives the bounds on the derivatives.
\end{proof}

The fluxes are logarithms of holonomies of the moved minimiser, and Cauchy's estimate in the data bounds their derivatives; the same argument applied to the action bounds the Hessian and its variation. We write $\partial_w\partial_uF(V) \coloneqq \partial_\sigma\partial_sF(e^{isu}e^{i\sigma w}V)|_{s=\sigma=0}$, so that $\langle u, \partial_w\nabla F(V)\rangle = \partial_w\partial_uF(V)$ with $\partial_w\nabla F(V) \coloneqq \frac{\dd}{\dd\sigma}\nabla F(e^{i\sigma w}V)|_{\sigma=0}$, and similarly $\partial_w\Delta(V) \coloneqq \frac{\dd}{\dd\sigma}\Delta(e^{i\sigma w}V)|_{\sigma=0}$. Let $\mathcal R'$ be the set of the $V_1 \in \mathcal R$ with $\norm{U_{K'}(V_1)(I_b) - 1} \leq r_G/2$ for $b \in \{ 0, 2, 3 \}$, a condition on the holonomy at the site $o_b$ of $\mathbb T_{j_q}$ which does not depend on the choice of the minimal configuration.

\begin{lemma}[Derivatives of the Fluxes and of the Action]\label{lem:flux-derivatives}
There are constants $C_\Xi$ and $C_\Delta$ of the construction such that the following hold for all real forms $u$, $w$ on $\mathcal B$.
\begin{enumerate}[label=(\roman*)]
\item At every $V_1 \in \mathcal R'$ the functions $X_{0,k}$, $k \in \{ 0, 2, 3 \}$, are real analytic, and $\norm{\partial_uX_{0,k}(V_1)} \leq C_\Xi P\norm u_\infty$, $\norm{\nabla X^\alpha(V_1)}, \norm{\nabla\Xi_\alpha(V_1)} \leq C_\Xi P$, and $\norm{\partial_w\nabla X^\alpha(V_1)}, \norm{\partial_w\nabla\Xi_\alpha(V_1)} \leq C_\Xi P^2\norm w_\infty$.
\item At every $V_1 \in \mathcal R$, $\norm{\Delta(V_1)} \leq C_\Delta N_q^4$ and $\norm{\partial_w\Delta(V_1)} \leq C_\Delta N_q^4\norm w_\infty$.
\item If $V_1 \in \mathcal R$ and $e^{-itw}V_1 \in \mathcal R$ for $0 \leq t \leq (2\mathsf A(V_1)/(C_\Delta N_q^4))^{1/2}$ and every $w$ with $\norm w = 1$, then $\norm{\nabla\mathsf A(V_1)} \leq (2C_\Delta N_q^4\mathsf A(V_1))^{1/2}$.
\end{enumerate}
\end{lemma}

\begin{proof}
(i) Let $U$ be a minimal configuration for $V_1$, and for complex $B$ with $\norm B_\infty < r_V$ let $H_k(B)$ be the holonomy of $\mathsf L_{0,k}$ at $e^{i\eta\mathcal H(B)}U$, with $\mathcal H$ of Lemma~\ref{lem:moved-minimiser} at the base $V_1$. The path $\mathsf L_{0,k}$ has $2kt_{\mathrm b} + 2t_{\mathrm b} + 2s_{\mathrm b} \leq P$ blocks, hence at most $P/\eta$ fine bonds, and each of its factors $e^{\pm i\eta\mathcal H_b}$ differs from $1$ by at most $e^{\eta\norm{\mathcal H_b}} - 1 \leq 2\eta\varepsilon_V\norm B_\infty/r_V$, by Lemma~\ref{lem:moved-minimiser}(ii). Expanding the product of the factors $(1 + a_l)U(b_l)^{\pm1}$ as in \eqref{eq:product-deviation}, with the unitary $U(b_l)$, gives $\norm{H_k(B) - H_k(0)} \leq e^{2P\varepsilon_V\norm B_\infty/r_V} - 1$. Let $\rho_P \coloneqq \min(r_V/2, r_V/(50P\varepsilon_V))$, so that $1/\rho_P \leq c_PP$ with $c_P \coloneqq (2 + 50\varepsilon_V)/r_V$. Since $H_k(0)$ is unitarily conjugate to $U(I_k)$ and $r_G \leq 1/10$, for $\norm B_\infty < \rho_P$ we get $\norm{H_k(B) - 1} \leq \frac1{20} + e^{1/25} - 1 \leq \frac1{10}$, so that $-i\log H_k(B)$ is analytic, with norm at most $\frac19$; and $\norm{\mathrm{Ad}(w_B(o))} \leq e^{144d\varepsilon_V} \leq 2$ for $\varepsilon_1$ small. By Lemma~\ref{lem:moved-minimiser}(i), $X_{0,k}(e^{iB}V_1) = \mathrm{Ad}(w_B(o))(-i\log H_k(B))$ for real $B$, and the right side is an analytic extension to $\norm B_\infty < \rho_P$, of norm at most $\frac14$. Hence $X^\alpha$ and $\Xi_\alpha$ are real analytic near $V_1$, with analytic extensions bounded by $\frac14$ and $\frac12$, since $\abs{t_\alpha \cdot Y} \leq \abs Y_2 \leq \norm Y$ for complex $Y$. Cauchy's estimate in $z$ on $\abs z < \rho_P/\norm u_\infty$, for the field $e^{izu}V_1$, gives $\norm{\partial_uX_{0,k}(V_1)} \leq \norm u_\infty/(4\rho_P)$, and with $\norm u_\infty \leq N^{1/2}\norm u$ the bounds on $\nabla X^\alpha$ and $\nabla\Xi_\alpha$. For the second derivatives, $e^{isu}e^{i\sigma w}V_1 = e^{iB(s,\sigma)}V_1$ with $B(s, \sigma)_b \coloneqq -i\log(e^{isu_b}e^{i\sigma w_b})$, which is analytic and has $\norm{B(s, \sigma)}_\infty \leq 2(\abs s\norm u_\infty + \abs\sigma\norm w_\infty)$ for $\abs s\norm u_\infty + \abs\sigma\norm w_\infty \leq \frac14$, by the series of the exponential and the logarithm; so Cauchy's estimate on the polydisc $\abs s\norm u_\infty, \abs\sigma\norm w_\infty \leq \rho_P/5$ gives $\abs{\partial_w\partial_u\Xi_\alpha(V_1)} \leq \frac{25}2\norm u_\infty\norm w_\infty/\rho_P^2$, and the same with $\frac{25}4$ for $X^\alpha$. The constant $C_\Xi \coloneqq 13N^{1/2}\max(c_P, c_P^2)$ covers all the bounds.

(ii) For real $s$ and $\sigma$ let $F(s, \sigma) \coloneqq \mathsf A(e^{isu}e^{i\sigma w}V_1)$, so that $\langle u, \Delta(V_1)u\rangle = \partial_s^2F(0, 0)$ and $\langle u, \partial_w\Delta(V_1)u\rangle = \partial_\sigma\partial_s^2F(0, 0)$. By Lemma~\ref{lem:moved-minimiser}(i), $F(s, \sigma) = A^{\mathbb C}(e^{i\eta\mathcal H(B(s,\sigma))}U)$ with $A^{\mathbb C}(U') \coloneqq \sum_p[1 - \frac12(\tr U'(\partial p) + \tr U'(\partial p)^{-1})]$, which equals the action at unitary configurations and is analytic. Let $p$ be a fine plaquette. Moving the transports of $U$ to the right, $(e^{i\eta\mathcal H}U)(\partial p)$ is the product of four exponentials $e^{i\eta A'_l}$, with $A'_l$ the transported values $\pm\mathcal H_{b_l}$, times $U(\partial p)$; by the Baker--Campbell--Hausdorff formula the four exponentials combine to $e^{iZ}$ with $Z$ equal to $\eta\sum_lA'_l$ up to $O(\eta^2\varepsilon_V^2)$, and $\eta\sum_lA'_l$ is $\eta^2$ times a combination of the covariant derivatives $\nabla^\eta\mathcal H$, up to a term bounded by $\norm{U(\partial p) - 1}\varepsilon_V\eta$. With \eqref{eq:moved-minimiser-bounds} and $\abs{U(\partial p) - 1} < B_3\varepsilon_1\eta^2$ \cite[Theorem~1]{BalabanVAR}, every plaquette of $e^{i\eta\mathcal H(B)}U$ is within $C_{\mathrm{pl}}\eta^2$ of $1$, $C_{\mathrm{pl}} \coloneqq 4\varepsilon_V + c_hB_3\varepsilon_1$, for $\varepsilon_1$ small. Its logarithm $Y_p$ has $\norm{Y_p} \leq 2C_{\mathrm{pl}}\eta^2 \leq 1$, and the term of $p$ in $A^{\mathbb C}$ is $\sum_{n \geq 1}(-1)^{n+1}\tr(Y_p^{2n})/(2n)!$, of modulus at most $\cosh\norm{Y_p} - 1 \leq \norm{Y_p}^2$. Summing over the $6N_q^4\eta^{-4}$ fine plaquettes, $\abs F \leq C_AN_q^4$ with $C_A \coloneqq 24C_{\mathrm{pl}}^2$ on the polydisc $\abs s\norm u_\infty, \abs\sigma\norm w_\infty \leq r_V/5$, and Cauchy's estimate there gives $\abs{\langle u, \Delta(V_1)u\rangle} \leq 50C_AN_q^4\norm u_\infty^2/r_V^2$ and $\abs{\langle u, \partial_w\Delta(V_1)u\rangle} \leq 250C_AN_q^4\norm u_\infty^2\norm w_\infty/r_V^3$. The operators are symmetric and $\norm u_\infty \leq N^{1/2}\norm u$, which gives (ii) with $C_\Delta \coloneqq 250NC_A\max(r_V^{-2}, r_V^{-3})$.

(iii) If $\mathsf A(V_1) = 0$, then $V_1$ is a minimum of $\mathsf A \geq 0$ and $\nabla\mathsf A(V_1) = 0$. Otherwise let $\norm w = 1$ and $t_0 \coloneqq (2\mathsf A(V_1)/(C_\Delta N_q^4))^{1/2}$. The second derivative of $t \mapsto \mathsf A(e^{-itw}V_1)$ is $\langle w, \Delta(e^{-itw}V_1)w\rangle \leq C_\Delta N_q^4$ by (ii), so that $0 \leq \mathsf A(e^{-it_0w}V_1) \leq \mathsf A(V_1) - t_0\langle w, \nabla\mathsf A(V_1)\rangle + \frac12t_0^2C_\Delta N_q^4$, that is, $\langle w, \nabla\mathsf A(V_1)\rangle \leq 2\mathsf A(V_1)/t_0 = (2C_\Delta N_q^4\mathsf A(V_1))^{1/2}$.
\end{proof}

The integration by parts itself is exact. On $\mathcal O$ we write $\mathcal A'(V) \coloneqq A_{K'}(U_{K'}(V)) + g^{-2}\mathsf A(V)$ for the logarithm of the density of $\ell_0$ without the Wilson action, and for $b \in \mathcal B$ and $\alpha \leq d_{\mathfrak g}$ we let $e_{b,\alpha}$ be the form equal to $t_\alpha$ at $b$ and to $0$ elsewhere, an orthonormal basis of $\mathfrak g^{\mathcal B}$.

\begin{lemma}[Integration by Parts]\label{lem:integration-by-parts}
Let $F$ be a $C^1$ function on $G^{\mathcal B}$ whose support lies in $\mathcal O$, and let $v$ be a $C^1$ map from a neighbourhood of $\operatorname{supp}F$ in $\mathcal O$ to $\mathfrak g^{\mathcal B}$, with the divergence $\operatorname{div}v \coloneqq \sum_{b,\alpha}\partial_{e_{b,\alpha}}(t_\alpha \cdot v_b)$. Then
\begin{equation}\label{eq:integration-by-parts}
    \bigl\langle F\langle v, \nabla\mathsf A\rangle\bigr\rangle = g^2\bigl\langle\langle v, \nabla F\rangle\bigr\rangle + g^2\bigl\langle F\operatorname{div}v\bigr\rangle + g^2\bigl\langle F\langle v, \nabla\mathcal A'\rangle\bigr\rangle\;.
\end{equation}
\end{lemma}

\begin{proof}
For a $C^1$ function $h$ on $G^{\mathcal B}$, $\int\partial_{e_{b,\alpha}}h(V)\,\dd V = \frac{\dd}{\dd s}\int h(e^{ise_{b,\alpha}}V)\,\dd V|_{s=0} = 0$, the Haar measure of the factor of $b$ being invariant under the left multiplication by $e^{ist_\alpha}$, and the derivative passing under the integral over the compact group. We apply this to $h \coloneqq (t_\alpha \cdot v_b)F\varrho_{\ell_0}$, which is $C^1$ on $G^{\mathcal B}$ when extended by $0$, because $\varrho_{\ell_0} = \exp(-g^{-2}\mathsf A + \mathcal A')$ is smooth on $\mathcal O$. Since $\partial_u$ is linear in $u$, $\sum_{b,\alpha}(t_\alpha \cdot v_b)\partial_{e_{b,\alpha}}f = \langle v, \nabla f\rangle$, and summing over $b$ and $\alpha$ gives $0 = \int\varrho_{\ell_0}[F\operatorname{div}v + \langle v, \nabla F\rangle + F\langle v, -g^{-2}\nabla\mathsf A + \nabla\mathcal A'\rangle]\,\dd V$. Dividing by $g^{-2}Z_{\ell_0}$ gives \eqref{eq:integration-by-parts}.
\end{proof}

The last term of \eqref{eq:integration-by-parts} contains the derivative of the effective action beyond the Wilson action. Its terms are of the order one per block, and their derivatives are bounded by Cauchy's estimate on the radius of analyticity of Ba{\l}aban's terms, which is larger than the size of the fields by a power of $\log\bar g^{-2}$.

\begin{lemma}[The Rest of the Density]\label{lem:density-rest}
Let $r_E \coloneqq \min(r_V/2, r_V\min(r_{0,K'}, r_{1,K'})/(8C_{Q^*}N^{1/2}\varepsilon_V))$, with the constant $C_{Q^*} = 2(1 + (2d + 1)L^{d-1}) \geq 1$ of Supplement, Section~15, and with $\varepsilon_V$ a common bound of $\mathcal H(B)$, of its first differences, and of $D^*D\mathcal H(B)$ on $\norm B_\infty < r_V$, in the units of \cite[(19), p.~281]{BalabanVAR}. There is a constant $M_E$ of the construction such that for $\bar g$ small, every $V \in \mathcal K$, and every real form $w$ on $\mathcal B$,
\begin{equation}\label{eq:density-rest}
    \bigl\lvert\langle w, \nabla\mathcal A'(V)\rangle\bigr\rvert \leq \frac{M_EN_q^4\norm w_\infty}{r_E}\;.
\end{equation}
\end{lemma}

\begin{proof}
Let $U \coloneqq U_{K'}(V)$, and for complex $t$ with $\abs t < r_E/\norm w_\infty$ let $U_t \coloneqq e^{i\eta\mathcal H(tw)}U$, with $\mathcal H$ of Lemma~\ref{lem:moved-minimiser} at the base $V$. On $\ell_0$ the function $\mathcal A'$ is, by \cite[(2.23), p.~258]{BalabanCRE} and Section~\ref{sec:identity}, with the coupling function $g_{K'}^2(x)$ of \cite[(2.24)]{BalabanCRE} equal to $g^2$ because $\phi_j = 1$ on the torus, the sum of the regular part $\mathbf E_{K'}$, the $\mathbf R$-part $\mathbf R_{K'}$, and the boundary part, each a gauge invariant function of the configuration, and of constants which do not depend on $V$. For real $t$ the configuration $U_t$ is a gauge transform of a minimal configuration for $e^{itw}V$, by Lemma~\ref{lem:moved-minimiser}(i), so that $\mathcal A'(e^{itw}V)$ is the sum of the three parts at $U_t$ and of the constants, and the parts at $U_t$ are analytic in $t$.

\emph{The class.} We show that $U_t$ lies in the regular subclass $\mathcal S_z$ of Supplement, Section~15, for every point $z$, at the reduced radii $r_{0,K'}$, $r_{1,K'}$ of the layer $K'$, the only layer on $\ell_0$. The configuration $U$ is $G$-valued with $\abs{U(\partial p) - 1} < B_3\varepsilon^\zeta\eta^2 \leq \frac12r_{0,K'}\eta^2$, and its current is at most $B_3\varepsilon^\zeta$ in units of the blocks by \cite[(2), (170), pp.~278, 305]{BalabanVAR}, for $\bar g$ small, because $\varepsilon^\zeta \preceq \bar g$ while $r_{0,K'} \geq c\bar g(\log\bar g^{-2})^{q_0}$. The field $A' \coloneqq \mathcal H(tw)$ is $\mathfrak g^{\mathbb C}$-valued with $\abs{A'}, \abs{\nabla^\eta A'} \leq N^{1/2}\varepsilon_V\abs t\norm w_\infty/r_V \leq \frac18\min(r_{0,K'}, r_{1,K'})/C_{Q^*}$ by Lemma~\ref{lem:moved-minimiser}(ii), and the same bound holds for $D^*DA'$, by the same Schwarz estimate applied to $D^*D\mathcal H$, which \cite[(19)]{BalabanVAR} bounds together with $\mathcal H$ and its first differences \cite[Proposition~9, p.~309]{BalabanVAR}, so that, by the argument of the proof of Lemma~\ref{lem:flux-derivatives}(ii), the plaquettes of $U_t$ are within $\frac12r_{0,K'}\eta^2 + \frac12r_{0,K'}\eta^2$ of $1$, and the current of $U_t$, which is the current of $U$ plus $D^*DA'$ plus terms of the second order in $A'$, is at most $B_3\varepsilon^\zeta + \frac18r_{0,K'}/C_{Q^*} + \frac18r_{0,K'}/C_{Q^*}$, which is below $r_{0,K'}/C_{Q^*}$ for $\bar g$ small because $\varepsilon^\zeta \preceq \bar g$. For real $t$ the restriction of a minimal configuration to the near domain $\square_0$ of a cube is, in a gauge on $\square_0$, the local minimiser of its own block data \cite[(3.18), p.~273]{BalabanRG1}; both sides of this identity, transported to $U_t$ by the gauge transformation of Lemma~\ref{lem:moved-minimiser}(i), are analytic in $t$, so it holds for complex $t$ by the identity theorem, in a complex gauge on $\square_0$.

\emph{The parts.} On $\ell_0$ every $\Lambda_j$ is the torus, so $\Lambda^0_j = \mathbb T^{(j)}$, the function $\phi_j$ of \cite[(2.25), p.~259]{BalabanCRE} is $1$, and the tents $h_z$ of \cite[p.~283]{BalabanCRE} sum to $1$; by \cite[(2.25)--(2.26)]{BalabanCRE},
\begin{equation}\label{eq:regular-part-points}
    \mathbf E_{K'}(U_t) = \sum_{j=1}^{K'}\sum_{z \in \mathbb T^{(j)}}\Bigl[\mathbf E^{(j)}(\mathbb T, U_t, z) - \mathbf E^{(j)}(\mathbb T, 1, z) - \beta_j(g_{j-1})A(h_z, U_t)\Bigr]\;.
\end{equation}
Supplement, Proposition~S.15(a), for the single family $\mathbf E^{(j)}$, that is with the second family $0$, $d = 1$ by \cite[(1.18), p.~263]{BalabanRG1}, and $N_j = 0$, bounds the bracket by $C_\alpha E_0(L^{j-K'})^{4+\beta'}$, since every $z$ lies in $\Omega_{K'} \setminus \Omega_{K'+1}$. The lattice $\mathbb T^{(j)}$ has $L^{4(K'-j)}N_q^4$ points, and $L^{\beta'} \geq 2$, so $\abs{\mathbf E_{K'}(U_t)} \leq C_\alpha E_0N_q^4\sum_{m \geq 0}L^{-\beta'm} \leq 2C_\alpha E_0N_q^4$. In the same way Supplement, Proposition~S.15(b), with the class norm $1$ by \cite[(2.31), p.~260]{BalabanCRE}, bounds the terms of \cite[(2.30)]{BalabanCRE} at a point of the scale $j$ by $C_\alpha g_j^{\kappa_0}(L^{j-K'})^4$, so that $\abs{\mathbf R_{K'}(U_t)} \leq C_\alpha N_q^4\sum_{j \leq K'}g_j^{\kappa_0} \leq C_\alpha N_q^4$, by $g_j^{-2} \geq g_{K'}^{-2} + b(K' - j) - C_W$ \eqref{eq:tube-window}, $\kappa_0 \geq 13$, and $\bar g$ small. The terms of the boundary part are analytic on the spaces of \cite[pp.~261--262]{BalabanCRE}, which contain $U_t$, with the bound \cite[(2.42)]{BalabanCRE}, and \cite[(2.47)--(2.48), p.~264]{BalabanCRE} bound their sum by $2B_\partial\sum_n\abs{\Gamma_n} = 2B_\partial N_q^4$. So $t \mapsto \mathcal A'(e^{itw}V)$ extends analytically to $\abs t < r_E/\norm w_\infty$, where it differs from a constant by at most $M_EN_q^4$, $M_E \coloneqq 2C_\alpha E_0 + C_\alpha + 2B_\partial$, and Cauchy's estimate at $t = 0$ gives \eqref{eq:density-rest}.
\end{proof}

We turn to the geometry of the fields of small action. Their plaquettes are close to $1$, so in a suitable gauge they are close to configurations which are constant off four hyperplanes, with values $h_1, \dots, h_4$ that nearly commute; how close the $h_\beta$ are to a commuting quadruple is a question about the real algebraic set of commuting quadruples, which {\L}ojasiewicz's inequality answers with an unknown exponent. The same inequality, applied to the eigenvalues of the operator $M(k)$ below on the strata of that set, bounds from below the twists of the components of a flat field which do not lie in the centraliser, and with them the spectral gap of $K$.

\begin{lemma}[{\L}ojasiewicz Inequalities for Commuting Holonomies]\label{lem:lojasiewicz}
Let $Z \subset G^4$ be the set of the commuting quadruples. For $k \in G^4$ let $\Phi(k) \coloneqq \max_{\beta<\gamma}\norm{k_\beta k_\gamma k_\beta^{-1}k_\gamma^{-1} - 1}$, let $M(k) \coloneqq \sum_\beta(\mathrm{Ad}(k_\beta) - 1)^*(\mathrm{Ad}(k_\beta) - 1)$ on $\mathfrak g$ with the product $X \cdot Y$, with the eigenvalues $0 \leq \mu_1(k) \leq \dots \leq \mu_{d_{\mathfrak g}}(k)$, and for a compact $S \subset G^4$ let $d_\infty(k, S) \coloneqq \min_{z \in S}\max_\beta\norm{k_\beta - z_\beta}$. Put $Z_{\geq 0} \coloneqq Z$ and $Z_{\geq j} \coloneqq \{ k \in Z \,|\, \mu_j(k) = 0 \}$ for $1 \leq j \leq d_{\mathfrak g}$. There are constants $C_{\mathrm{Lo}} > 0$, $\alpha_{\mathrm{Lo}} \in (0, 1]$, $\mathsf c_j > 0$, and $\varpi_j \geq 1$, depending on $G$ only, such that
\begin{enumerate}[label=(\roman*)]
\item $d_\infty(k, Z) \leq C_{\mathrm{Lo}}\Phi(k)^{\alpha_{\mathrm{Lo}}}$ for every $k \in G^4$;
\item $\mu_{j+1}(k) \geq \mathsf c_jd_\infty(k, Z_{\geq j+1})^{\varpi_j}$ for $0 \leq j < d_{\mathfrak g}$ and every $k \in Z_{\geq j}$.
\end{enumerate}
\end{lemma}

\begin{proof}
We use {\L}ojasiewicz's inequality in the form of \cite[Theorem~6.4, p.~34]{BierstoneMilman}: if $K$ is a subset of a real analytic manifold and $f, h \colon K \to \mathbb R$ are subanalytic functions with compact graphs and $f^{-1}(0) \subset h^{-1}(0)$, then $\abs{f} \geq c\abs h^\rho$ on $K$ for some $c, \rho > 0$. The manifold is $M_N(\mathbb C)^4$, a real vector space. The group $G$ is a closed subgroup of $U(N)$, hence a closed real analytic submanifold, so that $G^4$ and $Z = G^4 \cap \{ k_\beta k_\gamma = k_\gamma k_\beta \}$ are compact and semianalytic. By \cite[Remarks~3.11(2), p.~19]{BierstoneMilman}, $x \mapsto \min_{t \in T}f(x, t)$ is subanalytic if $T$ is compact and subanalytic and $f$ is continuous and subanalytic, and a maximum is $-\min(-\,\cdot\,)$. Now $\Phi(k)$ is the maximum, over $\beta < \gamma$ and over the pairs of unit vectors $x, y \in \mathbb C^N$, of the polynomial $\operatorname{Re}\langle y, (k_\beta k_\gamma k_\beta^*k_\gamma^* - 1)x\rangle$; in the same way $(k, z) \mapsto \max_\beta\norm{k_\beta - z_\beta}$ is subanalytic, and so is $d_\infty(\cdot, S)$ for compact subanalytic $S$. By the Courant--Fischer principle, $\mu_i(k)$ is the minimum over the orthonormal $i$-frames $(f_1, \dots, f_i)$ of $\mathfrak g$ of the maximum over the unit vectors $a \in \mathbb R^i$ of $\langle\sum_la_lf_l, M(k)\sum_la_lf_l\rangle$, a polynomial, so $\mu_i$ is continuous and subanalytic, and $Z_{\geq j} = Z \cap \mu_j^{-1}(0)$ is compact and subanalytic. All the functions are continuous on compact sets, so their graphs are compact.

(i) We take $K = G^4$, $f = \Phi$, and $h = d_\infty(\cdot, Z)$, which vanish exactly on $Z$. Since $\Phi \leq 2$, the bound $d_\infty(k, Z) \leq (\Phi(k)/c)^{1/\rho}$ gives (i) with $\alpha_{\mathrm{Lo}} \coloneqq \min(1, 1/\rho)$ and $C_{\mathrm{Lo}} \coloneqq c^{-1/\rho}2^{\max(0, 1/\rho - 1)}$.

(ii) We take $K = Z_{\geq j}$, $f = \mu_{j+1}$, and $h = d_\infty(\cdot, Z_{\geq j+1})$. On $Z_{\geq j}$ the function $\mu_{j+1}$ vanishes exactly on $Z_{\geq j+1}$, a closed set which contains $(1, 1, 1, 1)$, at which $M = 0$, and on which $h$ vanishes exactly. The theorem gives $c_j, \rho_j > 0$ with $\mu_{j+1} \geq c_jh^{\rho_j}$ on $Z_{\geq j}$, and since $h \leq 2$, (ii) holds with $\varpi_j \coloneqq \max(1, \rho_j)$ and $\mathsf c_j \coloneqq c_j2^{\min(0, \rho_j - 1)}$.
\end{proof}

We fix the hyperplanes once and for all. Let $\sigma \coloneqq o - 4(1, 1, 1, 1)$, so that each hyperplane $\{ x_\beta \equiv \sigma_\beta + \frac12 \}$ has distance at least $3$ from the rectangle bounded by $R_4$, as Lemma~\ref{lem:flat-fields} requires, because $N_q > 12t_{\mathrm b}$ and $N_q > 3s_{\mathrm b}$; let $x_0 \coloneqq \sigma + (1, 1, 1, 1)$, and let $\mathcal T_{\mathrm{tr}}$ be the set of the bonds $\langle x, x + e_\beta\rangle$ of $\mathbb T_{j_q}$ with $x_\gamma \equiv x_{0,\gamma}$ for $\gamma < \beta$ and $x_\beta \not\equiv \sigma_\beta$, a spanning tree, in which $\gamma_x$ denotes the path from $x_0$ to the site $x$. For a bond $b = \langle y, y + e_\beta\rangle$ we put $n(b) \coloneqq e_\beta$ if $y_\beta \equiv \sigma_\beta$ and $n(b) \coloneqq 0$ otherwise, and $c_\beta$ is the closed path of $N_q$ bonds from $x_0$ in the direction $e_\beta$.

\begin{lemma}[A Tree Gauge]\label{lem:tree-gauge}
Let $V \in G^{\mathcal B}$ with $\norm{V(\partial p') - 1} \leq \varepsilon$ for every plaquette $p'$ of $\mathbb T_{j_q}$. There is a gauge transformation $\omega$ of $\mathbb T_{j_q}$ with $V^\omega = 1$ on $\mathcal T_{\mathrm{tr}}$ such that, with $h_\beta \coloneqq V^\omega(c_\beta)$, $h^{e_\beta} \coloneqq h_\beta$, and $h^0 \coloneqq 1$, every bond has $\norm{V^\omega(b) - h^{n(b)}} \leq 50N_q^2\varepsilon$, and $\Phi(h) \leq 8N_q^2\varepsilon$.
\end{lemma}

\begin{proof}
Since $\mathcal T_{\mathrm{tr}}$ is a spanning tree, there is a unique gauge transformation $\omega$ with $\omega(x_0) = 1$ and $V^\omega = 1$ on its bonds, obtained by fixing $\omega$ along the paths $\gamma_x$ outwards from $x_0$. Then $V^\omega(b)$ is the holonomy at $V^\omega$ of the closed path $\ell_b \coloneqq \gamma_{b_-}b\gamma_{b_+}^{-1}$, all its other bonds being bonds of the tree, and it equals the holonomy of $\ell_b$ at $V$, because a gauge transformation conjugates the holonomy of a closed path by its value at the base point; in the same way $h_\beta$ is the value of $V^\omega$ on the last bond of $c_\beta$, the others being bonds of the tree. A closed path of length $m$ in $\mathbb Z^4$ reduces to the empty path by at most $m^2/2$ moves, each a transposition of two consecutive orthogonal steps or the cancellation of two consecutive opposite steps. Indeed, if a direction $e_\beta$ occurs, the signed steps in this direction sum to zero, so two of them which are consecutive among the steps of this direction have opposite signs; the steps between them are orthogonal to $e_\beta$, so at most $m - 2$ transpositions make the two adjacent, and one cancellation removes them, which gives at most $\frac m2(m - 1)$ moves in all. A transposition of two orthogonal steps replaces the holonomy $PXS$ by $PX'S$, where $X$ and $X'$ are the transports along the two sides of the plaquette which the steps span and $X'^{-1}X$ is a holonomy of that plaquette, so that it changes the holonomy by $\norm{X - X'} = \norm{X'^{-1}X - 1} \leq \varepsilon$, by the unitary invariance of the norm; a cancellation leaves it unchanged. The tree paths do not cross the hyperplanes, so $\ell_bc_\beta^{-1}$ for $n(b) = e_\beta$, and $\ell_b$ for $n(b) = 0$, wind around no cycle of the torus and lift to closed paths of $\mathbb Z^4$ of length at most $9N_q$, whose moves project to moves on the torus; and $c_\beta c_\gamma c_\beta^{-1}c_\gamma^{-1}$ lifts to a closed path of length $4N_q$. This gives the two bounds, since $(9N_q)^2/2 \leq 50N_q^2$ and $(4N_q)^2/2 = 8N_q^2$.
\end{proof}

In the gauge of Lemma~\ref{lem:flat-fields}(i) for the hyperplanes above, a flat field $W'$ with the commuting holonomies $\tilde h_\beta$ equals $\tilde h^{n(b)}$ on every bond $b$. We write $\mathfrak g^{\tilde h} \coloneqq \{ Y \in \mathfrak g \,|\, \mathrm{Ad}(\tilde h_\beta)Y = Y \text{ for every } \beta \}$, $\mathcal Z$ for the space of the real forms $c$ on $\mathcal B$ with values in $\mathfrak g^{\tilde h}$ which are constant in each direction, $c_b = c_\beta$ for the bonds of the direction $\beta$, with $\abs c \coloneqq \max_\beta\norm{c_\beta}$, $\Pi_{\mathcal Z}$ for the orthogonal projection onto $\mathcal Z$, and $\bar D$ for the covariant derivative at $W'$ of the $0$-forms of $\mathbb T_{j_q}$. Let $\varpi \coloneqq \max_j\varpi_j$,
\begin{equation}\label{eq:strata-exponents}
    \varsigma_j \coloneqq \frac{\alpha_{\mathrm{Lo}}}{2}(4\varpi)^{-j}\quad (0 \leq j \leq d_{\mathfrak g} + 1)\;,
    \qquad
    \varsigma_* \coloneqq \varsigma_{d_{\mathfrak g}+1}\;,
    \qquad
    \mathsf r_j \coloneqq \bar g^{\varsigma_j}\;,
\end{equation}
so that $\frac12 \geq \varsigma_0 > \varsigma_1 > \dots > \varsigma_* > 0$ and $\varpi_j\varsigma_{j+1} \leq \varsigma_j/4$, and let $c_\varkappa \coloneqq \frac14\min(1, \min_i\mathsf c_i2^{-\varpi_i}/(4\pi^2))$.

\begin{proposition}[Fields Near the Flat Ones]\label{prop:near-flat}
For $\bar g$ small and every $V \in \operatorname{supp}\zeta$ there are an index $j \in \{ 0, \dots, d_{\mathfrak g} \}$, a gauge transformation $\omega'$ of $\mathbb T_{j_q}$, commuting $\tilde h_1, \dots, \tilde h_4 \in G$ with the flat field $W'$ as above, a form $c \in \mathcal Z$, and a form $u_{\mathrm H} \in \mathfrak g^{\mathcal B}$ such that, with $B \coloneqq e^{ic}W'$ and $V_s \coloneqq e^{isu_{\mathrm H}}B$ for $s \in [0, 1]$, $V^{\omega'} = V_1$ and the following hold, where $K \coloneqq \Delta(W')$ and $\delta_* \coloneqq 2(\abs c + \norm{u_{\mathrm H}}_\infty)$.
\begin{enumerate}[label=(\roman*)]
\item The number $\varkappa \coloneqq \frac14N_q^{-2}\min(1, \mathsf c_j(\mathsf r_{j+1}/2)^{\varpi_j}/(4\pi^2))$, the second entry being omitted if $j = d_{\mathfrak g}$, satisfies $\varkappa \geq c_\varkappa\mathsf r_j^{1/4}N_q^{-2}$; moreover $\ker K = \operatorname{ran}\bar D \oplus \mathcal Z$, $K \geq \varkappa$ on $(\ker K)^\perp$, and $\bar D^*\bar D \geq 64\varkappa$ on $(\ker\bar D)^\perp$.
\item $\bar D^*u_{\mathrm H} = 0$ and $\Pi_{\mathcal Z}u_{\mathrm H} = 0$, so that $u_{\mathrm H} \perp \ker K$.
\item $\mathsf A(B) + \frac38\varkappa\norm{u_{\mathrm H}}^2 \leq \mathsf A(V)$, $\langle u_{\mathrm H}, Ku_{\mathrm H}\rangle \leq 3\mathsf A(V)$, and $\mathsf A(V_s) \leq 3\mathsf A(V)$ for $s \in [0, 1]$. Moreover $V_s = e^{iu'_s}W'$ with real forms $u'_s$, $\norm{u'_s}_\infty \leq \delta_*$, the fields $e^{i\sigma u'_s}W'$, $\sigma \in [0, 1]$, lie in $\mathcal R'$, and for all these fields $V'$, $\norm{\Delta(V') - K} \leq C_\Delta N_q^4\delta_* \leq \min(\varkappa/4, \lambda/2)$ and $\norm{(\Delta(V') + \lambda)^{-1}} \leq 2/\lambda$.
\item $\nabla\mathsf A(B) \in \mathcal Z$, and $(\Delta(B) + \lambda)^{-1}$ maps $\mathcal Z$ and $\mathcal Z^\perp$ into themselves.
\item $\abs c \preceq \bar g^{\varsigma_j}$, $\norm{u_{\mathrm H}}_\infty \leq N^{1/2}\norm{u_{\mathrm H}} \preceq \bar g^{1-\varsigma_j/8}$, and $\delta_* \preceq \bar g^{\varsigma_j}$.
\end{enumerate}
\end{proposition}

\begin{proof}
\emph{The flat field.} By Lemma~\ref{lem:tree-gauge} with $\varepsilon = \varepsilon_\zeta$ there is $\omega$ with $V^\omega = 1$ on the tree and holonomies $h$ with $\Phi(h) \leq 8N_q^2\varepsilon_\zeta$. Let $\delta_i \coloneqq d_\infty(h, Z_{\geq i})$ for $i \leq d_{\mathfrak g}$ and $\delta_{d_{\mathfrak g}+1} \coloneqq \infty$. By Lemma~\ref{lem:lojasiewicz}(i), $\delta_0 \leq C_{\mathrm{Lo}}(8N_q^2\varepsilon_\zeta)^{\alpha_{\mathrm{Lo}}} \preceq \bar g^{\alpha_{\mathrm{Lo}}}$, so $\delta_0 \leq \mathsf r_0$ for $\bar g$ small, and there is a least $j \leq d_{\mathfrak g}$ with $\delta_{j+1} > \mathsf r_{j+1}$; then $\delta_j \leq \mathsf r_j$. Let $\tilde h \in Z_{\geq j}$ with $d_\infty(h, \tilde h) = \delta_j$, and let $W'$ be the flat field with the holonomies $\tilde h$. Then $\norm{V^\omega(b) - W'(b)} \leq 50N_q^2\varepsilon_\zeta + \mathsf r_j \leq 2\mathsf r_j$ for every bond, so $V^\omega = e^{iu}W'$ bondwise with $u_b \in \mathfrak g$ and $\norm u_\infty \leq 4\mathsf r_j$, because $2\mathsf r_j \leq r_G$ and $\norm{\log(1 + x)} \leq 2\norm x$ for $\norm x \leq \frac12$.

\emph{The components.} The commuting unitary operators $\mathrm{Ad}(\tilde h_\beta)$ have a common orthonormal eigenbasis $(E_r)$ of $\mathfrak g^{\mathbb C}$, $\mathrm{Ad}(\tilde h_\beta)E_r = e^{2\pi i\theta_{r,\beta}}E_r$ with $\theta_r \in [0, 1)^4$, in which $M(\tilde h)E_r = \sum_\beta\abs{e^{2\pi i\theta_{r,\beta}} - 1}^2E_r$. So $\ker M(\tilde h) = \mathfrak g^{\tilde h}$ is spanned by the $E_r$ with $\theta_r = 0$, and every $r$ with $\theta_r \neq 0$ has $4\pi^2\sum_\beta\norm{\theta_{r,\beta}}^2 \geq \mu_{j'+1}(\tilde h)$, $\norm\cdot$ the distance to $\mathbb Z$ and $j' \coloneqq \dim\mathfrak g^{\tilde h}$, because $\abs{e^{2\pi ix} - 1} \leq 2\pi\norm x$. If $j < d_{\mathfrak g}$, the triangle inequality gives $d_\infty(\tilde h, Z_{\geq j+1}) \geq \delta_{j+1} - \delta_j > \mathsf r_{j+1} - \mathsf r_j \geq \mathsf r_{j+1}/2$ for $\bar g$ small, so $\tilde h \notin Z_{\geq j+1}$ and $j' = j$, and Lemma~\ref{lem:lojasiewicz}(ii) gives $\mu_{j+1}(\tilde h) \geq \mathsf c_j(\mathsf r_{j+1}/2)^{\varpi_j}$. Hence every component with $\theta_r \neq 0$ has $\sum_\beta\norm{\theta_{r,\beta}}^2 \geq 4N_q^2\varkappa$; if $j = d_{\mathfrak g}$ there is no such component.

\emph{(i).} Since $\mathsf r_{j+1}^{\varpi_j} = \bar g^{\varpi_j\varsigma_{j+1}} \geq \mathsf r_j^{1/4}$, $\varkappa \geq c_\varkappa\mathsf r_j^{1/4}N_q^{-2}$. On the component $r$ the translations $\mathcal T_\mu$ by one block, $(\mathcal T_\mu f)(x) \coloneqq f(x - e_\mu)$ with the phase of Lemma~\ref{lem:flat-fields}(ii) across the hyperplane, as in the proof of Proposition~\ref{prop:flat-covariance}, commute and have the joint eigenvalues $(e^{-2\pi i(n_\mu+\theta_{r,\mu})/N_q})_\mu$, $n \in \mathbb Z^4$, with the eigenvectors $x \mapsto e^{2\pi i(n+\theta_r) \cdot x/N_q}$ on the forms of each degree and, for $1$-forms, of each direction. Since $\abs{e^{i\phi} - 1} \geq (4/\pi)\operatorname{dist}(\phi/2, \pi\mathbb Z)$, $\abs{e^{2\pi i(n+\theta)/N_q} - 1}^2 \geq 16\norm\theta^2/N_q^2$, and this is at least $16/N_q^2$ if $\theta = 0$ and $n \notin N_q\mathbb Z$. Hence $\sum_\mu\norm{(\mathcal T_\mu - 1)u}^2 \geq 64\varkappa\norm u^2$ for the forms $u$, of degree $0$ or $1$, in a component with $\theta_r \neq 0$, and for those orthogonal to the constants in a component with $\theta_r = 0$. On $0$-forms $\bar D^*\bar D = \sum_\mu(\mathcal T_\mu - 1)^*(\mathcal T_\mu - 1)$, whose kernel consists of the constants of the components with $\theta_r = 0$, which gives the last claim. For $c \in \mathcal Z$ let $A_c$ be the fine $1$-form equal to $c_\beta$ on the fine bonds of the direction $\beta$. Since $\mathrm{Ad}(\tilde h_\beta)$ fixes the values of $c$, $A_c$ has the constant components of the components with $\theta_r = 0$, so $DA_c = 0$ and $Q^{\mathrm{dr}}A_c = c$ by Lemma~\ref{lem:flat-fields}(ii), the straight average of a constant being that constant; by Lemma~\ref{lem:flat-fields}(iii), $DHc$ is the point of least norm of $\{ DA \,|\, Q^{\mathrm{dr}}A = c \}$, which is $0$, so that $\mathcal Z \subset \ker DH = \ker K$. Next, $\operatorname{ran}\bar D \subset \ker K$, because $K\bar D = 0$ by the proof of Lemma~\ref{lem:coarse-coercivity}(i). Finally $\mathcal Z \perp \operatorname{ran}\bar D$: for a $0$-form $\omega$ and $c \in \mathcal Z$, $(\mathrm{Ad}(W'(b))\omega(b_+)) \cdot c_b = \omega(b_+) \cdot c_b$, so $\langle\bar D\omega, c\rangle = \sum_\beta\sum_b(\omega(b_+) - \omega(b_-)) \cdot c_\beta$, the inner sum over the bonds of the direction $\beta$, which vanishes because every site is the endpoint $b_+$ of exactly one of them and the endpoint $b_-$ of exactly one; thus $\bar D^*c = 0$. For $u \perp \operatorname{ran}\bar D \oplus \mathcal Z$, that is $\bar D^*u = 0$ and $\Pi_{\mathcal Z}u = 0$, Lemma~\ref{lem:coarse-coercivity}(ii) in each direction gives $64\langle u, Ku\rangle = 64\norm{DHu}^2 \geq \sum_\mu\norm{\nabla_\mu u}^2 \geq 64\varkappa\norm u^2$. So $\ker K = \operatorname{ran}\bar D \oplus \mathcal Z$ and $K \geq \varkappa$ on its orthogonal complement.

\emph{(ii), the slice.} For a $0$-form $\omega_1 \perp \ker\bar D$ and $c \in \mathcal Z$ we define $u_{\mathrm H}$ by $e^{iu_{\mathrm H,b}} \coloneqq e^{i\omega_1(b_-)}e^{iu_b}e^{-i\mathrm{Ad}(W'(b))\omega_1(b_+)}e^{-ic_b}$, so that $e^{iu_{\mathrm H,b}}e^{ic_b}W'(b) = e^{i\omega_1(b_-)}V^\omega(b)e^{-i\omega_1(b_+)}$, that is $(V^\omega)^{e^{i\omega_1}} = e^{iu_{\mathrm H}}e^{ic}W'$, and $\omega' \coloneqq e^{i\omega_1}\omega$. Then $u_{\mathrm H} = u - \bar D\omega_1 - c + \mathcal N(\omega_1, c)$, with $(\bar D\omega_1)_b = \mathrm{Ad}(W'(b))\omega_1(b_+) - \omega_1(b_-)$, and by the Baker--Campbell--Hausdorff series, which converges because the four exponents have norms at most $\frac1{10}$ on the set below for $\bar g$ small, with a constant $C_{\mathrm{BCH}}$ of $N$, the remainder satisfies $\norm{\mathcal N_b} \leq C_{\mathrm{BCH}}(\norm{\omega_1(b_-)} + \norm{\omega_1(b_+)} + \norm{c_b})(\norm{\omega_1(b_-)} + \norm{\omega_1(b_+)} + \norm{c_b} + \norm{u_b})$, every product of two of the exponents containing an exponent other than $u_b$, together with the corresponding Lipschitz bound. Since $\bar D^*c = 0$ and $\Pi_{\mathcal Z}\bar D\omega_1 = 0$, the conditions $\bar D^*u_{\mathrm H} = 0$ and $\Pi_{\mathcal Z}u_{\mathrm H} = 0$ say that $(\omega_1, c)$ is a fixed point of
\begin{equation}\label{eq:slice-map}
    \mathcal T_{\mathrm{sl}}(\omega_1, c) \coloneqq \Bigl((\bar D^*\bar D)^+\bar D^*\bigl(u + \mathcal N(\omega_1, c)\bigr),\, \Pi_{\mathcal Z}\bigl(u + \mathcal N(\omega_1, c)\bigr)\Bigr)\;.
\end{equation}
By (i), $\norm{(\bar D^*\bar D)^+\bar D^*} \leq (64\varkappa)^{-1/2}$. Let $R_u \coloneqq 8N_q^2\norm u_\infty$, which is at least $4\norm u$ because $\mathcal B$ has $4N_q^4$ bonds and $\abs\cdot_2 \leq \norm\cdot$, and $R_\omega \coloneqq R_u(64\varkappa)^{-1/2} \geq R_u$. On the set of the pairs with $\norm{\omega_1} \leq R_\omega$ and $\norm c \leq R_u$, every site lying on eight bonds and $\norm\cdot_\infty \leq N^{1/2}\norm\cdot$, the bound on $\mathcal N_b$ gives $\norm{\mathcal N} \leq 90NC_{\mathrm{BCH}}R_\omega^2$ and the Lipschitz constant $90NC_{\mathrm{BCH}}R_\omega$. If $C_{\mathrm{sl}}N_q^2\mathsf r_j \leq \varkappa$ with $C_{\mathrm{sl}} \coloneqq 180NC_{\mathrm{BCH}}$, then $90NC_{\mathrm{BCH}}R_\omega^2 \leq \frac34R_u$ and $2(64\varkappa)^{-1/2}\cdot90NC_{\mathrm{BCH}}R_\omega \leq \frac12$, so $\mathcal T_{\mathrm{sl}}$ maps the set into itself and is a contraction there, and it has a fixed point. The condition holds for $\bar g$ small, because $N_q^2\mathsf r_j/\varkappa \leq c_\varkappa^{-1}N_q^4\bar g^{3\varsigma_j/4}$. At the fixed point, $c$ is constant, so $\abs c \leq N^{1/2}R_u/N_q^2 = 8N^{1/2}\norm u_\infty \leq 32N^{1/2}\mathsf r_j$; and $\norm{\omega_1}_\infty \leq N^{1/2}R_\omega$ gives $\norm{u_{\mathrm H}}_\infty \leq C'_{\mathrm{sl}}N_q^2\mathsf r_j\varkappa^{-1/2}$ with a constant $C'_{\mathrm{sl}}$ of $N$. Finally $u_{\mathrm H} \perp \operatorname{ran}\bar D \oplus \mathcal Z = \ker K$.

\emph{(iii), the coercivity.} By the series of the exponential and the logarithm, $V_s = e^{iu'_s}W'$ with $u'_{s,b} \coloneqq -i\log(e^{isu_{\mathrm H,b}}e^{ic_b})$ and $\norm{u'_s}_\infty \leq 2(\norm{u_{\mathrm H}}_\infty + \abs c) = \delta_*$. The fields $V' = e^{i\sigma u'_s}W'$, $\sigma \in [0, 1]$, lie in $\mathcal R$: every bond value is within $\delta_*$ of that of the flat $W'$, so their plaquettes satisfy $\abs{V'(\partial p') - 1} \leq N^{1/2}\norm{V'(\partial p') - 1} \leq 4N^{1/2}\delta_* < \varepsilon_1$. They lie in $\mathcal R'$, because the holonomies of $W'$ around the $I_b$ are $1$, and by Lemma~\ref{lem:moved-minimiser}(i) at the base $W'$ and the bound on $\norm{H_k(B) - H_k(0)}$ in the proof of Lemma~\ref{lem:flux-derivatives}(i), with $B = \sigma u'_s$, those of $V'$ differ from $1$ by at most $e^{2P\varepsilon_V\delta_*/r_V} - 1 \leq r_G/2$. Since $\frac{\dd}{\dd\sigma}\Delta(e^{i\sigma u'_s}W') = \partial_{u'_s}\Delta(e^{i\sigma u'_s}W')$, Lemma~\ref{lem:flux-derivatives}(ii) gives $\norm{\Delta(V') - K} \leq C_\Delta N_q^4\delta_*$, with $K = \Delta(W')$ by Lemma~\ref{lem:flat-fields}(iv). With the slice bounds, $N_q^4\delta_*/\varkappa \preceq \bar g^{5\varsigma_j/8}$ and $N_q^4\delta_*/\lambda \preceq \bar g^{7\varsigma_j/8}$, so $C_\Delta N_q^4\delta_* \leq \min(\varkappa/4, \lambda/2)$ for $\bar g$ small; since $K \geq 0$, this gives $\Delta(V') + \lambda \geq \lambda/2$. Now $e^{i(s+h)u_{\mathrm H}}B = e^{ihu_{\mathrm H}}V_s$, so the function $s \mapsto \mathsf A(V_s)$, which is real analytic because $V_s \in \mathcal R$, has the derivatives $\langle u_{\mathrm H}, \nabla\mathsf A(V_s)\rangle$ and $\langle u_{\mathrm H}, \Delta(V_s)u_{\mathrm H}\rangle$; since $\mathsf A$ is invariant under the gauge transformations of $\mathbb T_{j_q}$ (see (iv)), Taylor's formula gives
\begin{equation}\label{eq:coercivity-taylor}
    \mathsf A(V) = \mathsf A(V_1) = \mathsf A(B) + \langle u_{\mathrm H}, \nabla\mathsf A(B)\rangle + \int_0^1(1 - s)\langle u_{\mathrm H}, \Delta(V_s)u_{\mathrm H}\rangle\,\dd s\;.
\end{equation}
The middle term vanishes, by (iv) below and $\Pi_{\mathcal Z}u_{\mathrm H} = 0$. By (i) and (ii), $\langle u_{\mathrm H}, Ku_{\mathrm H}\rangle \geq \varkappa\norm{u_{\mathrm H}}^2$, so that $\langle u_{\mathrm H}, \Delta(V_s)u_{\mathrm H}\rangle$ lies between $\langle u_{\mathrm H}, Ku_{\mathrm H}\rangle \mp \frac14\varkappa\norm{u_{\mathrm H}}^2$ and is at least $\frac34\varkappa\norm{u_{\mathrm H}}^2$. Hence \eqref{eq:coercivity-taylor} gives $\mathsf A(V) - \mathsf A(B) \geq \frac38\varkappa\norm{u_{\mathrm H}}^2$, and also $\frac12\langle u_{\mathrm H}, Ku_{\mathrm H}\rangle \leq \mathsf A(V) - \mathsf A(B) + \frac18\varkappa\norm{u_{\mathrm H}}^2 \leq \frac43(\mathsf A(V) - \mathsf A(B))$, so that $\langle u_{\mathrm H}, Ku_{\mathrm H}\rangle \leq \frac83\mathsf A(V)$. Taylor's formula on $[0, s]$ with the same upper bound for the integrand gives $\mathsf A(V_s) \leq \mathsf A(B) + \frac12(\langle u_{\mathrm H}, Ku_{\mathrm H}\rangle + \frac14\varkappa\norm{u_{\mathrm H}}^2) \leq \mathsf A(B) + \frac53(\mathsf A(V) - \mathsf A(B)) \leq 3\mathsf A(V)$, using $\mathsf A(B) \geq 0$.

\emph{(iv), the symmetry.} Let $\mathcal T$ be the translation of $\mathbb T_{j_q}$ by one block in a direction $\nu$, $(\mathcal TV)(b) \coloneqq V(b - e_\nu)$. With the convention $V^\upsilon(b) = \upsilon(b_-)V(b)\upsilon(b_+)^{-1}$, under which the gauge transformations act on $\mathfrak g^{\mathcal B}$ as in Section~\ref{sec:flat-covariance}, $\mathcal TW' = W'^{g_\nu}$, with $g_\nu$ equal to $\tilde h_\nu$ on the sites $x$ with $x_\nu \equiv \sigma_\nu + 1$ and to $1$ elsewhere: the bonds of the direction $\nu$ which enter this layer get the value $\tilde h_\nu\tilde h_\nu^{-1} = 1$, those which leave it the value $\tilde h_\nu$, and on the bonds inside the layer $g_\nu$ conjugates values in $\{ 1, \tilde h_\beta \}$, which commute with $\tilde h_\nu$. Since $\mathcal Tc = c$ and $\mathrm{Ad}(\tilde h_\nu)c_\beta = c_\beta$, also $\mathcal TB = B^{g_\nu}$. The function $\mathsf A$ is invariant under the gauge transformations of $\mathbb T_{j_q}$, because $U_{K'}(V^\upsilon) = U_{K'}(V)^{\bar\upsilon}$ \cite[(181), p.~307]{BalabanVAR} and the action is gauge invariant, and under the translations by whole blocks, because these commute with $M^{K'}$ and leave the action invariant, so that $\mathcal TU_{K'}(V)$ is minimal for $\mathcal TV$. With the orthogonal map $S_\nu u \coloneqq \mathrm{Ad}(g_\nu)^{-1}\mathcal Tu$, acting on the values at $b$ by $\mathrm{Ad}(g_\nu(b_-))^{-1}$, we get $(\mathcal T(e^{iu}B))^{g_\nu^{-1}} = e^{iS_\nu u}B$, hence $\mathsf A(e^{iS_\nu u}B) = \mathsf A(e^{iu}B)$ for all $u$. Differentiating, $S_\nu\nabla\mathsf A(B) = \nabla\mathsf A(B)$ and $S_\nu$ commutes with $\Delta(B)$. On the component $r$, $S_\nu$ is the twisted translation $\mathcal T_\nu$ of (i), so the joint fixed vectors of $S_1, \dots, S_4$ are the constants of the components with $\theta_r = 0$, and the real ones form $\mathcal Z$. This gives (iv).

\emph{(v), the orders.} By (iii) and $\mathsf A(V) \leq 2T_{\mathrm{ng}}g^2$, $\norm{u_{\mathrm H}}^2 \leq \frac{16}3T_{\mathrm{ng}}g^2/\varkappa \preceq \bar g^{2-\varsigma_j/4}$, by (i); and $\abs c \leq 32N^{1/2}\bar g^{\varsigma_j}$, so $\delta_* \preceq \bar g^{\varsigma_j}$, because $\varsigma_j \leq \frac12 < 1 - \varsigma_j/8$.
\end{proof}

With the decomposition in hand we can bound the defect $\mathfrak e_\alpha$. We follow it along the segment $s \mapsto V_s$ from the field $B$, which differs from the flat $W'$ only in the directions of $\mathcal Z$. At $B$ the gradient of the action lies in $\mathcal Z$, so only the component of $\nabla\Xi_\alpha$ in $\mathcal Z$ enters, and it is of the second order in $c$ because the fluxes vanish to the second order along $\mathcal Z$; the flux difference $\Xi_\alpha(B)$ itself is a commutator of a small transport with the flux of $I_0$. Along the segment the derivative of the defect is $-\lambda\langle u_{\mathrm H}, w_\alpha\rangle$ up to terms of a higher order, and at a flat field this is the term which Proposition~\ref{prop:flat-covariance}(ii) controls.

\begin{proposition}[The Defect of the Integration by Parts]\label{prop:ibp-error}
There is a constant $C_E$ of the kind of Section~\ref{sec:flux-moments} such that for $\bar g$ small, $K \geq K_{\mathrm w}$, and every $V \in \operatorname{supp}\zeta$, the operator $\Delta(V) + \lambda$ is at least $\lambda/2$, and
\begin{equation}\label{eq:ibp-error}
    \Bigl\lvert\sum_{\alpha=1}^{d_{\mathfrak g}}X^\alpha(V)\,\mathfrak e_\alpha(V)\Bigr\rvert \leq C_E\bar g^2\bigl(\log\bar g^{-2}\bigr)^{7r+1-p_\lambda/2} + \bar g^{2+\varsigma_*/2}\;.
\end{equation}
\end{proposition}

\begin{proof}
\emph{Reduction.} Under a gauge transformation $\omega$ of $\mathbb T_{j_q}$, the gradients and $\Delta$ are covariant under the orthogonal map $u_b \mapsto \mathrm{Ad}(\omega(b_-))u_b$, $\mathsf A$ is invariant, and the logarithms based at $o$ change by $\mathrm{Ad}(\omega(o))$; so, as for $c_\lambda$ in Section~\ref{sec:flat-covariance}, the sum in \eqref{eq:ibp-error} at $V^\omega$ is the same sum at $V$ for the orthonormal basis $(\mathrm{Ad}(\omega(o))^{-1}t_\alpha)$, and it does not depend on the basis, being the contraction of a bilinear expression in $t_\alpha$ with $\sum_\alpha t_\alpha \otimes t_\alpha$. The spectrum of $\Delta$ is invariant as well. We may therefore replace $V$ by $V_1 = V^{\omega'}$ of Proposition~\ref{prop:near-flat}, whose notation we use; there $\Delta(V_1) + \lambda \geq \lambda/2$ by Proposition~\ref{prop:near-flat}(iii). Since $\abs{X_{0,0}(V_1)}_2 \leq \norm{X_{0,0}(V_1)} \leq 2h^G$, the Cauchy--Schwarz inequality over $\alpha$ gives
\begin{equation}\label{eq:ibp-error-reduction}
    \Bigl\lvert\sum_\alpha X^\alpha\mathfrak e_\alpha\Bigr\rvert \leq 2d_{\mathfrak g}^{1/2}h^G\max_\alpha\abs{\mathfrak e_\alpha(V_1)}\;,
\end{equation}
and it remains to bound $\mathfrak e_\alpha(V_1)$ for a fixed $\alpha$. For the fields $V_s$, $s \in [0, 1]$, among which $V_0 = B$, we have $\mathsf A(V_s) \leq 3\mathsf A(V) \leq 6T_{\mathrm{ng}}g^2$, and $e^{-itw}V_s \in \mathcal R$ for $\norm w = 1$ and $0 \leq t \leq (2\mathsf A(V_s)/(C_\Delta N_q^4))^{1/2}$, a number $\preceq \bar g$, because the plaquettes of $V_s$ are within $4N^{1/2}\delta_*$ of $1$; so Lemma~\ref{lem:flux-derivatives}(iii) gives $\norm{\nabla\mathsf A(V_s)} \leq \gamma_A \coloneqq (12C_\Delta N_q^4T_{\mathrm{ng}})^{1/2}g$. On all the fields $V'$ of Proposition~\ref{prop:near-flat}(iii), Lemma~\ref{lem:flux-derivatives}(i) and the bound on the resolvent give $\norm{w_\alpha(V')} \leq 2C_\Xi P/\lambda$.

\emph{The flux at the fields of $\mathcal Z$.} For a real $z \in \mathcal Z$ with $\abs z < r_V$ let $B_z \coloneqq e^{iz}W'$, so that $B = B_c$, and let $\mathsf t_k(z) \in G$ be the transport from $o$ to $o_k$ along the bottom side of $R_4$ at a minimal configuration for $B_z$, which does not depend on its choice because both ends are sites of $\mathbb T_{j_q}$. By the proof of Lemma~\ref{lem:lasso}, $X_{0,k} = \mathrm{Ad}(\mathsf t_k)X_{k,k}$. We claim that $X_{k,k}(B_z) = X_{0,0}(B_z)$ for $k \leq 3$. The argument of Proposition~\ref{prop:near-flat}(iv) for the translation $\mathcal T^{-1}$ in the direction $\mu$ gives $\mathcal T^{-1}B_z = B_z^{g'}$, with $g'$ equal to $\tilde h_\mu^{-1}$ on the layer $x_\mu \equiv \sigma_\mu$ and to $1$ elsewhere, and by induction $\mathcal T^{-n}B_z = B_z^{G_n}$, where $G_n$ equals $1$ off the layers $x_\mu \equiv \sigma_\mu, \sigma_\mu - 1, \dots, \sigma_\mu - n + 1$. For $n = kt_{\mathrm b} \leq 3t_{\mathrm b}$ these layers do not contain $o$, because $o_\mu = \sigma_\mu + 4$ and $N_q > 3t_{\mathrm b} + 4$. Since $\mathcal T^{-n}$ commutes with $M^{K'}$, the holonomy of $I_0$ at $o$ for $\mathcal T^{-n}B_z$ is that of $I_k$ at $o_k$ for $B_z$, and by \cite[(181), p.~307]{BalabanVAR} it is also the holonomy of $I_0$ for $B_z$ conjugated by $G_n(o) = 1$; this is the claim. Hence
\begin{equation}\label{eq:flux-along-z}
    \hat\Xi(z) \coloneqq \Xi_\alpha(B_z) = t_\alpha \cdot \bigl(\mathrm{Ad}(\mathsf t_3(z)) - \mathrm{Ad}(\mathsf t_2(z))\bigr)X_{0,0}(B_z)\;.
\end{equation}
At $z = 0$ the minimal configuration is $U_{\mathrm s}$ (Lemma~\ref{lem:flat-fields}(i)), which is $1$ near $R_4$, so $\mathsf t_k(0) = 1$ and $X_{0,0}(B_0) = 0$; moreover $\frac{\dd}{\dd s}X^\beta(B_{sz})|_{s=0} = \langle\chi_0 \otimes t_\beta, DHz\rangle = 0$ by Lemma~\ref{lem:flat-fields}(iv) and $\mathcal Z \subset \ker DH$. By Lemma~\ref{lem:moved-minimiser}(i) at the base $W'$, the transports $\mathsf t_k(z)$ are real analytic in $z$, so in \eqref{eq:flux-along-z} the first factor vanishes to the first order at $z = 0$ and the second to the second order, and the Taylor polynomial of degree two of $\hat\Xi$ at $0$ vanishes. By the proof of Lemma~\ref{lem:flux-derivatives}(i) at the base $W' \in \mathcal R'$, $\hat\Xi$ extends analytically to the complex $z \in \mathcal Z^{\mathbb C}$ with $\abs z < \rho_P$, with modulus at most $\frac12$, and Schwarz's lemma for the function $\zeta \mapsto \hat\Xi(\zeta z/\abs z)$, which vanishes to the third order at $\zeta = 0$, gives $\abs{\hat\Xi(z)} \leq \frac12(\abs z/\rho_P)^3 \leq \frac12c_P^3P^3\abs z^3$.

\emph{The defect at $B$.} Since $\nabla\mathsf A(B) \in \mathcal Z$ and $(\Delta(B) + \lambda)^{-1}$ maps $\mathcal Z$ into itself (Proposition~\ref{prop:near-flat}(iv)), $\langle w_\alpha(B), \nabla\mathsf A(B)\rangle = \langle\Pi_{\mathcal Z}\nabla\Xi_\alpha(B), (\Delta(B) + \lambda)^{-1}\nabla\mathsf A(B)\rangle$. Let $m \coloneqq \abs c + \bar g$ and $z' \in \mathcal Z$ with $\abs{z'} = 1$. For complex $\tau$ with $\abs\tau \leq m$ the form $Z(\tau)$ with $e^{iZ(\tau)_b} = e^{i\tau z'_b}e^{ic_b}$ lies in $\mathcal Z^{\mathbb C}$, because $\mathfrak g^{\tilde h}$ is a Lie algebra, and $\abs{Z(\tau)} \leq 2(e^m - 1 + \abs c) \leq 5m < \rho_P$ for $\bar g$ small, by the series of the logarithm; so Cauchy's estimate for $\tau \mapsto \hat\Xi(Z(\tau))$ on $\abs\tau \leq m$, whose derivative at $0$ is $\partial_{z'}\Xi_\alpha(B)$, gives $\abs{\partial_{z'}\Xi_\alpha(B)} \leq \frac12c_P^3P^3(5m)^3/m \leq 63c_P^3P^3m^2$. Since $\norm{z'} \geq N_q^2N^{-1/2}$ for such $z'$, $\norm{\Pi_{\mathcal Z}\nabla\Xi_\alpha(B)} \leq \hat C_\Xi N_qm^2$ with $\hat C_\Xi \coloneqq 504N^{1/2}c_P^3$, using $P \leq 2N_q$. For $\Xi_\alpha(B)$ itself we use \eqref{eq:flux-along-z} at $z = c$: by unitarity, $\norm{(\mathrm{Ad}(\mathsf t_3) - \mathrm{Ad}(\mathsf t_2))Y} \leq 2\norm{\mathsf t_3 - \mathsf t_2}\norm Y$ and $\norm{\mathsf t_3 - \mathsf t_2}$ is the distance from $1$ of the transport from $o_2$ to $o_3$. By Lemma~\ref{lem:moved-minimiser} at the base $W'$, this transport is $w_c(o_2)\mathsf t'w_c(o_3)^{-1}$, up to the order of the factors, with $\mathsf t'$ the transport of $e^{i\eta\mathcal H(c)}U_{\mathrm s}$ along $t_{\mathrm b}/\eta$ fine bonds on which $U_{\mathrm s} = 1$, so that, all the factors being unitary, its distance from $1$ is at most $(t_{\mathrm b} + 144d)\varepsilon_V\abs c/r_V$. Moreover $\norm{X_{0,0}(B)} \leq \norm{X_{0,0}(V_1)} + C_\Xi P\norm{u_{\mathrm H}}_\infty \leq 2h^G + C_\Xi P\norm{u_{\mathrm H}}_\infty$, by Lemma~\ref{lem:flux-derivatives}(i) along the segment, whose fields lie in $\mathcal R'$. Together, with $\norm{(\Delta(B) + \lambda)^{-1}} \leq 2/\lambda$,
\begin{equation}\label{eq:defect-background}
    \abs{\mathfrak e_\alpha(B)} \leq \frac2\lambda\hat C_\Xi N_qm^2\gamma_A + \frac{2(P + 144d)\varepsilon_V}{r_V}\abs c\bigl(2h^G + C_\Xi P\norm{u_{\mathrm H}}_\infty\bigr)\;.
\end{equation}

\emph{The segment.} Let $w_s \coloneqq w_\alpha(V_s)$ and $\phi(s) \coloneqq \mathfrak e_\alpha(V_s)$. Since $V_{s+h} = e^{ihu_{\mathrm H}}V_s$, we have $\frac{\dd}{\dd s}\Xi_\alpha(V_s) = \langle u_{\mathrm H}, \nabla\Xi_\alpha(V_s)\rangle = \langle u_{\mathrm H}, (\Delta(V_s) + \lambda)w_s\rangle$, and $\frac{\dd}{\dd s}\langle w, \nabla\mathsf A(V_s)\rangle = \partial_\sigma\partial_t\mathsf A(e^{itw}e^{i\sigma u_{\mathrm H}}V_s)|_{t=\sigma=0}$ for a fixed $w$. By the Baker--Campbell--Hausdorff formula, $e^{itw}e^{i\sigma u_{\mathrm H}} = e^{i(tw + \sigma u_{\mathrm H} + \frac12t\sigma\,i[w, u_{\mathrm H}]) + O_3}$ bondwise, where $i[w, u_{\mathrm H}]$ is real and $O_3$ collects the terms of order three in $(t, \sigma)$, so that this mixed derivative is $\langle w, \Delta(V_s)u_{\mathrm H}\rangle + \frac12\langle i[w, u_{\mathrm H}], \nabla\mathsf A(V_s)\rangle$. Therefore
\begin{equation}\label{eq:defect-derivative}
    \phi'(s) = \bigl\langle\partial_sw_s, \nabla\mathsf A(V_s)\bigr\rangle - \lambda\langle u_{\mathrm H}, w_s\rangle + \frac12\bigl\langle i[w_s, u_{\mathrm H}], \nabla\mathsf A(V_s)\bigr\rangle\;,
\end{equation}
the terms with $\Delta(V_s)$ cancelling. We bound the three terms.

\emph{The first and the third term.} Differentiating the resolvent, $\partial_sw_s = -R_s(\partial_{u_{\mathrm H}}\Delta(V_s))R_s\nabla\Xi_\alpha(V_s) + R_s\partial_{u_{\mathrm H}}\nabla\Xi_\alpha(V_s)$ with $R_s \coloneqq (\Delta(V_s) + \lambda)^{-1}$, so Lemma~\ref{lem:flux-derivatives}(i), (ii) and $\norm{R_s} \leq 2/\lambda$ give $\norm{\partial_sw_s} \leq C_\Xi P\norm{u_{\mathrm H}}_\infty(4C_\Delta N_q^4/\lambda^2 + 2P/\lambda)$. Since $\abs{[a, b]}_2 \leq 2\norm b\abs a_2$, $\norm{i[w_s, u_{\mathrm H}]} \leq 2\norm{u_{\mathrm H}}_\infty\norm{w_s}$. Hence the first and the third term are at most
\begin{equation}\label{eq:defect-first-third}
    \gamma_AC_\Xi P\norm{u_{\mathrm H}}_\infty\Bigl(\frac{4C_\Delta N_q^4}{\lambda^2} + \frac{4P}{\lambda}\Bigr) \preceq \bar g^{2-\varsigma_j/8}\;,
\end{equation}
by Proposition~\ref{prop:near-flat}(v), $\gamma_A \preceq \bar g$, and $1/\lambda \preceq 1$.

\emph{The term of the regularisation.} We compare it with its value at the flat field. Writing $R_s - (K + \lambda)^{-1} = -(K + \lambda)^{-1}(\Delta(V_s) - K)R_s$,
\begin{equation}\label{eq:lambda-term-split}
    \begin{split}
        \lambda\langle u_{\mathrm H}, w_s\rangle &= \lambda\bigl\langle(K + \lambda)^{-1}u_{\mathrm H}, \nabla\Xi_\alpha(W')\bigr\rangle - \lambda\bigl\langle(K + \lambda)^{-1}u_{\mathrm H}, (\Delta(V_s) - K)w_s\bigr\rangle
        \\
        &\quad + \lambda\bigl\langle(K + \lambda)^{-1}u_{\mathrm H}, \nabla\Xi_\alpha(V_s) - \nabla\Xi_\alpha(W')\bigr\rangle\;.
    \end{split}
\end{equation}
Since $u_{\mathrm H} \perp \ker K$ and $K \geq \varkappa$ there, $\norm{\lambda(K + \lambda)^{-1}u_{\mathrm H}} \leq (\lambda/\varkappa)\norm{u_{\mathrm H}}$. By Proposition~\ref{prop:near-flat}(iii), $\norm{\Delta(V_s) - K} \leq C_\Delta N_q^4\delta_*$, and Lemma~\ref{lem:flux-derivatives}(i) along the fields $e^{i\sigma u'_s}W'$, which lie in $\mathcal R'$, gives $\norm{\nabla\Xi_\alpha(V_s) - \nabla\Xi_\alpha(W')} \leq C_\Xi P^2\delta_*$; so the last two terms of \eqref{eq:lambda-term-split} are at most
\begin{equation}\label{eq:lambda-term-rest}
    \frac{\norm{u_{\mathrm H}}\delta_*}{\varkappa}\bigl(2C_\Delta C_\Xi PN_q^4 + C_\Xi P^2\bigr) \preceq \bar g^{1+5\varsigma_j/8}\;,
\end{equation}
by Proposition~\ref{prop:near-flat}(i) and (v). In the first term we use $\nabla\Xi_\alpha(W') = H^*D^*((\chi_3 - \chi_2) \otimes t_\alpha)$ (Lemma~\ref{lem:flat-fields}(iv)) and Proposition~\ref{prop:flat-covariance}(ii), which give $\lambda\langle(K + \lambda)^{-1}u_{\mathrm H}, \nabla\Xi_\alpha(W')\rangle = \langle P_\lambda DHu_{\mathrm H}, (\chi_3 - \chi_2) \otimes t_\alpha\rangle$, where $P_\lambda$ denotes the operator $P$ of that proposition. We expand $t_\alpha = \sum_ra_rE_r$ in the basis of the components, with $\sum_r\abs{a_r}^2 = 1$; the operators $D$, $H$, $P_\lambda$ act on each component, so \eqref{eq:lambda-term-flat} and the Cauchy--Schwarz inequality over $r$ give
\begin{equation}\label{eq:lambda-term-main}
    \begin{split}
        \abs{\langle P_\lambda DHu_{\mathrm H}, (\chi_3 - \chi_2) \otimes t_\alpha\rangle} &\leq 2\pi t_{\mathrm b}\bigl(\lambda C_V(s_{\mathrm b} + 2)(t_{\mathrm b} + 2)\bigr)^{1/2}\norm{DHu_{\mathrm H}}
        \\
        &\leq 2\pi(6C_VT_{\mathrm{ng}})^{1/2}N_qg\bigl(\log\bar g^{-2}\bigr)^{-p_\lambda/2}\;,
    \end{split}
\end{equation}
because $\norm{DHu_{\mathrm H}}^2 = \langle u_{\mathrm H}, Ku_{\mathrm H}\rangle \leq 3\mathsf A(V) \leq 6T_{\mathrm{ng}}g^2$ by Proposition~\ref{prop:near-flat}(iii), $t_{\mathrm b} \leq N_q$, $(s_{\mathrm b} + 2)(t_{\mathrm b} + 2) \leq N_q^2$, and $\lambda^{1/2} = N_q^{-1}(\log\bar g^{-2})^{-p_\lambda/2}$.

\emph{Conclusion.} The right side of \eqref{eq:lambda-term-main} is at most $C\bar g(\log\bar g^{-2})^{3r+1/2-p_\lambda/2}$, by the sizes of Section~\ref{sec:flux-moments}. By Proposition~\ref{prop:near-flat}(v), $\abs c$, $m \preceq \bar g^{\varsigma_j}$ and $\norm{u_{\mathrm H}}_\infty \preceq \bar g^{1-\varsigma_j/8}$, so the two terms of \eqref{eq:defect-background} are $\preceq \bar g^{1+2\varsigma_j}$ and $\preceq \bar g^{1+7\varsigma_j/8}$. With \eqref{eq:defect-first-third}, \eqref{eq:lambda-term-rest}, and $2 - \varsigma_j/8 > 1 + 5\varsigma_j/8$, we obtain from $\mathfrak e_\alpha(V_1) = \phi(0) + \int_0^1\phi'(s)\,\dd s$ and $\phi(0) = \mathfrak e_\alpha(B)$ that $\abs{\mathfrak e_\alpha(V_1)} \leq C\bar g(\log\bar g^{-2})^{3r+1/2-p_\lambda/2} + \mathsf E$ with $\mathsf E \preceq \bar g^{1+5\varsigma_j/8}$. Since $h^G \leq C\bar g(\log\bar g^{-2})^{4r+1/2}$ and $\varsigma_j \geq \varsigma_*$, \eqref{eq:ibp-error-reduction} gives \eqref{eq:ibp-error}, the term with $\mathsf E$ being at most $\bar g^{2+\varsigma_*/2}$ for $\bar g$ small.
\end{proof}

We can now assemble the second moments. The integration by parts of Lemma~\ref{lem:integration-by-parts} with the vector fields $w_\alpha$ turns $\sum_\alpha\langle F_\alpha\Xi_\alpha\rangle$ into $g^2\langle\zeta c_\lambda\rangle$, up to the defect of Proposition~\ref{prop:ibp-error} and to three terms which carry a factor $g^2$ besides the factor $h^G$ of $F_\alpha$; and at a field near a flat one the number $c_\lambda$ is close to its value at the flat field, which Proposition~\ref{prop:flat-covariance}(i) bounds below.

\begin{proposition}[The Second Moments of the Fluxes]\label{prop:flux-moments}
For $\bar g$ sufficiently small, depending on $q$, $t/u$, $s/u$, $T/u$, $\pi$, and the constants of the construction, and for every $K \geq K_{\mathrm w}$, $\Psi_3 - \Psi_2 \geq 3c_O\bar g^2$, with $c_O$ of Theorem~\ref{thm:non-gaussian}.
\end{proposition}

\begin{proof}
\emph{The cutoff.} By \eqref{eq:flux-moments} and \eqref{eq:flux-coordinates}, $\Psi_3 - \Psi_2 = (C_2(\pi)/d_{\mathfrak g})\langle\mathbf 1_{\mathcal V}\sum_\alpha X^\alpha\Xi_\alpha\rangle$. The function $\zeta$ equals $1$ on $\mathcal V$ and vanishes outside $\mathcal V_3$, on which $\abs{\sum_\alpha X^\alpha\Xi_\alpha} \leq \abs{X_{0,0}}_2\abs{X_{0,3} - X_{0,2}}_2 \leq 8(h^G)^2$, and $\langle\mathbf 1_{\mathcal V^c}\rangle \leq \bar g^8$ by Lemma~\ref{lem:stopping-sizes}(ii); hence
\begin{equation}\label{eq:moments-cutoff}
    \Bigl\lvert\frac{d_{\mathfrak g}}{C_2(\pi)}(\Psi_3 - \Psi_2) - \sum_\alpha\langle F_\alpha\Xi_\alpha\rangle\Bigr\rvert \leq 8(h^G)^2\bar g^8\;.
\end{equation}

\emph{The integration by parts.} By Proposition~\ref{prop:ibp-error}, $\Delta + \lambda \geq \lambda/2$ on the compact set $\operatorname{supp}\zeta$, on which the holonomies of the $I_b$ are within $\sqrt3h^G < r_G/2$ of $1$; hence there is an open neighbourhood of $\operatorname{supp}\zeta$ contained in $\mathcal R'$ on which $\Delta + \lambda$ is invertible, and $w_\alpha$ is real analytic there, by Lemma~\ref{lem:flux-derivatives}. On $\operatorname{supp}\zeta$ we write $\Xi_\alpha = \langle w_\alpha, \nabla\mathsf A\rangle - \mathfrak e_\alpha$, and Lemma~\ref{lem:integration-by-parts} with $F = F_\alpha$, whose support lies in $\mathcal K \subset \mathcal O$, and $v = w_\alpha$ gives, after summation over $\alpha$ and with $\sum_\alpha\langle w_\alpha, \nabla X^\alpha\rangle = c_\lambda$,
\begin{equation}\label{eq:moments-ibp}
    \sum_\alpha\langle F_\alpha\Xi_\alpha\rangle = g^2\langle\zeta c_\lambda\rangle + g^2\sum_\alpha\Bigl(\bigl\langle X^\alpha\langle w_\alpha, \nabla\zeta\rangle\bigr\rangle + \bigl\langle F_\alpha\operatorname{div}w_\alpha\bigr\rangle + \bigl\langle F_\alpha\langle w_\alpha, \nabla\mathcal A'\rangle\bigr\rangle\Bigr) - \Bigl\langle\zeta\sum_\alpha X^\alpha\mathfrak e_\alpha\Bigr\rangle\;.
\end{equation}
On $\operatorname{supp}\zeta$ we have $\abs{X^\alpha} \leq 2h^G$, $\norm{w_\alpha} \leq 2C_\Xi P/\lambda$, and $\norm{\nabla\mathsf A} \leq (4C_\Delta N_q^4T_{\mathrm{ng}})^{1/2}g$, by Lemma~\ref{lem:flux-derivatives}(i) and (iii), whose condition holds there because the plaquettes are within $\varepsilon_\zeta$ of $1$.

\emph{The four error terms.} The gradient $\nabla\zeta = \vartheta'(\mathsf A/(T_{\mathrm{ng}}g^2))\nabla\mathsf A/(T_{\mathrm{ng}}g^2)$ vanishes on $\mathcal V$, so the first is at most $g^2d_{\mathfrak g}2h^G(2C_\Xi P/\lambda)\norm{\vartheta'}_\infty(4C_\Delta N_q^4/T_{\mathrm{ng}})^{1/2}g^{-1}\bar g^8 \preceq \bar g^{10}$. For the second, $\operatorname{div}w_\alpha = \sum_{b,\beta}\langle e_{b,\beta}, \partial_{e_{b,\beta}}w_\alpha\rangle$, and differentiating the resolvent as in the proof of Proposition~\ref{prop:ibp-error}, with $\norm{e_{b,\beta}}_\infty \leq N^{1/2}$,
\begin{equation}\label{eq:divergence-bound}
    \abs{\operatorname{div}w_\alpha} \leq 4N_q^4d_{\mathfrak g}N^{1/2}C_\Xi\Bigl(\frac{2P^2}{\lambda} + \frac{4C_\Delta N_q^4P}{\lambda^2}\Bigr)\;,
\end{equation}
so that the second term is $\preceq \bar g^3$. For the third, Lemma~\ref{lem:density-rest} and $\norm{w_\alpha}_\infty \leq N^{1/2}\norm{w_\alpha}$ bound it by $g^2d_{\mathfrak g}2h^GM_EN_q^4N^{1/2}(2C_\Xi P/\lambda)/r_E$, and $r_E \geq c\bar g(\log\bar g^{-2})^{q_0}$ by the bound on $r_{i,K'}$; with the sizes of Section~\ref{sec:flux-moments} it is at most $C\bar g^2(\log\bar g^{-2})^{11r+1/2+p_\lambda-q_0}$. The last term of \eqref{eq:moments-ibp} is bounded by Proposition~\ref{prop:ibp-error}, since $0 \leq \zeta \leq 1$.

\emph{The main term.} On $\operatorname{supp}\zeta$, in the gauge of Proposition~\ref{prop:near-flat}, $c_\lambda(V) = c_\lambda(V_1)$ and $V_1 = e^{iu'_1}W'$, and we follow $c_\lambda$ along the fields $e^{i\sigma u'_1}W'$, which lie in $\mathcal R'$ and satisfy $\norm{(\Delta + \lambda)^{-1}} \leq 2/\lambda$ by Proposition~\ref{prop:near-flat}(iii). Differentiating each of the three factors of \eqref{eq:c-lambda} by Lemma~\ref{lem:flux-derivatives}(i), (ii),
\begin{equation}\label{eq:c-lambda-variation}
    \abs{c_\lambda(V) - c_\lambda(W')} \leq d_{\mathfrak g}C_\Xi^2\Bigl(\frac{4P^3}{\lambda} + \frac{4C_\Delta N_q^4P^2}{\lambda^2}\Bigr)\delta_* \preceq \bar g^{\varsigma_*}\;.
\end{equation}
By Proposition~\ref{prop:flat-covariance}(i), whose last two terms tend to $0$ with $\bar g$ because $t_{\mathrm b} \geq (\log\bar g^{-2})^r$ and $\lambda t_{\mathrm b}^2(s_{\mathrm b} + 2)(t_{\mathrm b} + 2) \leq N_q^2(\log\bar g^{-2})^{-p_\lambda} \preceq (\log\bar g^{-2})^{2r-p_\lambda}$, we get $c_\lambda \geq d_{\mathfrak g}s^2t^3/(112e^{1/2}T^5) > 0$ on $\operatorname{supp}\zeta$ for $\bar g$ small. Since $\zeta \geq \mathbf 1_{\mathcal V}$ and $\langle\mathbf 1_{\mathcal V}\rangle \geq 1 - \bar g^8$, $g^2\langle\zeta c_\lambda\rangle \geq (1 - \bar g^8)g^2d_{\mathfrak g}s^2t^3/(112e^{1/2}T^5)$.

\emph{Conclusion.} The exponents $7r + 1 - p_\lambda/2$ and $11r + \frac12 + p_\lambda - q_0$ are negative by the choice of $p_\lambda$, so all the error terms in \eqref{eq:moments-cutoff} and \eqref{eq:moments-ibp} are $o(\bar g^2)$ as $\bar g \to 0$, uniformly in $K \geq K_{\mathrm w}$. With $g^2 \geq \bar g^2/2$,
\begin{equation}\label{eq:moments-conclusion}
    \Psi_3 - \Psi_2 \geq \frac{C_2(\pi)}{d_{\mathfrak g}}\Bigl((1 - \bar g^8)\frac{\bar g^2}{2}\cdot\frac{d_{\mathfrak g}s^2t^3}{112e^{1/2}T^5} - o(\bar g^2)\Bigr) \geq \frac{C_2(\pi)s^2t^3}{224e^2T^5}\bar g^2 = 3c_O\bar g^2
\end{equation}
for $\bar g$ small.
\end{proof}

\subsection{Proof of Theorem~\ref{thm:non-gaussian}}\label{sec:non-gaussianity-proof}

We now collect the pieces. The shape term is reduced to $\Psi_3 - \Psi_2$ by Propositions~\ref{prop:shape-stopping} and~\ref{prop:insertions-cancel} and Lemma~\ref{lem:lasso}, with errors of the orders $\bar g^3$ and $\bar g^4$ up to logarithms, Proposition~\ref{prop:flux-moments} bounds $\Psi_3 - \Psi_2$ below, and Theorem~\ref{thm:pair-terms} bounds the pair terms.

\begin{proof}[Proof of Theorem~\ref{thm:non-gaussian}]
Let $\bar g$ be so small that all the results of this section hold, in particular $\bar g \leq \bar g_1'$.

(i) This is Proposition~\ref{prop:normalisations-drop-out}(ii).

(ii) Let $K \geq K_{\mathrm w}$. By Proposition~\ref{prop:shape-stopping}, the decomposition $\xi_{R_n} = (f^{(R_n)}_{K'} - 1) + I(R_n)$, Proposition~\ref{prop:insertions-cancel}, and Lemma~\ref{lem:lasso}, in which $\Psi_3 - \Psi_2$ is real,
\begin{equation}\label{eq:shape-lower}
    \mathfrak y_S(W_K) \geq \Psi_3 - \Psi_2 - \abs{r_H} - 8\max_n\mathsf e(R_n) - \bar g^9 - 8\max_n\sigma(R_n)^2\;.
\end{equation}
Here $\abs{r_H} \leq 512\,\mathsf l^3 + 16\bar g^8$ with $\mathsf l = 1.05c_\pi h^G \leq C\bar g(\log\bar g^{-2})^{4r+1/2}$, $8\max_n\mathsf e(R_n) \leq C\bar g^3(\log\bar g^{-2})^{m'_\sigma}$, and $\sigma(R_n) \leq C_\sigma\bar g^2(\log\bar g^{-2})^{m_\sigma}$ by Lemma~\ref{lem:stopping-sizes}(iv), so that with Proposition~\ref{prop:flux-moments} the right side is at least $2c_O\bar g^2$ for $\bar g$ small, for every $K \geq K_{\mathrm w}$. By Proposition~\ref{prop:normalisations-drop-out}(i) and (ii), $\mathfrak y_S(W) = \lim_K\mathfrak y_S(N_K) = \lim_K\mathfrak y_S(W_K) \geq 2c_O\bar g^2$. By Theorem~\ref{thm:pair-terms}, $\abs{\mathfrak y_k(W)} \leq 16C_\sigma^2\bar g^4(\log\bar g^{-2})^{2m_\sigma} \leq \frac12c_O\bar g^2$ for $\bar g$ small, and \eqref{eq:witness-weights} gives $\mathfrak Y(W) = \mathfrak y_S(W) - \mathfrak y_2(W) + \mathfrak y_3(W) \geq c_O\bar g^2 \geq \frac12c_O\bar g^2$.

For a level $q' \geq q$ the eleven diagrams are ensembles at the level $q'$, with admissible loops, and $q' \geq q \geq q_\sharp(c_\pi, \bar g)$, so that Theorem~\ref{thm:main} gives the limits $W^{(q')}(\mathcal D)$. By \eqref{eq:level-change}, $W^{(q')}_K(\mathcal D)/W_K(\mathcal D)$ is a product of the $j_{q'} - j_q$ step factors $\kappa_k(C, \chi)$, $K - j_{q'} < k \leq K - j_q$, of the at most two loops of $\mathcal D$. By Definition~\ref{def:step-factor}, the exponent of such a factor is bounded by $\sum_vN^{(k-1)}_v(C)\abs{\kappa^{(k)}_v/\Pi_k + \mathbf 1_{\mathbf R}(k)\mu^{\mathbf R}_{k,v}} + \abs{\mathcal K(C)}\max_{v,w}\abs{\zeta^{(k)}_{v,w} + \mathbf 1_{\mathbf R}(k)\zeta^{\mathbf R}_{k,v,w}}$, where the coefficients are bounded by $C_{S'}\bar{\mathfrak a} + C^R_E$ and $C_\zeta\bar{\mathfrak a} + C^R_E$, by the proof of Proposition~\ref{prop:normalisation-structure}, and $N^{(k-1)}_v(C) \leq P_q(C)a_{K'}L^{K'-k+1} \leq P_q(C)a_{K'}L^{j_{q'}-j_q}$, which does not depend on $K$. So $\abs{W^{(q')}_K(\mathcal D)} \geq e^{-C'}\abs{W_K(\mathcal D)}$ with $C'$ independent of $K$, and $W^{(q')}(\mathcal D) \neq 0$ by (i). Proposition~\ref{prop:normalisations-drop-out}(iii) gives $\mathfrak Y(W^{(q')}) = \mathfrak Y(W)$.

(iii) By Proposition~\ref{prop:normalisations-drop-out}(ii) and part (ii), $\mathfrak Y^{\mathrm{sym}}(W) = \mathfrak Y(W) > 0$, and $W$ does not vanish on the eleven diagrams by (i). Proposition~\ref{prop:gaussian-models}, on the other hand, gives $\mathfrak Y^{\mathrm{sym}} = 0$ for the function of every Gaussian charge model and for every limit of such functions which does not vanish on the eleven diagrams, so $W$ is neither.
\end{proof}

\section{Notes on the Papers}\label{sec:notes-on-papers}

We record the places at which Ba{\l}aban's papers sketch or defer an argument which we use, the places at which we read a printed formula otherwise than it is printed, and the conventions in which our notation differs from his.

\emph{Deferred and sketched arguments.} The asymptotic behaviour of the couplings, \cite[Theorem~2, (0.31), p.~259]{BalabanRG1}, is stated with the remark that ``a proof of this theorem, based on perturbative calculations, will be given in a separate paper''; the bounds \cite[(2.6), (2.7), p.~255]{BalabanCRE} on the running couplings are stated to ``follow from the renormalization group equations (0.20) [I]''. We do not use these statements: Section~\ref{sec:coupling-flow} proves what we need for the construction of this paper. The bounds on the walk factors of the propagators are given ``rather briefly'' \cite[p.~409]{BalabanPBG}, and the bounds on the factors of the walk expansions are stated up to powers of the constants \cite[p.~413]{BalabanPBG}; we use them as printed, and the comparison of two fine lattices in Section~\ref{sec:walk-factor-difference} is proved in Supplement, Proposition~S.12. Ba{\l}aban sketches a second proof of \cite[Theorem~2]{BalabanCRE} by gauge invariance \cite[pp.~283--284]{BalabanCRE}; we use the first.

\emph{Formulae read otherwise than printed.} At a few places a printed reference or formula cannot be meant as it stands, because it contradicts another printed line or the computation into which it enters. At each of them exactly one reading agrees with the rest of the text, and we use that reading; for each we give the printed text, the lines which exclude it, and the place at which our argument depends on it.

\emph{References to displays.} Three references point to displays which do not exist or which say something else, and in each case the context names the display that is meant. Theorem~1 of \cite[p.~262]{BalabanCRE} assumes that the couplings satisfy ``the inequalities (I.0.33)'', and the second choice of the threshold of the first step replaces $g_0$ in the logarithmic factor by a bound $\gamma_0 \geq g_0$ ``coming from the inequality (0.33) [I]'' \cite[p.~246]{BalabanCRE}. The first paper of the series has no display (0.33), since the numbered displays of its Section~0 end with (0.32) on p.~260. The only one of them which bounds the sequence of the couplings is the two-sided bound (0.31) of \cite[Theorem~2, p.~259]{BalabanRG1}, whose left inequality at $k = 0$ gives $g_0 \leq \gamma_0$ once $g \leq \gamma$, and Ba{\l}aban refers to it elsewhere as ``the basic inequality (0.31)'', ``the result of Theorem~2'' \cite[p.~355]{BalabanLFR2}, there with the second paper of the series printed as its source. We therefore read both references as (0.31). Theorem~\ref{thm:balaban-inductive} is quoted with this hypothesis, and Supplement, Proposition~S.1, verifies it for our construction in the form of the inequalities (2.6)--(2.9) of \cite[pp.~255--256]{BalabanCRE}, which the proof uses in its place \cite[p.~255]{BalabanCRE}. Next, the estimates on the propagators in a background field use ``the weighted distance $d(y,y')$ defined by (2.36) in [4]'' \cite[p.~397]{BalabanPBG}, where [4] is \cite{BalabanPRT2}. The display (2.36) of that paper, on p.~229, is the partition of unity $\sum_\square h_\square^2 = 1$, whereas the weighted distance of two points of the set $\mathfrak B$ of the blocks of all scales, which is the set in which \cite{BalabanPBG} takes its points, is defined in (2.46), p.~231. We read (2.46); this is the distance in the bounds on the walks of Section~\ref{sec:walk-expansions} and in Supplement, Lemma~17.3. Finally, the gauge transformation $u_1$ of \cite[(124), p.~296]{BalabanVAR}, which is determined by $U_1$ through the conditions $\overline{R_0u_1}^j = 1$ on $\Lambda_j$, ``is given by the formula (106) [5]'' \cite[pp.~296--297]{BalabanVAR}, where [5] is \cite{BalabanPBG}. That paper numbers all its displays in the form (3.$n$), and its (3.106), p.~414, is the Neumann series of a propagator, while the display (106) of \cite[p.~294]{BalabanVAR} itself is an identity between quadratic forms. The formula meant is (106) of \cite[p.~33]{BalabanAVG}, the reference [4] of the same paper, which writes the gauge transformation determined by $U_1$ explicitly through the averages of $U_1$ \cite[pp.~32--33]{BalabanAVG}. We read it so in Supplement, Lemma~17.14, and in its application in Section~\ref{sec:flux-moments}.

\emph{The sections of [I].} The paper on the complete model cites the cluster expansion of the fluctuation integral as ``Sect.~7 [I]'', with the formulae (I.7.1)--(I.7.13), and the localisation of the action as ``Sect.~6 [I]'', with (I.6.34) and (I.6.41) \cite[pp.~275--278]{BalabanCRE}, although its reference [I] is the first paper of the series \cite{BalabanRG1}, which has no such sections, and it lists the second paper \cite{BalabanRG2} separately. The material is that of \cite{BalabanRG2}: its formulae (2.1), the Mayer expansion, (2.11)--(2.13), the polymer expansion and its exponentiated form, (1.34), the space of analyticity of the form of \cite[(3.43), p.~276]{BalabanCRE}, and (1.41), the localised action, are the formulae cited as (I.7.1), (I.7.11)--(I.7.13), (I.6.34), and (I.6.41), and the bound of \cite[(2.14)]{BalabanRG2} is printed there as a bound of ``(7.14)'' \cite[p.~15]{BalabanRG2}. We read a reference (I.7.$n$) as \cite[(2.$n$)]{BalabanRG2} and a reference (I.6.$n$) as \cite[(1.$n$)]{BalabanRG2}; Section~\ref{sec:marked-polymers} depends on this reading.

\emph{Thresholds of scaled fields.} Ba{\l}aban's fluctuation fields are scaled by the coupling and carry Gaussian weights of order one: the field $A$ of the first step enters the configuration as $\exp ig_0A$ \cite[(1.19), p.~250]{BalabanCRE}, and the field $A_k$ of the step $k + 1$ as $\exp i[g_kCA_k - h\tilde D(g_kCA_k)]$ \cite[(3.15), p.~268]{BalabanCRE}. A threshold $\delta$ for configurations is therefore the threshold $g^{-1}\delta$ for the field, and the first step prints it so, pairing $\abs{A(b)} < g_0^{-1}\delta_0$ with $\abs{\exp ig_0A(b)\,U_1(b)U_{1,\square'}^{-1}(b) - 1} < 2\delta_0$ \cite[(1.18), (1.19), p.~250]{BalabanCRE}. The decomposition of unity of the general step \cite[(3.16), p.~268]{BalabanCRE} prints $\abs{A_k(b)} < \delta_k$ instead, although $\delta_k = g_k(A_1/A_0)p_0(g_k)$ \cite[p.~265]{BalabanCRE} is a threshold for configurations, with which \cite[(3.19), p.~268]{BalabanCRE} compares $\abs{V^{(k)}(b)V^{(k)}_{\square'}(b)^{-1} - 1}$. As printed, the large-field function $\chi^{(k)c}(R_{k+1})$ of (3.16) would select every field above a positive power of $g_k$, where the Gaussian weight is close to $1$, whereas \cite[p.~380]{BalabanLFR2} draws from it the small factor $\exp(-\gamma_0g_j^{-2}\delta_j^2(LM_2R_j)^{-d}\abs{R_j})$, which is the Gaussian weight at the field $g_j^{-1}\delta_j$. We therefore read the threshold of (3.16) as $g_k^{-1}\delta_k$, which is our $\vartheta_k/g_k$, and Definition~\ref{def:construction}(ii) states it as the condition $\abs{g_kA_k(b)} < \vartheta_k$ on the unscaled field. The field $B_j$ of the $\mathbf R$-operation is in the same position. It carries the Gaussian weight $\exp(-\frac12\langle B_j, C^*\Delta^{(j)}CB_j\rangle)$ and enters the configuration as $\exp i[g_jCB_j - h\tilde D(g_jCB_j)]$ \cite[(1.60), pp.~188--189]{BalabanLFR1}, while $\delta'_j = g_jA_1p_1(g_j)$ \cite[p.~183]{BalabanLFR1} is a threshold for configurations \cite[(1.27), p.~182]{BalabanLFR1}. The restriction $\abs{B_j(b)} < \delta'_j$ printed in \cite[(1.55), p.~188]{BalabanLFR1} would confine the field to a ball whose radius is a positive power of $g_j$, so that the small-field function $\chi'^{(j)}$ would carry almost none of the Gaussian mass; the restriction of the field $A_j$ in the same paper is printed as $\abs{A_j(b)} < g_j^{-1}\delta_j$ \cite[(1.9), p.~178]{BalabanLFR1}, the condition paired with (1.55) is $\abs{\exp ig_jB_j(b) - 1} < 2\delta'_j$ \cite[(1.59), p.~188]{BalabanLFR1}, and \cite[p.~383]{BalabanLFR2} bounds a scaled field below by $g_j^{-1}\delta'_j = A_1p_1(g_j)$. We read the threshold of (1.55) as $g_j^{-1}\delta'_j$, which Definition~\ref{def:construction}(ii) states as $\abs{g_jB_j(b)} < \delta'_j$. The value of $\delta_0$ itself is printed as the chain $\delta_0 = g_0A_1p_0(g_0) = (A_1/A_0)\varepsilon_0$ \cite[p.~247]{BalabanCRE}, whose two members differ by the factor $A_0$ in the convention $p_0(g) = A_0(\log g^{-2})^{p_0}$ of \cite[p.~246]{BalabanCRE} and agree in the convention of \cite[pp.~380--383]{BalabanLFR2} discussed below. The second member is the one with which the text goes on: the next sentence assumes that $A_0$ is much larger than $A_1$, the general step defines $\delta_k = g_k(A_1/A_0)p_0(g_k) = (A_1/A_0)\varepsilon_k$ \cite[p.~265]{BalabanCRE}, and the estimates (3.7) and (3.9) of \cite[p.~266]{BalabanCRE} replace $\delta_k$ by $(A_1/A_0)\varepsilon_k$ and use that $A_1/A_0$ is small. We take $\delta_0 = (A_1/A_0)\varepsilon_0$, which is $\vartheta_0$ of Definition~\ref{def:tube-parameters}(ii) and fixes the radius $\vartheta_0/g_0$ of the fluctuation integral of the first step in Section~\ref{sec:first-step-proof}.

\emph{Two exponents.} In the displays of \cite[pp.~380--383]{BalabanLFR2} the symbol $p_0(g)$ stands for $(\log g^{-2})^{p_0}$, without the factor $A_0$ of \eqref{eq:balaban-threshold}, in the same way as $p_1(g) = (\log g^{-2})^{p_1}$ \cite[p.~183]{BalabanLFR1}. The constants are written out there: the factor of a component of a large-field region is $\exp(-\gamma_0\min\{\frac12B_3^{-2}A_0^2, 2A_1^2, A_1^2\}p_0^2(g_j)(d'_j + 1))$ \cite[p.~381]{BalabanLFR2}, in which $A_0$ comes from the plaquette thresholds $\varepsilon_j$ and $A_1$ from the thresholds $\delta_j$ of the factors on p.~380, and the bound of $\exp(-\frac14g_m^{-2}B_3^{-2}\frac12\varepsilon_m^2)$ by $\exp(-\frac18B_3^{-2}(1 + \beta_0)^{-2}A_0^2p_0^2(g_j))$ on p.~381 follows from $g_m^{-2}\varepsilon_m^2 = A_0^2(\log g_m^{-2})^{2p_0}$ and \cite[(2.7), p.~255]{BalabanCRE} in this reading only, since in the convention of \eqref{eq:balaban-threshold} its right side would carry the factor $A_0^4$ in the exponent and its left side the factor $A_0^2$. Supplement, Proposition~S.1, takes its conditions on $A_1$ and $A_0$ from these displays in this reading. The second exponent is printed $r_0$. It occurs twice, in the assumptions $\nu \geq 2p_0 + dr_0$ \cite[p.~363]{BalabanLFR2} and $2p_1 - (d + 5)r_0 > p_0$ \cite[p.~383]{BalabanLFR2}, and it is defined nowhere in the series. In both places it is the exponent of the cube sizes. On p.~383 the assumption serves to bound the factor $\exp(-R_j^{-d-5}p_1^2(g_j))$ by $\exp(-p_0(g_j))$, and since $(\log g_j^{-2})^r \leq R_j < L(\log g_j^{-2})^r$ by \cite[(2.5), p.~255]{BalabanCRE}, this bound holds for small $g_j$ when $2p_1 - (d + 5)r > p_0$ and fails when $2p_1 - (d + 5)r < p_0$; on p.~363 the factor $R_k^{d-1}$ of the bound is to be absorbed by $L^{-(d-1)N} \leq (\log g_k^{-2})^{-(d-1)\nu}$. We read $r_0$ as $r$, in the condition on $p_1$ of Appendix~\ref{app:constants} and in the condition on the depth of Section~\ref{sec:depth}.

\emph{Iterated averages.} The averaging operation of order $k$ enters \cite[(1.1), p.~260]{BalabanRG1} as $\bar U^k = M^k(U)$, while the first paper of the series defines only the average $M$ of order one \cite[(0.4), p.~253]{BalabanRG1}. The averages of higher order are defined in \cite[(43), p.~24]{BalabanAVG} recursively, the average of order $k + 1$ being the average of order one of the average of order $k$, so that $M^k$ is the $k$-fold iterate of $M$, and we use it in this sense.

\emph{Two averages of one step.} The linearised averages in the operators of \cite{BalabanPBG} are the linear parts of the average $\bar{\bar U}'$ of the fluctuation relative to the background \cite[(3.13), p.~392; (3.113), p.~418]{BalabanPBG}, \cite[(121)--(124), p.~36]{BalabanAVG}, while the constraints of the minimisers are those of the average $M$ itself \cite[(3), p.~278]{BalabanVAR}. The two linear parts differ by a gauge transformation of the coarse sites, which does not vanish at a flat background, and the papers pass from one to the other without a remark, which is harmless for gauge-invariant quantities. Where both occur we distinguish them, as $Q_1$ and $Q^{\mathrm{pl}}_1$ in Section~\ref{sec:galerkin} and as $Q^{\mathrm{dr}}$ and $Q$ in Section~\ref{sec:flat-covariance}; the relation between them is Supplement, Proposition~S.12(j).

\emph{The spaces of the regular terms.} The first paper of the series proves its inductive description on spaces $U^c_j(X, \alpha_0, \alpha_1)$ with numbers $\alpha_0$, $\alpha_1$ which do not depend on the coupling \cite[Theorem~3, p.~264]{BalabanRG1}, and the paper on the complete model states the properties of the regular terms on the smaller spaces with the radii \cite[(2.28), p.~259]{BalabanCRE}, which are proportional to the coupling up to logarithms. On the smaller spaces a second derivative at the trivial configuration is not bounded uniformly in the coupling by the bound of the term, and the coefficients of the flow are second derivatives. We use Theorem~3 of \cite{BalabanRG1} for the part of the regular terms which does not involve the $\mathbf R$-terms, along the couplings of a history, which need not satisfy the equation \cite[(0.20), p.~256]{BalabanRG1} of that paper; its proof treats one step at a time and uses of the couplings only that they are small (Supplement, Lemma~3.5).

\emph{Free choices.} The papers centre the cubes of the families $\pi_j$ on a lattice of spacing $M$ \cite[p.~257]{BalabanRG1} and leave the position of the other partitions free; we centre all of them at the origin (Section~\ref{sec:reference-gas}), which is compatible with the blocks of the averages \cite[(0.3), p.~252]{BalabanRG1} because all sides are odd, and which makes every partition invariant under the reflections $x_\mu \mapsto -x_\mu$. Every object of a step which the papers define by reference to a coordinate system, the axial contours of the fluctuation fields, the gauge-fixing graphs, and the designated bonds of the averages, is attached to its block or cube through the centre and the fixed orientation of the axes, so that a translation which preserves the partitions maps the object of a cube to the object of its image. This is the rule under which Lemma~\ref{lem:locality}(ii) holds; the reflections $\theta_\mu$ of Section~\ref{sec:reflection-positivity} map the objects of a cube to those of its image because the averages use the contours of all orders \cite[p.~252]{BalabanRG1}, and the axial contours, which start at the centre of the block and change the coordinates in a fixed order \cite[(1.7), p.~18]{BalabanPRT1}, \cite[p.~275]{BalabanRG1}, are mapped to the axial contours of the image block by a reflection of one coordinate, which does not change that order. A permutation of the axes changes that order, which is why the covariance of the normalisations is proved for the reflections and not for the exchange of two axes (Section~\ref{sec:symmetries-proof}).

\emph{Notation.} We write $\vartheta_k$ for the threshold of the fluctuation field which Ba{\l}aban writes $\delta_k$ \cite[p.~265]{BalabanCRE}, because $\delta_k$ is the difference of the couplings of the two runs in Sections~\ref{sec:walk-factor-difference}--\ref{sec:cauchy}. The constant of the boundary terms which \cite[(2.47), p.~264]{BalabanCRE} prints $B_1$ is our $B_\partial$. The exponent of the age bound \eqref{eq:balaban-theorem-two} is written $\beta < 1$ in \cite[Theorem~2, p.~263]{BalabanCRE}, and its proof gives it as $1 - \beta$ with an arbitrary $\beta > 0$, ``with $1 - \beta$, $\beta > 0$, instead of $\beta < 1$'' \cite[p.~283]{BalabanCRE}; we write the exponent as $1 - \beta_E$, with a small $\beta_E > 0$, because $\beta$ is also Ba{\l}aban's constant of the spaces of the boundary terms. In \cite[(2.5), p.~255]{BalabanCRE} the letter $r$ denotes both the exponent of the powers of $L$ among which $R_j$ is chosen and the fixed exponent of $\log g_j^{-2}$; we keep $r$ for the second. Our couplings, thresholds, cube sizes, and depths are those of the tube of a target (Definition~\ref{def:construction}), evaluated at $\hat g_j$, where Ba{\l}aban's are evaluated at $g_j$; Proposition~\ref{prop:construction-admissible} states that his bounds hold for them.

\appendix

\section{The Constants in the Order of Choices}\label{app:constants}

The estimates of the paper contain many constants, and their consistency rests on the order in which they are chosen: a smallness condition can be met only if the quantity which it constrains is chosen after the constants on which the condition depends. We fix the constants in seven stages, numbered from $0$ to $6$, so that a constant of a stage depends only on the constants of the earlier stages. Every condition of the paper is then of one of two kinds, a lower bound on a quantity of its own stage in terms of earlier ones, or an upper bound on a quantity of a later stage in terms of earlier ones, and conditions of both kinds are met by taking each quantity large or small enough when its stage comes. The \emph{constants of the construction} are those of the stages $0$ to $5$, together with the constants computed from them. The last stage chooses the bound $\gamma$ of Section~\ref{sec:tube}, and with it the coupling $\bar g$ and the target $(K, \bar g)$; every requirement that $\bar g$ be sufficiently small which involves no representation is a condition of this stage.

\emph{Stage $0$: the dimension, the group, and the exponents.} These numbers are fixed once and for all.
\begin{center}
\begin{tabular}{p{0.24\textwidth}p{0.68\textwidth}}
\hline
$d = 4$, $G \subset U(N)$ & the dimension and the gauge group, with the numbers $\kappa$ of \eqref{eq:killing-constant}, $c$ of \eqref{eq:one-loop-coefficient}, and $\dim G$, the radius $r_G$ of the logarithmic chart, the constant $c_G$ of the Haar measure of balls (Section~\ref{sec:large-field-weight}), and the constant $c_h$ which compares the two norms of Appendix~\ref{app:conventions} \\
$r = 2$ & the exponent of the cube sizes $R_j$, for which \cite[p.~245]{BalabanCRE} requires $r \geq 2$ \\
$p_0$ & the exponent of the thresholds, $p_0 \geq 5r$ \cite[p.~246]{BalabanCRE} \\
$p_1$ & the exponent of the thresholds of the $\mathbf R$-operation, $p_1 < p_0$ \cite[p.~183]{BalabanLFR1}, $p_0 - p_1 - 1 \geq \ceil{r}/2$ \cite[(1.94), p.~198]{BalabanLFR1}, and $2p_1 - (d + 5)r > p_0$ \cite[p.~383]{BalabanLFR2}, where the exponent printed $r_0$ is read as $r$, the only exponent of the cube sizes in the construction \\
$q_0 = q_1$ & the exponents of the radii of analyticity, integers greater than $1$ \cite[(2.28), p.~259]{BalabanCRE}, with $q_0 > p_0 + 3r$ and $q_0 > 25r + \frac52$ (Section~\ref{sec:flux-moments}) \\
$\nu_L$ & $p_0 + q_1 + 7r + 1$, the exponent of the depth (Definition~\ref{def:depth}) \\
$\beta_0 = 1/7$ & the constant of Ba{\l}aban's inequalities (2.6), which ``can be chosen arbitrarily small, if $g$ is sufficiently small'' \cite[p.~255]{BalabanCRE}; it satisfies $\beta_0 \leq \sqrt2 - 1$ (Section~\ref{sec:tube}) and $(1 + \beta_0)^8 \leq 3 < L/4$ (Section~\ref{sec:contraction}) \\
$\beta = 1/8$ & the constant of Ba{\l}aban's spaces of the boundary terms \cite[(2.34)--(2.39), p.~261]{BalabanCRE}, a small positive number ``but not too small'' \cite[p.~261]{BalabanCRE}; with it $\kappa_s = \beta/30 = 1/240$ in \eqref{eq:soft-width} \\
$\varepsilon_0$ & $0 < \varepsilon_0 < 1/125$, the loss in the exponents of the comparison of the two runs (Proposition~\ref{prop:minimiser-distance}) \\
$\beta_H$, $\beta_E$, $\beta'$, $\theta_V$ & $\beta_H \geq 19/20$, $\beta_E \leq 1/20$, and $\beta'$, $\theta_V$ of \eqref{eq:age-exponent} \\
H\"older exponent & a number in $[\frac34, 1)$, written $\beta$ in Section~\ref{sec:walk-factor-difference} only, with $\gamma_w$ of \eqref{eq:walk-defect} and $\gamma_H$ of \eqref{eq:gaussian-depth} \\
$\Lambda_*$ & a number $\Lambda_* \geq 2$, the bound on the chord--arc constants of the admissible loops (Definition~\ref{def:admissible}) \\
\hline
\end{tabular}
\end{center}

\emph{Stage $1$: the blocking factor.} The integer $L$ is odd and satisfies \eqref{eq:blocking-factor-condition}, which bounds it from below in terms of $d$, $\varepsilon_0$, and $\Lambda_*$. The constant $C_{\mathrm{dep}} = 16L^4(L + 1)^2$ of Lemma~\ref{lem:depth} is computed here.

\emph{Stage $2$: the constants computed from $d$, $L$, and $G$.} These are Ba{\l}aban's constants which depend on $d$ and $L$ only: the decay rate $\delta_0$ and the constants $B_0$, $B_0(\beta)$ of the propagators \cite[Theorem~3.1, (3.42), (3.43), p.~397]{BalabanPBG}, the constant $c_0$ of \cite[(2.58)]{BalabanPRT2}, and the constant $B_3$ of the minimisers \cite[Theorem~1, p.~279]{BalabanVAR}. With them come the constant $C_W$ of Theorem~\ref{thm:one-loop-window}, the numbers $b_* = c\ln L$ and $b^* = 3c\ln L$ of \eqref{eq:tube-class-constants}, the constant $C_M$ of Theorem~\ref{thm:lf-comparison}, which counts the ingredients of a term per point, and
\begin{equation}\label{eq:history-constant}
    C_P \coloneqq \max(2, 2C_M)\;,
\end{equation}
the constant of Proposition~\ref{prop:large-field-weight}; it satisfies the two requirements $C_P \geq 2$ and $C_P \geq 2C_M$ which Sections~\ref{sec:large-field-weight} and~\ref{sec:large-field-pairs} place on it.

\emph{Stage $3$: the decay rates and the sizes.} First Ba{\l}aban's decay rate $\kappa$ of the inductive bounds, large at fixed $\beta$ \cite[p.~260]{BalabanCRE}, and the exponent $\kappa_0 \geq 13$ of \eqref{eq:r-bound}, which ``can be chosen arbitrarily large, similarly as $\kappa$, if the other parameters are fixed'' \cite[p.~260]{BalabanCRE}; the comparison of two cutoffs uses two of its powers to pay the radius of the spaces of the tube (Section~\ref{sec:coefficient-difference}) and the others for the summability in $K$; then the rate $\kappa_1$ of the decoupling parameters of \cite{BalabanRG2}, which is bounded below in terms of $\kappa$ and $L$. Then the sizes, each bounded below in terms of the earlier constants and the sizes before it, in the order: the size $R_0M_0$ of \cite[p.~408]{BalabanPBG}, enlarged beyond Ba{\l}aban's smallest choice, which leaves his results valid because ``in [4] we have proved all theorems under the assumption that $R$, $M$ are sufficiently large'' \cite[p.~408]{BalabanPBG}; the integer $K$ and the block size $KR_0M_0$ of \cite{BalabanPBG}; the size $M_1$ of the cubes for which the theorems of \cite{BalabanPBG,BalabanVAR,BalabanREG} hold, with $M_2 = L^2M_1$ \cite[p.~245]{BalabanCRE}; the separation $R$ of the walk problems \cite[(2.2)]{BalabanPRT2}; and the cube size $M = L^m$, which is ``sufficiently large'' \cite[p.~245]{BalabanCRE}.

\emph{Stage $4$: the regularity parameters of the propagators.} The parameters $\alpha_0$ and $\alpha_1$ of the domains of \cite[(3.35)--(3.37), p.~396]{BalabanPBG}, with $\alpha_J \coloneqq KR_0M_0\alpha_0$ and the radius $\alpha'_* = \min(\alpha_1, \alpha_J/\hat x_J)$ of Supplement, Proposition~S.12(g), bounded above in terms of the stages $0$--$3$.

\emph{Stage $5$: the parameters of the variational problem.} The parameter $\varepsilon_3$ of \cite[Proposition~4, p.~292]{BalabanVAR}, and then the data parameter $\varepsilon_1$ of \cite[Theorem~1, p.~279]{BalabanVAR} at which its theorems are applied, bounded above in terms of the stages $0$--$4$. In Sections~\ref{sec:large-field-comparison} and~\ref{sec:cauchy} the letter $\varepsilon_1$ denotes the data parameter \eqref{eq:data-parameter} of a history, a function of $\bar g$ which lies below this fixed value once $\bar g$ is small.

\emph{Stage $6$: the coupling.} First Ba{\l}aban's large constants: $A_1$ and then $A_0$ of the thresholds \cite[p.~246, p.~265]{BalabanCRE}, with $A_1/A_0$ small \cite[p.~247]{BalabanCRE}, and $C_0^\alpha$, $C_1^\alpha$ of the radii \cite[(2.28), p.~259]{BalabanCRE}, each bounded below in terms of the stages $0$--$5$. Then the bound $\gamma$, small in terms of everything before it: the conditions \eqref{eq:gamma-conditions} and \eqref{eq:gamma-depth}, $\gamma \leq \gamma_1$ of Proposition~\ref{prop:coefficients}, the conditions (T1) of Appendix~\ref{app:conditions}, and the requirements ``$\bar g$ sufficiently small'' of the statements which involve no representation, where $\bar g \leq \gamma$ by \eqref{eq:tube-class}. The requirements which involve the representations are upper bounds on the number $\gamma_q$ of \eqref{eq:stopping-coupling}, and they are met by the choice of the level $q$ of the ensemble, which comes after the coupling. The integer $K$ and the tube are chosen last. The thresholds $\varepsilon_j$, $\vartheta_j$, $\delta'_j$, the radii $\alpha_{0,j}$, $\alpha_{1,j}$, the cube sizes $R_j$, and the depths $N_k$ of Definition~\ref{def:tube-parameters} and Definition~\ref{def:depth} are functions of the tube.

\emph{The derived constants.} Every other constant of the paper is computed from those above and is introduced where it first occurs. They fall into five groups, which differ in what they may depend on.
\begin{enumerate}[label=(\alph*)]
\item The constants of Ba{\l}aban's theorems, $E_0$, $E_1$, $R_1$, the constant $B_0$ of the boundary terms, and $B_\partial$ (Section~\ref{sec:density}), depend on the constants of the construction and not on the tube, the step, $K$, or the torus (Proposition~\ref{prop:construction-admissible}). The constants $E_\pm$ of the ultraviolet stability \eqref{eq:ultraviolet-stability} are not of this kind: they are independent of the lattice spacing and of the torus, and they depend on the tube and the step through the coupling $g_k$, as in Ba{\l}aban's statement. The constants $C_\Lambda$, $C_3$, $C_{\mathrm{rem}}$, $C_R$, and $\gamma_1$ of Proposition~\ref{prop:coefficients} are of the same kind.
\item The constants of the inductive class (Definition~\ref{def:dressed-loop}): the constants $C_\tau$, $c_\tau$, $C_\rho$, $C_t$ of \eqref{eq:class-parameters}, the box constant $C_{\mathrm{box}} = 2916$, the reserve $\nu$ of the first step \cite{WilsonFirstStep}, taken with $\nu \leq \eta(\mathcal C)$ of the proof of Theorem~\ref{thm:first-step}, the extent rate $\nu_e \coloneqq \nu/(2C_{\mathrm{box}})$ and the cross rate $\nu_\times \leq \nu_e$, the rates $\eta_w$, $c_{\mathrm{att}}$, $c'_*$, $\kappa_e$, $c_{\mathrm{th}}$, and the constants $C_\chi$, $C_{\mathrm{new}}$, $C_{\mathrm{rel}}$, $C_{\mathrm{cr}}$, $C_{S'}$. For ensembles with admissible loops they depend on the constants of the construction, among them $\Lambda_*$, and on the representations through $c_{\mathcal C}$, as nondecreasing functions of it, and on nothing else (Definition~\ref{def:dressed-loop}(v)).
\item The constants of the stopping scale, $C_{\mathcal E}$, $C_\partial$, $C_\lambda$, $C_{\mathrm{fs}}$, $C_A$, $C_e$, and $C_{\mathrm{ent}}$ of Section~\ref{sec:stopping-scale}, which are of the kind (a), and the constant $c_\pi$ of \eqref{eq:representation-constant}, which depends on the representation. The numbers $B_E(\mathcal C)$ and $I_*(\mathcal C)$ of Section~\ref{sec:uniform-bound} are not constants of a stage: they are bounds attached to an ensemble, which depend on $\bar g$ and not on $K$.
\item The constants of the comparison of two cutoffs: $C_{\mathrm{df}}(\beta)$ and $C_\square(\beta)$ of Section~\ref{sec:galerkin}, $C_w$ of Theorem~\ref{thm:walk-difference}, $C_G(\varepsilon_0)$ and $C_U$ of Proposition~\ref{prop:minimiser-distance}, $C_f$, $C_{\mathrm{sh}}$, $C_h$ of Section~\ref{sec:large-field-comparison}, $E'_1$, $E_V$, $C_\alpha$ of Proposition~\ref{prop:theorem-two-difference}, $C_\beta$ of Lemma~\ref{lem:coefficient-difference}, and $C_{\mathcal T}$, $C_s$, $C''_S$, $C_a$ of Section~\ref{sec:recursion}. They do not depend on $K$, the step, the lattice spacing, or the torus; some of them depend on $\bar g$ through the stopping index $j_q$ and through the bound $\bar\varepsilon_1$ of \eqref{eq:data-parameter}, as stated where they are introduced.
\item The rates of the sums over $K$: $\tilde\theta$ of \eqref{eq:theta-tilde}, which depends on $\gamma$ through $\varepsilon^*$, and $\ell_K$, $\mu_K$, $w_K$ of \eqref{eq:weight-w}.
\end{enumerate}

\section{Conventions}\label{app:conventions}

\emph{Lattices and units.} The lattice of the scale $j$ of a run with $K$ steps has the spacing $a_j = L^j\varepsilon$, with $\varepsilon = uL^{-K}$ the finest spacing, and a length ``in units of the scale $j$'' is a length divided by $a_j$. Distances are $\ell^\infty$-distances unless stated otherwise. The bonds and the plaquettes of a set are those which intersect it \cite[p.~251]{BalabanRG1}.

\emph{Cubes and partitions.} For a cube $\square$ of the family $\pi_j$, of side $M$, the cube $\tilde\square^n$ is the cube of side $(1 + 2n)M$ with the same centre \cite[p.~257]{BalabanRG1}. The operation $\sim$ applied to a set adds one layer of the cubes of the family in terms of which the set is defined \cite[p.~245]{BalabanCRE}, a cube being added when it has a point in common with the set, and $X^{\sim-1} \coloneqq ((X^c)^\sim)^c$ \cite[p.~247]{BalabanCRE}. Partitions are compatible, every cube of a coarser partition being a union of cubes of a finer one \cite[p.~245]{BalabanCRE}, and every partition consists of cubes centred at the points of $s\mathbb Z^4$, $s$ its side (Section~\ref{sec:reference-gas}).

\emph{Norms.} On matrices $\norm{\cdot}$ is the operator norm (Section~\ref{sec:coarse-loop}), which is invariant under multiplication by unitary matrices; on the insertions of the cross groups, which are tensors with one factor for each interval, it is the projective norm \eqref{eq:projective-norm} built from it; and $\abs{\cdot}$ is the norm in which Ba{\l}aban states the regularity of the fields; the constant $c_h \geq 1$ of Section~\ref{sec:coarse-loop} compares the two. The trace $\tr = N^{-1}\operatorname{Tr}$ is normalised.

\emph{Letters with several meanings.} A few letters keep the meanings which they have in the papers that we cite, and the context separates them.
\begin{enumerate}[label=(\roman*)]
\item $\beta$ is Ba{\l}aban's constant $1/8$ of the spaces of the boundary terms in Sections~\ref{sec:construction}--\ref{sec:stopping-scale} and in the coefficients $1 - \beta(1 - 2^{-(k-m)})$ of the radii; the exponent $\beta < 1$ of the age bound \eqref{eq:balaban-theorem-two}, which is $1 - \beta_E$, in Section~\ref{sec:construction} and in Lemma~\ref{lem:direct-bounds}; and the H\"older exponent of stage $0$ in Section~\ref{sec:walk-factor-difference}.
\item $M$ is the cube size $L^m$ of \cite[p.~245]{BalabanCRE} and \cite[p.~257]{BalabanRG1}, except in Section~\ref{sec:walk-expansions} and in the bounds \eqref{eq:walk-bound}, \eqref{eq:walk-difference}, where it is the block size $KR_0M_0$ of \cite[p.~408]{BalabanPBG}; $K$ is the number of steps, except in $KR_0M_0$.
\item $\kappa$ is the number of Lemma~\ref{lem:invariant-forms} in Section~\ref{sec:coupling-flow} and Ba{\l}aban's decay rate elsewhere.
\item $B_0$ is the constant of the boundary terms \cite[(2.42)]{BalabanCRE} in Sections~\ref{sec:construction} and~\ref{sec:large-field-comparison}, and the constant of the propagators \cite[(3.42)]{BalabanPBG} in Section~\ref{sec:walk-factor-difference}.
\item $\delta'$ without an index is the deviation of the coefficients in Theorem~\ref{thm:coupling-flow} and Proposition~\ref{prop:coefficients}; $\delta'_j$ is the threshold of the $\mathbf R$-operation.
\item $r$ is the exponent of stage $0$; $r(C)$ and $r_0(C)$ are the lengths of a loop of Section~\ref{sec:loops}, $r_k$ the radius of the fluctuation cutoff, and $r_G$ the radius of the logarithmic chart.
\end{enumerate}

\section{The Conditions and Their Sources}\label{app:conditions}

We list the conditions on the constants by the stage of Appendix~\ref{app:constants} to which they belong, with the place where each is used and, where it is Ba{\l}aban's, where he imposes it. Each of them bounds a quantity of its stage from below, or a quantity of a later stage from above, in terms of the earlier stages; so all of them hold once the stages are chosen in order, and none of them involves the number of loops, their sizes, or $K$.

\emph{Stage $1$.} The blocking factor $L$ is an odd integer, as in \cite[p.~251]{BalabanRG1}, with
\begin{equation}\label{eq:blocking-factor-condition}
    L \geq L_0(\Lambda_*)\;,
    \qquad
    L_0(\Lambda) \coloneqq \max\bigl(2^{52},\, 4\Lambda^2L_0^\times\bigr)\;,
    \qquad
    L_0^\times \coloneqq 4 \cdot 10^4 \cdot 640\bigl(2^{d+2} + 8\bigr)\;,
\end{equation}
so that $L_0^\times = 1.8432 \cdot 10^9$ at $d = 4$. The two terms have different origins. The term $2^{52} = 4^{26}$ gives the inequality $L^{5(1-\varepsilon_0)/128} \geq 4$, because $128/(5(1 - \varepsilon_0)) < 128 \cdot 125/620 < 26$ for $\varepsilon_0 < 1/125$; this inequality is used in Supplement, Proposition~S.14, for the variation of the terms of the effective action between the scales of a young history, where a factor $2^{n}$ from the radii of the analyticity domains has to be beaten by the decay $L^{-5(1-\varepsilon_0)n/128}$. The term $4\Lambda^2L_0^\times$ is the condition of the counts of the inductive step (Section~\ref{sec:contraction}): the shares of the costs of a group of the step in the reserve $\nu$ and in the extent rate $\nu_e$ of Definition~\ref{def:dressed-loop} grow linearly in the chord--arc constant $\Lambda$ through two factors, the counts of Lemma~\ref{lem:square-properties} and the extent of the insertions made by the $\mathbf R$-operation, and they remain within their margins when $L \geq 4\Lambda^2L_0^\times$, and this bound contains the lower bounds $5000\Lambda$, $3(1 + \Lambda)$, and $20\Lambda$ of the same counts. Since $L_0$ increases and every admissible loop has $\Lambda(C) \leq \Lambda_*$, the condition \eqref{eq:blocking-factor-condition} covers all of them at once. It covers as well the other lower bounds on $L$ of the paper and of the supplement, all of which are below $2^{52}$: $L > 11$ \cite[p.~251]{BalabanRG1}, $L \geq 3$ (Lemma~\ref{lem:depth}), $L \geq 5$ in \eqref{eq:dependence-range}, $(1 + \beta_0)^8 \leq L/4$ (Section~\ref{sec:contraction}), $L \geq 21$ (Lemma~\ref{lem:coarse-loop}), $L \geq 195$ (Lemma~\ref{lem:torus-condition}), and $L \geq 12d + 6$, $L^2 > 9d$, $L \geq 128(1 + \beta_0)$ (Supplement, Propositions~S.12--S.14).

\emph{Stage $3$.} The rate $\kappa_1$ satisfies $\frac12(\kappa_1 - 1) \geq 2L\kappa$ \cite[p.~20]{BalabanRG2}. The size $R_0M_0$ satisfies the two conditions of Lemma~\ref{lem:shell-reduction}, $\frac16\delta_0R_0M_0 > 8\log c_0(\frac16) + 1$ and $4\delta_0R_0M_0 \geq 3(1 + \beta)\log L$, with the H\"older exponent $\beta$, and the conditions of the same form $c\delta R_0M_0 \geq C\log L + C'$, with a rate $\delta$ of stage $2$ and numbers $c$, $C$, $C'$ depending on $d$ and on the H\"older exponent, of Supplement, Proposition~S.12; Ba{\l}aban's results hold for every such enlargement of $R_0M_0$ \cite[p.~408]{BalabanPBG}. The integer $K$ of \cite[p.~408]{BalabanPBG} is a power of $L$ with $K \geq 4$, and the separation of the walk problems is $R \geq 12dL$. The cube sizes satisfy Ba{\l}aban's lower bounds $M \geq M(\kappa)$ \cite[Theorem~3, p.~264]{BalabanRG1}, $\delta_0M \geq \kappa$ \cite[p.~21]{BalabanRG2}, $\delta_1M \geq \kappa_1$ \cite[p.~5]{BalabanRG2}, and $\delta(M/M_1) \geq 2$ \cite[p.~186]{BalabanLFR1}, and the lower bounds on $M_1$ and $M$ of Supplement, Propositions~S.12--S.14, each of which has the form $c\delta M_1 \geq C\log L + C'$ or $M \geq C(d, L)M_1$.

\emph{Stage $4$.} The parameters $\alpha_0$ and $\alpha_1$ are so small that ``$O(1)M\alpha_0$ is still a sufficiently small number'' \cite[p.~396]{BalabanPBG}, with $M$ the block size $KR_0M_0$, and that the upper bounds of Supplement, Proposition~S.12, hold; these involve the constants of the stages $0$--$3$ only.

\emph{Stage $5$.} The parameters of the variational problem satisfy $\varepsilon_3 \leq a_3$ \cite[Proposition~4, p.~292]{BalabanVAR} and $\varepsilon_1 \leq a_1$, and every cube of the parameter $M_V$ at which \cite[Theorem~1]{BalabanVAR} is applied has $M_V \leq M(\varepsilon_1) = R_1M_1a_1/\varepsilon_1$ \cite[p.~279]{BalabanVAR}, an upper bound on $\varepsilon_1$; the further upper bounds on $\varepsilon_3$ and $\varepsilon_1$ in terms of $\alpha_0$, $\alpha_1$, and the earlier constants are those of Supplement, Propositions~S.12 and~S.13.

\emph{Stage $6$.} The constants of the thresholds satisfy $A_1 \geq 48(15/8)^{1/2}$ and $A_0^2 \geq \frac{15}{8}A_1^2\max(8(1 + \beta_0)^2B_3^2,\, 4O(1)B_3M^2)$, with the constant $O(1)$ of \cite[p.~381]{BalabanLFR2}, which keep the comparisons of \cite[pp.~381--382]{BalabanLFR2} valid when the coefficient of the action is reduced by the factor $15/8$ (Supplement, Proposition~S.1). The bound $\gamma$ satisfies \eqref{eq:gamma-conditions} and \eqref{eq:gamma-depth}, $\gamma \leq \gamma_1$ of Proposition~\ref{prop:coefficients}, the bound $\delta' \leq (2/5)c\ln L$ of Theorem~\ref{thm:coupling-flow}, which holds for $\gamma$ small because $\delta' \to 0$ with $\gamma$, and $C_W\gamma^2 \leq 11/40$ (Supplement, Lemma~2.1). Moreover $\gamma$ satisfies the following conditions, which we call (T1): for every $g \leq (1 + \beta_0)\gamma$,
\begin{equation}\label{eq:t-one}
    \varepsilon^*(g)C_{\mathcal T} \leq \frac14\;,
    \qquad
    4\hat\varepsilon^*(g)N'\Bigl(\frac{g}{1 + \beta_0}\Bigr) \leq 1\;,
    \qquad
    4C_{\mathrm{ov}}\varepsilon^{*B}(g)N_0^+(g)W^\sharp \leq 1\;.
\end{equation}
Here $\varepsilon^*(g)$ is the coefficient $\varepsilon^*_k$ of \eqref{eq:recursion} as a function of the coupling, $C_{\mathcal T} \geq (1 - L^{-\beta'})^{-1}$ the constant of Supplement, Proposition~S.16(e), $\hat\varepsilon^*(g) \coloneqq \sup_{g' \leq 2g}\varepsilon^*(g')$, and $N'(g) \coloneqq \log_L(MR(g)) + 3$, with $R(g)$ the smallest power of $L$ which is at least $(\log g^{-2})^r$, the bound on the number of steps through which a recent $\mathbf R$-term is carried without contraction. In the third condition $\varepsilon^{*B}(g)$ is the coefficient of the recursion of the boundary terms, a constant times $gp_0(g)$, $N_0^+(g) \coloneqq 1 + 2(\log(32L(L + 1)) + 2\log\log g^{-2})/\log L$ bounds the depth, $C_{\mathrm{ov}} = 2$ counts the generations of created boundary terms at a point, and $W^\sharp$ is a constant of $L$ and $d$ (Supplement, Proposition~S.14). Each condition holds for $\gamma$ small: $\varepsilon^*(g)$ decays as a negative power of $\log g^{-2}$, $\varepsilon^{*B}(g)$ as $g$ times a power of it, while $N'$ and $N_0^+$ grow as double logarithms; the running couplings of the histories which the conditions concern are at most $(1 + \beta_0)\bar g \leq (1 + \beta_0)\gamma$ by \eqref{eq:tube-two-six}. None of the conditions on $\gamma$ involves a loop or a representation, so that the construction, and with it the step factors of the normalisation, is the same for all the ensembles.

Every statement of the paper which requires the coupling to be sufficiently small imposes one of two kinds of conditions (Section~\ref{sec:results}). If it involves no representation, it is an upper bound on $\bar g \leq \gamma'$ in terms of the constants of the construction and of $\Lambda_*$, and, where it concerns a given ensemble, of the data of the ensemble which the statement names. If it involves the representations of an ensemble $\mathcal C$, it concerns the steps $k \leq K'$ at which the loops are carried, and it is an upper bound on the number $\gamma_q$ of \eqref{eq:stopping-coupling} in terms of the constants of the construction, of $\Lambda_*$, and of $c_{\mathcal C}$: these are the conditions of Theorem~\ref{thm:first-step}, of Sections~\ref{sec:contraction} and~\ref{sec:large-field-step}, of Lemma~\ref{lem:root-bound}, Proposition~\ref{prop:torus-consistency}, and Corollary~\ref{cor:remainder-sum}, the condition \eqref{eq:attached-condition}, and the condition \eqref{eq:insertion-recursion-condition} on the coefficient $\varepsilon^Q(g)$ of the old vertices in the recursion of the insertions (Supplement, Proposition~S.16(f)), together with those of Supplement, Propositions~S.9 and~S.10. They are not conditions on the constants of a stage: they are met by the choice of the level $q$ in Theorem~\ref{thm:main}.

\endappendix
\bibliographystyle{amsplain}
\bibliography{refs}
\end{document}
